% !TEX program = xelatex
% 中文译本：Eleuterio F. Toro,
% Riemann Solvers and Numerical Methods for Fluid Dynamics, Third Edition.
\documentclass[11pt,a4paper,openany,fontset=none]{ctexbook}
\usepackage[margin=2.55cm,headheight=15pt]{geometry}
\usepackage{microtype}
\usepackage{fancyhdr}
\usepackage{hyperref}
\usepackage{xurl}
\usepackage{titlesec}
\usepackage{amsmath,amssymb,amsfonts,bm,mathtools,amsthm}
\usepackage[export]{adjustbox}
\usepackage{enumitem}
\usepackage{graphicx}
\usepackage{caption}
\usepackage{longtable,booktabs,array}
\usepackage{needspace}
\usepackage{listings}
\setmainfont[AutoFakeBold=2,AutoFakeSlant=.2]{Arial Unicode MS}
\setCJKmainfont{Songti SC}
\setCJKsansfont{Songti SC}
\setCJKmonofont{Songti SC}
\setlength{\parindent}{2em}
\setlength{\parskip}{0.62em}
\setlength{\emergencystretch}{3em}
\linespread{1.32}\selectfont
\xeCJKsetup{CJKglue=\hskip 0.08em plus 0.04em minus 0.03em}
\setlist[itemize]{leftmargin=2.8em,itemsep=0.32em,topsep=0.48em}
\setlist[enumerate]{leftmargin=3.0em,itemsep=0.32em,topsep=0.48em}
\titleformat{\chapter}[display]{\bfseries\Huge\centering}{第\thechapter 章}{0.75em}{}
\titleformat{\section}{\large\bfseries}{}{0pt}{}
\titleformat{\subsection}{\normalsize\bfseries}{}{0pt}{}
\titlespacing*{\section}{0pt}{1.45em}{0.72em}
\titlespacing*{\subsection}{0pt}{1.20em}{0.58em}
\pagestyle{fancy}
\fancyhf{}
\fancyhead[C]{\small\nouppercase{\leftmark}}
\fancyfoot[C]{\thepage}
\renewcommand{\headrulewidth}{0.3pt}
\hypersetup{hypertexnames=false,colorlinks=true,linkcolor=black,urlcolor=blue,
  pdftitle={黎曼求解器与流体动力学数值方法：实用导论},
  pdfauthor={Eleuterio F. Toro}}
\title{黎曼求解器与流体动力学数值方法\\[0.45em]
  \Large 实用导论（第三版）\\[0.65em]\large 中文译文}
\author{Eleuterio F. Toro 著\\\small AI 辅助翻译与排版}
\date{原著第三版：Springer，2009}
\clubpenalty=10000
\widowpenalty=10000
\displaywidowpenalty=10000
\raggedbottom
\allowdisplaybreaks[1]
\setlength{\abovedisplayskip}{0.82\baselineskip plus 0.2\baselineskip minus 0.1\baselineskip}
\setlength{\belowdisplayskip}{0.82\baselineskip plus 0.2\baselineskip minus 0.1\baselineskip}
\setlength{\abovedisplayshortskip}{0.58\baselineskip}
\setlength{\belowdisplayshortskip}{0.58\baselineskip}
\lstset{basicstyle=\ttfamily\small,breaklines=true,columns=fullflexible,
  frame=single,xleftmargin=0.8em,xrightmargin=0.8em,aboveskip=0.8em,belowskip=0.8em}
\newenvironment{sourceblock}
  {\par\addvspace{0.65\baselineskip}\noindent\begin{minipage}{\linewidth}\centering}
  {\end{minipage}\par\addvspace{0.48\baselineskip}}
% 参数：原PDF页、左下与右上坐标（bp）、最大高度比例。
\newcommand{\sourcecrop}[6]{%
  \includegraphics[page=#1,pagebox=mediabox,viewport={#2bp #3bp #4bp #5bp},clip,
    width=0.94\textwidth,max height=#6\textheight]{Toro.pdf}%
}
\newcommand{\eqnote}[1]{\par\smallskip\noindent\textit{式中：}#1\par}
\theoremstyle{definition}
\newtheorem{example}{例}[chapter]

\makeatletter
\renewcommand{\@pnumwidth}{2.4em}
\renewcommand{\@tocrmarg}{3.2em}
\makeatother

\begin{document}
\frontmatter
\maketitle
\chapter*{译本说明}
\markboth{译本说明}{译本说明}
本译文按正常教材阅读顺序重新排版，不与原 PDF 逐页对齐，也不复制原书带页码目录。
章节另起页，页眉显示当前章名；正文采用较宽松的字号、字距、行距和段间距。

译文按原书逐段覆盖，统一双曲守恒律、黎曼问题、有限体积法、激波捕捉与高分辨率格式的
专业术语。公式使用 LaTeX 重写；程序清单保留代码并配中文说明；只有真实曲线图、网格图
和数值结果图从原 PDF 紧裁保留，题注改为中文且不带入英文正文。译文供学习与检索使用，
涉及算法实现细节和数值常数时建议同时核对原著。

\begingroup\small\tableofcontents\endgroup

\chapter*{常用符号与中文释义}
\addcontentsline{toc}{chapter}{常用符号与中文释义}
\markboth{常用符号与中文释义}{常用符号与中文释义}
\begingroup\small
\renewcommand{\arraystretch}{1.20}
\begin{longtable}{>{\centering\arraybackslash}p{0.23\textwidth}p{0.67\textwidth}}
\toprule 符号 & 中文释义 \\ \midrule
\endfirsthead
\toprule 符号 & 中文释义 \\ \midrule
\endhead
$\bm U$；$\bm F(\bm U)$ & 守恒变量矢量；物理通量矢量 \\
$\bm A=\partial\bm F/\partial\bm U$ & 通量 Jacobian 矩阵 \\
$\lambda_i$；$\bm K^{(i)}$ & 第 $i$ 个特征速度（特征值）；右特征矢量 \\
$\bm W$；$\bm V$ & 特征变量；原始变量矢量 \\
$\bm U_L,\bm U_R$ & 黎曼问题的左、右初始状态 \\
$\bm U_*$；$p_*,u_*$ & 星形区状态；星形区压力与质点速度 \\
$\rho,u,v,w,p,E$ & 密度、速度分量、压力和单位体积总能量 \\
$e,H,a$ & 单位质量内能、总焓和声速 \\
$\gamma$ & 理想气体比热比 \\
$\Delta x,\Delta t$；$S=x/t$ & 网格间距、时间步长；自相似变量 \\
$\bm F_{i+1/2}$ & 单元界面数值通量 \\
$\nu=\Delta t/\Delta x$ & 网格比；与特征速度组合形成 Courant 数 \\
$C_{\mathrm{CFL}}$ & CFL 稳定性参数 \\
$TV(\bm U)$ & 数值解的总变差 \\
$\phi(r)$ & 通量限制器或斜率限制函数，依上下文而定 \\
$S_L,S_R,S_*$ & HLL/HLLC 左波、右波及接触波速度估计 \\
\bottomrule
\end{longtable}
\endgroup

\clearpage
\makeatletter\@mainmattertrue\makeatother
\pagenumbering{arabic}
% ===== 译文分片 1（已合并）=====

% [原PDF p6]
\chapter*{第一版前言}
\addcontentsline{toc}{chapter}{第一版前言}
\markboth{第一版前言}{第一版前言}
1917 年，英国科学家 L. F. Richardson 首次尝试用手工数值求解偏微分方程来预报天气。尽管计算没有成功，这项工作通常仍被视为计算流体力学（CFD）的开端。它已经包含现代 CFD 的四个支柱：需要解决的实际问题、把问题表示成偏微分方程组的数学模型、数值方法，以及执行计算的“计算机”——在 Richardson 的工作中，这台计算机就是人。八十年后，CFD 仍可定义为：计算代表流体流动现象的偏微分方程或积分方程之数值解的科学。

Richardson 之后，计算机性能以空前速度提升，成本则持续下降；同时，需要数值求解的实际问题大幅增加。除气象、海洋、物理与传统工程外，快速变化的制造业也逐步以定量设计取代“制造—试验—修补”的流程。新材料、环境约束下的决策，都不断提出新的数学建模、数值算法和高性能计算需求。连续介质力学的基本方程虽早已建立，热力学、状态方程、非平衡与多相流模型仍须发展。
Richardson 的经历还说明，CFD 从一开始就不是单独的一种算法，而是实际问题、数学模型、离散方法和计算工具的结合。任何一个环节失真，都可能使最终预测失去价值。计算机虽已从人工计算发展到大规模高性能机器，但这四个要素仍未改变；现代 CFD 的通常定义，正是求解描述流动现象的偏微分或积分模型之数值解。

应用需求的扩展不仅来自更快的计算机。传统气象、海洋、物理与工程问题之外，制造业正以全过程定量分析取代反复制造、试验和修补。环境约束下的设计决策以及新材料开发，也要求更可信的模型和算法。连续介质方程在 Richardson 之前一个多世纪已经存在，新的困难更多出现在热力学闭合、真实状态方程、非平衡过程和多相流建模上。因此，计算能力增加并不能代替模型改进与数值分析。

% [原PDF p7]
数值方法是过去八十年中最显著的成功之一。von Neumann 的稳定性研究解释并解决了 Richardson 遇到的困难；Courant、Friedrichs、Richtmyer、Lax、Oleinik、Wendroff、Godunov、Rusanov、van Leer、Harten、Roe、Osher、Colella、Yee 等人的工作，使今天能够使用更准确、高效、稳健且精细的算法处理雄心勃勃的实际问题。CFD 需求与算法复杂度的增长，也对科学家和工程师的教育提出了更高要求，本书正是对流体数值方法训练的一项贡献。

本书材料来自作者多年从事数值方法研究、工程应用和研究生教学的积累，起点是英国 Cranfield 航空学院计算流体力学硕士课程的一门模块；其后又用于英国 Cranfield、德国 Ernst Mach 研究所、日本东北大学冲击波研究中心、意大利 Trento 大学、智利 Federico Santa Maria 技术大学、德国 Aachen 技术大学以及英国 Manchester Metropolitan University 的课程和讲座。

本书讨论在一般多维几何中求解含光滑解和间断解的时变双曲守恒律之现代激波捕捉方法。表述力求全面、实用且主要采用非形式化方式；实现所需的信息尽量完整给出，方法均配有数值算例，并与精确解或可信实验数据比较。全书主体是 Godunov 方法：它以黎曼问题的解为基础。本书详述精确解及多种近似黎曼求解器，也包括随机选择法、通量矢量分裂法以及中心型激波捕捉格式。

% [原PDF p8]
中心格式比 Godunov 方法简单且廉价。本书为标量偏微分方程构造其高阶推广和总变差递减（TVD）版本，再把多数 TVD 方法推广到一维非线性方程组；同时讨论源项和一般几何中的多维系统。非线性系统的主要载体是时变气体动力学 Euler 方程。读完相关章节后，读者应能把这些技术推广到其他双曲系统以及对流—反应—扩散方程。

本书可用于自学。为突出概念清晰性，有时会牺牲部分数学严格性。预期读者包括工程、物理、应用数学或计算机专业的研究生，他们需要从基本原理实现对流—反应—扩散问题的数值算法；也包括刚开始算法研究的数值分析人员，以及工业界和研究机构的科学家、工程师。若系统学习全部内容，并实现、检验大多数方法，大约需要两个月。各章尽量自足且交叉引用充分，也允许从中间或最后若干章开始阅读。

本书亦可作为教学参考。数学或计算机课程可选第 1、2、5、13、15 章以及第 14、16 章的部分内容；工程或物理课程可选择第 3、4、6—12、15—17 章。研究生课程可覆盖大部分内容，但通常假定学生已掌握可压缩流体力学。工业短期培训也可从中选材。

\begin{flushright}
Eleuterio Toro\\
英国 Manchester，1997 年 3 月
\end{flushright}
本书所称“现代激波捕捉”并不要求预先追踪每一条间断曲线，而是在固定或一般网格上用守恒离散自然形成激波。作者对每种方法都尽量给出能够实际编程的细节，并用精确解或可靠实验数据检验结果。这样的安排使算法不只停留在形式推导，而能暴露波速估计、边界处理、源项耦合和多维几何中真正影响计算的环节。

Godunov 思想的核心是在相邻单元状态之间解局部黎曼问题。精确解被详细介绍，不只是为了气体动力学本身，也为了让读者能够把同样的波曲线分析迁移到其他双曲系统。各种近似黎曼求解器则在成本、分辨率和稳健性之间取舍。随机选择、通量矢量分裂和中心型格式提供了不同路线；中心格式结构较简单、价格较低，但其高阶扩展仍须满足非振荡或 TVD 约束。

全书既可连续自学，也可按背景选读。每章尽量自足并给出交叉引用，是有意允许研究人员直接进入与其问题最相关的章节。作者建议，完整自学若同时实现并测试多数算法，大约需要两个月；这一建议也表明，理解数值守恒律不能只阅读公式，还必须通过实现和标准算例检验对波传播、稳定性与守恒性的认识。

% [原PDF p9]
致谢中特别指出，这些访问与课程不仅提供行政支持，也直接塑造了书中材料的选择和讲授次序。作者逐一列出审稿者及所属机构，是为了说明书稿同时接受了气体动力学、数值分析、工程应用和软件实现等多方面检验。排版协助同样重要，因为大量矩阵、分段函数和算法步骤必须以可靠的 LaTeX 形式呈现。
\section*{第一版致谢}
作者感谢多年来安排其访问、研讨会、工作坊和短期课程的同事，这些活动深刻影响了本书。特别感谢德国 Freiburg 的 W. Heilig，意大利 Trento 的 V. Casulli 与 A. Armanini，日本 Sendai 的 K. Takayama，智利 Valparaiso 的 J. Figueroa，以及德国 Aachen 的 J. Ballmann。书稿最后排版阶段在 Manchester Metropolitan University 完成，作者感谢该校支持。

作者亦感谢历届博士生对课程与书稿的意见，以及他们对相关研究的贡献，特别包括 Stephen Billett、Caroline Lowe、Nikolaos Nikiforakis、Ed Boden、Wei Hu、Mathew Ivings 和 Richard Millington。感谢儿子 Marcelo 协助排版，David Ingram 与 John Nuttall 帮助解决 LaTeX 技术问题；感谢 Springer-Verlag 及 Heidelberg 的 Erdmuthe Raufelder 女士提供专业协助。Stephen Billett、John Clarke、Jean-Marc Moschetta、Claus-Dieter Munz、Jack Pike、Ning Qin、Tsutomu Saito、Charles Saurel、Trevor Randall、Peter Sweby、Marcelo Toro、Helen Yee 与 Clive Woodley 对书稿提供了有益反馈。Brigitte 长期不变的鼓励，是写作过程中极其珍贵的支持。

% [原PDF p10]
\chapter*{第三版前言}
\addcontentsline{toc}{chapter}{第三版前言}
\markboth{第三版前言}{第三版前言}
第一版出版十余年后，本书得到世界各地学生、教师和资深研究人员的广泛回应。这十年间，数值方法不仅在物理和工业过程中更普遍，也扩展到生物、经济、社会科学和跨学科研究。数学建模与数值模拟出现了新的趋势：数值方法正从工程设计工具转变为科学研究中理解各种现象的不可或缺的仪器。人们期望模拟真正代表原数学模型的解，这要求持续研究更高质量、更准确的算法，也要求大学加强本科与研究生训练。

第三版新增第 18—20 章。第 18 章把第 7 章一维系统的中心 FORCE 通量推广到多维。中心格式不按传统方式直接求解经典黎曼问题，因而对难以甚至无法得到黎曼解的复杂系统具有更广适用性。新的 FORCE 格式可用于一般结构与非结构网格上的二维、三维问题，并能在有限体积和间断 Galerkin 有限元框架中推广到时空高阶。

% [原PDF p11]
新章节的共同目标，是把经典二阶、分步式 Godunov 技术推进到时空统一的高阶单步方法。广义黎曼问题不只给出界面瞬时通量，还提供时间导数信息；在有限体积中可直接更新单元平均，在间断 Galerkin 中则更新多项式自由度。HLLC 修订保留接触波分辨能力，同时扩展稳健波速估计和复杂系统适用性。
第 19 章讨论高阶或广义黎曼问题：双曲平衡律的初值由原点间断所分隔的两个光滑向量构成。其解可作为构造时空任意阶、全离散、单步 Godunov 型格式的基本模块，是前面经典黎曼问题内容的推广。第 20 章则在有限体积框架中介绍这些高阶方法。

第三版还实质性修订了第 10 章 HLLC 黎曼求解器，以回应读者来信、新进展和更具挑战的科技应用；第 21 章亦有修改，并改正读者指出的错误。本书虽定位为入门书，却主要由研究人员使用，对本科教学可能过深。作者与 Enrico Bertolazzi、Gianluca Vignoli 正合作编写更基础、含大量练习、案例和小项目的教材；同时与 Claus-Dieter Munz、Vladimir Titarev、Michael Dumbser 合作准备面向研究的高阶有限体积与间断 Galerkin 专著。

作者感谢 Vladimir Titarev、Cristóbal Castro、Martin Käser、Michael Dumbser、Enrico Bertolazzi，以及访问学者 María Nofuentes 和 Arturo Hidalgo 对第三版材料的帮助。

% [原PDF p12]
\enlargethispage{3\baselineskip}
作者特别回顾第一版在英国 Cranfield University（原 Cranfield Institute of Technology）的酝酿阶段。20 世纪 80 年代末至 90 年代初，航空科学学院空气动力学系汇聚了背景多样、充满活力的科研群体。那段经历深刻影响本书的内容与范围。作者感谢 John Clarke、Phil Roe、Jack Pike、Smadar Karni、James Quirk、Nikos Nikiforakis、Stephen Billett、Alan Dowes、Caroline Lowe、William Speares、Michael Spruce、Ed Boden、Mauricio Cáceres 以及众多同事和学生，也感谢 Springer-Verlag 长期积极的鼓励和有效支持。

\begin{flushright}
Eleuterio Toro\\
意大利 Trento 大学，2008 年 12 月
\end{flushright}

% [原PDF p13--p22]
% 原书此处为带原始页码的英文目录；按用户要求不译为正文，最终中文目录由 LaTeX 自动生成。

% [原PDF p23]
\chapter{流体动力学方程}
本章给出气体等可压缩介质的动力学控制方程，并以状态方程作为闭合条件。状态方程描述材料性质，需要若干热力学概念。这里不企图穷尽或严格推导连续介质力学，而是以自足、可直接服务于黎曼求解器与数值方法的方式概述控制方程和热力学。

表述次序有意不同于传统教材：先介绍 Euler 方程微分形式、基本物理量以及导向状态方程的热力学关系，再加入黏性扩散和传热，最后给出通常作为起点的积分形式。本章还建立模型层次，包括时变 Euler 方程、定常超声速 Euler 方程、浅水方程和不可压 Navier--Stokes 的人工可压缩方程；这些双曲守恒律都可直接应用黎曼求解器和迎风方法。

采用笛卡尔坐标 $(x,y,z)$ 和时间 $t$。任意场写作 $\phi(x,y,z,t)$，偏导简记为 $\phi_t,\phi_x,\phi_y,\phi_z$。其他坐标系可按链式法则变换。向量点积、梯度和散度的定义如下页给出。
本章的目的，是为全书后续数值分析建立共同的物理语言。连续介质的质量、动量和总能量分别构成守恒变量；压力、温度和内能则通过材料的状态方程相联系。作者特别强调，这里给出的背景已足以支持后续各章，但它并不是连续介质力学的公理化论述。需要更深入理解应力、热传导或一般材料响应的读者，仍应查阅流体力学与热力学专著。

本章模型层次还具有方法论意义：复杂模型在不同假设下可约化成较简单的双曲系统。时变 Euler 方程描述无黏可压缩流；将一个空间方向作为“时间样”方向，可得到定常超声速推进问题；对自由表面流采用静水近似得到浅水方程；为不可压 Navier--Stokes 引入伪时间，则得到人工可压缩系统。它们的物理含义各异，却都具有有限传播速度和特征波结构，因此可在统一的黎曼框架中处理。

全文采用笛卡尔直角坐标。梯度、散度以及向量内积均按通常定义使用；当坐标或网格发生变换时，守恒形式必须连同 Jacobian 与度量项一起变换。书中先从微分方程出发，是为了尽快显露特征结构；随后回到积分形式，则是因为间断解与有限体积离散最终都依赖控制体收支。

% [原PDF p24]
本页原文还提醒，点积公式排印中每一分量必须对应相乘；梯度只作用于标量而给出向量，散度只作用于向量而给出标量。原始变量便于施加物理边界和解释状态，守恒变量则使控制体收支直接闭合。由守恒量恢复原始量时，先除以密度得到速度，再从总能量扣除动能求内能和压力。
对 $\boldsymbol A=(a_1,a_2,a_3)$、$\boldsymbol B=(b_1,b_2,b_3)$，
\begin{equation}
\boldsymbol A\cdot\boldsymbol B=a_1b_1+a_2b_2+a_3b_3,
\end{equation}
标量梯度与向量散度为
\begin{equation}
\nabla\phi=(\phi_x,\phi_y,\phi_z),\qquad
\nabla\cdot\boldsymbol A=(a_1)_x+(a_2)_y+(a_3)_z.
\end{equation}

\section{Euler 方程}
时变 Euler 方程是非线性双曲守恒律，描述忽略体积力、黏性应力和热流时的可压缩介质。原始变量为密度 $\rho$、压力 $p$ 和速度分量 $u,v,w$，速度向量 $\boldsymbol V=(u,v,w)$。守恒变量为质量密度、三个动量分量 $\rho u,\rho v,\rho w$ 以及单位体积总能量 $E$。它们分别来自质量守恒、Newton 第二定律和能量守恒。守恒变量形式是守恒型数值方法的基础。以下先假定解足够光滑，后面再允许激波等间断。

% [原PDF p25]
守恒向量与三个方向通量的完整矩阵表示还揭示了压力只加在法向动量通量上：在 $x$ 通量中是 $\rho u^2+p$，而 $y,z$ 动量只是 $\rho uv,\rho uw$。总能量 $E$ 是单位体积量，比内能 $e$ 则是单位质量量，二者不可混用。

三个通量向量的结构完全对称：质量通量是对应方向动量；法向动量通量多一个压力；两个切向动量仅随对流输运；能量通量为速度乘总焓密度。这个旋转不变结构意味着在任意控制面上，只需把速度分成法向和切向，即可复用一维黎曼求解器。
\subsection{守恒律形式}
五个控制方程为
\begin{align}
\rho_t+(\rho u)_x+(\rho v)_y+(\rho w)_z&=0,\tag{1.1}\\
(\rho u)_t+(\rho u^2+p)_x+(\rho uv)_y+(\rho uw)_z&=0,\tag{1.2}\\
(\rho v)_t+(\rho uv)_x+(\rho v^2+p)_y+(\rho vw)_z&=0,\tag{1.3}\\
(\rho w)_t+(\rho uw)_x+(\rho vw)_y+(\rho w^2+p)_z&=0,\tag{1.4}\\
E_t+[u(E+p)]_x+[v(E+p)]_y+[w(E+p)]_z&=0.\tag{1.5}
\end{align}
总能量密度
\begin{equation}
E=\rho\left(\frac12\boldsymbol V^2+e\right),\qquad
\frac12\boldsymbol V^2=\frac12(u^2+v^2+w^2),
\tag{1.6}
\end{equation}
其中 $e$ 为比内能。严格说只有三个动量方程叫 Euler 方程，但通常把式（1.1）—（1.5）统称为 Euler 方程。令守恒向量 $\boldsymbol U=(\rho,\rho u,\rho v,\rho w,E)^T$，并定义三个方向的通量 $\boldsymbol F,\boldsymbol G,\boldsymbol H$，即可写成
\begin{equation}
\boldsymbol U_t+\boldsymbol F(\boldsymbol U)_x+\boldsymbol G(\boldsymbol U)_y+\boldsymbol H(\boldsymbol U)_z=0.
\tag{1.7}
\end{equation}

% [原PDF p26]
非守恒形式的推导要反复使用连续方程消去速度乘质量残差。例如展开 $(\rho u)_t$ 与 $(\rho u^2)_x$ 后，含 $u[\rho_t+\nabla\cdot(\rho V)]$ 的项为零，留下 $\rho Du/Dt+p_x=0$。这一代数在光滑区严格成立；跨间断若分别离散展开项，则消去关系不再保证，可能给出错误激波速度。
\subsection{其他紧凑形式}
把通量对守恒变量求 Jacobian，$\boldsymbol A=\partial\boldsymbol F/\partial\boldsymbol U$、$\boldsymbol B=\partial\boldsymbol G/\partial\boldsymbol U$、$\boldsymbol C=\partial\boldsymbol H/\partial\boldsymbol U$，光滑解可写成准线性形式
\begin{equation}
\boldsymbol U_t+\boldsymbol A(\boldsymbol U)\boldsymbol U_x
+\boldsymbol B(\boldsymbol U)\boldsymbol U_y
+\boldsymbol C(\boldsymbol U)\boldsymbol U_z=0.\tag{1.12}
\end{equation}
但在间断处，通量导数和链式法则均无通常意义，所以数值守恒律应以散度形式为基础。用随体导数
\begin{equation}
\frac{D}{Dt}=\frac{\partial}{\partial t}+u\frac{\partial}{\partial x}
+v\frac{\partial}{\partial y}+w\frac{\partial}{\partial z},
\end{equation}
还可得到质量、速度和内能的非守恒形式。质量式为 $D\rho/Dt+\rho\nabla\cdot\boldsymbol V=0$，动量式为 $\rho D\boldsymbol V/Dt+\nabla p=0$，内能式为 $\rho De/Dt+p\nabla\cdot\boldsymbol V=0$。这些形式便于物理解释，却不能直接跨越激波使用。

% [原PDF p27]
压力、能量和温度的单位必须与状态方程常数一致。若代码使用无量纲变量，参考密度、速度、长度和温度会诱导参考压力与能量；所有 EOS 参数也必须按同一尺度无量纲化。热力学平衡假设并不要求流场全局均匀，只要求分子松弛时间远小于所解析流动时间。
\section{热力学考虑}
Euler 方程有五个方程，却包含 $\rho,u,v,w,E,p,e$ 等量，必须用材料状态关系闭合。热力学把状态变量分为广延量和强度量；本书主要使用单位质量量。SI 基本单位为米、千克、秒和 Kelvin。密度单位 $\mathrm{kg,m^{-3}}$，压力单位 Pascal 即 $\mathrm{N,m^{-2}}$，能量单位 Joule。量纲一致性不仅是物理要求，也是检查数值代码和状态方程实现错误的有效手段。

热力学平衡态可由两个独立变量确定。常用基本关系是比内能 $e=e(\rho,p)$ 或压力 $p=p(\rho,e)$。速度不属于热力学状态，只影响总能量中的动能部分。局部热力学平衡假设允许每个流体微团使用平衡状态方程，即使整个流场显著非均匀。
国际单位制中，长度、质量和时间的基本单位分别为 $mathrm{m}$、$mathrm{kg}$ 和 $mathrm{s}$。速度的量纲为 $mathrm{m,s^{-1}}$，加速度为 $mathrm{m,s^{-2}}$。一牛顿是使 $1,mathrm{kg}$ 质量产生 $1,mathrm{m,s^{-2}}$ 加速度所需的力，即
\[
1\,\mathrm N=1\,\mathrm{kg,m,s^{-2}}.
\]
压力是单位面积上的力，其单位 Pascal 满足 $1\,\mathrm{Pa}=1\,\mathrm{N,m^{-2}}$；常用换算为 $1\,\mathrm{bar}=10^5\,\mathrm{Pa}$、$1\,\mathrm{atm}=101325\,\mathrm{Pa}$。例如动量通量中的 $\rho u^2$ 与压力 $p$ 必须具有相同量纲，这种量纲一致性应在每一个物理等式中保持。

力沿其方向推动作用点移动一段距离所作的功等于力与位移之积，功和能量的单位均为 Joule，$1\,\mathrm J=1\,\mathrm{N,m}$。温度使用热力学绝对温标，单位为 Kelvin。与系统质量成比例的量称为广延性质，如总能量和总体积；不随系统质量改变的量称为强度性质，如温度和压力。广延量除以质量便得到对应的比量，例如比容 $v=V/m=1/\rho$。这里的 $v$ 与速度的 $y$ 分量使用同一字母，须由上下文区分。

% [原PDF p28]
状态曲面可在 $(p,\rho,e,T,s)$ 中任选两个独立坐标表示，但数值反演并非总是显式。表格式或多相 EOS 往往需要迭代求温度和压力；迭代必须返回热力学一致的声速导数。总能量守恒并不自动保证内能为正，高 Mach 数下动能远大于内能，减法误差尤其敏感。
\subsection{状态方程}
状态方程（EOS）概括材料的热性质。一般可写成
\begin{equation}
p=p(\rho,e),\qquad e=e(\rho,p),
\end{equation}
或用比容 $\tau=1/\rho$、温度 $T$、熵 $s$、比焓 $h=e+p/\rho$ 等变量表达。对数值流体力学，给定守恒量后必须先由
\begin{equation}
e=\frac{E}{\rho}-\frac12(u^2+v^2+w^2)
\end{equation}
恢复内能，再由 EOS 求压力。若恢复的密度或内能超出状态方程允许域，数值状态就没有物理意义，不能靠任意代数操作掩盖。

完全状态方程还可由热状态方程 $p=p(\rho,T)$ 和量热状态方程 $e=e(\rho,T)$ 组成。不同材料可能具有相同形式的控制守恒律，却因 EOS 不同而表现出完全不同的波速、压缩性和激波关系。
对 $p$--$v$--$T$ 系统，热平衡态的族可看作 $p$--$v$ 平面上的曲线，每条曲线由一个温度值标记。热状态方程可以写为
\begin{equation}
T=T(p,v),\qquad p=p(T,v),\qquad v=v(T,p).
\tag{1.14}
\end{equation}
具体函数关系随物质而变；热完全理想气体的关系最简单，为 $T=pv/R$。

热力学第一定律表明，在一个非绝热过程中，内能增量等于外界对系统所作的功与传入系统的热量之和。采用 $dW=-p\,dv$ 的符号约定，可写成
\begin{equation}
dQ=de+p\,dv.
\tag{1.16}
\end{equation}
内能还须由量热状态方程 $e=e(p,v)$ 给出，也可等价地表示为 $p=p(v,e)$ 或 $v=v(e,p)$。对量热完全理想气体，
\begin{equation}
e=\frac{pv}{\gamma-1}=\frac{p}{\rho(\gamma-1)}.
\tag{1.18}
\end{equation}
热状态方程与量热状态方程彼此约束，却不能互相完全替代。Euler 方程若只需从 $(\rho,e)$ 恢复压力，量热 EOS 已足够；若还要求温度，例如计算黏性或热传导系数，则必须显式提供热 EOS。

% [原PDF p29]
Gibbs 关系还给出 $dh=Tds+\tau dp$。等熵时 $dh=dp/\rho$，这在推导稀疏波黎曼不变量时直接使用。声速是小等熵扰动的传播速度，而不是等温导数；若错误采用 $(\partial p/\partial\rho)_T$，会同时改变特征值、时间步限制和激波关系。
\subsection{其他变量与关系}
比焓和总比焓定义为
\begin{equation}
h=e+\frac{p}{\rho},\qquad H=h+\frac12\boldsymbol V^2=\frac{E+p}{\rho}.
\end{equation}
能量通量因此可写成 $\rho uH,\rho vH,\rho wH$。Gibbs 关系为
\begin{equation}
T,ds=de+p,d\tau=de-\frac{p}{\rho^2},d\rho.
\tag{1.31}
\end{equation}
可逆绝热过程满足 $ds=0$，因而沿等熵曲线 $de=p,d\rho/\rho^2$。定容比热和定压比热分别为 $c_v=(\partial e/\partial T)_\rho$ 与 $c_p=(\partial h/\partial T)_p$；其比值常记为 $\gamma=c_p/c_v$。

热膨胀、等温压缩和等熵压缩系数描述状态对温度、压力的响应。声速是等熵压力对密度的导数，
\begin{equation}
c^2=\left(\frac{\partial p}{\partial\rho}\right)_s.
\tag{1.34}
\end{equation}
为使 $de+p\,dv$ 成为恰当微分，引入积分因子 $1/T$；热力学第二定律据此定义熵 $s$：
\begin{equation}
T\,ds=de+p\,dv.
\tag{1.19}
\end{equation}
任一过程中，熵的变化可分为从边界输入的部分和系统内部产生的部分，即 $\Delta s=\Delta s_0+\Delta s_i$。热传导、黏性以及激波内部结构都会产生熵；不可逆过程满足 $\Delta s_i>0$，只有可逆过程才有 $\Delta s_i=0$。

比焓定义为
\begin{equation}
h=e+pv.
\tag{1.20}
\end{equation}
由 Gibbs 关系得 $de=T\,ds-p\,dv$，故可把内能视为 $e=e(s,v)$。对焓求微分则有
\begin{equation}
dh=T\,ds+v\,dp,
\tag{1.23}
\end{equation}
从而可取 $h=h(s,p)$。这两个关系称为正则状态方程；与单独的热 EOS 或量热 EOS 不同，每一个正则关系都包含完整的热力学信息。例如比较 $de=e_s\,ds+e_v\,dv$ 可得
\begin{equation}
p=-\left(\frac{\partial e}{\partial v}\right)_s,
\qquad T=\left(\frac{\partial e}{\partial s}\right)_v.
\tag{1.25}
\end{equation}
下标指出求偏导时保持不变的独立变量，在变量频繁转换的热力学推导中不可省略。

% [原PDF p30]
若以 $p=p(\rho,e)$ 表示 EOS，令 $p_\rho=(\partial p/\partial\rho)_e$、$p_e=(\partial p/\partial e)_\rho$，结合 Gibbs 关系可得
\begin{equation}
c^2=p_\rho+\frac{p}{\rho^2}p_e.
\tag{1.36}
\end{equation}
该式适用于一般材料，是构造 Euler 方程特征速度和黎曼求解器的重要输入。基本导数必须在物理状态域内满足 $c^2>0$，否则方程失去通常的双曲性。温度可由熵表示或由量热状态方程反演；热力学相容性要求不同状态函数的混合偏导满足 Maxwell 关系。

流体微团的第一定律可写成内能变化等于热输入减去对外做功。对无热传导、无黏性且光滑的 Euler 流，随体熵保持不变；激波中不可逆性使熵增加，因此“等熵”只适用于光滑区域，不能跨越激波直接使用。
同理，由 $h=h(s,p)$ 和 $dh=T\,ds+v\,dp$ 得
\begin{equation}
T=\left(\frac{\partial h}{\partial s}\right)_p,
\qquad v=\left(\frac{\partial h}{\partial p}\right)_s.
\tag{1.26}
\end{equation}
Helmholtz 自由能定义为 $f=e-Ts$。若采用正则形式 $f=f(v,T)$，则
\begin{equation}
s=-\left(\frac{\partial f}{\partial T}\right)_v,
\qquad p=-\left(\frac{\partial f}{\partial v}\right)_T.
\tag{1.28}
\end{equation}

给定热 EOS $v=v(p,T)$ 后，可定义体膨胀系数和等温压缩系数
\begin{equation}
\alpha=\frac1v\left(\frac{\partial v}{\partial T}\right)_p,
\qquad
\beta=-\frac1v\left(\frac{\partial v}{\partial p}\right)_T.
\tag{1.29}
\end{equation}
由 Maxwell 关系可得 $(\partial s/\partial v)_T=(\partial p/\partial T)_v=\alpha/\beta$，进而
\begin{equation}
\left(\frac{\partial e}{\partial v}\right)_T
=\frac{\alpha T-\beta p}{\beta}.
\tag{1.30}
\end{equation}
这些恒等式把可测的热膨胀与压缩性质同量热状态函数联系起来，也是检验经验 EOS 热力学相容性的基本工具。

% [原PDF p31]
一般 EOS 的压力微分为 $dp=p_\rho,d\rho+p_e,de$。沿流体轨迹利用质量和内能方程，可推得压力方程
\begin{equation}
\frac{Dp}{Dt}+\rho c^2\nabla\cdot\boldsymbol V=0.
\tag{1.42}
\end{equation}
它与速度方程组合，揭示压缩扰动以声速传播。熵方程为 $Ds/Dt=0$，表明光滑 Euler 流中每个流体粒子的熵守恒，但空间上不同粒子可有不同熵。总焓在一般非定常流中不必沿粒子保持常数；定常绝热无功流动则有 Bernoulli 型关系。

基本热力学导数还决定基本导数或非线性参数，用来判断压缩波是否陡化为激波。对常见正常流体，声速随等熵压缩增大，特征场为真正非线性；复杂 EOS 可能出现非经典波，因此近似黎曼求解器不应暗中假设理想气体关系。
若传入热量 $dQ$ 使温度改变 $dT$，则 $c=dQ/dT$ 称为热容量。恒压过程中 $dQ=de+d(pv)=dh$，因此
\begin{equation}
c_p=\left(\frac{\partial h}{\partial T}\right)_p
=T\left(\frac{\partial s}{\partial T}\right)_p.
\tag{1.31}
\end{equation}
恒容热容量则为
\begin{equation}
c_v=\left(\frac{\partial e}{\partial T}\right)_v
=T\left(\frac{\partial s}{\partial T}\right)_v.
\tag{1.32}
\end{equation}

声速是本书最重要的状态量之一。对处于无约束热力学平衡的微小扰动，平衡声速定义为
\begin{equation}
a=\sqrt{\left(\frac{\partial p}{\partial\rho}\right)_s}.
\tag{1.34}
\end{equation}
它是等熵而非等温导数，因为声学扰动的时间尺度通常短到来不及与外界交换热量。若量热 EOS 写成 $h=h(p,\rho)$，把 $dh=T\,ds+dp/\rho$ 与 $h$ 的全微分比较，并令 $ds=0$，可把声速表示为 $h$ 的两个偏导之比。对热完全理想气体进一步得到
\begin{equation}
a=\sqrt{\gamma(T)RT}=\sqrt{\frac{\gamma p}{\rho}}.
\tag{1.35}
\end{equation}
在一般 EOS 中，声速必须从一致的热力学导数计算，不能用经验常数替代。

% [原PDF p32]
把原始变量取为 $\boldsymbol W=(\rho,u,v,w,p)^T$，Euler 方程可写成原始变量准线性系统。沿任意单位法向 $\boldsymbol n$，其特征速度为
\begin{equation}
V_n-c,\quad V_n,\quad V_n,\quad V_n,\quad V_n+c,
\end{equation}
其中 $V_n=\boldsymbol V\cdot\boldsymbol n$。两个声学特征对应压力与法向速度扰动，中间三个速度对应熵波和两个切向速度波。这个结构将在第 3、4 章系统推导，并贯穿精确与近似黎曼求解器。

热力学凸性保证声速实数及正常波结构。数值计算通常以守恒变量推进，却在原始或特征变量中重构；变量转换必须使用同一 EOS 及其导数，否则通量、波速与热力学闭合会互相矛盾。
若压力直接写成 $p=p(\rho,e)$，则声速可表示为
\begin{equation}
a^2=p_\rho+\frac{p}{\rho^2}p_e;
\end{equation}
若反过来使用 $e=e(\rho,p)$，则有
\begin{equation}
a^2=\frac{p/\rho^2-e_\rho}{e_p}.
\tag{1.36}
\end{equation}
这里每个下标都表示另一独立变量保持不变。两式在热力学一致时完全等价，并说明 EOS 接口除压力外还须可靠地返回偏导。

理想气体的宏观热状态方程为
\begin{equation}
pV=n\mathcal RT,
\tag{1.37}
\end{equation}
其中 $n$ 为 kilomole 数，$\mathcal R$ 为普适气体常数。一个 kilomole 含有 $N_A$ 个粒子，$N_A$ 为 Avogadro 常数；以相对分子质量 $\omega$ 除去物质的量后，得到单位质量形式
\begin{equation}
pv=RT,\qquad R=\frac{\mathcal R}{\omega}.
\tag{1.38}
\end{equation}
因此 $v=RT/p$，并有 $\alpha=1/T$、$\beta=1/p$。代入前述热力学恒等式即知 $(\partial e/\partial v)_T=0$：满足理想热 EOS 的气体，其内能只依赖温度。

% [原PDF p33]
理想气体的“热完全”只要求 $p=\rho RT$，允许 $c_v(T)$ 变化；“量热完全”进一步要求比热常数。后者才给出常数 $\gamma$ 和线性关系 $p=(\gamma-1)\rho e$。高温空气中振动自由度激活后 $\gamma$ 下降，若再发生解离，组分方程与化学能也必须进入模型。
\subsection{理想气体}
热完全理想气体满足
\begin{equation}
p=\rho RT,qquad e=e(T),
\tag{1.45}
\end{equation}
其中 $R$ 为比气体常数。若比热随温度不变，则为量热完全气体，
\begin{equation}
e=c_vT,\quad h=c_pT,\quad R=c_p-c_v,\quad
p=(\gamma-1)\rho e.
\tag{1.48}
\end{equation}
因此
\begin{equation}
p=(\gamma-1)\left[E-\frac{\rho}{2}(u^2+v^2+w^2)\right],
\end{equation}
声速 $c=\sqrt{\gamma p/\rho}$。空气在中等温度范围常取 $\gamma=1.4$；单原子气体取 $5/3$。温度范围过宽或存在振动、解离、电离时，常比热假设不再成立。
若进一步取
\begin{equation}
e=c_vT
\tag{1.41}
\end{equation}
且 $c_v$ 为常数，便得到量热完全或多方气体。一般热力学关系给出
\begin{equation}
c_p=c_v+\frac{\alpha^2Tv}{\beta};
\tag{1.42}
\end{equation}
对理想气体化为 Mayer 关系 $c_p-c_v=R$。热稳定性要求 $c_v>0$，机械稳定性要求 $\beta>0$，于是 $c_p>c_v>0$。定义比热比
\begin{equation}
\gamma=\frac{c_p}{c_v},
\tag{1.45}
\end{equation}
便有 $c_p=\gamma R/(\gamma-1)$、$c_v=R/(\gamma-1)$。量热完全气体的 $\gamma$ 为常数，热完全气体中它通常是温度的函数。

分子能量均分原理可解释这些关系。若一个分子有 $M$ 个有效自由度，每个二次自由度的平均能量为 $kT/2$，则单位质量内能为 $e=MRT/2$，从而 $c_v=MR/2$、$c_p=(M+2)R/2$，并有
\begin{equation}
\gamma=\frac{M+2}{M}.
\tag{1.47}
\end{equation}
单原子气体只有三个平动自由度，理论值十分准确；多原子气体还具有转动和振动自由度，其激活程度随温度改变，常比热近似因而可能偏离实验。

% [原PDF p34]
等熵关系只连接同一光滑熵水平上的状态，不能替代激波能量方程。激波两侧虽各自可视为平衡理想气体，却具有不同熵常数。数值中常用熵 $s\propto\ln(p/\rho^\gamma)$ 检查物理解：稀疏波保持熵，接触间断允许熵跳跃，压缩激波使熵增加。
量热完全气体的等熵关系为
\begin{equation}
p\rho^{-\gamma}=\text{常数},\qquad
T\rho^{1-\gamma}=\text{常数},
\tag{1.55}
\end{equation}
熵差可写成
\begin{equation}
s-s_0=c_v\ln\frac{p}{p_0}-c_p\ln\frac{\rho}{\rho_0}
=c_v\ln\left(\frac{p/p_0}{(\rho/\rho_0)^\gamma}\right).
\end{equation}
总焓为 $H=c_pT+\boldsymbol V^2/2$。这些闭式关系使理想气体 Euler 黎曼问题能够高效精确求解，后续章节多数公式据此建立。

比热比决定波的非线性和激波压缩极限。虽然程序常把 $\gamma$ 作为常数，实际混合气体的 $R,c_p,c_v$ 可随组分和温度改变；使用理想气体近似时，应明确适用范围，并保持能量零点的一致定义。
由 $M=2/(\gamma-1)$ 和 $e=Mpv/2$，可再次得到
\begin{equation}
e=\frac{pv}{\gamma-1}=\frac{p}{(\gamma-1)\rho}.
\tag{1.48}
\end{equation}
分子理论预测 $1<\gamma\le 5/3$；上界对应只有三个平动自由度的单原子气体，随着更多自由度参与储能，$M$ 增大而 $\gamma$ 向 $1$ 接近。多原子气体的实验比热随温度显著变化，因此在宽温区计算中应对 $c_p(T)$ 或组分热力学进行积分，而非在总能量关系中固定一个经验 $\gamma$。

理想气体的简洁关系并不意味着所有过程都等熵。可逆绝热光滑区满足 $p/\rho^\gamma$ 沿粒子轨迹不变；热传导、黏性或激波内部的不可逆耗散会产生熵。精确黎曼解中稀疏波沿等熵曲线积分，而激波必须使用质量、动量、能量的跳跃条件，这一区分贯穿全书。

% [原PDF p35]
\subsection{余容积气体与 van der Waals 气体}
高密度时分子有限体积不可忽略。余容积（covolume）EOS 为
\begin{equation}
p=\frac{(\gamma-1)\rho e}{1-b\rho},
\tag{1.60}
\end{equation}
$b$ 是余容积常数，物理域要求 $1-b\rho>0$。相应声速比理想气体大，并随密度接近 $1/b$ 急剧增加。该 EOS 保留理想气体形式的若干便利，常用于高压气体与爆轰模型。

van der Waals EOS 同时加入排斥体积和分子吸引，典型写法
\begin{equation}
\left(p+a\rho^2\right)\left(\frac1\rho-b\right)=RT.
\end{equation}
常数 $a$ 表示吸引效应，$b$ 表示排斥体积。它能表现临界现象，但在两相区需要 Maxwell 构造或更完整的相平衡模型。
余容积 EOS 更常见的热形式为
\begin{equation}
p(v-b)=RT,
\tag{1.49}
\end{equation}
其中 $b$ 的 SI 量纲为 $\mathrm{m^3\,kg^{-1}}$。它修正分子自身所占体积，使可供分子平动的有效比容从 $v$ 减为 $v-b$。该式也称 Noble--Abel EOS，适合描述推进系统中高密度燃烧产物。最简单模型取常数 $b$；实验表明许多固体推进剂的 $b$ 位于约 $0.9\times10^{-3}$ 至 $1.1\times10^{-3}\,\mathrm{m^3\,kg^{-1}}$，最佳拟合的误差可控制在百分之几以内。更精细模型允许 $b=b(\rho)$，并以表格或代数式给出。

常 $b$ 的量热关系与声速为
\begin{equation}
e=\frac{p(1-b\rho)}{\rho(\gamma-1)},
\qquad
a=\left[\frac{\gamma p}{(1-b\rho)\rho}\right]^{1/2}.
\tag{1.50}
\end{equation}
再考虑分子间吸引力，压力需减去 $c/v^2$，得到
\begin{equation}
\left(p+\frac{c}{v^2}\right)(v-b)=RT,
\tag{1.51}
\end{equation}
即真实气体的 van der Waals 方程。参数 $c$ 依气体而定；排斥与吸引两种修正共同造成临界与非凸现象。

% [原PDF p36]
van der Waals 等非凸 EOS 可产生基本导数变号、稀疏激波或复合波。此时以理想气体经验固定波型的求解器可能选错解。通用黎曼算法必须从实际 Hugoniot 曲线、等熵积分曲线和熵准则判断连接，同时避免进入机械不稳定区。
对 van der Waals 材料，热状态方程与量热关系必须一起给出，才能由 $(\rho,e)$ 唯一恢复 $p,T$。临界点由等温线上 $(\partial p/\partial\tau)_T=0$、$(\partial^2p/\partial\tau^2)_T=0$ 决定。低于临界温度的非凸区会导致经典双曲波结构发生变化；简单理想气体黎曼公式不能直接套用。

本书后续主要以量热完全理想气体说明方法，但强调算法结构可用于一般 EOS，只要能稳定计算压力、声速和所需热力学导数。实际代码应把 EOS 作为独立模块，而不是把 $p=(\gamma-1)\rho e$ 散布在通量、边界和黎曼求解器中。
常数余容积模型的合理性来自一个实用观察：在所关心的高压范围内，拟合得到的 $b$ 变化通常较小。若密度跨度足以使这一变化不可忽略，则应使用 $b(\rho)$，并在计算声速时把其导数一并计入，不能只把常数公式中的 $b$ 替换为函数值。

van der Waals 修正说明，理想气体与余容积气体都忽略了分子间吸引。吸引力使给定 $T,v$ 下的实际压力降低 $c/v^2$。这种 EOS 在临界区或两相区的数学性质与量热完全理想气体显著不同；后续以理想气体演示的波曲线公式并不能不加检查地外推。关于一般 EOS、燃烧材料和真实气体热力学，原书在本节末给出了进一步参考文献。

至此，Euler 动力学方程和最基本的状态关系已经齐备。下面加入黏性应力与热传导，随后再从物理守恒原理推导对间断仍有意义的积分形式。

% [原PDF p37]
流体总应力写成 $\boldsymbol S=-p\boldsymbol I+\boldsymbol\Pi$：第一项是球对称的热力学压力，第二项才是黏性应力。黏性本构和状态方程一样，都是对材料性质的近似描述；前者把应力与速度梯度联系起来，后者把压力、密度和内能闭合起来。

正应力含体积膨胀贡献，剪应力由速度对称梯度产生；刚体旋转的反对称梯度不产生黏性应力。Stokes 假设并非所有材料严格成立，但对许多气体是常用近似。若黏度依温度变化，散度展开时还会出现黏度梯度项，不能简单使用常系数 Laplacian。
\section{黏性应力}
Euler 模型忽略由速度梯度引起的内摩擦。Newton 流体的应力张量分为各向同性压力和黏性偏应力，
\begin{equation}
\boldsymbol\sigma=-p\boldsymbol I+\boldsymbol\tau.
\end{equation}
对各向同性 Newton 流体，
\begin{equation}
\tau_{ij}=\mu\left(\frac{\partial V_i}{\partial x_j}
+\frac{\partial V_j}{\partial x_i}\right)
+\lambda(\nabla\cdot\boldsymbol V)\delta_{ij},
\tag{1.68}
\end{equation}
$\mu$ 为剪切黏度，$\lambda$ 为第二黏度。Stokes 假设给出 $\lambda=-2\mu/3$。黏性应力在动量方程中以 $\nabla\cdot\boldsymbol\tau$ 出现，并通过 $\boldsymbol\tau\boldsymbol V$ 对能量作功。

% [原PDF p38]
原文逐项列出六个独立黏性应力分量，并区分剪切黏度 $\eta$ 与体积黏度 $\eta_b$。单原子气体的硬球分子理论给出 $\eta_b=0$，与实验较吻合；多原子气体通常需要由实验确定非零体积黏度。剪切黏度主要随温度变化而对压力不敏感。

加入黏性后，动量方程右端分别为黏性应力散度的三个分量；总能量方程还含黏性功。若 $\mu$ 为常数并采用 Stokes 假设，动量黏性项可写成
\begin{equation}
\mu\nabla^2\boldsymbol V+\frac{\mu}{3}\nabla(\nabla\cdot\boldsymbol V).
\end{equation}
黏性使方程包含二阶空间导数，完整 Navier--Stokes 系统不再是纯双曲系统，而具有双曲—抛物混合性质。高 Reynolds 数流中，黏性在大部分区域可能很小，却在边界层、剪切层和激波内部决定耗散。

数值激波捕捉通常不解析真实分子激波厚度，而由守恒离散和数值耗散得到正确弱解；这不等于可以在物理黏性问题中任意忽略黏性通量。

% [原PDF p39]
本节把能量通量 $\boldsymbol Q$ 的来源分为温度梯度导热、混合气体扩散和辐射三类，随后只保留第一类。Fourier 定律中的负号明确规定热流与温度梯度反向；在定压比热近似为常数时，$\mathrm{Pr}=c_p\eta/\kappa$ 衡量动量与热量扩散的相对强弱。

热流方向与温度梯度相反，正热导率保证热从高温流向低温并产生熵。总能量边界通量的符号取决于外法向：写成守恒左端时，导热通量通常以 $q\cdot n$ 出现。绝热壁规定 $q\cdot n=0$，等温壁规定温度；二者不能同时任意给定。
\section{热传导}
Fourier 定律给出热流
\begin{equation}
\boldsymbol q=-k\nabla T,
\tag{1.78}
\end{equation}
$k$ 为热导率。能量方程加入 $-\nabla\cdot\boldsymbol q=\nabla\cdot(k\nabla T)$。Prandtl 数 $\mathrm{Pr}=\mu c_p/k$ 比较动量扩散与热扩散；理想气体中常用 $k=\mu c_p/\mathrm{Pr}$。黏度和热导率可随温度变化，例如气体黏度常用 Sutherland 关系。

包含黏性应力与热传导的可压缩 Navier--Stokes 方程仍以质量、动量和总能量守恒为核心，只是通量分成无黏通量与黏性通量。边界条件必须与二阶项匹配，例如固壁需要速度和温度或热流条件。

% [原PDF p40]
无黏通量只依赖局部状态，适合黎曼求解；黏性通量依赖梯度，必须构造面梯度。对显式算法，双曲时间步按 $\Delta x/(|u|+c)$，扩散时间步按 $\Delta x^2/\nu$，细网格时后者更严格。积分形式随后为有限体积中“单元平均量加面通量”的更新提供直接物理解释。
完整微分形式可写成
\begin{equation}
\boldsymbol U_t+\nabla\cdot\boldsymbol F^{\rm inv}(\boldsymbol U)
=\nabla\cdot\boldsymbol F^{\rm visc}(\boldsymbol U,\nabla\boldsymbol U),
\end{equation}
其中黏性能量通量含 $\boldsymbol\tau\boldsymbol V+k\nabla T$。Euler 方程是令 $\mu=k=0$ 的极限。微分形式要求足够光滑；对含激波的无黏解，应回到控制体积分形式定义守恒。

\section{方程的积分形式}
取随时间固定或运动的控制体 $\Omega(t)$，边界为 $\partial\Omega$，外法向为 $\boldsymbol n$。积分守恒律直接陈述“控制体内总量变化率等于穿过边界的净通量及体源”，也是有限体积方法的起点。

% [原PDF p41]
输运定理中的边界速度允许统一固定网格、任意 Lagrange--Euler 网格和随体控制体。若边界速度不同于流速，相对通量应使用 $(V-V_b)\cdot n$。这一形式是 ALE 方法的基础，并要求几何守恒律：均匀流在移动网格上仍应保持均匀。
\subsection{时间导数}
对固定控制体，微分与积分可交换。对以边界速度 $\boldsymbol V_b$ 运动的控制体，Reynolds 输运定理为
\begin{equation}
\frac{d}{dt}\int_{\Omega(t)}\phi,dV
=\int_{\Omega(t)}\phi_t,dV
+\oint_{\partial\Omega(t)}\phi\boldsymbol V_b\cdot\boldsymbol n,dS.
\tag{1.88}
\end{equation}
若控制体随流体运动，$\boldsymbol V_b=\boldsymbol V$，左端就是某一物质体内总量的变化。Gauss 散度定理
\begin{equation}
\int_\Omega\nabla\cdot\boldsymbol a,dV
=\oint_{\partial\Omega}\boldsymbol a\cdot\boldsymbol n,dS
\end{equation}
连接局部散度形式与边界通量形式。
在进入积分守恒律前，先区分两种时间变化率。标量场 $\phi(x,y,z,t)$ 随以流速 $\boldsymbol V=(u,v,w)$ 运动的观察者所测得的变化率为
\begin{equation}
\frac{D\phi}{Dt}=\frac{\partial\phi}{\partial t}
+\boldsymbol V\cdot\nabla\phi.
\tag{1.64}
\end{equation}
它称为随体导数或物质导数。第一项是固定空间点看到的局部变化率，第二项是观察者穿过空间梯度产生的对流变化率。该算子逐分量作用于向量，特别地
\begin{equation}
\frac{D\boldsymbol V}{Dt}=\frac{\partial\boldsymbol V}{\partial t}
+\boldsymbol V\cdot\nabla\boldsymbol V,
\tag{1.65}
\end{equation}
正是流体微团的加速度。

设 $\mathcal V(t)$ 是由分段光滑表面包围、随物质运动的体积，并令
\begin{equation}
\Psi(t)=\iiint_{\mathcal V(t)}\phi(x,y,z,t)\,dV.
\tag{1.66}
\end{equation}
输运定理给出运动区域内总量的物质导数。它把局部变化与边界穿越贡献统一起来，是从粒子表述转到固定控制体表述的关键。对密度取 $\phi=\rho$，便可把“同一团物质质量不变”转写为固定控制面上的净质量通量平衡。

% [原PDF p42]
Reynolds 输运定理把随流体运动的物质体导数分成局部变化和运动边界通量。对标量场 $\phi$，边界项为 $\int_A(\boldsymbol n\cdot\phi\boldsymbol V)\,dA$；对向量场则相应成为张量通量。再用 Gauss 定理，曲面积分即可转为体积分。

质量守恒没有压力或应力项，因为质量只随物质穿过控制面输运。对封闭物质体通量为零；对固定体，流入与流出改变内部质量。离散层面若同一面两侧使用不同质量通量，全局质量将随时间漂移，所以共享数值通量是有限体积实现的基本规则。
\subsection{质量守恒}
物质体质量不变：
\begin{equation}
\frac{d}{dt}\int_{\Omega(t)}\rho,dV=0.
\end{equation}
用输运定理改写到固定控制体，得到
\begin{equation}
\frac{d}{dt}\int_\Omega\rho,dV
+\oint_{\partial\Omega}\rho\boldsymbol V\cdot\boldsymbol n,dS=0.
\tag{1.94}
\end{equation}
第一项是控制体内质量变化率，第二项是向外质量通量。若场光滑，应用散度定理和控制体任意性即可恢复局部连续方程；若含间断，积分式仍有明确意义，并导出跨间断的 Rankine--Hugoniot 条件。

% [原PDF p43]
压力表面力写为 $-p n$，移到守恒通量左侧成为 $+pn$。动量通量张量为 $\rho VV+pI$；与法向相乘得到 $\rho V(V\cdot n)+pn$。该张量形式自动包含法向压力和各方向对流动量，避免逐坐标推导时漏项。
\subsection{动量守恒}
Newton 第二定律说明物质体动量变化率等于表面力与体积力之和。忽略体积力时，
\begin{equation}
\frac{d}{dt}\int_{\Omega(t)}\rho\boldsymbol V,dV
=\oint_{\partial\Omega(t)}\boldsymbol\sigma\boldsymbol n,dS.
\end{equation}
对 Euler 流 $\boldsymbol\sigma=-p\boldsymbol I$，压力在控制面上向内作用。转换到固定体后，
\begin{equation}
\frac{d}{dt}\int_\Omega\rho\boldsymbol V,dV
+\oint_{\partial\Omega}[\rho\boldsymbol V(\boldsymbol V\cdot\boldsymbol n)+p\boldsymbol n],dS=0.
\tag{1.101}
\end{equation}
括号内是对流动量通量和压力通量。三个笛卡尔分量对应式（1.2）—（1.4）。
质量守恒的积分形式为
\begin{equation}
\frac{d}{dt}\iiint_{\mathcal V}\rho\,dV
=-\iint_{\mathcal A}\boldsymbol n\cdot(\rho\boldsymbol V)\,dA.
\tag{1.72}
\end{equation}
左端是控制体所含质量的变化率，右端是流出质量通量的负值；$\rho\boldsymbol V$ 是质量流向量，$\boldsymbol n$ 为外法向。若流场足够光滑，Gauss 定理把表面积分化为体积分，因控制体任意可恢复 $\rho_t+\nabla\cdot(\rho\boldsymbol V)=0$。但激波或接触面上微分不存在，而式（1.72）仍然成立，这正是积分形式更基本的原因。

动量守恒同样从 Newton 第二定律出发。物质体的动量变化率等于作用在其上的表面力和体积力。无黏流的表面应力为 $-p\boldsymbol n$；变换到固定控制体后，对流携带的动量与压力作用共同构成边界通量。若存在重力等体积力，则应在右端加入 $\iiint_{\mathcal V}\rho\boldsymbol g\,dV$。黏性流还要用完整应力张量替代 $-p\boldsymbol I$，其散度恢复 Navier--Stokes 方程中的黏性项。

% [原PDF p44]
加入体积力时，控制体总动量并非严格守恒，但“变化量等于已知源积分”仍成立。重力与地形、旋转源可能和压力梯度形成精确平衡，数值上需良平衡处理。黏性面对相邻单元也应成对抵消，只在外壁通过剪切交换总动量。
若保留黏性应力，则动量积分式右端加入 $\oint\boldsymbol\tau\boldsymbol n,dS$，或把 $-\boldsymbol\tau\boldsymbol n$ 作为总通量的一部分。重力等体积力写成 $\int_\Omega\rho\boldsymbol g,dV$。控制面上的压力和黏性力必须使用外法向约定，符号错误会同时破坏局部通量和全局动量守恒。

有限体积离散正是对每个网格单元近似这些面积分。相邻单元共享同一面的数值通量且符号相反，因此内部通量严格抵消，只剩外边界通量。这一守恒结构是激波速度正确的关键，而逐点非守恒差分一般不具备这种保证。

% [原PDF p45]
压力功和内能输运合并后得到 $(E+p)V$，这是总能量通量而非简单 $EV$。黏性应力做功可把动能转成内能，但总能量在绝热封闭域仍守恒；热传导只重新分配能量。若用内能方程替代总能量推进，必须额外确保激波处正确满足总能量跳跃。
\subsection{能量守恒}
物质体总能量变化来自表面力做功、热流和体积力做功。对无黏、绝热、无体积力 Euler 流，固定控制体形式为
\begin{equation}
\frac{d}{dt}\int_\Omega E,dV
+\oint_{\partial\Omega}(E+p)\boldsymbol V\cdot\boldsymbol n,dS=0.
\tag{1.111}
\end{equation}
通量中的 $E\boldsymbol V$ 是能量随质量输运，$p\boldsymbol V$ 是压力在控制面上做功。若有黏性与热传导，则边界还交换 $\boldsymbol\tau\boldsymbol V$ 和 $-\boldsymbol q$。

把总能量式减去动能方程可得内能式；但在间断处，可靠的弱解定义仍应使用总能量守恒，而不能单独以非守恒内能方程替代。

% [原PDF p46]
法向通量的旋转性质还说明，多维网格面上的黎曼问题可在局部正交基中求解：先把速度旋到法向—切向坐标，求一维波，再把动量通量旋回全局坐标。曲面网格若法向或面积计算不一致，会破坏均匀流保持，即便黎曼求解器本身完全正确。
五个 Euler 积分守恒律可统一写为
\begin{equation}
\frac{d}{dt}\int_\Omega\boldsymbol U,dV
+\oint_{\partial\Omega}\boldsymbol F_n(\boldsymbol U),dS=0,
\qquad
\boldsymbol F_n=n_x\boldsymbol F+n_y\boldsymbol G+n_z\boldsymbol H.
\tag{1.116}
\end{equation}
法向通量只依赖法向速度 $V_n$，切向速度随质量通量输运。旋转坐标使任意面法向问题归结为一维法向 Euler 通量，这是多维有限体积黎曼求解器的基础。

积分形式允许 $\boldsymbol U$ 分段光滑甚至有跳跃。对包围移动间断的薄控制体取极限，即得间断速度 $S$ 满足 $S[\boldsymbol U]=[\boldsymbol F_n]$。守恒只是必要条件；还需熵条件从可能的弱解中选择物理解。

% [原PDF p47]
能量积分式中的压力功、黏性功与热流分别组合为 $(E+p)\boldsymbol V$、$-\boldsymbol V\cdot\boldsymbol\Pi$ 和 $\boldsymbol Q$。若还有重力等体力，右端出现 $\rho\boldsymbol V\cdot\boldsymbol g$；依次忽略黏性、导热和体力，便恢复齐次 Euler 能量方程。

\section{子模型}
从完整方程施加几何、稳态、无黏、不可压等假设，可得到一系列常用子模型。总结形式包括：三维可压缩 Navier--Stokes、三维 Euler、带截面积变化的一维方程、轴对称或球对称 Euler、一维平面 Euler、定常可压缩流、黏性一维流、自由表面重力流、浅水方程、不可压黏性流及人工可压缩系统。

每次约化都应明确保留哪些守恒量和源项。几何源项并非外部物理力，而是把多维散度写成低维坐标形式产生；若离散不平衡，可破坏静止或定常解。后续黎曼求解器主要处理双曲齐次部分，源项另用一致方法耦合。

% [原PDF p48]
Navier--Stokes 的源项可包括重力、化学反应、辐射或旋转坐标力；不同源具有不同时间尺度。算子分裂简单但可能降低平衡精度或在刚性反应中不稳定。通量与源的一致耦合，是把本书无源黎曼方法应用到实际复杂方程时必须额外解决的问题。
完整 Navier--Stokes 的守恒向量仍为 $\boldsymbol U$，总通量分为无黏与黏性部分，并可含体源 $\boldsymbol S$：
\begin{equation}
\boldsymbol U_t+\boldsymbol F_x+\boldsymbol G_y+\boldsymbol H_z
=\boldsymbol F^v_x+\boldsymbol G^v_y+\boldsymbol H^v_z+\boldsymbol S.
\end{equation}
令黏性与热传导为零即得 Euler 双曲系统。高 Reynolds 数算法常以激波捕捉方法离散无黏通量，以中心或专门黏性离散处理二阶通量；两者的时间步限制不同。

模型层次的意义不仅是节省计算。简化模型具有更清楚的波结构，可用于验证黎曼求解器和数值通量；复杂模型则检验几何源项、扩散和边界条件。

% [原PDF p49]
喷管定常解中，喉部可能达到声速，随后在扩张段形成超声速流或激波。面积源使常状态通常不再是定常解；正确的平衡变量包括质量流率和总焓。几何突变处也可定义黎曼型连接，但面积与流体状态同时变化，不能直接当成普通均匀管道界面。
\subsection{截面积变化流动}
准一维管道截面积 $A(x)$ 变化时，对截面平均可得
\begin{align}
(A\rho)_t+(A\rho u)_x&=0,\\
(A\rho u)_t+[A(\rho u^2+p)]_x&=pA_x,\\
(AE)_t+[Au(E+p)]_x&=0.
\tag{1.128}
\end{align}
动量几何源 $pA_x$ 来自倾斜管壁压力。若把它错误地并入齐次通量，不能保持正确喷管定常解。守恒变量可选 $(A\rho,A\rho u,AE)^T$，通量显含给定几何 $A(x)$。

该系统用于收缩—扩张喷管和面积变化管道。局部声学特征仍为 $u-c,u,u+c$，但源项使波强和定常状态沿 $x$ 变化。

% [原PDF p50]
\subsection{轴对称流}
无旋向速度的圆柱轴对称流依赖 $(r,z,t)$。质量方程可写成
\begin{equation}
\rho_t+(\rho u)_z+(\rho v)_r=-\frac{\rho v}{r},
\end{equation}
其余方程亦有 $1/r$ 几何源。等价的乘 $r$ 守恒形式常更适合有限体积离散。轴线 $r=0$ 处表面上的奇异因子由对称正则性消除，径向速度为奇函数，标量和轴向速度为偶函数。

若存在旋向速度，还需角动量方程并出现离心源项。轴对称模型以二维成本描述三维旋转体几何，但不能表示非轴对称不稳定性。
因此，采用轴对称近似前必须确认主导流动确实保持旋转对称；一旦出现周向模态，该模型便会遗漏真实的三维增长机制。

% [原PDF p51]
径向方程统一写成 $\boldsymbol U_t+\boldsymbol F_r=\boldsymbol S$ 时，$\alpha=0,1,2$ 分别对应平面、圆柱和球对称。原式的几何源向量含质量、径向动量的对流部分与能量通量；压力几何贡献的放置取决于采用的等价守恒写法，离散时须与径向通量一致。

球对称爆炸中，波面面积随 $r^2$ 增长，振幅因几何扩散衰减；圆柱对称则随 $r$。源项在原点附近很大，但物理解由对称条件正则。平面一维 Euler 没有几何源，因而是研究波相互作用最纯粹、最常用的基准。
\subsection{圆柱与球对称流}
仅依赖半径 $r$ 的径向流可统一写成几何指数 $\alpha$ 的形式，$\alpha=1$ 对应圆柱对称，$\alpha=2$ 对应球对称：
\begin{equation}
\boldsymbol U_t+\boldsymbol F(\boldsymbol U)_r
=-\frac{\alpha}{r}\boldsymbol G(\boldsymbol U).
\end{equation}
质量源为 $-\alpha\rho u/r$，能量源为 $-\alpha u(E+p)/r$。动量的压力几何项须与通量散度一致处理。

\subsection{平面一维流}
令所有量只依赖 $x,t$ 且横向速度为零，得到经典一维 Euler 方程
\begin{equation}
\boldsymbol U_t+\boldsymbol F(\boldsymbol U)_x=0,\quad
\boldsymbol U=(\rho,\rho u,E)^T.
\tag{1.140}
\end{equation}

% [原PDF p52]
等熵近似令 $p=C\rho^\gamma$，于是能量方程冗余，系统缩为质量和动量两个守恒律。进一步取等温关系 $p=a^2\rho$，声速成为常数，但方程仍是非线性的；再作小扰动线性化，才得到常系数声学系统。这个模型层次在后续特征分析中反复使用。

一维通量为
\begin{equation}
\boldsymbol F=(\rho u,\rho u^2+p,u(E+p))^T,qquad
E=\rho(e+u^2/2).
\end{equation}
理想气体闭合 $p=(\gamma-1)(E-\rho u^2/2)$。这是本书分析黎曼问题、Godunov 方法和高分辨率格式的核心系统。尽管只有三个方程，它包含两族非线性声波和一族线性退化接触波。

若保留横向速度但空间仍一维，$\rho v,\rho w$ 作为额外守恒量以速度 $u$ 输运；它们跨声波可变化受约束，跨接触间断可以跳跃。多维面通量实际求解的就是这种旋转到法向的一维五方程问题。
这种法向旋转使同一套一维黎曼求解器能够复用于各坐标方向或任意网格面的通量计算。

% [原PDF p53]
\subsection{定常可压缩流}
令时间导数为零，二维 Euler 方程为
\begin{equation}
\boldsymbol F(\boldsymbol U)_x+\boldsymbol G(\boldsymbol U)_y=0.
\end{equation}
若主流超声速，可把一个空间方向视作类似时间的推进方向，系统呈双曲性质；若局部亚声速，则空间定常方程具有不同类型，不能用简单下游推进处理。Mach 数 $M=|\boldsymbol V|/c$ 决定扰动能否向上游传播。

定常激波仍满足法向 Rankine--Hugoniot 条件。直接求定常解和把非定常方程推进至稳态是两种不同策略；后者可统一处理跨声速区，但收敛速度受低频误差和边界条件影响。

% [原PDF p54]
定常三维方程可借方向分裂化成一系列增广二维问题，因此作者先强调二维形式。允许激波和滑移线后，微分式 $\boldsymbol F_x+\boldsymbol G_y=0$ 应替换为闭合曲线上的积分守恒式 $\oint(\boldsymbol F\,dy-\boldsymbol G\,dx)=0$，后者在间断处仍有意义。

定常二维超声速方程可选择 $x$ 为推进变量，写成 $\boldsymbol F_x+\boldsymbol G_y=0$，其关于 $\boldsymbol F$ 的特征速度与 Mach 角有关。流向速度必须保持足够大，使相应 Jacobian 可逆。对小扰动，信息沿 Mach 线传播；激波与膨胀扇则是其非线性对应。

工程空气动力学常用定常模型计算喷管、翼型或进气道，但强分离和涡脱落本质非定常。选择子模型时，不应仅因希望得到稳态结果就删除时间项；时间相关性应由物理解和稳定性判断。

% [原PDF p55]
一维可压缩 Navier--Stokes 的源向量中，动量黏性项为 $\tfrac43(\eta u_x)_x$，能量项同时含黏性功 $\tfrac43(\eta u u_x)_x$ 和导热 $-(\kappa T_x)_x$。Burgers 方程及其线性化对流—扩散式正是这一结构的标量简化。

\subsection{一维黏性可压缩流}
在平面一维中加入黏性和热传导，质量方程不变，动量通量含 $-(4/3)\mu u_x$，能量黏性通量含 $(4/3)\mu uu_x+kT_x$。二阶项把激波平滑为有限厚度，并需要额外边界条件。该模型可研究黏性激波结构，也是检验对流与扩散耦合离散的简化平台。

\subsection{自由表面重力流}
不可压液体具有自由表面 $z=\eta(x,y,t)$。完整三维自由表面问题包含内部速度场、运动学边界条件和表面动力条件；在水平尺度远大于水深时，可沿深度平均得到浅水模型。

% [原PDF p56]
自由表面流的几何由静止河床 $h(x,z)$ 与未知表面 $\eta(x,z,t)$ 决定。表面同时满足运动学条件和压力条件：流体质点不能穿过自由面，而在忽略表面张力时压力等于大气压。沿水深积分并采用静水压力近似后，得到后面的浅水方程。

自由表面运动学条件说明界面随流体运动，
\begin{equation}
\eta_t+u\eta_x+v\eta_y=w\quad\text{于 }z=\eta.
\end{equation}
底面不可穿透。若忽略表面张力且上方压力恒定，自由面压力给定。静水近似令竖直压力梯度平衡重力，$p=p_{\rm atm}+\rho g(\eta-z)$；水平速度近似随深度均匀。

沿水深积分质量与水平动量，定义水深 $h$ 和深度平均速度，即得到守恒型浅水方程。该约化保留重力波和水跃，却不解析竖直加速度、色散和三维涡结构。

% [原PDF p57]
浅水“压力”项来自水柱静水重力势能，因此 $gh^2/2$ 与气体压力在数学上扮演相似角色。Froude 数 $|u|/\sqrt{gh}$ 对应 Mach 数：小于一为缓流，大于一为急流。水跃是耗散的压缩间断，虽守恒质量和动量，却损失机械能。
\subsection{浅水方程}
一维无底坡浅水方程为
\begin{align}
h_t+(hu)_x&=0,\\
(hu)_t+\left(hu^2+\frac12gh^2\right)_x&=0.
\tag{1.158}
\end{align}
它与等熵气体动力学具有相似结构，重力波速 $c=\sqrt{gh}$，特征速度 $u\pm c$。水跃对应激波，稀疏水波对应稀疏波；干床使 $h\to0$，是黎曼求解器的特殊退化极限。

底形 $z_b(x)$ 产生源项 $-gh(z_b)_x$，摩阻和 Coriolis 力亦可加入。湖泊静止态 $u=0$、$h+z_b=\text{常数}$ 要求通量与地形源精确平衡。
完整自由表面问题在底面 $y=-h(x,z)$ 和自由面 $y=\eta(x,z,t)$ 上需要边界条件。自由面既是物质面又承受大气压力，故
\begin{equation}
\frac{D}{Dt}(\eta-y)=0,
\qquad p=p_{\rm atm}\quad\text{在 }y=\eta(x,z,t),
\tag{1.140}
\end{equation}
通常把 $p_{\rm atm}$ 取为压力零点。不可穿透底面满足
\begin{equation}
\frac{D}{Dt}(h+y)=0\quad\text{在 }y=-h(x,z).
\tag{1.141}
\end{equation}
即速度法向分量为零。即使已采用无黏、不可压等强假设，这个带未知移动边界的问题仍很难直接计算，因此需要进一步近似。

浅水近似忽略竖直加速度 $Dv/Dt$，竖直动量方程遂给出静水关系
\begin{equation}
p=\rho g(\eta-y),
\tag{1.142}
\end{equation}
并有 $p_x=\rho g\eta_x$、$p_z=\rho g\eta_z$，二者均与 $y$ 无关。因此若水平速度初始时沿深度均匀，它们将保持这一性质，水平动量方程化为
\begin{equation}
u_t+uu_x+wu_z=-g\eta_x,
\qquad
w_t+uw_x+ww_z=-g\eta_z.
\tag{1.145--1.146}
\end{equation}
再把连续方程沿竖直方向积分，并利用上下边界运动学条件，便得到以水深和深度平均速度为未知量的浅水守恒律。

% [原PDF p58]
二维系统中间特征速度为法向流速，携带切向动量，类似 Euler 接触波。旋转不变性允许沿任意面使用法向浅水黎曼问题。干湿界面处声速和通量退化，高阶插值若产生负水深会立即导致虚波速，必须使用正性保持限制。
二维浅水守恒向量与通量为
\begin{equation}
\boldsymbol U=(h,hu,hv)^T,\quad
\boldsymbol F=(hu,hu^2+gh^2/2,huv)^T,\quad
\boldsymbol G=(hv,huv,hv^2+gh^2/2)^T.
\end{equation}
沿任意法向的波结构与一维相同，另有以法向速度传播的切向速度波。因其是小型非线性双曲系统，浅水方程常用于检验 Euler 方法的多维推广。

物理变量要求 $h\ge0$。高阶重构和时间更新必须保持水深非负；否则波速和通量失去意义。处理岸线时还需兼顾守恒、正性和静水平衡。

% [原PDF p59]
浅水系统中的波速为 $u\pm\sqrt{gh}$，因此水深必须保持非负。静水稳态并非简单常状态：变底地形上 $u=0$ 且自由表面高程保持常数，通量梯度与地形源项恰好平衡。数值离散若不保持这种平衡，会在静水上生成伪波。

浅水方程可写成 $\boldsymbol U_t+\boldsymbol F_x+\boldsymbol G_y=\boldsymbol S$。无源齐次系统严格双曲，特征值 $u_n-c,u_n,u_n+c$。它说明黎曼求解器并非气体动力学专属：只要系统具有守恒形式和可识别的波结构，同样思想可用于水力学、交通流、磁流体等问题。

模型推导中的静水压力积分给出 $gh^2/2$；若密度变化、非静水或色散不可忽略，则需更完整模型。数值模型的精度上限受物理模型假设约束，细化网格无法弥补被删去的物理机制。

% [原PDF p60]
不可压压力方程来自对预测速度投影到无散空间。椭圆求解常占并行计算主要成本；人工可压缩法以伪时间局部双曲传播替代全局压力校正，但只适合求稳态或双时间方法。真实非定常不可压解不能把伪时间直接解释为物理时间。
\subsection{不可压黏性流}
常密度不可压 Navier--Stokes 方程为
\begin{equation}
\nabla\cdot\boldsymbol V=0,qquad
\boldsymbol V_t+\nabla\cdot(\boldsymbol V\boldsymbol V)+\nabla p
=\nu\nabla^2\boldsymbol V.
\end{equation}
这里压力不由状态方程给出，而是作为 Lagrange 乘子维持无散约束。标准投影法需求解椭圆压力方程，信息在数学上瞬时传遍全域，结构与可压缩双曲系统不同。

为使用双曲黎曼求解器，可引入人工可压缩性，把无散条件替换为含伪时间压力导数的方程；只有稳态极限与原不可压方程等价。

% [原PDF p61]
矩阵形式中的通量分别为 $(c^2u,u^2+p,uv)^T$ 和 $(c^2v,uv,v^2+p)^T$，源项第一分量为零，后两分量是速度 Laplacian。双曲主部具有人工声波和对流波，可用黎曼方法；黏性仍以抛物项存在。参数选择影响条件数与稳态收敛。
\subsection{人工可压缩方程}
二维形式取
\begin{align}
p_t+(uc^2)_x+(vc^2)_y&=0,\tag{1.177}\\
u_t+(u^2+p)_x+(uv)_y&=\alpha(u_{xx}+u_{yy}),\tag{1.178}\\
v_t+(uv)_x+(v^2+p)_y&=\alpha(v_{xx}+v_{yy}).\tag{1.179}
\end{align}
$c^2$ 是常数人工可压缩因子，$\alpha$ 为黏性系数。令 $\boldsymbol U=(p,u,v)^T$，可写成
\begin{equation}
\boldsymbol U_t+\boldsymbol F(\boldsymbol U)_x+\boldsymbol G(\boldsymbol U)_y
=\boldsymbol S(\boldsymbol U),
\tag{1.180}
\end{equation}
其中源向量含速度 Laplacian。伪时间达到稳态时 $p_t=0$，第一式恢复不可压约束。
二维不可压 Navier--Stokes 可用流函数--涡量变量改写。涡量输运方程
\begin{equation}
\zeta_t+u\zeta_x+v\zeta_y
=\nu(\zeta_{xx}+\zeta_{yy})
\tag{1.171}
\end{equation}
是抛物型对流--扩散方程，而流函数由椭圆型 Poisson 方程
\begin{equation}
\psi_{xx}+\psi_{yy}=-\zeta
\tag{1.172}
\end{equation}
恢复。因为不可压条件 $u_x+v_y=0$，对流部分也能写成守恒形式
\begin{equation}
\zeta_t+(u\zeta)_x+(v\zeta)_y
=\nu(\zeta_{xx}+\zeta_{yy}).
\tag{1.173}
\end{equation}
这种表述把原问题转化为每步求解一个抛物方程和一个椭圆方程。

Chorin 为定常不可压流提出人工可压缩法。先以参考速度 $V_\infty$、长度 $L$ 和压力尺度 $\rho_\infty V_\infty^2$ 无量纲化，黏性系数成为 $\alpha=1/Re_L$。把不可压约束乘以非零参数 $c^2$ 并加上伪时间项 $p_t$，就得到 $p_t+(uc^2)_x+(vc^2)_y=0$。在伪时间收敛后 $p_t\to0$，原散度约束被恢复；在迭代过程中，压力扰动则以有限的人工声速传播。

% [原PDF p62]
人工可压缩方程左端是非线性双曲系统，因此可定义并精确或近似求解黎曼问题，再应用 Godunov 型格式处理一般初值。参数 $c$ 控制伪声速和稳态收敛：过小会导致不可压约束收敛慢，过大则使显式时间步严苛。扩散源仍需相应离散。

本书不继续展开该系统的数值方法；精确与近似黎曼求解器可参见文献 [514]，第 18 章 18.2.3 节还有进一步说明。本章通过这一模型强调：黎曼求解器与守恒型高分辨率技术可超越标准可压缩 Euler 方程，只要问题的主部具有适当双曲结构。

% [原PDF p63]
双曲主部往往决定最快信息传播速度，因此显式算法的稳定时间步首先受它限制；黏性虽可能较小，却引入不同的扩散尺度。理论成熟并不意味着所有系统都有闭式黎曼解，实际方法经常只使用特征速度界限、守恒跳跃和少数重要中间波。后续章节将反复在“精确波结构”和“廉价稳健近似”之间权衡。
\chapter{双曲偏微分方程的基本概念}
本章选择双曲偏微分方程的若干初等性质，重点是分析流体方程和实现数值方法所必需的内容。全书离散技术强烈依赖方程背后的物理和数学结构，因此需要先建立偏微分方程基础。这里几乎只讨论双曲方程，特别是双曲守恒律，原因有三：忽略黏性和热传导后，可压缩流方程化为双曲 Euler 系统；流体方程中双曲项通常对离散提出最严格要求；双曲系统理论比完整 Navier--Stokes 等模型更成熟。

双曲问题还覆盖广泛科技领域，研究活动持续增长。一般 PDE 背景可参阅 John、Zachmanoglou 与 Thoe；守恒律可参阅 Godlewski--Raviart、Hörmander、Tveito--Winther 等。本章采取非形式化表述，主题由黎曼问题、高分辨率迎风和中心方法直接驱动。

% [原PDF p64]
方程中的 $t$ 只需是类时间变量；在定常超声速问题中，也可把某个空间坐标当作推进方向。初值规定整个空间轴上的状态，边界条件则只应从物理域边界的入射特征数目确定。若对双曲系统在所有变量上同时强加边界值，通常会过定；少给入射信息则欠定。
\section{准线性方程：基本概念}
考虑 $m$ 个未知量的一阶系统
\begin{equation}
\frac{\partial u_i}{\partial t}+
\sum_{j=1}^m a_{ij}(x,t,u_1,\ldots,u_m)
\frac{\partial u_j}{\partial x}+b_i(x,t,u_1,\ldots,u_m)=0,
\quad i=1,\ldots,m.\tag{2.1}
\end{equation}
矩阵形式为
\begin{equation}
\boldsymbol U_t+\boldsymbol A\boldsymbol U_x+\boldsymbol B=0.
\tag{2.2}
\end{equation}
若 $A,B$ 均常数，系统为常系数线性；若仅依赖 $x,t$，为变系数线性；若 $A=A(U)$，通常是非线性的准线性系统。$B=0$ 时称齐次。这里空间域取整条实线，故无需边界条件；时间从 $t=0$ 起，并给定初值。标量例子是 $u_t+au_x=0$ 和 Burgers 方程 $u_t+uu_x=0$。
式（2.1）含有 $m$ 个方程和 $m$ 个未知函数 $u_i(x,t)$。矩阵与向量明确写作
\begin{equation}
\boldsymbol U=(u_1,\ldots,u_m)^T,
\quad \boldsymbol B=(b_1,\ldots,b_m)^T,
\quad \boldsymbol A=(a_{ij})_{m\times m}.
\tag{2.3}
\end{equation}
若 $a_{ij}$ 与 $b_i$ 均为常数，系统是常系数线性系统；若只依赖独立变量 $x,t$，是变系数线性系统；$B$ 对 $U$ 线性时系统仍是线性的。若系数矩阵依赖未知状态 $A=A(U)$，则称准线性系统，整体一般仍为非线性。$B=0$ 时称齐次系统。

陈述偏微分方程问题还必须规定独立变量的范围。通常空间域是 $x_l<x<x_r$，其端点需要适当边界条件；本章为突出初值传播，取整条实线 $-\infty<x<\infty$，所以不指定空间边界。时间取 $0<t<\infty$，并在 $t=0$ 给出初始条件。两个基本标量例子为
\begin{equation}
u_t+a u_x=0,
\tag{2.4}
\end{equation}
和无黏 Burgers 方程
\begin{equation}
u_t+u u_x=0.
\tag{2.5}
\end{equation}
前者速度为给定常数，后者的传播速度就是未知量本身，因而会发生波形变陡和激波形成。

% [原PDF p65]
守恒律的每个分量代表一种局部收支，通量只依赖状态的情形称自治系统。准线性形式适合计算特征值，却不定义间断乘积；守恒形式即使在跳跃处仍可通过积分解释。严格双曲保证波速排序清楚，重特征值系统仍可能双曲，但波族的选择和数值分解更细致。
\paragraph{定义 2.1（守恒律）}
能写成
\begin{equation}
\boldsymbol U_t+\boldsymbol F(\boldsymbol U)_x=0
\tag{2.6}
\end{equation}
的方程组称为守恒律，$\boldsymbol U$ 是守恒变量，$\boldsymbol F$ 是通量。
光滑区用链式法则可写成 $\boldsymbol U_t+\boldsymbol A(U)\boldsymbol U_x=0$，其中 $\boldsymbol A=\partial\boldsymbol F/\partial\boldsymbol U$。但间断解必须保留守恒形式。

\paragraph{定义 2.2（双曲性）}
若 $A$ 有 $m$ 个实特征值和一组完备的线性无关右特征向量，系统称双曲；若特征值皆实且互异，称严格双曲。
双曲性是有限速度波传播的代数基础；严格双曲时各波族可清楚排序，而重特征值情形还需检查特征空间维数是否足够。
守恒形式和双曲性承担不同角色：前者保证任意控制体上的收支并规定间断跳跃，后者把局部扰动分解成具有实传播速度的波。一个系统可以写成守恒形式却不在全部状态域内双曲，因此实际求解前还必须界定密度、压力等物理变量所允许的状态集合。
线性化气体动力学的二方程系统以密度扰动和速度扰动为未知量，系数矩阵含参考密度 $\rho_0$ 与声速 $a$。其两个特征速度应为 $-a,+a$，表示左右传播声波。

% [原PDF p66]
右特征向量给出沿某波族状态如何改变，左特征向量把任意扰动投影到各波族。若 $K$ 由右特征向量作列组成，则 $K^{-1}$ 的各行是双正交左特征向量。特征值的正负相对于坐标方向定义；改变面法向时，法向速度和波速符号同时改变，迎风选择也随之改变。
矩阵特征值满足
\begin{equation}
\det(\boldsymbol A-\lambda\boldsymbol I)=0.
\tag{2.11}
\end{equation}
特征值在物理上代表信息传播速度，向 $x$ 增大方向为正。右特征向量满足 $AK^{(i)}=\lambda_iK^{(i)}$，左特征向量满足 $L^{(i)}A=\lambda_iL^{(i)}$；它们的缩放任意，只要不为零。对线性化声学矩阵，特征多项式给出 $\lambda^2-a^2=0$。

寻找右特征向量就是求 $A-\lambda_iI$ 的零空间。由于特征矩阵奇异，方程并非独立，解有任意缩放因子。通常选取使后续分解最方便的归一化，而不是机械规定单位长度。左、右特征向量可归一成 $L^{(i)}K^{(j)}=\delta_{ij}$，从而方便投影到特征量。

% [原PDF p67]
缩放因子不会改变特征方向，但会改变特征系数的数值。程序中若左右特征向量使用不同归一化，重构到物理变量时必须保持互逆。线性声学的两个向量分别表示压力—速度相位相反和相同的扰动，它们张成全部二维状态空间，因此任意小扰动都能唯一拆分。
对 $\lambda_1=-a$，取 $k_1=\alpha_1$ 得
\begin{equation}
K^{(1)}=\alpha_1(1,-a/\rho_0)^T;
\tag{2.17}
\end{equation}
对 $\lambda_2=a$，
\begin{equation}
K^{(2)}=\alpha_2(1,a/\rho_0)^T.
\tag{2.18}
\end{equation}
任意非零 $\alpha_i$ 都正确。两个向量线性无关，故该系统严格双曲。

本页还强调双曲、抛物和椭圆分类取决于主部。热方程的时间一阶、空间二阶结构为抛物型；Laplace 方程为椭圆型；波动传播方程可化成双曲一阶系统。混合型方程的分类可能随状态或空间位置改变，例如定常可压缩流跨越声速时会变型。

% [原PDF p68]
Cauchy--Riemann 方程没有实特征方向，局部值受周围边界整体影响；小扰动超声速方程则有两族实 Mach 线。$M_\infty=1$ 是退化点，系数趋于奇异，简单分类失效。跨声速流同时含椭圆与双曲区，数值边界和求解策略比纯超声速推进复杂得多。
\subsubsection*{例：Cauchy--Riemann 方程}
$u_x-v_y=0$、$v_x+u_y=0$ 写成 $U_x+AU_y=0$，其中
$A=\begin{bmatrix}0&-1\\1&0\end{bmatrix}$。特征方程 $\lambda^2+1=0$ 无实根，所以系统为椭圆型。

\subsubsection*{例：小扰动方程}
定常小扰动系统 $u_x-a^2v_y=0$、$v_x-u_y=0$，其中
$a^2=(M_\infty^2-1)^{-1}$。超声速 $M_\infty>1$ 时特征值实，方程双曲；亚声速时相应特征值非实，方程椭圆。这一例子直观说明：同一物理方程可由 Mach 数改变数学类型。
分类决定所需数据：双曲问题由初值及入射特征边界数据决定，椭圆问题则通常需要整个边界上的条件。

% [原PDF p69]
常系数系统的双曲性要求任意实方向上的系数矩阵具有实特征值和完备特征矢量。若特征值互异，称为严格双曲；重特征值仍可能双曲，但必须拥有足够多线性无关特征矢量。椭圆型或抛物型系统则不具有相同的有限实特征传播结构。

变系数平流只有在速度场无散时，非守恒形式才等同于无源守恒形式。常系数模型虽然简单，却是检验数值相位误差、幅值耗散和间断分辨率的标准问题。周期初值可精确平移一周后比较；光滑波检验收敛阶，方波检验振荡与数值扩散。
\section{线性平流方程}
三维变系数平流方程为 $u_t+a u_x+b u_y+c u_z=0$。若系数光滑，可写成带源守恒式
\begin{equation}
u_t+(au)_x+(bu)_y+(cu)_z=u(a_x+b_y+c_z).
\end{equation}
本节详论常系数一维初值问题
\begin{equation}
u_t+au_x=0,\quad -\infty<x<\infty,\quad t>0,\qquad u(x,0)=u_0(x).
\tag{2.32}
\end{equation}
$a$ 是波速。若初值足够光滑，经典解存在；即使初值分段连续，平流也只搬运间断而不改变形状。该模型包含迎风方向、有限传播速度和 CFL 条件等数值波传播的核心特征。

\subsection{特征线与一般解}
沿曲线 $x=x(t)$，链式法则给出 $du/dt=u_t+x'(t)u_x$。

% [原PDF p70]
图 2.1 中每条线都携带固定的 $u$，彼此不交换信息。特征脚点构成精确迎风规则：正速度从左侧追溯，负速度从右侧追溯。离散格式若时间步过大，真实脚点会越过数值模板边界，数值依赖域无法覆盖物理依赖域，这正是 CFL 限制的几何含义。
令特征线满足
\begin{equation}
\frac{dx}{dt}=a,qquad \frac{du}{dt}=0,
\tag{2.35--2.36}
\end{equation}
即可把 PDE 化为 ODE。通过初始点 $x_0$ 的特征线是
\begin{equation}
x=x_0+at.
\tag{2.38}
\end{equation}
图 2.1 在 $x$--$t$ 平面画出 $a>0$ 时的一族平行特征线；若以横轴 $x$、纵轴 $t$ 作图，线斜率为 $1/a$。$u$ 沿每条线保持初值 $u_0(x_0)$。

通过任意点 $(x,t)$ 向后追踪，可到达 $x_0=x-at$。这说明该点的解只依赖一个初始点，信息传播速度有限。数值格式若使用特征线影响域之外的数据，通常会引入不必要扩散或不稳定性。

% [原PDF p71]
黎曼数据在每个网格面自然出现：有限体积重构给出面左、右状态，局部黎曼解决定面通量。对线性平流，精确 Godunov 通量为 $f^*=au_L$（$a>0$）或 $au_R$（$a<0$）。因此最简单的一阶迎风格式并非经验偏置，而是精确求解每个局部线性黎曼问题的结果。
一般解为
\begin{equation}
u(x,t)=u_0(x-at).
\tag{2.39}
\end{equation}
所以 $a>0$ 时初始剖面向右平移，$a<0$ 时向左平移，形状和幅值不变。这与扩散方程不同：线性平流不会自动平滑间断。守恒形式 $u_t+(au)_x=0$ 在常 $a$ 时与非守恒形式相同。

\subsection{黎曼问题}
黎曼初值由原点两侧常状态构成：
\begin{equation}
u(x,0)=\begin{cases}u_L,&x<0,\\u_R,&x>0.
\end{cases}
\tag{2.40}
\end{equation}
其解仍是平移阶跃，间断位于 $x=at$。这是最简单的黎曼求解器：界面状态仅由特征速度符号决定，$a>0$ 取左状态，$a<0$ 取右状态。
由特征线方程 $dx/dt=a$ 可得通过初始点 $(x_0,0)$ 的直线 $x=x_0+at$。沿该线 $du/dt=0$，故
\begin{equation}
u(x,t)=u_0(x_0)=u_0(x-at).
\tag{2.39}
\end{equation}
所有特征线彼此平行，这是常系数线性双曲方程的典型性质。初始剖面的每一点都在时间 $t$ 内移动距离 $at$：$a>0$ 向右，$a<0$ 向左，波形不改变。这里的“波”是指扰动中可识别的特征以有限速度传播，而不必是正弦振荡。

最简单的分段常数初值问题称为黎曼问题：
\begin{equation}
u_t+a u_x=0,
\qquad
u(x,0)=\begin{cases}u_L,&x<0,\\u_R,&x>0.
\end{cases}
\tag{2.40}
\end{equation}
若 $u_L\ne u_R$，原点有一个初始间断。时间 $t$ 后它位于 $x=at$，该特征线把取值 $u_L$ 的左区与取值 $u_R$ 的右区分开：
\begin{equation}
u(x,t)=\begin{cases}u_L,&x-at<0,\\u_R,&x-at>0.
\end{cases}
\tag{2.41}
\end{equation}
这也是一般平移解 $u_0(x-at)$ 的直接特例。$u_L=u_R$ 时问题退化为全域常状态。

% [原PDF p72]
图 2.3 的间断线把平面分成两个常值区；其斜率由 $a$ 唯一决定。系统情形会从同一原点发出多条波，形成多个中间状态。后续所谓“星区”并非特殊物质区域，只是黎曼波之间的常状态；星号表示它由左右初值共同确定。
图 2.2 表示原点处分隔 $u_L,u_R$ 的初值；图 2.3 表示 $a>0$ 时解在 $x$--$t$ 平面的结构。除通过原点的特征线外，其他特征线两侧都不跨越初始跳跃；解只在 $x-at=0$ 处改变。若 $a=0$，间断静止。

\section{线性双曲系统}
把标量分析推广到常系数系统
\begin{equation}
\boldsymbol U_t+\boldsymbol A\boldsymbol U_x=0.
\tag{2.42}
\end{equation}
双曲性保证 $A$ 有实特征值和完备特征向量。目标是把耦合系统变换成若干相互独立的标量平流方程；每一个特征量沿自己的速度传播。线性系统的解是这些独立波的叠加。

% [原PDF p73]
对角化需要完备特征基。若矩阵有实特征值却缺少足够特征向量，会出现 Jordan 块，解可能含随时间增长的多项式因子，通常不称强双曲。数值上，特征矩阵条件数也重要：接近退化状态时投影可能放大舍入和重构误差。
\subsection{对角化与特征变量}
若
\begin{equation}
\boldsymbol A=\boldsymbol K\boldsymbol\Lambda\boldsymbol K^{-1},
\tag{2.43}
\end{equation}
$\Lambda=\operatorname{diag}(\lambda_1,\ldots,\lambda_m)$，$K$ 的第 $i$ 列为右特征向量 $K^{(i)}$，则系统可对角化。定义特征变量
\begin{equation}
\boldsymbol W=\boldsymbol K^{-1}\boldsymbol U,qquad
\boldsymbol U=\boldsymbol K\boldsymbol W.
\tag{2.44}
\end{equation}
因为 $K$ 常数，导数可直接交换，代回原方程并左乘 $K^{-1}$ 即可解耦。$K^{-1}$ 的各行就是适当归一的左特征向量；$w_i=L^{(i)}U$ 是第 $i$ 个特征量。
这一步同时解释了左右特征向量在高分辨率格式中的分工：左特征向量用于投影，右特征向量用于把各波族的贡献重新合成为物理变量。

% [原PDF p74]
特征变量由 $\boldsymbol W=\boldsymbol K^{-1}\boldsymbol U$ 定义，右特征矩阵的列是 $\boldsymbol A$ 的特征矢量。该变换把耦合系统对角化为若干独立平流方程；初始数据中每个特征分量以自己的 $\lambda_i$ 平移，最后叠加回物理变量。

式（2.53）清楚表明每个物理分量一般同时含所有波族。所谓特征重构，就是先用左特征向量把单元数据分解，再对每个 $w_i$ 按自身迎风方向限制或插值，最后用右特征向量合成。这样可减少在激波附近不同波族相互污染造成的非物理振荡。
标准或特征形式为
\begin{equation}
\boldsymbol W_t+\boldsymbol\Lambda\boldsymbol W_x=0,
\tag{2.46}
\end{equation}
即 $m$ 个方程
\begin{equation}
(w_i)_t+\lambda_i(w_i)_x=0,\quad i=1,\ldots,m.
\end{equation}
所以
\begin{equation}
w_i(x,t)=w_i^{(0)}(x-\lambda_it),qquad
U(x,t)=\sum_{i=1}^mw_i^{(0)}(x-\lambda_it)K^{(i)}.
\tag{2.53}
\end{equation}
每个初始特征分量保持形状，以对应特征速度独立平移。即使原变量强耦合，线性常系数系统在特征空间中只是多条标量波。

% [原PDF p75]
声学特征量也可看成线性 Riemann 不变量。系数中的 $\rho_0$ 把密度扰动换算成与速度相同的量纲，$a$ 则给出压力响应尺度。若参考状态随空间缓慢变化，冻结系数只在局部近似有效，特征线会弯曲，振幅也受系数梯度影响。
\subsubsection*{例：再次考察线性化气体动力学}
取前述 $2\times2$ 声学系统。由右特征向量组成 $K$ 并求逆，特征量可取
\begin{equation}
w_1=\frac{a\rho-\rho_0u}{2a\rho_0},qquad
w_2=\frac{a\rho+\rho_0u}{2a\rho_0}.
\end{equation}
它们分别满足速度 $-a$ 与 $+a$ 的平流方程。给定初始密度扰动 $\rho^{(0)}(x)$ 和速度扰动 $u^{(0)}(x)$，$w_1$ 从 $x+at$ 取值，$w_2$ 从 $x-at$ 取值。

该分解表明，任意小扰动都可拆成等幅或不同幅的左、右行声波。若初值恰满足 $a\rho^{(0)}-\rho_0u^{(0)}=0$，左行分量消失，只剩右行波。

% [原PDF p76]
解式中 $x+at$ 项来自左行波，因为从当前点沿速度 $-a$ 向后到初始轴会到达右侧位置；$x-at$ 项来自右行波。这个看似反直觉的符号可用反向追踪检查。若初始速度为零、密度有局部脉冲，它会分裂为两个向相反方向传播、各约一半幅值的声波。
变回物理变量得
\begin{align}
\rho(x,t)&=\frac{a\rho^{(0)}(x+at)-\rho_0u^{(0)}(x+at)}{2a}
+\frac{a\rho^{(0)}(x-at)+\rho_0u^{(0)}(x-at)}{2a},\\
u(x,t)&=-\frac{a\rho^{(0)}(x+at)-\rho_0u^{(0)}(x+at)}{2\rho_0}
+\frac{a\rho^{(0)}(x-at)+\rho_0u^{(0)}(x-at)}{2\rho_0}.
\end{align}
这就是 d'Alembert 型声学解。密度与速度在两列波中的符号关系不同：右行压缩波中二者同号，左行波中相反。线性叠加只适用于常系数系统；非线性系统的特征速度和向量随状态改变，波之间会相互作用并变形。

% [原PDF p77]
严格双曲排序保证任意固定 $t>0$ 时从左到右按 $\lambda_1,\ldots,\lambda_m$ 穿越波。若某些特征值相等，常状态区可能合并，仍需在重特征子空间内选择适当波。线性问题的波不会相互改变；非线性黎曼解虽仍自相似，但各波曲线和速度由未知中间状态耦合决定。
\subsection{线性系统的黎曼问题}
初值为
\begin{equation}
U(x,0)=\begin{cases}U_L,&x<0,\\U_R,&x>0,
\end{cases}
\tag{2.54}
\end{equation}
并令严格双曲特征值有序 $\lambda_1<\cdots<\lambda_m$。从原点发出 $m$ 条直线波 $x=\lambda_it$，把上半 $x$--$t$ 平面分成 $m+1$ 个常状态区。图 2.4 表示一般结构。由于系统线性，每条波只是一个沿 $K^{(i)}$ 方向的跳跃，波速不依赖左右状态。

把 $U_L,U_R$ 分别展开为特征向量线性组合，系数就是左右特征量。跨第 $i$ 条波时，只有第 $i$ 个系数从左值改成右值。

% [原PDF p78]
分别展开左右状态和只展开总跳跃在数学上等价。前者直接给每区特征系数，后者更贴近数值通量：从 $U_L$ 出发，依次加上各波强即可到 $U_R$。若只需要界面 $x/t=0$ 的状态，可以只累加速度小于零或大于零的一侧波，而无需构造整个时空解。
写成
\begin{equation}
U_L=\sum_i\alpha_iK^{(i)},\qquad
U_R=\sum_i\beta_iK^{(i)}.
\tag{2.56}
\end{equation}
每个 $w_i$ 解标量黎曼问题，
\begin{equation}
w_i(x,t)=\begin{cases}\alpha_i,&x-\lambda_it<0,\\
\beta_i,&x-\lambda_it>0.
\end{cases}
\tag{2.59}
\end{equation}
因此任一点状态由该点向后沿各特征线到达初始轴的左右位置决定。等价地，也可直接分解总跳跃
$\Delta U=U_R-U_L=\sum_i\Delta\omega_iK^{(i)}$，每个波强为 $\Delta\omega_i=\beta_i-\alpha_i$。
对常系数线性系统的黎曼数据，左右状态可在右特征向量基底上分别展开：
\begin{equation}
U_L=\sum_{i=1}^{m}\alpha_iK^{(i)},
\qquad U_R=\sum_{i=1}^{m}\beta_iK^{(i)}.
\tag{2.56}
\end{equation}
于是每个特征变量独立满足速度为 $\lambda_i$ 的标量黎曼问题，初值左、右分别为 $\alpha_i$ 和 $\beta_i$，其解为
\begin{equation}
w_i(x,t)=
\begin{cases}\alpha_i,&x-\lambda_it<0,\\
\beta_i,&x-\lambda_it>0.
\end{cases}
\tag{2.59}
\end{equation}
给定点 $(x,t)$，若 $\lambda_I<x/t<\lambda_{I+1}$，则穿过该点的各特征分量分别从初始轴的左侧或右侧取得数据，因而
\begin{equation}
U(x,t)=\sum_{i=I+1}^{m}\alpha_iK^{(i)}
+\sum_{i=1}^{I}\beta_iK^{(i)}.
\tag{2.60}
\end{equation}

对 $2\times2$ 系统，从原点发出速度为 $\lambda_1<\lambda_2$ 的两条波。第一条波左侧保持 $U_L$，第二条波右侧保持 $U_R$，中间楔形区域称为星形区，状态记为 $U_*$。从左状态跨过第一族波或从右状态反向跨过第二族波，都必须得到同一个 $U_*$；这一兼容性确定两族波强。此图景正是非线性 Euler 黎曼问题中“左波--接触波--右波”结构的线性原型。

% [原PDF p79]
图 2.5 星区公式中的混合系数有明确因果解释：慢波相对于星区点已经从右侧穿过，所以用 $\beta_1$；快波尚从左侧到达，所以用 $\alpha_2$。在一般 $m$ 系统中，位于第 $k$ 与 $k+1$ 波之间的状态取前 $k$ 个 $\beta$ 和后 $m-k$ 个 $\alpha$。
图 2.5 画出 $2\times2$ 系统：慢波 $\lambda_1$ 和快波 $\lambda_2$ 之间是星区 $U^*$。从星区点反向追踪，$\lambda_1$ 特征落在初始间断右侧，$\lambda_2$ 特征落在左侧，故
\begin{equation}
U^*=\beta_1K^{(1)}+\alpha_2K^{(2)}.
\tag{2.61}
\end{equation}
图 2.6 用固定时刻从左向右穿越波系说明同一规则：最左区全部取 $\alpha$；每跨过第 $i$ 波，就把第 $i$ 系数替换成 $\beta_i$。最右区自然恢复全部 $\beta$，即 $U_R$。

% [原PDF p80]
对二波系统，从左状态向右穿过第一条波时，仅第一特征系数由 $\alpha_1$ 变为 $\beta_1$；穿过第二波后第二系数也改变。一般 $m$ 波系统同理，所以第 $i$ 波的强度就是 $\beta_i-\alpha_i$，总状态跳跃是所有波强乘右特征矢量之和。

波强分解 $K\Delta\omega=\Delta U$ 可通过左乘 $K^{-1}$ 求得。数值黎曼求解器通常避免显式矩阵求逆，而用解析左特征向量或解小型线性系统。接近真空或特征退化时，特征矩阵可能病态；此时采用只需要极端波速的 HLL 类求解器往往更稳健。
总跳跃可写为
\begin{equation}
\Delta U=\sum_{i=1}^m(\beta_i-\alpha_i)K^{(i)}.
\end{equation}
因此另一种求解算法是只分解 $\Delta U$。对二方程系统，星区满足
$U^*=U_L+\Delta\omega_1K^{(1)}=U_R-\Delta\omega_2K^{(2)}$。这比对左右状态分别展开更直接，也是 Roe 型黎曼求解器的基本形式。

若某特征速度为负，该波向左；为正则向右。数值界面通量只需确定 $x/t=0$ 处落在哪个常状态区；这就是 Godunov 通量从黎曼解采样的思想。

% [原PDF p81]
星区密度由左右密度平均值减去速度散度贡献，星区速度由左右速度平均值减去密度跳跃产生的声学响应。若 $u_L=u_R$，密度跳跃会产生向两侧传播的压力波并在中间诱导速度；若 $\rho_L=\rho_R$，相向或相离速度则改变星区压缩程度。
\subsection{线性化气体动力学黎曼问题}
展开左右状态可得
\begin{equation}
\alpha_{1,2}=\frac{a\rho_L\mp\rho_0u_L}{2a\rho_0},qquad
\beta_{1,2}=\frac{a\rho_R\mp\rho_0u_R}{2a\rho_0}.
\end{equation}
代入星区组合，
\begin{align}
\rho^*&=\frac12(\rho_L+\rho_R)-\frac{\rho_0}{2a}(u_R-u_L),\\
u^*&=\frac12(u_L+u_R)-\frac{a}{2\rho_0}(\rho_R-\rho_L).
\tag{2.64}
\end{align}
左行波连接 $U_L$ 与 $U^*$，右行波连接 $U^*$ 与 $U_R$。这个公式也是许多线性声学边界条件和近似通量的基础。

% [原PDF p82]
依赖域和影响范围是互逆视角：前者从目标点向过去追踪，后者从初始点向未来展开。对变系数或非线性方程，特征边界不再直线，但只要解光滑即可由 ODE 追踪。激波出现后，多条入射特征终止于激波，弱解的因果结构仍由波速和熵条件控制。
\subsection{若干有用定义}
点 $P=(x^*,t^*)$ 的\textbf{依赖域}是初始数据中能影响 $P$ 的部分。标量平流只有一个脚点；$2\times2$ 系统的依赖域是由最慢、最快反向特征在初始轴截出的区间，图 2.7 给出示意。

一段初始区间的\textbf{决定域}是所有只依赖该区间数据的时空点集合，边界由从区间端点发出的极端特征线形成，如图 2.8。点 $Q=(x_0,0)$ 的\textbf{影响范围}则是会受到该点初值影响的所有时空点；对二方程系统，它位于从 $Q$ 发出的两条特征之间。

这些概念确定边界条件数目、局部时间步的因果范围，也解释显式离散为何必须满足数值依赖域覆盖物理依赖域的 CFL 要求。

% [原PDF p83]
守恒律理论关注的不是任意非线性 PDE，而是散度结构。通量 Jacobian 给局部波速，积分形式给跨间断收支，熵不等式给不可逆方向。三者缺一不可。数值算法也应相应满足一致性、守恒性和熵稳定或足够耗散，才能在网格加密时趋向正确弱解。
图 2.9 显示点 $Q$ 对 $2\times2$ 系统的影响范围。非线性系统中特征线可弯曲、汇聚或发散，但依赖域概念仍成立，边界由极端特征决定。

\section{守恒律}
本节概述发展和应用数值守恒律方法所需的数学性质。考虑
\begin{equation}
U_t+F(U)_x=0,
\tag{2.65}
\end{equation}
其准线性 Jacobian $A(U)=\partial F/\partial U$。与线性系统不同，特征值和特征向量依赖状态，导致波速随解改变、光滑波变形、特征交叉并形成激波。积分形式允许间断解，但也引入弱解不唯一问题，必须辅以熵条件。
一点 $Q=(x_0,0)$ 的影响范围，定义为解会受到该点初始数据影响的所有时空点集合。对具有速度 $\lambda_1<\lambda_2$ 的 $2\times2$ 线性系统，它是从 $Q$ 发出的两条极端特征线之间的区域。图 2.9 示意了这一锥形范围。依赖域回答某个解值需要哪些初始信息，影响域则回答某个初始扰动未来能传播到哪里；二者共同体现双曲方程的有限传播速度。

对非线性守恒律
\begin{equation}
U_t+F(U)_x=0,
\tag{2.65}
\end{equation}
若通量 Jacobian $A(U)=\partial F/\partial U$ 具有实特征值以及一组完备的线性无关右特征向量，系统即为双曲型。书中把特征值排序为 $\lambda_1<\cdots<\lambda_m$。此时特征值与特征向量均随状态改变；为简洁起见，后文有时省略它们的参数 $U$。严格不等式对应严格双曲性，特征值重合时则需另行检查特征向量是否仍完备。

% [原PDF p84]
对移动空间边界，积分形式还会出现边界速度乘状态；在激波随体控制体中，这一项正好产生 $S[U]$。时空控制体图形帮助理解数值通量：底边是旧时刻平均量，顶边是新时刻平均量，两侧边是时间积分通量。守恒更新就是四边总收支为零。
\subsection{守恒律的积分形式}
对一维固定控制区间 $[x_L,x_R]$，
\begin{equation}
\frac{d}{dt}\int_{x_L}^{x_R}U(x,t)\,dx
=F(U(x_L,t))-F(U(x_R,t)).
\tag{2.66}
\end{equation}
它说明区间内守恒量的变化等于左端流入减去右端流出。该式源于物理守恒原理，比微分形式要求更低的光滑性。对任意空间区间成立的积分关系定义弱解；光滑弱解可由微积分恢复微分方程。

在时空矩形 $[x_L,x_R]\times[t_1,t_2]$ 上再对时间积分，得到四条边通量的平衡。有限体积方法正是在每个时空控制体上近似此关系，而不是直接近似某一点的导数。
考虑时空平面中的矩形控制体 $\mathcal V=[x_L,x_R]\times[t_1,t_2]$。一维质量守恒首先给出
\begin{equation}
\frac{d}{dt}\int_{x_L}^{x_R}\rho(x,t)\,dx
=f(x_L,t)-f(x_R,t),\qquad f=\rho u.
\end{equation}
推广到完整系统即为式（2.66）。将其在 $t_1$ 至 $t_2$ 积分，得到
\begin{align}
\int_{x_L}^{x_R}U(x,t_2)\,dx
&=\int_{x_L}^{x_R}U(x,t_1)\,dx
+\int_{t_1}^{t_2}F(U(x_L,t))\,dt\\
&\quad-\int_{t_1}^{t_2}F(U(x_R,t))\,dt.
\tag{2.67}
\end{align}
作者称前一瞬时收支为积分形式 I，称这一时空矩形上的收支为积分形式 II。图 2.10 标出了控制体的四条边；顶、底边携带守恒量，左右边携带通量。

对任意时空域应用 Green 定理，还可写成
\begin{equation}
\oint_{\partial\mathcal V}[U\,dx-F(U)\,dt]=0,
\tag{2.68}
\end{equation}
积分沿边界逆时针进行。这是积分形式 III；矩形公式只是它的特例。任意形状的时空控制体形式在分析移动间断和推导跳跃条件时尤其便利。

% [原PDF p85]
式（2.67）不要求在 $t_1,t_2$ 之间处处可微，只需通量积分存在。若把空间区间取成单个网格单元并以时间平均数值通量近似两侧积分，就得到有限体积更新。高阶方法的核心是更准确地近似面通量随时间的积分，而不是改变守恒收支形式。
时空积分形式为
\begin{align}
\int_{x_L}^{x_R}U(x,t_2)\,dx
&=\int_{x_L}^{x_R}U(x,t_1)\,dx
+\int_{t_1}^{t_2}F(U(x_L,t))\,dt\\
&\quad-\int_{t_1}^{t_2}F(U(x_R,t))\,dt.
\tag{2.67}
\end{align}
也可对任意光滑测试函数积分并分部，得到分布意义的弱形式。间断线的贡献不会消失，而是给出 Rankine--Hugoniot 跳跃条件。

标量守恒律 $u_t+f(u)_x=0$ 的线性平流通量为 $f=au$，无黏 Burgers 通量为 $f=u^2/2$。通量曲率 $f''$ 决定特征速度 $f'(u)$ 如何随状态变化。
更偏数学的弱解定义把 $U$ 和 $F(U)$ 的导数转移到光滑测试函数上。若 $\phi(x,t)$ 连续可微且具有紧支集，则要求
\begin{equation}
\int_0^\infty\!\int_{-\infty}^{\infty}
[\phi_tU+\phi_xF(U)]\,dx\,dt
=-\int_{-\infty}^{\infty}\phi(x,0)U(x,0)\,dx.
\tag{2.69}
\end{equation}
“紧支集”意味着测试函数在某个有界集合外恒为零，因而分部积分时无穷远边界项消失。弱形式不要求 $U$ 可微，所以容纳激波和接触间断；若 $U$ 足够光滑，它又等价于原微分方程。式（2.66）至（2.69）适用于一般守恒律系统，并不限于 Euler 方程。

标量情形为
\begin{equation}
u_t+f(u)_x=0,
\tag{2.70}
\end{equation}
通量 $f(u)$ 必须是守恒量和给定问题参数的完全确定函数。线性平流取 $f(u)=au$，无黏 Burgers 方程取 $f(u)=u^2/2$。前者特征速度恒为 $a$，后者特征速度 $f'(u)=u$ 随状态而变，这一差别导致非线性波的压缩、展开和间断形成。

% [原PDF p86]
交通通量在 $\rho=0$ 和堵塞密度 $\rho_m$ 时均为零，在 $\rho_m/2$ 达到最大。低密度车辆扰动顺车流传播，高密度拥堵扰动可向后传播。Buckley--Leverett 的 $u$ 是某相饱和度，分式通量来自两相相对渗透率竞争；其拐点使一个黎曼解可能由稀疏波接激波组成。
交通流模型可写成
\begin{equation}
\rho_t+f(\rho)_x=0,qquad
f(\rho)=u_m\left(1-\frac{\rho}{\rho_m}\right)\rho,
\tag{2.71}
\end{equation}
$u_m$ 是自由流最大速度，$\rho_m$ 是拥堵密度。通量在中间密度最大，特征速度 $f'(\rho)$ 随密度降低；车辆密度高区可形成向上游传播的拥堵激波。

油藏中的 Buckley--Leverett 方程通量
\begin{equation}
f(u)=\frac{u^2}{u^2+b(1-u)^2}
\tag{2.72}
\end{equation}
通常非凸，需要复合激波与稀疏波。例 2.19 的等温气体动力学取 $p=a^2\rho$，守恒变量 $(\rho,\rho u)$。
线性系统也可直接写成守恒律。若
\begin{equation}
U_t+A U_x=0
\end{equation}
且 $A$ 为常矩阵，只需定义 $F(U)=AU$，便有
\begin{equation}
U_t+F(U)_x=0,
\tag{2.73}
\end{equation}
其通量 Jacobian 显然就是 $A$。

例 2.19 考察等温气体动力学。守恒变量和通量为
\begin{equation}
U=\begin{pmatrix}\rho\\ \rho u\end{pmatrix},
\qquad
F(U)=\begin{pmatrix}\rho u\\ \rho u^2+a^2\rho\end{pmatrix},
\tag{2.74}
\end{equation}
其中常数 $a>0$ 是等温声速。令 $u_1=\rho$、$u_2=\rho u$，则
\begin{equation}
F(U)=\begin{pmatrix}u_2\\u_2^2/u_1+a^2u_1\end{pmatrix},
\qquad
A(U)=\begin{pmatrix}
0&1\\ a^2-u^2&2u
\end{pmatrix}.
\end{equation}
这一显式改写提醒我们：Jacobian 必须对独立的守恒分量求导，不能把原始速度 $u$ 误当作独立变量保持不变。

% [原PDF p87]
等温模型没有独立能量方程，相当于温度由外界始终维持；等熵模型则允许温度随压缩变化但每个光滑粒子熵不变。两者都只有两族声波。完整 Euler 增加能量和接触场，使密度可以在压力、速度连续时发生跳跃。简化模型适用的热力学假设应与问题时间尺度匹配。
等温系统通量 Jacobian 的特征值和右特征向量为
\begin{equation}
\lambda_{1,2}=u\mp a,qquad
K^{(1)}=(1,u-a)^T,\quad K^{(2)}=(1,u+a)^T.
\tag{2.75--2.76}
\end{equation}
故在 $\rho>0$、$a>0$ 时严格双曲。等熵气体动力学则取 $p=p(\rho)$，声速 $a=\sqrt{p'(\rho)}$，波速同样为 $u\mp a$。若 $p=C\rho^\gamma$，$a=\sqrt{C\gamma\rho^{\gamma-1}}$。

这些二方程系统比完整 Euler 少了熵/接触波，却保留非线性声波，是研究激波、稀疏波、黎曼不变量和高分辨率算法的理想模型。物理域要求 $\rho>0$ 且 $p'(\rho)>0$。
容易验证，等温系统的特征值为
\begin{equation}
\lambda_1=u-a,
\qquad \lambda_2=u+a,
\tag{2.75}
\end{equation}
相应右特征向量可取
\begin{equation}
K^{(1)}=\begin{pmatrix}1\\u-a\end{pmatrix},
\qquad
K^{(2)}=\begin{pmatrix}1\\u+a\end{pmatrix}.
\tag{2.76}
\end{equation}
因此 $a>0$ 时两特征值为实数且不同，系统严格双曲。

例 2.20 的等熵气体动力学仍以 $(\rho,\rho u)^T$ 为守恒变量，但动量通量为 $\rho u^2+p$，闭合关系取
\begin{equation}
p=C\rho^\gamma,
\qquad C=\text{常数}.
\tag{2.78}
\end{equation}
练习要求求其 Jacobian、特征值和右特征向量，并推广到任意 $p=p(\rho)$。结果是系统双曲当且仅当 $p'(\rho)>0$；一般声速为 $a=\sqrt{p'(\rho)}$，而
\begin{equation}
\lambda_{1,2}=u\mp a,
\qquad
K^{(1,2)}=\begin{pmatrix}1\\u\mp a\end{pmatrix}.
\tag{2.79--2.80}
\end{equation}
压力随密度单调增加，正是等熵压缩具有实声速、方程保持双曲性的必要条件。

% [原PDF p88]
黏性正则化中的 $\varepsilon$ 决定激波厚度而不改变无黏极限的跳跃速度。令 $\varepsilon\to0^+$ 可导出熵解，符号必须为正；负扩散会放大高波数并导致病态。数值黏性扮演类似选择作用，但其大小随网格变化，理想情况是在光滑区高阶小耗散、间断区足以稳定。
\subsection{非线性与激波形成}
考虑标量初值问题
\begin{equation}
u_t+f(u)_x=0,qquad u(x,0)=u_0(x).
\tag{2.81}
\end{equation}
若只保留无黏守恒律，实际激波内部的黏性和热传导被压缩为零厚度间断。小黏性正则化 $u_t+f(u)_x=\varepsilon u_{xx}$ 可说明物理解选择，但本书主要直接使用积分守恒与熵条件。

光滑时写成
\begin{equation}
u_t+\lambda(u)u_x=0,qquad \lambda(u)=f'(u).
\tag{2.86}
\end{equation}
若 $f''>0$，通量凸；$f''<0$，通量凹；$f''$ 改号则非凸。非凸通量允许更复杂波序列。
对应的积分守恒关系为
\begin{equation}
\frac{d}{dt}\int_{x_L}^{x_R}u(x,t)\,dx
=f(u(x_L,t))-f(u(x_R,t)).
\tag{2.82}
\end{equation}
这里只让通量依赖 $u$ 本身，对某些真实物理过程并不充分。例如可压缩介质中的激波内部同时存在黏性耗散和热传导；一个更接近物理的标量模型会让通量依赖梯度，等价写成
\begin{equation}
u_t+f(u)_x=\alpha u_{xx},\qquad \alpha>0.
\tag{2.83}
\end{equation}
无黏式（2.81）则把有限厚度的耗散层视为间断面。

对光滑解使用链式法则，守恒律化成
\begin{equation}
u_t+\lambda(u)u_x=0,
\qquad \lambda(u)=\frac{df}{du}=f'(u).
\tag{2.84--2.85}
\end{equation}
$\lambda(u)$ 是标量方程的特征速度，对应系统中通量 Jacobian 的特征值。线性平流有常速度 $\lambda=a$；Burgers 方程有 $\lambda=u$；交通流方程则有 $\lambda=u_m(1-2u/\rho_m)$。

通量形状深刻决定解的行为，其中关键是特征速度对状态的单调性。若
\begin{equation}
\lambda'(u)=f''(u)>0,
\end{equation}
特征速度随 $u$ 增大，通量为凸函数；若 $f''(u)<0$，速度单调减小，通量为凹函数；若二阶导数在状态区间内变号，通量非凸。凸、凹情形的波序可通过左右状态大小统一判断，而非凸情形可能出现激波与稀疏波相接的复合结构。

由于 $du/dt=u_t+(dx/dt)u_x$，选择 $dx/dt=\lambda(u)$ 后便有 $du/dt=0$。因此在光滑阶段，每条特征线携带一个不变的状态，而该状态又决定本身的斜率。不同状态的特征线不再平行：若后方特征更快，它们会汇聚并在有限时间相交，经典光滑解随之失效，守恒弱解中出现激波；若后方特征更慢，特征线发散，间隙由连续稀疏扇填充。激波形成不是方程守恒失效，而是微分形式的光滑假设失效。积分关系仍控制跨间断的传播速度，熵条件则排除时间反向的非物理展开激波。

% [原PDF p89]
标量通量的凸凹性由 $f''(u)=\lambda'(u)$ 决定。Burgers 通量全凸，交通流例的通量全凹，而 Buckley--Leverett 通量具有拐点。非凸通量的黎曼解可能含激波与稀疏波的复合结构，不能只依据左右状态大小套用凸通量规则。

特征解隐式是因为速度要在未知脚点的初值上评价。求某点解可解标量方程 $x=x_0+f^{\prime}(u_0(x_0))t$；在特征交叉前解唯一。通量非凸时 $f^{\prime\prime}$ 可变号，同一初始剖面的某些区段压缩、另一些膨胀，交互后产生复合波，简单全局“左大激波、左小稀疏”的规则不再充分。
Burgers 通量满足 $f''=1>0$，交通流通量满足 $f''=-2u_m/\rho_m<0$；练习 2.23 要求分析 Buckley--Leverett 通量并证明其既非全凸也非全凹。

沿非线性特征线
\begin{equation}
\frac{dx}{dt}=\lambda(u),\qquad \frac{du}{dt}=0.
\end{equation}
所以通过初始点 $x_0$ 的特征为
\begin{equation}
x=x_0+\lambda(u_0(x_0))t,qquad u=u_0(x_0).
\tag{2.89--2.90}
\end{equation}
与线性平流不同，各条线的速度由各自携带的初值决定，不再平行。这是波形畸变和有限时激波形成的根源。

% [原PDF p90]
由隐式特征关系求导还得到 $x_{0x}=[1+\lambda'(u_0)u_0't]^{-1}$。只要分母非零，脚点 $x_0(x,t)$ 唯一且经典解成立；分母首次为零时，脚点映射失去单射性，梯度趋于无穷并标志波破裂。

验证导数时，由 $u=u_0(x_0)$ 得 $u_x=u_0^{\prime}(x_0)x_{0x}$、$u_t=u_0^{\prime}(x_0)x_{0t}$；隐式关系给 $x_{0t}=-\lambda(u_0)x_{0x}$，故 $u_t+\lambda u_x=0$。这一计算也显示梯度爆破完全来自坐标映射 Jacobian 为零，而状态幅值沿特征保持有界。
给定 $(x,t)$ 时需隐式求 $x_0$，因此
\begin{equation}
u(x,t)=u_0\!\left[x-\lambda(u_0(x_0))t\right].
\tag{2.91}
\end{equation}
图 2.11 画出一条从 $x_0$ 出发的典型特征。对解求导可验证它满足 PDE，只要映射 $x_0\mapsto x$ 可逆。其 Jacobian 为
\begin{equation}
\frac{\partial x}{\partial x_0}
=1+\lambda'(u_0)u_0'(x_0)t.
\end{equation}
当该量变为零时，多条特征相交，经典单值解失效。最早破裂时间由最负的 $\lambda'(u_0)u'_0$ 决定。

% [原PDF p91]
图示选取多个初始点比较传播速度：在凸通量下，高状态的特征更快。初值上升段的特征相互远离而形成膨胀区，下降段的特征相互追赶而形成压缩区；对凹通量，这两种判别正好反向。

图 2.12 中同一波形可同时含膨胀段和压缩段。凸通量下，上升坡展开、下降坡陡化；Burgers 正弦初值会在负斜率最陡处首先形成激波。破裂前解析特征解可用来精确验证数值格式，破裂后则必须与熵弱解比较，不能继续用多值经典公式。
由隐函数求导，
\begin{equation}
\frac{\partial x_0}{\partial x}
=\frac{1}{1+\lambda'(u_0(x_0))u'_0(x_0)t}.
\tag{2.95}
\end{equation}
故分母趋零时 $u_x$ 无界，尽管 $u$ 本身仍有界。图 2.12 以凸通量说明波陡化：若 $u$ 随 $x$ 增大且 $\lambda'(u)>0$，右侧特征更快，线束发散，剖面展开变平；若 $u$ 随 $x$ 减小，左侧高值更快，特征汇聚，剖面压缩变陡。

线性平流的所有特征速度相同，永不自发陡化。非线性使传播速度依赖振幅，即使初值完全光滑，也会在有限时间生成激波。

% [原PDF p92]
波破裂时间由初始剖面和特征速度的组合决定，而不是由状态振幅单独决定。若 $\min[\lambda'(u_0)u_0']<0$，最早交叉发生在该乘积最负的位置；若乘积处处非负，正时间内特征不会因压缩而交叉。

真实气体激波厚度通常远小于工程网格，直接以物理黏性解析需要极细尺度；激波捕捉则把间断铺在少数单元。只要格式守恒且收敛，铺展宽度可随网格缩小并保持正确跳跃。若非守恒离散得到漂亮而窄的剖面，却以错误速度移动，仍是不可接受的解。
对凹通量，上述压缩与膨胀区域的判据反向。特征相交后，直接延续隐式公式会得到折叠、三值解，显然不能代表单一物理状态。真实黏性把陡峭区平滑成很薄的激波层；无黏极限则以移动间断替代。

激波形成时间可写成
\begin{equation}
t_b=-\frac{1}{\min_{x_0}[\lambda'(u_0(x_0))u'_0(x_0)]},
\end{equation}
前提是分母最小值为负。若处处非负，则正向时间内不会由该光滑初值形成压缩激波。仅用人工平滑估计真实激波厚度可能非常不准，因为厚度取决于实际黏性、热传导和状态方程。

% [原PDF p93]
跨越移动间断对两侧积分分别使用 Leibniz 公式时，移动端点贡献带相反符号。将间断之外的光滑部分用守恒律消去，便得到 $S[ u]=[f]$。这一推导说明跳跃条件来自积分守恒，而不是把微分方程直接代入不连续点。

推导跳跃条件时，两侧体积分的移动上限、下限分别产生 $+Su_L$ 和 $-Su_R$；光滑区内积分用 PDE 化为端点通量，最终只剩间断项。矢量系统完全相同，只是除法不再可行。Rankine--Hugoniot 条件表达几何割线：通量跳跃必须与状态跳跃平行，比例就是波速。
设解除移动线 $x=s(t)$ 外光滑，间断速度 $S=ds/dt$。把积分守恒律应用到跨越间断的固定区间，并对两侧移动积分求导，得到
\begin{equation}
S(u_R-u_L)=f(u_R)-f(u_L),qquad
S=\frac{\Delta f}{\Delta u}.
\tag{2.99}
\end{equation}
这就是标量 Rankine--Hugoniot 条件。它只依赖左右状态和守恒通量，与激波内部细节无关。

对 Burgers 方程，$S=(u_L+u_R)/2$。但是跳跃条件允许压缩激波和非物理稀疏激波，因此守恒不足以保证唯一性；还必须要求熵增加或特征以正确方向进入激波。
换言之，跳跃条件决定传播速度，熵条件负责从所有弱解中选出物理解。

% [原PDF p94]
对 Burgers 压缩数据，激波速度是左右特征速度的算术平均，所以自动位于二者之间。质量从两侧特征流入激波，间断不需要额外初始数据。若数值格式的耗散过小，会在激波后产生振荡；过大则把激波涂宽。熵稳定通量旨在既选择压缩解，又控制非物理离散熵生成。
考虑 Burgers 黎曼数据。若 $u_L>u_R$，凸通量使 $\lambda_L=u_L>u_R=\lambda_R$，两侧特征立即交叉。物理解是
\begin{equation}
u(x,t)=\begin{cases}u_L,&x<St,\\u_R,&x>St,
\end{cases}\qquad S=\frac12(u_L+u_R).
\tag{2.101--2.102}
\end{equation}
图 2.13 显示特征从两侧进入激波。Lax 熵条件为
\begin{equation}
\lambda(u_L)>S>\lambda(u_R),
\tag{2.103}
\end{equation}
说明激波是压缩性的，信息被吸入而不是从间断无故发出。

% [原PDF p95]
稀疏激波虽满足守恒，却对初值微小平滑极不稳定：任意窄的连续过渡都会随时间展开成扇，而不是回到跳跃。小黏性解也趋向稀疏扇。这两种独立判据支持同一熵选择。对于系统，熵通常是凸标量函数 $\eta(U)$，配套熵通量 $q(U)$ 满足 $\eta_t+q_x\le0$ 的约定形式。
若 $u_L<u_R$，直接用 Rankine--Hugoniot 仍可构造同一形式的跳跃，但此时 $\lambda_L<S<\lambda_R$，特征从间断向外发散。图 2.14 称它为稀疏激波或违反熵条件的激波，应予拒绝。物理上，膨胀初值应立即展开成连续扇，而不是维持一个会向外发射信息的间断。

熵条件可通过小黏性极限、稳定性或热力学熵增加导出。对凸标量律，Oleinik 单侧 Lipschitz 条件和所有凸熵不等式给出等价选择。数值格式必须包含足够而一致的耗散，才能收敛到熵解，而不是任意弱解。

% [原PDF p96]
非中心稀疏波保留初始过渡的有限宽度，头、尾分别从 $x_R,x_L$ 出发；随着时间增大，初始宽度相对波扇宽度可忽略。中心扇是黎曼阶跃极限，因此只依赖 $x/t$。数值黎曼求解器可利用这种自相似性，在给定相似速度（通常零）处直接采样，无需显式时间积分整个局部解。
把原点跳跃先替换为 $[x_L,x_R]$ 内的平滑线性过渡。每个初值沿速度 $\lambda(u_0)$ 传播，左缘“尾部”走 $x=x_L+\lambda_Lt$，右缘“头部”走 $x=x_R+\lambda_Rt$。因 $\lambda_L<\lambda_R$，波区宽度随时间增加，梯度减小，图 2.15 给出非中心稀疏波。

令 $x_L,x_R\to0$，得到所有扇内特征从原点发出的中心稀疏波。这个极限对小幅初值扰动稳定，与稀疏激波截然不同。扇内状态由自相似变量 $\xi=x/t$ 唯一决定。

% [原PDF p97]
中心稀疏波的头部和尾部都沿特征传播，二者之间满足 $\lambda(u)=x/t$。若初始过渡区宽度趋于零，扇形结构仍保留有限角宽；这说明稀疏波的展开是非线性传播造成的，不是初始数据人为平滑的结果。

若 $\lambda(u)$ 单调，扇内反函数唯一；非凸通量中同一 $x/t$ 可能对应多个状态，必须用凸包或熵构造选取复合解。稀疏扇头尾是特征边界，不是激波，状态连续地接到常值区。高阶数值解在这些弱导数间断处也可能降低局部收敛阶。
图 2.16 的中心稀疏波解为
\begin{equation}
u(x,t)=\begin{cases}
u_L,&x/t\le\lambda_L,\\
\lambda^{-1}(x/t),&\lambda_L<x/t<\lambda_R,\\
u_R,&x/t\ge\lambda_R.
\end{cases}
\tag{2.110}
\end{equation}
尾特征携带 $u_L$，头特征携带 $u_R$。对 Burgers，$\lambda(u)=u$，故扇内 $u=x/t$。

稀疏波连续但在头、尾处导数可能跳跃。它不是物质扩散造成的平滑，而是不同振幅以不同特征速度展开。非线性 Riemann 解的自相似性来自方程和阶跃初值在 $(x,t)$ 同比缩放下不变。

% [原PDF p98]
允许弱间断以后，同一膨胀阶跃还存在满足 Rankine--Hugoniot 条件的非物理激波解。熵条件排除该解并选择中心稀疏扇；因此弱形式负责守恒，熵不等式负责唯一性。后续数值通量必须同时逼近这两项要求。

图 2.17 的波线从原点发出，激波是一条速度 $S$ 的单线，稀疏波是由连续特征填充的扇形。对黎曼问题，尺度不变性保证任意 $t>0$ 的剖面只是横坐标按 $t$ 拉伸；因此比较不同时间数值解时，应使用相似坐标并区分位置误差与波形误差。
同一膨胀初值既有违反熵的跳跃弱解，也有中心稀疏熵解，说明允许间断后弱解不唯一。熵条件选择后者。Burgers 黎曼问题完整解为：$u_L>u_R$ 时单激波；$u_L<u_R$ 时单稀疏波；两者相等时常状态。图 2.17 分别画出一般单波结构、激波和稀疏波。

对一般凸通量，激波速度是割线斜率 $[f]/[u]$，稀疏扇速度范围是 $f'(u_L)$ 到 $f'(u_R)$。几何上，凸性保证压缩激波割线斜率位于两端切线斜率之间，正好满足 Lax 条件。

% [原PDF p99]
线性退化意味着沿该特征向量方向移动时特征速度不变，所以有限强度跳跃可以与两侧特征平行传播；真正非线性意味着速度沿本族状态曲线单调变化，压缩产生激波、膨胀产生稀疏。某些复杂系统的场既非全线性退化也非全真正非线性，会出现非经典或复合基本波。
\subsection{特征场}
对系统 $U_t+F(U)_x=0$，第 $i$ 个特征值及右特征向量定义第 $i$ 特征场。
\paragraph{定义 2.24（线性退化）}
若
\begin{equation}
\nabla_U\lambda_i(U)\cdot K^{(i)}(U)=0
\tag{2.113}
\end{equation}
对所有物理状态成立，则该场线性退化。
\paragraph{定义 2.25（真正非线性）}
若上述内积处处非零，则该场真正非线性。
真正非线性场产生激波或稀疏波；线性退化场产生接触间断。图 2.17 的标量 Burgers 场真正非线性，而线性平流场线性退化。

% [原PDF p100]
检验真正非线性时，梯度必须对选定状态变量求导，特征向量也要用同一变量表示；换变量后内积的零或非零性质保持，但直接混用会算错。爆轰类比通过简化模型分离波传播和反应源的关键竞争，可用于研究数值分裂误差、点火阈值和反应区分辨率。
等温气体的 $\lambda_{1,2}=u\mp a$ 用守恒变量 $U=(\rho,\rho u)$ 表示后，
$\nabla\lambda_i\cdot K^{(i)}=\mp a/\rho\ne0$，所以两场都真正非线性。

例 2.29 介绍 Fickett 爆轰类比模型：它把无黏 Burgers 方程与反应进度模型结合，保留爆轰传播、反应放热和波耦合的若干基本特征，却比完整反应 Euler 方程易分析。这样的简化模型适合检验数值算法是否能处理激波与刚性源项相互作用，但所得定量结果不能替代完整物理模型。

% [原PDF p101]
对线性系统，$A\Delta U=S\Delta U$ 直接说明跳跃向量是特征向量；非线性系统中连接左右状态的是 Hugoniot 轨迹，而不是单个局部特征向量。波强很小时，Hugoniot 轨迹在基态处与右特征向量相切，这为 Roe 线性化和小扰动近似提供依据。
\paragraph{Rankine--Hugoniot 条件}
系统中速度为 $S_i$ 的单一间断连接 $U_L,U_R$ 时，
\begin{equation}
F(U_R)-F(U_L)=S_i(U_R-U_L).
\tag{2.120}
\end{equation}
它是 $m$ 个代数条件，通常不能像标量情况那样简单除法求 $S_i$。图 2.18 表示左右状态和相应通量跳跃。对常系数线性系统 $F=AU$，条件成为 $(A-S_iI)\Delta U=0$，所以非零跳跃必须沿某个右特征向量，且 $S_i=\lambda_i$。

练习 2.32 要求仅用跳跃条件重新求线性声学黎曼星区，从而验证特征展开、跳跃分解和 Rankine--Hugoniot 三种方法一致。

% [原PDF p102]
从左、右跳跃条件得到的两条直线可在 $(\rho,u)$ 平面相交于星状态。该几何构造与沿黎曼不变量曲线求交完全一致。线性系统中波曲线是直线；非线性气体动力学中，激波曲线与稀疏曲线形状不同，星压力通常需解非线性标量方程。
对左行波 $S_1=-a$ 与右行波 $S_2=a$ 分别应用跳跃条件，可得
\begin{equation}
u^*-u_L=-\frac{a}{\rho_0}(\rho^*-\rho_L),qquad
u_R-u^*=\frac{a}{\rho_0}(\rho_R-\rho^*).
\end{equation}
联立恢复
\begin{equation}
\rho^*=\frac{\rho_L+\rho_R}{2}-\frac{\rho_0}{2a}(u_R-u_L),\quad
u^*=\frac{u_L+u_R}{2}-\frac{a}{2\rho_0}(\rho_R-\rho_L).
\end{equation}
说明同一黎曼解可由波跳跃关系局部构造，无需先显式求全部特征系数。

% [原PDF p103]
广义不变量 ODE 的比例形式允许选择任一非零特征向量分量作参数；若某分量为零，应换用其他分量而不能除零。积分得到的常数组合数通常为 $m-1$，因为沿一条波曲线只有一个自由参数。它们可用于建立边界条件、验证数值波和减少黎曼未知量。
\paragraph{广义黎曼不变量}
沿第 $i$ 波族，若依赖变量为 $W=(w_1,\ldots,w_m)$，右特征向量为 $K^{(i)}$，则不变量由常微分关系
\begin{equation}
\frac{dw_1}{k^{(i)}_1}=\frac{dw_2}{k^{(i)}_2}
=\cdots=\frac{dw_m}{k^{(i)}_m}
\tag{2.123}
\end{equation}
给出。积分后得到在相应波上保持常数的 $m-1$ 个独立组合。它们把高维状态空间中的积分曲线与物理波连接起来。

对线性声学，左行波不变量 $u+a\rho/\rho_0$，右行波不变量 $u-a\rho/\rho_0$。注意“不变量”只针对指定波族，并非在整个流场处处常数。

% [原PDF p104]
练习要求用多种方法重复同一问题，不是冗余：特征展开适合线性代数实现，总跳跃分解适合数值通量，黎曼不变量适合非线性波曲线，Rankine--Hugoniot 适合间断。掌握它们之间的对应关系，才能在新方程组中选择最便于推导和编程的工具。
在线性声学黎曼问题中，跨左波应用
$u^*+a\rho^*/\rho_0=u_L+a\rho_L/\rho_0$，跨右波应用
$u^*-a\rho^*/\rho_0=u_R-a\rho_R/\rho_0$，联立即得同一星区式（2.64）。练习还要求用特征展开、总跳跃展开和广义不变量三种途径求小扰动方程黎曼问题，目的在于熟悉等价工具。

对非线性系统，特征向量随状态变化，必须沿其积分曲线积分；若不能得到闭式不变量，也可数值积分 ODE。近似黎曼求解器常通过线性化避免完整积分，但会牺牲部分非线性精确性。

% [原PDF p105]
一般 $m\times m$ 非线性严格双曲系统的黎曼解由 $m$ 个波族连接 $m+1$ 个常状态。每族波依据特征场性质表现为激波、接触间断或稀疏波；与线性系统不同，波速和波曲线都依赖未知中间状态，因此求解通常需要非线性迭代。

等熵不变量积分中 $a/\rho\propto\rho^{(\gamma-3)/2}$，积分后成为 $2a/(\gamma-1)$。当 $\gamma\to1$ 时应取对数极限，对应等温气体。完整 Euler 稀疏波中熵保持常数，所以同样关系适用于每一侧声学扇，但左右熵常数可不同。
等熵气体 $p=C\rho^\gamma$ 的声速 $a=\sqrt{C\gamma\rho^{\gamma-1}}$。沿左、右波积分得
\begin{equation}
I_L=u+\int\frac{a}{\rho},d\rho
=u+\frac{2a}{\gamma-1},qquad
I_R=u-\frac{2a}{\gamma-1}.
\tag{2.128}
\end{equation}
$I_L$ 跨 $u-a$ 波保持常数，$I_R$ 跨 $u+a$ 波保持常数。这些关系将在完整 Euler 稀疏波求解中再次出现。

\subsection{黎曼问题的基本波解}
一般严格双曲系统的黎曼解由 $m$ 个从原点发出的基本波组成，波间为常状态。这里考虑激波、接触间断和稀疏波。

% [原PDF p106]
一般黎曼解的中间状态数为 $m-1$，但若某波强为零，相邻状态重合。算法仍按固定波序处理，可自然覆盖退化情况。激波速度由跳跃条件确定，接触速度等于线性退化特征值，稀疏波则占据一段特征速度区间；近似求解器的差别主要在保留其中多少结构。
图 2.19 表示非线性系统的一般黎曼扇；图 2.20 分别画出单激波、单接触间断和单稀疏波。可先假设左右状态只由某一波族连接，其余波强为零，从而研究每种基本波；一般问题则把这些连接依次串联。

对第 $i$ 族激波，左右状态满足 Rankine--Hugoniot 条件和相应熵条件。Lax 压缩激波要求
\begin{equation}
\lambda_i(U_L)>S_i>\lambda_i(U_R),
\end{equation}
并对其他波族保持正确排序。图 2.20(a) 中两侧第 $i$ 特征都进入激波。激波曲线在状态空间中由跳跃条件定义，不等同于稀疏积分曲线，但在小波强极限二者相切。

% [原PDF p107]
接触间断没有质量穿过其自身移动面，两侧特征与波面平行；在 Euler 中它是物质界面，可携带密度或组分跳跃。稀疏波内部状态连续但非均匀，沿本族积分曲线变化。图 2.20 的三类波是后续精确 Euler 黎曼解的基本积木：左右各一条声学激波或稀疏波，中间一条接触间断。
\paragraph{接触波}
线性退化第 $i$ 场的接触间断满足
\begin{equation}
F(U_R)-F(U_L)=S_i(U_R-U_L),qquad
\lambda_i(U_L)=\lambda_i(U_R)=S_i,
\tag{2.132--2.134}
\end{equation}
并保持该族广义黎曼不变量。特征线在两侧与间断平行，不被吸入。Euler 接触间断上压力和法向速度连续，密度、熵和切向速度可跳跃。

\paragraph{稀疏波}
真正非线性场的稀疏波是光滑自相似过渡，状态沿 $K^{(i)}$ 积分曲线变化，并满足 $\lambda_i(U)=x/t$。头、尾速度分别为两端特征值；特征在扇内发散，广义不变量保持相应常数。

% [原PDF p108]
本章没有追求一般存在唯一性定理，而是建立数值实现所需的工作语言。遇到新的守恒系统，应依次求通量 Jacobian、确认物理域内双曲性、分类特征场、建立跳跃和不变量关系、确定熵条件，最后才设计黎曼求解器。若某一步失败，直接照搬 Euler 算法往往会隐藏模型本身的困难。
本章的入门概念将在第 3 章用于分析 Euler 方程性质；第 4 章精确求解一维时变 Euler 黎曼问题，第 9—12 章介绍近似黎曼求解器。进一步学习可参考 Lax、Whitham、Chorin--Marsden、Courant--Friedrichs、Smoller、LeVeque 以及 Lax 的相关论文。

核心逻辑可概括为：守恒积分形式允许间断；双曲特征结构给出有限传播速度和波族；非线性使特征汇聚或发散，产生激波与稀疏波；Rankine--Hugoniot 条件决定间断守恒关系；熵条件排除非物理解；黎曼不变量和特征场分类构成精确及近似黎曼求解器的理论基础。
% ===== 译文分片 2（已合并）=====

% [原PDF p109]
\chapter{欧拉方程的若干性质}

本章运用第2章介绍的数学工具，分析非定常欧拉方程的一些基本性质。正如第1章所述，欧拉方程是在可压缩介质中忽略黏性、热传导和体力后得到的。下面将证明这组方程构成双曲型守恒律系统，并讨论其数学结构；特别关注第4章求解黎曼问题所必需的性质。我们要分析方程的特征结构，即求出特征值与特征矢量；考察各特征场的性质；并建立跨越稀疏波、接触间断和激波的基本关系。求特征值与特征矢量通常需要相当多的代数运算，也要求熟悉基本物理量及其相互关系。对于很复杂的方程组，往往还须借助符号计算工具。本章所需的背景可参见第1、2章。

\section{一维欧拉方程}

这里研究采用理想气体状态方程的一维非定常欧拉方程，并分别使用守恒形式与非守恒形式。我们先概述黎曼问题解的基本结构，再详细研究解中出现的各类基本波，从而为第4章构造精确黎曼解奠定基础。

\subsection{守恒形式}

欧拉方程的微分守恒形式为
\begin{equation}
  \bm U_t+\bm F(\bm U)_x=0 . \tag{3.1}
\end{equation}
守恒变量矢量和物理通量矢量分别为
\begin{equation}
 \bm U=\begin{bmatrix}u_1\\u_2\\u_3\end{bmatrix}
 =\begin{bmatrix}\rho\\\rho u\\E\end{bmatrix},\qquad
 \bm F=\begin{bmatrix}f_1\\f_2\\f_3\end{bmatrix}
 =\begin{bmatrix}\rho u\\\rho u^2+p\\u(E+p)\end{bmatrix}. \tag{3.2}
\end{equation}
式中，$\rho$ 为密度，$p$ 为压力，$u$ 为质点速度，$E$ 为单位体积总能量，
\begin{equation}
 E=\rho\left(\frac12u^2+e\right), \tag{3.3}
\end{equation}
而单位质量内能 $e$ 由热量状态方程
\begin{equation}e=e(\rho,p) \tag{3.4}\end{equation}
给出。对理想气体，
\begin{equation}e=\frac{p}{(\gamma-1)\rho},\qquad \gamma=\frac{c_p}{c_v}. \tag{3.5}\end{equation}
结合第1章式(1.36)，声速为
\begin{equation}a=\sqrt{\frac{p/\rho^2-e_\rho}{e_p}}=\sqrt{\frac{\gamma p}{\rho}}. \tag{3.6}\end{equation}

% [原PDF p110]
本页从理想气体状态方程恢复压力 $p=(\gamma-1)[E-\rho u^2/2]$，再把三个通量分量完全写成 $(u_1,u_2,u_3)$ 的函数。Jacobian 的每一列分别是对密度、动量和总能量的偏导；保持这三个量独立，是推导矩阵而不混入原始变量导数的关键。

守恒律(3.1)--(3.2)也可写成拟线性形式
\begin{equation}\bm U_t+\bm A(\bm U)\bm U_x=0,\qquad
\bm A(\bm U)=\frac{\partial\bm F}{\partial\bm U}. \tag{3.7}\end{equation}
这里 $\bm A$ 是通量对守恒变量的 Jacobian 矩阵。

\textbf{命题3.1（Jacobian矩阵）}\quad 对理想气体，
\begin{equation}
\bm A=\begin{bmatrix}
0&1&0\\
\tfrac12(\gamma-3)u^2&(3-\gamma)u&\gamma-1\\
\tfrac12(\gamma-1)u^3-uH&H-(\gamma-1)u^2&\gamma u
\end{bmatrix}. \tag{3.8a}
\end{equation}

\textit{证明。}先把通量的每个分量写成 $u_1=\rho$、$u_2=\rho u$、$u_3=E$ 的函数。显然 $f_1=u_2$。由(3.3)和(3.5)，
\[
p=(\gamma-1)\left(u_3-\frac{u_2^2}{2u_1}\right),
\]
所以
\[
\bm F(\bm U)=\begin{bmatrix}
u_2\\[2pt]
\tfrac12(3-\gamma)u_2^2/u_1+(\gamma-1)u_3\\[2pt]
\gamma u_2u_3/u_1-\tfrac12(\gamma-1)u_2^3/u_1^2
\end{bmatrix}.
\]
逐项求偏导即可得到上述结果。

% [扩译 p110]

推导时最容易出错的是能量通量：必须先用 $p=(\gamma-1)[u_3-u_2^2/(2u_1)]$ 消去压力，再对三个独立守恒分量逐项求偏导。这样得到的 Jacobian 同时包含惯性、压力功和状态方程耦合；它不是把原始变量方程的系数矩阵简单改名。

% [原PDF p111]
\textbf{练习3.2}\quad 用声速 $a$ 和质点速度 $u$ 表示 $\bm A(\bm U)$。

\textbf{解3.3}\quad
\begin{equation}
\bm A=\begin{bmatrix}
0&1&0\\
\tfrac12(\gamma-3)u^2&(3-\gamma)u&\gamma-1\\
\tfrac12(\gamma-2)u^3-\dfrac{a^2u}{\gamma-1}&
\dfrac{a^2}{\gamma-1}+\tfrac12(3-2\gamma)u^2&\gamma u
\end{bmatrix}. \tag{3.8}
\end{equation}
Jacobian 也常用单位质量总焓 $H$ 表示。它与单位质量焓 $h$ 及其他变量的关系是
\begin{equation}H=\frac{E+p}{\rho}=\frac12u^2+h,qquad h=e+\frac p\rho. \tag{3.9}\end{equation}
于是也可写成
\begin{equation}
\bm A=\begin{bmatrix}
0&1&0\\
\tfrac12(\gamma-3)u^2&(3-\gamma)u&\gamma-1\\
(\gamma-1)u^3-\gamma uE/\rho&\gamma E/\rho-\tfrac32(\gamma-1)u^2&\gamma u
\end{bmatrix}. \tag{3.10}
\end{equation}

\textbf{命题3.4（齐次性）}\quad 采用理想气体状态方程(3.5)的欧拉方程满足
\begin{equation}\bm F(\bm U)=\bm A(\bm U)\bm U. \tag{3.11}\end{equation}
\textit{证明。}以(3.8)的 Jacobian 左乘(3.2)的 $\bm U$，逐项恰好重现通量 $\bm F$。

% [扩译 p111]

以 $H=(E+p)/\rho$ 表示第三行可显著缩短特征分析。理想气体通量还满足齐次关系 $\bm F(\bm U)=\bm A(\bm U)\bm U$：把 Jacobian 与守恒变量相乘会逐项恢复质量、动量和能量通量。该性质使非线性 Euler 方程在局部呈现类似常系数线性系统的代数结构，并成为后续通量矢量分裂方法的基础。

% [原PDF p112]
特征多项式最终分解为 $(\lambda-u)[(\lambda-u)^2-a^2]$。中间根 $u$ 对应随流体运动的接触/熵波，两侧根 $u\mp a$ 对应相对于流体以声速传播的声学波；这也解释了最大信号速度通常取 $|u|+a$。

这一性质十分重要，它构成第8章通量矢量分裂类格式的基础。欧拉方程中 $\bm F$、$\bm A$ 与 $\bm U$ 的关系，形式上与第2章常系数线性系统完全相同。稍比(3.5)一般的状态方程也具有这一性质，细节见 Steger 与 Warming [463]。

\textbf{命题3.5}\quad $\bm A$ 的特征值为
\begin{equation}\lambda_1=u-a,\qquad\lambda_2=u,\qquad\lambda_3=u+a, \tag{3.12}\end{equation}
相应右特征矢量可取为
\begin{equation}
\bm K^{(1)}=\begin{bmatrix}1\\u-a\\H-ua\end{bmatrix},\quad
\bm K^{(2)}=\begin{bmatrix}1\\u\\\tfrac12u^2\end{bmatrix},\quad
\bm K^{(3)}=\begin{bmatrix}1\\u+a\\H+ua\end{bmatrix}. \tag{3.13}
\end{equation}
\textit{证明。}将(3.8)代入特征方程 $|\bm A-\lambda\bm I|=0$。整理后，多项式含公因子 $\lambda-u$，故 $\lambda_2=u$。约去该因子后剩下
\[
\lambda^2-2u\lambda+u^2-a^2=0,
\]
其两个实根为 $u-a$ 与 $u+a$。再对每个 $\lambda_i$ 解 $(\bm A-\lambda_i\bm I)\bm K=0$，并适当选取归一化因子，便得到(3.13)。

% [原PDF p113]
三个特征值均为实数，且 $\bm K^{(1)},\bm K^{(2)},\bm K^{(3)}$ 构成一组完备的线性无关特征矢量。因此，理想气体的一维非定常欧拉方程是双曲型的。只要声速 $a>0$，三个特征值还彼此不同，故方程严格双曲。对更一般的状态方程，双曲性仍然成立；第4章的余容气体即是一例。

\subsection{非守恒形式}

欧拉方程也可使用守恒变量以外的变量来表述。对光滑解，各种表述彼此等价；但解若含激波，非守恒形式将给出错误的激波解。第3.3节将借浅水方程和等温气体方程说明这一点。尽管如此，非守恒形式在理论分析上有时比守恒形式方便。近年来，针对非守恒方程构造数值格式的思想也重新受到重视，参见 Karni [278] 和 Toro [508,517]。

\subsubsection{原始变量形式}

取原始变量（或物理变量）矢量 $\bm W=(\rho,u,p)^T$。展开质量守恒方程可得
\begin{equation}\rho_t+u\rho_x+\rho u_x=0. \tag{3.14}\end{equation}
展开动量方程，再利用(3.14)并除以 $\rho$，得到
\begin{equation}u_t+uu_x+\frac1\rho p_x=0. \tag{3.15}\end{equation}
同理重排能量方程，并使用(3.14)、(3.15)，可得
\begin{equation}p_t+\rho a^2u_x+up_x=0. \tag{3.16}\end{equation}

% [原PDF p114]
因此拟线性形式为
\begin{equation}\bm W_t+\bm A(\bm W)\bm W_x=0, \tag{3.17}\end{equation}
其中
\begin{equation}
\bm W=\begin{bmatrix}\rho\\u\\p\end{bmatrix},\qquad
\bm A(\bm W)=\begin{bmatrix}u&\rho&0\\0&u&1/\rho\\0&\rho a^2&u\end{bmatrix}. \tag{3.18}
\end{equation}

\textbf{命题3.6}\quad 系统(3.17)--(3.18)的特征值仍为
\begin{equation}\lambda_1=u-a,\qquad\lambda_2=u,\qquad\lambda_3=u+a, \tag{3.19}\end{equation}
相应右特征矢量为
\begin{equation}
\bm K^{(1)}=\alpha_1\begin{bmatrix}1\\-a/\rho\\a^2\end{bmatrix},\quad
\bm K^{(2)}=\alpha_2\begin{bmatrix}1\\0\\0\end{bmatrix},\quad
\bm K^{(3)}=\alpha_3\begin{bmatrix}1\\a/\rho\\a^2\end{bmatrix}, \tag{3.20}
\end{equation}
其中 $\alpha_i$ 是尺度或归一化因子。左特征矢量可写成
\begin{equation}
\bm L^{(1)}=\beta_1(0,1,-1/\rho a),\quad
\bm L^{(2)}=\beta_2(1,0,-1/a^2),\quad
\bm L^{(3)}=\beta_3(0,1,1/\rho a). \tag{3.21}
\end{equation}
证明留给读者。

\textbf{练习3.7}\quad 适当选取 $\alpha_i$ 和 $\beta_i$，验证左右特征矢量满足双正交归一关系
\begin{equation}\bm L^{(j)}\!\cdot\bm K^{(i)}=\begin{cases}1,&i=j,\\0,&i\ne j.\end{cases} \tag{3.22}\end{equation}

% [扩译 p114]

左右特征矢量的尺度并非任意无关：可选择 $\alpha_i,\beta_i$ 使 $\bm L^{(j)}\!\cdot\bm K^{(i)}=\delta_{ij}$。此双正交归一化让任意扰动能稳定地投影到三族特征波，再从各波幅重构回物理变量；不同归一化不改变特征速度和波族的物理意义。

% [原PDF p115]
第二条特征关系 $dp-a^2d\rho=0$ 沿 $dx/dt=u$ 成立，并与熵的物质输运相容。把 $p$ 换成熵 $s$ 后，第三个方程成为 $s_t+us_x=0$；光滑无黏流中熵随质点保持不变，激波处则必须增加。

\subsubsection{特征方程}

特征值 $u-a,u,u+a$ 分别定义特征方向 $dx/dt=\lambda_i$。对(3.17)，沿第 $i$ 个特征方向有 $\bm L^{(i)}\cdot d\bm W=0$，即
\begin{equation}\bm L^{(i)}\cdot(d\rho,du,dp)^T=0. \tag{3.23}\end{equation}
展开得到三条特征关系
\begin{equation}
\begin{aligned}
dp-\rho a\,du&=0,&&dx/dt=u-a,\\
dp-a^2d\rho&=0,&&dx/dt=u,\\
dp+\rho a\,du&=0,&&dx/dt=u+a.
\end{aligned}\tag{3.24}
\end{equation}
这些微分关系沿相应特征方向成立。数值计算中，将其线性化便可近似求解欧拉方程的黎曼问题，见第9章第9.3节。

\subsubsection{熵变量形式}

熵可写成
\begin{equation}s=c_v\ln\!\left(\frac p{\rho^\gamma}\right)+C_0, \tag{3.25}\end{equation}
等价地
\begin{equation}p=C_1\rho^\gamma e^{s/c_v}. \tag{3.26}\end{equation}
若在原始变量形式中以熵代替压力，则未知量为
\begin{equation}\bm W=(\rho,u,s)^T. \tag{3.27}\end{equation}

\textbf{命题3.8}\quad 熵满足
\begin{equation}s_t+us_x=0. \tag{3.28}\end{equation}
\textit{证明。}由(3.25)和 $a^2=\gamma p/\rho$，
\[
s_t=\frac{c_v}{p}(p_t-a^2\rho_t),\qquad s_x=\frac{c_v}{p}(p_x-a^2\rho_x).
\]
再代入(3.14)、(3.16)，便得到所述结果。

% [原PDF p116]
式(3.28)说明在光滑流区
\begin{equation}\frac{Ds}{Dt}=s_t+us_x=0, \tag{3.29}\end{equation}
即沿质点轨线 $dx/dt=u$ 熵保持不变。因而轨线上满足等熵关系
\begin{equation}p=C\rho^\gamma, \tag{3.30}\end{equation}
其中 $C=C(s_0)$ 由初始熵决定；只要流动保持光滑，它在同一条轨线上不变，但一般会随轨线而异。黎曼问题的初始状态分片常值，所以可直接由左右初始数据计算熵。若整个流场中 $C$ 都是同一个常数，则称为等熵流，有时也称均熵流。

熵变量形式为
\begin{equation}\bm W_t+\bm A(\bm W)\bm W_x=0,\qquad
\bm A=\begin{bmatrix}u&\rho&0\\a^2/\rho&u&(1/\rho)\,\partial p/\partial s\\0&0&u\end{bmatrix}. \tag{3.31--3.32}
\end{equation}
其特征值仍为
\begin{equation}\lambda_1=u-a,\quad\lambda_2=u,\quad\lambda_3=u+a, \tag{3.33}\end{equation}
右特征矢量可取
\begin{equation}
\bm K^{(1)}=\begin{bmatrix}1\\-a/\rho\\0\end{bmatrix},\quad
\bm K^{(2)}=\begin{bmatrix}-\partial p/\partial s\\0\\a^2\end{bmatrix},\quad
\bm K^{(3)}=\begin{bmatrix}1\\a/\rho\\0\end{bmatrix}. \tag{3.34}
\end{equation}
证明留给读者。

\subsection{黎曼问题的基本波解}

下面把欧拉方程黎曼解描述为稀疏波、接触间断和激波等基本波的组合，并逐一建立跨波关系。第4章将利用这些关系，把所有未知状态同给定左右状态联系起来。

一维非定常欧拉方程、左右数据 $(\bm U_L,\bm U_R)$ 的黎曼问题为
\begin{equation}
\bm U_t+\bm F(\bm U)_x=0,\qquad
\bm U(x,0)=\begin{cases}\bm U_L,&x<0,\\\bm U_R,&x>0.\end{cases} \tag{3.35}
\end{equation}

% [原PDF p117]
其物理对应物是气体动力学中的激波管问题；激波管两侧初始速度通常为零，隔膜被理想化为初始间断。一个多世纪以来，激波管及其问题一直是流体力学研究的重要工具。

相似解 $\bm U(x/t)$ 包含与三个特征场对应的三族波。两条外波在类型尚不明确时用从原点发出的两条射线表示，中间波用虚线表示。三族波把平面分为四个常状态：从左到右依次为左初态 $\bm U_L$、第一和第二族波之间的 $\bm U_{*L}$、第二和第三族波之间的 $\bm U_{*R}$，以及右初态 $\bm U_R$。中间两状态合称星形区。解中的波只有三类：稀疏波、接触间断和激波。为判断各波的类型，需要考察 $\bm K^{(i)}$ 所对应的特征场。

\textbf{图3.1}\quad 一维非定常欧拉方程黎曼解在 $x$--$t$ 平面中的结构。三族波分别对应特征值 $u-a$、$u$ 与 $u+a$。

\textbf{命题3.10}\quad $\bm K^{(2)}$ 特征场线性退化，而 $\bm K^{(1)}$、$\bm K^{(3)}$ 特征场真正非线性。

\textit{证明。}对第二特征场，
\[
\nabla\lambda_2=[-u/\rho,1/\rho,0],\qquad
\nabla\lambda_2\cdot\bm K^{(2)}=0,
\]
故其线性退化。第一和第三特征场真正非线性的证明留给读者。

% [原PDF p118]
第二特征场对应接触间断；第一和第三特征场则各自对应一条稀疏波（光滑）或激波（间断）。事先通常不知道两条外波的类型，唯一确定的是中间波总为接触间断。例如，左波可为稀疏波，中间为接触间断，右波为激波。跨越左稀疏波时
\[
\lambda_1(\bm U_L)\leq\lambda_1(\bm U_{*L}),
\]
$\lambda_1$ 从左向右单调增大，波扇两侧的特征线向外发散。跨越右激波时，
\[
\lambda_3(\bm U_{*R})>S_3>\lambda_3(\bm U_R),
\]
特征线汇入激波，这正是熵条件；$S_3$ 是第三族激波速度。接触间断满足
\[
\lambda_2(\bm U_{*L})=\lambda_2(\bm U_{*R})=S_2,
\]
相应特征线与接触间断平行。

\textbf{图3.2}\quad 一种典型黎曼解：左侧为稀疏波，中间为接触间断，右侧为激波；箭头表示各波两侧相应的特征方向。

\subsubsection{接触间断}

可利用广义黎曼不变量分析接触间断两侧哪些量发生变化。对一般 $m\times m$ 双曲系统，若 $\bm W=(w_1,\ldots,w_m)^T$，第 $i$ 个右特征矢量为 $\bm K^{(i)}=(k_1^{(i)},\ldots,k_m^{(i)})^T$，则第 $i$ 族广义黎曼不变量由 $m-1$ 条常微分关系给出：
\[
\frac{dw_1}{k_1^{(i)}}=\frac{dw_2}{k_2^{(i)}}=\cdots=\frac{dw_m}{k_m^{(i)}}.
\]

% [原PDF p119]
对守恒变量形式的第二族波，
\begin{equation}\frac{d\rho}{1}=\frac{d(\rho u)}u=\frac{dE}{u^2/2}. \tag{3.36}\end{equation}
整理这些关系得到
\[
p=\text{常数},\qquad u=\text{常数}.
\]
从原始变量右特征矢量 $\bm K^{(2)}$ 也可直接看出这一结论：跨波跳跃与特征矢量各分量成比例，而速度、压力对应分量均为零。密度一般发生有限跳跃，所有依赖密度的量——单位质量内能、温度、声速、熵等——也会相应跳变。因此，接触间断是一条压力和质点速度连续、密度不连续的间断波。

\subsubsection{稀疏波}

欧拉方程中的稀疏波属于第一和第三特征场。由原始变量特征矢量可知，$\rho,u,p$ 都跨波连续变化。使用熵变量形式的广义黎曼不变量，可得如下结果。

\textbf{命题3.11}\quad 跨越第一、第三族稀疏波的广义黎曼不变量分别为
\begin{align}
I_L(u,a)&=u+\frac{2a}{\gamma-1}=\text{常数},&s&=\text{常数},&&\lambda_1=u-a, \tag{3.37}\\
I_R(u,a)&=u-\frac{2a}{\gamma-1}=\text{常数},&s&=\text{常数},&&\lambda_3=u+a. \tag{3.38}
\end{align}

% [原PDF p120]
\textit{证明。}跨第一族波有
\[
\frac{d\rho}{1}=\frac{du}{-a/\rho}=\frac{ds}{0},
\]
于是得到
\begin{equation}du+\frac a\rho d\rho=0,qquad ds=0. \tag{3.39}\end{equation}
第三族波同理给出
\begin{equation}du-\frac a\rho d\rho=0,qquad ds=0. \tag{3.40}\end{equation}
由于熵不变，可使用 $p=C\rho^\gamma$。结合 $a^2=\gamma p/\rho$，积分得
\[
\int\frac a\rho\,d\rho=\frac{2a}{\gamma-1},
\]
从而重现(3.37)--(3.38)。总之，稀疏波是第一或第三族的光滑波；$\rho,u,p$ 均在其中变化。它呈波扇形，头部与尾部由两条特征线界定，波内完整解将在第4章给出。

\subsubsection{激波}

激波的物理背景可见 Becker [35]、Anderson [10]、Landau 与 Lifshitz [297]，专著 Zeldovich 与 Raizer [599] 亦值得参考。对一维欧拉方程，激波是属于真正非线性第一或第三特征场的间断波，$\rho,u,p$ 都发生跳变。

% [扩译 p120]

第一族与第三族稀疏波分别保持 $u+2a/(\gamma-1)$ 和 $u-2a/(\gamma-1)$ 不变，同时熵沿波扇保持常数。因而跨稀疏波的状态曲线由一个初态唯一参数化，特征线从波扇内部向外发散；这正是容许膨胀波而排除膨胀激波的几何原因。

% [原PDF p121]
考虑以常速 $S_3$ 向右传播的第三族激波。激波前状态记为 $\bm W_R=(\rho_R,u_R,p_R)^T$，激波后状态记为 $\bm W_*=(\rho_*,u_*,p_*)^T$。为推导跨激波关系，把参考系变换到随激波运动的坐标系，使激波在新参考系中静止。密度与压力不变，相对速度为
\begin{equation}\hat u_*=u_*-S_3,qquad\hat u_R=u_R-S_3. \tag{3.41}\end{equation}
在静止激波参考系中，Rankine--Hugoniot 条件为
\begin{align}
\rho_*\hat u_*&=\rho_R\hat u_R, \tag{3.42}\\
\rho_*\hat u_*^2+p_*&=\rho_R\hat u_R^2+p_R, \tag{3.43}\\
\hat u_*(\hat E_*+p_*)&=\hat u_R(\hat E_R+p_R). \tag{3.44}
\end{align}
\textbf{图3.3}\quad 右行激波：（a）固定参考系中激波速度为 $S_3$；（b）以 $S_3$ 运动的参考系中激波静止。

% [扩译 p121]

在随激波移动的参考系中，守恒律成为稳态通量相等：质量通量、动量通量和总焓通量分别跨激波守恒。相对速度的符号必须与激波传播方向一致，否则质量通量 $Q$ 的符号和后续激波速度公式会同时出错。把推导放在静止激波系中，可把移动间断的 Rankine--Hugoniot 条件化为普通代数关系。

% [原PDF p122]
定义单位质量焓
\begin{equation}h_*=e_*+p_*/\rho_*,\qquad h_R=e_R+p_R/\rho_R. \tag{3.45}\end{equation}
利用质量与能量跃迁关系可得
\begin{equation}\frac12\hat u_*^2+h_*=\frac12\hat u_R^2+h_R. \tag{3.46}\end{equation}
再把质量关系代入动量关系，经整理得到
\begin{align}
\hat u_*^2&=\frac{\rho_R}{\rho_*}\frac{p_R-p_*}{\rho_R-\rho_*}, \tag{3.47}\\
\hat u_R^2&=\frac{\rho_*}{\rho_R}\frac{p_R-p_*}{\rho_R-\rho_*}. \tag{3.48}
\end{align}
将其代入(3.46)，有
\begin{equation}h_*-h_R=\frac12(p_*-p_R)\frac{\rho_*+\rho_R}{\rho_*\rho_R}. \tag{3.49}\end{equation}
以 $h=e+p/\rho$ 重写，则
\begin{equation}e_*-e_R=\frac12(p_*+p_R)\frac{\rho_*-\rho_R}{\rho_*\rho_R}. \tag{3.50}\end{equation}
到这里尚未对一般热量状态方程作任何特殊假设。对理想气体，把(3.5)代入(3.50)，得到
\begin{equation}
\frac{\rho_*}{\rho_R}=
\frac{p_*/p_R+(\gamma-1)/(\gamma+1)}{[(\gamma-1)/(\gamma+1)]p_*/p_R+1}. \tag{3.51}
\end{equation}

% [扩译 p122]

消去两侧相对速度后得到 Hugoniot 能量关系：内能跳跃等于平均压力乘比容变化的一半。它与具体参考系无关，并同质量、动量跳跃共同定义激波曲线。对理想气体代入 $e=p/[(\gamma-1)\rho]$，即可把密度比完全写成压力比的有理函数；因此只要给定激波前状态和一个压力比，激波后热力学状态便已确定。

% [原PDF p123]
式（3.53）还给出理想气体激波压缩比的上界：相对 Mach 数趋于无穷时，$\rho_*/\rho_R\to(\gamma+1)/(\gamma-1)$。压力比则随相对 Mach 数平方增长。这个有限密度压缩上界是检查强激波计算结果的重要物理约束。

引入来流 Mach 数和激波 Mach 数
\begin{equation}M_R=u_R/a_R,qquad M_S=S_3/a_R. \tag{3.52}\end{equation}
由(3.48)、(3.51)和(3.52)可将跨激波密度比、压力比表示为相对 Mach 数 $M_R-M_S$ 的函数：
\begin{align}
\frac{\rho_*}{\rho_R}&=\frac{(\gamma+1)(M_R-M_S)^2}{(\gamma-1)(M_R-M_S)^2+2}, \tag{3.53}\\
\frac{p_*}{p_R}&=\frac{2\gamma(M_R-M_S)^2-(\gamma-1)}{\gamma+1}. \tag{3.54}
\end{align}
由压力比可得右行激波速度
\begin{equation}
S_3=u_R+a_R\sqrt{\frac{\gamma+1}{2\gamma}\frac{p_*}{p_R}+\frac{\gamma-1}{2\gamma}}. \tag{3.55}
\end{equation}
当激波强度趋于零时 $p_*/p_R\to1$，于是 $S_3\to u_R+a_R$，正如特征速度所要求的。激波后质点速度由质量关系给出：
\begin{equation}u_*=(1-\rho_R/\rho_*)S_3+u_R\rho_R/\rho_*. \tag{3.56}\end{equation}

\textbf{例3.12（激波）}\quad 设 Mach 数 $M_S=3$ 的激波在大气中传播，激波前 $\rho_R=1.225\,\mathrm{kg/m^3}$、$u_R=0$、$p_R=101325\,\mathrm{Pa}$，取 $\gamma=1.4$。由声速公式得 $a_R=340.294\,\mathrm{m/s}$，故 $S_3=1020.882\,\mathrm{m/s}$。再由(3.53)、(3.54)、(3.56)依次得到 $\rho_*=4.725\,\mathrm{kg/m^3}$、$p_*=1\,047\,025\,\mathrm{Pa}$、$u_*=756.2089\,\mathrm{m/s}$。

% [扩译 p123]

压力比随相对 Mach 数的平方增长，而密度比在强激波极限趋于 $(\gamma+1)/(\gamma-1)$，不会无限增大。弱激波极限中传播速度趋近相应声学特征速度，说明有限振幅激波曲线与无穷小特征方向平滑衔接。熵条件要求第三族特征从两侧进入右行激波。

% [原PDF p124]
\textbf{注3.13}\quad 给定激波前状态 $(\rho_R,u_R,p_R)^T$ 以及激波 Mach 数 $M_S$（或等价的速度 $S_3$），关系(3.53)、(3.54)、(3.56)唯一确定第三族激波后的状态。这些关系很适合构造单激波测试问题，用来检验数值方法。

左行第一族激波的分析完全类似。激波前左状态记为 $\bm W_L=(\rho_L,u_L,p_L)^T$，激波后右状态为 $\bm W_*=(\rho_*,u_*,p_*)^T$。相对速度和 Mach 数为
\begin{align}
\hat u_L&=u_L-S_1,&\hat u_*&=u_*-S_1, \tag{3.57}\\
M_L&=u_L/a_L,&M_S&=S_1/a_L. \tag{3.58}
\end{align}
密度--压力比关系为
\begin{equation}
\frac{\rho_*}{\rho_L}=
\frac{p_*/p_L+(\gamma-1)/(\gamma+1)}{[(\gamma-1)/(\gamma+1)]p_*/p_L+1}. \tag{3.59}
\end{equation}
等价地，
\begin{align}
\frac{\rho_*}{\rho_L}&=\frac{(\gamma+1)(M_L-M_S)^2}{(\gamma-1)(M_L-M_S)^2+2}, \tag{3.60}\\
\frac{p_*}{p_L}&=\frac{2\gamma(M_L-M_S)^2-(\gamma-1)}{\gamma+1}. \tag{3.61}
\end{align}

% [扩译 p124]

第一族激波是第三族结果的镜像：它向左传播，质量通量符号随参考方向改变，但压力比、密度比与相对 Mach 数的函数形式不变。数值实现中最好为左右波分别保留明确公式，而不要只靠符号替换，以免在 $S_1$、$S_3$ 和 $Q_L$、$Q_R$ 的正负号上产生隐蔽错误。

% [原PDF p125]
多维通量的结构由旋转不变性决定：沿任意面法向，质量通量是 $\rho u_n$，法向动量通量为 $\rho u_n^2+p$，两个切向动量以 $u_n$ 对流，能量通量为 $u_n(E+p)$。因此任意方向问题都可旋转成一维法向问题。

第一族激波速度与激波后速度分别为
\begin{align}
S_1&=u_L-a_L\sqrt{\frac{\gamma+1}{2\gamma}\frac{p_*}{p_L}+\frac{\gamma-1}{2\gamma}}, \tag{3.62}\\
u_*&=(1-\rho_L/\rho_*)S_1+u_L\rho_L/\rho_*. \tag{3.63}
\end{align}
弱激波极限 $p_*/p_L\to1$ 下，$S_1\to u_L-a_L$。给定左侧激波前状态和 $M_S$（或 $S_1$），式(3.60)、(3.61)、(3.63)确定第一族激波后的完整状态。

\section{多维欧拉方程}

三维欧拉方程的微分守恒形式为
\begin{equation}\bm U_t+\bm F(\bm U)_x+\bm G(\bm U)_y+\bm H(\bm U)_z=0, \tag{3.64}\end{equation}
其中
\begin{equation}
\bm U=\begin{bmatrix}\rho\\\rho u\\\rho v\\\rho w\\E\end{bmatrix},\quad
\bm F=\begin{bmatrix}\rho u\\\rho u^2+p\\\rho uv\\\rho uw\\u(E+p)\end{bmatrix},\quad
\bm G=\begin{bmatrix}\rho v\\\rho uv\\\rho v^2+p\\\rho vw\\v(E+p)\end{bmatrix},\quad
\bm H=\begin{bmatrix}\rho w\\\rho uw\\\rho vw\\\rho w^2+p\\w(E+p)\end{bmatrix}. \tag{3.65}
\end{equation}
这里 $u,v,w$ 是三个坐标方向的速度分量，
\begin{equation}E=\rho\left[\frac12(u^2+v^2+w^2)+e\right]. \tag{3.66}\end{equation}

% [扩译 p125]

三维守恒变量含两个额外横向动量分量，三个方向的通量都包括压力对法向动量的贡献和总焓输运。虽然本章先沿坐标方向分析，实际网格面可以任意取向；后续旋转不变性表明，只要把速度分解为法向与切向分量，仍可复用一维黎曼波结构。

% [原PDF p126]
对应的积分守恒形式为
\begin{equation}
\frac{d}{dt}\int_V\bm U\,dV+\int_A\bm{\mathcal H}\cdot\bm n\,dA=0,
\qquad \bm{\mathcal H}=(\bm F,\bm G,\bm H). \tag{3.67}
\end{equation}
其中 $V$ 为控制体，$A=\partial V$ 为其边界，$\bm n$ 是指向控制体外的单位法向，$dA$ 是面积元，$\bm{\mathcal H}\cdot\bm n$ 是穿过边界的法向通量。该式说明：控制体内守恒量的变化率完全由穿过控制面的总通量决定。第16章介绍的有限体积法正是以此积分形式为基础；其推导见第1章第1.5、1.6.1节。

\subsection{二维方程的守恒形式}

二维欧拉方程为
\begin{equation}\bm U_t+\bm F(\bm U)_x+\bm G(\bm U)_y=0, \tag{3.68}\end{equation}
其中
\begin{equation}
\bm U=\begin{bmatrix}\rho\\\rho u\\\rho v\\E\end{bmatrix},\quad
\bm F=\begin{bmatrix}\rho u\\\rho u^2+p\\\rho uv\\u(E+p)\end{bmatrix},\quad
\bm G=\begin{bmatrix}\rho v\\\rho uv\\\rho v^2+p\\v(E+p)\end{bmatrix}. \tag{3.69}
\end{equation}

\subsubsection{特征结构}

\textbf{命题3.14}\quad $x$ 方向分裂方程的通量 Jacobian 为
\begin{equation}
\bm A=\begin{bmatrix}
0&1&0&0\\
-u^2+\tfrac12(\gamma-1)V^2&(3-\gamma)u&-(\gamma-1)v&\gamma-1\\
-uv&v&u&0\\
u[\tfrac12(\gamma-1)V^2-H]&H-(\gamma-1)u^2&-(\gamma-1)uv&\gamma u
\end{bmatrix}, \tag{3.70}
\end{equation}
其中 $V^2=u^2+v^2$。其特征值为
\begin{equation}\lambda_1=u-a,qquad\lambda_2=\lambda_3=u,qquad\lambda_4=u+a, \tag{3.71}\end{equation}
右特征矢量可取
\begin{equation}
\bm K^{(1)}=\begin{bmatrix}1\\u-a\\v\\H-au\end{bmatrix},\quad
\bm K^{(2)}=\begin{bmatrix}1\\u\\v\\V^2/2\end{bmatrix},\quad
\bm K^{(3)}=\begin{bmatrix}0\\0\\1\\v\end{bmatrix},\quad
\bm K^{(4)}=\begin{bmatrix}1\\u+a\\v\\H+au\end{bmatrix}. \tag{3.72}
\end{equation}
证明作为练习留给读者。

% [原PDF p127]
\subsubsection{旋转不变性}

旋转不变性既可用于证明二维方程关于时间的双曲性，也可在数值计算中处理与 Cartesian 坐标不对齐的边界。二维控制面外法向写成
\begin{equation}\bm n=(n_1,n_2)=(\cos\theta,\sin\theta),\qquad0\leq\theta\leq2\pi, \tag{3.73}\end{equation}
所以法向通量为
\begin{equation}(\bm F,\bm G)\cdot\bm n=\cos\theta\,\bm F(\bm U)+\sin\theta\,\bm G(\bm U). \tag{3.74}\end{equation}
\textbf{图3.4}\quad $x$--$y$ 平面内的控制体：边界为 $A$，外法向为 $\bm n$，$\theta$ 是 $x$ 方向与法向之间的逆时针夹角。

\textbf{命题3.15（旋转不变性）}\quad 对任意 $\theta$ 和 $\bm U$，二维欧拉方程满足
\begin{equation}
\cos\theta\,\bm F(\bm U)+\sin\theta\,\bm G(\bm U)
=\bm T^{-1}\bm F(\bm T\bm U), \tag{3.75}
\end{equation}
其中旋转矩阵及其逆为
\begin{equation}
\bm T=\begin{bmatrix}1&0&0&0\\0&c&s&0\\0&-s&c&0\\0&0&0&1\end{bmatrix},\quad
\bm T^{-1}=\begin{bmatrix}1&0&0&0\\0&c&-s&0\\0&s&c&0\\0&0&0&1\end{bmatrix},
\quad c=\cos\theta, s=\sin\theta. \tag{3.76}
\end{equation}

% [扩译 p127]

旋转矩阵只混合两个动量分量，密度和总能量保持不变。恒等式 $n_1\bm F+n_2\bm G=\bm T^{-1}\bm F(\bm T\bm U)$ 表明沿任意单位法向的通量等于先旋转状态、沿局部 $x$ 方向求通量、再旋转回去。该结果把一般边界面上的特征分解化为标准一维问题。

% [原PDF p128]
\textit{证明。}先计算 $\widehat{\bm U}=\bm T\bm U=(\rho,\rho\hat u,\rho\hat v,E)^T$，其中
\[
\hat u=u\cos\theta+v\sin\theta,qquad
\hat v=-u\sin\theta+v\cos\theta.
\]
于是
\[
\bm F(\widehat{\bm U})=
(\rho\hat u,\rho\hat u^2+p,\rho\hat u\hat v,\hat u(E+p))^T.
\]
左乘 $\bm T^{-1}$ 后，第一与第四分量显然等于 $\cos\theta\bm F+\sin\theta\bm G$ 的相应分量；展开第二、第三分量也完全相同，命题得证。

\subsubsection{关于时间的双曲性}

\textbf{定义3.16}\quad 若对所有可容许状态 $\bm U$ 和实角度 $\theta$，矩阵
\begin{equation}
\bm A(\bm U,\theta)=\cos\theta\,\bm A(\bm U)+\sin\theta\,\bm B(\bm U) \tag{3.77}
\end{equation}
均可对角化，则称二维系统关于时间是双曲的；$\bm A$ 与 $\bm B$ 分别是 $\bm F$ 与 $\bm G$ 的 Jacobian。

\textbf{命题3.17}\quad 二维非定常欧拉方程关于时间双曲。

\textit{证明。}需要证明存在对角矩阵 $\bm\Lambda$ 和非奇异矩阵 $\bm K$，使
\begin{equation}\bm A(\bm U,\theta)=\bm K(\bm U,\theta)\bm\Lambda(\bm U,\theta)\bm K^{-1}(\bm U,\theta). \tag{3.78}\end{equation}
对旋转不变关系(3.75)关于 $\bm U$ 求导，得到
\[
\bm A(\bm U,\theta)=\bm T^{-1}\bm A(\bm T\bm U)\bm T.
\]

% [扩译 p128]

证明双曲性时，对每个法向 $\bm n$ 考察组合 Jacobian $\bm A(\bm U,\theta)=n_1\bm A+n_2\bm B$。旋转恒等式使它与局部法向状态上的 $x$ 向 Jacobian 相似，因此两者具有相同实特征值并拥有完整特征基。双曲性由此与坐标系选择无关。

% [原PDF p129]
由于 $x$ 方向 Jacobian 可由(3.72)的四个线性无关右特征矢量对角化，
\[
\bm A(\bm U)=\bm K(\bm U)\bm\Lambda(\bm U)\bm K^{-1}(\bm U),
\]
故
\[
\bm A(\bm U,\theta)=
[\bm T^{-1}\bm K(\bm T\bm U)]\bm\Lambda(\bm T\bm U)
[\bm T^{-1}\bm K(\bm T\bm U)]^{-1}.
\]
取 $\bm K(\bm U,\theta)=\bm T^{-1}\bm K(\bm T\bm U)$、$\bm\Lambda(\bm U,\theta)=\bm\Lambda(\bm T\bm U)$ 即得结论。

\subsubsection{特征场}

\textbf{命题3.18}\quad 第一、第四特征场真正非线性，第二、第三特征场线性退化。

\textit{证明。}由于
\[
\nabla\lambda_2=\nabla\lambda_3=(-u/\rho,1/\rho,0,0),
\]
直接与(3.72)中的 $\bm K^{(2)}$、$\bm K^{(3)}$ 点乘便得零。另一方面，
\[
\nabla\lambda_1\cdot\bm K^{(1)}=-\frac{(\gamma+1)a}{2\rho}\ne0,qquad
\nabla\lambda_4\cdot\bm K^{(4)}=\frac{(\gamma+1)a}{2\rho}\ne0.
\]
在黎曼问题中，跨第二、第三族波压力 $p$ 与法向速度 $u$ 都不变。第二族是密度跳变的接触间断，第三族是切向速度跳变的剪切波；第一、第四族则为激波或稀疏波。

% [扩译 p129]

二维系统的两个重特征值都等于法向速度：一个对应熵或密度扰动，另一个对应切向速度扰动；它们都是线性退化场。两列外侧声学波是真正非线性的。重特征值并不破坏双曲性，因为其特征子空间仍有两个线性无关的特征矢量。

% [原PDF p130]
\subsection{三维方程的守恒形式}

以下把二维结果推广到三维；证明只涉及基础但冗长的代数，故省略。令 $\widehat\gamma=\gamma-1$，则 $x$ 方向通量 Jacobian 为
\begin{equation}
\bm A=\frac{\partial\bm F}{\partial\bm U}=
\begin{bmatrix}
0&1&0&0&0\\
\widehat\gamma H-u^2-a^2&(3-\gamma)u&-\widehat\gamma v&-\widehat\gamma w&\widehat\gamma\\
-uv&v&u&0&0\\
-uw&w&0&u&0\\
\tfrac12u[(\gamma-3)H-a^2]&H-\widehat\gamma u^2&-\widehat\gamma uv&-\widehat\gamma uw&\gamma u
\end{bmatrix}, \tag{3.79}
\end{equation}
其中
\begin{equation}H=\frac{E+p}{\rho}=\frac12V^2+\frac{a^2}{\gamma-1},qquad V^2=u^2+v^2+w^2. \tag{3.80}\end{equation}
该分裂系统的实特征值是
\begin{equation}\lambda_1=u-a,\quad\lambda_2=\lambda_3=\lambda_4=u,\quad\lambda_5=u+a. \tag{3.81}\end{equation}
相应右特征矢量矩阵（各列依次为特征矢量）为
\begin{equation}
\bm K=\begin{bmatrix}
1&1&0&0&1\\u-a&u&0&0&u+a\\v&v&1&0&v\\w&w&0&1&w\\H-ua&V^2/2&v&w&H+ua
\end{bmatrix}. \tag{3.82}
\end{equation}

% [扩译 p130]

三维 $x$ 分裂 Euler 系统有五个特征速度 $u-a,u,u,u,u+a$。中间三重特征值分别携带接触/熵波与两个切向剪切波；外侧两族为声学波。对任意网格面，只需建立包含法向和两个正交切向的局部基，就能得到相同谱结构。

% [原PDF p131]
三维旋转矩阵可分解为绕两个坐标方向的旋转，且只作用于三个动量分量。关系 $n_x\bm F+n_y\bm G+n_z\bm H=\bm T^{-1}\bm F(\bm T\bm U)$ 表明：先旋转状态、求 $x$ 向通量、再逆旋转，等价于直接求法向通量。

其逆矩阵可由直接消元获得；它把守恒变量增量投影到五个特征变量上。三维欧拉方程也具有旋转不变性。

\textbf{命题3.19}\quad 令法向的两个方向角为 $\theta_y,\theta_z$，则
\begin{multline}
\cos\theta_y\cos\theta_z\,\bm F(\bm U)
+\cos\theta_y\sin\theta_z\,\bm G(\bm U)
+\sin\theta_y\,\bm H(\bm U)\\
=\bm T^{-1}\bm F(\bm T\bm U), \tag{3.84}
\end{multline}
其中 $\bm T$ 只旋转三个动量分量，密度与总能量保持不变。证明与二维情形同样通过旋转速度分量并直接验证法向通量完成。

\subsection{三维原始变量形式}

取
\begin{equation}\bm W=(\rho,u,v,w,p)^T. \tag{3.87}\end{equation}
三维欧拉方程的原始变量形式为
\begin{align}
\rho_t+u\rho_x+v\rho_y+w\rho_z+\rho(u_x+v_y+w_z)&=0,\\
u_t+uu_x+vu_y+wu_z+p_x/\rho&=0,\\
v_t+uv_x+vv_y+wv_z+p_y/\rho&=0,\\
w_t+uw_x+vw_y+ww_z+p_z/\rho&=0,\\
p_t+up_x+vp_y+wp_z+\rho a^2(u_x+v_y+w_z)&=0. \tag{3.88}
\end{align}

% [原PDF p132]
推导步骤与一维式(3.14)--(3.16)完全相同：展开守恒形式中的各阶导数，先用质量方程消去密度导数，再由动量方程和能量方程整理。拟线性形式为
\begin{equation}\bm W_t+\bm A(\bm W)\bm W_x+\bm B(\bm W)\bm W_y+\bm C(\bm W)\bm W_z=0. \tag{3.89}\end{equation}
三个系数矩阵分别是
\begin{equation}
\bm A=\begin{bmatrix}u&\rho&0&0&0\\0&u&0&0&1/\rho\\0&0&u&0&0\\0&0&0&u&0\\0&\rho a^2&0&0&u\end{bmatrix},
\quad
\bm B=\begin{bmatrix}v&0&\rho&0&0\\0&v&0&0&0\\0&0&v&0&1/\rho\\0&0&0&v&0\\0&0&\rho a^2&0&v\end{bmatrix}, \tag{3.90--3.91}
\end{equation}
\begin{equation}
\bm C=\begin{bmatrix}w&0&0&\rho&0\\0&w&0&0&0\\0&0&w&0&0\\0&0&0&w&1/\rho\\0&0&0&\rho a^2&w\end{bmatrix}. \tag{3.92}
\end{equation}
矩阵 $\bm A$ 的特征值为
\begin{equation}\lambda_1=u-a,\quad\lambda_2=\lambda_3=\lambda_4=u,\quad\lambda_5=u+a. \tag{3.93}\end{equation}

% [扩译 p132]

原始变量形式直观显示：密度由速度散度压缩，三个速度分量由压力梯度加速，压力以 $\rho a^2\nabla\cdot\bm u$ 与膨胀耦合。它适于特征分析和边界条件推导，但遇到激波时必须回到守恒形式计算通量，因为乘积展开在分布意义下不再等价。

% [原PDF p133]
方向分裂的三维黎曼解仍有两条非线性外波，但中间特征值 $u$ 的重数为三。它们携带熵/密度跳跃及两个切向速度跳跃，并以同一速度重合传播。因而压力和法向速度跨中间波连续，切向速度却不必连续。

一组相应右特征矢量为
\begin{equation}
\bm K^{(1)}=\begin{bmatrix}\rho\\-a\\0\\0\\\rho a^2\end{bmatrix},\quad
\bm K^{(2)}=\begin{bmatrix}1\\0\\v\\w\\0\end{bmatrix},\quad
\bm K^{(3)}=\begin{bmatrix}\rho\\0\\1\\w\\0\end{bmatrix},\quad
\bm K^{(4)}=\begin{bmatrix}\rho\\0\\v\\1\\0\end{bmatrix},\quad
\bm K^{(5)}=\begin{bmatrix}\rho\\a\\0\\0\\\rho a^2\end{bmatrix}. \tag{3.94}
\end{equation}

\subsection{方向分裂的三维黎曼问题}

当前多数迎风型多维欧拉数值方法都需要求解方向分裂黎曼问题。沿 $x$ 方向分裂的三维问题为
\begin{equation}
\bm U_t+\bm F(\bm U)_x=0,qquad
\bm U(x,0)=\begin{cases}\bm U_L,&x<0,\\\bm U_R,&x>0,
\end{cases} \tag{3.95}
\end{equation}
其中
\begin{equation}
\bm U=(\rho,\rho u,\rho v,\rho w,E)^T,\quad
\bm F=(\rho u,\rho u^2+p,\rho uv,\rho uw,u(E+p))^T. \tag{3.96}
\end{equation}
其相似解结构几乎与一维情况相同。星形区中，压力和法向速度 $u$ 跨中间波保持不变。由于特征值 $u$ 为三重根，还出现两条分别使切向速度 $v,w$ 跳变的剪切波；相应第三、第四特征场线性退化。第一、第五场真正非线性，仍对应激波或稀疏波。由(3.94)可知，无论跨越第一、第五族稀疏波，$v,w$ 均不改变。

% [原PDF p134]
这一结论对激波同样成立。考虑第五族右行激波，以速度 $S$ 传播。在激波静止参考系中，除一维关系(3.42)--(3.44)外，还有切向动量跃迁关系
\begin{align}
\rho_*(u_*-S)&=\rho_R(u_R-S), \tag{3.97}\\
\rho_*(u_*-S)v_*&=\rho_R(u_R-S)v_R, \tag{3.98}\\
\rho_*(u_*-S)w_*&=\rho_R(u_R-S)w_R. \tag{3.99}
\end{align}
用(3.97)约去质量通量，立即得到 $v_*=v_R$、$w_*=w_R$。左行激波完全类似。因此，不论第一、第五族非线性波是激波还是稀疏波，切向速度都跨波保持不变。

\textbf{图3.5}\quad 三维方向分裂黎曼解的结构：星形区压力和法向速度唯一，密度可跨接触间断跳变，两个切向速度分别跨剪切波跳变。

所以，求三维分裂黎曼问题在本质上与求一维黎曼问题相同。额外变量 $v,w$ 的解非常简单：它们只在相应剪切波处从左值 $v_L,w_L$ 跳到右值 $v_R,w_R$。某些近似黎曼求解器会错误地模拟这种本来很简单的切向速度行为。

\section{守恒形式与非守恒形式的比较}

本节首先强调：即使假定解光滑，所谓守恒与非守恒表述也并不唯一。必须仔细审查一个数学上的守恒形式究竟保存什么量，以及由它导出的“守恒”是否具有物理意义。其次，在存在激波时，只在数学形式上守恒、却不对应正确物理守恒量的方程会给出错误的激波速度，从而产生错误解。

% [原PDF p135]
以一维浅水方程为例，正确的质量和动量守恒律是
\begin{equation}
\begin{bmatrix}\phi\\\phi u\end{bmatrix}_t+
\begin{bmatrix}\phi u\\\phi u^2+\tfrac12\phi^2\end{bmatrix}_x=0. \tag{3.100}
\end{equation}
在光滑解假设下展开导数，得到原始变量形式
\begin{align}
\phi_t+u\phi_x+\phi u_x&=0, \tag{3.101}\\
u_t+uu_x+\phi_x&=0. \tag{3.102}
\end{align}
将第二式改写成
\begin{equation}u_t+(\tfrac12u^2+\phi)_x=0, \tag{3.103}\end{equation}
便会得到另一个表面上的“守恒形式”
\begin{equation}
\begin{bmatrix}\phi\\u\end{bmatrix}_t+
\begin{bmatrix}\phi u\\\tfrac12u^2+\phi\end{bmatrix}_x=0. \tag{3.104}
\end{equation}
数学上它确实是一组守恒律，但第二式声称质点速度守恒，物理上并无意义。只有在解处处光滑时，(3.100)与(3.104)才等价；一旦出现激波，两者给出的解便不同。

\textbf{命题3.22}\quad 设右行激波前状态为 $(\phi_R,u_R)$。正确方程(3.100)给出的激波速度为
\begin{equation}
S=u_R+Q/\phi_R,qquad Q=\left[\tfrac12(\phi_*+\phi_R)\phi_*\phi_R\right]^{1/2}, \tag{3.105}
\end{equation}
而错误形式(3.104)给出
\begin{equation}
\widehat S=u_R+\widehat Q/\phi_R,qquad
\widehat Q=\left[\frac{2}{\phi_*+\phi_R}\right]^{1/2}\phi_*\phi_R. \tag{3.106}
\end{equation}
证明留作练习，可直接应用第2章和本章第3.1.3节的跃迁关系。

% [扩译 p135]

浅水方程例子说明，同一组光滑解可以对应多个形式上的“守恒形式”，但它们跨间断给出不同 Rankine--Hugoniot 条件。物理守恒量必须从积分平衡定律确定，不能由链式法则随意制造。等温气体动力学也有相同警示：非守恒变量离散若未规定路径，激波速度和波后状态通常错误。

% [原PDF p136]
这一例子强调：从光滑原始变量方程积分出一个“守恒形式”并不保证它表达真实守恒量。同一准线性系统可能对应多个形式守恒律，而它们跨激波给出不同跳跃速度；只有由物理质量、动量和能量平衡导出的守恒形式才是正确弱解基础。

\textbf{注3.23}\quad 仅在 $\phi_*=\phi_R$、即激波退化时，$S$ 与 $\widehat S$ 才相等。一般有
\begin{equation}\widehat S\le S, \tag{3.107}\end{equation}
因此由错误“守恒律”得到的激波比正确浅水方程的激波传播得慢；并且类似(3.104)的形式还不唯一。

再看声速 $a$ 为常数的等温气体方程。物理正确的质量、动量守恒形式为
\begin{equation}
\begin{bmatrix}\rho\\\rho u\end{bmatrix}_t+
\begin{bmatrix}\rho u\\\rho u^2+a^2\rho\end{bmatrix}_x=0. \tag{3.108}
\end{equation}
对光滑解展开为
\begin{align}
\rho_t+u\rho_x+\rho u_x&=0, \tag{3.109}\\
u_t+uu_x+\frac{a^2}{\rho}\rho_x&=0. \tag{3.110}
\end{align}
从这组原始变量方程又可构造
\begin{equation}
\begin{bmatrix}\rho\\u\end{bmatrix}_t+
\begin{bmatrix}\rho u\\\tfrac12u^2+a^2\ln\rho\end{bmatrix}_x=0. \tag{3.111}
\end{equation}
它在数学上也是守恒律，却错误地宣称质量与速度守恒，因此在物理上无用。

\textbf{练习3.24}\quad 利用第3.1.3节关于孤立激波的结果，比较(3.108)和(3.111)的激波解：哪一种给出的激波更快？并为(3.109)--(3.110)寻找其他形式上的守恒表述。

% [原PDF p137]
\chapter{欧拉方程的黎曼问题}

Godunov 在1959年的经典论文[216]中，把 Courant、Isaacson 与 Rees [144] 的一阶迎风格式以守恒方式推广到非线性双曲守恒律系统；其中的核心正是黎曼问题的解。本章详细给出一维非定常欧拉方程对理想气体和余容气体的完整精确黎曼解，包括真空情形；同样的方法还可用于其他双曲系统。

精确黎曼解有多方面用途。它是双曲守恒律在最简单而非平凡初值下的解；初值虽简单，解却包含该方程组最基本的物理和数学特征。一般初值问题的解可视为许多局部黎曼解的非线性叠加[212]。对欧拉方程而言，黎曼问题还包括气体动力学的激波管问题。精确解是评估数值方法、检查程序早期实现正确性的宝贵基准；精确或近似的黎曼解也可局部用于 Godunov 方法及其高阶推广，这是本书最主要的用途。透彻理解精确解，也是使用、评价和发展近似黎曼求解器的基础。

即便理想气体的欧拉方程也没有封闭形式的黎曼解，甚至更简单的等熵、等温模型亦然。不过可以构造迭代算法，把解求到任意实际需要的精度。设计精确求解器时须决定：选用哪些变量和方程、迭代未知量的数目、迭代技术与初猜，以及如何处理负压力等非物理迭代值。

% [原PDF p138]
Godunov [216] 首先提出欧拉方程的精确黎曼求解器；按今天的标准，其第一种算法冗长且效率低。后来他又提出第二种求解器[218]，所用方程更简单、变量更适合计算，迭代过程也更精巧。Chorin [110] 独立改进了第一种算法；van Leer [559] 在1979年作了另一种改进，所得方法与 Godunov 的第二种求解器相似。Smoller [451] 提出了不同思路，Dutt [179] 后来给出实用实现。Gottlieb 与 Groth [222] 的理想气体算法在其比较中效率最高；Toro [498] 给出的理想气体与余容气体算法具有相近效率。Schleicher [430]、Pike [384] 后来给出的算法似乎更快。一般状态方程可参阅 Colella 与 Glaz [135] 的开创性工作，以及 Menikoff 与 Plohr [349]、Saurel、Larini 与 Loraud [429]。

本章先对理想气体详述求解步骤，再处理真空并给出三种可能真空情形的精确解，最后推广到余容气体。尤其重视如何在任意点抽样相似解，因为这既用于绘制完整解，也用于 Godunov 方法和 Glimm 随机选择法。所需背景见第1--3章。

\section{求解策略}

一维欧拉方程的黎曼问题是
\begin{equation}
\bm U_t+\bm F(\bm U)_x=0,\quad
\bm U=\begin{bmatrix}\rho\\\rho u\\E\end{bmatrix},\quad
\bm F=\begin{bmatrix}\rho u\\\rho u^2+p\\u(E+p)\end{bmatrix}, \tag{4.1}
\end{equation}
配以初值
\begin{equation}\bm U(x,0)=\begin{cases}\bm U_L,&x<0,\\\bm U_R,&x>0.
\end{cases} \tag{4.2}\end{equation}

% [原PDF p139]
关心的区域为 $-\infty<x<\infty,t>0$；实际计算常取包含原点的有限区间 $[x_L,x_R]$。求解时通常使用原始变量 $\bm W=(\rho,u,p)^T$。初始数据只有被 $x=0$ 间断分开的两个常状态 $\bm W_L$ 与 $\bm W_R$。激波管问题是其特例：两侧气体初始静止，隔膜破裂后产生近似中心化的稀疏波、接触间断和激波。黎曼问题允许 $u_L,u_R$ 非零，但波系结构不变。

材料状态方程不仅决定黎曼解的结构，也决定方程的数学性质。本章考虑理想气体
\begin{equation}e=\frac{p}{(\gamma-1)\rho}, \tag{4.3}\end{equation}
以及余容气体
\begin{equation}e=\frac{p(1-b\rho)}{(\gamma-1)\rho}, \tag{4.4}\end{equation}
其中 $\gamma$ 为定比热比，$b$ 为余容常数。

\textbf{图4.1}\quad 一维非定常欧拉方程黎曼解在 $x$--$t$ 平面中的一般结构。

% [扩译 p139]

黎曼问题是局部初值只有一个间断的 Cauchy 问题。由于 Euler 方程无固有长度尺度，其熵解在 $t>0$ 时依赖 $S=x/t$；这使二维时空问题化为沿自相似轴的波型判定。状态方程既提供声速和能量闭合，也决定压力函数、稀疏曲线以及真空条件。

% [原PDF p140]
无真空时，解中有三族波，分别与 $u-a,u,u+a$ 相关，但波速一般并不等于某一固定状态的特征速度。三族波分开四个常状态：$\bm W_L,\bm W_{*L},\bm W_{*R},\bm W_R$。左右非线性波之间的星形区被接触间断分成左右两部分。中间波总是接触间断；左右两波各自可能为激波或稀疏波，因此共有四种组合：左稀疏/右激波、左激波/右稀疏、双稀疏和双激波。若某一稀疏波恰为声速波，还会出现两种变体；它们主要在 Godunov 型格式的界面抽样中重要。

由第3章的特征分析，星形区内压力 $p_*$ 和质点速度 $u_*$ 跨接触间断保持常值，而密度取两个值 $\rho_{*L},\rho_{*R}$。因此主要未知量是 $p_*,u_*,\rho_{*L},\rho_{*R}$。下面利用 $p_*,u_*$ 的连续性，把问题化为关于 $p_*$ 的单个非线性代数方程。

\textbf{图4.2}\quad 黎曼解的四种基本波型：（a）左稀疏、接触、右激波；（b）左激波、接触、右稀疏；（c）双稀疏；（d）双激波。

% [原PDF p141]
压力函数的两段定义分别来自激波 Hugoniot 关系和稀疏波 Riemann 不变量。分支条件直接由试探压力与初始压力比较：$p>p_K$ 表示压缩激波，$p\le p_K$ 表示膨胀稀疏波。两段在 $p=p_K$ 处连续且一阶导数相接。

\section{压力与质点速度方程}

\textbf{命题4.1（$p_*$ 与 $u_*$ 的解）}\quad 理想气体黎曼问题的星形区压力 $p_*$ 是方程
\begin{equation}
f(p)\equiv f_L(p,\bm W_L)+f_R(p,\bm W_R)+\Delta u=0,qquad
\Delta u=u_R-u_L \tag{4.5}
\end{equation}
的根，其中
\begin{equation}
f_L(p)=\begin{cases}
(p-p_L)\sqrt{A_L/(p+B_L)},&p>p_L\quad(\text{激波}),\\[2pt]
\dfrac{2a_L}{\gamma-1}\left[(p/p_L)^{(\gamma-1)/(2\gamma)}-1\right],&p\le p_L\quad(\text{稀疏波}),
\end{cases} \tag{4.6}
\end{equation}
\begin{equation}
f_R(p)=\begin{cases}
(p-p_R)\sqrt{A_R/(p+B_R)},&p>p_R\quad(\text{激波}),\\[2pt]
\dfrac{2a_R}{\gamma-1}\left[(p/p_R)^{(\gamma-1)/(2\gamma)}-1\right],&p\le p_R\quad(\text{稀疏波}),
\end{cases} \tag{4.7}
\end{equation}
而
\begin{equation}
A_K=\frac{2}{(\gamma+1)\rho_K},\qquad
B_K=\frac{\gamma-1}{\gamma+1}p_K,qquad K=L,R. \tag{4.8}
\end{equation}
压力求出后，星形区质点速度为
\begin{equation}u_*=\frac12(u_L+u_R)+\frac12[f_R(p_*)-f_L(p_*)]. \tag{4.9}\end{equation}

% [扩译 p141]

压力函数的每一侧都按试探压力相对于本侧初始压力选择分支：较高压力意味着激波，较低压力意味着稀疏波。根 $p_*$ 一旦确定，同一压力和速度跨接触间断连续，而左右星区密度可不同。因而三波四状态结构最终只剩一个非线性标量未知量。

% [原PDF p142]
一旦由(4.5)求得 $p_*$，式(4.9)立即给出 $u_*$，其余未知量可由第3章的标准气体动力学关系获得。$f_L$ 描述左非线性波，把未知 $u_*$ 与左数据联系；$f_R$ 同理描述右波。两函数具体采用激波支还是稀疏支，由试探压力 $p$ 与相应初始压力比较决定。为简洁起见，后文常省略数据参数。对物理可容许数据，压力函数行为十分简单，且 $f(p)=0$ 有唯一解。

\textbf{图4.3}\quad 通过压力函数求解黎曼问题的策略：$f_L$、$f_R$ 分别把星形区质点速度同左右初态联系起来。

\subsection{左激波的函数 $f_L$}

设左波为速度 $S_L$ 的激波，激波前为 $\rho_L,u_L,p_L$，激波后为 $\rho_{*L},u_*,p_*$。在随激波运动的参考系中
\begin{equation}\hat u_L=u_L-S_L,qquad\hat u_*=u_*-S_L. \tag{4.10}\end{equation}
Rankine--Hugoniot 条件为
\begin{align}
\rho_L\hat u_L&=\rho_{*L}\hat u_*, \tag{4.11}\\
\rho_L\hat u_L^2+p_L&=\rho_{*L}\hat u_*^2+p_*, \tag{4.12}\\
\hat u_L(\hat E_L+p_L)&=\hat u_*(\hat E_{*L}+p_*). \tag{4.13}
\end{align}

% [原PDF p143]
定义质量通量
\begin{equation}Q_L=\rho_L\hat u_L=\rho_{*L}\hat u_*. \tag{4.14}\end{equation}
由动量跃迁关系，
\begin{equation}Q_L=-\frac{p_*-p_L}{\hat u_* -\hat u_L}
=-\frac{p_*-p_L}{u_*-u_L}, \tag{4.15--4.16}\end{equation}
从而
\begin{equation}u_*=u_L-\frac{p_*-p_L}{Q_L}. \tag{4.17}\end{equation}
把 $\hat u_L=Q_L/\rho_L$、$\hat u_*=Q_L/\rho_{*L}$ 代回，得到
\begin{equation}Q_L^2=-\frac{p_*-p_L}{1/\rho_{*L}-1/\rho_L}. \tag{4.18}\end{equation}
激波后的密度由
\begin{equation}
\rho_{*L}=\rho_L\frac{(\gamma-1)/(\gamma+1)+p_*/p_L}{[(\gamma-1)/(\gamma+1)]p_*/p_L+1} \tag{4.19}
\end{equation}
给出，故
\begin{equation}Q_L=\sqrt{\frac{p_*+B_L}{A_L}}, \tag{4.20}\end{equation}
代入(4.17)即有
\begin{equation}u_*=u_L-f_L(p_*),\qquad
f_L=(p_*-p_L)\sqrt{\frac{A_L}{p_*+B_L}}. \tag{4.21}
\end{equation}

% [扩译 p143]

左激波函数来自在激波参考系中依次消去星区密度和质量通量。常数 $A_L,B_L$ 只依赖左初态与 $\gamma$，故 $f_L(p)$ 可在每次压力迭代中廉价求值。符号约定使 $u_*=u_L-f_L(p_*)$：左行激波提高压力时，$f_L>0$，星区速度相对左态向右侧调整。

% [二次扩译 p143]
激波支中的平方根始终取正值，质量通量的传播方向则由外层符号体现。将推导写成 $f_K(p)$ 后，可把热力学跃迁和速度修正封装在一个连续函数中。

% [原PDF p144]
\subsection{左稀疏波的函数 $f_L$}

若左波为稀疏波，未知状态 $\bm W_{*L}$ 通过等熵关系和广义黎曼不变量与 $\bm W_L$ 相连。跨稀疏波有
\begin{equation}p=C\rho^\gamma, \tag{4.22}\end{equation}
在左初态求得 $C=p_L/\rho_L^\gamma$，因此
\begin{equation}\rho_{*L}=\rho_L(p_*/p_L)^{1/\gamma}. \tag{4.23}\end{equation}
左稀疏波的黎曼不变量给出
\begin{equation}u_L+\frac{2a_L}{\gamma-1}=u_*+\frac{2a_{*L}}{\gamma-1}. \tag{4.24}\end{equation}
由(4.23)和声速定义，
\begin{equation}a_{*L}=a_L(p_*/p_L)^{(\gamma-1)/(2\gamma)}. \tag{4.25}\end{equation}
所以
\begin{equation}
u_*=u_L-\frac{2a_L}{\gamma-1}\left[(p_*/p_L)^{(\gamma-1)/(2\gamma)}-1\right]
=u_L-f_L(p_*). \tag{4.26--4.27}
\end{equation}

% [扩译 p144]

左稀疏支由等熵关系和第一族广义黎曼不变量直接积分得到。它在 $p=p_L$ 处与激波支具有相同函数值和一阶导数，使组合压力函数连续可微；这也是 Newton 迭代跨越波型切换点时仍能稳定工作的关键。

% [原PDF p145]
左稀疏波的等熵关系先给出 $\rho_{*L}=\rho_L(p_*/p_L)^{1/\gamma}$，声速继而满足 $a_{*L}=a_L(p_*/p_L)^{(\gamma-1)/(2\gamma)}$。代入左 Riemann 不变量后得到 $u_*=u_L-f_L(p_*)$，与左激波分支具有同一统一形式。

\subsection{右激波的函数 $f_R$}

对右行激波，同样在随激波运动的参考系中应用跃迁条件。定义
\begin{equation}Q_R=\rho_R\hat u_R=\rho_{*R}\hat u_*, \tag{4.28}\end{equation}
则动量关系给出
\begin{equation}u_*=u_R+\frac{p_*-p_R}{Q_R}, \tag{4.29}\end{equation}
且
\begin{equation}Q_R=\sqrt{\frac{p_*+B_R}{A_R}}. \tag{4.30}\end{equation}
因此
\begin{equation}u_*=u_R+f_R(p_*),\qquad
f_R=(p_*-p_R)\sqrt{\frac{A_R}{p_*+B_R}}. \tag{4.31--4.32}
\end{equation}

\subsection{右稀疏波的函数 $f_R$}

跨右稀疏波有等熵关系和第三族广义黎曼不变量
\begin{equation}
u_R-\frac{2a_R}{\gamma-1}=u_*-\frac{2a_{*R}}{\gamma-1}. \tag{4.33}
\end{equation}
又有
\begin{equation}a_{*R}=a_R(p_*/p_R)^{(\gamma-1)/(2\gamma)}, \tag{4.34}\end{equation}
所以
\begin{equation}u_*=u_R+f_R(p_*),\quad
f_R=\frac{2a_R}{\gamma-1}\left[(p_*/p_R)^{(\gamma-1)/(2\gamma)}-1\right]. \tag{4.35--4.36}
\end{equation}

% [扩译 p145]

右侧关系镜像为 $u_*=u_R+f_R(p_*)$。当 $p_*>p_R$ 时采用激波质量通量，当 $p_*\le p_R$ 时采用第三族稀疏不变量。左右函数虽符号相反地进入速度关系，却都随试探压力单调增加，最终在唯一 $p_*$ 处给出同一个接触速度。

% [原PDF p146]
左、右两种关系都必须给出同一个 $u_*$，故
\[
u_L-f_L(p_*)=u_R+f_R(p_*),
\]
立即得到压力方程(4.5)。两式相加则得到速度公式(4.9)。至此命题4.1得证。

\section{压力的数值求解}

压力方程一般不能显式求根，必须迭代。下面先分析压力函数的单调性、凹凸性与根的存在唯一性，再给出实用迭代格式和初值估计。

% [扩译 p146]

令左右速度表达式相等即可消去 $u_*$，得到标量压力方程；两式相加再除以二则给出星区速度。这个化简把原来的七个未知状态量拆成“先求压力根、再代数回代”的顺序，是精确求解器以及多数近似求解器设计的核心。

% [原PDF p147]
合并左右关系得到 $f(p)=0$ 后，函数严格递增且向下凹。Newton 迭代因而通常快速，但初值过小可能越过 $p=0$；实现应同时设正压力下限和相对残差终止准则。求出 $p_*$ 后，两侧给出的 $u_*$ 理论上相同，可用其差作为迭代质量检查。

\subsection{压力函数的性态}

\textbf{命题4.3}\quad 对 $p>0$，$f(p)$ 连续且一阶连续可导，并满足
\begin{equation}f'(p)>0,qquad f''(p)<0. \tag{4.37}\end{equation}
也就是说，压力函数严格单调递增且向下凹。

\textit{证明。}由于 $f=f_L+f_R+\Delta u$，只需分别检查每个波函数。对激波支 $K=L,R$，
\begin{align}
f_K(p)&=(p-p_K)\sqrt{\frac{A_K}{p+B_K}},\\
f'_K(p)&=\sqrt{\frac{A_K}{p+B_K}}\left[1-\frac{p-p_K}{2(p+B_K)}\right], \tag{4.38}\\
f''_K(p)&=-\frac14\sqrt{\frac{A_K}{(p+B_K)^5}}\,[4B_K+3p+p_K]<0. \tag{4.39}
\end{align}
可见 $f'_K>0$、$f''_K<0$。稀疏支同样直接求导：
\begin{align}
f'_K(p)&=\frac{1}{\rho_Ka_K}(p/p_K)^{-(\gamma+1)/(2\gamma)}>0, \tag{4.40}\\
f''_K(p)&=-\frac{\gamma+1}{2\gamma p\rho_Ka_K}(p/p_K)^{-(\gamma+1)/(2\gamma)}<0. \tag{4.41}
\end{align}

% [扩译 p147]

激波支和稀疏支均满足 $f_K'(p)>0$、$f_K''(p)<0$。严格单调性保证根至多一个，向下凹性决定 Newton 切线相对于根的位置并可用于设计安全初值。实现时应按波型使用解析导数，避免数值差分在 $p=p_K$ 附近损失精度。

% [二次扩译 p147]
在分支连接点 $p=p_K$，两侧二阶导数一般不必相同，但函数与一阶导数连续已经足以保证标准 Newton 更新。严格递增还意味着残差符号能直接判断当前压力位于根的哪一侧。

% [原PDF p148]
激波支与稀疏支在 $p=p_K$ 处函数值和一阶导数相同，因而整体 $C^1$ 连续。于是 $f$ 在正压力范围严格递增、向下凹。其极限为
\begin{equation}
\lim_{p\to\infty}f(p)=+\infty,qquad
f(0)=u_R-u_L-\frac{2(a_L+a_R)}{\gamma-1}. \tag{4.42}
\end{equation}
因此正根存在且唯一的充要条件是
\begin{equation}u_R-u_L<\frac{2(a_L+a_R)}{\gamma-1}. \tag{4.43}\end{equation}
若等号成立，$p_*=0$；若左侧更大，则两股气体分离得过快，中间生成真空。该判据将在第4.6节详细讨论。

压力函数的图像会在 $p_L,p_R$ 处切换分支，却保持连续光滑。根相对左右压力的位置直接决定波型：$p_*>p_L$ 为左激波，否则为左稀疏；$p_*>p_R$ 为右激波，否则为右稀疏。

% [扩译 p148]

考察 $f(0)$ 与 $p\to\infty$ 的极限可判定根是否存在。若两侧分离速度超过两列最大稀疏波所能连接的范围，正压力根消失并生成真空。该判断必须先于压力迭代；将负迭代值简单钳成小正数并不能修复错误的波系假设。

% [原PDF p149]
双稀疏近似把两条非线性波都按等熵扇处理，能显式给出压力初值；双激波近似则在参考压力附近线性化两侧激波函数。PVRS 初值适合弱波。自适应选择依据左右压力比与预测压力所在区间，在稳健性和额外计算之间折中。

\subsection{求压力的迭代格式}

最直接的 Newton--Raphson 迭代为
\begin{equation}
p^{(i+1)}=p^{(i)}-\frac{f(p^{(i)})}{f'(p^{(i)})}. \tag{4.44}
\end{equation}
若初值 $p^{(0)}>0$ 合理，压力函数单调且凹，迭代通常快速收敛。计算时应以相对变化量
\begin{equation}
\frac{2|p^{(i+1)}-p^{(i)}|}{p^{(i+1)}+p^{(i)}}<\varepsilon \tag{4.45}
\end{equation}
作为停止准则。不得接受负压力迭代值；实用实现可在迭代步骤中检测，并用严格为正的安全初值重新开始。

两稀疏波近似可显式给出压力初猜
\begin{equation}
p_{TR}=\left[
\frac{a_L+a_R-\tfrac12(\gamma-1)(u_R-u_L)}
{a_L/p_L^z+a_R/p_R^z}
\right]^{1/z},qquad z=\frac{\gamma-1}{2\gamma}. \tag{4.46}
\end{equation}
这一表达式在真实波型为双稀疏时实际上就是精确星区压力。

% [扩译 p149]

Newton 法通常二次收敛，但初值过小或数据接近真空时仍可能越过 $p=0$。稳健实现应同时检查相对压力变化和残差，限制最大迭代次数，并在失败时报告左右状态。双稀疏估计在真实波型确为两稀疏时就是精确根，因此是低压力问题的优良起点。

% [原PDF p150]
简单的原始变量线性化估计（PVRS）是
\begin{equation}
p_{PV}=\frac12(p_L+p_R)-\frac18(u_R-u_L)(\rho_L+\rho_R)(a_L+a_R), \tag{4.47}
\end{equation}
并取 $p_{PV}\leftarrow\max(0,p_{PV})$ 作为估计值。若压力比不大、左右状态接近，它尤其准确。两激波近似则取
\begin{equation}
p_{TS}=\frac{g_L(p_{PV})p_L+g_R(p_{PV})p_R-(u_R-u_L)}{g_L(p_{PV})+g_R(p_{PV})},
\quad g_K(p)=\sqrt{\frac{A_K}{p+B_K}}. \tag{4.48}
\end{equation}

实用混合初猜先计算 $p_{PV}$，并以 $p_{\min}=\min(p_L,p_R)$、$p_{\max}=\max(p_L,p_R)$、$q=p_{\max}/p_{\min}$ 判断：若 $q$ 不大且 $p_{PV}$ 位于两压力之间，采用 $p_{PV}$；若 $p_{PV}<p_{\min}$，采用 $p_{TR}$；否则采用 $p_{TS}$。这样既利用近邻状态的廉价线性估计，又在强稀疏或强激波时选用合适非线性估计。

% [扩译 p150]

PVRS、双稀疏和双激波三个估计分别适合弱波、强膨胀和强压缩。自适应策略以左右压力比及 $p_{\rm PV}$ 是否落在 $[p_{\min},p_{\max}]$ 内决定分支，既降低迭代次数，又避免在强激波问题上用过低初值。它们只提供初猜，最终仍求解同一精确压力方程。

% [二次扩译 p150]
双激波估计把激波函数中的平方根系数冻结在预测压力处，从而得到一个可显式求解的线性速度匹配式。对于极强压力比，它通常比 PVRS 更接近真实根。

% [原PDF p151]
精确算法在迭代前首先检查真空产生条件 $u_R-u_L\ge2(a_L+a_R)/(\gamma-1)$。若成立，星形区正压力根不存在，应进入真空解分支；继续对 $f(p)$ 作 Newton 迭代只会产生负压力或不收敛，掩盖真实波系。

\subsection{数值试验}

下面用五组标准黎曼数据检验压力迭代。这些算例涵盖弱、强稀疏波，强激波以及左右镜像的冲击问题。取 $\gamma=1.4$，左右状态依次为：
\begin{center}
\begin{tabular}{c|ccc|ccc|c}
试验&$\rho_L$&$u_L$&$p_L$&$\rho_R$&$u_R$&$p_R$&输出时刻\\\hline
1&1.0&0&1.0&0.125&0&0.1&0.25\\
2&1.0&-2&0.4&1.0&2&0.4&0.15\\
3&1.0&0&1000&1.0&0&0.01&0.012\\
4&1.0&0&0.01&1.0&0&100&0.035\\
5&5.99924&19.5975&460.894&5.99242&-6.19633&46.0950&0.035
\end{tabular}
\end{center}
试验1是经典 Sod 激波管；试验2产生两个强稀疏波但尚未形成真空；试验3、4分别是向右和向左的强激波问题；试验5包含两个相向运动的激波。精确解的密度、速度、压力和单位质量内能曲线可用来核对波的位置、跳跃幅度以及稀疏扇内的连续变化。

% [扩译 p151]

五组算例覆盖 Sod 管、对称双稀疏、左高压强激波、其镜像问题以及高 Mach 碰撞。表中不仅比较 $p_*,u_*$，还记录迭代次数和不同初值的表现。镜像算例应给出对称的压力与相反速度，是检查左右分支符号是否一致的有效单元测试。

% [原PDF p152]
迭代的性能取决于初始压力。双稀疏估计对稀疏主导问题非常好，两激波估计适于强压缩问题，PVRS 则对数值方法中最常见的相邻近状态既便宜又准确。第4.9节程序的子程序 \texttt{STARTE} 正是把 $p_{PV},p_{TR},p_{TS}$ 组合成混合方案。在实际 Godunov 计算中，绝大多数界面黎曼问题的左右状态相近，因此简单的 $p_{PV}$ 已足够精确。

\textbf{图4.6}\quad 试验1在 $t=0.25$ 时的密度、速度、压力和单位质量内能精确曲线。

% [原PDF p153]
\textbf{图4.7}\quad 试验2在 $t=0.15$ 时的精确解。两股气体背向运动，产生左右稀疏扇，中心压力和密度很低但仍为正。

\textbf{图4.8}\quad 试验3在 $t=0.012$ 时的精确解。左侧高压驱动强右行激波，并伴随左稀疏波和接触间断。

求得 $p_*$ 后，$u_*$ 由(4.9)给出，星形区密度由各波对应关系计算。五个试验的关键精确值为
\begin{center}
\begin{tabular}{c|rrrr}
试验&$p_*$&$u_*$&$\rho_{*L}$&$\rho_{*R}$\\\hline
1&0.30313&0.92745&0.42632&0.26557\\
2&0.00189&0&0.02185&0.02185\\
3&460.894&19.5975&0.57506&5.99924\\
4&46.0950&-6.19633&5.99242&0.57511\\
5&1691.64&8.68975&14.2823&31.0426
\end{tabular}
\end{center}
\textbf{表4.3}\quad 五组试验的精确星形区压力、速度与左右密度，可用于程序初步验证。

% [原PDF p154]
\section{完整解}

到此已有算法求出星形区的 $p_*,u_*$，但尚须求 $\rho_{*L},\rho_{*R}$ 并完全确定左右波。把 $p_*$ 分别与 $p_L,p_R$ 比较即可判断外波类型。若为激波，只需求波后密度和激波速度；若为稀疏波，还需波后密度、波头与波尾速度以及扇形区内部解。

先考虑接触间断左侧。

\subsubsection{左激波}

条件 $p_*>p_L$ 表明左波为激波。密度为
\begin{equation}
\rho_{*L}=\rho_L
\frac{p_*/p_L+(\gamma-1)/(\gamma+1)}{[(\gamma-1)/(\gamma+1)]p_*/p_L+1}. \tag{4.50}
\end{equation}

% [扩译 p154]

完整解的回代顺序必须固定：先由 $p_*$ 与本侧初始压力比较确定波型，再求星区密度、波速及必要的扇内状态。接触两侧共享 $p_*,u_*$，但密度由各自的激波或等熵关系独立计算，因此一般有 $\rho_{*L}\ne\rho_{*R}$。

% [原PDF p155]
由相对速度和质量通量，激波速度为
\begin{equation}S_L=u_L-Q_L/\rho_L, \tag{4.51}\end{equation}
即
\begin{equation}
S_L=u_L-a_L\sqrt{\frac{\gamma+1}{2\gamma}\frac{p_*}{p_L}+\frac{\gamma-1}{2\gamma}}. \tag{4.52}
\end{equation}
于是接触间断左侧的解完全确定：$S=x/t<S_L$ 时为左初态，$S_L<S<u_*$ 时为激波后星区态。

\subsubsection{左稀疏波}

条件 $p_*\le p_L$ 表明左波为稀疏波。由等熵关系，
\begin{align}
\rho_{*L}&=\rho_L(p_*/p_L)^{1/\gamma}, \tag{4.53}\\
a_{*L}&=a_L(p_*/p_L)^{(\gamma-1)/(2\gamma)}. \tag{4.54}
\end{align}
稀疏波头、尾速度分别为
\begin{equation}S_{HL}=u_L-a_L,qquad S_{TL}=u_*-a_{*L}. \tag{4.55}\end{equation}

% [扩译 p155]

左激波只有一条传播射线，抽样点按 $S_L$ 分成初态和左星态；左稀疏波则有波头、连续扇和波尾三个区域。波头以左初态特征速度传播，波尾以星区特征速度传播，二者次序同时体现了稀疏波特征发散的熵条件。

% [原PDF p156]
抽样程序先判断相似速度 $S=(x-x_0)/t$ 位于接触波左侧还是右侧，再按 $p_*$ 与初始压力比较确定激波或稀疏波。激波只需比较单一波速；稀疏波还要依次比较头、尾速度，并在扇内用自相似解析式求状态。

对扇内任一点，经过原点的特征射线斜率满足 $S=x/t=u-a$，同时左黎曼不变量满足
\[
u_L+\frac{2a_L}{\gamma-1}=u+\frac{2a}{\gamma-1}.
\]
联立求 $u,a$，再用声速定义和等熵律，得到左稀疏扇内解
\begin{equation}
\bm W_{L\mathrm{fan}}(S)=
\begin{bmatrix}
\rho_L\left[\dfrac{2}{\gamma+1}+\dfrac{\gamma-1}{(\gamma+1)a_L}(u_L-S)\right]^{2/(\gamma-1)}\\[6pt]
\dfrac{2}{\gamma+1}\left[a_L+\dfrac{\gamma-1}{2}u_L+S\right]\\[6pt]
p_L\left[\dfrac{2}{\gamma+1}+\dfrac{\gamma-1}{(\gamma+1)a_L}(u_L-S)\right]^{2\gamma/(\gamma-1)}
\end{bmatrix}. \tag{4.56}
\end{equation}

\textbf{图4.12}\quad 接触间断左侧的抽样区域：（a）左激波；（b）左稀疏波，图中标出波头和波尾特征。

% [原PDF p157]
五组标准激波管覆盖双稀疏、强激波、近真空和不同波型组合。验证不应只看最终曲线，还要核对 $p_*,u_*$、各波速度、星形区密度和采样分支；这样才能区分压力迭代错误、波型判断错误与绘图采样错误。

接着考虑接触间断右侧。

\subsubsection{右激波}

若 $p_*>p_R$，则右波为激波。其星区密度和波速为
\begin{align}
\rho_{*R}&=\rho_R
\frac{p_*/p_R+(\gamma-1)/(\gamma+1)}{[(\gamma-1)/(\gamma+1)]p_*/p_R+1}, \tag{4.57}\\
S_R&=u_R+Q_R/\rho_R \tag{4.58}\\
&=u_R+a_R\sqrt{\frac{\gamma+1}{2\gamma}\frac{p_*}{p_R}+\frac{\gamma-1}{2\gamma}}. \tag{4.59}
\end{align}

\subsubsection{右稀疏波}

若 $p_*\le p_R$，则
\begin{align}
\rho_{*R}&=\rho_R(p_*/p_R)^{1/\gamma}, \tag{4.60}\\
a_{*R}&=a_R(p_*/p_R)^{(\gamma-1)/(2\gamma)}, \tag{4.61}\\
S_{HR}&=u_R+a_R,qquad S_{TR}=u_*+a_{*R}. \tag{4.62}
\end{align}

% [扩译 p157]

右侧抽样是左侧的镜像，但比较号方向相反。右激波速度大于右态质点速度；右稀疏波头为 $u_R+a_R$，波尾为 $u_*+a_{*R}$。实现时将左右分支写成清晰的嵌套判断，可避免把右扇内部误用左扇公式。

% [二次扩译 p157]
右星区密度仍由压力比决定，而接触间断只允许密度和熵跳跃；压力与法向速度必须连续。这一约束是检验返回状态是否自洽的首要条件。

% [原PDF p158]
右稀疏扇内解为
\begin{equation}
\bm W_{R\mathrm{fan}}(S)=
\begin{bmatrix}
\rho_R\left[\dfrac{2}{\gamma+1}-\dfrac{\gamma-1}{(\gamma+1)a_R}(u_R-S)\right]^{2/(\gamma-1)}\\[6pt]
\dfrac{2}{\gamma+1}\left[-a_R+\dfrac{\gamma-1}{2}u_R+S\right]\\[6pt]
p_R\left[\dfrac{2}{\gamma+1}-\dfrac{\gamma-1}{(\gamma+1)a_R}(u_R-S)\right]^{2\gamma/(\gamma-1)}
\end{bmatrix}. \tag{4.63}
\end{equation}
\textbf{图4.13}\quad 接触间断右侧的抽样区域：（a）右激波；（b）右稀疏波。

% [扩译 p158]

右稀疏扇公式由 $S=u+a$ 与第三族黎曼不变量联立得到。扇内 $u,a,\rho,p$ 都连续随 $S$ 变化，并在 $S_{TR}$ 与右星态、在 $S_{HR}$ 与右初态精确衔接。检查这两个端点是验证解析实现和指数符号的直接方法。

% [二次扩译 p158]
扇内声速括号必须保持非负；接近真空边界时它趋于零，密度与压力以相应幂次衰减。若因舍入得到负括号，应视为抽样区间判断错误而不是直接取绝对值。

% [原PDF p159]
物质与真空不能由有限强度激波直接相连，因为将真空状态代入质量和动量跳跃条件会迫使物质侧压力同时为零。合法连接是稀疏波，其密度、压力和声速在扇尾连续趋零；真空边界速度正是该尾特征速度。

\section{解的抽样}

上述结果可在任意 $(x,t)$ 求出完整精确解。由于解是相似解，只需使用自相似速度 $S=x/t$ 抽样。若给定时刻 $t$，则解剖面只随 $x$ 变化。首先判断抽样点位于接触间断哪一侧。

\subsection{接触间断左侧：$S\le u_*$}

若左波是激波，
\begin{equation}
\bm W(S)=\begin{cases}
\bm W_L,&S\le S_L,\\
\bm W_{*L}^{\mathrm{sho}},&S_L<S\le u_*.
\end{cases} \tag{4.64}
\end{equation}
若左波是稀疏波，
\begin{equation}
\bm W(S)=\begin{cases}
\bm W_L,&S\le S_{HL},\\
\bm W_{L\mathrm{fan}}(S),&S_{HL}<S<S_{TL},\\
\bm W_{*L}^{\mathrm{fan}},&S_{TL}\le S\le u_*.
\end{cases} \tag{4.65}
\end{equation}
其中各星区状态和扇内状态均由上一节公式给出。算法上依次比较 $S$ 与 $S_L$，或与 $S_{HL},S_{TL}$，即可决定返回哪个状态。

% [扩译 p159]

自相似抽样无需显式构造整条剖面：给定一个 $S$，先由 $S\le u_*$ 决定接触间断哪一侧，再判断该侧波型，最后同一个或两个波速比较。有限体积 Godunov 通量只关心界面 $S=0$，因此精确求解器可在得到 $p_*,u_*$ 后只返回这一点的状态。

% [二次扩译 p159]
每个返回状态都应在对应区间端点连续匹配；只有接触间断和激波允许跳跃。逐个测试所有分支，比仅绘制一条典型激波管曲线更容易发现边界比较号错误。

% [原PDF p160]
当左右物质同时向外快速分离时，两列稀疏波之间形成真空。左扇尾速度为 $u_L+2a_L/(\gamma-1)$，右扇尾速度为 $u_R-2a_R/(\gamma-1)$；前者不大于后者时，中间区无物质，二者相等对应临界真空。

\subsection{接触间断右侧：$S\ge u_*$}

若右波是激波，
\begin{equation}
\bm W(S)=\begin{cases}
\bm W_{*R}^{\mathrm{sho}},&u_*\le S<S_R,\\
\bm W_R,&S\ge S_R.
\end{cases} \tag{4.66}
\end{equation}
若右波是稀疏波，
\begin{equation}
\bm W(S)=\begin{cases}
\bm W_{*R}^{\mathrm{fan}},&u_*\le S\le S_{TR},\\
\bm W_{R\mathrm{fan}}(S),&S_{TR}<S<S_{HR},\\
\bm W_R,&S\ge S_{HR}.
\end{cases} \tag{4.67}
\end{equation}
\textbf{图4.14}\quad 任意 $S=x/t$ 处抽样精确黎曼解的流程图：先按 $S\lessgtr u_*$ 选左右支，再判断该侧是激波还是稀疏波，并与相应波速比较。

% [扩译 p160]

抽样返回原始变量后，再按状态方程组装守恒变量和物理通量。所有分支都应保证 $\rho>0,p>0$；若事先判定为真空问题，则必须使用下一节专门的自由边界结构，而不是让常规星区公式在零密度处分母失效。

% [原PDF p161]
初始一侧已经是真空时，只剩从物质侧发出的一列稀疏波和一个自由边界，另一条非线性波不存在。自由边界上 $a=\rho=p=0$，速度由相应 Riemann 不变量确定；数值抽样须把真空区显式返回为零密度、零压力。

\section{含真空的黎曼问题}

流动物质与真空相邻的情形在自由膨胀、空化和高速气体动力学中很重要。这里以 $\rho=0$ 表示真空，并相应取单位体积总能量 $E=0$。连续介质模型在真空内部本身不再有效，本节求解的是物质区域与真空区域之间自由边界的运动。需要区分三类初值：左侧物质、右侧真空；左侧真空、右侧物质；以及左右均有物质、但相互背离的膨胀波在中间生成真空。

含真空解不再具有通常的两个星形区。若仍把式（4.5）的压力方程或非真空迭代器生硬用于该问题，算法实际上假定了不存在的波系结构；以很小的正密度和压力冒充真空，也会改变局部波结构。尤其要注意，激波不能直接与真空相邻。设物质态 $\bm W_L=(\rho_L,u_L,p_L)^T$ 通过速度为 $S$ 的间断连接真空态 $\bm W_0=(0,u_0,p_0)^T$，Rankine--Hugoniot 条件给出
\begin{align}
 \rho_Lu_L-\rho_0u_0&=S(\rho_L-\rho_0), \tag{4.68}\\
 \rho_Lu_L^2+p_L-(\rho_0u_0^2+p_0)&=S(\rho_Lu_L-\rho_0u_0), \tag{4.69}\\
 u_L(E_L+p_L)-u_0(E_0+p_0)&=S(E_L-E_0). \tag{4.70}
\end{align}

% [原PDF p162]
代入 $\rho_0=E_0=0$ 并假设 $u_0$ 有限，三式只能给出
\begin{equation}
 u_L=u_0=S,\qquad p_L=p_0. \tag{4.71}
\end{equation}
因此有限强度激波不可能一侧为真空；接触间断却可以成为物质边界。边界速度 $u_0$ 是连接物质态与真空态的波系中最大质点速度，称为\textbf{逸出速度}。为确定边界状态，采用等熵型状态方程
\begin{equation}p=p(\rho), \tag{4.72}\end{equation}
并要求
\begin{equation}p(0)=0,\quad p'(0)=0,\quad p'(\rho)>0,\quad p''(\rho)>0. \tag{4.73}\end{equation}
于是记 $\bm W_0=(0,u_0,0)^T$。可能出现的基本波只剩接触间断和稀疏波。

\subsection{情形一：右侧为真空}
初值为
\begin{equation}
 \bm W(x,0)=\begin{cases}\bm W_L\ne\bm W_0,&x<0,\\\bm W_0,&x>0.\end{cases} \tag{4.74}
\end{equation}
它可看成装有均匀气体的激波管右端活塞突然向右运动：活塞速度超过逸出速度时，活塞与气体之间形成真空；等于逸出速度对应刚好开始空化。

% [二次扩译 p162]
接触间断与真空前沿重合时，物质一侧的压力和密度趋于零，但前沿速度仍有限。所谓真空速度不是可在真空中定义的流体速度，而是自由边界的传播速度。状态方程条件 $p'(0)=0$ 正好使边界声速消失。

% [原PDF p163]
活塞以大于临界速度离开均匀气体时，活塞与气体之间形成真空。解中的接触前沿与稀疏波尾重合，跨该前沿压力和速度连续，而密度降至零。真空侧没有介质承载右行声波，这正是通常三波结构少一条外波的物理原因。

解由一列左行稀疏波和与其尾部重合的真空边界组成，通常黎曼图中的右非线性波消失，因为真空中没有传播该波的介质。沿左广义黎曼不变量有
\begin{equation}
 u_0+\frac{2a_0}{\gamma-1}=u_L+\frac{2a_L}{\gamma-1}. \tag{4.75}
\end{equation}
上述状态方程条件使边界处 $a_0=0$，故自由边界速度为
\begin{equation}
 S_{*L}=u_0=u_L+\frac{2a_L}{\gamma-1}. \tag{4.76}
\end{equation}
完整自相似解为
\begin{equation}
 \bm W_{L0}(x,t)=\begin{cases}
 \bm W_L,&x/t\le u_L-a_L,\\
 \bm W_{L\mathrm{fan}},&u_L-a_L<x/t<S_{*L},\\
 \bm W_0,&x/t\ge S_{*L},
 \end{cases} \tag{4.77}
\end{equation}
其中扇内状态使用式（4.56）。真空前沿正好是稀疏扇的尾特征，压力、密度和声速沿扇连续降到零。

% [二次扩译 p163]
图4.15中，左稀疏波头从原点以 $u_L-a_L$ 传播，扇内质点不断向真空加速；到尾部时声速降为零，物质速度达到最大值 $u_0$。不存在右行声波，因为其右侧没有物质承载压力扰动。

% [原PDF p164]
\subsection{情形二：左侧为真空}
反向情形的初值为
\begin{equation}
 \bm W(x,0)=\begin{cases}\bm W_0,&x<0,\\\bm W_R\ne\bm W_0,&x>0.\end{cases} \tag{4.78}
\end{equation}
此时左非线性波缺失，真空边界同右稀疏波尾部重合。由右广义黎曼不变量得到
\begin{equation}S_{*R}=u_R-\frac{2a_R}{\gamma-1}, \tag{4.79}\end{equation}
相应解为
\begin{equation}
 \bm W_{R0}(x,t)=\begin{cases}
 \bm W_0,&x/t\le S_{*R},\\
 \bm W_{R\mathrm{fan}},&S_{*R}<x/t<u_R+a_R,\\
 \bm W_R,&x/t\ge u_R+a_R.
 \end{cases} \tag{4.80}
\end{equation}

\subsection{情形三：生成真空}
若 $\bm W_L$ 与 $\bm W_R$ 都非真空，但两侧物质以足够大的相对速度背离，则两列稀疏波之间生成有限宽真空区。两条真空边界的速度分别为 $S_{*L}$ 和 $S_{*R}$。

% [二次扩译 p164]
左真空情形可由空间反射 $x\mapsto-x$ 从右真空情形得到。此时最小质点速度出现在左自由边界，右初态保持到右稀疏波头之外。公式中的负号反映物质向左侧真空膨胀。

% [原PDF p165]
完整解可组合前两种单边真空解写成
\begin{equation}
 \bm W(x,t)=\begin{cases}
 \bm W_{L0}(x,t),&x/t\le S_{*L},\\
 \bm W_0,&S_{*L}<x/t<S_{*R},\\
 \bm W_{R0}(x,t),&x/t\ge S_{*R}.
 \end{cases} \tag{4.81}
\end{equation}
要使中间真空区存在，必须满足 $S_{*L}\le S_{*R}$，即
\begin{equation}
 (\Delta u)_{\mathrm{crit}}
 =\frac{2a_L}{\gamma-1}+\frac{2a_R}{\gamma-1}
 \le u_R-u_L. \tag{4.82}
\end{equation}
这与压力正性条件（4.40）完全一致。Godunov 型程序应在进入非线性压力迭代以前检验该条件。近似黎曼求解器若使用自身的真空判据，其临界条件一般不会与精确式（4.82）相同。

\section{余体积气体的黎曼问题}
余体积状态方程是理想气体关系的一个小修正。声速为
\begin{equation}
 a=\sqrt{\frac{p/\rho^2-e_\rho}{e_p}}
  =\sqrt{\frac{\gamma p}{\rho(1-b\rho)}}. \tag{4.83}
\end{equation}
波系结构与理想气体完全相同；主要区别是稀疏扇内部不再有封闭解析式，需要增加一次标量迭代。

% [二次扩译 p165]
生成真空时，两稀疏扇尾部之间的距离按 $(S_{*R}-S_{*L})t$ 线性增长。临界等号只在一个自相似射线上出现零密度状态；严格不等式才产生具有有限宽度的真空区。

% [原PDF p166]
余体积状态方程要求 $1-b\rho>0$，并把声速改为 $a^2=\gamma p/[\rho(1-b\rho)]$。黎曼解的三波拓扑与理想气体相同，激波压力函数也保持相似形式；差别主要进入等熵关系、系数 $A_K$ 以及稀疏扇内部的非显式求解。

\subsection{星形区压力和质点速度}
星形区压力仍是单调方程
\begin{equation}
 f(p)=f_L(p,\bm W_L)+f_R(p,\bm W_R)+u_R-u_L=0. \tag{4.84}
\end{equation}
对 $K=L,R$，压力函数统一写成
\begin{equation}
 f_K(p)=\begin{cases}
 (p-p_K)\sqrt{A_K/(p+B_K)},&p>p_K,\\[2pt]
 \dfrac{2a_K(1-b\rho_K)}{\gamma-1}
 \left[(p/p_K)^{(\gamma-1)/(2\gamma)}-1\right],&p\le p_K,
 \end{cases} \tag{4.85--4.86}
\end{equation}
其中
\begin{equation}
 A_K=\frac{2(1-b\rho_K)}{(\gamma+1)\rho_K},\qquad
 B_K=\frac{\gamma-1}{\gamma+1}p_K. \tag{4.87}
\end{equation}
求出 $p_*$ 后，星形区速度取两侧表达式的平均：
\begin{equation}
 u_*=\frac12(u_L+u_R)+\frac12\bigl[f_R(p_*)-f_L(p_*)\bigr]. \tag{4.88}
\end{equation}

% [二次扩译 p166]
余体积参数 $b$ 表示分子不可占据的体积修正，并要求 $1-b\rho>0$。当 $b\to0$ 时，声速、$A_K$ 以及稀疏支全部连续退化为理想气体公式，因此可用这一极限检查实现。
与理想气体相比，给定压力和密度时因可用体积减小，声速会增大；数值程序必须把 $b$ 的影响一致地用于初态声速、压力函数和星区恢复。

% [原PDF p167]
左激波的质量通量 $Q_L$ 仍把压力跳跃与速度跳跃联系起来，但 $A_L$ 含因子 $1-b\rho_L$。左稀疏波不变量则成为 $u+2a(1-b\rho)/(\gamma-1)$。令 $b\to0$，所有公式连续退化为理想气体结果，可作为实现回归测试。

若左波为激波，Rankine--Hugoniot 条件给出星区密度
\begin{equation}
 \rho_{*L}=\rho_L
 \frac{p_*/p_L+\dfrac{\gamma-1}{\gamma+1}
       \dfrac{\gamma-1+2b\rho_L}{\gamma+1}}
      {\dfrac{\gamma-1}{\gamma+1}p_*/p_L+
       \dfrac{\gamma+1-2b\rho_L}{\gamma+1}}, \tag{4.89}
\end{equation}
质量通量与速度关系为
\begin{equation}
 Q_L=-\frac{p_*-p_L}{u_*-u_L}
     =\sqrt{\frac{p_*+B_L}{A_L}},\qquad
 u_*=u_L-f_L(p_*). \tag{4.90--4.92}
\end{equation}
若左波为稀疏波，则沿左黎曼不变量
\begin{equation}
 u_L+\frac{2a_L}{\gamma-1}(1-b\rho_L)
 =u_*+\frac{2a_{*L}}{\gamma-1}(1-b\rho_{*L}), \tag{4.93}
\end{equation}
并结合余体积气体等熵关系求 $\rho_{*L}$。最终仍得到 $u_*=u_L-f_L(p_*)$，但此处 $f_L$ 应采用稀疏支。

% [二次扩译 p167]
左激波密度式同时体现压力压缩和余体积上限。质量通量 $Q_L$ 的平方根形式仍与理想气体相同，只是 $A_L$ 含因子 $1-b\rho_L$。在稀疏支中，这一因子进入黎曼不变量与等熵关系。
所有根号量都应保持正值，且星区密度不能越过 $1/b$。

% [原PDF p168]
右侧推导与左侧镜像对称。右激波满足
\begin{equation}
 Q_R=\frac{p_*-p_R}{u_*-u_R}=\sqrt{\frac{p_*+B_R}{A_R}},
 \qquad u_*=u_R+f_R(p_*), \tag{4.97--4.99}
\end{equation}
右稀疏波则利用
\begin{equation}
 u_R-\frac{2a_R}{\gamma-1}(1-b\rho_R)
 =u_*-\frac{2a_{*R}}{\gamma-1}(1-b\rho_{*R}), \tag{4.100}
\end{equation}
与等熵关系确定星区密度，得到同一速度公式（4.102）。因此，理想气体精确求解器的总体控制结构可以保留，只需替换声速、压力函数、星区密度以及扇内抽样公式。

% [二次扩译 p168]
右侧推导中的 $a_{*R}$ 必须由右星区密度和余体积状态方程计算，不能直接沿用理想气体压力幂式。只有星形区边界量可以通过等熵关系联立求得；扇内一般点还需下一小节的密度迭代。

% [原PDF p169]
余体积压力迭代仍利用 $f'(p)>0$ 的单调性。双稀疏初值中，每侧声速贡献乘以 $1-b\rho_K$；PVRS 形式表面不变，但声速已含余体积修正。迭代和后续状态恢复都必须检查 $\rho<1/b$，否则状态方程出现奇点。

\subsection{压力方程的数值求解}
式（4.84）仍可用 Newton--Raphson 法。双稀疏初值为
\begin{equation}
 p_{\mathrm{TR}}=\left[
 \frac{(1-b\rho_L)a_L+(1-b\rho_R)a_R-\tfrac12(\gamma-1)(u_R-u_L)}
 {(1-b\rho_L)a_L/p_L^z+(1-b\rho_R)a_R/p_R^z}
 \right]^{1/z},\quad z=\frac{\gamma-1}{2\gamma}. \tag{4.103}
\end{equation}
若两侧确为稀疏波，此值就是精确解。原始变量线性化给出的 PVRS 初值为
\begin{equation}
 p_0=\max(\mathrm{TOL},p_{\mathrm{PV}}),\qquad
 p_{\mathrm{PV}}=\frac12(p_L+p_R)-\frac18(u_R-u_L)(\rho_L+\rho_R)(a_L+a_R). \tag{4.104}
\end{equation}
其形式与理想气体相同，余体积修正通过 $a_L,a_R$ 进入。求得 $p_*,u_*$ 后，激波速度由质量通量得到；稀疏波边界仍为 $u\mp a$。

% [二次扩译 p169]
压力迭代本身仍是一维求根，因此计算开销主要增加在需要抽样稀疏扇内部时。若最终只在 $S=0$ 取界面状态，而且零射线落在常状态区，便不需要执行额外的扇内 Newton 迭代。
双激波初猜也可沿用理想气体的自适应框架，只需使用余体积定义的 $A_K,B_K$。因此求解器控制结构不变，差异集中在热力学函数层。

% [原PDF p170]
左稀疏扇内以射线速度 $S=x/t=u-a$ 和广义 Riemann 不变量建立方程。由于 $a$ 同时依赖 $p,\rho$ 及 $1-b\rho$，不能像理想气体那样直接消元成幂函数；作者把问题归结为单个密度方程 $Z_L(\rho)=0$ 再作 Newton 迭代。

\subsection{稀疏扇内部的解}
余体积气体不能像理想气体那样直接写出扇内状态。对左扇内射线 $S=x/t$，有
\begin{equation}S=u-a, \tag{4.105}\end{equation}
并由左黎曼不变量得到
\begin{equation}
 a\left[1+\frac{2}{\gamma-1}(1-b\rho)\right]
 =I_L(u_L,a_L)-S. \tag{4.106}
\end{equation}
等熵关系为
\begin{equation}
 p=p_L\left(\frac{1-b\rho_L}{\rho_L}\right)^\gamma
       \left(\frac{\rho}{1-b\rho}\right)^\gamma. \tag{4.107}
\end{equation}
消去 $u,p,a$ 后，只需对密度求解
\begin{equation}
 Z_L(\rho)=\rho^{\gamma-1}(\gamma+1-2b\rho)^2
 -\beta_L(1-b\rho)^{\gamma+1}=0, \tag{4.108}
\end{equation}
其中
\begin{equation}
 \beta_L=\frac{\{(\gamma-1)[I_L(u_L,a_L)-S]\}^2}
 {\gamma p_L[(1-b\rho_L)/\rho_L]^\gamma}. \tag{4.109}
\end{equation}

% [二次扩译 p170]
求解 $Z_L(\rho)=0$ 时，物理解必须同时满足 $0<\rho<1/b$ 并落在左初态与左星态密度之间。利用这个区间作括区或步长保护，比允许 Newton 跨越余体积奇点更加可靠。
每次迭代都应由当前密度重新计算 $Z_L$ 及解析导数；收敛后再恢复压力、声速和速度，避免在未收敛状态上混合更新多个变量。

% [原PDF p171]
求解 $Z_K(\rho)=0$ 时，可用相应星形区与初始态密度作为物理区间，并保证 $0<\rho<1/b$。密度确定后，压力、声速和速度按固定顺序恢复。右扇公式只需交换左右并改变特征号，代码中适合共享一个带方向符号的例程。

Newton 迭代所需导数为
\begin{equation}
 Z_L'(\rho)=(\gamma+1)\left[b\beta_L(1-b\rho)^\gamma
 +(\gamma+1-2b\rho)(\gamma-1-2b\rho)\rho^{\gamma-2}\right]. \tag{4.110}
\end{equation}
密度收敛后，由式（4.107）求压力，由式（4.83）求声速，再由式（4.105）求速度。右稀疏扇完全类似：
\begin{align}
 Z_R(\rho)&=\rho^{\gamma-1}(\gamma+1-2b\rho)^2
 -\beta_R(1-b\rho)^{\gamma+1}=0, \tag{4.111}\\
 \beta_R&=\frac{\{(\gamma-1)[I_R(u_R,a_R)-S]\}^2}
 {\gamma p_R[(1-b\rho_R)/\rho_R]^\gamma}. \tag{4.112}
\end{align}
这样即可在任意 $(x,t)$ 抽样余体积气体的精确黎曼解。练习要求读者据此画出完整抽样流程图。

\section{方向分裂的多维情形}
为了在多维 Godunov 方法中使用黎曼解，把 $x$ 方向理解为网格面的法向。三维 $x$ 分裂系统为
\begin{equation}
 \bm U_t+\bm F(\bm U)_x=0,\quad
 \bm U=(\rho,\rho u,\rho v,\rho w,E)^T,\quad
 \bm F=(\rho u,\rho u^2+p,\rho uv,\rho uw,u(E+p))^T. \tag{4.113}
\end{equation}

% [二次扩译 p171]
右扇函数只是把左黎曼不变量换成右黎曼不变量并改变特征方向。密度根确定后，所有其他量都由显式关系恢复，因此额外迭代仍是标量问题，不需要联立求解 $\rho,u,p$。

% [原PDF p172]
两个切向速度虽都跨接触波跳跃，却不影响 $p_*$、$u_*$ 和外侧声学波的求解。它们相当于以法向速度被动输运的变量。近似黎曼求解器若不能分辨接触波，会同时抹宽密度、切向速度和材料标记，这是多维剪切层耗散的重要来源。

初值仍为左右常状态：
\begin{equation}\bm U(x,0)=\begin{cases}\bm U_L,&x<0,\\\bm U_R,&x>0.\end{cases} \tag{4.114}\end{equation}
密度、法向速度 $u$ 和压力的解与真正的一维问题完全相同。特征值 $u$ 的三重性使三个中间波重合，因此表面上仍只有左非线性波、接触波和右非线性波。两个切向速度 $v,w$ 只跨接触波从左初值跳到右初值：
\begin{equation}
 q(x,t)=\begin{cases}q_L,&x/t<u_*,\\q_R,&x/t>u_*,\end{cases}
 \qquad q=v\ \text{或}\ w. \tag{4.115}
\end{equation}
这说明多维面通量计算应先把速度旋转到法向--切向坐标，求解一维法向黎曼问题，再把动量通量旋转回全局坐标。

% [二次扩译 p172]
三个重合中间波具有相同传播速度，却对应三个独立特征方向；在一维剖面上它们无法按空间分开。其综合效果是压力与法向速度连续，而密度以及两个切向速度都允许跨中间波跳跃。

% [原PDF p173]
程序清单把精确求解分为三层：\texttt{STARPU} 只求星形区压力速度，\texttt{SAMPLE} 只按给定 $S$ 抽样波结构，主程序负责网格点循环与输出。这样的职责分离便于分别验证非线性迭代和波型逻辑，也便于在 Godunov 通量中复用中心抽样。

任意被流体被动输运的标量 $q$ 满足
\begin{equation}q_t+uq_x+vq_y+wq_z=0. \tag{4.116}\end{equation}
与连续方程合并可得守恒形式
\begin{equation}
 (\rho q)_t+(u\rho q)_x+(v\rho q)_y+(w\rho q)_z=0. \tag{4.117}
\end{equation}
所以浓度、反应进度变量和流体界面标识等被动量也按式（4.115）跨接触波跳跃。精确解在这一点十分简单，但某些近似黎曼求解器会过度耗散接触波，甚至给出错误的切向速度或组分分布。该方法还可推广到真实气体、可压缩液体和稳态超声速 Euler 方程；后续各章将系统构造近似求解器。

\section{精确黎曼求解器的 FORTRAN 程序}
本节程序求一维理想气体 Euler 方程的精确黎曼解。输入文件为 \texttt{exact.ini}，输出为 \texttt{exact.out}。主程序先调用 \texttt{STARPU} 迭代求 $p_*,u_*$，其中初值由 \texttt{GUESSP} 给出；随后对每个输出位置计算 $S=(x-x_0)/t$，调用 \texttt{SAMPLE} 返回 $\rho,u,p$。Godunov 一阶方法只需在 $S=0$ 抽样，而 Glimm 方法在随机射线上抽样。

% [二次扩译 p173]
对多种被动标量，可为每个 $q_j$ 增加守恒量 $\rho q_j$。精确接触波不混合左右组分；若近似通量把中间波合并成单个 HLL 平均态，组分界面会被额外扩散，这也是 HLLC 恢复接触波的动机之一。

% [原PDF p174]
程序首先声明网格数、气体常数、左右状态和输出参数，并把 $\gamma$ 相关的八个常数集中预计算。这一做法既减少循环中的重复运算，也使各子程序共享完全相同的指数和系数。原程序的逻辑骨架为：
\begin{lstlisting}[language=Fortran,caption={精确黎曼求解器主程序的结构}]
read input: domain, diaphragm, cells, gamma, time
read left state (DL,UL,PL) and right state (DR,UR,PR)
compute G1,...,G8 and sound speeds CL,CR
test vacuum-generation condition
call STARPU(PM,UM)
for every cell centre:
    S = (XPOS-DIAPH)/TIMEOUT
    call SAMPLE(PM,UM,S,DS,US,PS)
    write XPOS,DS,US,PS,internal_energy
end for
\end{lstlisting}
这里 \texttt{PM,UM} 分别代表星形区压力与速度，\texttt{DS,US,PS} 是某条抽样射线上的状态。

% [原PDF p175]
输入依次包括计算域长度、初始间断位置、单元数、比热比、输出时刻、左右两侧的密度--速度--压力，以及压力归一化常数。程序随后设置
\begin{align*}
G_1&=(\gamma-1)/(2\gamma),&G_2&=(\gamma+1)/(2\gamma),\\
G_3&=2\gamma/(\gamma-1),&G_4&=2/(\gamma-1),\\
G_5&=2/(\gamma+1),&G_6&=(\gamma-1)/(\gamma+1),\\
G_7&=(\gamma-1)/2,&G_8&=\gamma-1,
\end{align*}
以及 $C_L=\sqrt{\gamma P_L/D_L}$、$C_R=\sqrt{\gamma P_R/D_R}$。进入压力迭代前先以 $G_4(C_L+C_R)\le U_R-U_L$ 检测是否生成真空，避免 Newton 法在不存在星形区的结构上失败。

% [二次扩译 p175]
八个 $G_i$ 只是常用组合的缓存，并不引入新的物理参数。将它们集中定义还能避免不同子程序对指数和激波系数采用不一致的舍入表达式。输入压力归一化常数只用于输出单位换算，不参与黎曼解。

% [原PDF p176]
若数据生成真空，原程序输出诊断并停止；否则调用 \texttt{STARPU}。输出网格采用单元中心 $x_i=(i-1/2)\Delta x$，每点的自相似速度为 $S_i=(x_i-x_0)/t_{\rm out}$。\texttt{SAMPLE} 返回的密度、速度、压力与单位质量内能写入文件，其中内能由 $p/[\rho(\gamma-1)]$ 计算。这个循环也展示了精确解如何作为有限体积程序的验证基准：数值解与精确解必须在相同单元中心和相同时刻比较。

% [原PDF p177]
\texttt{STARPU} 先调用 \texttt{GUESSP} 取得 $p_0$，再执行最多 20 次 Newton--Raphson 迭代。每次分别调用 \texttt{PREFUN} 得到 $f_L,f_L',f_R,f_R'$，更新
\begin{equation}
 p^{(n+1)}=p^{(n)}-
 \frac{f_L(p^{(n)})+f_R(p^{(n)})+u_R-u_L}
 {f_L'(p^{(n)})+f_R'(p^{(n)})}.
\end{equation}
相对变化量采用 $2|p^{(n+1)}-p^{(n)}|/(p^{(n+1)}+p^{(n)})$；小于 $10^{-6}$ 时停止。若更新值为负，程序把它置为正容差以维持下一次函数求值；若达到最大迭代次数则报告不收敛。

% [二次扩译 p177]
若达到最大迭代次数，后续抽样不应继续使用最后一次压力；正确做法是把未收敛状态作为显式错误返回。原清单打印发散信息，工程实现还应附带左右初态、残差与迭代次数，便于定位病态数据。

% [原PDF p178]
\texttt{GUESSP} 中先计算 $p_{\rm PV}$、$p_{\min}$、$p_{\max}$ 和压力比。若满足弱波判据就采用 PVRS；若预测压力低于两侧压力则用双稀疏公式；否则用双激波近似。初值选择只影响迭代成本，不应改变最终精确根。

压力收敛后，\texttt{STARPU} 用
\begin{equation}u_*=\tfrac12(u_L+u_R+f_R-f_L)\end{equation}
计算星形区速度并打印结果。\texttt{GUESSP} 实现第 9 章介绍的自适应初值策略：先计算原始变量线性化压力 $p_{\rm PV}$，再依据左右压力比和 $p_{\rm PV}$ 所在区间选择 PVRS、双稀疏或双激波近似。阈值 \texttt{QUSER=2} 表示当 $p_{\max}/p_{\min}\le2$ 且 $p_{\rm PV}$ 位于两侧压力之间时直接采用 PVRS。

% [原PDF p179]
若 $p_{\rm PV}<p_{\min}$，初值选择双稀疏近似：先由左右声速和压力比估计 $u_*$，再分别计算两侧等熵压力并取平均。否则选择双激波近似，以 $p_{\rm PV}$ 代入两侧激波函数系数后解线性组合。该分支逻辑不是在近似最终解，而只是为精确压力方程提供靠近根的起点；最终精度仍由 Newton 迭代容差决定。

% [原PDF p180]
\texttt{PREFUN} 对单侧状态计算 $f_K(p)$ 及导数。若 $p\le p_K$，使用稀疏支
\begin{equation}
 f_K=\frac{2a_K}{\gamma-1}\left[(p/p_K)^{(\gamma-1)/(2\gamma)}-1\right],\quad
 f_K'=\frac{1}{\rho_Ka_K}(p/p_K)^{-(\gamma+1)/(2\gamma)}.
\end{equation}
若 $p>p_K$，使用激波支
\begin{equation}
 f_K=(p-p_K)\sqrt{\frac{A_K}{p+B_K}},\quad
 f_K'=\left[1-\frac{p-p_K}{2(p+B_K)}\right]\sqrt{\frac{A_K}{p+B_K}}.
\end{equation}
同一子程序通过输入的 $\rho_K,p_K,a_K$ 同时服务左右两侧。

% [二次扩译 p180]
激波支导数中的第二项来自平方根系数对压力的导数，遗漏它会破坏 Newton 法的收敛速度。稀疏支导数在 $p\to0$ 时增大，这解释了近真空问题对初值和正性判断格外敏感。

% [原PDF p181]
左稀疏扇抽样的三段判断必须按空间顺序执行：$S\le S_{HL}$ 是左初态，$S\ge S_{TL}$ 是左星形态，中间才是扇内解析态。若头尾比较号写反，会在正常状态上计算分数幂并产生看似随机的密度或压力错误。

\texttt{SAMPLE} 首先比较 $S$ 与 $u_*$。若 $S\le u_*$，抽样点位于接触间断左侧；再以 $p_*\le p_L$ 判断左波是否为稀疏波。对左稀疏波，程序依次比较 $S$ 与波头 $u_L-a_L$、波尾 $u_*-a_{*L}$：波头外返回左初值，波尾内返回左星形区，二者之间按式（4.56）求扇内状态。这一判定顺序避免了在不属于稀疏扇的位置计算分数幂。

% [原PDF p182]
若左波为激波，程序计算
\begin{equation}
 S_L=u_L-a_L\sqrt{\frac{\gamma+1}{2\gamma}\frac{p_*}{p_L}
 +\frac{\gamma-1}{2\gamma}}.
\end{equation}
$S\le S_L$ 时返回左初值，否则返回左星形区，后者密度采用激波压缩比。若 $S>u_*$，程序转入接触间断右侧；先判断是否为右激波，并用对应 $S_R$ 区分右初值与右星形区。整个抽样过程只包含有限次比较和代数运算，适合作为数值通量中的局部基准。

% [原PDF p183]
右侧抽样与左侧镜像，但比较号相反：从右向左依次是右初态、右稀疏扇、右星形态。清单末尾把扇内注释误写成 left fan，公式中的 $-a_R$ 和 $u_R-S$ 明确表明它实际是右扇；中文说明已按控制逻辑纠正。

右稀疏波分支与左侧镜像：波头为 $S_{HR}=u_R+a_R$，波尾为 $S_{TR}=u_*+a_{*R}$。若 $S\ge S_{HR}$ 返回右初值；若 $S\le S_{TR}$ 返回右星形区；处于扇内时使用式（4.63）计算 $u,a,\rho,p$。原清单注释中将该区域误写为 ``inside left fan''，从控制流和公式可知实际应为右稀疏扇。

% [原PDF p184]
最后两个 \texttt{ENDIF} 依次闭合右波类型和接触间断左右侧判断，随后子程序返回。至此精确求解器完整实现了：真空预检、星形区压力迭代、自适应压力初值、左右激波/稀疏波函数，以及任意自相似速度上的状态抽样。程序虽采用 FORTRAN 77 写法，但算法可直接移植到 C++、Fortran 90 或 GPU 核函数；移植时应保留波型判断顺序、压力正性诊断和收敛失败报告，不能以静默钳位掩盖错误的波系结构。
% ===== 译文分片 3（已合并）=====

% [原PDF p185]
\chapter{数值方法的若干概念}

这里假定读者已经熟悉偏微分方程数值方法的基本概念，特别是截断误差、精度阶、一致性、修正方程、稳定性和收敛性。一般背景可参阅 Smith、Anderson 等、Mitchell 与 Griffiths、Roache、Richtmyer 与 Morton、Hoffmann 以及 Fletcher 的教材；与本书主题最接近的参考书还包括 Sod、Holt、Hirsch、LeVeque、Godlewski 与 Raviart、Kröner 和 Thomas。

本章专门为后续双曲守恒律的高分辨率迎风格式与中心格式提供准备。首要目标是在简单模型上说明 Godunov 方法，使它推广到非线性守恒律系统时更容易理解。绝对初学者可先学习有限差分、稳定性与线性双曲方程的基础，再回到本章。

\section{离散化：入门概念}
数值方法把由偏微分方程表示的连续问题替换为有限个离散量。首先用网格把定义域划分成有限个点或有限体积，然后在网格上离散方程。有限差分法把未知量看作网格点上的点值，有限体积法则把它看作有限体积上的平均值。本书主要采用后一种观点，但为介绍基本概念，也使用有限差分模型。

% [原PDF p186]
考虑常系数线性平流方程
\begin{equation}
u_t+a u_x=0,\qquad u=u(x,t),\quad a\ne0,
\end{equation}
其中 $a$ 为恒定传播速度。方程含时间导数 $u_t$ 和空间导数 $u_x$。

\subsection{导数的近似}
设 $g(x)$ 足够光滑。Taylor 定理给出
\begin{equation}
g(x_0+\Delta x)=g(x_0)+\sum_{k=1}^{\infty}\frac{(\Delta x)^k}{k!}g^{(k)}(x_0).\tag{5.1}
\end{equation}
截断此级数既可近似邻点函数值，也可近似导数。对等距三点 $x_0-\Delta x,x_0,x_0+\Delta x$，保留到二阶项有
\begin{align}
g(x_0+\Delta x)&=g(x_0)+\Delta x g'(x_0)+\frac{\Delta x^2}{2}g''(x_0)+O(\Delta x^3),\tag{5.2}\\
g(x_0-\Delta x)&=g(x_0)-\Delta x g'(x_0)+\frac{\Delta x^2}{2}g''(x_0)+O(\Delta x^3).\tag{5.3}
\end{align}
因此前向差分为
\begin{equation}
g'(x_0)=\frac{g(x_0+\Delta x)-g(x_0)}{\Delta x}+O(\Delta x).\tag{5.4}
\end{equation}

Taylor 展开中的余项阶数决定近似精度。前向式只舍弃了以 $\Delta x^2$ 起始的项，除以 $\Delta x$ 后误差为 $O(\Delta x)$；把正、负两式相减时偶次项抵消，除以 $2\Delta x$ 后首个未保留项为 $O(\Delta x^2)$，所以中心差分自然提高一阶。这个“借助对称消去低阶误差”的机制在后续高阶重构中反复出现。另一方面，较高形式精度并不保证波动问题的稳定性，导数近似还必须与时间推进和特征方向整体分析。

% [原PDF p187]
相应的后向与中心差分为
\begin{align}
g'(x_0)&=\frac{g(x_0)-g(x_0-\Delta x)}{\Delta x}+O(\Delta x),\tag{5.5}\\
g'(x_0)&=\frac{g(x_0+\Delta x)-g(x_0-\Delta x)}{2\Delta x}+O(\Delta x^2).\tag{5.6}
\end{align}
前两式是一阶单边近似，后一式是二阶中心近似，也可理解为前向和后向差分的算术平均。对波传播问题，单边还是中心并非次要细节：迎风格式本质上使用信息来向一侧的单边差分，而选哪一侧必须由传播方向决定。

\begin{sourceblock}
\sourcecrop{187}{52}{314}{388}{487}{0.25}
\captionof{figure}{函数一阶导数的后向、中心和前向差分示意。切线表示真实导数，离散割线表示三类近似。}
\end{sourceblock}

\subsection{偏微分方程的有限差分近似}
在线性平流方程的矩形区域 $[0,L]\times[0,T]$ 上，初边值问题写成
\begin{equation}
u_t+a u_x=0,\qquad u(x,0)=u_0(x),\qquad
u(0,t)=u_l(t),\quad u(L,t)=u_r(t).\tag{5.7}
\end{equation}
求解就是由初值出发，在边界条件约束下推进 $u(x,t)$。边界条件与具体物理密切相关，第 6 章将对物理系统作更有意义的讨论。

图 5.1 的几何意义是：前向与后向割线只取目标点一侧的信息，中心割线同时取两侧信息。对静态函数求导，中心式通常更精确；但对含明确因果传播方向的双曲方程，未来状态只能依赖上游数据。若格式从下游取值，就会把尚未传播到目标点的信息错误地带入更新。因而“高阶中心”不能单独作为选择格式的理由，必须先尊重偏微分方程的依赖域。

% [原PDF p188]
取规则网格，空间和时间步长分别为 $\Delta x$、$\Delta t$。若 $[0,L]$ 用 $M$ 个等距点离散，则
\begin{equation}\Delta x=\frac{L}{M-1}.\tag{5.8}\end{equation}
网格点为 $(x_i,t^n)=(i\Delta x,n\Delta t)$，离散值记作
\[
u_i^n\equiv u(i\Delta x,n\Delta t)=u(x_i,t^n).
\]
上标 $n$ 是时间层，下标 $i$ 是空间位置。相同符号也用于表示精确网格值的数值近似，必要时再作区分。已知第 $n$ 层所有数据后，目标是求第 $n+1$ 层；端点值由边界条件确定。

\begin{sourceblock}
\sourcecrop{188}{52}{160}{389}{410}{0.34}
\captionof{figure}{$x$--$t$ 平面上的规则有限差分网格；$\Delta x$ 与 $\Delta t$ 分别为空间和时间间距。}
\end{sourceblock}

这里的初边值问题还揭示边界条件数量与传播方向的联系。常速 $a>0$ 时信息从 $x=0$ 进入、从 $x=L$ 离开，连续问题原则上只需要左端入流条件；同时指定两端可能造成过定约束。此处暂把 $u_l(t),u_r(t)$ 都列出，是为了说明离散端点必须由某种规则赋值，真正的物理边界处理将在系统的特征结构下决定。空间索引与时间索引的分工也很重要：显式格式仅由已知第 $n$ 层计算第 $n+1$ 层，不需求解全局代数方程。

% [原PDF p189]
在 $(x_i,t^n)$ 处用一阶时间前向差分和二阶空间中心差分，
\begin{equation}
u_t\approx\frac{u_i^{n+1}-u_i^n}{\Delta t},\qquad
u_x\approx\frac{u_{i+1}^n-u_{i-1}^n}{2\Delta x}.\tag{5.9--5.10}
\end{equation}
代入方程得显式格式
\begin{equation}
u_i^{n+1}=u_i^n-\frac{c}{2}(u_{i+1}^n-u_{i-1}^n),
\qquad c=\frac{a\Delta t}{\Delta x},\tag{5.12--5.13}
\end{equation}
其中 $c$ 是 Courant 数，即物理波速 $a$ 与网格速度 $\Delta x/\Delta t$ 之比。这个构造看似合理，却无条件不稳定。令试探解 $u_i^n=A^n e^{\mathrm{i}i\theta}$，放大因子为 $G=1-\mathrm{i}c\sin\theta$，故
\[
|G|^2=1+c^2\sin^2\theta\ge1.
\]
除平凡模态外误差都会增长，说明“时间前向、空间中心”的组合不能用于纯平流。

放大因子分析把任意网格误差分解成 Fourier 模态。相角 $\theta=P\Delta x$，$P$ 为波数，波长为 $2\pi/P$；$\mathrm{i}=\sqrt{-1}$。对每个模态，格式一步只把振幅乘以 $G$。稳定要求所有可解析相角均有 $|G|\le1$。本例除 $\sin\theta=0$ 的特殊模态外都有 $|G|>1$，而且缩小时间步只能让增长率靠近 1，不能使它降到 1 以下，所以称“无条件不稳定”。这与截断误差的一致性并不矛盾：一致但不稳定的格式不会收敛。

% [原PDF p190]
\section{若干差分格式}
\subsection{一阶迎风格式}
补救方法是把中心空间差分换成一阶单边差分：
\begin{equation}
u_x\approx\frac{u_i^n-u_{i-1}^n}{\Delta x},\qquad
u_x\approx\frac{u_{i+1}^n-u_i^n}{\Delta x}.\tag{5.14--5.15}
\end{equation}
只有与波速符号相符的一项有用。若 $a>0$，信息从左向右传播，应取后向差分，得到 Courant--Isaacson--Rees（CIR）格式
\begin{equation}
u_i^{n+1}=u_i^n-c(u_i^n-u_{i-1}^n).\tag{5.16}
\end{equation}
von Neumann 分析给出条件稳定范围
\begin{equation}0\le c\le1.\tag{5.17}\end{equation}
在给定物理波速与网格尺度后，这一条件限制时间步长。所谓“迎风”，就是空间差分取自信息流入方向；$a>0$ 时上游在左，$a<0$ 时上游在右。

对式（5.16），放大因子满足
\[
|G|^2=(1-c)^2+2c(1-c)\cos\theta+c^2.
\]
当 $0\le c\le1$ 时可写成凸组合，振幅不增长；若 $c$ 超出范围，就存在相角使误差放大。给定 $a$ 后，网格 $\Delta x$ 多由分辨率和计算资源选择，剩余可调量是 $\Delta t$。因此 CFL 限制不是任意经验规则，而是离散因果传播速度必须覆盖物理传播速度的条件。它也说明加密空间网格时，显式时间步必须按比例缩小。

% [原PDF p191]
若 $a>0$ 却使用右侧数据，则下风格式
\begin{equation}u_i^{n+1}=u_i^n-c(u_{i+1}^n-u_i^n)\tag{5.18}\end{equation}
无条件不稳定。为统一处理正负波速，定义
\begin{equation}
a^+=\max(a,0)=\frac{a+|a|}{2},\qquad
a^-=\min(a,0)=\frac{a-|a|}{2},\tag{5.19}
\end{equation}
以及 $c^{\pm}=a^{\pm}\Delta t/\Delta x$。统一迎风格式为
\begin{equation}
u_i^{n+1}=u_i^n-c^+(u_i^n-u_{i-1}^n)-c^-(u_{i+1}^n-u_i^n),\tag{5.21}
\end{equation}
稳定条件是 $0\le |c|\le1$。这种正、负部分分裂的思想可直接推广到守恒律系统。

三点模板的底边构成格式的数值依赖域，而偏微分方程的真实依赖域由通过目标点的特征线决定。稳定格式必须让数值依赖域包含真实依赖域。

符号分裂满足 $a^+\ge0$、$a^-\le0$。当 $a\ge0$ 时 $a^+=a,a^-=0$，统一式自动删除右侧差分；当 $a\le0$ 时恰好相反。绝对值关系 $a=a^++a^-$ 与 $|a|=a^+-a^-$ 后续会变成矩阵正负分裂。格式模板由 $(i-1,n),(i,n),(i+1,n)$ 和 $(i,n+1)$ 组成，但某一速度符号下只使用真正上游的两个旧层点。这个模板随波向而改变，正是迎风法区别于固定中心模板之处。

% [原PDF p192]
对 $a>0$，由 $(x_i,t^{n+1})$ 回溯的特征线在第 $n$ 层落于 $x_p\in[x_{i-1},x_i]$，精确解满足
\begin{equation}u(x_i,t^{n+1})=u(x_p,t^n).\tag{5.23}\end{equation}
网格只给出 $u_{i-1}^n,u_i^n$，于是用两点线性插值估计 $u(x_p,t^n)$。因 $x_i-x_p=a\Delta t=c\Delta x$，插值结果恰为
\[
\widetilde u(x_p)=u_i^n-c(u_i^n-u_{i-1}^n),
\]
即 CIR 格式。故它也可理解为沿特征回溯后的线性插值。

\begin{sourceblock}
\sourcecrop{192}{58}{328}{386}{545}{0.30}
\captionof{figure}{$a>0$ 时 CIR 格式的模板。数值依赖域包住特征线给出的真实依赖点。}
\end{sourceblock}

当 $c=1$ 时，特征足点正好落在 $x_{i-1}$，CIR 更新成为 $u_i^{n+1}=u_{i-1}^n$，精确搬运一个网格间距；当 $0<c<1$ 时，足点落在两个节点之间，线性插值的两个权重为 $c$ 和 $1-c$，均非负。若 $c>1$，足点跑到数值模板之外，所谓插值就变成外推，至少一个权重为负并引起不稳定。这给出 CFL 条件最直观的几何解释：真实依赖点必须位于格式能看到的旧层区间内。

% [原PDF p193]
稳定条件也可写成 $0\le a/(\Delta x/\Delta t)\le1$，即网格速度不小于物理传播速度。截断误差分析表明 CIR 格式对时间、空间均为一阶，其修正方程为
\begin{equation}
q_t+a q_x=\alpha_{\mathrm{CIR}}q_{xx},\qquad
\alpha_{\mathrm{CIR}}=\frac12\Delta x\,a(1-|c|).\tag{5.24--5.25}
\end{equation}
$\alpha_{\mathrm{CIR}}q_{xx}$ 是人工黏性；它使间断被涂抹、极值被削平。只有 $|c|=1$ 或 $a=0$ 时该黏性消失，前者对多波速非线性系统通常无法处处满足，后者则说明静止间断可以精确保持。

更一般的线性格式可写作
\begin{equation}
u_i^{n+1}=H(u_{i-l_L}^n,\ldots,u_{i+l_R}^n)
=\sum_{k=-l_L}^{l_R}b_k u_{i+k}^n.\tag{5.26}
\end{equation}
CIR 格式的系数为 $b_{-1}=c^+$、$b_0=1-|c|$、$b_1=-c^-$，稳定范围内全部非负。

修正方程不是原偏微分方程的简单近似，而是数值解以更高阶精度实际满足的连续方程。其额外二阶导数项像物理扩散一样平滑梯度，所以称数值黏性。$\alpha_{\mathrm{CIR}}$ 与 $\Delta x$ 成正比，网格细化时消失，因而保持一致性；但有限网格上它不可避免。非线性系统同时含多种波速，不可能让所有局部 Courant 数都等于 1，所以一阶迎风法在实际计算中总会涂抹某些波。随机选择法是少数不靠这种扩散稳定间断的一阶方法。

% [原PDF p194]
\textbf{定义 5.1（单调格式）}\quad 若式（5.26）的所有 $b_k\ge0$，则称格式单调；等价地，更新函数满足 $\partial H/\partial u_k^n\ge0$。单调格式构成现代守恒律格式的基础，但其精度至多一阶；高阶扩展需要第 13、14、16 章中的限制器与 TVD 理论。

\subsection{其他著名格式}
Lax--Friedrichs 格式把不稳定格式中的 $u_i^n$ 换成邻点平均，得到
\begin{equation}
u_i^{n+1}=\frac12(1+c)u_{i-1}^n+\frac12(1-c)u_{i+1}^n.\tag{5.30}
\end{equation}
它在 $|c|\le1$ 时稳定且一阶准确，修正方程的人工黏性系数为
\begin{equation}
\alpha_{\mathrm{LF}}=\frac{a\Delta x}{2c}(1-c^2).\tag{5.31}
\end{equation}
对 $0<c\le1$，$\alpha_{\mathrm{LF}}/\alpha_{\mathrm{CIR}}=(1+c)/c\ge2$，所以它比 CIR 更耗散。其三个系数为 $(1+c)/2,0,(1-c)/2$，稳定范围内非负，因而也是单调格式。

单调性的第二个定义把 $H$ 看成旧层各数据的函数：任何输入增加都不能使输出反向减少。对线性格式，这恰好等价于所有系数非负。再结合一致性所需的系数和为 1，更新就是旧值的凸组合，不会产生超出邻域最小值和最大值的新极值。这个性质提供强稳健性，也解释 Godunov 阶数障碍：线性单调格式不能超过一阶。现代高分辨率方法必须成为非线性格式，让局部系数随解变化，在光滑区达到高阶、在间断附近退回单调行为。

% [原PDF p195]
Lax--Wendroff 格式在空间、时间均为二阶。可把空间导数看成迎风与下风一阶差分的加权组合，权重 $\beta_1=(1+c)/2$、$\beta_2=(1-c)/2$，结果为
\begin{equation}
u_i^{n+1}=\frac12c(1+c)u_{i-1}^n+(1-c^2)u_i^n
-\frac12c(1-c)u_{i+1}^n.\tag{5.34}
\end{equation}
尽管构造中只用了形式上一阶的导数近似，整体格式却是二阶，并在 $|c|\le1$ 稳定。这提醒我们：格式的精度阶不能仅由各偏导近似的局部阶数直接判断。

由于 $b_1=-c(1-c)/2$ 在一般情况下为负，Lax--Wendroff 不是单调格式，在间断附近会产生非物理振荡。Warming--Beam 二阶迎风格式在 $a>0$ 时为
\begin{equation}
u_i^{n+1}=\frac12c(c-1)u_{i-2}^n+c(2-c)u_{i-1}^n
+\frac12(c-1)(c-2)u_i^n,\tag{5.35}
\end{equation}
它完全使用模板中心及其左侧点。

Lax--Wendroff 的权重中包含负系数，因此它不是旧层值的凸组合。光滑解中，负权重帮助抵消一阶误差，正是二阶精度的来源；间断处 Taylor 展开失效，同样的负权重却制造新极值。Warming--Beam 完全单边，适合明确迎风方向，并因更宽的上游模板允许 $|c|\le2$；它仍有负系数，所以不能仅凭“迎风”二字推断无振荡。精度、稳定性、单调性是三个不同概念，必须分别检验。

% [原PDF p196]
Warming--Beam 格式的稳定范围扩大到 $0\le|c|\le2$，允许更大的时间步，但并不单调。Fromm 格式对 $a>0$ 为
\begin{equation}
u_i^{n+1}=-\frac14(1-c)c u_{i-2}^n+\frac14(5-c)c u_{i-1}^n
+\frac14(1-c)(4+c)u_i^n-\frac14(1-c)c u_{i+1}^n,\tag{5.37}
\end{equation}
其稳定条件为 $|c|\le1$，同样非单调。Lax--Wendroff、Warming--Beam 与 Fromm 的主导修正项均具有色散形式
\begin{equation}q_t+a q_x=\beta q_{xxx}.\tag{5.38}\end{equation}
这解释了它们在陡峭梯度处的相位误差与振荡。

\section{守恒方法}
含激波的解对控制方程形式和离散方法都提出严格要求。非守恒变量形式跨越激波会给出错误的跃迁条件、激波强度、速度和位置。Hou 与 Le Floch 证明，存在激波时非守恒格式一般不能收敛到正确解；Lax--Wendroff 定理则指出，收敛的守恒格式必收敛到守恒律的弱解。因此，除激波拟合等特殊技术外，激波捕捉应以守恒格式为基础。

Fromm 格式可看成 Lax--Wendroff 与 Warming--Beam 的某种组合，在相位误差上取得折中。二阶格式的修正方程以三阶导数为主，属于色散而非扩散：不同波数以不同数值相速传播，波形因而出现前导或滞后振铃。对平滑周期波，这种误差可能比一阶削峰小；对激波，振铃会产生负密度、负压力等不可接受状态。因此后续高阶格式需要限制器，只在光滑区域保留二阶修正。

% [原PDF p197]
\subsection{基本定义}
考虑标量守恒律
\begin{equation}u_t+f(u)_x=0.\tag{5.39}\end{equation}
为容纳弱解，使用控制体积分式。对任意矩形 $[x_1,x_2]\times[t_1,t_2]$，
\begin{equation}
\int_{x_1}^{x_2}u(x,t_2)\,dx=\int_{x_1}^{x_2}u(x,t_1)\,dx
+\int_{t_1}^{t_2}f(u(x_1,t))\,dt-\int_{t_1}^{t_2}f(u(x_2,t))\,dt.\tag{5.41}
\end{equation}

\textbf{定义 5.2（守恒格式）}\quad 守恒格式具有通量差形式
\begin{equation}
u_i^{n+1}=u_i^n+\frac{\Delta t}{\Delta x}
\left(f_{i-1/2}-f_{i+1/2}\right),\tag{5.42}
\end{equation}
其中 $f_{i+1/2}$ 是依赖有限模板的数值通量。它必须满足一致性
\begin{equation}f_{i+1/2}(v,\ldots,v)=f(v).\tag{5.44}\end{equation}
练习要求验证取左右物理通量平均时，在线性平流情形正好恢复无条件不稳定的时间前向、空间中心格式。

微分守恒形式对光滑解与非守恒形式代数等价，但弱解中的乘积与链式法则不能任意展开。积分式直接表达控制体内总量变化等于边界净通量，不要求解在内部连续。数值方法若以质量、动量、总能量等守恒变量更新，相邻单元共享的界面通量以相反符号出现，内部交换自动抵消。正是这种离散守恒结构，结合一致性与收敛性，保证极限解满足正确 Rankine--Hugoniot 激波速度。

% [原PDF p198]
守恒格式使用有限体积离散。区间 $[0,L]$ 分成 $M$ 个计算单元
\begin{equation}
I_i=[x_{i-1/2},x_{i+1/2}],\qquad
x_{i-1/2}=(i-1)\Delta x,\quad x_{i+1/2}=i\Delta x,\tag{5.45}
\end{equation}
其中 $\Delta x=L/M$。时间层 $t^n$ 上的单元平均定义为
\begin{equation}
u_i^n=\frac{1}{\Delta x}\int_{x_{i-1/2}}^{x_{i+1/2}}u(x,t^n)\,dx.\tag{5.47}
\end{equation}
尽管单元内部的精确解可有空间变化，这个积分平均是单个常数，通常放在单元中心，形成单元中心有限体积法。所有单元平均共同定义一段分片常数数据；实际计算存储的是这些平均的近似值。

\begin{sourceblock}
\sourcecrop{198}{52}{312}{389}{540}{0.30}
\captionof{figure}{有限体积网格：单元界面 $x_{i\pm1/2}$ 上定义数值通量，单元中心存储平均值。}
\end{sourceblock}

界面数值通量可依赖多个邻点，但无论模板多宽，同一 $f_{i+1/2}$ 必须同时用于左单元的流出和右单元的流入。若两个单元各自计算不同的“界面通量”，局部误差就会积累成全局非守恒。条件（5.44）则保证均匀状态下数值通量退化为真实物理通量；否则即便常值解也会被格式改变。一致性是必要条件，却仍需稳定性才能保证收敛。

% [原PDF p199]
\subsection{Godunov 一阶迎风方法}
Godunov 首次把 CIR 格式成功推广为非线性守恒律系统的守恒方法。时间层 $n$ 的数据是分片常数的，相邻状态 $(u_i^n,u_{i+1}^n)$ 在界面 $x_{i+1/2}$ 处形成局部黎曼问题
\[
u_t+f(u)_x=0,\qquad
u(x,0)=\begin{cases}u_i^n,&x<0,\\u_{i+1}^n,&x>0.\end{cases}
\]
Godunov 法以这些局部黎曼问题的解构造界面通量。

第一种解释是：对单元 $I_i$ 两侧的 $RP(u_{i-1}^n,u_i^n)$ 与 $RP(u_i^n,u_{i+1}^n)$ 求精确解，在新时间对落入单元的两半解作积分平均，并把结果赋给 $u_i^{n+1}$。在线性平流且 $a>0$ 时，两组解都只是以速度 $a$ 传播的间断。

有限体积中的 $u_i^n$ 与有限差分点值概念不同。它包含单元内部全部细节的积分信息，即使真实解在单元中含激波，平均仍有明确定义。分片常数重构只是在当前时间层用这个平均代表单元内部，并不声称真实解恒定。一步演化后局部黎曼解产生空间结构，再重新平均回有限个自由度。这种“重构--演化--平均”的循环是现代有限体积法的基本框架，高阶方法只需把常数重构换成受限多项式并提高时间积分精度。

% [原PDF p200]
两个局部黎曼解为
\begin{equation}
u_{i-1/2}(x/t)=\begin{cases}u_{i-1}^n,&x/t<a,\\u_i^n,&x/t>a,
\end{cases}\quad
u_{i+1/2}(x/t)=\begin{cases}u_i^n,&x/t<a,\\u_{i+1}^n,&x/t>a.
\end{cases}\tag{5.48--5.49}
\end{equation}
更新平均写成
\begin{equation}
u_i^{n+1}=\frac1{\Delta x}\left[
\int_0^{\Delta x/2}u_{i-1/2}(x/\Delta t)\,dx+
\int_{-\Delta x/2}^{0}u_{i+1/2}(x/\Delta t)\,dx\right].\tag{5.50}
\end{equation}
若要求 $c=a\Delta t/\Delta x\le1/2$，两侧波在这段时间内不会相互作用，积分可直接按几何长度计算。

第一版本的几何构造中，每个黎曼问题都有自己的局部原点，分别位于 $x_{i-1/2}$ 和 $x_{i+1/2}$；式（5.50）中的两个积分也分别使用这两个局部坐标。只取左问题的右半解、右问题的左半解，才能覆盖单元 $I_i$。条件 $c\le1/2$ 使两条传播前沿在单元中心之前不相遇，积分区域可简单分段。若忽略局部坐标差异，容易把同一条射线错误用于两个界面，导致更新方向和通量符号出错。

% [原PDF p201]
第一项等于 $c u_{i-1}^n+(1/2-c)u_i^n$，第二项等于 $u_i^n/2$，所以
\begin{equation}u_i^{n+1}=u_i^n+c(u_{i-1}^n-u_i^n),\tag{5.52}\end{equation}
恰好是 $a>0$ 时的 CIR 格式。由于更新是守恒律精确局部解的积分平均，守恒性显而易见；$a<0$ 时用同样几何论证也得到相应 CIR 格式。

第二种解释直接导出便于实现的通量形式，并去掉过严的半单元 CFL 限制。把两侧局部解拼成 $\widetilde u(x,t)$，在时空控制体 $[x_{i-1/2},x_{i+1/2}]\times[0,\Delta t]$ 上应用守恒律积分式，单元平均变化就等于左右边界通量的时间积分之差。

从积分长度可清楚看到守恒：左黎曼问题把长度 $c\Delta x$ 的 $u_{i-1}^n$ 搬入单元，其余部分仍为 $u_i^n$；右黎曼问题在 $a>0$ 时没有向左传播信息，贡献半单元的 $u_i^n$。二者相加得到原单元量减去右流出再加左流入。对 $a<0$，几何图镜像，流入来自 $i+1$。Godunov 的贡献不是单纯“使用精确黎曼解”，而是把局部物理解与全局守恒平均一致结合。

% [原PDF p202]
由此得到
\begin{align}
u_i^{n+1}&=u_i^n+\frac{\Delta t}{\Delta x}(f_{i-1/2}-f_{i+1/2}),\tag{5.55}\\
f_{i+1/2}&=\frac1{\Delta t}\int_0^{\Delta t}f(\widetilde u(x_{i+1/2},t))\,dt.
\tag{5.56}
\end{align}
黎曼解自相似，沿局部 $t$ 轴即 $x/t=0$ 的状态恒定，因此 Godunov 通量为
\begin{equation}f^{\mathrm{God}}_{i+1/2}=f(u_{i+1/2}(0)),\tag{5.59}\end{equation}
其中 $u_{i+1/2}(0)$ 是 $RP(u_i^n,u_{i+1}^n)$ 在界面射线上的精确状态。在线性平流且 $a>0$ 时，$f_{i-1/2}=a u_{i-1}^n$、$f_{i+1/2}=a u_i^n$，再次恢复 CIR。实际计算主要采用这一通量版本。

第二版本不需要真正计算新时间的整段黎曼剖面或空间积分，只需界面射线上的状态。因为自相似解沿 $x/t=0$ 不随局部时间变化，时间平均通量退化为单点物理通量。这把复杂度从“每单元两段解积分”降到“每界面一次抽样”，并把允许 CFL 从 $1/2$ 放宽到 1。公式仍源自积分守恒律，所以不是经验性通量拼接。对非线性系统，精确求界面状态可能昂贵，后续近似黎曼求解器就是近似这一关键抽样或直接近似其通量。

% [原PDF p203]
\subsection{Burgers 方程的 Godunov 法}
对无黏 Burgers 方程
\begin{equation}u_t+f(u)_x=0,\qquad f(u)=\frac12u^2,\tag{5.61}\end{equation}
局部黎曼问题在 $u_L>u_R$ 时产生速度 $S=(u_L+u_R)/2$ 的激波，在 $u_L\le u_R$ 时产生稀疏扇：
\begin{equation}
u(x/t)=\begin{cases}
u_L,&u_L>u_R\ \text{且}\ x/t<S,\\
u_R,&u_L>u_R\ \text{且}\ x/t>S,\\
u_L,&u_L\le u_R\ \text{且}\ x/t\le u_L,\\
x/t,&u_L<x/t<u_R,\\
u_R,&u_L\le u_R\ \text{且}\ x/t\ge u_R,
\end{cases}
\end{equation}
其中激波情形按 $x/t\lessgtr S$ 取左右状态。Godunov 通量只需 $x/t=0$ 的抽样值，因此激波按 $S$ 的符号选状态；稀疏波则有全向左、全向右和跨声速三种情况。

五种波型分别是：正速激波、负速激波、全向左稀疏波、全向右稀疏波和跨声速稀疏波。跨声速扇中存在 $u=0$ 的声速点，界面抽样状态为零。

Burgers 通量是凸函数，特征速度 $f'(u)=u$ 随状态变化。$u_L>u_R$ 时左侧特征更快，互相汇聚形成激波，速度由 Rankine--Hugoniot 条件
\[
S=\frac{f(u_R)-f(u_L)}{u_R-u_L}=\frac{u_L+u_R}{2}
\]
确定；$u_L<u_R$ 时特征发散，产生连续稀疏扇，扇内状态满足 $u=x/t$。因此界面抽样不只是看 $u_L,u_R$ 的符号，还要先判波型。图 5.7 的五类情况穷尽凸标量守恒律在零射线处的可能关系。

% [原PDF p204]
完整界面状态可写成
\begin{equation}
u_{i+1/2}(0)=
\begin{cases}
u_i^n,&u_i^n>u_{i+1}^n\ \text{且}\ S\ge0,\\
u_{i+1}^n,&u_i^n>u_{i+1}^n\ \text{且}\ S<0,\\
u_i^n,&u_i^n\le u_{i+1}^n\ \text{且}\ u_i^n\ge0,\\
0,&u_i^n<0<u_{i+1}^n,\\
u_{i+1}^n,&u_i^n\le u_{i+1}^n\ \text{且}\ u_{i+1}^n\le0.
\end{cases}\tag{5.63}
\end{equation}
精确黎曼解也可换成适当近似解。实现时还须处理边界和时间步。内部单元直接使用两个界面通量；边界单元可直接给定物理通量，也可在域外设置虚单元，通过边界黎曼问题求缺失通量。对 Burgers 方程，简单透射条件取 $u_0^n=u_1^n$、$u_{M+1}^n=u_M^n$，使波尽量无反射地离开。左右端所需规则必须在每个时间层重新施加。

式（5.63）中跨声速稀疏波的 $u_*=0$ 使 Godunov 通量为零，满足正确熵解。若把稀疏波错误地当作以平均速度传播的间断，就可能得到“膨胀激波”，违反熵条件。虚单元边界法的优势是把边界物理转化为同一个黎曼求解接口；但虚状态不能任意复制，必须根据入流、出流、固壁等物理特征设置。本例的复制条件只适用于希望波无扰动穿出的简单一维边界。

% [原PDF p205]
非线性问题每一时间层存在多个波速。令 $S_{\max}^n$ 为全域最大绝对波速，定义
\begin{equation}
C_{\mathrm{CFL}}=\frac{\Delta t S_{\max}^n}{\Delta x},\qquad
\Delta t=C_{\mathrm{CFL}}\frac{\Delta x}{S_{\max}^n},\quad 0<C_{\mathrm{CFL}}\le1.\tag{5.65--5.66}
\end{equation}
最大波速估计至关重要；低估会使时间步过大，即使格式本身很精细也可能崩溃。Burgers 方程界面速度可取
\begin{equation}
S_{i+1/2}^n=\begin{cases}
\left|\frac12(u_i^n+u_{i+1}^n)\right|,&\text{激波},\\
\max(|u_i^n|,|u_{i+1}^n|),&\text{稀疏波},
\end{cases}\tag{5.67}
\end{equation}
再令 $S_{\max}^n=\max_i S_{i+1/2}^n$。对线性化稳定界限 $|c|\le1$ 的格式，常取经验值 $C_{\mathrm{CFL}}=0.9$。

每个界面的最大速度必须包含稀疏波的两条边缘特征，而不能只看扇中某一平均速度；激波则使用跃迁速度。遍历 $i=0,\ldots,M$ 时应包括左右边界黎曼问题，因为反射壁或强入流可能产生比内部更快的波。选定 $C_{\mathrm{CFL}}$ 后，式（5.66）使最快波一步不超过一个单元。若下一步会越过指定输出时刻，可再缩短 $\Delta t$，这不会破坏稳定性。实际代码还应处理 $S_{\max}=0$ 的静止常值情形，避免除零。

% [原PDF p206]
\subsection{差分格式的守恒形式}
许多传统差分格式都能重写为通量差形式。Lax--Friedrichs 格式对应的界面通量为
\begin{equation}
f_{i+1/2}=\frac{1+c}{2c}f(u_i^n)+\frac{c-1}{2c}f(u_{i+1}^n).\tag{5.70}
\end{equation}
另一种解释是，在单元 $i$ 内对 $RP(u_{i-1}^n,u_{i+1}^n)$ 于局部时间 $\Delta t/2$ 的解作积分平均：
\begin{equation}
u_i^{n+1}=\frac1{\Delta x}\int_{x_{i-1/2}}^{x_{i+1/2}}
\widetilde u(x,\Delta t/2)\,dx.\tag{5.71}
\end{equation}
在线性平流中，黎曼解只是以速度 $a$ 传播的左右常值，精确积分即恢复原格式。

Lax--Friedrichs 的黎曼积分解释使用邻居 $u_{i-1},u_{i+1}$ 而不使用中心 $u_i$，并在半时间步取平均。这一“跨中心”结构导致相邻网格值成对的现象。尽管可以对非线性系统精确求该黎曼解并积分，守恒积分式表明结果只依赖左右状态与通量，根本无需知道内部波型。由此得到的式（5.77）包含中心通量和与 $\Delta x/\Delta t$ 成正比的耗散，后者正是其稳健但过度扩散的来源。

% [原PDF p207]
这种解释显示 Lax--Friedrichs 更新是邻单元状态所定义黎曼解的加权平均，上游项总有较大权重。对非线性系统
\begin{equation}\bm U_t+\bm F(\bm U)_x=0,\tag{5.73}\end{equation}
可直接对 $RP(\bm U_{i-1}^n,\bm U_{i+1}^n)$ 的精确解积分，但这种版本比常规格复杂、昂贵，结果却几乎无差别，因而没有实际优势。若把积分改为随机抽样，则得到第 7 章的随机选择类方法。

在半时间步控制体上运用守恒积分式，可消去黎曼解，导出系统的常规守恒更新
\begin{align}
\bm U_i^{n+1}&=\bm U_i^n+\frac{\Delta t}{\Delta x}
(\bm F_{i-1/2}-\bm F_{i+1/2}),\tag{5.76}\\
\bm F_{i+1/2}^{\mathrm{LF}}&=\frac12(\bm F_i^n+\bm F_{i+1}^n)
+\frac12\frac{\Delta x}{\Delta t}(\bm U_i^n-\bm U_{i+1}^n).\tag{5.77}
\end{align}

对系统直接实现式（5.74）不仅要解 $RP(\bm U_{i-1},\bm U_{i+1})$，还要在单元上积分可能含多个激波、接触和稀疏扇的解，计算量远大于常规通量，却没有可见精度收益。随机评价积分则只抽一个状态，改变了误差性质并引出第 7 章方法。式（5.75）的推导说明，只要使用积分守恒律，黎曼解对控制体四条边的贡献可精确化简为初始平均与边界通量，不必解析内部细节。

% [原PDF p208]
Lax--Wendroff 格式也能写成守恒形式。在线性平流中，界面通量是左右物理通量的 Courant 加权平均。推广到非线性系统，可先由
\begin{equation}
\bm U_{i+1/2}^{n+1/2}=\frac1{\Delta x}
\int_{-\Delta x/2}^{\Delta x/2}\bm U_{i+1/2}(x,\Delta t/2)\,dx
\end{equation}
求半时间界面平均，再令 $\bm F_{i+1/2}=\bm F(\bm U_{i+1/2}^{n+1/2})$。这一路线称为加权平均通量（WAF）法。

再次用守恒积分式消去黎曼解，则
\begin{equation}
\bm U_{i+1/2}^{n+1/2}=\frac12(\bm U_i^n+\bm U_{i+1}^n)
+\frac12\frac{\Delta t}{\Delta x}(\bm F_i^n-\bm F_{i+1}^n),\tag{5.79}
\end{equation}
这就是非线性系统的两步 Richtmyer 型 Lax--Wendroff 格式。

WAF 版本保留黎曼波结构：半时间的界面平均由每道波在积分区间中所占比例决定，因此适合随后加入通量限制器。Richtmyer 版本则只用左右状态与物理通量预估半步界面状态，形式更简单，但没有显式波分解。两者在线性方程上等价，非线性系统上是不同推广。这个例子强调：一个线性格式往往存在多种非线性扩展，它们都能在小扰动极限恢复原格式，却可能在强激波下表现迥异。

% [原PDF p209]
Lax--Friedrichs 与 Lax--Wendroff 的两种推广彼此平行：保留积分平均就得到基于黎曼问题的版本；用守恒积分式代换后，则分别得到常规 Lax--Friedrichs 通量和 Richtmyer 两步通量。前者的黎曼版本通常无实用价值，后者的 WAF 版本和 Richtmyer 版本都实际使用。

Fromm 格式在线性 $a>0$ 情形的守恒通量为
\[
f_{i+1/2}^{\mathrm{FR}}=-\frac14(1-c)f_{i-1}+f_i+\frac14(1-c)f_{i+1},
\]
Warming--Beam 通量为
\[
f_{i+1/2}^{\mathrm{WB}}=\frac12(c-1)f_{i-1}+\frac12(3-c)f_i.
\]
练习要求验证这些关系，并在负波速 $-2\le c\le-1$ 下由黎曼积分解释推导 Warming--Beam 通量。

\section{线性系统的迎风格式}
考虑常系数双曲系统 $\bm U_t+\bm A\bm U_x=0$。

练习 5.13--5.18 的目的，是训练从点值更新反推共享界面通量。把所求通量代入 $f_{i-1/2}-f_{i+1/2}$，相邻项应望远镜式消去并恢复原差分系数。Warming--Beam 在 $1\le c\le2$ 的黎曼积分解释还说明，当一步传播超过一个单元却不超过两个单元时，上游模板必须扩展到 $i-2$；负速度情形完全镜像。格式稳定范围和模板宽度都可以从特征在一步内跨越的网格数理解。

% [原PDF p210]
设 $m\times m$ 常矩阵 $\bm A$ 的实特征值依次为 $\lambda_1<\cdots<\lambda_m$，右特征矢量为 $\bm K^{(j)}$。若特征值有正有负，对原变量作同一侧单边差分只会对其中一部分波迎风、对另一部分波下风。解决方法是按特征速度符号分裂矩阵。

由双曲性，
\begin{equation}
\bm A=\bm K\bm\Lambda\bm K^{-1},\qquad
\bm\Lambda=\operatorname{diag}(\lambda_1,\ldots,\lambda_m),\tag{5.81--5.82}
\end{equation}
特征变量 $\bm V=\bm K^{-1}\bm U$ 满足解耦系统
\begin{equation}
\bm V_t+\bm\Lambda\bm V_x=0,
\qquad (v_j)_t+\lambda_j(v_j)_x=0.\tag{5.84--5.85}
\end{equation}
于是每个分量都是常速线性平流方程。

特征变换把耦合系统分解为 $m$ 个独立标量波。矩阵 $\bm K$ 的第 $j$ 列是右特征矢量，$\bm K^{-1}$ 把物理变量投影为特征幅度。若所有 $\lambda_j>0$，整个系统可统一用左侧差分；若全负则用右侧；混合符号时，必须在特征空间逐波选择。直接在守恒变量每个分量上根据某个单一“流速”判断方向，会把声波、接触波等不同速度混在一起，无法同时迎风。

% [原PDF p211]
\subsection{CIR 格式}
定义
\begin{equation}
\lambda_j^+=\frac12(\lambda_j+|\lambda_j|),\qquad
\lambda_j^-=\frac12(\lambda_j-|\lambda_j|),\tag{5.86}
\end{equation}
则 $\lambda_j=\lambda_j^++\lambda_j^-$、$|\lambda_j|=\lambda_j^+-\lambda_j^-$。对每个特征分量使用 CIR 格式：
\begin{equation}
(v_j)_i^{n+1}=(v_j)_i^n-\frac{\Delta t}{\Delta x}\lambda_j^+
[(v_j)_i^n-(v_j)_{i-1}^n]-\frac{\Delta t}{\Delta x}\lambda_j^-
[(v_j)_{i+1}^n-(v_j)_i^n].\tag{5.88}
\end{equation}
令 $\bm\Lambda^\pm=\operatorname{diag}(\lambda_1^\pm,\ldots,\lambda_m^\pm)$，并定义
\begin{equation}\bm A^\pm=\bm K\bm\Lambda^\pm\bm K^{-1}.\tag{5.90}\end{equation}
于是 $\bm A=\bm A^++\bm A^-$，$|\bm A|=\bm A^+-\bm A^-$。

对角矩阵的正负部分是逐个特征值分裂，而不是逐个物理分量截断。回到原变量后，$\bm A^+$ 和 $\bm A^-$ 一般是满矩阵；它们分别携带所有非负与非正特征场。式（5.92）在特征空间对每个 $v_j$ 采用正确单边差分，左乘 $\bm K$ 后得到式（5.93）。稳定条件相应要求所有特征 Courant 数 $|\lambda_j|\Delta t/\Delta x\le1$，即由谱半径最大的波决定时间步。

% [原PDF p212]
变回原变量，系统 CIR 格式为
\begin{equation}
\bm U_i^{n+1}=\bm U_i^n-\frac{\Delta t}{\Delta x}\bm A^+
(\bm U_i^n-\bm U_{i-1}^n)-\frac{\Delta t}{\Delta x}\bm A^-
(\bm U_{i+1}^n-\bm U_i^n).\tag{5.93}
\end{equation}
第一差分对 $\bm A^+$ 所含正速波迎风，第二差分对 $\bm A^-$ 所含负速波迎风。因此，矩阵分裂保留了标量迎风法“按信息传播方向选边”的核心原则。

\subsection{Godunov 方法}
把系统写成守恒形式
\begin{equation}\bm U_t+\bm F(\bm U)_x=0,\qquad \bm F=\bm A\bm U.\tag{5.94}\end{equation}
Godunov 法在界面求 $RP(\bm U_i^n,\bm U_{i+1}^n)$，取
\begin{equation}\bm F_{i+1/2}=\bm F(\bm U_{i+1/2}(0)).\tag{5.95}\end{equation}
把左右初值分别展开为 $\sum\alpha_j\bm K^{(j)}$ 和 $\sum\beta_j\bm K^{(j)}$，即可显式构造界面状态。

在线性系统中，黎曼解由 $m$ 道互不改变形状的特征波组成。左右状态展开系数 $\alpha_j,\beta_j$ 的差给出第 $j$ 波强度，跨越该波的状态跳跃为 $(\beta_j-\alpha_j)\bm K^{(j)}$。沿 $S=0$ 抽样时，从左态出发只跨过速度小于或等于零的波；从右态反向看则只撤去正速度波。图 5.8 中零射线位于 $\lambda_I$ 与 $\lambda_{I+1}$ 之间，这使界面状态无须构造全部时空剖面即可得到。

% [原PDF p213]
若 $I$ 满足 $\lambda_I\le0\le\lambda_{I+1}$，则界面状态既可从左状态越过所有非正速度波得到，也可从右状态反向越过所有正速度波得到：
\begin{align}
\bm U_{i+1/2}(0)&=\bm U_i^n+\sum_{j=1}^{I}(\beta_j-\alpha_j)\bm K^{(j)},\tag{5.98}\\
&=\bm U_{i+1}^n-\sum_{j=I+1}^{m}(\beta_j-\alpha_j)\bm K^{(j)}.\tag{5.99}
\end{align}
合并两式，
\begin{equation}
\bm U_{i+1/2}(0)=\frac12(\bm U_i^n+\bm U_{i+1}^n)
-\frac12\sum_{j=1}^{m}\operatorname{sign}(\lambda_j)(\beta_j-\alpha_j)\bm K^{(j)}.\tag{5.100}
\end{equation}
利用 $\bm A\bm K^{(j)}=\lambda_j\bm K^{(j)}$，通量也可写成左右物理通量加波强度修正。

将界面状态代入 $\bm F=\bm A\bm U$ 后，每个波跳跃只需乘对应 $\lambda_j$。从左通量加所有负速波修正得到式（5.102），从右通量减所有正速波修正得到式（5.103）；取二者平均便得到绝对特征速度形式（5.104）。该式清楚区分中心部分与上风耗散部分：波速越大，状态跳跃受到的耗散越强；零速波不被耗散，这也是无缺陷 Godunov 通量能精确保持静止间断的原因。

% [原PDF p214]
三种等价 Godunov 通量形式中，对称形式最紧凑：
\begin{equation}
\bm F_{i+1/2}=\frac12(\bm F_i^n+\bm F_{i+1}^n)
-\frac12\sum_{j=1}^m(\beta_j-\alpha_j)|\lambda_j|\bm K^{(j)}.\tag{5.104}
\end{equation}
由于
\[
\sum_{j=1}^m(\beta_j-\alpha_j)\bm K^{(j)}=\bm U_{i+1}^n-\bm U_i^n,
\]
且 $|\bm A|=\bm K|\bm\Lambda|\bm K^{-1}$，故
\begin{equation}
\boxed{\bm F_{i+1/2}=\frac12(\bm F_i^n+\bm F_{i+1}^n)
-\frac12|\bm A|(\bm U_{i+1}^n-\bm U_i^n).}\tag{5.105}
\end{equation}
第一项是中心通量，第二项是按各特征速度绝对值加权的耗散项。

证明式（5.105）时，利用 $\lambda_j^+-\lambda_j^-=|\lambda_j|$ 以及 $\bm A^\pm\bm K^{(j)}=\lambda_j^\pm\bm K^{(j)}$，即可把逐波求和重新合并成 $|\bm A|(\bm U_R-\bm U_L)$。这也是许多近似黎曼求解器的标准“中心通量减矩阵耗散”结构。实现时不一定显式形成 $|\bm A|$；可投影到特征空间，对每个波强度乘 $|\lambda_j|$，再叠加回物理空间，数值上往往更清晰。

% [原PDF p215]
同一通量还可写成分裂形式
\begin{equation}\boxed{\bm F_{i+1/2}=\bm A^+\bm U_i^n+\bm A^-\bm U_{i+1}^n.}\tag{5.106}\end{equation}
正通量分量只从左状态取值，负通量分量只从右状态取值，这正是系统迎风性的代数表达。练习以线性化气体动力学系统
\[
\bm U_t+\bm A\bm U_x=0,
\quad \bm U=(\rho,u)^T,
\quad \bm A=\begin{bmatrix}0&\rho_0\\a^2/\rho_0&0\end{bmatrix}
\]
为例，要求求出 $\bm\Lambda^\pm,\bm A^\pm$，展开分量格式，并利用星形区显式解直接计算 Godunov 通量。由于该线性系统只有两条定速特征，界面抽样所需波型远少于非线性 Euler 方程。

分裂式（5.106）还可以看作通量矢量分裂的线性原型：$\bm A^+\bm U_L$ 从左侧携带右行信息，$\bm A^-\bm U_R$ 从右侧携带左行信息。对线性化气体动力学，特征值为 $\pm a$，所以无论平均流速如何，两个声学扰动分别从两侧进入界面。星形区显式式展示密度跳跃与速度跳跃怎样相互耦合；界面通量只有一个波型组合，不需要非线性 Euler 的激波/稀疏分支判断。

% [原PDF p216]
\section{数值结果示例}
\subsection{线性平流}
四种格式用于求解 $u_t+(au)_x=0$。测试 1 的光滑初值为
\begin{equation}u(x,0)=\alpha e^{-\beta x^2},\tag{5.110}\end{equation}
参数取 $a=1$、$\alpha=1$、$\beta=8$、$C_{\mathrm{CFL}}=0.8$，区间为 $[-1,1]$。分别在 $t=1,10,100$ 比较精确解与 Godunov、Lax--Friedrichs、Lax--Wendroff、Warming--Beam 数值解。

在 $t=1$，一阶 Godunov/CIR 的数值黏性削低峰值，耗散更强的 Lax--Friedrichs 削峰更明显；这里 $\alpha_{\mathrm{CIR}}=0.1\Delta x$、$\alpha_{\mathrm{LF}}=0.225\Delta x$。两种二阶格式更准确，但已有相位误差：Lax--Wendroff 波峰滞后，Warming--Beam 波峰超前，原因是修正方程中的色散项。

数值实验使用周期性边界，使波越过右端后从左端重新进入，因此可比较很多传播周期而不改变形状的精确解。$t=1$ 对应 125 步，说明每步平均 $\Delta t=0.008$；随着时间增加，微小的每步振幅与相位误差系统累积。一阶格式的误差主要表现为峰值衰减和波宽增长，二阶格式主要表现为相位漂移。这种诊断比只比较某一时刻的 $L^2$ 误差更能揭示算法用于声学、涡输运等长时问题时的实际局限。

% [原PDF p217]
演化到 $t=10$，一阶格式的数值扩散已使峰值只剩真实值的约 $30\%$--$40\%$；二阶格式虽仍较好，却出现清楚的色散相位误差，同时也开始受扩散影响。到 $t=100$，一阶格式几乎把波压平，二阶格式的波位误差和削峰也不可接受。声学计算常同时要求解析极弱信号和很长传播距离，因此会淘汰大量短时看似合格的格式。

测试 2 使用方波初值
\begin{equation}
u(x,0)=\begin{cases}0,&x\le0.3,\\1,&0.3<x<0.7,\\0,&x\ge0.7.
\end{cases}\tag{5.111}
\end{equation}
一阶格式把间断涂抹到许多单元，Lax--Friedrichs 还出现相邻点成对；二阶格式减小涂抹，却在间断两侧产生超调和欠调。随时间增长，这两类误差都持续累积。第 13、14 章将用 TVD 与限制器抑制这些振荡。

方波包含两个传播方向相同、跃迁符号相反的间断。线性方程没有特征汇聚来压缩跃迁，所以任何人工黏性都会持续扩大过渡区；二阶色散则在每个跃迁的不同侧产生符号相反的振铃。Lax--Friedrichs 的相邻点配对来自其更新完全跳过中心旧值。到 $t=100$，即使格式稳定且守恒，解也可能在工程意义上毫无用处，这再次说明“稳定”只保证误差不爆炸，不保证长时精度。

% [原PDF p218]
\subsection{无黏 Burgers 方程}
测试 3 在 $[0,3/2]$ 上求解
\begin{equation}u_t+\left(\frac12u^2\right)_x=0,
\qquad u(x,0)=\begin{cases}-1/2,&x\le1/2,\\1,&1/2<x\le1,\\0,&x>1.
\end{cases}\tag{5.112--5.113}
\end{equation}
取 $M=75$、$\Delta x=0.02$、$C_{\mathrm{CFL}}=0.8$。在 $t=0.5$，右侧间断是满足熵条件的激波，两侧特征汇入激波，压缩机制有利于尖锐捕捉。Godunov 法约用 3 个单元解析激波，而 Lax--Friedrichs 约需 10 个单元。$x=1/2$ 的跨声速稀疏波处出现熵瑕疵；Godunov 法受影响，而耗散更强的 Lax--Friedrichs 不明显。如何构造满足熵条件的格式是关键问题。

\subsection*{精度、稳定性、守恒性与迎风性的联系}
本章列出的格式说明，评价一个离散方法不能只用单一标签。所谓一致性，是把足够光滑的精确解代入差分方程后，局部截断误差随 $\Delta x,\Delta t\to0$ 消失；精度阶说明消失速度。稳定性控制离散扰动在有限时间内的放大；线性常系数问题可用 Fourier 模态求放大因子，非线性系统则还依赖能量估计、单调性或总变差分析。Lax 等价定理把三者联系起来：对适定线性初值问题，一致性加稳定性等价于收敛性。

守恒性是另一层结构要求。格式即使稳定、一致，若更新不是共享数值通量之差，含激波时也可能收敛到错误传播速度。有限体积守恒式在每个单元写净通量，内部界面自动抵消；全域总量只受物理边界通量改变。Lax--Wendroff 定理所保证的是：一个收敛的守恒格式，其极限满足弱守恒律。它并不保证自动选中唯一熵解，所以还需要熵条件、熵修正或单调/TVD 性质排除非物理解。

迎风性描述离散信息方向与特征传播方向一致。标量常速方程只有一个速度 $a$，正速取左、负速取右；线性系统要先投影到特征场，每个 $\lambda_j$ 分别选边；非线性系统的速度随状态变化，Godunov 法通过局部黎曼解自动决定界面上哪些波从左来、哪些从右来。中心格式不显式选择波向，通常必须加入足够耗散才能稳定。

单调格式把新值写成旧邻域值的非负组合，因而保持最大值原理，不产生新极值。对标量线性格式，这是无振荡和稳健性的强保证，但 Godunov 定理限制其最多一阶。Lax--Wendroff 等二阶线性格式必须出现负系数，因而在间断处振荡。高分辨率方法的核心不是推翻这一定理，而是使用非线性数据相关系数：光滑区采用高阶修正，检测到陡梯度时限制修正，使局部退回单调一阶通量。

修正方程把误差机制翻译为连续项。一阶 CIR 与 LF 的主误差是正的二阶扩散，表现为振幅衰减、间断展宽和极值削平；LF 系数更大，所以更耗散。二阶 LW、WB、Fromm 的主误差通常是三阶色散，表现为不同波数相速不同、波峰超前或滞后，并在间断附近产生振铃。真实算法常同时含更高阶扩散和色散，所以长时结果既会移位也会削峰。

CFL 数 $c=a\Delta t/\Delta x$ 可解释为一步物理传播距离与单元宽度之比。对 CIR，$|c|\le1$ 使特征足点落在上游模板内；对 Warming--Beam，更宽模板允许传播最多两个单元；对第一版 Godunov 和常规 RCM，半单元限制来自要求左右局部波在单元内不相交。不同格式的 CFL 上限不能只由“显式”或“一阶”猜测，必须结合模板与推导。

Godunov 通量的定义包含两个操作：先解相邻常状态的黎曼问题，再沿固定界面射线抽样。对线性系统，结果等价于中心通量减 $|\bm A|$ 加权的状态跳跃；对 Burgers 方程，必须区分激波和稀疏波；对 Euler 方程，还要区分左右非线性波与接触。近似黎曼求解器可近似全部波结构，也可只近似界面通量，但应保持一致、守恒并正确覆盖主要传播速度。

数值通量不是物理通量的简单平均。LF 通量在中心平均外加入统一耗散；Godunov 通量取精确界面状态的物理通量；Richtmyer 通量先预测半时间状态；WAF 通量对黎曼解作半时间空间平均。相同的线性格式可以有不止一种非线性推广，所以程序或论文只说“使用 Lax--Wendroff”仍可能不够，必须说明具体通量构造。

边界处理与内部格式同样决定结果。透射边界复制邻近状态只是一维近似，无限域中的出射波在多维会以不同角度到达边界，简单复制会反射。反射壁则必须翻转法向速度而保留切向速度、密度和压力，再通过边界黎曼问题产生壁压。无论何种边界，由其产生的最快波都应纳入 CFL 估计。

三组示例还说明测试设计的重要性。光滑高斯波主要测量幅值与相位误差；线性方波把扩散和 Gibbs 振荡暴露出来；非线性 Burgers 数据同时包含稀疏、跨声速点与熵激波，可检验格式是否正确选择弱解。对 Euler 方程还要增加接触、近真空、强激波和慢波测试。单一算例通过并不能证明方法普遍可靠。

在实现层面，建议分别保存新旧时间层，先计算所有界面通量再更新；记录每步全域守恒量、最小密度和压力、最大波速及实际 CFL；输出时同时比较多种原始量和守恒量。网格收敛应至少使用三组网格，检查误差阶、间断位置和非物理振荡随 $\Delta x$ 的变化。这样才能把公式中的稳定性、守恒性与实际代码行为对应起来。

\section{Godunov 方法的 FORTRAN 程序}
下面的程序用精确 Burgers 黎曼求解器与一阶 Godunov 更新。输入文件提供 CFL 系数、区域长度、测试编号、单元数、屏幕输出频率、最大步数和目标时刻；输出文件保存数值解。

% [原PDF p219]
程序名为 \texttt{HL-B1GOD}，属于作者的 NUMERICA 教学与研究代码库。主程序声明公共数据区和时间步，随后依次读取参数、初始化场，并进入显式时间推进循环。原程序采用 FORTRAN 77 风格，数组容量和输入输出文件名均在源码中固定。
\begin{lstlisting}[language=Fortran]
      IMPLICIT NONE
      INTEGER ITEST, CELLS, N, NFREQ, NTMAXI
      REAL CFLCOE, DOMLEN, DT, TIME, TIMEOUT, TIMETO
      COMMON /DATAIN/ CFLCOE, DOMLEN, ITEST, CELLS,
     &                NFREQ, NTMAXI, TIMEOUT
      COMMON /DELTAT/ DT
      DATA TIMETO /1.0E-07/
      CALL READER
      CALL INITIA(DOMLEN, ITEST, CELLS)
\end{lstlisting}
符号说明：\texttt{CFLCOE} 是 CFL 系数，\texttt{DOMLEN} 是区域长度，\texttt{CELLS} 是计算单元数，\texttt{TIMEOUT} 是目标输出时刻，\texttt{DT} 是当前时间步。

% [原PDF p220]
时间循环的计算顺序体现了有限体积程序的典型结构：先设置虚单元边界值，再按当前最大特征速度确定稳定时间步，更新物理时间，计算全部界面通量，最后用通量差更新单元平均。
\begin{lstlisting}[language=Fortran]
      TIME = 0.0
      DO 10 N = 1, NTMAXI
         CALL BCONDI(CELLS)
         CALL CFLCON(CFLCOE, CELLS, TIME, TIMEOUT)
         TIME = TIME + DT
         CALL FLUXES(CELLS)
         CALL UPDATE(CELLS)
         IF(MOD(N,NFREQ).EQ.0) WRITE(6,20) N,TIME
\end{lstlisting}
边界必须在通量计算前设置，因为首、末界面的黎曼问题都依赖虚单元；CFL 时间步也必须使用当前层状态，不能在更新后才估计。

% [原PDF p221]
达到目标时刻后，程序调用 \texttt{OUTPUT} 写出解并终止。为避免最后一步越过输出时刻，\texttt{CFLCON} 会在必要时缩短 $\Delta t$。读入子程序 \texttt{READER} 从 \texttt{b1god.ini} 顺序读取七个参数，因此输入文件的行次必须与程序一致。
\begin{lstlisting}[language=Fortran]
         IF(ABS(TIME-TIMEOUT).LE.TIMETO) THEN
            CALL OUTPUT(CELLS)
            WRITE(6,*) 'Number of time steps = ',N
            STOP
         ENDIF
   10 CONTINUE
   20 FORMAT(I12,6X,F12.7)
      END

      SUBROUTINE READER
      OPEN(UNIT=1,FILE='b1god.ini',STATUS='UNKNOWN')
\end{lstlisting}
\texttt{TIMETO} 是比较浮点时间的容差，避免直接用严格相等判断。

% [原PDF p222]
读入顺序为 CFL 系数、区域长度、测试编号、单元数、屏幕输出频率、最大时间步数及输出时刻。随后把参数回显到标准输出，便于确认算例配置。
\begin{lstlisting}[language=Fortran]
      READ(1,*) CFLCOE
      READ(1,*) DOMLEN
      READ(1,*) ITEST
      READ(1,*) CELLS
      READ(1,*) NFREQ
      READ(1,*) NTMAXI
      READ(1,*) TIMEOUT
      CLOSE(1)
\end{lstlisting}
初始化子程序 \texttt{INITIA} 同时清零通量数组与解数组。\texttt{DX} 表示空间网格宽度，\texttt{XPOS} 是单元采样位置，\texttt{XLEFT}、\texttt{XMIDDL}、\texttt{XRIGHT} 是分段初值的三个分界位置。数组下标从 0 到 \texttt{CELLS+1}，两端额外元素留作虚单元。

% [原PDF p223]
规则单元宽度为 \texttt{DX=DOMLEN/REAL(CELLS)}。程序先把所有数组元素置零。若 \texttt{ITEST=1}，则在 $[-1,1]$ 上初始化光滑高斯分布 $u=\exp(-8x^2)$：
\begin{lstlisting}[language=Fortran]
      DX = DOMLEN/REAL(CELLS)
      DO 10 I = 0, IDIM+1
         FLUX(I)=0.0
         U(I)=0.0
   10 CONTINUE
      IF(ITEST.EQ.1) THEN
         XPOS=-1.0
         DO 20 I=1,CELLS
            XPOS=XPOS+2.0/REAL(CELLS)
            U(I)=EXP(-8.0*XPOS*XPOS)
   20    CONTINUE
\end{lstlisting}
这对应前述线性平流光滑测试的轮廓；在 Burgers 方程中它会因非线性速度随 $u$ 改变而逐渐变形。

% [原PDF p224]
若不是测试 1，程序构造四段分片常数数据：在 $0.1L$、$0.5L$、$0.9L$ 处分隔，状态依次为 $-1,1,0,-1$。
\begin{lstlisting}[language=Fortran]
      ELSE
         XLEFT=0.1*DOMLEN
         XMIDDL=0.5*DOMLEN
         XRIGHT=0.9*DOMLEN
         DO 30 I=1,CELLS
            XPOS=(REAL(I)-1.0)*DX
            IF(XPOS.LT.XLEFT) U(I)=-1.0
            IF(XPOS.GE.XLEFT.AND.XPOS.LE.XMIDDL) U(I)=1.0
            IF(XPOS.GT.XMIDDL.AND.XPOS.LE.XRIGHT) U(I)=0.0
            IF(XPOS.GT.XRIGHT) U(I)=-1.0
   30    CONTINUE
      ENDIF
\end{lstlisting}
接下来的 \texttt{BCONDI} 子程序将根据周期边界复制虚单元状态。

% [原PDF p225]
周期边界条件令左虚单元等于最右物理单元，右虚单元等于最左物理单元：
\begin{lstlisting}[language=Fortran]
      SUBROUTINE BCONDI(CELLS)
      U(0)=U(CELLS)
      U(CELLS+1)=U(1)
      RETURN
      END
\end{lstlisting}
这样穿过右边界的波会从左边界重新进入。CFL 子程序遍历含虚单元的全部状态，寻找 Burgers 特征速度 $|u|$ 的最大值；因为 $f'(u)=u$，这里的最大特征速度正是 $\max_i|U(i)|$。
\begin{lstlisting}[language=Fortran]
      SMAX=-1.0E+06
      DO 10 I=0,CELLS+1
         IF(ABS(U(I)).GT.SMAX) SMAX=ABS(U(I))
   10 CONTINUE
\end{lstlisting}

% [原PDF p226]
稳定时间步按
\[
\Delta t=C_{\mathrm{CFL}}\frac{\Delta x}{S_{\max}}
\]
计算。如果 $t+\Delta t$ 将越过目标输出时刻，就把本步缩短为 \texttt{TIMEOUT-TIME}，从而恰好落到输出时刻。
\begin{lstlisting}[language=Fortran]
      DT=CFLCOE*DX/SMAX
      IF((TIME+DT).GT.TIMEOUT) THEN
         DT=TIMEOUT-TIME
      ENDIF
\end{lstlisting}
\texttt{UPDATE} 使用显式守恒公式，先计算 \texttt{DTODX=DT/DX}，再逐单元以左、右界面通量之差更新。这里 \texttt{FLUX(I-1)} 对应 $f_{i-1/2}$，\texttt{FLUX(I)} 对应 $f_{i+1/2}$。

% [原PDF p227]
更新核完整写作
\begin{lstlisting}[language=Fortran]
      DO 10 I=1,CELLS
         U(I)=U(I)+DTODX*(FLUX(I-1)-FLUX(I))
   10 CONTINUE
\end{lstlisting}
它逐单元保持通量成对抵消，因此除边界净通量外，全域离散总量守恒。输出子程序打开 \texttt{numer.out}，把每个单元的位置和数值写成两列：
\begin{lstlisting}[language=Fortran]
      DO 10 I=1,CELLS
         XPOS=REAL(I)*DX
         WRITE(1,20) XPOS,U(I)
   10 CONTINUE
   20 FORMAT(2(4X,F10.5))
\end{lstlisting}
注意这里输出坐标取 $i\Delta x$，而严格单元中心通常是 $(i-1/2)\Delta x$；使用该数据绘图时应知道这一约定。

% [原PDF p228]
\texttt{FLUXES} 对每个界面 $i=0,\ldots,M$ 设置左、右状态 \texttt{UL=U(I)}、\texttt{UR=U(I+1)}，调用精确黎曼子程序得到界面抽样状态 \texttt{USTAR}，再计算 Godunov 通量 $f=\tfrac12u_*^2$。
\begin{lstlisting}[language=Fortran]
      DO 10 I=0,CELLS
         UL=U(I)
         UR=U(I+1)
         CALL RIEMANN(UL,UR,USTAR)
         FLUX(I)=0.5*USTAR*USTAR
   10 CONTINUE
\end{lstlisting}
边界界面与内部界面使用完全相同的流程，差异只在虚单元状态的构造。把边界处理与通量核分开，是这段教学代码中值得保留的模块化思想。

% [原PDF p229]
精确 Burgers 黎曼求解器先判断 \texttt{UL.GT.UR}。若成立，解为激波，速度 $S=(u_L+u_R)/2$；在界面 $x/t=0$ 抽样时，$S\ge0$ 取左状态，否则取右状态。
\begin{lstlisting}[language=Fortran]
      IF(UL.GT.UR) THEN
         S=0.5*(UL+UR)
         IF(S.GE.0.0) THEN
            USTAR=UL
         ELSE
            USTAR=UR
         ENDIF
      ELSE
\end{lstlisting}
若 $u_L\le u_R$，解为稀疏波；程序随后依据左右特征速度相对零的位置区分三类波型。这里的 \texttt{USTAR} 不是整块星形区，而只是沿界面射线抽到的单个状态。

% [原PDF p230]
稀疏波若 $u_L\ge0$，整个扇位于界面右侧，取 $u_*=u_L$；若 $u_R\le0$，整个扇位于左侧，取 $u_*=u_R$；若 $u_L\le0\le u_R$，界面穿过跨声速扇，取声速点 $u_*=0$。
\begin{lstlisting}[language=Fortran]
         IF(UL.GE.0.0) USTAR=UL
         IF(UR.LE.0.0) USTAR=UR
         IF(UL.LE.0.0.AND.UR.GE.0.0) USTAR=0.0
      ENDIF
      RETURN
      END
\end{lstlisting}
至此程序闭合了初始化、周期边界、CFL 步长、精确黎曼通量、守恒更新与输出的完整循环。

\subsection*{程序结构与注释的中文释义}
原清单开头说明：程序用途是“联合精确黎曼求解器，以 Godunov 一阶迎风格式求解无黏 Burgers 方程”；输入为 \texttt{b1god.ini}，数值输出为 \texttt{numer.out}。主驱动中的“读取问题参数”“设置初始条件”“施加边界条件”“施加 Courant--Friedrichs--Lewy 条件”“计算界面数值通量”“按守恒公式更新解”“在指定时刻输出”等注释，正好对应一个有限体积时间步的各阶段。

变量表中，\texttt{ITEST} 选择测试；\texttt{N} 是时间步编号；\texttt{NFREQ} 控制终端报告频率；\texttt{NTMAXI} 是最大步数，防止因配置错误无限循环；\texttt{TIME} 是当前物理时间；\texttt{TIMEOUT} 是目标时间；\texttt{TIMETO} 是输出时刻比较容差。公共块 \texttt{/DATAIN/} 传递输入参数，\texttt{/DELTAT/}、\texttt{/DELTAX/} 分别传递时间步和网格宽度，\texttt{/SOLUTI/}、\texttt{/FLUXFS/} 保存解和界面通量。现代语言可把这些公共块封装为显式数据结构，以避免隐藏依赖，但数值算法并不改变。

\texttt{READER} 的注释含义是“读入问题初始参数”；它按固定顺序读取，随后把所有输入回显到屏幕。\texttt{INITIA} 的用途是“设置解 $U$ 的初值并初始化其他变量”，两组初值由 \texttt{ITEST} 决定。局部变量中，\texttt{DX} 为空间网格宽度，\texttt{FLUX} 为界面通量数组，\texttt{U} 为数值解数组，\texttt{XPOS} 为 $x$ 轴位置，三个分界坐标用于构造方波。清零数组不仅避免未初始化值，也确保虚单元在首次边界赋值前处于确定状态。

\texttt{BCONDI} 的用途是“施加边界条件”。此清单采用周期边界：左虚单元复制最右单元，右虚单元复制最左单元。若改为透射边界，应分别复制最近的物理单元；若改为规定入流，则应由外部函数设置相应虚状态。边界子程序不直接算通量，而只提供完整黎曼初值，这使 \texttt{FLUXES} 对边界和内部保持同一循环。

\texttt{CFLCON} 的用途是“应用 CFL 条件计算稳定时间步”。循环寻找最大特征速度，随后由 CFL 系数乘 $\Delta x/S_{\max}$。最后一步若超过目标时刻，则重新计算时间步。实际稳健实现还应检查 \texttt{SMAX} 是否为正的有限数；若全域 $u=0$，方程当前静止，可直接走到输出时刻，而不是进行除零。

\texttt{UPDATE} 的用途是“用显式守恒公式推进到新时间层”。它原地覆盖 $U(i)$，但通量数组已在更新前由整个旧状态计算完，所以单元遍历顺序不会改变结果。所有内部界面通量一次加到右单元、一次从左单元减去，离散总量严格望远镜抵消。\texttt{OUTPUT} 在指定时刻输出位置和解；若把 $U(i)$ 解释为单元平均，更一致的绘图坐标是单元中心 $x_i=(i-1/2)\Delta x$。

\texttt{FLUXES} 的用途是“由 Godunov 一阶迎风法和精确黎曼求解器计算界面通量”。每次循环先定义局部黎曼问题的左、右状态，再由 \texttt{RIEMANN} 返回零射线状态，最后代入物理通量 $f(u_*)=u_*^2/2$。注释“解存入 FLUX”严格说应理解为通量而非状态；状态暂存在 \texttt{USTAR} 中。

\texttt{RIEMANN} 的用途是“精确求解无黏 Burgers 黎曼问题”。局部变量 \texttt{UL}、\texttt{UR} 是左右初态，\texttt{S} 是激波速度，\texttt{USTAR} 是抽样状态。激波分支先计算 $S$，再沿 $t$ 轴抽样；稀疏波分支依次处理全向右、全向左与跨声速三种情况。三个独立 \texttt{IF} 在边界等号处可能重复赋同一个相容值，不造成歧义；现代写法也可用互斥 \texttt{ELSE IF} 明确表达波型分区。

整段程序虽短，却完整呈现 Godunov 法的数学对象：单元平均是状态自由度，虚单元提供边界黎曼数据，CFL 条件约束局部波传播，精确黎曼解确定共享界面通量，通量差完成守恒推进。把 Burgers 求解器替换为 Euler 求解器时，外层时间循环、边界--CFL--通量--更新顺序均可保留；主要变化是标量数组变成守恒矢量，\texttt{RIEMANN} 需要处理三波结构，最大速度也要取所有左右外层波。

% [原PDF p231]
\begin{sourceblock}
\sourcecrop{231}{45}{385}{395}{610}{0.28}\\[0.45em]
\sourcecrop{231}{45}{115}{395}{337}{0.28}
\captionof{figure}{测试 1 的长时间传播：四种格式在 $t=1$（125 步）与 $t=10$（1250 步）的数值解（符号）和精确解（实线）。图中 Position 为位置。}
\end{sourceblock}
Godunov 和 Lax--Friedrichs 的峰值随时间明显降低；Lax--Wendroff 与 Warming--Beam 保峰较好，但分别出现滞后与超前。图中英文格式名为方法的通用专名，曲线物理量均为无量纲 $u$。

% [原PDF p232]
\begin{sourceblock}
\sourcecrop{232}{45}{385}{395}{610}{0.28}\\[0.45em]
\sourcecrop{232}{45}{115}{395}{337}{0.28}
\captionof{figure}{上：测试 1 在 $t=100$（12499 步）的结果；下：方波测试 2 在 $t=1$（125 步）的结果。横轴均为位置。}
\end{sourceblock}
经过极长时间，光滑波也会被一阶耗散压平或被二阶色散移位。对方波，一阶解单调但间断很宽；二阶解的跃迁较窄，却出现明显超调与欠调。这两组结果直接展示“耗散过强”和“色散振荡”两种典型误差。

% [原PDF p233]
\begin{sourceblock}
\sourcecrop{233}{45}{385}{395}{610}{0.28}\\[0.45em]
\sourcecrop{233}{45}{115}{395}{337}{0.28}
\captionof{figure}{方波测试 2 在 $t=10$（1250 步）与 $t=100$（12499 步）的结果；数值解为符号，精确解为实线。}
\end{sourceblock}
误差随传播时间积累：Godunov 与 Lax--Friedrichs 的平台逐渐降低并展宽，Lax--Wendroff 和 Warming--Beam 的振荡区扩大且相位偏移。因而评价波动算法不能只看一两个时间步，还应检查与应用相关的长时传播尺度。

% [原PDF p234]
\begin{sourceblock}
\sourcecrop{234}{48}{385}{393}{610}{0.29}
\captionof{figure}{无黏 Burgers 测试 3 在 $t=0.5$（32 步）的结果。上为 Godunov，下为 Lax--Friedrichs；符号是数值解，实线是精确解。}
\end{sourceblock}
非线性压缩使激波比线性平流间断更易保持陡峭。Godunov 法的激波过渡约 3 个单元，而 Lax--Friedrichs 约 10 个单元；在跨声速稀疏波处，前者能看到小的熵瑕疵，后者的强耗散则把它掩盖。

% [原PDF p235]
\chapter{非线性系统的 Godunov 方法}

Godunov 在约四十年前把 Courant--Isaacson--Rees 一阶迎风格式守恒地推广到非线性双曲守恒律系统。第 5 章已在标量方程和常系数线性系统上介绍基本思想，本章处理一般非线性系统，并详细说明一维非定常 Euler 方程。

方法的核心是黎曼问题解，可以是精确解，也可以是适当近似解。本章采用精确解；第 9--12 章将介绍近似黎曼求解器，谨慎使用可提高效率。基本 Godunov 法仅一阶准确，实际问题若要充分分辨往往需要极细网格；第 13、14、16 章将给出二阶、三阶和 TVD 扩展。

\section{Godunov 方法的基础}
考虑一般非线性双曲守恒律初边值问题
\begin{equation}
\bm U_t+\bm F(\bm U)_x=0,\quad
\bm U(x,0)=\bm U^{(0)}(x),\quad
\bm U(0,t)=\bm U_l(t),\quad \bm U(L,t)=\bm U_r(t).\tag{6.1}
\end{equation}

% [原PDF p236]
$\bm U$ 是守恒变量矢量，$\bm F$ 是通量矢量，空间域为 $[0,L]$。为容纳间断解，对任意时空控制体采用积分守恒式
\begin{equation}
\int_{x_1}^{x_2}\bm U(x,t_2)dx=\int_{x_1}^{x_2}\bm U(x,t_1)dx
+\int_{t_1}^{t_2}\bm F(\bm U(x_1,t))dt
-\int_{t_1}^{t_2}\bm F(\bm U(x_2,t))dt.\tag{6.2}
\end{equation}
把 $[0,L]$ 分成 $M$ 个规则单元 $I_i=[x_{i-1/2},x_{i+1/2}]$，
\begin{equation}
x_{i-1/2}=(i-1)\Delta x,\quad x_i=(i-\tfrac12)\Delta x,
\quad x_{i+1/2}=i\Delta x.\tag{6.3}
\end{equation}
非线性系统的波速随时空变化，所以时间步通常在推进过程中动态选择。给定 $t^n$ 的一般数据，先取单元平均，形成分片常数表示。

时间区间 $[0,T]$ 的终点 $T$ 只是输出时刻，不是时间方向的边界。空间网格可先固定，时间网格则因 $S_{\max}$ 随解变化而非均匀。对一般数据 $\widetilde{\bm U}(x,t^n)$ 取平均后，每个单元只保留一个守恒矢量；图 6.1 仅画其中一个分量。这样的投影会损失单元内变化，正是基本方法一阶空间精度的来源，但它保证单元总质量、动量和能量与投影前一致。

% [原PDF p237]
单元平均及其分片常数重构为
\begin{equation}
\bm U_i^n=\frac1{\Delta x}\int_{x_{i-1/2}}^{x_{i+1/2}}
\widetilde{\bm U}(x,t^n)dx,
\qquad \bm U(x,t^n)=\bm U_i^n\quad(x\in I_i).\tag{6.4--6.5}
\end{equation}
守恒量可转换成便于解黎曼问题的原始量；对 Euler 方程，$\bm W=(\rho,u,p)^T$。

相邻常状态在每个界面形成 $RP(\bm U_i^n,\bm U_{i+1}^n)$，局部解记为 $\bm U_{i+1/2}(\bar x/\bar t)$。在时间步足够小、局部波未越过相邻界面的条件下，可把这些局部自相似解拼成时间条带上的全局精确解：
\begin{equation}
\widetilde{\bm U}(x,t)=\bm U_{i+1/2}(\bar x/\bar t),
\qquad x\in[x_i,x_{i+1}].\tag{6.6}
\end{equation}

每个界面初值严格说应写成 $\bm U_i^n$（左）和 $\bm U_{i+1}^n$（右）。Euler 精确黎曼求解通常先转成 $(\rho,u,p)$，解出 $p_*,u_*$ 及波型，再恢复各区域守恒量。不同界面的局部坐标原点不同，但解都只依赖 $\bar x/\bar t$。若时间步足够小，界面发出的波在本步内不会穿过相邻界面，因此可在相邻单元中心之间用该局部解定义全局拼接解，且拼接处相容。

% [原PDF p238]
全局坐标与每个界面黎曼问题的局部坐标关系为
\begin{equation}
\bar x=x-x_{i+1/2},\qquad \bar t=t-t^n,
\quad -\frac{\Delta x}{2}\le\bar x\le\frac{\Delta x}{2},
\quad0\le\bar t\le\Delta t.\tag{6.7}
\end{equation}
找到局部解后，Godunov 法通过新的单元平均把解推进到 $t^{n+1}=t^n+\Delta t$。

\section{Godunov 格式}
第一种版本直接定义
\begin{equation}
\bm U_i^{n+1}=\frac1{\Delta x}\int_{x_{i-1/2}}^{x_{i+1/2}}
\widetilde{\bm U}(x,t^{n+1})dx.\tag{6.8}
\end{equation}
为保证单元内左右波系不相互作用，需 $\Delta t\le\Delta x/(2S_{\max}^n)$。此时只有左界面向右传播的波和右界面向左传播的波影响单元。

图 6.3 的全局时间 $t^n$ 对应每个局部问题的 $\bar t=0$，界面 $x_{i+1/2}$ 对应 $\bar x=0$。写公式时常省略横线，直接用 $(x,t)$ 表示局部坐标，读者必须根据上下文辨认。Godunov 步的两个操作有明确含义：局部精确演化让分片常数数据按照真实方程产生波，重新取平均则把无限维解压回有限个单元自由度。前者引入物理传播，后者引入一阶耗散。

% [原PDF p239]
上述平均可写成两半局部黎曼解的积分：
\begin{equation}
\bm U_i^{n+1}=\frac1{\Delta x}\int_0^{\Delta x/2}
\bm U_{i-1/2}(x/\Delta t)dx
+\frac1{\Delta x}\int_{-\Delta x/2}^{0}
\bm U_{i+1/2}(x/\Delta t)dx.\tag{6.10}
\end{equation}
它可直接实现，但半单元 CFL 限制过严，且含稀疏扇时积分复杂。

\textbf{命题 6.1}\quad Godunov 法可写成
\begin{equation}
\bm U_i^{n+1}=\bm U_i^n+\frac{\Delta t}{\Delta x}
(\bm F_{i-1/2}-\bm F_{i+1/2}),\qquad
\bm F_{i+1/2}=\bm F(\bm U_{i+1/2}(0)),\tag{6.11--6.12}
\end{equation}
只需满足 $\Delta t\le\Delta x/S_{\max}^n$。证明是在单元时空控制体上对局部精确解应用积分守恒律。

第一版本的积分需要知道稀疏扇内状态随 $x/t$ 的连续函数，并对每个分量积分；复杂状态方程下代价更高。更严重的是条件（6.9）只允许最快波走半个单元，使时间步比通常稳定性要求小一半。命题 6.1 的第二版本通过时空守恒直接计算单元收支，不需要内部空间积分。其证明并未近似积分式，而是利用局部黎曼解沿界面射线恒定这一精确性质。

% [原PDF p240]
界面上的自相似状态恒定：
\begin{equation}
\widetilde{\bm U}(x_{i\pm1/2},t)=\bm U_{i\pm1/2}(0).\tag{6.15}
\end{equation}
因此时间积分直接化为 $\Delta t$ 乘界面物理通量，除以 $\Delta x$ 即得式（6.11）。第二版本允许一条波在一步内最多走完整单元宽度，比第一版本允许更大时间步，实际计算均采用它。

即使单元内部发生波相互作用，只要相互作用没有使波加速到穿过下一界面，界面通量推导仍成立；这是一种局部线性化假设。CFL 条件真正约束的是在当前时间步内界面射线保持属于同一个局部状态。

\section{Euler 方程的 Godunov 方法}
已知时间层 $n$ 的 $\{\bm U_i^n\}$，用守恒公式推进需要左右界面通量。每个通量来自相应 Euler 黎曼问题在 $S=x/t=0$ 的抽样状态。

% [原PDF p241]
\subsection{界面通量的计算}
第 4 章给出了 Euler 黎曼问题的完整精确解与原始变量抽样过程。一般黎曼解有左非线性波、接触间断和右非线性波；左右非线性波各可为激波或稀疏波，界面相对接触的位置又取决于星形区速度 $u_*\gtrless0$。因此，计算 Godunov 通量时仍须区分十种波型，令 $S=0$ 并不会消除这些分支。

\begin{sourceblock}
\sourcecrop{241}{52}{235}{388}{530}{0.38}
\captionof{figure}{Euler 方程计算 Godunov 通量时的十种波型：左列对应 $u_*\ge0$，右列对应 $u_*\le0$。粗线为激波，虚线为接触，扇形为稀疏波。}
\end{sourceblock}

式（6.14）左端是新层单元总守恒量，右端第一项是旧层总量，后两项分别是左界面流入和右界面流出。条件（6.13）保证在积分时间内每个界面的状态不被来自下一界面的波改变，因此式（6.15）可把时间积分精确化为常值通量。把初末空间积分分别除以 $\Delta x$，就成为单元平均。这一推导也说明通量差中的正负号由控制体外法向决定，不能按“左、右状态”随意改变。

% [原PDF p242]
例如 $u_*>0$ 时，界面在接触左侧。若左波为激波，先用激波关系求左星区状态 $\bm W_{*L}$ 和激波速度 $S_L$；若 $S_L\ge0$，界面尚在左初态，取 $\bm W_{i+1/2}(0)=\bm W_L$，若 $S_L<0$，取 $\bm W_{*L}$。左波为稀疏波时，还需判断界面位于波头外、星形区或扇内。$u_*<0$ 的右侧五种情况完全对称。

十种情况的抽样值可概括为：正星速侧依次可能取 $\bm W_L,\bm W_{*L},\bm W_L,\bm W_{*L},\bm W_{L\mathrm{fan}}$；负星速侧相应取 $\bm W_R,\bm W_{*R},\bm W_R,\bm W_{*R},\bm W_{R\mathrm{fan}}$。确定原始状态后，转换为守恒量并代入物理通量函数。练习要求绘制完整流程图，并对等熵气体动力学列出所有界面波型。

图 6.5 的两大类首先按接触速度 $u_*$ 的符号划分，因为接触正好分隔左右星区；再按所处一侧的非线性波是激波还是稀疏波细分。稀疏波还要比较零射线与波头、波尾速度，因此有初态、星态、扇内态三种抽样位置。表 6.1 只列返回哪个区域，实际计算还需用第 4 章公式求该区域的 $\rho,u,p$。抽样器必须覆盖等号情形，避免波速接近零时因分支间隙返回未定义状态。

% [原PDF p243]
\subsection{时间步长}
时间步取
\begin{equation}
\Delta t=C_{\mathrm{CFL}}\frac{\Delta x}{S_{\max}^n},
\qquad0<C_{\mathrm{CFL}}\le1.\tag{6.17--6.18}
\end{equation}
$C_{\mathrm{CFL}}$ 越接近 1，推进越高效。可靠的一维 Euler 最大波速估计是遍历所有界面黎曼解的左右最快波：
\begin{equation}
S_{\max}^n=\max_i\{|S_{L,i+1/2}|,|S_{R,i+1/2}|\}.\tag{6.19}
\end{equation}
稀疏波取波头速度，激波取激波速度，并且必须包含边界产生的波。由于这些量在通量求解时已经得到，额外代价很小。

常见替代式 $S_{\max}^n=\max_i(|u_i^n|+a_i^n)$ 易推广到多维，但可能低估新生激波速度。例如静止激波管初值的 $u_i=0$，只剩声速贡献，初始几步尤其危险。

式（6.19）使用每个黎曼解的两条最外层波，因为接触速度总夹在它们之间。对左稀疏波，最快绝对速度可能是波头 $u_L-a_L$；对右稀疏波则是 $u_R+a_R$；激波速度由跃迁关系给出。虽然这些波速比简单的 $|u|+a$ 更可靠，仍应在强非理想状态方程或近真空中检查求解器输出。任何 NaN、负密度或负压力都不应被 CFL 逻辑静默跳过，否则可能产生虚假的“大时间步”。

% [原PDF p244]
若 $S_{\max}$ 估计可靠，可取 $C_{\mathrm{CFL}}$ 接近 1；实际常用 0.9。若使用可能低估的 $|u|+a$ 估计，则应更保守。多维计算虽常不得不采用后者，也必须监测稳定性。

\subsection{边界条件}
区间 $[0,L]$ 的首、末单元更新还需要边界通量 $\bm F_{1/2}$ 与 $\bm F_{M+1/2}$。可直接规定通量，也可建立虚单元 $I_0,I_{M+1}$，构造 $RP(\bm U_0^n,\bm U_1^n)$ 与 $RP(\bm U_M^n,\bm U_{M+1}^n)$，再用内部界面相同的 Godunov 通量算法。边界本质上是物理问题，数值实现必须谨慎。本章讨论反射壁和透射边界。

若初始激波管两侧静止，式（6.20）只看到初态声速，却看不到高压差形成后极快的激波。前几步减小 CFL 系数能让波系先建立，再使 $|u|+a$ 估计更接近真实最大速度，但这只是实用策略，不取代正确波速分析。多维中常把各坐标方向的谱半径组合进 CFL 条件，而不是逐面精确求全域最外波；第 16 章会给出相应规则。

% [原PDF p245]
固定不渗透右壁采用镜像原始状态
\begin{equation}
\rho_{M+1}^n=\rho_M^n,qquad u_{M+1}^n=-u_M^n,qquad
p_{M+1}^n=p_M^n.\tag{6.21}
\end{equation}
边界黎曼解在 $u_M>0$ 时为两道激波，在 $u_M\le0$ 时为两道稀疏波，两者都有 $u_*=0$，恰好满足固定壁无穿透条件；边界通量仅动量分量含星区压力 $p_*$。

若壁以给定速度 $u_{\mathrm{wall}}$ 移动，则
\begin{equation}
\rho_{M+1}^n=\rho_M^n,\quad
u_{M+1}^n=-u_M^n+2u_{\mathrm{wall}},\quad
p_{M+1}^n=p_M^n.\tag{6.22}
\end{equation}
黎曼解关于移动壁轨迹对称，接触波与壁重合，因而 $u_*=u_{\mathrm{wall}}$。若流体追不上外移壁则形成两稀疏波；若流体撞向壁则形成两激波。

镜像速度保证边界黎曼初值关于壁速度反对称。固定壁时若内部流体向右撞墙，镜像态向左撞墙，中央停滞区由两道激波夹住；若内部速度向左离壁，则两侧相背运动，形成两道稀疏波。移动壁把参考速度从 0 改为 $u_{\mathrm{wall}}$，在壁随动坐标中仍是同样对称问题。因而接触速度精确等于壁速，质量和能量通量满足移动边界运动学条件，动量通量由壁压给出。

% [原PDF p246]
令 $C_M=u_M-u_{\mathrm{wall}}$。两稀疏波情形 $u_M\le u_{\mathrm{wall}}$ 的壁面压力为
\begin{equation}
p_*=p_M\left[1+\frac{\gamma-1}{2}\frac{C_M}{a_M}\right]^{2\gamma/(\gamma-1)},\tag{6.23}
\end{equation}
两激波情形 $u_M>u_{\mathrm{wall}}$ 为
\begin{equation}
p_*=p_M+\frac{C_M}{2A_M}\left[C_M+\sqrt{C_M^2+4A_M(B_M+p_M)}\right],\tag{6.24}
\end{equation}
其中 $A_M=2/[(\gamma+1)\rho_M]$、$B_M=(\gamma-1)p_M/(\gamma+1)$，且两者均有 $u_*=u_{\mathrm{wall}}$。这些闭式边界解也可与近似黎曼求解器配合，用来避免近似器在强壁面问题上的误差。

透射边界令虚单元完全复制邻近内部单元：$\rho_{M+1}=\rho_M$、$u_{M+1}=u_M$、$p_{M+1}=p_M$。边界黎曼问题于是平凡，不产生有限强度的入射波。多维非反射边界则远复杂得多，通常称开放端、远场、辐射或透明边界条件。

两稀疏波公式中的括号必须保持正值；若壁退开过快使其趋于零，就接近真空形成条件。两激波压力式由对称压力函数化为二次方程后取物理解根。$A_M,B_M,C_M$ 的量纲分别使根号内各项一致，编程时可用它们减少重复运算。无论采用精确还是近似内部通量，壁面闭式解都可单独用于边界，因为边界压力对载荷和稳定性极敏感。

% [原PDF p247]
实现 Godunov 法所需的三项内容至此完备：界面通量、用于计算 $\Delta t$ 的最大波速，以及边界条件。边界黎曼问题产生的速度也必须纳入 $S_{\max}^n$。练习要求为一端固定不渗透、另一端透射的一维等截面管编写完整流程图。

\section{数值结果与讨论}
五个具有精确解的 Euler 黎曼问题用于比较 Godunov、Lax--Friedrichs 和 Richtmyer 方法。初值为在 $x=x_0$ 分隔的左右常状态 $\bm W_L=(\rho_L,u_L,p_L)^T$、$\bm W_R=(\rho_R,u_R,p_R)^T$，取 $\gamma=1.4$，区域 $[0,1]$，$M=100$，$C_{\mathrm{CFL}}=0.9$，边界透射。前五步把 CFL 系数再乘 0.2，以弥补 $\max(|u|+a)$ 初始波速估计可能偏小的问题。

五组状态依次为：
\begin{center}\small
\begin{tabular}{c*{6}{r}}
\toprule 测试&$\rho_L$&$u_L$&$p_L$&$\rho_R$&$u_R$&$p_R$\\\midrule
1&1&0.75&1&0.125&0&0.1\\2&1&-2&0.4&1&2&0.4\\
3&1&0&1000&1&0&0.01\\4&5.99924&19.5975&460.894&5.99242&-6.19633&46.095\\
5&1&-19.59745&1000&1&-19.59745&0.01\\\bottomrule
\end{tabular}
\end{center}

反射与透射条件只是两类基本示例。真实入流/出流边界应按入射特征数给定相应物理量，不能总把全部原始量复制或指定。特别在亚声速边界，一部分特征进入、一部分离开，边界状态需要内外信息共同确定。Godunov 框架的便利之处是最终仍可把处理后的边界状态与内部状态送入同一黎曼求解器。与此同时，边界波速必须参与时间步估计，否则强反射会在首个单元违反 CFL 条件。

% [原PDF p248]
测试 1 是修正的 Sod 问题，含右行激波、右行接触和左行跨声速稀疏波，可检验熵性质。测试 2 是“123 问题”，含两道对称稀疏波和静止平凡接触，星形区接近真空，适合检验低密度稳健性。测试 3 含强激波、接触和左稀疏波，用于考察强激波精度与稳健性。测试 4 的三道强间断全部右行。测试 5 在测试 3 上叠加均匀负背景速度，使接触几乎静止，专门检验缓慢接触波分辨率。

\subsection{Godunov 方法的结果}
测试 1 的激波被涂抹约 $4\Delta x$，即四个单元；对一阶格式这已属典型甚至较好结果，而且激波附近保持单调。接触间断在密度和内能图中展宽约 20 个单元，因为接触两侧特征与波平行，缺少激波两侧特征汇聚所提供的压缩机制。尽管厚度不同，激波和接触的平均位置均正确，这是守恒性的直接结果。

五个测试覆盖不同数值困难，而不是只比较一个 Sod 问题。测试 1 检查正常激波、接触和跨声速熵；测试 2 检查接近真空时正性和内能恢复；测试 3、4 检查极强波与多间断相互作用；测试 5 通过 Galilean 平移把接触速度调到近零，检查格式是否对缓慢线性波产生过量耗散。所有测试保持相同网格和 CFL，才便于把差异归因于波型与算法。

% [原PDF p249]
稀疏波整体近似良好，但波头和波尾存在导数间断，较难捕捉；跨声速扇内部还出现小的“熵瑕疵”。Godunov 法理论上满足熵条件，因此网格细化时该跳跃应趋于零，$M=500$ 的结果确实显示收敛趋势。

测试 2 的 $\rho,u,p$ 尚可，但比内能 $e=p/[(\gamma-1)\rho]$ 误差明显；压力和密度都接近零时，二者小误差会被比值放大。低密度问题足以使某些近似黎曼求解器下的 Godunov 法失败，Richtmyer 法在此测试也无法给出解。

测试 3 的激波压力比 $p_*/p_R=46000$，激波 Mach 数约 198。理想气体模型在如此强激波下未必真实，但作为数学稳健性测试仍有效。接触过度涂抹使激波后密度平台未完全达到，稀疏波尾附近速度还有轻微超调；Richtmyer 法同样失败。

激波宽度通常以从前态过渡到后态所跨单元数衡量。Godunov 法对移动激波并非精确零厚度，因为一步中的平均把界面移动量分摊到相邻单元；接触没有压缩机制，所以误差每步像平流扩散一样累积。熵瑕疵只在稀疏扇穿过声速点时出现，其幅度随网格缩小。用内能图检查低密度解很有价值：压力曲线可能看似平滑正确，而 $p/\rho$ 会把相关误差暴露出来。

% [原PDF p250]
测试 4 的左右激波和接触均向右传播，来自两道反向强激波的相互作用。最快右激波约展宽 5 个单元；左侧缓慢激波只占约 2 个单元，却在内能中伴随低频非物理振荡。Godunov 法能精确保持孤立静止激波和接触（前提是黎曼求解器无缺陷），但缓慢移动激波的振荡仍是困难问题。

测试 5 的接触几乎静止，Godunov 法把它解析得非常尖锐，明显优于 Lax--Friedrichs。该问题在后续检验其他方法时仍将使用。

\subsection{其他方法的结果}
Lax--Friedrichs 在测试 1 中出现相邻单元配对，激波约占 8 个单元，稀疏波和接触都很差。测试 2 的内能误差尤大；测试 3 无法达到激波后密度；测试 4 也明显不如 Godunov。测试 5 最清楚地区分中心格式与 Godunov 型格式：线性接触波的保持是当今数值方法最困难的任务之一。

测试 3 的强度超出许多真实理想气体应用范围，但数值测试并不要求模型在该极限完全物理，只要求方程与精确解定义明确。一个稳健方法应在这种高压比下保持密度和压力为正并正确传播激波。测试 4 中缓慢激波的低频振荡不同于 Lax--Wendroff 的高频 Gibbs 振铃，它甚至会出现在精确 Godunov 通量中，与激波相对网格的缓慢移动有关。已有若干修正方案，但作者指出该问题尚未彻底解决。

% [原PDF p251]
二阶 Richtmyer 法仅在测试 1 成功：光滑区比一阶方法准确，间断更尖锐，但附近出现非物理振荡；测试 2--5 均计算失败。本章使用的是精确黎曼求解器，第 9--12 章将以近似求解器降低成本，第 13、14、16 章则提高时空精度。

\begin{sourceblock}
\sourcecrop{251}{45}{190}{395}{400}{0.30}
\captionof{figure}{Godunov 法用于测试 1：密度、速度、压力与比内能。符号为数值解，实线为精确解；$x_0=0.3$，输出时刻 $t=0.2$。}
\end{sourceblock}
图中 Density、Velocity、Pressure、Internal energy 分别表示密度、速度、压力和比内能，Position 表示位置。

% [原PDF p252]
\begin{sourceblock}
\sourcecrop{252}{45}{385}{395}{610}{0.28}\\[0.45em]
\sourcecrop{252}{45}{115}{395}{337}{0.28}
\captionof{figure}{Godunov 法的测试 2（上）与测试 3（下）。测试 2：$x_0=0.5,t=0.15$；测试 3：$x_0=0.5,t=0.012$。}
\end{sourceblock}
测试 2 的近真空星形区使内能比值异常敏感，但基本物理量仍保持正且轮廓合理；测试 3 中极强右激波仍稳定传播，接触面则因一阶耗散而显著展宽。

% [原PDF p253]
\begin{sourceblock}
\sourcecrop{253}{45}{385}{395}{610}{0.28}\\[0.45em]
\sourcecrop{253}{45}{115}{395}{337}{0.28}
\captionof{figure}{Godunov 法的测试 4（上）与测试 5（下）。测试 4：$x_0=0.4,t=0.035$；测试 5：$x_0=0.8,t=0.012$。}
\end{sourceblock}
测试 4 的三道右行强间断彼此靠近，慢激波尖锐但有轻微振荡；测试 5 的几乎静止接触在密度和内能中保持为窄跃迁，显示精确 Godunov 通量对零速线性波的优势。

% [原PDF p254]
\begin{sourceblock}
\sourcecrop{254}{45}{385}{395}{610}{0.28}\\[0.45em]
\sourcecrop{254}{45}{115}{395}{337}{0.28}
\captionof{figure}{Lax--Friedrichs 法的测试 1（上）与测试 2（下），参数与 Godunov 对照相同。}
\end{sourceblock}
中心耗散使所有波型更宽，并出现相邻点配对。测试 2 中压力与密度的小偏差被比内能计算放大，说明只看压力曲线会高估数值解质量。

% [原PDF p255]
\begin{sourceblock}
\sourcecrop{255}{45}{385}{395}{610}{0.28}\\[0.45em]
\sourcecrop{255}{45}{115}{395}{337}{0.28}
\captionof{figure}{Lax--Friedrichs 法的测试 3（上）与测试 4（下）。}
\end{sourceblock}
测试 3 的强激波后平台没有恢复到精确值；测试 4 的多道强间断均受到严重涂抹。虽然格式简单、稳健且无明显高频振荡，它的人工黏性会损失关键波结构。

% [原PDF p256]
\begin{sourceblock}
\sourcecrop{256}{45}{385}{395}{610}{0.28}\\[0.45em]
\sourcecrop{256}{45}{115}{395}{337}{0.28}
\captionof{figure}{上：Lax--Friedrichs 法用于测试 5；下：Richtmyer 法用于测试 1。}
\end{sourceblock}
Lax--Friedrichs 无法保持测试 5 的近静止接触；Richtmyer 在测试 1 的光滑区更准确、间断更窄，却在跃迁附近产生色散振荡。这是中心一阶耗散格式与无非线性限制二阶格式的典型取舍。

% [原PDF p257]
\begin{sourceblock}
\sourcecrop{257}{45}{385}{395}{610}{0.29}
\captionof{figure}{Godunov 法在测试 1 上使用细网格 $M=500$ 的结果；符号为数值解，实线为精确解，$t=0.2$。}
\end{sourceblock}
网格由 100 单元细化到 500 单元后，激波、接触与稀疏波端点都显著变尖，跨声速稀疏波内的熵瑕疵也减小。这验证了该一阶守恒方法虽在粗网格上耗散较大，但随 $\Delta x\to0$ 向熵解收敛。

\section*{本章方法的实现顺序与结果解读}
把本章公式组织成可执行算法时，可按以下逻辑理解。首先由给定初值计算每个单元的守恒平均 $\bm U_i^0$。Euler 方程常从原始量开始给定，因此要用
\[
\bm U=(\rho,\rho u,E)^T,\qquad
E=\frac{p}{\gamma-1}+\frac12\rho u^2
\]
转换。每个时间层开始时，先根据物理边界构造虚单元：透射边界复制状态，固定反射壁反转法向速度，移动壁则关于壁速镜像。边界状态必须物理可接受，密度和压力均为正。

随后对所有界面建立 $RP(\bm U_i^n,\bm U_{i+1}^n)$。精确 Euler 黎曼求解器通常先解非线性压力方程得到 $p_*$，再求 $u_*$，判断左右波是激波还是稀疏波，计算各自外层速度。为了形成 Godunov 通量，并不需要输出整条空间曲线，只需令 $S=0$ 调用抽样器。抽样器先比较 $S$ 与 $u_*$ 判断接触哪一侧，再比较相应激波速度或稀疏波头、波尾。返回 $(\rho_0,u_0,p_0)$ 后，形成
\[
\bm F_{i+1/2}=\left(\rho_0u_0,;\rho_0u_0^2+p_0,;u_0(E_0+p_0)\right)^T.
\]
同一界面结果必须共享给两侧单元。

最大波速最好在上述黎曼循环中同时累积。对激波记录其 $|S|$，对稀疏波记录外侧波头绝对速度；取全部界面（包括边界）的最大值后，用式（6.17）选 $\Delta t$。如果通量计算必须先知道时间步，可把黎曼结构计算和通量抽样分成两遍，或在第一遍缓存星区及波速。精确 Godunov 通量本身不显含 $\Delta t$，所以也可先求通量与最大速度，再推进。到达指定输出时刻前缩短最后一步。

更新阶段对每个物理单元使用
\[
\bm U_i^{n+1}=\bm U_i^n-\frac{\Delta t}{\Delta x}
(\bm F_{i+1/2}-\bm F_{i-1/2}).
\]
更新后由新守恒量恢复原始量用于诊断和下一轮黎曼求解。对理想气体，$\rho>0$ 且
\[
p=(\gamma-1)\left(E-\frac{(\rho u)^2}{2\rho}\right)>0
\]
是状态物理性的基本要求。若出现非法状态，应首先检查 CFL、边界状态、黎曼分支和通量索引，而不能用任意截断掩盖错误。

图 6.8--6.19 的四个纵向变量各有不同诊断作用。密度最清楚地显示接触间断，因为跨接触压力和速度连续而密度可跳跃；压力最清楚地显示激波，却往往是最容易“看起来正确”的变量；速度能揭示稀疏波尾部超调；比内能 $e=p/[(\gamma-1)\rho]$ 在低密度区放大小误差，因此是严格质量指标。只画压力会漏掉接触涂抹与低密度能量误差。

测试 1 的结构按从左到右依次为左稀疏波、接触、右激波。Godunov 法的激波速度和平均位置正确，因为离散守恒保证 Rankine--Hugoniot 收支；激波有限厚度来自单元平均。接触比激波宽得多，并非黎曼求解不精确，而是一阶平流数值黏性累积。跨声速稀疏波中的熵瑕疵会随网格细化而减小，因此应通过收敛序列判断，而不是仅凭一张粗网格图断言错误熵解。

测试 2 的两个稀疏波向外传播，中间密度、压力接近真空。这里 Newton 压力迭代、声速计算、守恒量恢复都最容易受舍入误差影响。精确黎曼 Godunov 仍能给出正状态，但某些近似求解器可能产生负密度。测试 3 的极端压力比则检验强激波波速、后态压缩与能量通量；测试 4 同时考验多个强间断和缓慢激波振荡；测试 5 通过静止接触检验数值通量对线性退化波的分辨率。

与 Lax--Friedrichs 相比，Godunov 通量把耗散按真实波结构分配，所以激波、接触均较尖锐，特别能精确保留孤立零速间断。Lax--Friedrichs 的耗散尺度由统一的 $\Delta x/\Delta t$ 决定，对所有特征场一视同仁，因此简单稳健却过度扩散。Richtmyer 的二阶预测在光滑区精度高，但未加限制时会在间断附近产生振荡，并可能在强波或近真空中给出非物理中间状态而失败。

基本 Godunov 法的价值更多在于建立基准与架构，而非直接作为所有工程计算的最终格式。后续近似黎曼求解器保留“局部波决定界面通量”的思想，降低精确压力迭代成本；高阶重构则用单元内斜率或多项式替换分片常数数据。在任何扩展中，守恒通量共享、物理边界、可靠波速与正状态检查仍是不可省略的骨架。

% [原PDF p258]
\chapter{随机选择法及相关方法}
\section{引言}
Glimm 于 1965 年在一类非线性双曲守恒律解的构造性存在证明中提出随机选择法（Random Choice Method, RCM）；Chorin 于 1976 年把其修改成可计算的气体动力学 Euler 方程算法。实现 RCM 需要两项核心工具：局部黎曼问题的精确解，以及从这些解中选取下一时间层状态的随机抽样过程。

RCM 与 Godunov 法都求解局部黎曼问题，差别在于如何利用解：Godunov 法对局部解作积分平均，RCM 则从解中随机选择一个状态。随机序列的统计性质会显著影响精度。Chorin 将方法推广到燃烧，Sod 处理柱对称和球对称 Euler 方程并引入代数源项，Colella 系统研究其优缺点并改进随机抽样；随后该法还用于油藏两相流、浅水方程、水下爆炸、定常超声速流、激波聚焦与衍射以及反应流。

% [原PDF p259]
RCM 可直接用于任意维标量问题和只有两个独立变量的非线性系统，例如一维非定常 Euler、二维定常超声速 Euler 和一维浅水方程。借助第 15 章的算子分裂，还能加入代数源项乃至黏性扩散项；但分裂法通常无法把 RCM 的核心优点成功推广到多维非线性系统。

RCM 的吸引力在于能把激波和材料界面保持为真正的零厚度间断，而多数方法会把它们涂抹若干单元，接触面尤其严重。代价是间断位置带随机误差，光滑稀疏波也含随机噪声；有源项时这种噪声可能被放大。

本章首先介绍常规非交错 RCM，又给出一种类似 Lax--Friedrichs 的新随机选择格式，以及从交错网格 RCM 重新解释而来的确定性一阶中心 FORCE 格式。以下均以一维非定常 Euler 方程说明，学习前宜复习第 4 章精确黎曼解。

\section{非交错网格上的 RCM}

% [原PDF p260]
考虑一般初边值问题
\begin{equation}
\bm U_t+\bm F(\bm U)_x=0,\quad
\bm U(x,0)=\bm U^{(0)}(x),\quad
\bm U(0,t)=\bm U_l(t),\quad \bm U(L,t)=\bm U_r(t).\tag{7.1}
\end{equation}
RCM 只有在黎曼问题的激波 Rankine--Hugoniot 条件中必须使用守恒变量；其他步骤用原始变量更方便。Euler 方程取 $\bm W=(\rho,u,p)^T$。

\subsection{非线性系统的格式}
把 $[0,L]$ 分成 $M$ 个规则单元，
\begin{equation}
x_{i-1/2}=(i-1)\Delta x,\quad x_i=(i-\tfrac12)\Delta x,
\quad x_{i+1/2}=i\Delta x.\tag{7.2}
\end{equation}
非交错网格每一步都在同一组单元中心更新；时间步因波速变化而可逐步改变。

在 RCM 中使用原始变量并不会破坏激波条件，因为真正求黎曼解时仍由守恒律和 Rankine--Hugoniot 关系决定波后状态；只是抽到一个完整物理状态后，以 $(\rho,u,p)$ 存储和传递更方便。与有限体积平均不同，RCM 单元值可理解为被选中的点状态，因此不必强制先计算单元积分平均。这个区别也解释它为什么能保持零厚度间断，却不具备逐步严格的离散总量守恒。

% [原PDF p261]
给定 $t^n$ 的数据，令每个单元内 $\bm W(x,t^n)=\bm W_i^n$ 为常数。相邻状态定义局部黎曼问题 $RP(\bm W_i^n,\bm W_{i+1}^n)$，其自相似解为 $\bm W_{i+1/2}(x/t)$。第 4 章的抽样程序可在任意给定 $S=x/t$ 返回密度、速度与压力。

在时间 $t^*$，黎曼解通常由两道非线性波和中间接触组成。RCM 不对整个抽样区积分，而是在区间 $[x_l,x_r]$ 中随机选择 $x^*$，调用精确抽样器求 $\bm W(x^*/t^*)$。为了从 $\bm W_i^n$ 推进到 $\bm W_i^{n+1}$，需要同时使用单元左、右两个界面的黎曼解。

图 7.2 中每一对相邻常状态都产生独立局部波系。抽样区必须位于一步内不发生相邻波系相互作用的范围中，否则某一点的真实状态不再由单个局部黎曼解决定。精确抽样程序输入左右状态和射线速度 $S$，内部先求 $p_*,u_*$，再按接触侧、激波或稀疏波以及波头波尾位置返回状态。RCM 的随机性只来自射线选择，黎曼问题本身仍以确定方式精确求解。

% [原PDF p262]
算法分两步：
\begin{enumerate}
\item 求解 $RP(\bm W_{i-1}^n,\bm W_i^n)$ 与 $RP(\bm W_i^n,\bm W_{i+1}^n)$，得到 $\bm W_{i-1/2}(x/t)$、$\bm W_{i+1/2}(x/t)$。
\item 在局部时间 $\Delta t$ 于单元 $I_i$ 内随机抽样。给定 $0\le\theta_n\le1$，更新为
\end{enumerate}
\begin{equation}
\bm W_i^{n+1}=\begin{cases}
\bm W_{i-1/2}(\theta_n\Delta x/\Delta t),&0\le\theta_n\le1/2,\\
\bm W_{i+1/2}((\theta_n-1)\Delta x/\Delta t),&1/2<\theta_n\le1.
\end{cases}\tag{7.3}
\end{equation}
前一分支在左黎曼解的右半部抽样，后一分支在右黎曼解的左半部抽样。所选状态直接赋给整个单元的下一时间层。

式（7.3）的两个抽样坐标分别覆盖单元中心左右半区：第一支 $S\in[0,\Delta x/(2\Delta t)]$，第二支 $S\in[-\Delta x/(2\Delta t),0)$。使用同一个 $\theta_n$ 更新所有单元会让全域抽样线在该时间层保持一致，减少单元间无结构噪声；若每个单元独立抽数，局部随机误差更强。时间层之间仍使用序列的不同项，从长期统计上再现积分平均。

% [原PDF p263]
当 $0\le\theta_n\le1/2$ 时，抽样射线为 $S=\theta_n\Delta x/\Delta t$，落在左界面黎曼问题右侧；$1/2<\theta_n\le1$ 时则取右界面左侧。程序可按多种顺序组织“求解黎曼问题”和“抽样”，但 Chorin 指出每个完整时间层只需一个共同随机数 $\theta_n$；Glimm 的存在性证明则允许每步每单元一个随机数。

与 Godunov 法对比，后者在同一区域取守恒量的积分平均
\begin{equation}
\bm U_i^{n+1}=\frac1{\Delta x}\int_0^{\Delta x/2}\bm U_{i-1/2}(x/\Delta t)dx
+\frac1{\Delta x}\int_{-\Delta x/2}^{0}\bm U_{i+1/2}(x/\Delta t)dx.\tag{7.4}
\end{equation}
RCM 不是逐步严格守恒的，但可视为统计意义上的守恒方法。

Godunov 法的积分平均是确定性的，并严格保留单元总量，但每次平均都会把移动间断分摊到部分单元，形成数值扩散。RCM 选择某一侧完整状态，间断从一个界面跳到另一个界面，永不产生中间混合值；它以随机位置误差换取无限锐度。在许多次抽样或统计期望意义下，被选中左右状态的比例接近它们在积分区间所占比例，因此可以说“统计守恒”，但单次实现的全域总量会有小幅随机波动。

% [原PDF p264]
\subsection{边界条件与时间步长}
首、末单元的更新需要虚状态。透射边界取
\begin{equation}
(\rho_0,u_0,p_0)=(\rho_1,u_1,p_1),\qquad
(\rho_{M+1},u_{M+1},p_{M+1})=(\rho_M,u_M,p_M).\tag{7.5}
\end{equation}
若左右反射壁速度分别为 $u_{wl},u_{wr}$，则
\begin{equation}
u_0=-u_1+2u_{wl},\qquad u_{M+1}=-u_M+2u_{wr},\tag{7.6}
\end{equation}
密度、压力镜像复制。如此可像内部一样求边界黎曼解并随机抽样。

为避免一个单元内左右波系相互作用，时间步必须满足
\begin{equation}
\Delta t\le\frac{\Delta x}{2S_{\max}^n}
=C_{\mathrm{CFL}}\frac{\Delta x}{S_{\max}^n},\qquad
0<C_{\mathrm{CFL}}\le\frac12.\tag{7.7--7.9}
\end{equation}

移动反射壁条件与第 6 章完全相同：速度关于壁速镜像，密度和压力复制。若只在一侧有壁，另一侧可使用透射或按物理入流条件设置。由于 RCM 更新可能在左右两个黎曼解之间选择，首单元必须同时具备左边界问题与右内部问题，末单元同理。虚单元使这一逻辑不需特殊分支，但必须在每个新时间层抽样前重新赋值。

% [原PDF p265]
因此常规 RCM 的稳定上限只有 Godunov 法的一半。最大波速最好直接取全部局部黎曼解中的真实最快波。经验表明 RCM 对低估 $S_{\max}$ 不像确定性格式那样立即表现为爆炸式不稳定，而更可能让计算波以明显错误的速度传播；这仍然是不允许的误差。

\section{Lax--Friedrichs 型随机选择格式}
这一路线从 Lax--Friedrichs 格式的黎曼解积分平均解释出发。若用守恒积分式变换平均，就恢复普通确定性 Lax--Friedrichs 格式并消去黎曼问题；若保留黎曼解而把平均解释为随机抽样，则得到 Lax--Friedrichs 型 RCM。

\subsection{Lax--Friedrichs 格式回顾}
线性平流方程的格式
\begin{equation}
u_i^{n+1}=\frac{1+c}{2}u_{i-1}^n+\frac{1-c}{2}u_{i+1}^n\tag{7.11}
\end{equation}
等价于在单元 $i$ 内对 $RP(u_{i-1}^n,u_{i+1}^n)$ 于时间 $\Delta t/2$ 的解作积分平均。

常规 RCM 的 $C_{\mathrm{CFL}}\le1/2$ 不是通过 Fourier 稳定性得出，而是保证抽样几何有效：最快波从任一界面出发不能超过半个单元，否则左右波系可能在抽样线前相遇。低估波速时，算法仍会从某个局部黎曼解返回有限状态，所以未必爆炸，但所选解已不代表真实局部演化，尤其会让激波以错误频率跳格。显然的错速和隐蔽的统计偏差都要求使用可靠的界面波速。

% [原PDF p266]
即
\begin{equation}
u_i^{n+1}=\frac1{\Delta x}\int_{x_{i-1/2}}^{x_{i+1/2}}
\widehat u_i(x,\Delta t/2)dx,\tag{7.12}
\end{equation}
其中
\begin{equation}
\widehat u_i(x/t)=\begin{cases}u_{i-1}^n,&x/t<a,\\u_{i+1}^n,&x/t>a.
\end{cases}\tag{7.13}
\end{equation}
这一解释有两个特殊点：黎曼初值使用 $i-1$ 与 $i+1$，不含中心值 $u_i^n$；平均发生在半个完整时间步。上游状态在积分区间中占的长度更大，因而权重更高。

\subsection{格式}
对非线性系统，把式（7.12）推广为
\begin{equation}
\bm U_i^{n+1}=\frac1{\Delta x}\int_{x_{i-1/2}}^{x_{i+1/2}}
\widehat{\bm U}_i(x,\Delta t/2)dx,\tag{7.15}
\end{equation}
其中 $\widehat{\bm U}_i$ 是 $RP(\bm U_{i-1}^n,\bm U_{i+1}^n)$ 的解。

图 7.6 中，时间 $\Delta t/2$ 的间断位于 $a\Delta t/2=(c/2)\Delta x$。它把积分区间分成长为 $(1+c)\Delta x/2$ 的左状态部分和 $(1-c)\Delta x/2$ 的右状态部分，除以 $\Delta x$ 就得到式（7.11）的两个权重。这一几何推导还表明稳定要求间断仍留在单元内，即 $|c|\le1$。虽然中心值未参与，本次更新仍是左右两状态的凸组合。

% [原PDF p267]
有三种处理方式：（1）精确求黎曼解并直接积分，结果与普通 Lax--Friedrichs 几乎不可区分，却更昂贵；（2）用守恒积分式消去黎曼解，恢复普通格式；（3）保留黎曼解并随机化积分。第三种取 $-1/2\le\theta_n\le1/2$，更新为
\begin{equation}
\boxed{\bm U_i^{n+1}=\widehat{\bm U}_i(\theta_n\Delta x,\Delta t/2).}\tag{7.18}
\end{equation}
它在半时间处对跨越整个单元的黎曼解随机抽样。稳定范围为
\begin{equation}0<C_{\mathrm{CFL}}\le1,\tag{7.19}\end{equation}
是常规 RCM 的两倍。实现只需用隔一个单元的左右状态建立黎曼问题，并把随机抽样时间设为 $\Delta t/2$。

随机式（7.18）可视为用一个样本估计积分平均：被抽状态是 $\widehat{\bm U}_i$ 在位置 $x=\theta_n\Delta x$、时间 $\Delta t/2$ 的值。若 $\theta_n$ 在 $[-1/2,1/2]$ 均匀分布，长期期望等于空间平均。因为波只需在半时间内不跑出半单元，所以条件变成 $S_{\max}(\Delta t/2)\le\Delta x/2$，即 $C_{\mathrm{CFL}}\le1$。稳定区扩大不是凭空获得，而来自对半时间黎曼解的不同抽样几何。

% [原PDF p268]
\section{交错网格上的 RCM}
\subsection{非线性系统的格式}
交错版本通过两个半步骤把 $\bm U_i^n$ 更新到 $\bm U_i^{n+1}$。第一步求 $RP(\bm U_{i-1}^n,\bm U_i^n)$ 与 $RP(\bm U_i^n,\bm U_{i+1}^n)$，在稳定时间 $\Delta t^{n+1/2}$ 随机抽样，得到交错位置状态
\begin{equation}
\bm U_{i\pm1/2}^{n+1/2}=\widehat{\bm U}_{i\pm1/2}^{n+1/2}
(\theta_n\Delta x,\Delta t^{n+1/2}).\tag{7.21}
\end{equation}
第二步以这两个交错状态建立 $RP(\bm U_{i-1/2}^{n+1/2},\bm U_{i+1/2}^{n+1/2})$，再以新随机数 $\theta_{n+1}$ 抽样，得到中心状态 $\bm U_i^{n+1}$。两个半步可使用不同稳定时间步，但都须满足 RCM 的 CFL 上限。

\subsection{确定性一阶中心 FORCE 格式}
把两次随机抽样换成黎曼解的积分平均，并令两个半步时间均为 $\Delta t/2$，可导出确定性、守恒、一阶且中心的 FORCE 格式。

交错算法第一步把整数中心的数据转移到半整数位置，第二步再从半整数位置转回整数中心。每次都解相邻状态黎曼问题并随机抽样，因此一个完整步需要两轮黎曼求解。若两半步分别重新计算最大速度，$\Delta t^{n+1/2}$ 与 $\Delta t^{n+1}$ 可以不同；为推导 FORCE，作者把它们都设成总步长的一半。此后改用守恒变量，是因为积分平均必须直接保持守恒量。

% [原PDF p269]
第一阶段的两个交错平均为
\begin{align}
\bm U_{i-1/2}^{n+1/2}&=\frac12(\bm U_{i-1}^n+\bm U_i^n)
+\frac{\Delta t}{2\Delta x}(\bm F_{i-1}^n-\bm F_i^n),\tag{7.25}\\
\bm U_{i+1/2}^{n+1/2}&=\frac12(\bm U_i^n+\bm U_{i+1}^n)
+\frac{\Delta t}{2\Delta x}(\bm F_i^n-\bm F_{i+1}^n).\tag{7.26}
\end{align}
第二阶段对这两个状态所定义黎曼解再作平均。应用积分守恒式得到
\begin{equation}
\bm U_i^{n+1}=\frac12(\bm U_{i-1/2}^{n+1/2}+\bm U_{i+1/2}^{n+1/2})
+\frac{\Delta t}{2\Delta x}(\bm F_{i-1/2}^{n+1/2}-\bm F_{i+1/2}^{n+1/2}).\tag{7.28}
\end{equation}
半时间通量 $\bm F(\bm U_{i+1/2}^{n+1/2})$ 正是 Richtmyer 通量。

式（7.25）、（7.26）其实就是对相邻两状态在半时间内的黎曼解应用控制体守恒，不需要真正求解黎曼问题。第二阶段的式（7.28）同理。把三次平均代数合并后，所有中间黎曼结构都消失，只剩左右状态和物理通量。FORCE 虽由随机选择法的确定性替代推导而来，最终却完全是中心有限体积格式，这也是名称中“centred”的由来。

% [原PDF p270]
整理为标准守恒形式
\begin{equation}
\bm U_i^{n+1}=\bm U_i^n+\frac{\Delta t}{\Delta x}
(\bm F_{i-1/2}^{\mathrm{FORCE}}-\bm F_{i+1/2}^{\mathrm{FORCE}}),\tag{7.30}
\end{equation}
可得一个出人意料而简洁的结果：
\begin{equation}
\boxed{\bm F_{i+1/2}^{\mathrm{FORCE}}=
\frac12\left(\bm F_{i+1/2}^{\mathrm{RI}}+
\bm F_{i+1/2}^{\mathrm{LF}}\right).}\tag{7.32}
\end{equation}
即 FORCE（First-Order Centred，一阶中心）通量就是 Richtmyer 通量和 Lax--Friedrichs 通量的算术平均。它无需黎曼求解器，结构简单而且保持守恒。

\subsection{FORCE 格式分析}
对线性平流方程，FORCE 更新为
\begin{equation}u_i^{n+1}=b_{-1}u_{i-1}^n+b_0u_i^n+b_1u_{i+1}^n,\tag{7.33}\end{equation}
其中 $b_{-1}=(1+c)^2/4$、$b_0=(1-c^2)/2$、$b_1=(1-c)^2/4$。

Richtmyer 通量由预测的半时间界面状态求物理通量，偏二阶而色散；Lax--Friedrichs 通量含较强扩散。二者平均后，FORCE 用 LF 的耗散抑制 RI 振荡，又以 RI 成分把 LF 的黏性减半。它仍是一阶格式，因为在间断与一般非线性问题上主导截断误差为 $O(\Delta x)$。无需特征分解使其适合复杂系统，但代价是不能像 Godunov 通量那样精确识别接触和剪切波。

% [原PDF p271]
\textbf{命题 7.1}\quad FORCE 格式在 $0\le|c|\le1$ 时稳定并单调，修正方程为
\begin{equation}
q_t+a q_x=\alpha_{\mathrm{FO}}q_{xx},\qquad
\alpha_{\mathrm{FO}}=\frac14a\Delta x\frac{1-c^2}{c}
=\frac12\alpha_{\mathrm{LF}}.\tag{7.37}
\end{equation}
稳定性可由 von Neumann 分析证明，也可用 Billett 对三点格式的判据：$b_0(b_{-1}+b_1)\ge0$ 且 $b_0(b_{-1}+b_1)+4b_{-1}b_1\ge0$。在上述 Courant 范围内三个系数全部非负，因此格式单调。其人工黏性恰为 Lax--Friedrichs 的一半，预示在相同成本下分辨率更好。高阶扩展将在第 13、14 章介绍，多维推广可结合第 16 章分裂技术。

\section{随机数}
RCM 解的质量关键取决于随机序列 $\{\theta_n\}$。研究表明，序列越接近通常意义上的“完全随机”，计算结果反而可能越差；Colella 因此引入低差异的 van der Corput 伪随机序列。

三点系数在 $|c|\le1$ 时显然非负且和为 1，所以 FORCE 更新是凸组合，不能产生新极值。Billett 判据给出与 von Neumann 分析相同的稳定范围，说明这里单调性与线性稳定性一致。人工黏性 $\alpha_{\mathrm{FO}}$ 在 $|c|=1$ 时消失，$c$ 很小时则因 $\Delta t$ 过小而相对增大；这也是固定网格上盲目使用极小 CFL 不一定提高一阶显式格式质量的例子。

% [原PDF p272]
\subsection{van der Corput 伪随机数}
取互素整数 $k_1>k_2>0$，$(k_1,k_2)$ van der Corput 序列定义为
\begin{align}
\theta_n&=\sum_{i=0}^{m}A_i k_1^{-(i+1)},\tag{7.38}\\
A_i&=k_2a_i\pmod{k_1},\tag{7.39}\\
n&=\sum_{i=0}^{m}a_i k_1^i.\tag{7.40}
\end{align}
先把非负整数 $n$ 按基数 $k_1$ 展开得到数字 $a_i$；再将 $k_2a_i$ 除以 $k_1$ 的余数作为 $A_i$；最后把这些数字按负幂“反写到小数点后”。例如 $k_1=2$ 时是二进制反位序；$n=3=1\times2^0+1\times2^1$，故 $(2,1)$ 序列的 $\theta_3=1/2+1/4=0.75$。

练习要求求 $(2,1)$ 与 $(5,3)$ 序列前十项。$(2,1)$ 的前十项为
\[
\begin{gathered}
0.5,\ 0.25,\ 0.75,\ 0.125,\ 0.625,\\
0.375,\ 0.875,\ 0.0625,\ 0.5625,\ 0.3125.
\end{gathered}
\]

反位序构造使序列逐层填补区间中尚未覆盖的位置。以二进制序列为例，前两项落在 $1/2,1/4$，第三项补到 $3/4$，随后依次补八分点；任意 $2^m$ 长的前缀都近似均匀覆盖。一般 $(k_1,k_2)$ 中，$k_2$ 对基数数字作模乘排列，改变访问子区间的次序而不破坏覆盖。这里称“伪随机”强调序列完全确定、可复现，却具有抽样所需的均匀分布特征。

% [原PDF p273]
$(5,3)$ 序列前十项为 $0.6,0.2,0.8,0.4,0.12,0.72,0.32,0.92,0.52,0.04$。这些数并非独立随机样本，而是有意在 $[0,1]$ 内均匀填充空隙。

\subsection{统计性质}
对 $N$ 个数，算术平均与样本标准差定义为
\begin{equation}
x_{\mathrm{ar}}=\frac1N\sum_{n=1}^N\theta_n,qquad
x_{\mathrm{sd}}=\left[\frac1{N-1}\sum_{n=1}^N(\theta_n-x_{\mathrm{ar}})^2\right]^{1/2}.\tag{7.41--7.42}
\end{equation}
均匀序列应有 $x_{\mathrm{ar}}\approx1/2$。把 $[0,1]$ 分成 $D$ 个子区间 $R_i$，记其中实际落点数为 $c(R_i)$、理论概率为 $p_i$，卡方统计量为
\begin{equation}
x_{\mathrm{sq}}^2=\sum_{i=1}^{D}\frac{[c(R_i)-Np_i]^2}{Np_i}.\tag{7.43}
\end{equation}
均匀覆盖越好，此量越小，RCM 结果通常也越好。

标准差只衡量样本围绕平均的总体散布，许多排列不同的序列会得到同一值，无法反映区间空洞和聚集。卡方统计先把区间分箱，再比较每箱实际计数与均匀期望 $N/D$，因而更敏感于覆盖质量。式（7.43）原文求和上限应理解为子区间数 $D$，不是随机数个数 $N$。若每箱计数完全相同，$x_{\mathrm{sq}}^2=0$；数值越大，偏离均匀分布越严重。

% [原PDF p274]
对 200 个数、20 个等宽子区间的比较表明，六种 van der Corput 序列的平均约 $0.494$--$0.500$，标准差约 $0.288$--$0.289$；主要差别在 $x_{\mathrm{sq}}^2$。$(5,3)$ 序列的值为 0.2，优于多数比较序列的 1.0 或 1.4。Olivier 与 Grönig 的结论是：平均和标准差相近时，卡方均匀度才是影响 RCM 的关键统计量。Anderson--Gottlieb 序列也有相似品质，其示例统计量为 $x_{\mathrm{ar}}=0.4981$、$x_{\mathrm{sd}}=0.2891$、$x_{\mathrm{sq}}^2=0.6$。

\subsection{单一激波的传播}
考虑正速度 $S$ 的单激波连接常状态 $\bm U_L,\bm U_R$：
\begin{equation}
\bm U(x,t)=\begin{cases}\bm U_L,&x/t<S,\\\bm U_R,&x/t>S.
\end{cases}\tag{7.45}
\end{equation}
在非交错 RCM 中取基于激波速度的 Courant 数 $c\le1/2$，激波始终保持为真间断，但其离散位置由随机比较决定。

表 7.3 的六个序列平均都接近 $1/2$，标准差也近似均匀分布的理论值 $1/\sqrt{12}\approx0.288675$，仅卡方值有明显区分。$(5,3)$ 的 0.2 表示 200 个数在 20 个箱中的计数极均匀，这与其 RCM 计算通常较好相符。作者的标准差结果与 Olivier--Grönig 文献所报值存在差异，但两者关于卡方指标最有判别力的结论一致。

% [原PDF p275]
假设时间层开始时激波位于界面 $x_{i-1/2}$，一步后真实位置向右移动 $c\Delta x$。随机点为 $\theta_n\Delta x$：若 $\theta_n\le c$，抽到激波后状态 $\bm U_L$，整个单元被更新，激波离散位置右移一个 $\Delta x$；若 $\theta_n>c$，抽到 $\bm U_R$，激波位置不动。因此 $n$ 步后
\begin{equation}
X_s=\sum_{j=1}^{n}\Delta x P_j,qquad
P_j=\begin{cases}1,&\theta_j\le c,\\0,&\theta_j>c.
\end{cases}\tag{7.46--7.47}
\end{equation}
表 7.4 手算了 $(2,1)$、$(5,3)$ 序列在 $c=1/4,1/2$ 时前十步的位置误差。$(5,3)$ 通常更准；$c=1/2$ 时每隔一步甚至精确一次。这显示低差异序列如何控制随机位置误差。

单激波例子把 RCM 位置误差化为简单的 Bernoulli 累加。精确位移每步是 $c\Delta x$，离散位移则是 0 或 $\Delta x$；第 $n$ 步位置误差等于 $\Delta x(\sum_{j=1}^nP_j-nc)$。低差异序列保证任意较长前缀中 $\theta_j\le c$ 的比例接近 $c$，从而使累计误差有界在少数网格宽度，而普通独立随机数的偏差通常按 $\sqrt n$ 增长。间断锐度无限并不等于位置完全精确，二者必须分别评价。

% [原PDF p276]
练习要求对 $(3,2)$ 和 $(7,3)$ 序列、$c=1/4$ 与 $1/3$ 重复位置误差计算。

\section{数值结果}
比较三种方法：与 Godunov 法对应的非交错常规 RCM（记作 RCMG）、式（7.18）的 Lax--Friedrichs 型随机选择法（RCMLF），以及式（7.30）、（7.32）的确定性 FORCE 格式。仍使用第 6 章五个一维 Euler 精确黎曼测试和相同左右状态。区域 $[0,1]$ 划分为 $M=100$；随机法取 $C_{\mathrm{CFL}}=0.45$，FORCE 取 0.9。精确解、RCMG 与 FORCE 分别由 NUMERICA 代码 \texttt{HE-E1RPEXACT}、\texttt{HE-E1RCM} 和 \texttt{HE-E1FOCENT} 计算。

表 7.4 中误差按 $\Delta x$ 归一化。$(2,1)$ 在 $c=1/4$ 时十步最大误差为一个网格宽度，$(5,3)$ 最大约 $3/4$；在 $c=1/2$ 时，$(5,3)$ 每两步抽样次数恰好匹配期望位移，偶数步位置精确。练习换用其他序列和 Courant 数，是为了说明序列优劣还可能与阈值 $c$ 的算术关系有关，不能只依据有限个全局统计量机械选择。

% [原PDF p277]
RCMG 在测试 1 上非常准确：激波与接触保持真正间断，稀疏波头尾的导数间断也解析良好；跨声速稀疏波中的熵瑕疵比 Godunov 法更小。在无跃迁而近真空的测试 2 中，光滑区出现随机噪声，暴露 RCM 的主要弱点。测试 3、4 含多道强间断，RCM 的无限间断分辨率优势最明显。

RCMLF 的结果稍逊于常规 RCM，但稳定区间加倍；曲线可看到与普通 Lax--Friedrichs 类似的相邻状态配对。FORCE 在五个测试中一般不如 Godunov，测试 5 差距尤其明显；不过 FORCE 无需黎曼求解器，结构与成本接近 Lax--Friedrichs，而结果显著优于后者。

\section{结论性评述}
常规 RCM 可直接处理只有两个独立变量的双曲系统。一维非定常 Euler 的两个变量是 $(x,t)$；二维定常超声速 Euler 则把主流方向 $x$ 看作时间样变量，另一个变量为横向空间 $y$。

三类数值结果应与第 6 章同一测试、同一网格比较。RCMG 和 RCMLF 的 CFL 取 0.45，是接近各自常规抽样设置中的安全值；FORCE 的确定性稳定上限为 1，取 0.9。RCM 的计算成本与精确 Godunov 法相近，因为每界面仍需求完整精确黎曼结构；FORCE 只计算状态平均和物理通量，明显更便宜。评价时因此不仅看曲线锐度，还要结合稳定步长与每步成本。

% [原PDF p278]
RCM 最大优点是把间断保持为零厚度，最大缺点是光滑区的随机性。混合策略可只在间断附近使用 RCM，在其余区域换用光滑高阶格式。对非齐次系统
\begin{equation}\bm U_t+\bm F(\bm U)_x=\bm S(\bm U),\tag{7.48}\end{equation}
算子分裂可处理柱/球对称代数源项，甚至含高阶空间导数的黏性项。源项不破坏间断的无限分辨率，却可能增强光滑区噪声；若特殊情况下能把源项纳入黎曼解，噪声会明显减轻。

把 RCM 的核心特征推广到三个以上独立变量尚未成功。普通维度分裂会破坏多维波几何；若加入人工黏性或采用混合方法，则可计算多维问题，但必须放弃激波的无限分辨率。Toro 的二维混合方法在接触与剪切波处使用 RCM，其他位置使用高分辨率激波捕捉，因此仍能精确保持极难解析的材料界面；Loh 等的三维定常超声速方法也放弃激波零厚度，却以 Lagrange 思路保留滑移面。

% [原PDF p279]
目前，RCM 与第 15 章算子分裂结合，仍是求解一维非线性非齐次系统的高精度选择。

\section*{三类方法的关系与使用要点}
从统一的局部演化观点看，Godunov、RCM 与 FORCE 都从分片常数数据出发。Godunov 在每个界面求黎曼解，并把一步后的解在固定单元上作守恒平均；RCM 求同样的黎曼解，却用一个离散样本代替平均；FORCE 的推导曾借用交错 RCM，但最终用控制体积分式消去黎曼解，只保留中心状态和物理通量。这三条路线分别突出守恒平均、随机无扩散抽样和低成本中心通量。

对单个移动间断，Godunov 每步把部分单元填成左右状态的平均，位置均值正确但界面逐渐展宽；RCM 每步要么把激波推进整格，要么完全不动，始终只有一个网格跳跃，但实际位置在精确轨迹附近随机摆动。若研究激波与材料界面的细节，零厚度非常有吸引力；若研究光滑声波或需要可重复的小扰动，随机噪声可能比平均扩散更麻烦。方法选择取决于希望控制哪类误差。

低差异序列并不是为了模拟物理随机性，而是为了用少量确定样本尽可能均匀地代表积分区间。普通随机数会出现长串聚集，使激波连续多步不动或连续跳格；van der Corput 序列有规律地填补空隙，使任何前缀的抽样比例更接近期望。一个时间层全域共用同一 $\theta_n$ 还能避免相邻单元各自随机选择造成不相容的高频噪声。为了复现实验，程序应固定序列定义，而不是依赖平台随机数发生器。

常规非交错 RCM 的实现步骤是：设置虚单元；为每个界面求精确黎曼结构；由所有波求 $S_{\max}$ 并选 $C_{\mathrm{CFL}}\le1/2$；取得本层序列值 $\theta_n$；对每个单元按 $\theta_n\le1/2$ 与否选择左或右界面解，并在式（7.3）的射线上抽样；把返回原始状态写入新数组。新旧数组必须分开，不能原地更新后又拿新值建立尚未处理单元的黎曼问题。

Lax--Friedrichs 型 RCM 则用 $RP(\bm U_{i-1}^n,\bm U_{i+1}^n)$，在半时间 $\Delta t/2$、位置 $\theta_n\Delta x$ 抽样，允许 $C_{\mathrm{CFL}}\le1$。它跳过中心状态，因而保留普通 Lax--Friedrichs 的奇偶配对倾向。尽管一步数可约减半，计算每个新单元仍需精确黎曼解，且曲线质量通常低于 RCMG；是否更高效要同时考虑目标误差而非只看稳定上限。

交错 RCM 每完整步需要两轮求解和抽样。它的主要意义之一是揭示确定性 FORCE 格式：把随机样本替换为积分平均，再用守恒律计算平均，最终得到 RI 与 LF 通量的均值。实现 FORCE 时先计算
\[
\bm U_{i+1/2}^{n+1/2}=\frac12(\bm U_i^n+\bm U_{i+1}^n)
-\frac{\Delta t}{2\Delta x}(\bm F_{i+1}^n-\bm F_i^n),
\]
由它得到 Richtmyer 通量；再与 Lax--Friedrichs 通量平均。整个过程不需要特征值、特征矢量或波型分支，适合状态方程复杂、黎曼解困难的系统。

FORCE 的单调性结论在本章针对线性平流三点格式。对一般非线性系统，实际稳健性仍受物理通量、状态正性和时间步影响，不能把线性证明无限外推。它的一阶人工黏性虽只有 LF 的一半，却仍会显著涂抹接触。第 13、14 章的高阶 FORCE/WAF 类方法将通过重构和限制器提高光滑区精度，同时在间断附近控制振荡。

含源项时，最直接做法是把齐次双曲步与源项常微分方程步分裂。RCM 的齐次步产生随机位置误差，源项若对状态非线性敏感，会把这种误差转换成幅值噪声；柱对称几何源项靠近轴线时尤需注意。若能构造含源项的广义黎曼问题，局部演化会更平衡，随机性增强较小。否则可采用混合检测器，只在强间断附近使用 RCM，光滑区换成确定性高阶方法。

多维推广的困难也源于局部波几何。沿坐标方向分裂的一维 RCM 会分别随机移动间断，斜激波和涡片的几何位置无法保持一致，抽样误差在方向切换中累积。加入人工黏性可以稳定计算，却放弃 RCM 最独特的“无限分辨率”。混合方法选择只对接触、剪切或材料界面采用 Lagrange/随机跟踪，而让激波由常规高分辨率守恒格式捕捉，是更现实的折中。

解读图 7.10--7.22 时应分别观察跃迁宽度、位置误差、光滑区噪声和整体守恒趋势。RCMG 的符号可能围绕精确稀疏曲线散布，但间断几乎垂直；RCMLF 除随机散布外还有偶奇配对；FORCE 曲线平滑、单调，却把接触扩展成宽带。没有一种方法在所有指标上占优，这组图的意义正是把不同误差机制并列展示。
\begin{sourceblock}
\sourcecrop{279}{45}{325}{395}{555}{0.31}
\captionof{figure}{RCMG 用于测试 1：密度、速度、压力和比内能；符号为数值解，实线为精确解，$x_0=0.3,t=0.2$。}
\end{sourceblock}
激波与接触几乎只有一个网格跳跃，稀疏波也保持清楚。图中 density、velocity、pressure、internal energy、position 分别为密度、速度、压力、比内能和位置。

% [原PDF p280]
\begin{sourceblock}
\sourcecrop{280}{45}{385}{395}{610}{0.28}\\[0.45em]
\sourcecrop{280}{45}{115}{395}{337}{0.28}
\captionof{figure}{RCMG 的测试 2（上）与测试 3（下）。测试 2 为近真空双稀疏波，测试 3 含极强右激波。}
\end{sourceblock}
测试 2 的光滑低密度区可见随机散点，这是随机抽样代替积分平均的直接后果；测试 3 的强激波和接触仍作为尖锐跳跃传播，显示该方法处理多间断相互作用的能力。

% [原PDF p281]
\begin{sourceblock}
\sourcecrop{281}{45}{385}{395}{610}{0.28}\\[0.45em]
\sourcecrop{281}{45}{115}{395}{337}{0.28}
\captionof{figure}{上：RCMG 用于测试 4；下：RCMLF 用于测试 1。}
\end{sourceblock}
常规 RCM 在测试 4 中保持三道强间断的零厚度；RCMLF 在测试 1 中也相当尖锐，但相邻单元出现配对，光滑区随机误差更明显。后者的优势是可采用两倍 CFL 上限。

% [原PDF p282]
\begin{sourceblock}
\sourcecrop{282}{45}{385}{395}{610}{0.28}\\[0.45em]
\sourcecrop{282}{45}{115}{395}{337}{0.28}
\captionof{figure}{RCMLF 的测试 2（上）与测试 3（下）。}
\end{sourceblock}
近真空测试中的抽样噪声和强激波测试中的状态配对均比 RCMG 更明显。曲线坐标依次为密度、速度、压力、比内能与位置；原图测试 3 第二小图的纵轴文字重复为 density，应按物理变量理解为速度。

% [原PDF p283]
\begin{sourceblock}
\sourcecrop{283}{45}{385}{395}{610}{0.28}\\[0.45em]
\sourcecrop{283}{45}{115}{395}{337}{0.28}
\captionof{figure}{上：RCMLF 用于测试 4；下：FORCE 格式用于测试 1。}
\end{sourceblock}
RCMLF 仍把多道间断保持得很窄，但曲线带有离散随机偏移；确定性 FORCE 没有随机噪声，也不像二阶无限制格式产生强振荡，不过一阶人工黏性使接触和激波展宽。

% [原PDF p284]
\begin{sourceblock}
\sourcecrop{284}{45}{385}{395}{610}{0.28}\\[0.45em]
\sourcecrop{284}{45}{115}{395}{337}{0.28}
\captionof{figure}{FORCE 格式的测试 2（上）与测试 3（下）。}
\end{sourceblock}
FORCE 在近真空问题中保持稳定、曲线平滑，但比 RCM 更耗散；在极强激波问题中也能给出解，激波后平台与接触分辨率低于精确 Godunov 法。其优势是通量只用相邻状态和物理通量，无需特征分解或黎曼求解器。

% [原PDF p285]
\begin{sourceblock}
\sourcecrop{285}{45}{385}{395}{610}{0.28}\\[0.45em]
\sourcecrop{285}{45}{115}{395}{337}{0.28}
\captionof{figure}{FORCE 格式的测试 4（上）与测试 5（下）。}
\end{sourceblock}
测试 4 的三道右行间断均被单调涂抹；测试 5 的近静止接触展宽尤为明显，说明中心耗散格式难以保持零速线性波。与 Lax--Friedrichs 相比，FORCE 的数值黏性减半，因此在相近复杂度下仍有更好的整体分辨率。

% [原PDF p286]
\chapter{通量矢量分裂方法}
\section{引言}
迎风数值方法的特征是：网格上的方程离散要根据信息传播方向选择模板，从而把控制方程所描述的关键物理传播性质嵌入离散格式。识别迎风方向主要有两条路线：第 6 章的 Godunov 路线，以及本章将讨论的通量矢量分裂（Flux Vector Splitting, FVS）路线；二者也分别称黎曼路线和 Boltzmann 路线，相应方法称通量差分裂与通量矢量分裂。

FVS 判定迎风方向所需工作少于 Godunov 型方法，格式较简单、效率稍高，尤其适合隐式算法和定常空气动力学计算，因此应用广泛。但它对波结构的描述较粗，间断分辨率通常更差，特别是静止接触波和剪切波。Navier--Stokes 应用中，van Leer、Thomas 与 Roe 报告其 FVS 格式明显不如配 Roe 近似黎曼求解器的 Godunov 法准确。与 FVS 密切相关但本书不展开的还有 KFVS 和其他动理学格式。

“通量差分裂”从左右状态差或黎曼波跳跃出发，把界面通量写成中心部分加上沿特征场传播的修正；“通量矢量分裂”则先在每个状态处把物理通量本身分为正向和负向两部分，再在界面组合左侧正通量与右侧负通量。在线性常系数系统中，两条路线都可化为 $\bm A^+\bm U_L+\bm A^-\bm U_R$，看起来相同；对非线性系统，矩阵和特征速度依赖状态，先取状态差还是先分裂每侧通量会产生不同数值方法。

Godunov 路线的几何对象是局部黎曼问题。它尝试解析左右状态相互作用后形成的激波、稀疏波、接触和剪切波，所以能按不同波型分配耗散。FVS 路线更接近分子向不同方向运动所形成的宏观通量：正速成分来自左状态，负速成分来自右状态，不必显式求中间星形区。这种局部、可微且结构规则的分裂容易写成隐式 Jacobian，也便于在定常迭代中构造近似因子分解。

简化的代价主要落在零速或缓慢线性波上。接触间断跨越时，压力与法向速度相等而密度不同；理想黎曼求解器可识别这是一道单独接触波，若其速度为零则通量不应把两侧密度混合。许多 FVS 格式却分别从完整状态计算正负通量，产生与声学速度相关的耗散，即使接触静止也会被涂抹。剪切波具有相似退化特征，因此边界层和黏性计算中的切向分辨率会受影响。

所谓 KFVS（kinetic flux-vector splitting，动理学通量矢量分裂）从速度分布函数的正、负分子速度半空间积分构造通量。它与宏观 FVS 共享按传播方向拆分的思想，但分裂依据来自 Boltzmann 或模型动理学描述。Pullin、Perthame、Mandal 与 Deshpande、Xu 与 Prendergast 等人的工作给出了不同动理学实现；本书只列为相关路线，不在本章展开。

比较两类方法时，不能只看每个界面的浮点操作数。较耗散格式可能需要更细网格才能达到同样误差，抵消单步便宜的优势；精确或近似黎曼求解虽复杂，却可能用较少单元解析接触与剪切。另一方面，在大型定常隐式空气动力学程序中，易线性化、易构造稳定隐式算子可能比单个间断的最优锐度更重要。方法优劣必须结合目标波型、稳态或非稳态、显式或隐式推进以及允许网格规模共同判断。

下一页起将利用 Euler 通量的齐次性与特征值符号构造具体 FVS。读者可把第 5 章线性分裂关系作为参照：正、负特征速度决定两部分通量的方向，而非线性推广的关键是怎样在声速点附近平滑地定义分裂、怎样保持一致性，以及怎样避免过量耗散或熵问题。

本章沿用的符号中，$\bm F^+$ 表示向右传播的信息通量，$\bm F^-$ 表示向左传播的信息通量；上标正负指传播方向，不表示通量各分量数值本身必为正或负。界面组装始终取左状态的正向部分和右状态的负向部分，这一点与标量迎风格式完全对应。
% ===== 译文分片 4（已合并）=====

% [原PDF p287]
FVS 的历史发展说明，这类方法的价值不只在于避免精确黎曼迭代，还在于把传播方向直接编码到物理通量本身。Sanders--Prendergast、Steger--Warming 与 van Leer 的方案在航空计算中很早就获得应用，随后研究集中于提高音速区光滑性、降低接触耗散以及扩展到不满足齐次性的系统。齐次性缺失并不意味着无法迎风，而是不能简单用 $\bm A^\pm\bm U$ 定义分裂通量；必须直接构造满足和式、符号与超声速一致性的 $\bm F^\pm$。

阅读本章时应区分三个对象：原方程 Jacobian 的特征值 $\lambda_i$，构造矩阵 $\bm A^\pm$ 的谱，以及真正分裂通量 Jacobian 的特征值 $\widehat\lambda_i^\pm$。数值传播方向由最后一类决定。若只检查 $\lambda_i^\pm$ 的符号而忽略 $\partial\bm F^\pm/\partial\bm U$，可能得到形式上分裂、实际却不满足单向传播要求的通量。

第3章第3.1.1节已经指出，Euler 方程具有通量齐次性
$\bm F(\bm U)=\bm A(\bm U)\bm U$；浅水方程等重要系统却不满足这一性质。即使如此，也可通过适当改写绕开齐次性限制，将通量矢量分裂（flux-vector splitting，FVS）推广到这些系统，相关构造见 V\'azquez--Cend\'on [568]。Sanders 与 Prendergast [424]、Steger 与 Warming [463]、van Leer [560,561] 的奠基工作之后，FVS 得到了广泛应用并不断改进。

本章给出 FVS 的基础介绍。第8.2、8.3节先在简单模型方程上说明其思想，第8.4节依次推导 Steger--Warming、van Leer 以及 Liou--Steffen 三种 Euler 方程分裂，第8.5节比较数值结果。由 FVS 构造高阶格式的方法见第13、14章；含源项和多维问题分别见第15、16章。

\section{通量矢量分裂方法}
\subsection{迎风差分}
先考虑小扰动定常超声速方程
\begin{equation}
 u_x-a^2v_y=0,\qquad v_x-u_y=0,\qquad
 a=\sqrt{M_\infty^2-1},\quad M_\infty>1, \tag{8.1--8.2}
\end{equation}
写成
\begin{equation}
 \bm U_x+\bm A\bm U_y=0,\qquad
 \bm U=(u,v)^T. \tag{8.3}
\end{equation}

% [原PDF p288]
这里的“迎风”必须按每个特征场分别理解，而不能仅按总体流速判断。矩阵同时含正、负特征值时，一个固定方向的单边差分必然至少对一族信息采用顺风取点；线性稳定性分析会显示该分量随迭代增长。把矩阵谱分成正、负两部分，实质上是先投影到特征变量，各自选择上游点，再变换回守恒变量。这样既尊重有限传播方向，又保留系统耦合。

系数矩阵及其特征结构为
\begin{equation}
 \bm A=\begin{bmatrix}0&-a^2\\-1&0\end{bmatrix},\qquad
 \lambda_1=-a,\quad\lambda_2=a,\qquad
 \bm K^{(1)}=\begin{bmatrix}a\\1\end{bmatrix},\quad
 \bm K^{(2)}=\begin{bmatrix}a\\-1\end{bmatrix}. \tag{8.4--8.6}
\end{equation}
两个特征速度符号相反，因此单独采用任一侧单边差分都不能同时迎风。例如
\begin{align}
 \bm U_i^{n+1}&=\bm U_i^n-\frac{\Delta x}{\Delta y}\bm A(\bm U_i^n-\bm U_{i-1}^n), \tag{8.7}\\
 \bm U_i^{n+1}&=\bm U_i^n-\frac{\Delta x}{\Delta y}\bm A(\bm U_{i+1}^n-\bm U_i^n). \tag{8.8}
\end{align}
前式对 $\lambda_2>0$ 迎风，却对 $\lambda_1<0$ 顺风而不稳定；后式情形正好相反。CIR 方法把
\begin{equation}
 \bm A=\bm A^++\bm A^-, \tag{8.9}
\end{equation}
其中 $\bm A^+$ 的特征值非负、$\bm A^-$ 的特征值非正，于是得到
\begin{equation}
 \bm U_i^{n+1}=\bm U_i^n-\frac{\Delta x}{\Delta y}\bm A^+(\bm U_i^n-\bm U_{i-1}^n)
 -\frac{\Delta x}{\Delta y}\bm A^-(\bm U_{i+1}^n-\bm U_i^n). \tag{8.10}
\end{equation}
Chakravarthy 等人的分裂系数矩阵格式把这一过程推广到非守恒形式的非线性系统。对常系数线性系统，定义 $\bm F=\bm A\bm U$ 即可把 CIR 格式写成守恒形式。

% [原PDF p289]
守恒形式的重要性不只是书写方便。对含激波的弱解，只有界面通量在相邻单元间以相反符号出现，才能保证离散总量严格收支；直接离散非守恒的矩阵乘积可能给出错误激波速度。在线性常系数情形，$\bm F=\bm A\bm U$ 使矩阵分裂和通量分裂完全等价，也解释了式(8.14)为何与精确 Godunov 通量一致。非线性推广的困难正是如何保留这种方向性与守恒性。

矩阵分裂自然诱导
\begin{equation}
 \bm F=\bm F^++\bm F^-,\qquad
 \bm U_i^{n+1}=\bm U_i^n-\frac{\Delta t}{\Delta x}
 (\bm F_{i+1/2}-\bm F_{i-1/2}), \tag{8.12--8.13}
\end{equation}
其中
\begin{equation}
 \bm F_{i+1/2}=\bm F_i^+(\bm U_i^n)+\bm F_{i+1}^-(\bm U_{i+1}^n). \tag{8.14}
\end{equation}
对线性问题，这恰好是 Godunov 单元界面通量；FVS 就是把这一思想推广到守恒形式的非线性系统。

\subsection{FVS 的一般构造}
考虑 $m$ 个非线性双曲守恒律
\begin{equation}
 \bm U_t+\bm F(\bm U)_x=0,\qquad
 \bm A(\bm U)=\frac{\partial\bm F}{\partial\bm U}
 =\bm K\bm\Lambda\bm K^{-1}. \tag{8.15--8.17}
\end{equation}
其中 $\bm\Lambda=\operatorname{diag}(\lambda_1,\ldots,\lambda_m)$，$\bm K$ 的第 $i$ 列是与 $\lambda_i$ 对应的右特征矢量。FVS 寻找
\begin{equation}\bm F(\bm U)=\bm F^+(\bm U)+\bm F^-(\bm U).\tag{8.20}\end{equation}

% [原PDF p290]
分裂后的 $\widehat{\bm A}^{\pm}$ 是分裂通量自身对状态的导数，不应同构造中使用的 $\bm A^\pm$ 混淆。二者只有在很特殊的线性情况下相同。符号条件施加在前者的特征值上，目的是保证 $\bm F^+$ 只携带向右信息、$\bm F^-$ 只携带向左信息；超声速一致性条件则防止在全部波同向时仍人为混入下游状态。齐次性使矩阵谱分裂能够直接转成通量谱分裂。

分裂通量的 Jacobian
\begin{equation}
 \widehat{\bm A}^{\pm}=\frac{\partial\bm F^{\pm}}{\partial\bm U}
\end{equation}
必须分别只有非负和非正特征值，即 $\widehat\lambda_i^+\ge0$、$\widehat\lambda_i^-\le0$。当原 Jacobian 的全部特征值同号时，分裂还必须退化为普通迎风：全为非负时 $\bm F^+=\bm F,\bm F^-=0$；全为非正时 $\bm F^+=0,\bm F^-=\bm F$。

若系统另满足齐次性
\begin{equation}\bm F(\bm U)=\bm A(\bm U)\bm U,\tag{8.24}\end{equation}
只需分裂特征值
\begin{equation}
 \lambda_i=\lambda_i^++\lambda_i^-,\qquad \lambda_i^+\ge0,\quad\lambda_i^-\le0, \tag{8.25}
\end{equation}
再令
\begin{equation}
 \bm\Lambda=\bm\Lambda^++\bm\Lambda^-,\qquad
 \bm A^\pm=\bm K\bm\Lambda^\pm\bm K^{-1},\qquad
 \bm A=\bm A^++\bm A^-. \tag{8.26--8.29}
\end{equation}
于是 $\bm F^\pm=\bm A^\pm\bm U$。

% [原PDF p291]
绝对值分裂可写成 $\bm A^\pm=(\bm A\pm|\bm A|)/2$，其中 $|\bm A|=\bm K|\bm\Lambda|\bm K^{-1}$。因此正、负通量的和恢复物理通量，差则对应特征耗散。每个 $\lambda_i$ 在零点的绝对值折角会传入数值通量，这正是 Steger--Warming 分裂在音速区光滑性不足的根源。练习8.1要求逐级验证标量、对角矩阵和完整矩阵三层恒等式。

Steger--Warming 分裂取
\begin{equation}
 \lambda_i^+=\frac12(\lambda_i+|\lambda_i|),\qquad
 \lambda_i^-=\frac12(\lambda_i-|\lambda_i|), \tag{8.32}
\end{equation}
显然 $\lambda_i^+\ge0$、$\lambda_i^-\le0$，并有
\begin{equation}
 \lambda_i=\lambda_i^++\lambda_i^-,\quad |\lambda_i|=\lambda_i^+-\lambda_i^-,\quad
 |\bm A|=\bm A^+-\bm A^-. \tag{8.35}
\end{equation}

\section{等温气体方程的 FVS}
等温气体动力学方程为
\begin{equation}
 \bm U_t+\bm F(\bm U)_x=0,\qquad
 \bm U=\begin{bmatrix}\rho\\\rho u\end{bmatrix},\quad
 \bm F=\begin{bmatrix}\rho u\\\rho u^2+\rho a^2\end{bmatrix}, \tag{8.36--8.37}
\end{equation}
其中声速 $a>0$ 为常数。其 Jacobian、特征值和右特征矩阵是
\begin{equation}
 \bm A=\begin{bmatrix}0&1\\a^2-u^2&2u\end{bmatrix},\quad
 \lambda_{1,2}=u\mp a,\quad
 \bm K=\begin{bmatrix}1&1\\u-a&u+a\end{bmatrix}. \tag{8.38--8.40}
\end{equation}
读者可直接验证该系统满足齐次性。

% [原PDF p292]
实际推导中最费事的步骤是求 $\bm K^{-1}$ 并完成三个矩阵相乘；系统规模增大后，手工展开非常容易出错。等温二方程例子把这一过程完整展示出来：先形成任一符号分支的 $\bm A^\alpha$，再乘守恒状态。最终式(8.44)只含两个特征速度分支，说明复杂矩阵代数可以压缩成按波族加权的通量。后续 Euler 公式沿用同一组织方式。

\subsection{分裂通量}
对任意特征值分裂，$\bm\Lambda^\alpha=\operatorname{diag}(\lambda_1^\alpha,\lambda_2^\alpha)$，其中 $\alpha$ 表示 $+$ 或 $-$。由于
\begin{equation}
 \bm K^{-1}=\frac1{2a}\begin{bmatrix}u+a&-1\\a-u&1\end{bmatrix}, \tag{8.42}
\end{equation}
可得
\begin{equation}
 \bm A^\alpha=\frac1{2a}
 \begin{bmatrix}
 \lambda_1^\alpha(u+a)-\lambda_2^\alpha(u-a)&\lambda_2^\alpha-\lambda_1^\alpha\\
 (u^2-a^2)(\lambda_1^\alpha-\lambda_2^\alpha)&
 \lambda_2^\alpha(u+a)-\lambda_1^\alpha(u-a)
 \end{bmatrix}. \tag{8.43}
\end{equation}
因而
\begin{equation}
 \bm F^\alpha=\bm A^\alpha\bm U=
 \frac{\rho}{2}\begin{bmatrix}
 \lambda_1^\alpha+\lambda_2^\alpha\\
 \lambda_1^\alpha(u-a)+\lambda_2^\alpha(u+a)
 \end{bmatrix}. \tag{8.44}
\end{equation}
令 $\alpha=+$ 或 $-$ 就分别得到正、负通量。练习8.5要求在亚声速情形求 $\widehat{\bm A}^{\pm}$ 及其特征值。

% [原PDF p293]
练习8.5揭示了这一差别：即便构造时令某个 $\lambda_i^+=0$，求导后的 $\widehat{\bm A}^+$ 也可能两个特征值都严格为正。因此“分裂通量退化”不是从原矩阵分裂自动继承的性质。若零特征场消失，静止或缓慢传播结构会被正负分裂同时携带并产生额外数值黏性。van Leer 方案明确把亚声速退化作为设计约束，正是为了让某些中间结构不受无关传播速度影响。

在有限体积实现中，式(8.48)可看成一个无需解局部非线性波扇的界面函数：先分别在 $\bm U_i$、$\bm U_{i+1}$ 上计算热力学量与分裂通量，再相加。它不需要左右状态共同的平均声速或星形压力，所以代码路径短、并行友好；但两侧独立分裂也意味着它无法自动利用接触波上压力和法向速度连续等跨界面关系。

对正通量，练习的结果为
\begin{equation}
 \widehat{\bm A}^+=\begin{bmatrix}
 a/2&1/2\\(a^2-u^2)/2&u+a
 \end{bmatrix}. \tag{8.45a}
\end{equation}
其特征多项式和特征值为
\begin{align}
 \lambda^2-(3a/2+u)\lambda+\tfrac14(u+a)^2&=0,\\
 \widehat\lambda_{1,2}^+&=\frac a4\left(2M+3\mp\sqrt{4M+5}\right),\qquad M=u/a. \tag{8.45b}
\end{align}
应注意 $\widehat{\bm A}^+\ne\bm A^+$、$\widehat\lambda_i^+\ne\lambda_i^+$；并且这里两个 $\widehat\lambda_i^+$ 都严格为正，没有零特征值。这一性质在数值上并不理想，而且会延续到 Euler 方程。后续分裂会消除该缺点。

\subsection{FVS 数值格式}
显式守恒更新为
\begin{equation}
 \bm U_i^{n+1}=\bm U_i^n-\frac{\Delta t}{\Delta x}
 (\bm F_{i+1/2}-\bm F_{i-1/2}),\qquad
 \bm F_{i+1/2}=\bm F_i^+(\bm U_i^n)+\bm F_{i+1}^-(\bm U_{i+1}^n). \tag{8.47--8.48}
\end{equation}
界面通量由左单元向右传播的正通量和右单元向左传播的负通量组成。对某单元，Steger--Warming 分裂须区分向左超声速、向右超声速与亚声速三种波型。

% [原PDF p294]
图8.1的物理含义是：每个单元内部先把通量分成向右和向左两束，界面只接收左单元的向右束与右单元的向左束。图8.2进一步按局部 Mach 数分类。在完全向左或向右的超声速单元中，一束为零而另一束等于完整物理通量；亚声速单元中，两支特征波分居零的两侧，因而两束都非零。实现时这种分支判断应使用同一时层的局部状态。

向左超声速时 $u_i^n+a_i^n\le0$，全部波向左，故
\begin{equation}\bm F_i^+=0,\qquad\bm F_i^-=\bm F_i^n.\tag{8.49}\end{equation}
向右超声速时 $u_i^n-a_i^n\ge0$，全部波向右，故
\begin{equation}\bm F_i^+=\bm F_i^n,\qquad\bm F_i^-=0.\tag{8.50}\end{equation}
亚声速时 $u_i^n-a_i^n\le0\le u_i^n+a_i^n$，只有 $\lambda_2$ 为正，因此
\begin{equation}
 \bm F_i^+=\frac{\rho_i^n}{2}\begin{bmatrix}u_i^n+a_i^n\\(u_i^n+a_i^n)^2\end{bmatrix},\qquad
 \bm F_i^-=\frac{\rho_i^n}{2}\begin{bmatrix}u_i^n-a_i^n\\(u_i^n-a_i^n)^2\end{bmatrix}. \tag{8.52}
\end{equation}
这三种情形说明，分裂在超声速极限自动恢复普通单边迎风，而在亚声速区由界面两侧共同贡献。

\section{Euler 方程的 FVS}
本节给出三种用于非定常 Euler 方程的通量矢量分裂方案。

% [原PDF p295]
这里再次列出 Euler 特征结构，是为了强调 FVS 并不求解左右状态之间的黎曼问题，而是在每个单元自身状态上分裂物理通量。三条一维特征速度分别对应左声学波、质点对流波和右声学波；$H$ 把压力功并入能量输运。尽管界面通量也使用左右数据，其组合方式是两侧独立分裂后相加，不像 Godunov 方法那样先构造统一波扇。

\subsection{方程回顾}
一维 Euler 方程为
\begin{equation}
 \bm U_t+\bm F(\bm U)_x=0,\quad
 \bm U=\begin{bmatrix}\rho\\\rho u\\E\end{bmatrix},\quad
 \bm F=\begin{bmatrix}\rho u\\\rho u^2+p\\u(E+p)\end{bmatrix}. \tag{8.53--8.54}
\end{equation}
其特征值和右特征矩阵为
\begin{equation}
 \lambda_1=u-a,\quad\lambda_2=u,\quad\lambda_3=u+a,
\end{equation}
\begin{equation}
 \bm K=\begin{bmatrix}1&1&1\\u-a&u&u+a\\H-ua&u^2/2&H+ua\end{bmatrix},\qquad
 H=\frac{E+p}{\rho}=\frac12u^2+\frac{a^2}{\gamma-1}. \tag{8.56--8.58}
\end{equation}
三维 Euler 方程可先把通量旋转到单元面法向；笛卡尔网格上只需逐坐标方向处理。一般曲线网格则使用第3章的旋转矩阵。

% [原PDF p296]
三维情况下，法向速度 $u$ 决定全部五族波的传播方向，$v,w$ 作为切向量随中间三重特征速度输运。对任意朝向的单元面，可通过正交旋转把法向对齐局部 $x$ 轴，调用同一一维通量后再旋回全局坐标。这个旋转不变性使下述分裂可用于一般有限体积几何，而无需为每一种面方向重新推导特征系统。

沿 $x$ 方向分裂的三维系统写作
\begin{equation}
 \bm U=\begin{bmatrix}\rho\\\rho u\\\rho v\\\rho w\\E\end{bmatrix},\qquad
 \bm F=\begin{bmatrix}\rho u\\\rho u^2+p\\\rho uv\\\rho uw\\u(E+p)\end{bmatrix}. \tag{8.59--8.60}
\end{equation}
记 $V^2=u^2+v^2+w^2$、$\widehat\gamma=\gamma-1$，则
\begin{equation}H=\frac12V^2+\frac{a^2}{\gamma-1}.\tag{8.62}\end{equation}
五个特征值为
\begin{equation}\lambda_1=u-a,\quad \lambda_2=\lambda_3=\lambda_4=u,\quad\lambda_5=u+a.\tag{8.63}\end{equation}
对应的右特征矢量分别描述两支声学波、熵波以及两支切向剪切波。三维分裂系统也满足 $\bm F=\bm A\bm U$；因此一维推导可逐面推广，切向动量仅随法向速度对流。

% [原PDF p297]
三维逆特征矩阵虽比一维长，最终通量式仍只出现左右声学特征值和一个代表三重中间场的特征值。这是因为同速的熵波、两支剪切波可以合并权重。质量和法向动量由声学与中间场共同贡献；切向动量的每一项只是对应速度乘同一质量型权重；能量项含 $H\mp ua$ 与 $V^2$。按这一结构编程比逐项存储完整 $5\times5$ 矩阵更省运算，也更容易核对旋转对称性。

Steger--Warming 公式对任意符号分支 $\alpha$ 都相同，只有 $\lambda_k^\alpha$ 不同。因此可先用 $\lambda_k^\pm=(\lambda_k\pm|\lambda_k|)/2$ 计算两个谱数组，再调用同一通量装配函数。音速处 $|\lambda|$ 的不可微性会使通量导数出现折角，可能在跨音速膨胀中造成熵相关误差；这与 Roe 线性化的扩张激波机制不同，但都提示零特征速度需要特别关注。

\subsection{Steger--Warming 分裂}
一维右特征矩阵的逆为
\begin{equation}
 \bm K^{-1}=\frac{\gamma-1}{2a^2}
 \begin{bmatrix}
 \tfrac12u^2+ua/(\gamma-1)&-u-a/(\gamma-1)&1\\
 2a^2/(\gamma-1)-u^2&2u&-2\\
 \tfrac12u^2-ua/(\gamma-1)&a/(\gamma-1)-u&1
 \end{bmatrix}. \tag{8.66}
\end{equation}
由 $\bm F^\alpha=\bm K\bm\Lambda^\alpha\bm K^{-1}\bm U$ 得
\begin{equation}
 \bm F^\alpha=\frac{\rho}{2\gamma}
 \begin{bmatrix}
 \lambda_1^\alpha+2(\gamma-1)\lambda_2^\alpha+\lambda_3^\alpha\\
 (u-a)\lambda_1^\alpha+2(\gamma-1)u\lambda_2^\alpha+(u+a)\lambda_3^\alpha\\
 (H-ua)\lambda_1^\alpha+(\gamma-1)u^2\lambda_2^\alpha+(H+ua)\lambda_3^\alpha
 \end{bmatrix}. \tag{8.67}
\end{equation}
三维公式具有相同结构：中间三重特征值的贡献分别携带法向动量、两项切向动量和总能量；将 $u^2$ 换成 $V^2$ 即得能量分量。其优势是计算直接，缺点是分裂通量在声速点附近不光滑，且难以保持静止接触的锐利分辨。

% [原PDF p298]
van Leer 的两个附加要求直接针对 Steger--Warming 的数值缺点。分裂 Jacobian 连续可避免 Mach 数穿过零或一时通量导数突跳；亚声速退化则让某个中间方向没有人为传播速度，从而减少静止或慢速结构的耗散。构造从质量通量的二次 Mach 多项式开始，再用动量、能量一致性确定其余分量，确保两支通量相加严格恢复原 Euler 通量。

\subsection{van Leer 分裂}
van Leer 要求：（I）$\widehat{\bm A}^{\pm}=\partial\bm F^\pm/\partial\bm U$ 连续；（II）亚声速时分裂退化，即每个分裂 Jacobian 至少有一个零特征值。以 $M=u/a$ 表示原通量：
\begin{equation}
 \bm F(\rho,a,M)=\begin{bmatrix}
 \rho aM\\ \rho a^2(M^2+1/\gamma)\\
 \rho a^3M(M^2/2+1/(\gamma-1))
 \end{bmatrix}. \tag{8.70}
\end{equation}
亚声速区的质量分裂为
\begin{equation}f_{\rm mas}^+=\tfrac14\rho a(1+M)^2,\qquad
f_{\rm mas}^-=-\tfrac14\rho a(1-M)^2.\tag{8.71}\end{equation}
动量与能量通量则取
\begin{align}
 f_{\rm mom}^{\pm}&=f_{\rm mas}^{\pm}\frac{2a}{\gamma}
 \left(\frac{\gamma-1}{2}M\pm1\right), \tag{8.72}\\
 f_{\rm ene}^{\pm}&=\frac{\gamma^2}{2(\gamma^2-1)}
 \frac{(f_{\rm mom}^{\pm})^2}{f_{\rm mas}^{\pm}}. \tag{8.73}
\end{align}
它们在 $M=\pm1$ 与完全迎风通量连续衔接，避免了 Steger--Warming 分裂在音速转换处的某些缺陷。

% [原PDF p299]
式(8.76)--(8.77)中切向速度直接乘分裂质量通量，表示切向动量没有独立声学压力项。稳定上限依赖 Mach 数：静止亚声速区最严格，接近声速时逐渐放宽到1。这种依赖意味着全局统一 CFL 数必须按整个网格最不利状态选取；若局部时间步，则各单元可使用自身上限。相比之下，Godunov 通量的允许 CFL 通常更接近由最快物理波决定的自然界限。

对沿 $x$ 分裂的三维方程，正负通量为
\begin{equation}
 \bm F^\pm=\pm\frac14\rho a(1\pm M)^2
 \begin{bmatrix}
 1\\ \dfrac{2a}{\gamma}\left(\dfrac{\gamma-1}{2}M\pm1\right)\\v\\w\\
 \dfrac{2a^2}{\gamma^2-1}\left(\dfrac{\gamma-1}{2}M\pm1\right)^2+\dfrac12(v^2+w^2)
 \end{bmatrix}. \tag{8.76--8.77}
\end{equation}
van Leer 给出的实用稳定条件为
\begin{equation}
 C_{\rm CFL}=\frac{\Delta t}{\Delta x}(|u|+a)
 \le \frac{2\gamma+|M|(3-\gamma)}{\gamma+3}. \tag{8.78}
\end{equation}
当 $\gamma=1.4$ 时，$|M|=1$ 给出上限1，而 $M=0$ 时最小值约为0.636。它比 Godunov 方法接近1的 CFL 限制更严格。

\subsection{Liou--Steffen（AUSM）格式}
Liou 与 Steffen 把通量分成对流部分和压力部分：
\begin{equation}
 \bm F=\underbrace{\begin{bmatrix}\rho u\\\rho u^2\\\rho uv\\\rho uw\\\rho uH\end{bmatrix}}_{\bm F^{(c)}}+
 \underbrace{\begin{bmatrix}0\\p\\0\\0\\0\end{bmatrix}}_{\bm F^{(p)}}. \tag{8.80}
\end{equation}

% [原PDF p300]
AUSM 的关键不是再对整个通量作特征分解，而是把质量随流输运和压力声学作用分开。界面 Mach 数控制密度、三项动量的对流部分及焓通量从哪一侧取值；界面压力只进入法向动量。亚声速多项式保证 $M^++M^-=M$、$p^++p^-=p$，并在 $|M|=1$ 与完全迎风分支连续连接。两种压力多项式在低 Mach 区给出不同耗散和光滑性。

由于 $\bm F^{(c)}=M\widehat{\bm F}^{(c)}$，界面通量取
\begin{equation}
 \bm F_{i+1/2}=\bm F_{i+1/2}^{(c)}+\bm F_{i+1/2}^{(p)},\qquad
 \bm F_{i+1/2}^{(c)}=M_{i+1/2}[\widehat{\bm F}^{(c)}]_{i+1/2}, \tag{8.82--8.83}
\end{equation}
方括号中的状态按 $M_{i+1/2}$ 的符号迎风选取。这就是“对流迎风分裂方法”（AUSM）名称的来源。界面 Mach 数与压力分别分裂为
\begin{equation}M_{i+1/2}=M_i^++M_{i+1}^-,\qquad p_{i+1/2}=p_i^++p_{i+1}^-.\tag{8.85--8.86}\end{equation}
Mach 数分裂为
\begin{equation}
 M^\pm=\begin{cases}\displaystyle \pm\frac14(M\pm1)^2,&|M|\le1,\\
 \displaystyle\frac12(M\pm|M|),&|M|>1,
 \end{cases} \tag{8.87}
\end{equation}
压力可用线性多项式 $p^\pm=p(1\pm M)/2$，也可用更平滑的三次形式
\begin{equation}
 p^\pm=\frac14p(M\pm1)^2(2\mp M),\qquad |M|\le1, \tag{8.89}
\end{equation}
超声速区则恢复按速度方向的完全迎风压力。AUSM 把对流传播和声学压力作用分开，兼具 FVS 与 Godunov 型方法的一些特点。

% [原PDF p301]
五个测试覆盖的不是同一类难点。测试1检验跨音速熵行为；测试2把星形区推向真空，检验正密度和正压力保持；测试3包含极强单向激波；测试4由相向强激波产生多重强间断；测试5通过 Galilean 平移使接触近乎静止，专门暴露线性波耗散。统一使用相同网格和边界条件，才能把差别主要归因于数值通量而非分辨率或边界处理。

\section{数值结果}
下面比较 $\gamma=1.4$ 的一维 Euler 方程上三种格式：Steger--Warming、van Leer 和 AUSM。数值解由 NUMERICA 的 HE--E1FVS 程序获得，精确解由 HE--E1RPEXACT 生成。五个黎曼测试均在 $0\le x\le1$ 上用100个单元计算，透射边界条件；一般取 $C_{\rm CFL}=0.9$，van Leer 格式取0.6，且最初5步进一步乘0.2。

初始数据如下：
\begin{center}\small
\begin{tabular}{c|rrr|rrr}
\toprule 测试&$\rho_L$&$u_L$&$p_L$&$\rho_R$&$u_R$&$p_R$\\\midrule
1&1.0&0.75&1.0&0.125&0&0.1\\
2&1.0&-2.0&0.4&1.0&2.0&0.4\\
3&1.0&0&1000&1.0&0&0.01\\
4&5.99924&19.5975&460.894&5.99242&-6.19633&46.0950\\
5&1.0&-19.59745&1000&1.0&-19.59745&0.01\\\bottomrule
\end{tabular}
\end{center}
测试5相当于给测试3叠加均匀负背景速度。所有结果都应与第6章 Godunov 格式的图6.8--6.12比较。

% [原PDF p302]
在测试1中，音速点是左稀疏扇内 $u-a=0$ 的位置；若分裂函数在零特征值附近不合适，就会形成折点、平台或额外扩散。测试2的中心区同时具有极低密度和压力，计算比内能时又要从总能量扣除动能，因而很容易放大误差。三种 FVS 能通过该测试，说明按单侧物理状态分裂的形式对某些近真空数据比线性化平均状态更耐受。

\subsection{测试1与测试2}
测试1是修正的 Sod 问题，包含向右激波、向右接触间断和向左含音速点的稀疏波，特别适合检查熵条件。Steger--Warming 对激波和接触的分辨与 Godunov 接近，但稀疏波头尾扩散，音速点处理明显不正确。van Leer 对稀疏波和接触几乎与 Godunov 一致，激波稍宽；AUSM 的激波更尖锐、接触相近，但音速点附近的稀疏波较差。

测试2是“123问题”：两支对称稀疏波围出近真空星形区，零速接触波本身是平凡的。三种 FVS 都能运行并给出总体良好的低密度解。Steger--Warming 与 Godunov 几乎一致；van Leer 的稀疏波头略扩散；AUSM 总体略优于 van Leer，但接触附近密度、压力趋零时，速度和内能曲线略显杂乱。许多线性化黎曼求解器会在该问题上失效，因此 FVS 的这一稳健性值得注意。

% [原PDF p303]
测试3的强度使微小通量失配在激波后平台上清晰可见，因此“能运行”并不等同于准确。测试4的慢激波则检验格式在波速接近零时是否产生异常耗散或振荡；高速右激波上的锐利表现不能补偿慢波附近的误差。测试5进一步把线性接触速度调到零，使任何与最大声速成比例、却不随接触速度消失的耗散都直接表现为持续抹宽。

\subsection{测试3、测试4与测试5}
测试3用于检验强健性和精度，解包含 Mach 数约198的向右强激波、接触面与向左稀疏波。Steger--Warming 比 Godunov 误差大，激波后出现非物理凹陷；van Leer 有所改善但仍不及 Godunov。作者实现的 AUSM 即使把 CFL 数降至0.1也无法完成这一严苛算例。

测试4由两支相向强激波产生，最终形成三条向右传播的强间断，其中左激波移动很慢。Steger--Warming 总体接近 Godunov，慢激波更扩散，却削弱了 Godunov 中的低频振荡。van Leer 更准确，慢激波跨两个内部单元，仍可见微弱低频伪振荡。AUSM 对快速右激波最锐利但产生过冲，慢激波又比 van Leer 稍宽。

测试5给测试3叠加负匀速，使接触间断静止；其解由左稀疏波、缓慢向右激波和静止接触组成，主要考验对低速接触的分辨能力。

% [原PDF p304]
本章比较还表明，“完整/不完整波模型”和“近真空正性”是两个相对独立的维度。精确或完整 Godunov 求解器善于保持接触，但某些线性化版本在近真空失败；FVS 不显式构造统一星形区，可能更耐受低密度，却会扩散接触。AUSM 通过分离对流与压力改善接触，同时在极强测试中暴露振荡或失效。选择通量时不能以单一 Sod 问题代替这种多维度评估。

从应用角度看，多维剪切层、涡和材料界面往往长时间以低法向速度沿网格存在，接触耗散会累计成显著误差；强激波附近则优先要求稳健和正性。现代方法常在不同区域混合或修正基础通量。无论采用何种增强，高阶重构都不能自动弥补基础一阶通量缺失的波族：高阶只在光滑区降低截断误差，间断附近仍受底层黎曼模型控制。

在选定输出时刻，右激波经过81步仅移动5个单元。Steger--Warming 与 van Leer 都把接触严重抹宽；AUSM 的接触明显更锐利，但在激波、接触甚至稀疏波尾附近出现非物理振荡。使用精确黎曼求解器的 Godunov 结果明显优于三者。

总体而言，FVS 通常能达到与 Godunov 相近的质量，差别集中在慢速或静止接触波。多维计算中，这直接影响剪切层、材料界面和涡流的精度。AUSM 对接触更敏锐，却在测试3崩溃、测试5产生强振荡，稳健性仍有疑问。另一方面，采用线性化黎曼求解器的 Godunov 方法可能在测试2的低密度区失效，而三种 FVS 均给出可接受解；若使用忽略线性波的不完备黎曼求解器，Godunov 对接触的分辨也会退化。可见黎曼求解器的选择决定了 Godunov 方法的表现。

高阶 FVS 见第13、14章；源项处理见第15章；多维扩展见第16章。

% [原PDF p305]
两组曲线采用相同坐标范围，因而可直接比较间断过渡占用的网格数。密度图最清楚显示接触与激波，速度图显示稀疏扇的连续展开，压力在接触处理论上连续，比内能则会跨接触跳跃。Steger--Warming 的稀疏扇头、扇尾更圆钝；van Leer 的音速区跟随精确线更自然。逐物理量对照比只看密度更能定位误差属于哪一特征场。

图8.3和图8.4分别给出测试1的 Steger--Warming 与 van Leer 结果，初始间断位于 $x_0=0.3$，输出时刻 $t=0.2$。四个子图依次比较密度、速度、压力和比内能；离散符号表示数值解，实线表示精确解。前者在音速稀疏波处较扩散，后者显著改善稀疏波轮廓，但激波厚度略有增加。

% [原PDF p306]
AUSM 曲线说明较锐利的激波并不必然伴随同样好的稀疏波。测试2中心区的压力、密度接近零，纵轴很小的绝对变化会对应显著相对误差；比内能由压力与密度之比形成，对二者误差尤其敏感。因此近真空算例应同时检查密度、压力、速度与内能，不能以单一曲线是否光滑判断稳健性。

图8.5为测试1的 AUSM 结果，条件同上；其激波更陡，音速点附近的扇形轮廓稍差。图8.6为测试2的 Steger--Warming 结果，$x_0=0.5$、$t=0.15$。对称稀疏波使中心区域接近真空，密度和压力降到很小，速度在两扇形之间迅速变化。

% [原PDF p307]
两幅图的精确线关于 $x=0.5$ 对称，任何明显左右不对称都可能来自实现或舍入。van Leer 的主要差异集中在稀疏扇外缘；AUSM 在中心的速度和比内能符号点有散乱，但密度、压力总体仍保持非负。近真空测试还应报告最小密度、最小压力和异常处理，绘图范围内看似平滑并不能完全证明正性。

图8.7和图8.8分别比较测试2的 van Leer 与 AUSM 结果，均为 $x_0=0.5$、$t=0.15$。van Leer 在两稀疏波头部略有扩散；AUSM 的整体轮廓稍锐利，但中心近真空区的速度、比内能出现轻微不规则。这些曲线同时说明，近真空条件下仅比较密度并不足以判断格式质量。

% [原PDF p308]
测试3的压力从约 $10^{-2}$ 到 $10^3$，比内能也跨越多个数量级。激波后的凹陷在速度和压力图最明显，说明它不是单纯密度过渡变宽，而是星形平台状态失真。van Leer 减小但未消除偏差。输出时间很短，是为了让极强激波仍位于计算域内并排除边界反射的干扰。

图8.9、8.10给出测试3在 $x_0=0.5$、$t=0.012$ 的结果。Steger--Warming 与 van Leer 都捕获了强激波位置和星形区平台，但前者在激波后速度与压力中出现更明显的非物理下凹；van Leer 抑制了该误差，仍比精确 Godunov 解更扩散。AUSM 在此算例中没有可用结果。

% [原PDF p309]
测试4中波相互作用后主要间断都向右，但速度相差悬殊。慢激波附近的低频振荡表现为平台上跨多个单元的缓慢波纹，不同于局部 Gibbs 尖峰。Steger--Warming 的额外耗散压低这类波纹，却同时增宽慢激波；van Leer 在锐利度与振荡幅度之间取得不同平衡。

图8.11、8.12给出测试4在 $x_0=0.4$、$t=0.035$ 的 Steger--Warming 和 van Leer 结果。两者均捕获向右传播的三个强间断。van Leer 对缓慢左激波的过渡约占两个网格单元，分辨优于 Steger--Warming；低频振荡虽弱，但在密度和压力平台上仍可辨认。

% [原PDF p310]
AUSM 在右激波上的过冲表示数值值越过精确平台，并非更高分辨的无代价收益。测试5给全部状态叠加相同负速度，使接触相对网格静止；连续 Euler 方程具有 Galilean 不变性，但离散耗散通常依赖相对网格速度。因此测试3与测试5的对比也检验格式在参考系变换下的表现。

图8.13为同一测试4的 AUSM 结果：高速右激波最锐利，却伴随过冲。图8.14为测试5的 Steger--Warming 结果，$x_0=0.8$、$t=0.012$；接触间断几乎不移动，却在短时间内被数值黏性显著抹宽，显示传统 FVS 对静止线性波的不利特性。

% [原PDF p311]
静止接触上精确质量通量为零，压力和速度连续；任何密度及内能混合都来自数值耗散。AUSM 的对流速度分裂可使该耗散显著变小，但其他位置振荡说明压力--对流耦合在强背景流下不够稳健。多维材料界面应用还需同时检查压力平衡、组分守恒与界面附近的非物理速度。

图8.15和图8.16分别为测试5的 van Leer 与 AUSM 结果。van Leer 与 Steger--Warming 相似，静止接触被严重扩散；AUSM 使接触跃变陡得多，但其附近以及缓慢激波和稀疏波尾部出现明显非物理振荡。锐利度与稳健性之间的取舍由此十分直观。

% [原PDF p312]
精确黎曼 Godunov 通量利用接触两侧共同的 $p_*,u_*$，所以静止接触的耗散只来自一阶更新和复合波影响，而不是缺失特征场。相同网格与时间步下，求解器内部波模型会显著改变结果。若换用忽略接触的 HLL 或传统 FVS，外侧激波位置可能仍正确，却不能据此推断界面分辨同样可靠。

图8.17给出测试5使用精确黎曼求解器的 Godunov 结果，仍为 $x_0=0.8$、$t=0.012$。它既能保持静止接触的锐利程度，又没有 AUSM 那样显著的伪振荡，因而在该测试上明显最佳。不过，这一结论依赖所用黎曼求解器；若换成线性化或不完备求解器，其近真空稳健性或接触分辨会相应下降。

% [原PDF p313]
\chapter{近似状态黎曼求解器}
“近似状态”路线的目标不是复现整个黎曼扇，而是用最少代价得到 Godunov 采样所需的信息。对 Euler 方程，这些信息集中在 $p_*$、$u_*$、两侧星形密度及外波类型。若 $S=0$ 位于常状态区，以上量足够；若位于音速稀疏扇，再调用无迭代扇内精确公式。于是昂贵的非线性压力迭代可被少量闭式运算替代，同时保留与精确采样相同的波型逻辑。

简单求解器“不够稳健”通常表现为星形压力或密度变成非物理值，或者强波速度严重失准。混合方案没有把这些结果钳制后继续使用，而是根据压力比和估计压力位置识别近似假设失效，再切换到与双稀疏或双激波结构一致的方法。这样的自适应是算法模型选择，而非掩盖失败的数值兜底。

\section{引言}
Godunov 方法及其高阶扩展需要反复求解黎曼问题；实际计算中这一操作可能执行数十亿次，因而通常是最昂贵的单项工作。第4章给出了理想气体和余容气体的精确求解器，但它总要迭代。若状态方程代数结构复杂，或待解方程组本身庞大，迭代代价会急剧增加。

数值方法可用两条路线提取近似黎曼信息：一是直接近似界面数值通量（第10--12章）；二是先近似界面所需状态，再把该状态代入物理通量。本章采用后一条路线，构造无迭代的星形区近似。最简单的公式虽然便宜，却未必足够稳健；自适应混合求解器会在光滑区使用廉价近似，在强间断附近切换到更可靠方法。

这些显式近似还有多种用途：可用于 MUSCL 特征限制；为 Roe 求解器的 Harten--Hyman 熵修正提供状态；为 Osher 路径积分给出交点；为 HLL 方法估计最小、最大信号速度；也可服务于前沿追踪。相同思路容易推广到浅水方程、定常超声速 Euler 方程、人工可压缩方程以及一般状态方程的 Euler 系统。

% [原PDF p314]
近似星形状态并非只能用于一阶通量。MUSCL 方法在特征空间限制斜率时需要波方向与波强，熵修正需要估计跨音速扇两侧状态，HLL 需要最外信号速度，Osher 路径积分需要中间交点。这些信息都可由本章闭式公式廉价提供。因此，即便最终采用别的黎曼通量，PVRS、TRRS、TSRS 仍是实用的辅助模块。其推广通常只要求掌握目标系统的原始变量特征关系或精确波曲线的近似。

本章先回顾 Godunov 通量和黎曼解，再给出原始变量线性化求解器（PVRS）、基于精确压力函数的双稀疏波和双激波近似，以及两种自适应组合，最后比较数值结果。

\section{黎曼问题与 Godunov 通量}
一般初边值问题写作
\begin{equation}
 \bm U_t+\bm F(\bm U)_x=0,\quad \bm U(x,0)=\bm U^{(0)}(x),\quad
 \bm U(0,t)=\bm U_l(t),\quad\bm U(L,t)=\bm U_r(t), \tag{9.1}
\end{equation}
显式守恒更新和 Godunov 通量为
\begin{equation}
 \bm U_i^{n+1}=\bm U_i^n+\frac{\Delta t}{\Delta x}
 (\bm F_{i-1/2}-\bm F_{i+1/2}),\qquad
 \bm F_{i+1/2}=\bm F(\bm U_{i+1/2}(0)). \tag{9.2--9.3}
\end{equation}
这里 $\bm U_{i+1/2}(0)$ 是局部黎曼问题自相似解在 $x/t=0$ 的取值。

% [原PDF p315]
图9.1把左右数据、两侧星形状态以及三类不同特征速度画在同一 $x$--$t$ 平面。法向速度和压力跨中间波连续，密度与切向速度可以跳跃；因此星形区不是单一状态。求通量只需知道竖直 $t$ 轴穿过哪一块常状态或哪一个稀疏扇，却必须先判断两支外波的类型和速度。十种波型来自星形速度正负、外波为激波或稀疏波，以及稀疏扇是否跨越 $S=0$ 的组合。

对沿 $x$ 分裂的三维 Euler 方程，
\begin{equation}
 \bm U=(\rho,\rho u,\rho v,\rho w,E)^T,\qquad
 \bm F=(\rho u,\rho u^2+p,\rho uv,\rho uw,u(E+p))^T. \tag{9.5}
\end{equation}
左右原始状态为 $\bm W_K=(\rho_K,u_K,v_K,w_K,p_K)^T$，$K=L,R$。精确解由两支非线性声学波和一支中间接触波组成，包含十种可能的采样波型。本章把任务分为三部分：
\begin{enumerate}
\item 求星形区 $p_*,u_*,\rho_{*L},\rho_{*R}$；
\item 求整个波系中的切向速度 $v,w$；
\item 当 $t$ 轴穿过音速稀疏扇时，求扇内的 $\rho,u,p$。
\end{enumerate}
后两项有直接精确公式，本章其余部分集中近似第一项。

% [原PDF p316]
式(9.8)体现了接触/剪切特征场的严格结构：$v,w$ 不受声学波改变，只随速度 $u_*$ 的中间波选择左右值。在多组分流中，每个无反应质量分数也按同一规则传播。如果近似求解器把整个星形区压成一个平均状态，这些量就会被过度混合。音速稀疏波则相反，不能只选扇头或扇尾常状态，因为 $S=0$ 真正落在连续自相似解内部。

\subsection{切向速度分量}
切向速度跨两支非线性波保持不变，只跨中间波发生跳跃。因此，一旦有 $u_*$，便有
\begin{equation}
 (v,w)(x,t)=\begin{cases}(v_L,w_L),&x/t\le u_*,\\(v_R,w_R),&x/t>u_*.
 \end{cases} \tag{9.8}
\end{equation}
除了 $u_*$ 本身是近似量外，这一切向解是精确的。任何随流体对流的被动标量都具有同样规律，例如多组分流中的组分浓度。因此近似求解器应保持中间线性波，而不能把所有中间结构都并入一个平均状态。

\subsection{音速稀疏波}
当 $t$ 轴穿过左或右稀疏扇时，建议直接使用第4章的无迭代精确扇内解。左扇内令 $S=x/t$，
\begin{align}
 \rho&=\rho_L\left[\frac2{\gamma+1}+\frac{\gamma-1}{(\gamma+1)a_L}(u_L-S)\right]^{2/(\gamma-1)},\\
 u&=\frac2{\gamma+1}\left[a_L+\frac{\gamma-1}{2}u_L+S\right],\\
 p&=p_L\left[\frac2{\gamma+1}+\frac{\gamma-1}{(\gamma+1)a_L}(u_L-S)\right]^{2\gamma/(\gamma-1)}. \tag{9.9}
\end{align}

% [原PDF p317]
图9.2列出的采样模式可以理解成一棵判定树：先看 $u_*$ 的符号决定从左或右星形侧采样，再看对应外波是激波还是稀疏波，最后比较波头、波尾速度与零。只有零落在扇内的两个末端分支需要式(9.9)、(9.10)；其余情形使用数据态或星形常态。把这套判定与星形近似分离，有助于保证不同求解器共享同一可靠采样程序。

右稀疏扇内相应为
\begin{align}
 \rho&=\rho_R\left[\frac2{\gamma+1}-\frac{\gamma-1}{(\gamma+1)a_R}(u_R-S)\right]^{2/(\gamma-1)},\\
 u&=\frac2{\gamma+1}\left[-a_R+\frac{\gamma-1}{2}u_R+S\right],\\
 p&=p_R\left[\frac2{\gamma+1}-\frac{\gamma-1}{(\gamma+1)a_R}(u_R-S)\right]^{2\gamma/(\gamma-1)}. \tag{9.10}
\end{align}
求 Godunov 通量时取 $S=0$。其余八种波型只需星形值和切向速度即可采样。

\section{原始变量黎曼求解器（PVRS）}
用原始变量 $\bm W=(\rho,u,v,w,p)^T$，Euler 方程的拟线性形式是
\begin{equation}\bm W_t+\bm A(\bm W)\bm W_x=0.\tag{9.11}\end{equation}

% [原PDF p318]
PVRS 的假设是整个局部波扇只是某个常状态附近的小扰动，因此把随状态变化的系数冻结。冻结系统仍严格双曲，其五个特征方向与原 Euler 方程相同，只是速度和声速取参考值。它不是在通量上作任意 Taylor 展开，而是精确求解一个定义明确的线性黎曼问题。误差来自左右态距离参考态过大，尤其是强激波、强膨胀和近真空，而不是线性系统求解本身。

其中
\begin{equation}
 \bm A=\begin{bmatrix}
 u&\rho&0&0&0\\0&u&0&0&1/\rho\\0&0&u&0&0\\0&0&0&u&0\\0&\rho a^2&0&0&u
 \end{bmatrix}. \tag{9.12}
\end{equation}
特征值为 $u-a,u,u,u,u+a$。困难在于 $\bm A$ 依赖未知状态。若左右数据和解都靠近某常状态 $\overline{\bm W}$，可冻结
\begin{equation}\overline{\bm A}=\bm A(\overline{\bm W})\tag{9.15}\end{equation}
并精确求解常系数线性黎曼问题
\begin{equation}\bm W_t+\overline{\bm A}\bm W_x=0.\tag{9.16}\end{equation}
把它看成通量 $\overline{\bm A}\bm W$ 的守恒系统，跨每个冻结特征波应用 Rankine--Hugoniot 条件即可。

% [原PDF p319]
算术平均 $\overline\rho,\overline a$ 的选择简单对称，但并不满足 Roe 那样的完整通量差性质。值得重视的是孤立接触的精确识别：当 $p_L=p_R$、$u_L=u_R$ 而密度或切向速度跳跃时，声学关系给出 $p_*=p_L$、$u_*=u_L$，密度和切向量保持各自一侧值。因此 PVRS 在光滑区和线性波附近既便宜又可靠，这正是混合求解器让它覆盖大多数界面的依据。

跨左右声学波得到
\begin{equation}
\begin{aligned}
 \overline\rho(u_*-u_L)+\overline a(\rho_{*L}-\rho_L)&=0,&
 (p_*-p_L)/\overline\rho+\overline a(u_*-u_L)&=0,\\
 \overline\rho(u_R-u_*)-\overline a(\rho_R-\rho_{*R})&=0,&
 (p_R-p_*)/\overline\rho-\overline a(u_R-u_*)&=0.
\end{aligned}\tag{9.19}
\end{equation}
故
\begin{align}
 p_*&=\tfrac12(p_L+p_R)+\tfrac12(u_L-u_R)\overline\rho\,\overline a,\\
 u_*&=\tfrac12(u_L+u_R)+\frac{p_L-p_R}{2\overline\rho\,\overline a},\\
 \rho_{*L}&=\rho_L+(u_L-u_*)\overline\rho/\overline a,\qquad
 \rho_{*R}=\rho_R+(u_*-u_R)\overline\rho/\overline a. \tag{9.20}
\end{align}
可取简单算术平均 $\overline\rho=(\rho_L+\rho_R)/2$、$\overline a=(a_L+a_R)/2$。若左右态只由孤立接触间断或剪切波相连，此解不依赖平均方式而完全精确，这对界面分辨非常重要。

也可由特征关系 $dp\mp\rho a\,du=0$、$dp-a^2d\rho=0$ 推导另一版本。

% [原PDF p320]
特征积分版本用各自特征脚点的阻抗 $C_K=\rho_Ka_K$，可看作左右声学信息的阻抗加权匹配。压力公式中 $C_Rp_L+C_Lp_R$ 是交叉权重，速度跳跃通过 $C_LC_R$ 转为压力修正；速度公式则由左右速度与压力差共同决定。若两侧阻抗相同，它退化为简单平均 PVRS。两种表达都可能在强膨胀时给出负压力，因此不能脱离自适应判据盲目使用。

令 $C_K=\rho_Ka_K$，从左右数据沿对应声学特征积分：
\begin{equation}p_*+C_Lu_*=p_L+C_Lu_L,\qquad p_*-C_Ru_*=p_R-C_Ru_R.\tag{9.26--9.27}\end{equation}
解得
\begin{align}
 p_*&=\frac{C_Rp_L+C_Lp_R+C_LC_R(u_L-u_R)}{C_L+C_R},\\
 u_*&=\frac{C_Lu_L+C_Ru_R+p_L-p_R}{C_L+C_R},\\
 \rho_{*L}&=\rho_L+(p_*-p_L)/a_L^2,qquad
 \rho_{*R}=\rho_R+(p_*-p_R)/a_R^2. \tag{9.28}
\end{align}
它不需人为选平均值；若 $C_L=C_R=\overline\rho\,\overline a$，便与前式一致。PVRS 极其便宜，也容易推广到复杂方程组；不过在很强波系中可能不够稳健，因此实际应用宜采用第9.5节的混合策略。

% [原PDF p321]
精确压力函数的两个分支具有明确物理来源：稀疏分支来自等熵 Riemann 不变量，激波分支来自 Rankine--Hugoniot 关系。函数对 $p$ 单调，精确根唯一。TRRS 和 TSRS 的核心不是对任意函数拟合，而是分别假定两侧都落在同一种真实分支；这使近似在相应极限具有正确渐近行为，也便于用 $p_*$ 与左右压力比较识别实际波型。

\section{基于精确求解器的近似}
精确压力方程为
\begin{equation}f(p)=f_L(p,\bm W_L)+f_R(p,\bm W_R)+u_R-u_L=0,\tag{9.29}\end{equation}
其中
\begin{equation}
 f_K(p)=\begin{cases}
 (p-p_K)\sqrt{A_K/(p+B_K)},&p>p_K\quad(\text{激波}),\\
 \dfrac{2a_K}{\gamma-1}\left[(p/p_K)^z-1\right],&p\le p_K\quad(\text{稀疏波}),
 \end{cases}\tag{9.30}
\end{equation}
且 $z=(\gamma-1)/(2\gamma)$、$A_K=2/[(\gamma+1)\rho_K]$、$B_K=(\gamma-1)p_K/(\gamma+1)$。

\subsection{双稀疏波黎曼求解器（TRRS）}
先验假定两支非线性波均为稀疏波，压力方程可以显式求解：
\begin{equation}
 p_*=\left[\frac{a_L+a_R-\tfrac12(\gamma-1)(u_R-u_L)}{a_L/p_L^z+a_R/p_R^z}\right]^{1/z}. \tag{9.32}
\end{equation}

% [原PDF p322]
直接由式(9.32)求压力要计算三次分数幂；先由式(9.35)求速度只需一次，再用一侧或两侧平均恢复压力，通常更高效。若实际一侧为激波，纯等熵密度式就不再正确，所以改进版本在获得近似压力后重新比较 $p_*$ 与 $p_K$：稀疏侧用等熵律，激波侧用压缩比关系。这样保留闭式压力估计的速度，又恢复真实波曲线的主要非线性。

星形速度可由任一稀疏波关系计算：
\begin{equation}
 u_*=u_L-\frac{2a_L}{\gamma-1}\left[(p_*/p_L)^z-1\right]
 =u_R+\frac{2a_R}{\gamma-1}\left[(p_*/p_R)^z-1\right]. \tag{9.33--9.34}
\end{equation}
更省运算的闭式形式是
\begin{equation}
 u_*=\frac{P_{LR}u_L/a_L+u_R/a_R+2(P_{LR}-1)/(\gamma-1)}{P_{LR}/a_L+1/a_R},
 \qquad P_{LR}=(p_L/p_R)^z. \tag{9.35}
\end{equation}
随后可用左右稀疏关系的平均恢复 $p_*$，星形密度由等熵律
\begin{equation}\rho_{*L}=\rho_L(p_*/p_L)^{1/\gamma},\qquad
\rho_{*R}=\rho_R(p_*/p_R)^{1/\gamma}.\tag{9.37}\end{equation}
给出。若获得 $p_*$ 后再按真实的激波/稀疏分支计算 $u_*$ 和密度，精度还会提高。条件
\begin{equation}f(p_{\min})\ge0,\qquad p_{\min}=\min(p_L,p_R)\tag{9.39}\end{equation}
可预先识别真实的双稀疏波情形，此时 TRRS 是精确的；它在近真空问题上尤其稳健。

% [原PDF p323]
把 $g_K(p)$ 冻结在 $p_0$，相当于对双激波压力方程作一次定点式线性化。式(9.42)分母是左右激波函数斜率之和，分子则把两侧压力和速度跳跃组合起来。选择 $p_0=\max(0,p_{\rm pvrs})$ 既利用廉价线性估计，又避免向平方根函数传入负压力。TSRS 的密度公式保持激波压缩比随压力比增加而趋向有限上限，这是简单线性密度近似没有的性质。

\subsection{双激波黎曼求解器（TSRS）}
假定两支非线性波均为激波，则
\begin{equation}
 (p-p_L)g_L(p)+(p-p_R)g_R(p)+u_R-u_L=0,qquad
 g_K(p)=\sqrt{\frac{A_K}{p+B_K}}. \tag{9.40--9.41}
\end{equation}
该式没有直接闭式解。用二次方程近似常产生多根选择或复根，甚至在精确问题有解时也会失败。较可靠的做法是先给压力估计 $p_0$，在 $g_K$ 中冻结它，得到
\begin{equation}
 p_*=\frac{g_L(p_0)p_L+g_R(p_0)p_R-(u_R-u_L)}{g_L(p_0)+g_R(p_0)}. \tag{9.42}
\end{equation}
一致的速度与激波后密度为
\begin{align}
 u_*&=\tfrac12(u_L+u_R)+\tfrac12[(p_*-p_R)g_R(p_0)-(p_*-p_L)g_L(p_0)], \tag{9.43}\\
 \rho_{*K}&=\rho_K\frac{p_*/p_K+(\gamma-1)/(\gamma+1)}{(\gamma-1)(p_*/p_K)/(\gamma+1)+1}. \tag{9.44}
\end{align}
推荐取 $p_0=\max(0,p_{\rm pvrs})$。

% [原PDF p324]
TSRS 的优势来自对强压缩波曲线形状的保留。虽然 $g_K$ 在 $p_0$ 冻结，但压力仍通过左右激波阻抗非对称加权；相比 PVRS 的线性声学阻抗，它能更好响应大压力比。近真空时实际波往往是强稀疏，双激波假设自然不如 TRRS。因而两者不是简单的“贵/便宜”排序，而是各自对应压力函数的不同区域。

典型流场中强间断只占少数界面，这使自适应求解器的平均成本主要由廉价分支决定。即使稳健分支单次成本高出数倍，只要触发比例很低，总体开销仍接近 PVRS。反过来，在强反射、爆炸初期等大范围非光滑阶段，切换比例会升高；这时额外成本正用于避免最危险的近似失效，是合理的计算投入。

TSRS 比 TRRS 更省，且只比 PVRS 略贵；除近真空外，它在更宽流动范围内通常比 PVRS 和 TRRS 准确。和 TRRS 一样，可在得到近似 $p_*$ 后改用真实波关系求 $u_*$ 与密度，以进一步提高精度。

\section{自适应黎曼求解器}
实际流场中绝大多数局部黎曼问题表示光滑流，只有激波、接触面、剪切波和其他尖锐结构附近的左右态差异很大，而近似求解器恰在这些区域最可能致命失准。混合格式的原则是：光滑区和孤立线性波附近用简单求解器，其他区域自适应切换到复杂稳健的求解器。二维 Euler、定常超声速 Euler 以及结合自适应网格的 MUSCL 计算都已验证了这一思想。

\subsection{自适应迭代求解器（AIRS）}
AIRS 在 PVRS 与第4章精确迭代求解器之间切换。仅当
\begin{equation}
 Q=p_{\max}/p_{\min}<Q_{\rm user},\qquad
 p_{\min}<p_{\rm pvrs}<p_{\max} \tag{9.46--9.48}
\end{equation}
同时满足时使用 PVRS，否则使用精确求解器。

% [原PDF p325]
两个切换条件缺一不可。仅有压力比 $Q$ 时，$p_L=p_R$ 也可能因两侧高速相撞产生双强激波，或因高速背离产生双强稀疏；仅有 $p_{\rm pvrs}$ 位于压力区间的条件，又不能排除极端压力不对称。$Q_{\rm user}=2$ 是经过大量测试的保守选择，而非严格稳定界限。提高阈值虽让更多界面走 PVRS，却很少带来显著总成本收益，因为典型流场本来已有九成以上界面平滑。

第一条件限制左右压力比；第二条件实质上限制 $u_R-u_L$。即使 $p_L\approx p_R$，很大的负速度跳跃也会产生 $p_*>p_{\max}$ 的双强激波，很大的正跳跃则产生 $p_*<p_{\min}$ 的双强稀疏波。这两类情形常出现在反射边界，不能仅凭压力比误用 PVRS。

大量测试表明 $Q_{\rm user}=2$ 足够稳健而高效；典型网格中超过90\%的局部问题仍由廉价 PVRS 处理，因此整体成本接近最便宜求解器，稳健性接近精确求解器。缺点是强分支仍含迭代。若硬件不适合变次数迭代，可固定迭代次数；得益于良好初值，一次迭代通常已能给出很准的压力。

\subsection{自适应非迭代求解器（ANRS）}
ANRS 以 PVRS 为廉价分支，以 TRRS 和 TSRS 为稳健分支。若式(9.46)--(9.47)成立，使用 PVRS；否则根据
\begin{equation}p_{\rm pvrs}\le p_{\min}\tag{9.49}\end{equation}
选择 TRRS，若不成立则选择 TSRS。

% [原PDF p326]
ANRS 流程图中的判断次序有实际意义。先计算压力比，可快速排除明显强不对称；再检查 PVRS 压力是否位于两侧压力之间，识别温和混合波型。若不满足，则 $p_{\rm pvrs}\le p_{\min}$ 指向膨胀主导，选择 TRRS；否则选择 TSRS。所有分支返回同一组星形变量接口，因此后续切向速度、音速扇判断和通量采样完全共用。

由于 ANRS 不含迭代，单个界面的指令数与控制流上界明确，适合向量化、GPU 和固定吞吐硬件。分支发散仍可能影响效率，但比可变迭代次数更容易管理。若实现为批处理，可先计算判据掩码，再分别对三类界面执行相应公式，最后合并结果；这种组织也便于统计各分支触发率和定位异常流区。

这一逻辑与精确压力函数的形状一致：低于两侧压力的星形估计对应稀疏主导，反之由激波近似更可靠。流程可以概括为：先求 $p_{\rm pvrs}$ 与压力比 $Q$；若压力比小且估计位于 $[p_{\min},p_{\max}]$ 内，直接用 PVRS；否则按式(9.49)分流到 TRRS 或 TSRS；最后以同一套切向速度、音速扇和 $S=0$ 采样规则形成 Godunov 通量。ANRS 完全无迭代，作者推荐用于实际计算。

\section{数值结果}
仍使用 $\gamma=1.4$ 的五个一维 Euler 测试、100个均匀单元、透射边界和 $C_{\rm CFL}=0.9$；因最大波速简式在启动阶段不够可靠，前5步将 CFL 数再乘0.2。初始数据与表8.1相同：修正 Sod 问题、近真空123问题、强激波问题、两强激波碰撞问题，以及把强激波问题平移到静止接触参考系的测试5。

% [原PDF p327]
比较时所有方法共享网格、CFL、边界和波速启动策略。若不同方法使用不同时间步或边界外推，曲线差异就不能单纯归因于黎曼求解器。TSRS 在全部界面使用双激波近似；ANRS 则动态选择 PVRS、TRRS 或 TSRS。两者最终都按真实波型采样，而不是直接把某个平均状态当通量，因此能保留接触和音速扇处理。

测试1检查熵满足与音速稀疏波；测试2检查低密度和近真空稳健性；测试3、4检查强间断下的稳健性与精度；测试5专门检查缓慢或静止接触的分辨。作者用 NUMERICA 的 HE--E1GODSTATE 分别运行 TSRS 和 ANRS，并与精确黎曼求解器的 Godunov 结果比较。

图9.5--9.9对应 TSRS，图9.10--9.14对应 ANRS。在绘图精度内，两组近似结果都与第6章精确 Godunov 结果没有可辨差异，因此 TSRS 和 ANRS 均适合实际应用。相较于上一章的纯 FVS，它们既保留完整黎曼波型的接触分辨，又避免每个界面都进行精确迭代。高阶扩展、源项和多维推广分别见第13--16章。

% [原PDF p328]
TSRS 数值符号与精确线重合还说明，星形压力的小误差未必直接转化成明显通量误差：波型判断、$S=0$ 采样和物理通量映射共同决定最终结果。测试1中正确处理音速扇尤为关键；若只用星形常态而不检查扇内采样，即使 $p_*$ 很准也会在零特征速度附近出错。

图9.5给出 TSRS 对测试1的结果：$x_0=0.3$、$t=0.2$。密度、速度、压力与比内能曲线的数值符号几乎落在精确实线上，向右激波、接触间断及含音速点的左稀疏扇均被正确定位。它说明即使先验采用双激波近似，随后按真实波型采样仍可得到高质量 Godunov 通量。

% [原PDF p329]
测试2中“双激波”先验与真实双稀疏结构相反，但改用真实波关系和采样后仍给出良好结果；测试3更符合压缩主导情形，TSRS 对高压星形区的估计优于线性 PVRS。两图说明求解器名称描述的是压力近似的构造，不等于最终强制采用同一种波型。

图9.6为 TSRS 对测试2的近真空结果，$x_0=0.5$、$t=0.15$；两支稀疏波和中心低密度区均与精确解吻合。图9.7为测试3，$x_0=0.5$、$t=0.012$；Mach 数极高的右激波、接触平台和左稀疏波均无上一章 FVS 中明显的激波后凹陷。

% [原PDF p330]
测试4的慢波和测试5的静止接触要求波速估计在零附近准确。状态近似求解器先得到 $u_*$，再按完整波型选择左右状态，因此接触速度不会被声速级包络取代。数值过渡仍占少量单元，因为一阶更新对移动间断有固有耗散；孤立且严格静止的接触则可达到更锐利极限。

图9.8为 TSRS 对测试4的结果，$x_0=0.4$、$t=0.035$；三个强间断的位置和平台值与精确解一致。图9.9为测试5，$x_0=0.8$、$t=0.012$；静止接触虽受一阶 Godunov 数值黏性影响，却没有传统 FVS 那样严重的额外抹宽，慢激波也保持稳定。

% [原PDF p331]
ANRS 的优势不能从光滑区曲线直接看出，因为那里三种分支结果接近；关键是它在绝大多数界面避免复杂运算，又在稀疏极端自动切换。测试2中心区的 TRRS 分支使用等熵关系，不需迭代就接近真实近真空状态。切换边界没有可见折痕，说明各求解器在温和状态下具有良好一致性。

图9.10、9.11分别为 ANRS 对测试1和测试2的结果。前者在普通 Sod 型状态上多数界面走 PVRS 廉价分支，仅强波附近切换；后者在近真空区域自动选择 TRRS。两幅图都表明切换没有产生可见的通量不连续或伪振荡，数值点与精确曲线在绘图尺度上重合。

% [原PDF p332]
强波附近 TSRS 分支只占少数界面，随激波移动；其余平台继续由 PVRS 处理。若输出每步分支掩码，可以看到成本集中在间断窄邻域。图中没有因分支改变产生额外振荡，说明星形状态接口与采样逻辑统一。高阶重构扩大局部数据变化范围后，切换比例可能与一阶方法不同。

图9.12、9.13分别展示 ANRS 对强激波测试3和碰撞激波测试4的结果。压力比和 PVRS 星形压力判据会在强间断附近选择 TSRS，而光滑平台仍使用 PVRS。密度、速度、压力、比内能均正确再现精确波速与星形状态，说明局部切换同时获得了效率和强健性。

% [原PDF p333]
测试5中接触两侧压力比接近1，PVRS 星形压力位于区间内，因此接触附近通常使用能精确识别孤立线性波的廉价分支；强激波附近再切换 TSRS。这种按局部波型分工，正是混合方法优于全域采用单一昂贵或单一脆弱求解器的原因。

图9.14为 ANRS 对测试5的结果，$x_0=0.8$、$t=0.012$。缓慢激波与静止接触附近的数值点紧随精确曲线，没有明显过冲。结合前四个测试可见，单纯比较平均误差不足以评价黎曼求解器；还必须检查近真空、音速点、强激波和低速接触等不同特征场。

% [原PDF p334]
\chapter{HLL 与 HLLC 黎曼求解器}
HLL 的推导只依赖最外信号包络，因此常被视为复杂双曲系统的稳健起点：即便完整特征分解难以获得，也常能估计最快左右速度。代价是包络内部所有结构平均化。HLLC 展示了一种通用改进范式：从系统物理中挑选最重要的缺失中间波，增加相应星形状态和跳跃条件，在不恢复全部精确解的情况下显著降低特定耗散。

对 Euler 方程，中间速度场同时承载密度接触、两项切向剪切和被动组分，所以恢复一支速度 $S_*$ 的收益很大。对 MHD 或多相系统，不同中间结构可能以多个不同速度传播，仅增加一支“接触”波未必完备。评价某个 HLLC 扩展时，应明确它保留了哪些真实特征场、忽略了哪些场，而不能只凭名称假定性能与 Euler HLLC 相同。

Harten、Lax 与 van Leer 于1983年提出 HLL 近似求解器。它只需估计从初始间断发出的最左和最右信号速度，以两支波、三个常状态近似真实波系。Toro 等人于1992年提出 HLLC，在模型中恢复第三支中间波，因此能更准确地分辨接触间断和剪切波。

\section{引言}
HLL 与第4、9章不同，它不先构造某个近似状态再代入物理通量，而是直接由守恒律积分形式推导界面数值通量。给定两支包络波速后，控制体积分可闭式确定中间平均状态和通量。Davis 与 Einfeldt 分别提出实用波速估计；采用 Einfeldt 估计的两波格式通常称为 HLLE。HLL 类方法高效而稳健。

但两波假设只对含两个方程的系统（例如一维浅水方程）完全符合真实结构。Euler 方程以及沿某方向分裂的二维浅水方程含有额外的线性中间波，纯 HLL 会把这些波平均掉，造成接触面、材料界面和剪切层过度扩散。HLLC 名称中的 C 表示恢复 contact（接触）波；它在保持 HLL 稳健框架的同时补回缺失中间结构。

% [原PDF p335]
两波模型的中间状态是整个真实波扇的积分平均，因此它对总守恒量十分可靠，却不知道平均内部由哪些波分隔。对快速外声学波，这种信息损失通常不严重；对速度接近零的接触、剪切或材料界面，数值耗散会在许多时间步中持续作用而波本身几乎不离开原单元，最终形成很宽的混合层。HLLM/HLLEM 用线性修正降低这部分耗散，HLLC 则通过显式恢复中间波从结构上解决问题。

对于相对网格静止的接触、剪切或材料界面，HLL 的抹宽尤其不可接受。Einfeldt 曾提出 HLLM/HLLEM：以线性分布修正单一中间状态，并用参数控制对中间波的过量耗散；在特定参数和波速选择下，它可化为修正 Roe 格式。

Toro、Spruce 与 Speares 采取另一条路线，直接在 HLL 的两支包络波之间加入速度 $S_*$ 的中间波，于是星形区分成左右两个状态。这一 HLLC 构造随后用于定常二维超声速 Euler 方程、非定常二维浅水方程、隐式湍流计算等。对 Euler 方程，空间维数增加并不增加不同特征速度的数目：仍只有 $u_n-a$、$u_n$、$u_n+a$ 三个不同场，所以三波 HLLC 是“完备”的。若系统有三个以上不同特征场，普通 HLLC 仍是不完备求解器，低速中间波会受额外耗散；改善方法是建立包含该系统全部主要特征场的波模型。

% [原PDF p336]
HLL 推导的优点是只使用守恒律积分形式，不需逐一知道完整特征矢量。对复杂系统，只要能可靠包络最左、最右信号，就能迅速得到稳健通量。HLLC 多加一个中间波，仍保持这种简洁性。以下所有公式都从时空控制体的收支出发，所以中间状态不是随意插值，波速改变时也必须同时按一致性关系改变状态与通量。

本章针对三维非定常 Euler 方程推导 HLL/HLLC，但原理可直接迁移到其他双曲系统。以下先回顾黎曼问题与积分关系，再推导 HLL 和 HLLC，讨论波速估计、算法汇总、被动标量与数值结果。

\section{黎曼问题}
在 $0\le x\le L$ 上考虑
\begin{equation}
 \bm U_t+\bm F(\bm U)_x=0,\quad \bm U(x,0)=\bm U^{(0)}(x),\quad
 \bm U(0,t)=\bm U_l(t),\quad\bm U(L,t)=\bm U_r(t), \tag{10.1}
\end{equation}
用守恒格式
\begin{equation}
 \bm U_i^{n+1}=\bm U_i^n-\frac{\Delta t}{\Delta x}
 (\bm F_{i+1/2}-\bm F_{i-1/2}). \tag{10.2}
\end{equation}

\subsection{Godunov 通量}
精确 Godunov 通量为
\begin{equation}\bm F_{i+1/2}=\bm F(\bm U_{i+1/2}(0)),\tag{10.3}\end{equation}
即对左右常状态局部黎曼解在 $S=x/t=0$ 采样。沿 $x$ 分裂的三维 Euler 守恒量与通量仍是式(9.5)。

% [原PDF p337]
控制体左右边界必须位于最外信号之外，才能保证在 $x=x_L$、$x=x_R$ 处状态始终分别为初始左、右态。时间 $T$ 可以任意取正；由于黎曼解自相似，积分平均与 $T$ 无关。这个推导没有假定内部是激波还是稀疏波，也没有使用理想气体 EOS，因而 HLL 公式具有很强的系统通用性。

精确波系有五族波：$u-a$，三重的 $u$，以及 $u+a$。第9章先近似采样状态再计算物理通量；本章沿 Harten--Lax--van Leer 思路直接近似通量。

\subsection{积分关系}
取足以包住整个波扇的时空控制体 $[x_L,x_R]\times[0,T]$，满足
\begin{equation}x_L\le TS_L,qquad x_R\ge TS_R.\tag{10.7}\end{equation}
守恒律积分形式给出
\begin{equation}
 \int_{x_L}^{x_R}\bm U(x,T)\,dx=x_R\bm U_R-x_L\bm U_L+T(\bm F_L-\bm F_R). \tag{10.9}
\end{equation}
这称为一致性条件，且 $\bm F_K=\bm F(\bm U_K)$。

% [原PDF p338]
式(10.11)把波扇内未知积分完全换成已知端点状态、端点通量和包络速度。除以扇宽后得到的 $\bm U_{\rm hll}$ 是时刻 $T$ 的空间平均，也等价于自相似坐标 $S$ 上的平均。它一般不等于真实解在 $S=0$ 的状态；HLL 的思想是让一个常状态代表整个平均，再通过跨外波守恒确定界面通量。

把积分区间按两支最外波分成三段，左右外区分别保持 $\bm U_L$ 与 $\bm U_R$，比较后得到
\begin{equation}
 \int_{TS_L}^{TS_R}\bm U(x,T)\,dx
 =T(S_R\bm U_R-S_L\bm U_L+\bm F_L-\bm F_R). \tag{10.11}
\end{equation}
因此波扇内精确解的积分平均为
\begin{equation}
 \overline{\bm U}=\frac1{T(S_R-S_L)}\int_{TS_L}^{TS_R}\bm U(x,T)\,dx
 =\frac{S_R\bm U_R-S_L\bm U_L+\bm F_L-\bm F_R}{S_R-S_L}
 \equiv\bm U_{\rm hll}. \tag{10.12--10.13}
\end{equation}
只要最外信号速度已知，这个平均状态就是已知常量。

% [原PDF p339]
左右半控制体公式还说明了为什么数值通量必须从积分关系而非 $\bm F(\bm U_{\rm hll})$ 取得。后者通常不同时满足两侧 Rankine--Hugoniot 条件，而式(10.18)、(10.19)在一致的 $\bm U_{\rm hll}$ 下完全相等。选择任一表达实现都应得到相同结果；若浮点结果差异显著，通常意味着波速次序或状态量出现问题。

在左、右半控制体上分别积分，还可写出 $t$ 轴通量
\begin{align}
 \bm F_{0L}&=\bm F_L-S_L\bm U_L-\frac1T\int_{TS_L}^{0}\bm U(x,T)\,dx, \tag{10.15}\\
 \bm F_{0R}&=\bm F_R-S_R\bm U_R+\frac1T\int_{0}^{TS_R}\bm U(x,T)\,dx. \tag{10.16}
\end{align}
令两者相等就恢复一致性条件。至此没有作近似。

\section{HLL 近似黎曼求解器}
HLL 用三个常状态表示解：
\begin{equation}
 \widetilde{\bm U}(x,t)=\begin{cases}
 \bm U_L,&x/t\le S_L,\\
 \bm U_{\rm hll},&S_L\le x/t\le S_R,\\
 \bm U_R,&x/t\ge S_R.
 \end{cases}\tag{10.17}
\end{equation}
它把全部中间波和状态合并成单一 $\bm U_{\rm hll}$。注意界面通量不是简单的 $\bm F(\bm U_{\rm hll})$，而应由积分关系得到。

% [原PDF p340]
式(10.21)自动覆盖三种传播格局：若 $S_L\ge0$，整个波扇在界面右侧，界面只看左物理通量；若 $S_R\le0$，整个波扇在左侧，取右物理通量；只有包络跨越零时才用 HLL 混合式。实现应先检查两个超声速分支，避免在 $S_R-S_L$ 很小时无意义地做除法。

在非平凡的亚声速情形 $S_L\le0\le S_R$，
\begin{equation}
 \bm F_{\rm hll}=\bm F_L+S_L(\bm U_{\rm hll}-\bm U_L)
 =\bm F_R+S_R(\bm U_{\rm hll}-\bm U_R), \tag{10.18--10.19}
\end{equation}
即
\begin{equation}
 \bm F_{\rm hll}=\frac{S_R\bm F_L-S_L\bm F_R+S_LS_R(\bm U_R-\bm U_L)}{S_R-S_L}. \tag{10.20}
\end{equation}
完整界面通量为
\begin{equation}
 \bm F_{i+1/2}^{\rm hll}=\begin{cases}
 \bm F_L,&0\le S_L,\\
 \bm F_{\rm hll},&S_L\le0\le S_R,\\
 \bm F_R,&0\ge S_R.
 \end{cases}\tag{10.21}
\end{equation}
给定 $S_L,S_R$ 就形成一个近似 Godunov 方法。Harten、Lax 与 van Leer 证明，在适当条件下该格式收敛到满足熵条件的弱解；关键要求之一是近似解仍满足积分一致性。

% [原PDF p341]
“与积分形式一致”意味着把近似解代回同一时空控制体后，右端由初始数据决定的净通量不能改变。这个条件使近似波系即使省略内部细节，也保持正确的总体收支。HLLC 的目标不是破坏 HLL 平均，而是在平均约束之内重新分配星形区，使左右子平均由一支中间波分开，因此仍继承原 HLL 的一致性基础。

HLL 的缺点来自单一星形状态：积分平均只保留波扇总量，不保留中间空间变化。Euler 方程中接触、两支剪切波均对应多重特征值 $u$，会被 HLL 额外耗散。修复的自然办法就是恢复缺失波，这导向 HLLC。

\section{HLLC 近似黎曼求解器}
\subsection{若干有用关系}
在最外波 $S_L,S_R$ 之间加入中间波 $S_*$。把波扇平均积分在 $S_*$ 处分开，定义
\begin{align}
 \bm U_{*L}&=\frac1{T(S_*-S_L)}\int_{TS_L}^{TS_*}\bm U(x,T)\,dx,\\
 \bm U_{*R}&=\frac1{T(S_R-S_*)}\int_{TS_*}^{TS_R}\bm U(x,T)\,dx. \tag{10.23}
\end{align}

% [原PDF p342]
式(10.24)是 HLLC 的核心约束：$\bm U_{*L}$、$\bm U_{*R}$ 按各自自相似宽度加权后，必须恰好给出 $\bm U_{\rm hll}$。因此增加中间波并没有凭空改变波扇总平均。三条跳跃条件中任意两条配合一致性可推出第三条；实际推导保留全部三式，便于检查每支近似波都满足守恒。

一致性要求
\begin{equation}
 \frac{S_*-S_L}{S_R-S_L}\bm U_{*L}
 +\frac{S_R-S_*}{S_R-S_L}\bm U_{*R}=\bm U_{\rm hll}. \tag{10.24}
\end{equation}
HLLC 近似解为四个常状态
\begin{equation}
 \widetilde{\bm U}=\begin{cases}
 \bm U_L,&S\le S_L,\\\bm U_{*L},&S_L\le S\le S_*,\\
 \bm U_{*R},&S_*\le S\le S_R,\\\bm U_R,&S\ge S_R,
 \end{cases}\tag{10.25}
\end{equation}
对应通量按 $S=0$ 所落区间从 $\bm F_L,\bm F_{*L},\bm F_{*R},\bm F_R$ 中选择。跨三支波的 Rankine--Hugoniot 条件为
\begin{align}
 \bm F_{*L}&=\bm F_L+S_L(\bm U_{*L}-\bm U_L),\tag{10.27}\\
 \bm F_{*R}&=\bm F_{*L}+S_*(\bm U_{*R}-\bm U_{*L}),\tag{10.28}\\
 \bm F_{*R}&=\bm F_R+S_R(\bm U_{*R}-\bm U_R).\tag{10.29}
\end{align}

% [原PDF p343]
由式(10.33)--(10.34)逐分量求解时，质量方程先给出星形密度比例，法向动量方程再给星形压力，切向动量方程证明 $v_{*K}=v_K,w_{*K}=w_K$ 与闭合假设一致，能量方程最后确定星形总能量。这个顺序揭示了式(10.39)每一分量的来源，也提供实现单元测试：各星形状态代回跨 $S_K$ 的跳跃关系后，五个守恒分量残差都应接近舍入误差。

中间波上的 Rankine--Hugoniot 条件要求 $\bm F_{*R}-\bm F_{*L}=S_*(\bm U_{*R}-\bm U_{*L})$。在压力和法向速度连续时，它对质量、法向动量和能量自动成立；切向分量允许左右跳跃并随 $S_*$ 传播。由此可见，HLLC 恢复的不是一个人为“修正项”，而是一个满足守恒、具有 Euler 接触波不变量的近似波。

Euler 精确接触波上法向速度和压力连续，而密度、切向速度可跳变。HLLC 采用这些真实性质作为闭合条件，并把中间波速直接设为星形法向速度。这不是假设两侧星形状态完全相同；恰恰相反，左右密度和切向量分别保留。式(10.36)还表明星形压力由每侧质量通量因子 $\rho_K(S_K-u_K)$ 与速度改变量共同确定。

\subsection{Euler 方程的 HLLC 通量}
仿照精确接触波条件，施加
\begin{equation}p_{*L}=p_{*R}=p_*,\quad u_{*L}=u_{*R}=u_*,\quad
v_{*K}=v_K,\quad w_{*K}=w_K,\quad S_*=u_*.\tag{10.30--10.32}\end{equation}
又有
\begin{equation}\bm F(\bm U)=u\bm U+p\bm D,\qquad \bm D=(0,1,0,0,u)^T.\tag{10.35}\end{equation}
由左右波的质量、动量关系可得
\begin{equation}
 p_{*K}=p_K+\rho_K(S_K-u_K)(S_*-u_K),\qquad K=L,R. \tag{10.36}
\end{equation}
令左右星形压力相等就能只用 $S_L,S_R$ 表出 $S_*$。

% [原PDF p344]
中间速度公式的分母是左右波的有效质量通量之差。对物理波速估计，通常有 $S_L<u_L$、$S_R>u_R$，使分母符号稳定。式(10.39)的前因子就是星形密度与原密度之比；能量分量同时计入法向动能改变和压力功。实现时应从守恒状态恢复合法的 $p_K,H_K$，并保持左右索引 $K$ 一致，尤其不能把式中某一侧密度误写为固定左值。

具体为
\begin{equation}
 S_*=\frac{p_R-p_L+\rho_Lu_L(S_L-u_L)-\rho_Ru_R(S_R-u_R)}
 {\rho_L(S_L-u_L)-\rho_R(S_R-u_R)}. \tag{10.37}
\end{equation}
于是
\begin{equation}
 \bm F_{*K}=\bm F_K+S_K(\bm U_{*K}-\bm U_K), \tag{10.38}
\end{equation}
\begin{equation}
 \bm U_{*K}=\rho_K\frac{S_K-u_K}{S_K-S_*}
 \begin{bmatrix}
 1\\S_*\\v_K\\w_K\\
 E_K/\rho_K+(S_*-u_K)\left[S_*+\dfrac{p_K}{\rho_K(S_K-u_K)}\right]
 \end{bmatrix}. \tag{10.39}
\end{equation}
最后仍按 $S_L,S_*,S_R$ 的符号选取通量。另一等价实现可由 $\bm D_*=(0,1,0,0,S_*)^T$ 直接写出星形通量；它可放松左右近似压力严格相等的条件，更贴合式(10.36)分别给出的近似。

% [原PDF p345]
三个 HLLC 变体的差异集中在星形压力处理。第一种严格施加 $p_{*L}=p_{*R}$；第二种保留式(10.36)两侧各自近似；第三种用算术平均 $P_{LR}$。在波速估计足够准确时差异很小，但复杂 EOS 或极端数据下，代数形式可能影响舍入误差与稳健性。无论选哪种，都必须按同一 Rankine--Hugoniot 关系构造状态和通量。

还可用左右星形压力的算术平均
\begin{equation}
 P_{LR}=\tfrac12[p_L+p_R+\rho_L(S_L-u_L)(S_*-u_L)+\rho_R(S_R-u_R)(S_*-u_R)] \tag{10.42}
\end{equation}
并令
\begin{equation}
 \bm U_{*K}=\frac{S_K\bm U_K-\bm F_K+P_{LR}\bm D_*}{S_K-S_*},\qquad
 \bm F_{*K}=\frac{S_*(S_K\bm U_K-\bm F_K)+S_KP_{LR}\bm D_*}{S_K-S_*}. \tag{10.43--10.44}
\end{equation}
以上代数不依赖理想气体状态方程；EOS 只通过最外波速估计进入。

\subsection{多维与多组分流}
一般多维网格先把速度分解成界面法向分量 $u_n$ 和两个切向分量，方程形式便与沿 $x$ 分裂完全相同。对 $m$ 个由载体流输运的组分浓度 $q_l$，非守恒对流式 $q_{l,t}+u q_{l,x}=0$ 与连续方程合并为
\begin{equation}(\rho q_l)_t+(\rho u q_l)_x=0,\qquad l=1,\ldots,m.\tag{10.45a}\end{equation}

% [原PDF p346]
组分守恒式由非守恒对流方程与连续方程组合得到，这一点保证了跨密度间断时组分质量仍严格守恒。附加每个 $q_l$ 只增加中间速度 $u$ 的重数，不产生新声学速度。HLLC 恢复这一中间场后，所有组分都按载体质量通量迎风；HLL 的单平均态则会把左右浓度混合。对无反应多组分问题，这种差异直接决定材料界面的锐利度。

扩充后的中间特征速度 $u$ 的重数为 $m+3$。只需把 $\rho q_l$ 和 $\rho uq_l$ 分别附加到守恒量、通量矢量，HLLC 星形状态中相应分量就是
\begin{equation}
 (\rho q_l)_{*K}=\rho_K\frac{S_K-u_K}{S_K-S_*}(q_l)_K. \tag{10.46}
\end{equation}
因此无需特殊补丁，HLLC 就像精确求解器一样分辨这些附加中间场。切向速度 $v,w$ 正是被动标量的特殊情况。

\section{波速估计}
HLL 与 HLLC 都需要 $S_L,S_R$；HLLC 的 $S_*$ 随后由式(10.37)算出。可直接估计最小、最大波速，也可先估计星形压力，再由精确波关系恢复包络速度。

% [原PDF p347]
简单 Davis 估计未必真正包住非线性强激波后的最快信号，因为只看初始态声速。取两侧最小/最大特征值虽稍保守，仍可能在极端问题中不足。Roe 平均速度和声速包含左右状态耦合，常能更好表示星形区；Einfeldt 估计又加入速度跳跃贡献，目的在于扩大包络并保持密度、内能正性。波速选择决定 HLL 类通量的耗散量和安全性。

\subsection{直接波速估计}
Davis 的简单建议包括
\begin{equation}S_L=u_L-a_L,\qquad S_R=u_R+a_R,\tag{10.47}\end{equation}
以及更保守的
\begin{equation}S_L=\min(u_L-a_L,u_R-a_R),\qquad
S_R=\max(u_L+a_L,u_R+a_R).\tag{10.48}\end{equation}
它们只用初始数据，虽便宜却不推荐作为通用实用选择。Davis 与 Einfeldt 还建议 Roe 平均
\begin{equation}S_L=\widetilde u-\widetilde a,qquad S_R=\widetilde u+\widetilde a,\tag{10.49}\end{equation}
其中
\begin{equation}
 \widetilde u=\frac{\sqrt{\rho_L}u_L+\sqrt{\rho_R}u_R}{\sqrt{\rho_L}+\sqrt{\rho_R}},\quad
 \widetilde H=\frac{\sqrt{\rho_L}H_L+\sqrt{\rho_R}H_R}{\sqrt{\rho_L}+\sqrt{\rho_R}},\quad
 \widetilde a=\sqrt{(\gamma-1)(\widetilde H-\widetilde u^2/2)}. \tag{10.50--10.51}
\end{equation}
Einfeldt 的 HLLE 估计采用 $S_{L,R}=\overline u\mp\overline d$，其中 $\overline d$ 还包含左右速度跳跃项，用以增强正性和强健性。

% [原PDF p348]
正性分析提醒，较窄波速包络虽然降低耗散，却可能让星形平均状态落出物理可接受集合。Einfeldt 型估计用更保守的信号范围换取密度和内能正性，这在强稀疏和近真空尤其重要。HLLC 在 HLL 星形平均内重新分配左右状态，通常继承良好稳健性，但具体正性仍依赖波速与状态公式；不能因为名称含 HLL 就省略极端状态测试。

Rusanov 通量可解释为假设所有特征场都以同一最大速度传播。它完全避免特征分解，适合复杂系统的基准实现和强健后备分析，但会对慢波施加声速级耗散。Lax--Friedrichs 又用全局网格时间比代替局部物理最大速度，通常更耗散。三者形成从局部双边包络、局部对称包络到全局对称包络的层级。

对称包络把 HLL 化成局部 Lax--Friedrichs/Rusanov 型通量，这揭示两种方法的关系：中心平均负责一致性，$S^+$ 乘状态跳跃提供耗散。$S^+$ 取太小会漏掉真实信号并破坏稳定性，取太大则把所有波过度抹宽。用 $\max(|u|+a)$ 比逐项取四个绝对特征值更简洁且至少同样保守。全局 $\Delta x/\Delta t$ 则进一步退化为耗散很强的 Lax--Friedrichs。

所谓正性守恒求解器，是指对物理可接受初始态，计算中密度与内能保持正。波速包络是否足够宽对此至关重要。

若取对称速度 $S_L=-S^+$、$S_R=S^+$，HLL 通量化为 Rusanov 通量
\begin{equation}
 \bm F_{i+1/2}=\tfrac12(\bm F_L+\bm F_R)-\tfrac12S^+(\bm U_R-\bm U_L), \tag{10.55}
\end{equation}
稳健选择为
\begin{equation}S^+=\max(|u_L|+a_L,|u_R|+a_R).\tag{10.56}\end{equation}
若令 $S^+=S_{\max}^n=C_{\rm CFL}\Delta x/\Delta t$，且 $C_{\rm CFL}=1$，则得到 Lax--Friedrichs 通量
\begin{equation}
 \bm F_{i+1/2}=\tfrac12(\bm F_L+\bm F_R)-\frac{\Delta x}{2\Delta t}(\bm U_R-\bm U_L). \tag{10.58}
\end{equation}

% [原PDF p349]
压力型估计先回答每侧外波是压缩还是膨胀。若 $p_*\le p_K$，该侧为稀疏波，扇头速度就是数据特征速度，所以 $q_K=1$；若 $p_*>p_K$，该侧为激波，$q_K>1$ 扩大相对于质点速度的传播幅度。即使 $p_*$ 只是近似，这种分支仍比固定使用数据特征速度更能适应强压缩。

\subsection{基于压力的波速估计}
先取得星形压力近似 $p_*$，再选
\begin{equation}S_L=u_L-a_Lq_L,\qquad S_R=u_R+a_Rq_R,\tag{10.59}\end{equation}
其中
\begin{equation}
 q_K=\begin{cases}1,&p_*\le p_K,\\
 \sqrt{1+\dfrac{\gamma+1}{2\gamma}(p_*/p_K-1)},&p_*>p_K.
 \end{cases}\tag{10.60}
\end{equation}
若第 $K$ 支波是稀疏波，$S_K$ 就是稀疏扇头的精确最快信号；若是激波，则使用精确激波速度关系，但其中的星形压力比是近似的。

最简单的 PVRS 压力为
\begin{equation}
 p_{\rm pvrs}=\tfrac12(p_L+p_R)-\tfrac12(u_R-u_L)\overline\rho\,\overline a,qquad
 \overline\rho=\tfrac12(\rho_L+\rho_R),\quad\overline a=\tfrac12(a_L+a_R). \tag{10.61--10.62}
\end{equation}
也可用第9章 TRRS 或 TSRS 压力。

% [原PDF p350]
完整 ANRS 压力判据可在普通区用 PVRS、近真空膨胀用 TRRS、强压缩用 TSRS；若追求最小实现，截断到非负的 PVRS 压力通常已足够形成 HLLC 波速。这里的 $\max(0,p_{\rm pvrs})$ 是波速估计中的算法选择，不等于在物理解更新后静默钳制压力；后者会掩盖失效，前者只选择包络速度公式的输入。

计算实践中可采用第9.5.2节 ANRS 自适应选择 $p_*$。对理想气体，甚至 $p_*=\max(0,p_{\rm pvrs})$ 这一简化就通常足够稳健。

\section{HLLC 通量算法汇总}
推荐算法如下。

\textbf{步骤I：压力估计。}
\begin{equation}
 p_*=\max(0,p_{\rm pvrs}),\quad
 p_{\rm pvrs}=\tfrac12(p_L+p_R)-\tfrac12(u_R-u_L)\overline\rho\,\overline a. \tag{10.67}
\end{equation}
\textbf{步骤II：波速。}用式(10.59)--(10.60)求 $S_L,S_R$，再由
\begin{equation}
 S_*=\frac{p_R-p_L+\rho_Lu_L(S_L-u_L)-\rho_Ru_R(S_R-u_R)}
 {\rho_L(S_L-u_L)-\rho_R(S_R-u_R)} \tag{10.70}
\end{equation}
求中间速度。

% [原PDF p351]
算法顺序不可颠倒：$p_*$ 决定 $q_K$，$q_K$ 决定外波速，外波速又决定 $S_*$ 和星形状态。通量选择必须在完成三支速度后进行。四个分支互斥且覆盖全部次序正确的波速情形。若计算得到 $S_L>S_R$ 或星形密度非物理，应诊断波速估计或输入状态，而不能任意交换速度以继续计算。

\textbf{步骤III：通量选择。}
\begin{equation}
 \bm F_{i+1/2}^{\rm hllc}=\begin{cases}
 \bm F_L,&0\le S_L,\\
 \bm F_{*L},&S_L\le0\le S_*,\\
 \bm F_{*R},&S_*\le0\le S_R,\\
 \bm F_R,&0\ge S_R,
 \end{cases}\tag{10.71}
\end{equation}
其中 $\bm F_{*K}=\bm F_K+S_K(\bm U_{*K}-\bm U_K)$，$\bm U_{*K}$ 用式(10.39)。也可直接采用式(10.43)--(10.44)的两个变体；在本章测试条件和相同波速估计下，三种实现给出相同结果。

\section{接触波与被动标量}

% [原PDF p352]
被动标量例子把 HLL 的中间波耗散定量化。声速 $a$ 本不应影响纯对流跳跃的精确传播，但 HLL 用 $u\pm a$ 包络整个 Euler 波扇，于是标量通量含 $1/M$；低 Mach 数越小，额外耗散越强。HLLC 恰好把标量放回速度 $u$ 的中间波，因此耗散随实际平移速度变化，静止时不扩散。这也是低 Mach 材料界面计算偏好完整求解器的原因。

考虑随流速输运的被动标量 $q$：
\begin{equation}(\rho q)_t+(\rho qu)_x=0.\tag{10.77}\end{equation}
若 $p,\rho,u$ 均为常数，而 $q$ 在原点有跳跃，则精确解只是平移
\begin{equation}q(x,t)=q_0(x-ut).\tag{10.79}\end{equation}
取 $S_L=u-a,S_R=u+a$，亚声速下 HLL 数值通量为
\begin{equation}
 f_{i+1/2}^{\rm hll}=\tfrac12(1+1/M)f_i+\tfrac12(1-1/M)f_{i+1},\qquad M=u/a. \tag{10.80}
\end{equation}
由于 $|M|<1$ 时 $|1/M|>1$，它比精确 Godunov 通量更耗散；若某一参数组合使它等价于 Lax--Friedrichs，耗散尤其大，静止标量是最坏极限。相同结论适用于孤立接触以及三维切向速度 $v,w$。

HLLC 则像精确求解器一样传播中间波；静止接触在理想离散中可保持无限分辨率。因而材料界面、剪切波和涡结构上，HLL 会增加过量黏性，HLLC 明显更合适。

% [原PDF p353]
七个测试把强健性与中间波精度分开考察。测试1--5含声学波和接触的组合，测试6、7去掉所有声学扰动，只留下密度跳跃，使压力、速度恒定。后一对测试若出现密度变化，就只能来自数值通量本身。HLLC 三个变体在这些理想气体算例上结果一致，说明主要收益来自恢复中间波，而非某个星形压力公式的细小差别。

\section{数值结果}
本节以 $\gamma=1.4$、100个单元、$C_{\rm CFL}=0.9$ 和透射边界比较 HLL、HLLC 与 Rusanov。HLLC 使用式(10.38)--(10.39)，两者的 $S_L,S_R$ 均由第9章自适应压力估计形成。测试1--5与前两章相同；另加测试6：$(\rho_L,u_L,p_L)=(1.4,0,1)$、$(\rho_R,u_R,p_R)=(1,0,1)$，表示静止孤立接触；测试7把两侧速度同时改为0.1，表示缓慢右移接触。

精确解由 HE--E1RPEXACT 产生，数值解由 HE--E1GODFLUX 产生。测试1检查音速稀疏波熵性质；测试2检查近真空；测试3、4检查强激波；测试5检查复杂波系中的静止接触；测试6、7专门隔离接触传播误差。

% [原PDF p354]
测试5中的接触并非孤立：它与稀疏波、慢激波同处一个黎曼演化，所以 HLLC 的分辨虽好，也受一阶更新和邻近波影响。测试6的静止孤立接触更纯粹，精确通量在两侧完全相同；HLLC 因中间波速为零且压力连续，界面更新不改变密度跃变。测试7只给接触一个很小正速度，数值扩散随实际 Courant 数出现，而不再与声速直接相连。

音速稀疏波结果也揭示“精确黎曼求解器”并不保证离散曲线在每个指标上最优：一阶 Godunov 采样、时间步和网格共同决定数值扩散。某些近似通量的平滑波速过渡可能在有限网格上看起来更好，但这不改变其连续黎曼模型的近似性质。评价应同时考虑收敛性、熵、正性和不同波族，而非只看单幅图。

HLL 和 HLLC 对强声学波相近，是因为它们共享同样的最外波速估计；差异主要在两外波之间。Rusanov 又把左右包络压成一个对称最大速度，所以对中间场更耗散，也可能对音速稀疏波施加不必要的对称黏性。观察测试6、7能避免被强激波图中的视觉相似误导：求解器是否完备，最清楚地体现在孤立慢波上。

HLL 与 HLLC 的总体结果都接近精确黎曼求解器的 Godunov 方法。两者对测试1的音速稀疏波甚至比这里采用的精确求解器采样更平滑。测试5中 HLLC 对静止接触的分辨与精确 Godunov 相当，而 HLL 像 Steger--Warming 与 van Leer FVS 一样把接触抹宽。Rusanov 与 HLL 总体相近，也扩散低速接触；但在含音速稀疏波的测试1中，Rusanov 明显不如 HLL。

测试6、7进一步证实：HLLC 对静止接触保持无限分辨，对慢接触的数值耗散也很小；HLL 则产生不可接受的宽化。相同结论适用于带组分方程的扩充系统，以及多维剪切波和涡。

% [原PDF p355]
多维推广时，HLLC 每个面都在局部法向坐标计算。法向速度决定三支波，两个切向速度连同被动组分作为中间场携带；得到通量后再旋回全局动量分量。若旋转矩阵、法向方向或左右态约定不一致，会破坏相邻单元共享通量的反对称性。结构和非结构网格虽在几何数据上不同，黎曼核心可以完全复用。

对极高速流撞击固壁，镜像边界会形成极强压缩黎曼问题，一些常用平均波速可能低估真实反射激波速度。此时应采用能包络信号并通过专门壁面测试验证的估计。对更多特征场系统，提升 HLLC 的正确方向是明确增加波和状态，而不是在三波公式外随意调整耗散参数；后者可能改善某些图，却不恢复缺失物理结构。

\section{结束语}
HLLC 在三维和多组分 Euler 流中能够像精确黎曼求解器一样处理法向速度为 $u_n$ 的中间特征场；切向速度和各组分浓度都随这一场传播，所以一支恢复的接触波足以正确捕获接触、剪切和组分界面。HLL 则缺少此场。

这些通量可直接用于第6章一阶 Godunov 法，第13、14章给出 TVD 二阶扩展，第15章处理源项，第16章推广至三维。HLLC 还常作为结构或非结构网格上的半离散、全离散有限体积与间断 Galerkin 高阶方法基本模块。需要特别留意两个问题：极高速度撞击静止固壁时，一些常用波速估计可能失效；对多相流、MHD 等具有三个以上不同特征场的系统，应建立包含更多真实波的 HLLC 型模型，而不能默认三波结构已经完备。

% [原PDF p356]
HLLC 在测试1、2中的曲线验证了压力型波速估计能够同时处理跨音速稀疏和近真空双稀疏。最外速度若过窄会漏掉扇头并失稳，过宽则增加耗散。图中外波位置与精确解一致，说明 $q_K$ 分支有效；接触分辨的提升来自恢复中间波，并非通过危险地缩窄外包络获得。

图10.5、10.6给出 HLLC 对测试1和测试2的结果。测试1为 $x_0=0.3,t=0.2$，音速稀疏波、接触与激波位置准确；测试2为 $x_0=0.5,t=0.15$，近真空中心区保持稳定。四个子图依次是密度、速度、压力、比内能，符号为数值解，实线为精确解。

% [原PDF p357]
强激波测试中 $q_K>1$，外波速扩展到数据声速之外，使控制体包住真实压缩波。测试4的复杂全局解仍由各界面局部三波模型组合产生，求解器无需预先知道波相互作用。星形平台准确说明跨外波的质量、动量和能量跳跃在大压力比下仍保持一致。

图10.7、10.8给出 HLLC 对测试3和测试4的结果，输出时刻分别为0.012和0.035。Mach 数198的强激波以及三条向右强间断均被稳健捕获，星形区平台与精确解重合，未见显著激波后振荡。

% [原PDF p358]
HLL 对测试1的声学部分良好，因为两支外波本来就保留；接触差异来自单一中间平均。测试5中这种平均混合密度和比内能跃变，但压力、速度仍较平滑，所以只看压力图可能误判 HLL 与 HLLC 等价。材料界面问题应重点检查密度、组分和切向速度等中间场变量。

图10.9为 HLLC 对测试5的结果，$x_0=0.8,t=0.012$，静止接触仍较锐利。图10.10为 HLL 对测试1：外侧声学波与 HLLC 相近，但接触过渡更宽。这正是两波模型把中间场并入平均状态的直接表现。

% [原PDF p359]
近真空测试里 HLL 的单平均状态具有稳健优势：积分平均与保守包络避免某些线性化负状态。强激波测试中，外波速度和总体平台正确，接触扩散在巨大压力尺度下不突出。这说明 HLL 可作为复杂系统的安全基线，即使最终仍需 HLLEM 或 HLLC 恢复中间分辨。

图10.11、10.12分别给出 HLL 对近真空测试2和强激波测试3。HLL 在低密度区保持稳健，在强激波位置也正确；其主要差异不是外波速度，而是接触平台附近的额外扩散。

% [原PDF p360]
测试4同时包含声学波和中间结构，HLL 对最强外波较好，对慢中间波较差；测试5进一步放大差异。若目标主要是激波位置和整体守恒，HLL 可能足够；若目标涉及混合层厚度、涡量或界面传质，其附加耗散会改变物理解读。

图10.13、10.14为 HLL 对测试4、5。强声学波仍得到合理分辨，但测试5中的静止接触明显抹宽，其程度与传统 FVS 相近。因此仅凭激波算例判断黎曼求解器会掩盖中间场缺失问题。

% [原PDF p361]
Rusanov 的单速度耗散无需区分左右包络，实施极简且易稳定。但跨音速稀疏波内部不同特征速度变号，统一 $S^+$ 不能利用这种方向变化，因而比 HLL 圆钝。近真空测试仍能运行，反映宽保守耗散对正性的帮助；这种稳健性必须与解析度代价一起评价。

图10.15、10.16给出 Rusanov 对测试1、2的结果。该通量用单一最大速度对所有特征场施加相同耗散，结构简单且近真空稳健，但测试1音速稀疏波的分辨比 HLL 明显差，接触也受到强烈扩散。

% [原PDF p362]
在强激波上，Rusanov 的平台与波位置可接受，说明最大信号速度覆盖传播范围；过渡宽度则比特征选择性求解器大。多维计算中，这种统一耗散还会削弱切向速度与涡量。高阶重构能降低光滑区误差，却在限制器启动的间断附近重新显露基础通量的耗散特性。

图10.17、10.18为 Rusanov 对测试3、4。格式能承受极强压力比并定位主要间断，代价是所有波都比 HLLC 更宽。Rusanov 的稳定性来自宽包络速度和标量耗散矩阵，并不意味着它能区分不同特征场。

% [原PDF p363]
图10.20上下两行分别对应静止和慢移接触，左右列是 HLL 和 HLLC。静止时 HLLC 符号与精确阶跃重合，HLL 已形成宽斜坡；慢移时两者都有一阶对流扩散，但 HLLC 显著较小。该对照排除了声学波和压力比因素，是验证接触实现最简洁的回归算例。

图10.19为 Rusanov 对测试5，静止接触被明显抹平。图10.20把 HLL（左）与 HLLC（右）用于孤立接触测试6、7：$t=2$ 时 HLL 的静止与慢移密度跃变都形成宽过渡带；HLLC 对静止跃变保持初始一个界面，对慢移跃变也显著更陡。该图是恢复接触波价值的最直接证据。

% [原PDF p364]
\chapter{Roe 黎曼求解器}
Roe 方法在 CFD 中长期流行，还因为其通量可自然写成中心项减去特征耗散项，便于与 MUSCL、TVD、ENO/WENO 等高阶重构结合。每个重构界面只需左右状态，计算 Roe 平均后投影波强即可。与 HLLC 相比，它显式使用完整 Euler 特征基；与 Osher 相比，它不沿非线性积分曲线积分；与精确求解器相比，它无压力迭代。

这些优点伴随两类结构性风险。第一，常系数线性化把稀疏波也表示为间断，跨音速时会选择违反熵的扩张激波；第二，Roe 平均中间状态不保证密度和内能正，双强稀疏可产生真空异常。可靠实现必须把熵修正和低密度策略作为明确算法组成，并在输出中暴露触发，而不能依靠事后裁剪物理量。

Roe 于1981年公开发表的近似黎曼求解器，是最著名、应用最广的方法之一。Roe--Pike 改进不必显式形成 Roe 平均 Jacobian，因此对新型复杂双曲守恒律或复杂介质更易推广；Glaister 由此把 Roe 方法扩展到一般状态方程 Euler 系统。

原始 Roe 格式会把跨音速稀疏波计算成“稀疏激波”，违反熵条件；Harten--Hyman、Roe--Pike、Dubois--Mehlman 等提出了熵修正。Einfeldt 等又指出真空附近的低密度异常，并说明它实际影响所有线性化黎曼求解器。Roe 方法仍被广泛用于 MHD、固体爆轰、自适应网格、二维前沿追踪、可压与不可压 Navier--Stokes、多相流和并行气体动力学。以下从线性化守恒律的构造原则出发，分别推导原始 Roe 与 Roe--Pike 形式，并讨论熵修正和数值表现。

% [原PDF p365]
Roe 求解器的核心吸引力在于“按波族线性化”：它不像 HLL 只保留包络，也不像精确求解器迭代求非线性星形压力，而是构造一个特殊常矩阵，使有限左右跳跃在通量层面仍精确守恒。由此得到的每个特征波都有独立速度和强度，接触与剪切不会被最大声速统一耗散。代价是平均状态必须精心选择，并且线性波系不能天然表示展开稀疏扇。

Roe 方法还曾通过相分裂用于推进系统中的反应多相流，Sainsaulieu 则把它推广到无需相分裂的一类多相模型。本章的具体对象是沿某界面法向分裂的三维非定常 Euler 方程；数值方法只需这样的方向子系统。原始 Roe 路线和 Roe--Pike 路线都先在较简单的等温气体方程上演示，再用于完整 Euler 系统。

\section{Roe 方法的基础}
\subsection{精确黎曼问题与 Godunov 通量}
仍从初边值问题
\begin{equation}
 \bm U_t+\bm F(\bm U)_x=0,\quad \bm U(x,0)=\bm U^{(0)}(x),\quad
 \bm U(0,t)=\bm U_l(t),\quad\bm U(L,t)=\bm U_r(t) \tag{11.1}
\end{equation}
和显式守恒更新
\begin{equation}
 \bm U_i^{n+1}=\bm U_i^n+\frac{\Delta t}{\Delta x}
 (\bm F_{i-1/2}-\bm F_{i+1/2}) \tag{11.2}
\end{equation}
出发。精确 Godunov 通量是局部黎曼解在 $S=0$ 的物理通量。

% [原PDF p366]
沿 $x$ 分裂的三维 Euler 方程仍有
\begin{equation}
 \bm U=(\rho,\rho u,\rho v,\rho w,E)^T,qquad
 \bm F=(\rho u,\rho u^2+p,\rho uv,\rho uw,u(E+p))^T. \tag{11.5}
\end{equation}
星形区包含未知状态，而 $t$ 轴穿过的具体状态决定界面通量。第4章用精确非线性波关系求解，第9章近似星形状态再代入 $\bm F$；Roe 方法则把非线性方程组替换成保留左右数据的常系数线性系统，并直接构造通量贡献。

\subsection{近似守恒律}
定义精确 Jacobian
\begin{equation}\bm A(\bm U)=\partial\bm F/\partial\bm U,\tag{11.7}\end{equation}
光滑区有 $\bm U_t+\bm A(\bm U)\bm U_x=0$。Roe 用依赖左右态的常矩阵
\begin{equation}\widetilde{\bm A}=\widetilde{\bm A}(\bm U_L,\bm U_R)\tag{11.9}\end{equation}
替换它，从而把原非线性黎曼问题化为常系数问题
\begin{equation}\bm U_t+\widetilde{\bm A}\bm U_x=0\tag{11.10}\end{equation}
而初始左右状态保持不变。

% [原PDF p367]
一般 $m$ 方程双曲系统的 Roe 矩阵需满足三条性质：
\begin{enumerate}
\item \textbf{双曲性：}具有有序实特征值 $\widetilde\lambda_1\le\cdots\le\widetilde\lambda_m$ 和完备线性无关右特征矢量 $\widetilde{\bm K}^{(i)}$；
\item \textbf{与精确 Jacobian 一致：}$\widetilde{\bm A}(\bm U,\bm U)=\bm A(\bm U)$；
\item \textbf{跨间断守恒：}
\begin{equation}
 \bm F(\bm U_R)-\bm F(\bm U_L)=\widetilde{\bm A}(\bm U_R-\bm U_L). \tag{11.15}
\end{equation}
\end{enumerate}
第一条保持方程的数学类型，第二条保证左右态趋同时恢复原方程，第三条是离散 Rankine--Hugoniot 一致性。若左右态恰由一支孤立间断相连，第三条使近似黎曼求解器识别该波；但这不表示有限网格的一阶 Roe--Godunov 更新一定对任意孤立间断给出全局精确解。

% [原PDF p368]
性质(C)还可理解为：在有限状态跳跃上，$\widetilde{\bm A}$ 充当通量函数的精确割线，而性质(B)要求左右态合并时割线趋向切线 Jacobian。这与标量函数的平均导数类似，但系统还要同时保持实谱和完备特征矢量。任意算术平均 Jacobian 往往只能满足一致性，不能满足精确通量差，因此会错认孤立间断或破坏守恒。

对一般系统同时构造满足三条性质的矩阵可能极其复杂。Euler 方程存在较简洁的参数化，而 Roe--Pike 进一步避免显式形成矩阵。

\subsection{近似黎曼解与界面通量}
把数据跳跃投影到 Roe 右特征矢量：
\begin{equation}
 \Delta\bm U=\bm U_R-\bm U_L=
 \sum_{i=1}^m\widetilde\alpha_i\widetilde{\bm K}^{(i)}, \tag{11.16}
\end{equation}
其中 $\widetilde\alpha_i$ 是第 $i$ 族波强。线性黎曼解在 $S=0$ 的状态可从任一侧累加：
\begin{equation}
 \bm U(0)=\bm U_L+\sum_{\widetilde\lambda_i\le0}
 \widetilde\alpha_i\widetilde{\bm K}^{(i)}
 =\bm U_R-\sum_{\widetilde\lambda_i\ge0}
 \widetilde\alpha_i\widetilde{\bm K}^{(i)}. \tag{11.17--11.18}
\end{equation}
负速波从右向左越过 $t$ 轴，正速波则留在右侧，这正是特征迎风分解。

% [原PDF p369]
线性近似系统的形式通量 $\widetilde{\bm A}\bm U$ 只确定到一个常矢量：给它加任意常量并不改变 PDE 的空间导数。正因为这个“通量零点”不唯一，直接代入 $\bm U(0)$ 会丢失与原物理通量的基准匹配。左右积分公式用 $\bm F_L$ 或 $\bm F_R$ 固定该常量，最终才得到全向右时精确等于左通量、全向左时精确等于右通量的迎风公式。

不能简单取 $\widetilde{\bm A}\bm U(0)$ 作为界面通量；例如全向右超声速时，它未必等于真实左通量 $\bm F_L$。正确做法是把线性近似解代入第10章左右半控制体积分关系，并以真实端点通量 $\bm F_L,\bm F_R$ 固定基准。得到
\begin{equation}
 \bm F_{i+1/2}=\bm F_L+
 \sum_{\widetilde\lambda_i\le0}\widetilde\alpha_i\widetilde\lambda_i
 \widetilde{\bm K}^{(i)}, \tag{11.27}
\end{equation}
或
\begin{equation}
 \bm F_{i+1/2}=\bm F_R-
 \sum_{\widetilde\lambda_i\ge0}\widetilde\alpha_i\widetilde\lambda_i
 \widetilde{\bm K}^{(i)}. \tag{11.28}
\end{equation}
这两式分别从左、右物理通量出发，只加上真正穿越界面的特征波通量跳跃。

% [原PDF p370]
三种 Roe 通量写法在精确算术中等价。左、右单侧形式在特征速度有明确符号时直观；中心形式适合统一实现并显示耗散矩阵 $|\widetilde{\bm A}|$。零特征值对应相对网格静止的波，其耗散项为零，这解释了 Roe 对孤立静止接触可达到极高分辨率，也预示跨音速稀疏波处绝对值函数需要熵修正。

对称形式为
\begin{equation}
 \bm F_{i+1/2}=\frac12(\bm F_L+\bm F_R)
 -\frac12\sum_{i=1}^m\widetilde\alpha_i|\widetilde\lambda_i|
 \widetilde{\bm K}^{(i)}. \tag{11.29}
\end{equation}
它清楚显示 Roe 通量由中心平均加特征选择性耗散构成：每一波族的耗散系数是自身速度绝对值，而非 Rusanov 那样统一使用最大速度。上述公式适用于任意双曲系统和任意线性化；实际计算只需 $\widetilde\lambda_i$、$\widetilde{\bm K}^{(i)}$、$\widetilde\alpha_i$，不一定要显式存储 $\widetilde{\bm A}$。

\section{原始 Roe 方法}
原始方法先构造满足三条性质的 Roe 矩阵，再求其特征结构。关键是引入参数矢量 $\bm Q$，使守恒量和通量都成为 $\bm Q$ 的二次函数：
\begin{equation}\bm U=\bm U(\bm Q),\qquad\bm F=\bm F(\bm Q).\tag{11.30}\end{equation}
如此可把 $\Delta\bm U$、$\Delta\bm F$ 线性表示为 $\Delta\bm Q$，系数只含 $\bm Q_L,\bm Q_R$ 的算术平均。这个参数化避免了平均值定理构造的非唯一性与复杂性。

% [原PDF p371]
参数矢量使用 $\sqrt\rho$ 而不是 $\rho$，是因为两端二次差满足
$q_R^2-q_L^2=(q_R+q_L)(q_R-q_L)$，可用算术平均精确因式分解。密度平方根把动量和能量中的密度乘积都转成 $\bm Q$ 的二次项，于是有限跳跃精确线性化。该技巧并非对任意状态方程都成立，这也是 Roe--Pike 直接平均标量方法在复杂介质中更有吸引力的原因。

\subsection{等温气体方程}
等温系统为
\begin{equation}
 \bm U=\begin{bmatrix}\rho\\\rho u\end{bmatrix},\qquad
 \bm F=\begin{bmatrix}\rho u\\\rho u^2+a^2\rho\end{bmatrix}, \tag{11.32}
\end{equation}
其中 $a$ 为常声速。精确特征结构是
\begin{equation}
 \bm A=\begin{bmatrix}0&1\\a^2-u^2&2u\end{bmatrix},\quad
 \lambda_{1,2}=u\mp a,\quad
 \bm K^{(1,2)}=\begin{bmatrix}1\\u\mp a\end{bmatrix}. \tag{11.33}
\end{equation}
选取
\begin{equation}
 \bm Q=\sqrt\rho\begin{bmatrix}1\\u\end{bmatrix}
 =\begin{bmatrix}q_1\\q_2\end{bmatrix}, \tag{11.34}
\end{equation}
则
\begin{equation}
 \bm U=\begin{bmatrix}q_1^2\\q_1q_2\end{bmatrix},\qquad
 \bm F=\begin{bmatrix}q_1q_2\\q_2^2+a^2q_1^2\end{bmatrix}. \tag{11.35--11.36}
\end{equation}
取 $\widetilde{\bm Q}=(\bm Q_L+\bm Q_R)/2$，寻找 $\Delta\bm U=\widetilde{\bm B}\Delta\bm Q$、$\Delta\bm F=\widetilde{\bm C}\Delta\bm Q$。

% [原PDF p372]
$\widetilde{\bm B}$ 表示守恒量对参数跳跃的离散导数，$\widetilde{\bm C}$ 表示通量对参数跳跃的离散导数。由 $\Delta\bm F=\widetilde{\bm C}\Delta\bm Q$ 和 $\Delta\bm U=\widetilde{\bm B}\Delta\bm Q$ 消去 $\Delta\bm Q$，自然得到 $\Delta\bm F=\widetilde{\bm C}\widetilde{\bm B}^{-1}\Delta\bm U$。这一步直接保证性质(C)，而矩阵形式与精确 Jacobian 相同又使双曲性检查十分简单。

由二次差的恒等式可得
\begin{equation}
 \widetilde{\bm B}=\begin{bmatrix}2\widetilde q_1&0\\\widetilde q_2&\widetilde q_1\end{bmatrix},\qquad
 \widetilde{\bm C}=\begin{bmatrix}\widetilde q_2&\widetilde q_1\\2a^2\widetilde q_1&2\widetilde q_2\end{bmatrix}, \tag{11.41}
\end{equation}
故
\begin{equation}
 \widetilde{\bm A}=\widetilde{\bm C}\widetilde{\bm B}^{-1}
 =\begin{bmatrix}0&1\\a^2-\widetilde u^2&2\widetilde u\end{bmatrix},\qquad
 \widetilde u=\frac{\sqrt{\rho_L}u_L+\sqrt{\rho_R}u_R}{\sqrt{\rho_L}+\sqrt{\rho_R}}. \tag{11.42--11.43}
\end{equation}
这与精确 Jacobian 形式完全相同，只把 $u$ 换为密度平方根加权平均；因 $a$ 为常数，不需平均声速或密度。

平均特征值、特征矢量为 $\widetilde u\mp a$ 和 $(1,\widetilde u\mp a)^T$。波强为
\begin{equation}
 \widetilde\alpha_1=\frac{\Delta\rho(\widetilde u+a)-\Delta(\rho u)}{2a},\qquad
 \widetilde\alpha_2=\frac{-\Delta\rho(\widetilde u-a)+\Delta(\rho u)}{2a}. \tag{11.46}
\end{equation}

% [原PDF p373]
Euler 的五个特征矢量选择并非唯一，任何非零尺度倍数都表示同一方向；波强会反向缩放，使乘积 $\alpha_i\bm K^{(i)}$ 不变。这里的归一化让密度分量在声学和熵波中等于1，便于由质量跳跃求波强。两支剪切矢量的密度与法向动量分量为零，明确表示剪切只改变一项切向动量及相应动能输运。

\subsection{Euler 方程}
沿 $x$ 的三维 Euler 系统有五个特征值
\begin{equation}\lambda_1=u-a,\quad\lambda_2=\lambda_3=\lambda_4=u,\quad\lambda_5=u+a.\tag{11.48}\end{equation}
右特征矢量可取
\begin{align}
 \bm K^{(1)}&=(1,u-a,v,w,H-ua)^T,\\
 \bm K^{(2)}&=(1,u,v,w,V^2/2)^T,\\
 \bm K^{(3)}&=(0,0,1,0,v)^T,\quad
 \bm K^{(4)}=(0,0,0,1,w)^T,\\
 \bm K^{(5)}&=(1,u+a,v,w,H+ua)^T. \tag{11.49}
\end{align}
其中 $H=(E+p)/\rho$、$E=\rho V^2/2+\rho e$、$V^2=u^2+v^2+w^2$，理想气体 $e=p/[(\gamma-1)\rho]$。第一、第五矢量是左右声学波，第二是熵/接触波，第三、第四是两支剪切波。

% [原PDF p374]
Roe 参数矢量选为
\begin{equation}
 \bm Q=\sqrt\rho\,(1,u,v,w,H)^T. \tag{11.54}
\end{equation}
这样 $\bm U$、$x$ 向通量乃至完整三维的其他方向通量，每个分量都是 $q_i$ 的二次式。按等温例同样构造 $\widetilde{\bm B}$、$\widetilde{\bm C}$，便得到 Roe 矩阵及其特征结构。完整矩阵表达虽然较长，其物理含义很简单：精确 Euler Jacobian 的函数形式保持不变，只把速度、总焓与声速替换成 Roe 平均。

% [原PDF p375]
平均总焓而不是直接平均压力或声速，是满足能量通量割线关系的关键。先平均 $H$ 和三个速度，再由理想气体热力学关系恢复 $\widetilde a$，可以保证平均特征结构内部一致。若直接算术平均 $a_L,a_R$，一般会破坏通量跳跃性质。根号内 $\widetilde H-\widetilde V^2/2$ 代表平均比焓的热部分，对合法理想气体数据应为正。

平均量为
\begin{align}
 \widetilde u&=\frac{\sqrt{\rho_L}u_L+\sqrt{\rho_R}u_R}{\sqrt{\rho_L}+\sqrt{\rho_R}},&
 \widetilde v&=\frac{\sqrt{\rho_L}v_L+\sqrt{\rho_R}v_R}{\sqrt{\rho_L}+\sqrt{\rho_R}},\\
 \widetilde w&=\frac{\sqrt{\rho_L}w_L+\sqrt{\rho_R}w_R}{\sqrt{\rho_L}+\sqrt{\rho_R}},&
 \widetilde H&=\frac{\sqrt{\rho_L}H_L+\sqrt{\rho_R}H_R}{\sqrt{\rho_L}+\sqrt{\rho_R}},\\
 \widetilde a&=\sqrt{(\gamma-1)\left(\widetilde H-\tfrac12\widetilde V^2\right)}.
 \tag{11.60}
\end{align}
其中 $\widetilde V^2=\widetilde u^2+\widetilde v^2+\widetilde w^2$。平均特征值是 $\widetilde u-\widetilde a$、三重 $\widetilde u$、$\widetilde u+\widetilde a$；平均特征矢量由式(11.49)逐量加波浪号得到。平方根加权并非经验插值，而是保证通量跳跃等于平均 Jacobian 乘状态跳跃的代数结果。

% [原PDF p376]
式(11.69)先分离切向动能贡献，再求熵波与两支声学波强，避免直接求解五阶线性方程。程序中 $\Delta u_j$ 指守恒分量跳跃，不应与速度跳跃混淆；$\Delta\widehat u_5$ 只是扣除剪切能量后的辅助量。若改用 Roe--Pike 的原始变量波强式，则同一物理分解可写得更紧凑，两套公式在一致的 Roe 平均上等价。

把守恒跳跃 $\Delta u_j=(u_j)_R-(u_j)_L$ 投影到五个平均特征矢量。剪切波强先直接得到
\begin{equation}
 \widetilde\alpha_3=\Delta u_3-\widetilde v\Delta u_1,qquad
 \widetilde\alpha_4=\Delta u_4-\widetilde w\Delta u_1. \tag{11.68}
\end{equation}
把它们从能量跳跃中扣除，记
\begin{equation}
 \Delta\widehat u_5=\Delta u_5-widetilde v(\Delta u_3-\widetilde v\Delta u_1)
 -\widetilde w(\Delta u_4-\widetilde w\Delta u_1), \tag{11.70}
\end{equation}
则
\begin{align}
 \widetilde\alpha_2&=\frac{\gamma-1}{\widetilde a^2}
 [\Delta u_1(\widetilde H-\widetilde u^2)+\widetilde u\Delta u_2-\Delta\widehat u_5],\\
 \widetilde\alpha_1&=\frac{\Delta u_1(\widetilde u+\widetilde a)-\Delta u_2-\widetilde a\widetilde\alpha_2}{2\widetilde a},\\
 \widetilde\alpha_5&=\Delta u_1-(\widetilde\alpha_1+\widetilde\alpha_2). \tag{11.69}
\end{align}
算法依次计算平均量、特征值、特征矢量、波强，再代入式(11.27)--(11.29)。纯一维时 $v=w=0$，剪切波强消失。实用计算还必须加入第11.4节的熵修正。

% [原PDF p377]
Roe--Pike 的“先局部、后全局”思路很实用：局部参考态线性化只负责发现波强应由哪些原始变量跳跃组合；随后不再把参考态固定为任意平均，而是用有限跳跃的状态与通量重构方程反求真正平均。这样既利用解析特征结构，又把守恒约束直接施加在最终用于通量的量上。

\section{Roe--Pike 方法}
原始路线先求 $\widetilde{\bm A}$；Roe--Pike 直接寻找一组标量平均量，然后把解析特征值、特征矢量和波强公式在平均状态上求值。它假定原系统的特征结构已知，并分两步进行。

\subsection{方法思想}
第一步，在参考态 $\widehat{\bm U}$ 附近作局部线性化，假定左右态与参考态的差为小量，并忽略二阶项，得到
\begin{equation}\bm U_t+\widehat{\bm A}\bm U_x=0.\tag{11.71}\end{equation}
在此常系数系统中分解 $\Delta\bm U$，求得波强的简单解析表达 $\widehat\alpha_i$。这一步不是 Roe 的最终线性化，只是为了获得可在尚未知的平均状态上评估的公式。

% [原PDF p378]
两个代数条件的角色不同。状态重构保证各波增量之和连接 $\bm U_L$ 与 $\bm U_R$；通量重构保证每个波增量再乘自身速度后连接 $\bm F_L$ 与 $\bm F_R$。只满足前者的普通线性化可以得到某种波分解，却不一定跨有限间断守恒。Roe--Pike 平均必须让两者同时成立，这通常形成超定但含恒等式的非线性代数系统。

第二步选择平均原始变量矢量 $\widetilde{\bm W}$，令
\begin{equation}
 \widetilde\alpha_i=\widehat\alpha_i(\widetilde{\bm W}),\quad
 \widetilde\lambda_i=\lambda_i(\widetilde{\bm W}),\quad
 \widetilde{\bm K}^{(i)}=\bm K^{(i)}(\widetilde{\bm W}), \tag{11.74}
\end{equation}
并要求它同时精确满足
\begin{equation}
 \Delta\bm U=\sum_i\widetilde\alpha_i\widetilde{\bm K}^{(i)},\qquad
 \Delta\bm F=\sum_i\widetilde\alpha_i\widetilde\lambda_i\widetilde{\bm K}^{(i)}. \tag{11.75--11.76}
\end{equation}
第一式重构状态跳跃，第二式重构通量跳跃，后者就是 Roe 守恒性质。解这个代数问题即可确定平均状态，而无需写出整个平均 Jacobian。

\subsection{等温方程示例}
对等温方程，在参考原始状态 $(\widehat\rho,\widehat u)$ 上，特征值与右特征矢量是 $\widehat u\mp a$ 与 $(1,\widehat u\mp a)^T$。

% [原PDF p379]
乘积跳跃公式中的 $O(\Delta^2)$ 项明确写出了左右态偏离参考态的乘积。因此参考态附近波强式是一阶准确的；它只用于导出解析形式，最终平均态将通过有限跳跃条件消除这种局部限制。等温例的两个波强分别把密度跳跃与阻抗型速度跳跃相加、相减，正好对应左右传播的声学信息。

将数据跳跃分解为
\begin{equation}
 \Delta\rho=\widehat\alpha_1+\widehat\alpha_2,qquad
 \Delta(\rho u)=\widehat\alpha_1(\widehat u-a)+\widehat\alpha_2(\widehat u+a). \tag{11.82--11.83}
\end{equation}
用乘积跳跃线性化
\begin{equation}\Delta(\rho u)=\widehat\rho\,\Delta u+\widehat u\,\Delta\rho+O(\Delta^2),\tag{11.84}\end{equation}
忽略二阶项后解得
\begin{equation}
 \widehat\alpha_1=\tfrac12(\Delta\rho-\widehat\rho\Delta u/a),\qquad
 \widehat\alpha_2=\tfrac12(\Delta\rho+\widehat\rho\Delta u/a). \tag{11.86}
\end{equation}
这些公式在参考态附近同时重构状态和通量跳跃到二阶小量精度。

% [原PDF p380]
验证通量重构时，质量通量方程与动量状态方程重合，是因为第一物理通量就是第二守恒分量。动量通量还包含 $\rho u^2+a^2\rho$，给出真正新增的约束。把波强表达代入后，所有特征分解细节都化成密度、动量和二阶动量跳跃的平均恒等式。

把帽号参考量换成待求平均量，要求两组重构关系对有限左右跳跃精确成立。令
\begin{equation}
 \widetilde\alpha_{1,2}=\tfrac12(\Delta\rho\mp\widetilde\rho\Delta u/a),\quad
 \widetilde\lambda_{1,2}=\widetilde u\mp a,\quad
 \widetilde{\bm K}^{(1,2)}=(1,\widetilde u\mp a)^T. \tag{11.91--11.93}
\end{equation}
质量跳跃方程自动成立，动量状态跳跃与质量通量跳跃实际上是同一方程；剩余条件可以化为两个关于 $\widetilde\rho,\widetilde u$ 的代数关系。

% [原PDF p381]
二次方程的两个根说明，仅凭有限跳跃代数可能存在表面上都满足重构的分支，必须再用一致性和连续极限筛选。物理解在左右密度相等时平滑趋于速度平均；伪根的分母则趋零。数值实现不应直接用带 $\Delta\rho$ 分母的二次根公式，而应使用已经化简的平方根密度加权式，避免接近等密度时灾难性消去。

这两个关系是
\begin{equation}
 \Delta(\rho u)=\widetilde\rho\Delta u+\widetilde u\Delta\rho,qquad
 \Delta(\rho u^2)=2\widetilde u\widetilde\rho\Delta u+\widetilde u^2\Delta\rho. \tag{11.98--11.99}
\end{equation}
消去 $\widetilde\rho$ 得
\begin{equation}
 \Delta\rho\,\widetilde u^2-2\Delta(\rho u)\widetilde u+\Delta(\rho u^2)=0. \tag{11.100}
\end{equation}
判别式化为 $\rho_L\rho_R(\Delta u)^2$，故有两个根。只有取适当符号的根在 $\rho_L\to\rho_R$ 时保持良性。

% [原PDF p382]
$\widetilde\rho=\sqrt{\rho_L\rho_R}$ 是几何平均，而速度使用平方根密度加权平均。前者进入 Roe--Pike 原始变量波强，后者进入特征速度。它们各自承担不同代数作用，不能互换为统一算术平均。把结果代回式(11.98)--(11.99)可以直接验证两条有限跳跃恒等式，从而确认与原始参数法的一致性。

物理解为
\begin{equation}
 \widetilde u=\frac{\sqrt{\rho_L}u_L+\sqrt{\rho_R}u_R}{\sqrt{\rho_L}+\sqrt{\rho_R}},qquad
 \widetilde\rho=\sqrt{\rho_L\rho_R}. \tag{11.103--11.104}
\end{equation}
另一个代数根
\begin{equation}
 \widetilde u=\frac{\sqrt{\rho_R}u_R-\sqrt{\rho_L}u_L}{\sqrt{\rho_R}-\sqrt{\rho_L}} \tag{11.105}
\end{equation}
在 $\rho_L=\rho_R$、$u_L\ne u_R$ 时甚至无定义，故必须舍弃。所得 $\widetilde u$ 与原始 Roe 方法完全相同，说明平方根加权同时由参数化和守恒代数两条路线唯一选出。

\subsection{Euler 方程}
对三维 Euler 方程，先在参考态计算式(11.48)--(11.49)的特征结构，并把 $\Delta\bm U$ 投影到五个右特征矢量。

% [原PDF p383]
原始变量波强式还显示压力和法向速度跳跃如何分配到两支声学波：$\Delta p$ 对左右波同号，$\rho a\Delta u$ 对两者反号。纯接触满足 $\Delta p=\Delta u=0$，所以声学波强为零，只剩 $\alpha_2=\Delta\rho$；纯切向剪切则只激活 $\alpha_3$ 或 $\alpha_4$。这些极限关系是验证 Roe 投影代码最有效的单元测试。

对于一般面法向，应把 $\Delta u$ 理解为法向速度跳跃，把 $\Delta v,\Delta w$ 理解为任意两正交切向分量。只要局部基正交，回转后的耗散与切向基具体选择无关。若在全局笛卡尔分量上直接套用沿 $x$ 公式，而面并不对齐 $x$，则声学和剪切波会错误耦合。

用乘积跳跃法则 $\Delta(rs)=\widehat r\Delta s+\widehat s\Delta r+O(\Delta^2)$，可把波强写成更有物理意义的原始变量跳跃：
\begin{equation}
\begin{aligned}
 \widehat\alpha_1&=\frac{\Delta p-\widehat\rho\widehat a\Delta u}{2\widehat a^2},&
 \widehat\alpha_2&=\Delta\rho-\frac{\Delta p}{\widehat a^2},\\
 \widehat\alpha_3&=\widehat\rho\Delta v,&
 \widehat\alpha_4&=\widehat\rho\Delta w,\\
 \widehat\alpha_5&=\frac{\Delta p+\widehat\rho\widehat a\Delta u}{2\widehat a^2}.
\end{aligned}\tag{11.113}
\end{equation}
第一、第五项分别是左右声学波强，第二项是熵/接触波强，第三、第四项是两支切向剪切波强。该表达比守恒分量差公式直观，也常用于程序实现。

% [原PDF p384]
Euler 代数问题虽有五个方程和更多热力学耦合，最终仍给出同样的根号密度权重。平均密度只出现在波强，平均总焓决定声速，三个速度分量分别按相同权重处理。这种形式也保持坐标旋转的一致性：把速度旋转到法切坐标后平均，与先平均笛卡尔速度再旋转得到同一结果。

第二步要求平均状态 $\widetilde{\bm W}=(\widetilde\rho,\widetilde u,\widetilde v,\widetilde w,\widetilde a)^T$ 同时精确重构 $\Delta\bm U$ 和 $\Delta\bm F$。代数解为
\begin{equation}
 \widetilde\rho=\sqrt{\rho_L\rho_R},\qquad
 (\widetilde u,\widetilde v,\widetilde w,\widetilde H)
 =\frac{\sqrt{\rho_L}(u_L,v_L,w_L,H_L)+\sqrt{\rho_R}(u_R,v_R,w_R,H_R)}
 {\sqrt{\rho_L}+\sqrt{\rho_R}}, \tag{11.118a}
\end{equation}
\begin{equation}
 \widetilde a=\sqrt{(\gamma-1)(\widetilde H-\widetilde V^2/2)}. \tag{11.118b}
\end{equation}
这与原始 Roe 参数法所得平均完全一致。然后在此平均状态计算式(11.113)的波强、解析特征值和特征矢量，再代入 Roe 通量。

% [原PDF p385]
熵修正只针对非线性声学族，因为接触和剪切的精确解本来就是间断，不应将它们展开。跨音速检测需要波两侧的特征速度，而 Roe 线性化只给单一平均速度；因此必须额外估计星形速度和声速。这些估计可以比最终通量粗略，但必须正确判断零是否位于扇头、扇尾之间。

另一种常见熵修正会直接平滑 $|\lambda|$ 在零附近的绝对值。本章采用 Harten--Hyman 的波拆分解释，更清楚地保持积分守恒并使用物理扇边速度。二者目标相近，参数与触发范围却不同；实现和文档应明确采用哪一种，避免把经验小量平滑误称为完整 Harten--Hyman 构造。

Roe--Pike 实现步骤是：先求式(11.118)平均；再求五个平均特征值；把解析右特征矢量在平均态求值；用式(11.113)求波强；最后选用单侧或对称 Roe 通量。它的优势是避开显式平均矩阵，尤其便于给新方程组编写求解器。

\section{熵修正}
线性常系数黎曼解只由跳跃组成。对接触和激波，这至少保持了波的间断性质；对随时间展开、内部连续的稀疏波则是根本失真。多数情况下数值黏性尚能把离散跳跃抹成可接受扇形，但当稀疏波跨越音速点，线性化会产生违反熵条件的非物理“稀疏激波”，必须显式修正。

\subsection{熵问题}
以测试1为例，精确解含左跨音速稀疏波、速度 $u_*$ 的接触间断和右激波。左特征值 $\lambda_1=u-a$ 从扇头的
$u_L-a_L<0$ 变到扇尾的 $u_*-a_{*L}>0$，扇内必有 $u=a$ 的音速点。

% [原PDF p386]
熵错误可从特征几何看得更清楚：真实稀疏波中同族特征线向外张开，状态沿扇连续变化；线性 Roe 波却把整个变化压到一条速度固定的跳线上。当扇的特征速度不跨零时，界面采样常仍落在同一侧，错误不一定显著；跨零时，压缩成单跳跃会在 $t$ 轴附近选择错误熵分支，因此形成稳定但非物理的离散激波。

未修正 Roe 解会在连续稀疏扇内形成一个数值间断。物理可接受的激波应满足特征线从两侧汇入间断，即 Lax 熵条件；这里却是特征线从数值间断向外发散，故该跳跃违反熵条件。Harten--Hyman 熵修正的基本思想，是把跨音速族的单一 Roe 跳跃拆成位于音速点两侧的两个较小跳跃，使其积分通量等价，同时让零特征速度附近获得适当耗散。Euler 方程只需检查左右两支非线性声学波；中间接触和剪切场无需此类修正。

% [原PDF p387]
拆分跳跃并不是试图精确重建整个稀疏扇，而是只修正界面所见的零速附近部分。新中间状态 $\bm U_{SL}$ 由两段波速加权，使两小跳跃的总积分通量等于原 Roe 跳跃。只要扇头和扇尾速度估计合理，修正就能让特征速度在零附近连续过渡，消除扩张激波而保留 Roe 在其他波族上的低耗散。

\subsection{Harten--Hyman 熵修正}
对左波，先估计
\begin{equation}\lambda_1^L=u_L-a_L,qquad\lambda_1^R=u_*-a_{*L}.\tag{11.119}\end{equation}
若
\begin{equation}\lambda_1^L<0<\lambda_1^R,\tag{11.120}\end{equation}
则左波是跨音速稀疏波。原 Roe 单跳跃
\begin{equation}\bm U_{*L}-\bm U_L=\widetilde\alpha_1\widetilde{\bm K}^{(1)}\tag{11.121}\end{equation}
以 $\widetilde\lambda_1$ 传播；修正把它拆为 $\bm U_{SL}-\bm U_L$ 与 $\bm U_{*L}-\bm U_{SL}$，速度分别为 $\lambda_1^L$、$\lambda_1^R$。保持积分守恒得到
\begin{equation}
 \bm U_{SL}=\frac{(\widetilde\lambda_1-\lambda_1^L)\bm U_L+(\lambda_1^R-\widetilde\lambda_1)\bm U_{*L}}
 {\lambda_1^R-\lambda_1^L}. \tag{11.123}
\end{equation}

% [原PDF p388]
修正速度 $\lambda_1^{\rm ent}$ 不是物理扇头或扇尾速度，而是左侧小跳跃对界面通量的等效系数。其分数因子表示原 Roe 跳跃中位于零速左侧的比例。随着跨音速区间缩小，它平滑回到未修正贡献；若不满足符号跨越条件，则应继续使用原 $\widetilde\lambda_1$，避免在普通激波或非跨音速稀疏波上增加耗散。

从左侧单向通量公式看，只有速度 $\lambda_1^L<0$ 的左小跳跃穿越 $t$ 轴。由上式可写成
\begin{equation}
 \bm U_{SL}-\bm U_L=\frac{\lambda_1^R-\widetilde\lambda_1}{\lambda_1^R-\lambda_1^L}
 \widetilde\alpha_1\widetilde{\bm K}^{(1)}. \tag{11.125--11.126}
\end{equation}
因此只需把原特征速度在该波贡献中替换为
\begin{equation}
 \lambda_1^{\rm ent}=\lambda_1^L
 \frac{\lambda_1^R-\widetilde\lambda_1}{\lambda_1^R-\lambda_1^L}, \tag{11.128}
\end{equation}
通量变为 $\bm F_{i+1/2}=\bm F_L+\lambda_1^{\rm ent}\widetilde\alpha_1\widetilde{\bm K}^{(1)}$。实现上无需显式构造中间音速状态。

对右跨音速稀疏波，定义 $\lambda_5^L=u_*+a_{*R}$、$\lambda_5^R=u_R+a_R$。

% [原PDF p389]
右波推导与左波镜像，但通量基准取 $\bm F_R$ 并减去向右族贡献。实现时最容易出现的错误是把左右扇边速度或符号条件照抄而未镜像。稳妥做法是分别计算 $(u_L-a_L,u_*-a_{*L})$ 与 $(u_*+a_{*R},u_R+a_R)$，只对严格跨零的那一族替换等效速度。

若 $\lambda_5^L<0<\lambda_5^R$，则类似定义
\begin{equation}
 \lambda_5^{\rm ent}=\lambda_5^R
 \frac{\widetilde\lambda_5-\lambda_5^L}{\lambda_5^R-\lambda_5^L}, \tag{11.133}
\end{equation}
并用右侧单向通量
\begin{equation}
 \bm F_{i+1/2}=\bm F_R-\lambda_5^{\rm ent}
 \widetilde\alpha_5\widetilde{\bm K}^{(5)}. \tag{11.134}
\end{equation}
若使用对称形式(11.29)，也可把同一修正转写为相应的特征耗散系数。修正只在零速度落于稀疏扇内部时启动，不应无差别增加所有波的耗散。

\subsection{星形区速度与声速估计}
第一种方案用 Roe 线性化星形状态：
\begin{equation}
 \bm U_{*L}=\bm U_L+\widetilde\alpha_1\widetilde{\bm K}^{(1)},\qquad
 \bm U_{*R}=\bm U_R-\widetilde\alpha_5\widetilde{\bm K}^{(5)}. \tag{11.135,11.138}
\end{equation}

% [原PDF p390]
各种星形近似的用途仅是判断熵修正区间，不必达到精确黎曼解的全部精度。不过若估计在近真空区产生负密度或压力，声速便无定义，修正会失效。TRRS 在双膨胀极限最可靠，TSRS 在强压缩更合适，自适应方案据压力区间选择两者。固定采用昂贵 TRRS 会增加每个界面的分数幂成本，而只用 PVRS 又牺牲极端状态稳健性。

从左右星形守恒量恢复 $\rho_{*K},u_*,p_*$，再用
$a_{*K}=\sqrt{\gamma p_*/\rho_{*K}}$。另一种方案直接采用第9章 PVRS：
\begin{align}
 p_*&=\tfrac12(p_L+p_R)+\tfrac12(u_L-u_R)\overline\rho\,\overline a,\\
 u_*&=\tfrac12(u_L+u_R)+\frac{p_L-p_R}{2\overline\rho\,\overline a},\\
 \rho_{*L}&=\rho_L+(u_L-u_*)\overline\rho/\overline a,qquad
 \rho_{*R}=\rho_R+(u_*-u_R)\overline\rho/\overline a, \tag{11.140}
\end{align}
其中 $\overline\rho=(\rho_L+\rho_R)/2$、$\overline a=(a_L+a_R)/2$；为避免负压力，建议把 $p_*$ 替换为 $\max(0,p_*)$。

TRRS 则给
\begin{equation}
 p_*=\left[\frac{a_L+a_R-\tfrac12(\gamma-1)(u_R-u_L)}{a_L/p_L^z+a_R/p_R^z}\right]^{1/z},\quad z=\frac{\gamma-1}{2\gamma}, \tag{11.142}
\end{equation}
并有
\begin{equation}
 a_{*K}=a_K(p_*/p_K)^z,\quad
 u_*=u_L+\frac{2(a_L-a_{*L})}{\gamma-1}
 =u_R+\frac{2(a_{*R}-a_R)}{\gamma-1}. \tag{11.143--11.144}
\end{equation}
PVRS 和 Roe 星形状态在低密度附近可能失败；TRRS 对双稀疏波精确且稳健，但要计算多次分数次幂。TSRS 或第9章自适应无迭代求解器通常在成本与稳健性间更合适。

% [原PDF p391]
表11.1沿用统一五组数据，使 Roe 可同前面 FVS、状态近似、HLL/HLLC 直接横向比较。测试结果必须连同“是否加熵修正”和“是否加低密度正性修正”一起报告；只写“Roe”会掩盖算法实质差别。这里基本 Roe 对测试2失败是公开的适用性边界，而不是通过隐藏钳位得到一条表面可画的曲线。

\section{数值结果与讨论}
本节把一阶 Godunov 更新与 Roe 求解器结合。五个理想气体测试仍取 $\gamma=1.4$，区域 $[0,1]$、100个单元、$C_{\rm CFL}=0.9$ 和透射边界。数据与前章相同：测试1含跨音速稀疏波；测试2含双强稀疏与近真空星形区；测试3含 Mach 数198强激波；测试4含三个向右强间断；测试5含左稀疏、慢右激波和静止接触。精确解与数值解分别由 NUMERICA 的 HE--E1RPEXACT 和 HE--E1GODFLUX 产生。

这些算例不是重复展示同一误差，而是分别触发 Roe 的不同机制：测试1揭示熵错误，测试2揭示线性化正性缺陷，测试3、4检查强波幅值与速度，测试5检查低速线性场分辨。因而任何声称“Roe 格式稳健”的实现都至少应覆盖这些类型。

% [原PDF p392]
Roe 对测试3、4的高精度来自其波强和波速能分别表示声学、接触结构，而不是统一平均整个波扇。测试5中接触不如 HLLC 锐利，说明有限网格结果还受波速近似、相邻波相互作用和具体通量形式影响；“完备特征分解”并不自动意味着所有复合接触都无限分辨。

测试2失败则必须单独对待。若线性化星形状态落到负密度或负内能区域，声速和后续通量失去物理意义。Einfeldt 修正通过更保守的波模型维持正性，可能增加耗散，却是明确的算法代价。把负值改成小正数虽然让程序继续运行，会改变守恒和波速，并隐藏求解器不适用的事实。

未经熵修正的测试1结果明显错误，稀疏扇内部形成非物理跳跃。加入 Harten--Hyman 修正后，曲线在绘图精度内几乎与精确 Godunov 重合，音速点甚至略优于书中精确求解器版本，也优于 van Leer FVS；不过第10章 HLL/HLLC 在音速点仍最准确。

测试2中 Roe 按上述基本形式会失败，这是所有线性化求解器在低密度区域可能出现的共同正性问题，必须采用 Einfeldt 等人的修正，而不能以压力或密度钳位掩盖。测试3、4结果几乎与精确 Godunov 完全相同。测试5也很接近精确解；非孤立静止接触不如 HLLC 锐利，但优于 Steger--Warming 和 van Leer FVS。

\section{扩展}
本章已用等温方程和三维分裂 Euler 方程说明原始 Roe 与 Roe--Pike。定常三维超声速 Euler 的推导见 Roe 原论文。一维实现还需第6章的时间步和边界处理；二阶 TVD 标量及非线性系统扩展见第13、14章，源项见第15章，多维方向分裂与旋转见第16章。

% [原PDF p393]
未经修正的扩张激波可能在残差和全局守恒上看似正常，因为它仍是某个弱解；错误在于不满足物理熵选择。只监视守恒残差无法发现它，必须用含跨音速稀疏的精确基准检查。四个变量在同一位置同步出现跳跃，是同一错误特征波影响全部守恒量的表现。

图11.4给出测试1未加熵修正的 Roe 结果，$x_0=0.3,t=0.2$。密度、速度、压力和比内能在左稀疏扇的音速点都出现同位置的非物理折跃。这不是网格不足造成的一般扩散，而是线性化允许了违反熵条件的扩张激波；单纯加密不会选择正确弱解。

% [原PDF p394]
熵修正后稀疏扇头、扇尾和音速点均恢复正确，同时右激波与接触几乎不受影响，说明修正具有波族选择性。测试3没有跨音速膨胀触发，Roe 保持原低耗散优势。若某实现加修正后所有间断都明显变宽，应检查是否错误地对全部特征值统一平滑。

图11.5为加入 Harten--Hyman 修正后的测试1，稀疏扇恢复连续并与精确曲线重合。图11.6为强激波测试3，$x_0=0.5,t=0.012$；Roe 的特征选择性耗散使激波、接触和稀疏波均保持较高分辨，未出现传统 FVS 的明显激波后凹陷。

% [原PDF p395]
测试4、5显示 Roe 特征分解对强波和低速接触总体有效。其静止接触不如 HLLC 的差异还受复合波环境与一阶离散影响，并非缺少接触特征。实际选型时，Roe 需要额外熵与正性处理；HLLC 不需完整特征矢量且通常更稳健，两者在高阶框架中的成本与可微性也不同。

图11.7为测试4，$x_0=0.4,t=0.035$，三个强间断位置和星形平台准确。图11.8为测试5，$x_0=0.8,t=0.012$；慢激波稳定，静止接触比 HLLC 稍宽，但明显优于 HLL 和传统 FVS。这些结果说明 Roe 方法的优势在完整特征分解，而主要风险是熵违反与低密度正性。

% [原PDF p396]
\chapter{Osher 黎曼求解器}
Osher 通量的“平滑”是指它随左右状态变化通常没有 Roe 绝对特征值那样尖锐的切换；这对隐式迭代和定常收敛很有吸引力。它通过沿相空间路径累积正、负特征贡献，自然处理特征值穿零的音速点并满足熵条件。另一方面，路径必须穿过由不同特征积分曲线连接的中间状态，因而比 Roe 或 HLLC 需要更多热力学和不变量计算。

P 次序按物理黎曼解从左波族到右波族排列路径，O 次序反向排列。两条路径端点相同，却经过不同交点，非线性系统的积分结果会不同。广义黎曼不变量用于确定沿某一特征积分曲线保持不变的组合，从而求交点；若路径上特征值变号，还必须定位音速状态，把积分的正负部分正确分开。

Osher 近似求解器是最早的一批方法之一，其基础由 Engquist--Osher（1981）和 Osher--Solomon（1982）提出，后来由 Osher--Chakravarthy 用于 Euler 方程。它在定常空气动力学 CFD 中尤其流行，吸引力包括数值通量随状态平滑、满足熵条件、良好处理音速流。对缓慢移动激波，它又表现出独特性质。

Osher 格式与第8章 FVS 密切相关，也同 Godunov 方法一样是常系数线性 CIR 迎风格式的非线性推广。它在状态相空间中沿路径积分，需要选择积分路径、交点和音速点。路径取作右特征矢量对应的积分曲线，目前主要有两种次序：物理次序（P-ordering）按物理空间真实波族排列；Osher 原始次序（O-ordering）恰为其逆序。交点和音速点通过广义黎曼不变量计算。

% [原PDF p397]
本章将先对一般非线性双曲守恒律给出可直接实现的 Osher 原理，再依次用于标量方程、等熵气体动力学和理想气体三维分裂 Euler 方程。一维应用所需的其余时间推进与边界资料见第6章。

\section{一般系统的 Osher 格式}
\subsection{数学基础}
Osher 方法通过若干特殊状态上的物理通量 $\bm F(\bm U_k)$ 近似 Godunov 通量；这些状态包括左右数据、积分曲线交点和音速点。考虑
\begin{equation}\bm U_t+\bm F(\bm U)_x=0\tag{12.1}\end{equation}
及守恒更新
\begin{equation}
 \bm U_i^{n+1}=\bm U_i^n+\frac{\Delta t}{\Delta x}
 (\bm F_{i-1/2}-\bm F_{i+1/2}). \tag{12.2}
\end{equation}
目标是构造 $\bm F_{i+1/2}$。先假定系统严格双曲，特征值满足
\begin{equation}\lambda_1(\bm U)<\lambda_2(\bm U)<\cdots<\lambda_m(\bm U),\tag{12.3}\end{equation}
并具有相应完备右特征矢量。

% [原PDF p398]
通量 Jacobian 可对角化：
\begin{equation}
 \bm A(\bm U)=\frac{\partial\bm F}{\partial\bm U}
 =\bm K(\bm U)\bm\Lambda(\bm U)\bm K(\bm U)^{-1}, \tag{12.4--12.5}
\end{equation}
其中
\begin{equation}
 \bm K=[\bm K^{(1)},\bm K^{(2)},\ldots,\bm K^{(m)}],\qquad
 \bm\Lambda=\operatorname{diag}(\lambda_1,\lambda_2,\ldots,\lambda_m). \tag{12.6}
\end{equation}
这组特征矢量定义状态空间中的积分曲线；Osher 通量沿由这些曲线拼接的路径对 $\bm A\,d\bm U=d\bm F$ 的正、负谱部分积分。后续推导将说明路径次序为何影响交点，却仍能保持一致性和熵性质。
% ===== 译文分片 5（已合并）=====

% [原PDF p399]
这里的路径积分不是守恒变量各分量的逐项普通积分，而是沿相空间中指定曲线对矩阵一形式积分。路径一旦改变，交点、声速点和最终通量组合都可能改变；因此排序属于求解器定义的一部分。等温模型把声速固定为常数，因而适合作为检查路径排序的最小系统。它清楚显示 P 排序近似真实星形区，而 O 排序只是代数积分路径；不能把后者的交点变量直接解释成物理解。声速点不是额外的物理间断，而是积分路径上某一特征速度穿过零的位置。其状态由声速条件和相应广义黎曼不变量共同确定，随后只以物理通量值进入分段积分公式。

由特征值的正、负部分定义
\begin{equation}
 \bm A^+(\bm U)=\bm K\bm\Lambda^+\bm K^{-1},\qquad
 \bm A^-(\bm U)=\bm K\bm\Lambda^-\bm K^{-1}, \tag{12.14--12.15}
\end{equation}
便有 Jacobian 分裂
\begin{equation}\bm A(\bm U)=\bm A^+(\bm U)+\bm A^-(\bm U).\tag{12.16}\end{equation}
其中 $\bm A^+$ 只有非负特征值，$\bm A^-$ 只有非正特征值。这正是第5章常系数线性系统 Jacobian 分裂向非线性系统的直接推广；亦可与第8章的通量矢量分裂比较。

\subsection{Osher 数值通量}
Osher 方法假定存在矢量函数 $\bm F^+(\bm U)$ 与 $\bm F^-(\bm U)$，满足
\begin{equation}
 \bm F=\bm F^++\bm F^-,\qquad
 \frac{\partial\bm F^+}{\partial\bm U}=\bm A^+,
 \quad\frac{\partial\bm F^-}{\partial\bm U}=\bm A^-. \tag{12.17--12.18}
\end{equation}
把界面黎曼数据记作 $\bm U_0\equiv\bm U_L=\bm U_i^n$、$\bm U_1\equiv\bm U_R=\bm U_{i+1}^n$，则
\begin{equation}\bm F_{i+1/2}=\bm F^+(\bm U_0)+\bm F^-(\bm U_1).\tag{12.20}\end{equation}
利用通量差与 Jacobian 的路径积分关系，可写成单侧形式
\begin{equation}\bm F_{i+1/2}=\bm F(\bm U_0)+\int_{\bm U_0}^{\bm U_1}\bm A^-(\bm U)\,d\bm U.\tag{12.21}\end{equation}

% [原PDF p400]
符号下标零和一始终表示界面左右数据，分数下标表示部分路径交点，字母S表示特征速度为零的声速状态。实现时应先由原始量构造这些状态，再统一转换成守恒量并计算物理通量。两种排序在温和数据上可能给出近似相同的结果，但强双稀疏或强双激波会放大压力公式中速度差符号的影响。数值算法应在构造交点后检查密度、压力和声速的可定义性，并对失效作明确诊断。数值结果应同时比较密度、速度、压力和内能，因为单个变量看似合理并不能保证状态正确。特别是测试五中压力与速度平台、激波传播方向和接触面位置共同暴露了 P 排序的异常。

另外两种完全等价的形式为
\begin{align}
\bm F_{i+1/2}&=\bm F(\bm U_1)-\int_{\bm U_0}^{\bm U_1}\bm A^+(\bm U)\,d\bm U,\tag{12.22}\\
\bm F_{i+1/2}&=\tfrac12[\bm F(\bm U_0)+\bm F(\bm U_1)]
-\tfrac12\int_{\bm U_0}^{\bm U_1}|\bm A(\bm U)|\,d\bm U.\tag{12.23}
\end{align}
积分是在相空间 $\mathbb R^m$ 内进行的，其中 $\bm U=[u_1,\ldots,u_m]^T$。一般来说，结果依赖连接两端状态的路径；Osher 方法的关键正是选择能够解析或便捷数值积分的特殊路径。

\subsection{单波情形的 Osher 通量}
一般黎曼解包含 $m$ 个波族。先设除第 $j$ 波外其余波均为平凡波，即 $\bm U_0$ 与 $\bm U_1$ 由单一波相连。对任一特征矢量场 $\bm K^{(k)}(\bm U)$，取其积分曲线 $I_k$ 为积分路径；曲线上每点的切向均与该特征矢量平行。应注意，选作路径的第 $k$ 特征场未必就是实际非平凡的第 $j$ 波族。

若以 $\bm U(\xi)$ 参数化路径，$\bm U(0)=\bm U_0$、$\bm U(\xi_1)=\bm U_1$，并令
\begin{equation}\frac{d\bm U}{d\xi}=\bm K^{(k)}[\bm U(\xi)],\tag{12.26}\end{equation}
则变量代换给出
\begin{equation}
\int_{\bm U_0}^{\bm U_1}\bm A^-d\bm U
=\int_0^{\xi_1}\lambda_k^-[\bm U(\xi)]\bm K^{(k)}[\bm U(\xi)]\,d\xi.\tag{12.27}
\end{equation}

% [原PDF p401]
正负 Jacobian 的作用是把沿路径传播的特征信息按方向分开。对每一特征场，正速度贡献进入右行部分，负速度贡献进入左行部分；变号时必须显式插入声速点，不能用端点符号粗略代替。表格中的多分支表达式本质上都来自同一条规则：负特征部分只积累向左传播的信息。使用端点与声速点的物理通量组合，可以避免对路径积分再做数值求积，因而计算成本仍较低。对欧拉方程，左、右声场真正非线性，中间接触场线性退化。星形区的速度和压力跨接触波连续，密度可以跳跃；切向速度也在接触面两侧分别保持左右数据，这些关系必须在多维分裂通量中保存。

若第 $k$ 特征场线性退化，则 $\lambda_k$ 沿 $I_k$ 为常数。当 $\lambda_k\ge0$ 时，$\lambda_k^-=0$，因此路径积分为零；当 $\lambda_k\le0$ 时，$\bm A^-d\bm U=d\bm F$，故
\begin{equation}
\int_{\bm U_0}^{\bm U_1}\bm A^-d\bm U=
\begin{cases}0,&\lambda_k\ge0,\\
\bm F(\bm U_1)-\bm F(\bm U_0),&\lambda_k<0.
\end{cases}\tag{12.28--12.29}
\end{equation}
代入式(12.21)，得到纯迎风选择
\begin{equation}
\bm F_{i+1/2}=\begin{cases}\bm F(\bm U_0),&\lambda_k\ge0,\\
\bm F(\bm U_1),&\lambda_k<0.
\end{cases}\tag{12.30}
\end{equation}
若实际单波也属于该线性退化场，即 $j=k$，它与 Godunov 通量完全相同。

对真正非线性场，$\lambda_k$ 沿积分曲线单调，因此至多变号一次。若不变号，仍用式(12.30)；若在 $\xi_S$ 处变号，令 $\bm U_S=\bm U(\xi_S)$，即声速点。当符号由正变负时，积分只保留 $[\xi_S,\xi_1]$ 段，等于 $\bm F(\bm U_1)-\bm F(\bm U_S)$；由负变正时则等于 $\bm F(\bm U_S)-\bm F(\bm U_0)$。

% [原PDF p402]
Osher 通量的核心优点是连续地使用特征结构，而不是只用单一平均 Jacobian。它通常能平滑处理跨声速稀疏波，并对缓慢激波表现出较少的奇偶振荡，但路径选择仍会影响极端状态下的鲁棒性。在编程中可把单波规则写成统一函数：输入两端特征速度和三个物理通量，先判断同号，再判断符号变化方向，最后返回相应的通量差。三波系统只需依次调用三次并累加。双稀疏波近似用等熵关系闭合交点状态。它在稀疏波占主导时很自然，但若真实波为强激波，等熵延拓不再精确；本章的极端算例正用于揭示这种近似与路径排序叠加后的后果。

综上，单波路径的负特征部分积分为
\begin{equation}
\int_{\bm U_0}^{\bm U_1}\bm A^-d\bm U=
\begin{cases}
0,&\lambda_k\ge0,\\
\bm F_1-\bm F_0,&\lambda_k\le0,\\
\bm F_1-\bm F_S,&\lambda_k(\bm U_0)\ge0\ge\lambda_k(\bm U_1),\\
\bm F_S-\bm F_0,&\lambda_k(\bm U_0)\le0\le\lambda_k(\bm U_1),
\end{cases}\tag{12.31a}
\end{equation}
其中变号两行均满足 $\lambda_k(\bm U_S)=0$。于是
\begin{equation}
\bm F_{i+1/2}=
\begin{cases}
\bm F_0,&\lambda_k\ge0,\\
\bm F_1,&\lambda_k\le0,\\
\bm F_0+\bm F_1-\bm F_S,&\lambda_k(\bm U_0)\ge0\ge\lambda_k(\bm U_1),\\
\bm F_S,&\lambda_k(\bm U_0)\le0\le\lambda_k(\bm U_1).
\end{cases}\tag{12.31}
\end{equation}
这些结论可直接用于任意标量非线性守恒律 $u_t+f(u)_x=0$。

\subsection{无黏 Burgers 方程的 Osher 通量}
取 $f(u)=u^2/2$。当 $u_0>u_1$ 时精确黎曼解是速度 $S=(u_0+u_1)/2$ 的激波；当 $u_0\le u_1$ 时为由 $u_0$ 到 $u_1$ 的稀疏扇。

% [原PDF p403]
这里的路径积分不是守恒变量各分量的逐项普通积分，而是沿相空间中指定曲线对矩阵一形式积分。路径一旦改变，交点、声速点和最终通量组合都可能改变；因此排序属于求解器定义的一部分。等温模型把声速固定为常数，因而适合作为检查路径排序的最小系统。它清楚显示 P 排序近似真实星形区，而 O 排序只是代数积分路径；不能把后者的交点变量直接解释成物理解。声速点不是额外的物理间断，而是积分路径上某一特征速度穿过零的位置。其状态由声速条件和相应广义黎曼不变量共同确定，随后只以物理通量值进入分段积分公式。

唯一特征速度为 $\lambda=u$，故式(12.31)立即给出
\begin{equation}
f^{\rm osh}_{i+1/2}=\begin{cases}
f(u_0),&u_0,u_1>0,\\
f(u_1),&u_0,u_1<0,\\
f(u_0)+f(u_1),&u_0\ge0\ge u_1,\\
0,&u_0\le0\le u_1.
\end{cases}\tag{12.36}
\end{equation}
声速点 $u_S=0$，所以 $f(u_S)=0$。除跨声速激波 $u_0\ge0\ge u_1$ 外，Osher 通量与精确 Godunov 通量相同；在该例外中，Godunov 通量按激波速度 $S$ 的符号选取 $f(u_0)$ 或 $f(u_1)$。这一差异揭示了路径积分近似与精确熵解之间的关系。

\subsection{一般多波情形}
对 $m$ 个波族，Osher 选择一组分段路径 $I_k$，使 $I_k$ 与 $\bm K^{(k)}$ 相切，相邻两段在相空间交于
\begin{equation}\bm U_{k/m}=I_k\cap I_{k+1},\qquad I=\bigcup_{k=1}^m I_k.\tag{12.37--12.38}\end{equation}
于是每一段都能套用单波结论。对 $3\times3$ 系统，四个状态 $\bm U_0,\bm U_{1/3},\bm U_{2/3},\bm U_1$ 对应物理空间黎曼解的四个常值区；两个非线性波族还可能产生声速点 $\bm U_{S0},\bm U_{S1}$。

% [原PDF p404]
符号下标零和一始终表示界面左右数据，分数下标表示部分路径交点，字母S表示特征速度为零的声速状态。实现时应先由原始量构造这些状态，再统一转换成守恒量并计算物理通量。两种排序在温和数据上可能给出近似相同的结果，但强双稀疏或强双激波会放大压力公式中速度差符号的影响。数值算法应在构造交点后检查密度、压力和声速的可定义性，并对失效作明确诊断。数值结果应同时比较密度、速度、压力和内能，因为单个变量看似合理并不能保证状态正确。特别是测试五中压力与速度平台、激波传播方向和接触面位置共同暴露了 P 排序的异常。

只要交点与声速点已知，总积分就可逐段计算。交点的求取需要黎曼问题的部分信息，传统做法是利用广义黎曼不变量；不同的路径排序正由此产生。

\section{等温方程的 Osher 通量}
等温气体方程写作
\begin{equation}
\bm U_t+\bm F(\bm U)_x=0,\quad
\bm U=\begin{bmatrix}\rho\\\rho u\end{bmatrix},\quad
\bm F=\begin{bmatrix}\rho u\\\rho u^2+a^2\rho\end{bmatrix}.\tag{12.39--12.40}
\end{equation}
声速 $a$ 为常数。其 Jacobian、特征值与右特征矢量是
\begin{equation}
\bm A=\begin{bmatrix}0&1\\a^2-u^2&2u\end{bmatrix},\quad
\lambda_{1,2}=u\mp a,\quad
\bm K^{(1,2)}=\begin{bmatrix}1\\u\mp a\end{bmatrix}.\tag{12.41--12.43}
\end{equation}

% [原PDF p405]
正负 Jacobian 的作用是把沿路径传播的特征信息按方向分开。对每一特征场，正速度贡献进入右行部分，负速度贡献进入左行部分；变号时必须显式插入声速点，不能用端点符号粗略代替。表格中的多分支表达式本质上都来自同一条规则：负特征部分只积累向左传播的信息。使用端点与声速点的物理通量组合，可以避免对路径积分再做数值求积，因而计算成本仍较低。对欧拉方程，左、右声场真正非线性，中间接触场线性退化。星形区的速度和压力跨接触波连续，密度可以跳跃；切向速度也在接触面两侧分别保持左右数据，这些关系必须在多维分裂通量中保存。

黎曼解由左右两个非线性波及中间状态 $\bm U_{1/2}$ 组成，左右积分路径分别记为 $I_1,I_2$；潜在声速点为 $\bm U_{S0},\bm U_{S1}$。

\subsection{P 排序}
P 排序遵循物理空间中真实波族的次序。沿左波与右波的黎曼不变量分别为
\begin{equation}
u_{1/2}+a\ln\rho_{1/2}=u_0+a\ln\rho_0,qquad
u_{1/2}-a\ln\rho_{1/2}=u_1-a\ln\rho_1.\tag{12.44--12.45}
\end{equation}
联立可得
\begin{align}
u_{1/2}&=\tfrac12(u_0+u_1)+\tfrac12a\ln(\rho_0/\rho_1),\tag{12.46}\\
\rho_{1/2}&=\sqrt{\rho_0\rho_1}\exp[-(u_1-u_0)/(2a)].\tag{12.47}
\end{align}
交点即 $\bm U_{1/2}=[\rho_{1/2},\rho_{1/2}u_{1/2}]^T$。

% [原PDF p406]
Osher 通量的核心优点是连续地使用特征结构，而不是只用单一平均 Jacobian。它通常能平滑处理跨声速稀疏波，并对缓慢激波表现出较少的奇偶振荡，但路径选择仍会影响极端状态下的鲁棒性。在编程中可把单波规则写成统一函数：输入两端特征速度和三个物理通量，先判断同号，再判断符号变化方向，最后返回相应的通量差。三波系统只需依次调用三次并累加。双稀疏波近似用等熵关系闭合交点状态。它在稀疏波占主导时很自然，但若真实波为强激波，等熵延拓不再精确；本章的极端算例正用于揭示这种近似与路径排序叠加后的后果。

左、右声速点分别由 $u_{S0}=a$ 与 $u_{S1}=-a$ 再配合对应黎曼不变量确定：
\begin{align}
\rho_{S0}&=\rho_0\exp[(u_0-a)/a],&u_{S0}&=a,\tag{12.49}\\
\rho_{S1}&=\rho_1\exp[-(u_1+a)/a],&u_{S1}&=-a.\tag{12.50}
\end{align}
沿 $I_1$ 积分时，若两端 $u-a$ 均非负，积分为零；均非正时为 $\bm F_{1/2}-\bm F_0$；若从正变负为 $\bm F_{1/2}-\bm F_{S0}$；若从负变正则为 $\bm F_{S0}-\bm F_0$。这四种情况完整覆盖左波的局部特征构型。

% [原PDF p407]
这里的路径积分不是守恒变量各分量的逐项普通积分，而是沿相空间中指定曲线对矩阵一形式积分。路径一旦改变，交点、声速点和最终通量组合都可能改变；因此排序属于求解器定义的一部分。等温模型把声速固定为常数，因而适合作为检查路径排序的最小系统。它清楚显示 P 排序近似真实星形区，而 O 排序只是代数积分路径；不能把后者的交点变量直接解释成物理解。声速点不是额外的物理间断，而是积分路径上某一特征速度穿过零的位置。其状态由声速条件和相应广义黎曼不变量共同确定，随后只以物理通量值进入分段积分公式。

沿 $I_2$ 的规则完全类似，只需考察 $u+a$：两端均非负时为零，均非正时为 $\bm F_1-\bm F_{1/2}$，正变负时为 $\bm F_1-\bm F_{S1}$，负变正时为 $\bm F_{S1}-\bm F_{1/2}$。

值得指出，式(12.46)--(12.47)只在左右波均为稀疏波时严格精确；一般情形是近似，实质上就是第9章欧拉方程的双稀疏波近似 TRRS。数值通量采用
\begin{equation}
\bm F_{i+1/2}=\bm F_0+
\int_{\bm U_0}^{\bm U_{1/2}}\bm A^-d\bm U+
\int_{\bm U_{1/2}}^{\bm U_1}\bm A^-d\bm U.\tag{12.51--12.52}
\end{equation}
理论上两个四分支组合产生十六种情况，但其中四种因 $a>0$ 而相互矛盾，故只有十二种可实现。

% [原PDF p408]
符号下标零和一始终表示界面左右数据，分数下标表示部分路径交点，字母S表示特征速度为零的声速状态。实现时应先由原始量构造这些状态，再统一转换成守恒量并计算物理通量。两种排序在温和数据上可能给出近似相同的结果，但强双稀疏或强双激波会放大压力公式中速度差符号的影响。数值算法应在构造交点后检查密度、压力和声速的可定义性，并对失效作明确诊断。数值结果应同时比较密度、速度、压力和内能，因为单个变量看似合理并不能保证状态正确。特别是测试五中压力与速度平台、激波传播方向和接触面位置共同暴露了 P 排序的异常。

例如要求 $u_{1/2}-a\ge0$ 且 $u_{1/2}+a\le0$ 显然不可能。其余十二种中，最终通量通常是端点、交点和声速点物理通量的线性组合。相比之下，精确 Godunov 通量以及第9--11章的近似状态求解器通常只返回单个状态的通量，因此后者在程序逻辑与计算量上往往更简单。

\subsection{O 排序}
原始 Osher 方法采用 O 排序，即把分路径次序倒置。此时从 $\bm U_0$ 到 $\bm U_{1/2}$ 的路径与 $\bm K^{(2)}$ 相切，从 $\bm U_{1/2}$ 到 $\bm U_1$ 的路径与 $\bm K^{(1)}$ 相切。交点也必须用同一倒序关系求取：
\begin{equation}
u_{1/2}-a\ln\rho_{1/2}=u_0-a\ln\rho_0,\quad
u_{1/2}+a\ln\rho_{1/2}=u_1+a\ln\rho_1.
\end{equation}

% [原PDF p409]
正负 Jacobian 的作用是把沿路径传播的特征信息按方向分开。对每一特征场，正速度贡献进入右行部分，负速度贡献进入左行部分；变号时必须显式插入声速点，不能用端点符号粗略代替。表格中的多分支表达式本质上都来自同一条规则：负特征部分只积累向左传播的信息。使用端点与声速点的物理通量组合，可以避免对路径积分再做数值求积，因而计算成本仍较低。对欧拉方程，左、右声场真正非线性，中间接触场线性退化。星形区的速度和压力跨接触波连续，密度可以跳跃；切向速度也在接触面两侧分别保持左右数据，这些关系必须在多维分裂通量中保存。

因此
\begin{align}
u_{1/2}&=\tfrac12(u_0+u_1)+\tfrac12a\ln(\rho_1/\rho_0),\tag{12.54}\\
\rho_{1/2}&=\sqrt{\rho_0\rho_1}\exp[(u_1-u_0)/(2a)].\tag{12.55}
\end{align}
声速点则为
\begin{align}
u_{S0}&=-a,&\rho_{S0}&=\rho_0\exp[-(u_0+a)/a],\tag{12.56}\\
u_{S1}&= a,&\rho_{S1}&=\rho_1\exp[(u_1-a)/a].\tag{12.57}
\end{align}
O 排序的交点不能解释为真实星形区状态。例如真实解若为双稀疏波，星形区密度应降低，而式(12.55)却可能给出升高的密度。这种非物理解读在欧拉方程中会产生更严重后果。

% [原PDF p410]
Osher 通量的核心优点是连续地使用特征结构，而不是只用单一平均 Jacobian。它通常能平滑处理跨声速稀疏波，并对缓慢激波表现出较少的奇偶振荡，但路径选择仍会影响极端状态下的鲁棒性。在编程中可把单波规则写成统一函数：输入两端特征速度和三个物理通量，先判断同号，再判断符号变化方向，最后返回相应的通量差。三波系统只需依次调用三次并累加。双稀疏波近似用等熵关系闭合交点状态。它在稀疏波占主导时很自然，但若真实波为强激波，等熵延拓不再精确；本章的极端算例正用于揭示这种近似与路径排序叠加后的后果。

沿倒序路径积分时，第一段考察 $\lambda_2=u+a$：同号时分别为零或 $\bm F_{1/2}-\bm F_0$，跨越零点时引入 $\bm F_{S0}$。第二段考察 $\lambda_1=u-a$：同号时为零或 $\bm F_1-\bm F_{1/2}$，变号时引入 $\bm F_{S1}$。把两张四分支表组合，同样只有十二种特征构型可实现。

% [原PDF p411]
这里的路径积分不是守恒变量各分量的逐项普通积分，而是沿相空间中指定曲线对矩阵一形式积分。路径一旦改变，交点、声速点和最终通量组合都可能改变；因此排序属于求解器定义的一部分。等温模型把声速固定为常数，因而适合作为检查路径排序的最小系统。它清楚显示 P 排序近似真实星形区，而 O 排序只是代数积分路径；不能把后者的交点变量直接解释成物理解。声速点不是额外的物理间断，而是积分路径上某一特征速度穿过零的位置。其状态由声速条件和相应广义黎曼不变量共同确定，随后只以物理通量值进入分段积分公式。

O 排序的十二个通量表达式仍由 $\bm F_0,\bm F_{1/2},\bm F_1,\bm F_{S0},\bm F_{S1}$ 组合得到。它与 P 排序的区别不在积分公式本身，而在路径、交点以及声速点均按相反波序确定。

\section{欧拉方程的 Osher 格式}
一维欧拉方程为
\begin{equation}
\bm U_t+\bm F_x=0,\quad
\bm U=\begin{bmatrix}\rho\\\rho u\\E\end{bmatrix},\quad
\bm F=\begin{bmatrix}\rho u\\\rho u^2+p\\u(E+p)\end{bmatrix}.\tag{12.58--12.59}
\end{equation}
守恒更新仍为
\begin{equation}\bm U_i^{n+1}=\bm U_i^n+\frac{\Delta t}{\Delta x}(\bm F_{i-1/2}-\bm F_{i+1/2}).\tag{12.60}\end{equation}
特征值 $u-a,u,u+a$ 分别对应左声波、接触波和右声波。

% [原PDF p412]
符号下标零和一始终表示界面左右数据，分数下标表示部分路径交点，字母S表示特征速度为零的声速状态。实现时应先由原始量构造这些状态，再统一转换成守恒量并计算物理通量。两种排序在温和数据上可能给出近似相同的结果，但强双稀疏或强双激波会放大压力公式中速度差符号的影响。数值算法应在构造交点后检查密度、压力和声速的可定义性，并对失效作明确诊断。数值结果应同时比较密度、速度、压力和内能，因为单个变量看似合理并不能保证状态正确。特别是测试五中压力与速度平台、激波传播方向和接触面位置共同暴露了 P 排序的异常。

相应右特征矢量可取
\begin{equation}
\bm K^{(1)}=\begin{bmatrix}1\\u-a\\H-ua\end{bmatrix},\quad
\bm K^{(2)}=\begin{bmatrix}1\\u\\u^2/2\end{bmatrix},\quad
\bm K^{(3)}=\begin{bmatrix}1\\u+a\\H+ua\end{bmatrix}.\tag{12.62}
\end{equation}
求 Osher 通量分两步：先确定交点 $\bm U_{1/3},\bm U_{2/3}$，并把它们对应于黎曼解中的 $\bm U_{*L},\bm U_{*R}$；再沿三条部分路径分别积分。P 排序下，$I_1,I_2,I_3$ 按物理波从左到右排列。

\subsection{欧拉方程的 P 排序}
P 排序利用物理上正确的广义黎曼不变量
\begin{equation}I_L=u+\frac{2a}{\gamma-1},\qquad I_R=u-\frac{2a}{\gamma-1}.\tag{12.63--12.64}\end{equation}
若两侧非线性波均为稀疏波，这些关系精确；含激波时则构成双稀疏波近似。

% [原PDF p413]
正负 Jacobian 的作用是把沿路径传播的特征信息按方向分开。对每一特征场，正速度贡献进入右行部分，负速度贡献进入左行部分；变号时必须显式插入声速点，不能用端点符号粗略代替。表格中的多分支表达式本质上都来自同一条规则：负特征部分只积累向左传播的信息。使用端点与声速点的物理通量组合，可以避免对路径积分再做数值求积，因而计算成本仍较低。对欧拉方程，左、右声场真正非线性，中间接触场线性退化。星形区的速度和压力跨接触波连续，密度可以跳跃；切向速度也在接触面两侧分别保持左右数据，这些关系必须在多维分裂通量中保存。

星形区速度与压力在接触间断两侧连续：$u_{1/3}=u_{2/3}=u_*$、$p_{1/3}=p_{2/3}=p_*$。结合等熵律，令 $z=(\gamma-1)/(2\gamma)$，得到
\begin{align}
p_*&=\left[\frac{a_0+a_1-(u_1-u_0)(\gamma-1)/2}
{a_0/p_0^z+a_1/p_1^z}\right]^{1/z},\tag{12.72}\\
u_*&=\frac{Hu_0/a_0+u_1/a_1+2(H-1)/(\gamma-1)}{H/a_0+1/a_1},
\quad H=(p_0/p_1)^z.\tag{12.73}
\end{align}
两侧星形密度为
\begin{equation}\rho_{1/3}=\rho_0(p_*/p_0)^{1/\gamma},\qquad
\rho_{2/3}=\rho_1(p_*/p_1)^{1/\gamma}.\tag{12.74}\end{equation}

% [原PDF p414]
Osher 通量的核心优点是连续地使用特征结构，而不是只用单一平均 Jacobian。它通常能平滑处理跨声速稀疏波，并对缓慢激波表现出较少的奇偶振荡，但路径选择仍会影响极端状态下的鲁棒性。在编程中可把单波规则写成统一函数：输入两端特征速度和三个物理通量，先判断同号，再判断符号变化方向，最后返回相应的通量差。三波系统只需依次调用三次并累加。双稀疏波近似用等熵关系闭合交点状态。它在稀疏波占主导时很自然，但若真实波为强激波，等熵延拓不再精确；本章的极端算例正用于揭示这种近似与路径排序叠加后的后果。

左、右声速点通过 $u-a=0$、$u+a=0$ 与对应不变量求出：
\begin{align}
u_{S0}&=\frac{(\gamma-1)u_0+2a_0}{\gamma+1},&a_{S0}&=u_{S0},\\
u_{S1}&=\frac{(\gamma-1)u_1-2a_1}{\gamma+1},&a_{S1}&=-u_{S1},\tag{12.75--12.76}
\end{align}
且 $\rho_S=\rho_{0,1}(a_S/a_{0,1})^{2/(\gamma-1)}$、$p_S=p_{0,1}(\rho_S/\rho_{0,1})^\gamma$。

总积分分解为
\begin{equation}
\bm F_{i+1/2}=\bm F_0+
\int_{\bm U_0}^{\bm U_{1/3}}\bm A^-d\bm U+
\int_{\bm U_{1/3}}^{\bm U_{2/3}}\bm A^-d\bm U+
\int_{\bm U_{2/3}}^{\bm U_1}\bm A^-d\bm U.\tag{12.77}
\end{equation}
左右场真正非线性，特征值至多变号一次；中间场线性退化，$u_*$ 沿路径恒定。

% [原PDF p415]
这里的路径积分不是守恒变量各分量的逐项普通积分，而是沿相空间中指定曲线对矩阵一形式积分。路径一旦改变，交点、声速点和最终通量组合都可能改变；因此排序属于求解器定义的一部分。等温模型把声速固定为常数，因而适合作为检查路径排序的最小系统。它清楚显示 P 排序近似真实星形区，而 O 排序只是代数积分路径；不能把后者的交点变量直接解释成物理解。声速点不是额外的物理间断，而是积分路径上某一特征速度穿过零的位置。其状态由声速条件和相应广义黎曼不变量共同确定，随后只以物理通量值进入分段积分公式。

三段积分分别按以下规则选取。左段考察 $u-a$：正正为 $0$，负负为 $\bm F_{1/3}-\bm F_0$，正负为 $\bm F_{1/3}-\bm F_{S0}$，负正为 $\bm F_{S0}-\bm F_0$。中段考察 $u_*$：$u_*\ge0$ 为 $0$，$u_*<0$ 为 $\bm F_{2/3}-\bm F_{1/3}$。右段考察 $u+a$：正正为 $0$，负负为 $\bm F_1-\bm F_{2/3}$，正负为 $\bm F_1-\bm F_{S1}$，负正为 $\bm F_{S1}-\bm F_{2/3}$。这些规则就是表12.7的全部内容。

% [原PDF p416]
符号下标零和一始终表示界面左右数据，分数下标表示部分路径交点，字母S表示特征速度为零的声速状态。实现时应先由原始量构造这些状态，再统一转换成守恒量并计算物理通量。两种排序在温和数据上可能给出近似相同的结果，但强双稀疏或强双激波会放大压力公式中速度差符号的影响。数值算法应在构造交点后检查密度、压力和声速的可定义性，并对失效作明确诊断。数值结果应同时比较密度、速度、压力和内能，因为单个变量看似合理并不能保证状态正确。特别是测试五中压力与速度平台、激波传播方向和接触面位置共同暴露了 P 排序的异常。

形式上三段共有 $4\times2\times4=32$ 个组合，但特征速度的相对次序排除了其中一半。表12.8把可实现的十六种情形整理为对端点状态与交点状态的判据。实际程序不必存储整张表，只需依次判断三个路径上的特征值符号并累加相应通量差。

\subsection{欧拉方程的 O 排序}
原始 Osher 格式把路径顺序完全反转：左侧物理波使用第3特征场，右侧物理波使用第1特征场；交点、声速点以及路径积分均须保持这一倒序，不能混合使用两种排序。

% [原PDF p417]
正负 Jacobian 的作用是把沿路径传播的特征信息按方向分开。对每一特征场，正速度贡献进入右行部分，负速度贡献进入左行部分；变号时必须显式插入声速点，不能用端点符号粗略代替。表格中的多分支表达式本质上都来自同一条规则：负特征部分只积累向左传播的信息。使用端点与声速点的物理通量组合，可以避免对路径积分再做数值求积，因而计算成本仍较低。对欧拉方程，左、右声场真正非线性，中间接触场线性退化。星形区的速度和压力跨接触波连续，密度可以跳跃；切向速度也在接触面两侧分别保持左右数据，这些关系必须在多维分裂通量中保存。

因此 $\bm U_0$ 到 $\bm U_{1/3}$ 使用右黎曼不变量，$\bm U_{2/3}$ 到 $\bm U_1$ 使用左黎曼不变量：
\begin{equation}
u_*-\frac{2a_{1/3}}{\gamma-1}=u_0-\frac{2a_0}{\gamma-1},\quad
u_*+\frac{2a_{2/3}}{\gamma-1}=u_1+\frac{2a_1}{\gamma-1}.\tag{12.78--12.79}
\end{equation}
这相当于把 $\lambda_3,\bm K^{(3)}$ 配给物理左波，而把 $\lambda_1,\bm K^{(1)}$ 配给物理右波。

% [原PDF p418]
Osher 通量的核心优点是连续地使用特征结构，而不是只用单一平均 Jacobian。它通常能平滑处理跨声速稀疏波，并对缓慢激波表现出较少的奇偶振荡，但路径选择仍会影响极端状态下的鲁棒性。在编程中可把单波规则写成统一函数：输入两端特征速度和三个物理通量，先判断同号，再判断符号变化方向，最后返回相应的通量差。三波系统只需依次调用三次并累加。双稀疏波近似用等熵关系闭合交点状态。它在稀疏波占主导时很自然，但若真实波为强激波，等熵延拓不再精确；本章的极端算例正用于揭示这种近似与路径排序叠加后的后果。

O 排序给出的压力和速度为
\begin{align}
p_*&=\left[\frac{a_0+a_1+(u_1-u_0)(\gamma-1)/2}
{a_0/p_0^z+a_1/p_1^z}\right]^{1/z},\tag{12.82}\\
u_*&=\frac{Hu_0/a_0+u_1/a_1-2(H-1)/(\gamma-1)}{H/a_0+1/a_1}.\tag{12.85}
\end{align}
密度仍由等熵律 $\rho_{1/3}=\rho_0(p_*/p_0)^{1/\gamma}$、$\rho_{2/3}=\rho_1(p_*/p_1)^{1/\gamma}$ 得到。声速点则改为
\begin{align}
u_{S0}&=\frac{(\gamma-1)u_0-2a_0}{\gamma+1},&a_{S0}&=-u_{S0},\\
u_{S1}&=\frac{(\gamma-1)u_1+2a_1}{\gamma+1},&a_{S1}&=u_{S1}.\tag{12.87--12.88}
\end{align}

% [原PDF p419]
这里的路径积分不是守恒变量各分量的逐项普通积分，而是沿相空间中指定曲线对矩阵一形式积分。路径一旦改变，交点、声速点和最终通量组合都可能改变；因此排序属于求解器定义的一部分。等温模型把声速固定为常数，因而适合作为检查路径排序的最小系统。它清楚显示 P 排序近似真实星形区，而 O 排序只是代数积分路径；不能把后者的交点变量直接解释成物理解。声速点不是额外的物理间断，而是积分路径上某一特征速度穿过零的位置。其状态由声速条件和相应广义黎曼不变量共同确定，随后只以物理通量值进入分段积分公式。

积分的三个区间仍相同，但分别沿 $I_3,I_2,I_1$ 计算：
\begin{equation}
\int_{\bm U_0}^{\bm U_1}\bm A^-d\bm U=
\int_{I_3}\bm A^-d\bm U+\int_{I_2}\bm A^-d\bm U+\int_{I_1}\bm A^-d\bm U.\tag{12.90--12.93}
\end{equation}
例如全部特征值为正时三项皆为零；而同时要求 $u_*\le0$ 与 $u_*-a_{2/3}\ge0$ 因 $a_{2/3}>0$ 不可能成立。用这种逻辑可从三十二个形式组合中筛出十六个可实现组合。

% [原PDF p420]
符号下标零和一始终表示界面左右数据，分数下标表示部分路径交点，字母S表示特征速度为零的声速状态。实现时应先由原始量构造这些状态，再统一转换成守恒量并计算物理通量。两种排序在温和数据上可能给出近似相同的结果，但强双稀疏或强双激波会放大压力公式中速度差符号的影响。数值算法应在构造交点后检查密度、压力和声速的可定义性，并对失效作明确诊断。数值结果应同时比较密度、速度、压力和内能，因为单个变量看似合理并不能保证状态正确。特别是测试五中压力与速度平台、激波传播方向和接触面位置共同暴露了 P 排序的异常。

倒序三段的积分规则为：A 段考察 $u+a$，B 段考察 $u_*$，C 段考察 $u-a$。每段仍依“同号取端点通量差、变号引入声速点通量”的单波原则。表12.9给出了所有局部规则，表12.10则按数据状态和交点状态汇总最终界面通量。这样组织判断既保持数学定义，也便于程序实现。

% [原PDF p421]
正负 Jacobian 的作用是把沿路径传播的特征信息按方向分开。对每一特征场，正速度贡献进入右行部分，负速度贡献进入左行部分；变号时必须显式插入声速点，不能用端点符号粗略代替。表格中的多分支表达式本质上都来自同一条规则：负特征部分只积累向左传播的信息。使用端点与声速点的物理通量组合，可以避免对路径积分再做数值求积，因而计算成本仍较低。对欧拉方程，左、右声场真正非线性，中间接触场线性退化。星形区的速度和压力跨接触波连续，密度可以跳跃；切向速度也在接触面两侧分别保持左右数据，这些关系必须在多维分裂通量中保存。

P 排序和 O 排序的十六种通量组合实际上适用于任意 $3\times3$ 非线性双曲系统。具体方程的性质只进入交点和声速点的计算。例如，二维浅水方程按方向分裂后可直接套用这些通量公式。

\subsection{关于路径排序的讨论}
把两种压力近似写成
\begin{align}
p_*^{(P)}&=\left[\frac{a_0+a_1-\Delta u(\gamma-1)/2}{a_0/p_0^z+a_1/p_1^z}\right]^{1/z},\tag{12.94}\\
p_*^{(O)}&=\left[\frac{a_0+a_1+\Delta u(\gamma-1)/2}{a_0/p_0^z+a_1/p_1^z}\right]^{1/z},\tag{12.95}
\end{align}
其中 $\Delta u=u_1-u_0$。当 $\Delta u=0$ 时二者相同，否则行为迥异。

% [原PDF p422]
Osher 通量的核心优点是连续地使用特征结构，而不是只用单一平均 Jacobian。它通常能平滑处理跨声速稀疏波，并对缓慢激波表现出较少的奇偶振荡，但路径选择仍会影响极端状态下的鲁棒性。在编程中可把单波规则写成统一函数：输入两端特征速度和三个物理通量，先判断同号，再判断符号变化方向，最后返回相应的通量差。三波系统只需依次调用三次并累加。双稀疏波近似用等熵关系闭合交点状态。它在稀疏波占主导时很自然，但若真实波为强激波，等熵延拓不再精确；本章的极端算例正用于揭示这种近似与路径排序叠加后的后果。

当 $\Delta u\gg0$ 时，两侧流体彼此远离，真实解包含强稀疏波，极限情形出现初生空化乃至真空。P 排序会正确反映低压、低密度趋势；O 排序却给出异常大的交点压力，更像强激波状态，因而对低密度流会严重失准甚至失败。

当 $\Delta u\ll0$ 时，两股流体相向运动并形成强激波。P 排序仍给出随压缩增强的高压；O 排序则可能使压力表达式无定义，其解析失效条件为
\begin{equation}a_0+a_1+\tfrac12\Delta u(\gamma-1)\le0.\tag{12.96}\end{equation}
一种补救思路是舍弃 O 排序，或以廉价的近似状态求解器判断波型，在局部自适应切换排序。

\subsection{方向分裂的三维情形}
向二维、三维扩展是直接的。以 $x$ 方向分裂为例，还需通量分量 $\rho uv,\rho uw$，因此守恒变量必须含 $\rho v,\rho w$，交点和声速点也要补齐这些分量。

% [原PDF p423]
这里的路径积分不是守恒变量各分量的逐项普通积分，而是沿相空间中指定曲线对矩阵一形式积分。路径一旦改变，交点、声速点和最终通量组合都可能改变；因此排序属于求解器定义的一部分。等温模型把声速固定为常数，因而适合作为检查路径排序的最小系统。它清楚显示 P 排序近似真实星形区，而 O 排序只是代数积分路径；不能把后者的交点变量直接解释成物理解。声速点不是额外的物理间断，而是积分路径上某一特征速度穿过零的位置。其状态由声速条件和相应广义黎曼不变量共同确定，随后只以物理通量值进入分段积分公式。

黎曼问题中切向速度的精确结构为
\begin{equation}
q(x,t)=\begin{cases}q_L,&x/t<u_*,\\q_R,&x/t>u_*,\end{cases}
\qquad q=v,w.\tag{12.97}
\end{equation}
因此计算左声速点和 $\bm U_{1/3}$ 时取 $v_L,w_L$，计算右声速点和 $\bm U_{2/3}$ 时取 $v_R,w_R$。Osher 格式由此能良好分辨接触波、剪切波以及 Navier--Stokes 应用中的剪切层；这也是精确求解器、近似状态求解器、HLLC 与 Roe 求解器共同具备的性质。

\section{数值结果与讨论}
以下用一阶 Godunov 更新结合 Osher 近似黎曼求解器，计算理想气体 $\gamma=1.4$ 的六个一维测试。计算域为 $0\le x\le1$，使用 $M=100$ 个单元、$C_{\rm CFL}=0.9$ 和透射边界。左右初态 $\bm W=[\rho,u,p]^T$ 为：测试1 $(1,0.75,1)/(0.125,0,0.1)$；测试2 $(1,-2,0.4)/(1,2,0.4)$；测试3 $(1,0,1000)/(1,0,0.01)$；测试4 $(5.99924,19.5975,460.894)/(5.99242,-6.19633,46.095)$；测试5 $(1,-19.59745,1000)/(1,-19.59745,0.01)$；测试6 $(1,2,0.1)/(1,-2,0.1)$。

% [原PDF p424]
符号下标零和一始终表示界面左右数据，分数下标表示部分路径交点，字母S表示特征速度为零的声速状态。实现时应先由原始量构造这些状态，再统一转换成守恒量并计算物理通量。两种排序在温和数据上可能给出近似相同的结果，但强双稀疏或强双激波会放大压力公式中速度差符号的影响。数值算法应在构造交点后检查密度、压力和声速的可定义性，并对失效作明确诊断。数值结果应同时比较密度、速度、压力和内能，因为单个变量看似合理并不能保证状态正确。特别是测试五中压力与速度平台、激波传播方向和接触面位置共同暴露了 P 排序的异常。

测试1是修正的 Sod 问题，含右行激波、右行接触波和左侧跨声速稀疏波，可检验熵条件。测试2含两个对称稀疏波，中间区接近真空，用于检验低密度鲁棒性。测试3含激波马赫数约198的强右激波、接触面和左稀疏波。测试4的三个强间断均向右运动。测试5包含左稀疏波、右行激波与静止接触间断，着重检验缓慢接触波的分辨。测试6模拟两股均匀流相撞，形成两个对称强激波。

O 排序对测试1、3、4、5给出结果，但在测试2和6失败。测试1与精确 Godunov 求解器几乎一致，声速点附近也相当。测试3、4总体准确，而且测试4中的缓慢激波没有其他许多黎曼求解器常见的伪振荡，这是 Osher 方法的显著优点。

% [原PDF p425]
正负 Jacobian 的作用是把沿路径传播的特征信息按方向分开。对每一特征场，正速度贡献进入右行部分，负速度贡献进入左行部分；变号时必须显式插入声速点，不能用端点符号粗略代替。表格中的多分支表达式本质上都来自同一条规则：负特征部分只积累向左传播的信息。使用端点与声速点的物理通量组合，可以避免对路径积分再做数值求积，因而计算成本仍较低。对欧拉方程，左、右声场真正非线性，中间接触场线性退化。星形区的速度和压力跨接触波连续，密度可以跳跃；切向速度也在接触面两侧分别保持左右数据，这些关系必须在多维分裂通量中保存。

P 排序能够运行测试1--6，对前四题的整体精度与精确求解器相当，并克服 O 排序在测试2和6中的失效；不过测试6的激波附近出现伪振荡。测试5则得到错误解：密度、速度和压力均有巨大非物理过冲，预期的右行激波甚至没有传播，初始间断右侧没有任何信号。独立程序也复现了这一现象。

这些结果说明可考虑自适应排序：通常使用 O 排序，仅在黎曼解包含双激波或双稀疏波时改用其他排序。第9章任一近似状态求解器都能低成本而可靠地识别这种波型。

\section{扩展}
本章近似黎曼求解器可直接与第6章一阶迎风 Godunov 方法结合。标量问题的二阶 TVD 扩展见第13章，非线性一维系统见第14章；带源项问题见第15章，多维推广见第16章。

% [原PDF p426]
Osher 通量的核心优点是连续地使用特征结构，而不是只用单一平均 Jacobian。它通常能平滑处理跨声速稀疏波，并对缓慢激波表现出较少的奇偶振荡，但路径选择仍会影响极端状态下的鲁棒性。在编程中可把单波规则写成统一函数：输入两端特征速度和三个物理通量，先判断同号，再判断符号变化方向，最后返回相应的通量差。三波系统只需依次调用三次并累加。双稀疏波近似用等熵关系闭合交点状态。它在稀疏波占主导时很自然，但若真实波为强激波，等熵延拓不再精确；本章的极端算例正用于揭示这种近似与路径排序叠加后的后果。

图12.5比较 O 排序对测试1在 $x_0=0.3,t=0.2$ 的数值解与精确解；密度、速度、压力和内能曲线均能正确捕捉左稀疏波、接触间断和右激波。图12.6对应测试3，取 $x_0=0.5,t=0.012$；尽管压力比高达十万量级，强激波位置及星形区状态仍得到稳定分辨。图中的离散符号代表数值解，实线代表精确解。

% [原PDF p427]
这里的路径积分不是守恒变量各分量的逐项普通积分，而是沿相空间中指定曲线对矩阵一形式积分。路径一旦改变，交点、声速点和最终通量组合都可能改变；因此排序属于求解器定义的一部分。等温模型把声速固定为常数，因而适合作为检查路径排序的最小系统。它清楚显示 P 排序近似真实星形区，而 O 排序只是代数积分路径；不能把后者的交点变量直接解释成物理解。声速点不是额外的物理间断，而是积分路径上某一特征速度穿过零的位置。其状态由声速条件和相应广义黎曼不变量共同确定，随后只以物理通量值进入分段积分公式。

图12.7给出 O 排序测试4在 $x_0=0.4,t=0.035$ 的结果，三个高速向右传播间断的位置吻合良好。图12.8为测试5在 $x_0=0.8,t=0.012$ 的结果；O 排序能维持静止接触间断，并正确给出接触两侧的压力和速度平台。四个物理量采用不同纵轴量级，比较时必须留意尺度。

% [原PDF p428]
符号下标零和一始终表示界面左右数据，分数下标表示部分路径交点，字母S表示特征速度为零的声速状态。实现时应先由原始量构造这些状态，再统一转换成守恒量并计算物理通量。两种排序在温和数据上可能给出近似相同的结果，但强双稀疏或强双激波会放大压力公式中速度差符号的影响。数值算法应在构造交点后检查密度、压力和声速的可定义性，并对失效作明确诊断。数值结果应同时比较密度、速度、压力和内能，因为单个变量看似合理并不能保证状态正确。特别是测试五中压力与速度平台、激波传播方向和接触面位置共同暴露了 P 排序的异常。

图12.9是 P 排序测试1，参数与图12.5相同，结果同样接近精确解。图12.10是双稀疏波测试2，取 $x_0=0.3,t=0.2$；P 排序能够在接近真空的中央区域保持计算并再现低压、低密度状态，这正体现式(12.94)的物理优势。

% [原PDF p429]
正负 Jacobian 的作用是把沿路径传播的特征信息按方向分开。对每一特征场，正速度贡献进入右行部分，负速度贡献进入左行部分；变号时必须显式插入声速点，不能用端点符号粗略代替。表格中的多分支表达式本质上都来自同一条规则：负特征部分只积累向左传播的信息。使用端点与声速点的物理通量组合，可以避免对路径积分再做数值求积，因而计算成本仍较低。对欧拉方程，左、右声场真正非线性，中间接触场线性退化。星形区的速度和压力跨接触波连续，密度可以跳跃；切向速度也在接触面两侧分别保持左右数据，这些关系必须在多维分裂通量中保存。

图12.11与图12.12分别给出 P 排序的测试3和测试4。强激波、接触面及稀疏波的位置与精确线基本一致，密度与内能在间断处的数值扩散约为数个网格单元。测试4同时存在极高速度与较大压比，仍未出现 O 排序的解析失效。

% [原PDF p430]
Osher 通量的核心优点是连续地使用特征结构，而不是只用单一平均 Jacobian。它通常能平滑处理跨声速稀疏波，并对缓慢激波表现出较少的奇偶振荡，但路径选择仍会影响极端状态下的鲁棒性。在编程中可把单波规则写成统一函数：输入两端特征速度和三个物理通量，先判断同号，再判断符号变化方向，最后返回相应的通量差。三波系统只需依次调用三次并累加。双稀疏波近似用等熵关系闭合交点状态。它在稀疏波占主导时很自然，但若真实波为强激波，等熵延拓不再精确；本章的极端算例正用于揭示这种近似与路径排序叠加后的后果。

图12.13显示 P 排序在测试5的异常：密度峰值、压力过冲和错误速度平台远离精确解，右侧应传播的激波消失。图12.14为对撞流测试6，在 $x_0=0.5,t=0.8$ 比较数值与精确解；P 排序虽然能够完成计算并给出两个对称激波，但激波邻域可见非物理振荡。两组结果共同表明，单一固定排序无法对所有极端波型同时保持最佳鲁棒性与精度。

% [原PDF p431]
\chapter{标量方程的高阶方法与 TVD 方法}
以下所有阶数结论都假定解足够光滑，且把空间步长与时间步长按固定 Courant 数同时缩小。间断附近经典 Taylor 展开失效，所以高阶格式必须借助非线性限制，而不能仅依赖截断误差阶数。当 Courant 数超过一时，一条特征线会跨越多个单元，单界面黎曼问题不再足以描述积分区域。大时间步 WAF 通过扩大数据支撑处理这种情形，但非线性系统的波相互作用使推广明显更复杂。模板是否迎风取决于其点位配置是否随特征速度符号改变，而不只是某个系数较大。中心格式也可能带有迎风偏置，但它的支撑方向不随波速反转，这是后续中心TVD构造的基础。

本章要解决数值方法中的一对核心矛盾：既希望在光滑区获得高阶精度，又希望在大梯度或间断附近不出现虚假的非物理振荡。常系数高阶线性格式在大梯度附近会振荡，而单调格式虽不振荡，却至多只有一阶精度。这一限制由稍后讨论的 Godunov 定理概括。绕开它的基本途径是使格式成为非线性的；总变差递减（Total Variation Diminishing，TVD）方法正是最重要的一类非线性格式。

\section{引言}
本章先在一维标量问题这一最清楚的环境中建立 TVD 概念，再为非线性系统和多维问题的推广准备工具。TVD 理论严格成立的范围主要是一维标量守恒律；然而大量实际计算表明，它也是构造复杂系统格式的可靠经验准则。其思想可追溯到 Godunov、van Leer 以及 Boris--Book 的通量修正输运工作。

% [原PDF p432]
线性格式的系数矩条件既能检查一致性，也能自动识别误差阶。零阶矩保证常数解保持，一阶矩给出正确传播速度，二阶及更高矩则控制波形的色散和耗散误差。单元平均是守恒有限体积格式的基本未知量。重构可以改变单元内部形状，却必须保持单元积分不变；否则即使界面通量形式守恒，也会在重构阶段悄然破坏总守恒。稳定性、一致性和收敛性不能混为一谈。线性理论中稳定加一致推出收敛；非线性问题则还要检查熵条件、正性和网格加密行为，因此本章把线性分析作为构造指南而非最终保证。

总变差稳定格式已被证明具有收敛性，TVD 是其中要求数值解总变差随时间不增的一类。标量守恒律的精确熵解本身就具有这一性质，因此 TVD 格式是在离散层面模仿精确解，而不只是人为抑制振荡。

TVD 与传统人工黏性方法关系密切。二者都试图构造 $p>1$ 阶方法，同时消除或控制大梯度附近的振荡，实际机制都是增加适量耗散。人工黏性法把附加项显式写入偏微分方程；TVD 方法则把耗散嵌入数值通量，并用非线性限制器精细激活：光滑区尽量保留高阶格式，间断区自动靠近单调的一阶格式。TVD 思想最初多用于迎风格式，后来也扩展到中心格式，从而为人工黏性提供更统一、更有理论依据的解释。

% [原PDF p433]
模板是否迎风取决于其点位配置是否随特征速度符号改变，而不只是某个系数较大。中心格式也可能带有迎风偏置，但它的支撑方向不随波速反转，这是后续中心TVD构造的基础。公式中的通量权重可以直接转成更新系数，反之亦然。实际核对格式时应同时检查两种表示：通量形式便于守恒和边界实现，更新系数形式便于精度、稳定性与单调性分析。以下所有阶数结论都假定解足够光滑，且把空间步长与时间步长按固定 Courant 数同时缩小。间断附近经典 Taylor 展开失效，所以高阶格式必须借助非线性限制，而不能仅依赖截断误差阶数。

本章首先回顾格式的一致性、稳定性和精度，然后介绍两类高阶构造：加权平均通量（WAF）与 MUSCL 变量外推。随后讨论单调性、Godunov 定理与总变差，再分别建立通量限制器和斜率限制器方法，并给出线性平流算例。非线性系统的推广留到第14章和第16章。

\section{若干格式的基本性质}
以线性平流方程
\begin{equation}u_t+f(u)_x=0,qquad f(u)=au,\tag{13.1}\end{equation}
为模型，$a$ 是常波速。对 $a>0$，Godunov 一阶迎风格式为 $u_i^{n+1}=u_i^n-c(u_i^n-u_{i-1}^n)$；对 $a<0$ 则使用右侧差分。Lax--Wendroff 格式为
\begin{equation}
u_i^{n+1}=\tfrac12c(1+c)u_{i-1}^n+(1-c^2)u_i^n-\tfrac12c(1-c)u_{i+1}^n,
\quad c=\frac{a\Delta t}{\Delta x}.\tag{13.3--13.4}
\end{equation}

% [原PDF p434]
稳定性、一致性和收敛性不能混为一谈。线性理论中稳定加一致推出收敛；非线性问题则还要检查熵条件、正性和网格加密行为，因此本章把线性分析作为构造指南而非最终保证。WAF 通量把半时间层黎曼解中的各状态按其占据长度加权，因此二阶时间信息由波在半步内传播的几何位置自然产生。对线性常速问题，该几何平均恰好还原 Lax--Wendroff。线性格式的系数矩条件既能检查一致性，也能自动识别误差阶。零阶矩保证常数解保持，一阶矩给出正确传播速度，二阶及更高矩则控制波形的色散和耗散误差。

许多线性格式可统一写成
\begin{equation}u_i^{n+1}=\sum_{k=-k_L}^{k_R}b_k u_{i+k}^n.\tag{13.5}\end{equation}
常系数 $b_k$ 完全决定格式。支撑宽度为 $k_L+k_R+1$，使用点数少者称紧致格式。Godunov 迎风格式只用两个点，Fromm 格式则用四个点。

模板的偏置也很重要。$a>0$ 时迎风方向在左，$a<0$ 时在右，Godunov 模板随波速符号而翻转。Lax--Friedrichs、FORCE、Godunov 一阶中心格式与 Lax--Wendroff 的模板构型不随波速符号改变，故称中心格式；不过中心格式的迎风系数仍可大于下风系数。评价格式的四项基本性质是一致性、稳定性、收敛性和精度。

% [原PDF p435]
以下所有阶数结论都假定解足够光滑，且把空间步长与时间步长按固定 Courant 数同时缩小。间断附近经典 Taylor 展开失效，所以高阶格式必须借助非线性限制，而不能仅依赖截断误差阶数。当 Courant 数超过一时，一条特征线会跨越多个单元，单界面黎曼问题不再足以描述积分区域。大时间步 WAF 通过扩大数据支撑处理这种情形，但非线性系统的波相互作用使推广明显更复杂。模板是否迎风取决于其点位配置是否随特征速度符号改变，而不只是某个系数较大。中心格式也可能带有迎风偏置，但它的支撑方向不随波速反转，这是后续中心TVD构造的基础。

\subsection{精度}
局部截断误差的首项确定格式阶数。Roe 给出一个十分便捷的判据：式(13.5)在空间和时间均为 $p$ 阶，当且仅当
\begin{equation}
\sum_{k=-k_L}^{k_R}k^q b_k=(-c)^q,qquad 0\le q\le p.\tag{13.6}
\end{equation}
$q=0$ 正是格式的一致性条件。以 $a>0$ 的 Godunov 格式为例，$b_{-1}=c,b_0=1-c$；$q=0,1$ 的矩条件成立，而 $q=2$ 时左端为 $c\ne c^2$，所以它通常只有一阶精度。

% [原PDF p436]
线性格式的系数矩条件既能检查一致性，也能自动识别误差阶。零阶矩保证常数解保持，一阶矩给出正确传播速度，二阶及更高矩则控制波形的色散和耗散误差。单元平均是守恒有限体积格式的基本未知量。重构可以改变单元内部形状，却必须保持单元积分不变；否则即使界面通量形式守恒，也会在重构阶段悄然破坏总守恒。稳定性、一致性和收敛性不能混为一谈。线性理论中稳定加一致推出收敛；非线性问题则还要检查熵条件、正性和网格加密行为，因此本章把线性分析作为构造指南而非最终保证。

Lax--Wendroff 的三个系数满足 $q=0,1,2$ 的矩条件，故为空间、时间二阶；但 $q=3$ 时左端为 $-c$，一般不等于 $-c^3$，所以不是三阶。相同判据可验证 Lax--Friedrichs、FORCE 和 Godunov 中心格式是一阶，Warming--Beam 与 Fromm 是二阶。还可证明，在三点中心模板上，任何时空二阶线性格式必然就是 Lax--Wendroff 格式。

\subsection{稳定性}
对线性问题，Lax 等价定理指出：一致格式收敛当且仅当稳定。非线性问题缺少同样完整的理论，实际中常结合线性稳定性分析、网格加密试验和可用的精确解判断。只在单一网格上获得平滑曲线并不足以证明稳定或收敛。

% [原PDF p437]
模板是否迎风取决于其点位配置是否随特征速度符号改变，而不只是某个系数较大。中心格式也可能带有迎风偏置，但它的支撑方向不随波速反转，这是后续中心TVD构造的基础。公式中的通量权重可以直接转成更新系数，反之亦然。实际核对格式时应同时检查两种表示：通量形式便于守恒和边界实现，更新系数形式便于精度、稳定性与单调性分析。以下所有阶数结论都假定解足够光滑，且把空间步长与时间步长按固定 Courant 数同时缩小。间断附近经典 Taylor 展开失效，所以高阶格式必须借助非线性限制，而不能仅依赖截断误差阶数。

对线性平流方程的三点中心格式
\begin{equation}u_i^{n+1}=b_{-1}u_{i-1}^n+b_0u_i^n+b_1u_{i+1}^n,\end{equation}
稳定性的充分必要条件可写成
\begin{equation}
b_0(b_{-1}+b_1)\ge0,qquad
b_0(b_{-1}+b_1)+4b_{-1}b_1\ge0.\tag{13.7}
\end{equation}
代入 Godunov 迎风系数，两式均化为 $c(1-c)\ge0$，结合波速符号得
\begin{equation}0\le |c|\le1.\tag{13.8}\end{equation}
Lax--Friedrichs 和 FORCE 也满足同一 CFL 区间；Godunov 一阶中心格式的限制更严，为 $|c|\le1/(2\sqrt2)$。以下高阶方法均采用显式全离散、基于黎曼问题、便于推广到非线性守恒律的构造。

% [原PDF p438]
稳定性、一致性和收敛性不能混为一谈。线性理论中稳定加一致推出收敛；非线性问题则还要检查熵条件、正性和网格加密行为，因此本章把线性分析作为构造指南而非最终保证。WAF 通量把半时间层黎曼解中的各状态按其占据长度加权，因此二阶时间信息由波在半步内传播的几何位置自然产生。对线性常速问题，该几何平均恰好还原 Lax--Wendroff。线性格式的系数矩条件既能检查一致性，也能自动识别误差阶。零阶矩保证常数解保持，一阶矩给出正确传播速度，二阶及更高矩则控制波形的色散和耗散误差。

\section{WAF 型高阶格式}
加权平均通量（Weighted Average Flux，WAF）把 Lax--Wendroff、Godunov 一阶迎风以及 Warming--Beam 的思想统一推广到非线性系统。它源于随机通量方法，但 WAF 是确定性的全离散显式二阶格式。

对凸通量标量守恒律
\begin{equation}u_t+f(u)_x=0,\qquad
u_i^{n+1}=u_i^n+\frac{\Delta t}{\Delta x}(f_{i-1/2}-f_{i+1/2}),\tag{13.10--13.11}
\end{equation}
$u_i^n$ 表示单元 $I_i=[x_{i-1/2},x_{i+1/2}]$ 上的积分平均：
\begin{equation}u_i^n=\frac1{\Delta x}\int_{x_{i-1/2}}^{x_{i+1/2}}u(x,t^n)\,dx.\tag{13.12}\end{equation}

% [原PDF p439]
以下所有阶数结论都假定解足够光滑，且把空间步长与时间步长按固定 Courant 数同时缩小。间断附近经典 Taylor 展开失效，所以高阶格式必须借助非线性限制，而不能仅依赖截断误差阶数。当 Courant 数超过一时，一条特征线会跨越多个单元，单界面黎曼问题不再足以描述积分区域。大时间步 WAF 通过扩大数据支撑处理这种情形，但非线性系统的波相互作用使推广明显更复杂。模板是否迎风取决于其点位配置是否随特征速度符号改变，而不只是某个系数较大。中心格式也可能带有迎风偏置，但它的支撑方向不随波速反转，这是后续中心TVD构造的基础。

\subsection{基本 WAF 格式}
界面通量定义为局部黎曼解在半时间层上的空间平均：
\begin{equation}
f^{\rm waf}_{i+1/2}=\frac1{\Delta x}\int_{-\Delta x/2}^{\Delta x/2}
f\!\left[u_{i+1/2}(x,\Delta t/2)\right]dx.\tag{13.13}
\end{equation}
这里 $u_{i+1/2}$ 是以 $(u_i^n,u_{i+1}^n)$ 为分段常数初值的黎曼解。对线性平流，它只有两个常状态，由速度 $a$ 的直线分隔：
\begin{equation}u_{i+1/2}(x,t)=\begin{cases}u_i^n,&x/t<a,\\u_{i+1}^n,&x/t>a.
\end{cases}\tag{13.14}\end{equation}

% [原PDF p440]
线性格式的系数矩条件既能检查一致性，也能自动识别误差阶。零阶矩保证常数解保持，一阶矩给出正确传播速度，二阶及更高矩则控制波形的色散和耗散误差。单元平均是守恒有限体积格式的基本未知量。重构可以改变单元内部形状，却必须保持单元积分不变；否则即使界面通量形式守恒，也会在重构阶段悄然破坏总守恒。稳定性、一致性和收敛性不能混为一谈。线性理论中稳定加一致推出收敛；非线性问题则还要检查熵条件、正性和网格加密行为，因此本章把线性分析作为构造指南而非最终保证。

精确计算上述空间积分可得
\begin{equation}
f^{\rm waf}_{i+1/2}=\tfrac12(1+c)au_i^n+\tfrac12(1-c)au_{i+1}^n.\tag{13.16}
\end{equation}
这正是 Lax--Wendroff 通量。$a>0$ 时第一项为迎风通量，权重 $(1+c)/2$；第二项为下风通量，权重 $(1-c)/2$。$a<0$ 时两项角色交换，但迎风权重仍较大。由于其中包含 Godunov 通量，WAF 也可看作 Godunov 一阶方法的高阶推广。

% [原PDF p441]
模板是否迎风取决于其点位配置是否随特征速度符号改变，而不只是某个系数较大。中心格式也可能带有迎风偏置，但它的支撑方向不随波速反转，这是后续中心TVD构造的基础。公式中的通量权重可以直接转成更新系数，反之亦然。实际核对格式时应同时检查两种表示：通量形式便于守恒和边界实现，更新系数形式便于精度、稳定性与单调性分析。以下所有阶数结论都假定解足够光滑，且把空间步长与时间步长按固定 Courant 数同时缩小。间断附近经典 Taylor 展开失效，所以高阶格式必须借助非线性限制，而不能仅依赖截断误差阶数。

线性平流的各种界面通量都可写成
\begin{equation}f_{i+1/2}=\sum_{k=-1}^{2}\beta_k(au_{i+k}^n).\tag{13.17}\end{equation}
系数 $\beta_k$ 使 Lax--Friedrichs、FORCE、迎风 Godunov、中心 Godunov、Lax--Wendroff、Warming--Beam、Fromm 与 SLIC 能在同一框架比较。对无黏 Burgers 方程也可直接把激波或稀疏波的精确黎曼解代入式(13.13)，但积分区间内的波型与声速点必须分别处理。

\subsection{WAF 格式的推广}
非均匀网格上，用相邻单元宽度的平均替代 $\Delta x$，积分上下限分别为 $-(\Delta x)_i/2$ 与 $(\Delta x)_{i+1}/2$。更一般的时空平均定义为
\begin{equation}
f^{\rm waf}_{i+1/2}=\frac1{t_2-t_1}\frac1{x_2-x_1}
\int_{t_1}^{t_2}\int_{x_1}^{x_2} f[\widehat u_{i+1/2}(x,t)]\,dx\,dt.\tag{13.19}
\end{equation}

% [原PDF p442]
稳定性、一致性和收敛性不能混为一谈。线性理论中稳定加一致推出收敛；非线性问题则还要检查熵条件、正性和网格加密行为，因此本章把线性分析作为构造指南而非最终保证。WAF 通量把半时间层黎曼解中的各状态按其占据长度加权，因此二阶时间信息由波在半步内传播的几何位置自然产生。对线性常速问题，该几何平均恰好还原 Lax--Wendroff。线性格式的系数矩条件既能检查一致性，也能自动识别误差阶。零阶矩保证常数解保持，一阶矩给出正确传播速度，二阶及更高矩则控制波形的色散和耗散误差。

取 $x_{1,2}=\mp\Delta x/2$、$t_1=0,t_2=\Delta t$，以中点公式近似时间积分，并令 $\widehat u$ 为单一黎曼问题解，在 $|c|<1$ 时便退化为基本 WAF。若时间方向精确积分，$c\le1/2$ 时仍得到 Lax--Wendroff，而 $1/2\le c\le1$ 时得到另一二阶方法。

放宽 CFL 限制后，积分区域会穿过多个黎曼问题。线性平流可由此构造任意大时间步显式格式；非线性推广则困难得多。特别地，WAF--CFL2 格式允许 $|c|\le2$：$|c|<1$ 使用 Lax--Wendroff 通量；$1\le|c|\le2$ 使用 Warming--Beam 通量。例如 $a>0$ 时
\begin{equation}f^{\rm WB}_{i+1/2}=-\tfrac12(1-c)au_{i-1}^n+\tfrac12(3-c)au_i^n.\tag{13.21}\end{equation}

% [原PDF p443]
以下所有阶数结论都假定解足够光滑，且把空间步长与时间步长按固定 Courant 数同时缩小。间断附近经典 Taylor 展开失效，所以高阶格式必须借助非线性限制，而不能仅依赖截断误差阶数。当 Courant 数超过一时，一条特征线会跨越多个单元，单界面黎曼问题不再足以描述积分区域。大时间步 WAF 通过扩大数据支撑处理这种情形，但非线性系统的波相互作用使推广明显更复杂。模板是否迎风取决于其点位配置是否随特征速度符号改变，而不只是某个系数较大。中心格式也可能带有迎风偏置，但它的支撑方向不随波速反转，这是后续中心TVD构造的基础。

当 $a>0,1\le c\le2$ 时，半时间层的特征线越过一个完整单元，归一化权重为 $\beta_{-1}=-(1-c)/2$、$\beta_0=(3-c)/2$；$a<0$ 时相应为 $\beta_1=(3+c)/2$、$\beta_2=-(1+c)/2$。该方法已推广到二维非定常欧拉方程，并配有适合的大时间步 TVD 条件。

若空间用中点公式而时间精确积分，则得到 CFL 上限为2的一阶迎风格式：$|c|\le1$ 时为普通 Godunov 迎风通量；$1\le c\le2,a>0$ 时为 $[(c-1)au_{i-1}^n+au_i^n]/c$，负速度时取对称形式。它在 $|c|=1$ 和 $2$ 时都精确地把数据平移整数个网格。

% [原PDF p444]
MUSCL 的边界外推值不是新的独立未知量，而是由单元平均和受限斜率重构得到。半步预测把单元内部通量差对两端值施加相同修正，使界面黎曼数据近似位于时间中点。半离散线法把空间重构和时间积分解耦，适合高阶 Runge--Kutta、ENO 和 WENO；全离散 Hancock 预测则把二阶时间精度紧凑地嵌入一步更新。两者的 CFL 与边界处理需要分别分析。斜率参数控制左右差分的偏置。未限制时可选择参数获得二阶或三阶形式；限制后局部阶数会随数据变化，光滑单调区趋向高阶，极值与间断附近则主动降为一阶。

WAF 的 TVD 修正将在第13.7节构造；非线性系统、源项与多维推广分别见第14、15和16章。下面转向另一种提高阶数的方法。

\section{MUSCL 型高阶方法}
van Leer 提出的 MUSCL（Monotone Upstream-centred Scheme for Conservation Laws，守恒律单调迎风中心格式）首先把 Godunov 方法的分段常数数据改为高阶重构。其两项核心是：用重构获得高阶精度；约束重构以避免伪振荡。本节先只讨论提高精度，TVD 限制稍后引入。

\subsection{数据重构}
最简单的选择是在每个单元内以线性函数替代常数平均值：
\begin{equation}u_i(x)=u_i^n+\frac{x-x_i}{\Delta x}\Delta_i,\tag{13.23}\end{equation}
其中 $\Delta_i/\Delta x$ 是单元斜率。

% [原PDF p445]
PLM、GRP 与 MUSCL--Hancock 在线性平流上得到同一通量，并不意味着它们在非线性系统中完全等价。差别来自非线性波速、广义黎曼问题的时间导数以及所采用的近似闭合。隐式格式可以放宽显式稳定性时间步限制，却不会自动提高间断分辨率或保证熵正确。每步的非线性代数求解还可能占据主要成本，故它特别适合以时间推进求稳态而非所有瞬态问题。中心 SLIC 用低阶中心通量替代界面黎曼求解器，其高阶信息仍来自重构和半步预测。它在复杂方程组中实现简单，但耗散与波分辨率更依赖所选低阶通量及限制器。

线性重构在单元左右边界的外推值为
\begin{equation}u_i^L=u_i^n-\tfrac12\Delta_i,qquad u_i^R=u_i^n+\tfrac12\Delta_i.\tag{13.24}\end{equation}
重构函数的单元积分与原平均值相同，因而保持守恒。相邻单元在界面通常给出不同的 $u_i^R$ 与 $u_{i+1}^L$，于是形成广义黎曼问题（GRP）：初值在两侧不再是常数，而是各自的线性函数；其波轨迹在 $x$--$t$ 平面中通常为曲线。

% [原PDF p446]
斜率参数控制左右差分的偏置。未限制时可选择参数获得二阶或三阶形式；限制后局部阶数会随数据变化，光滑单调区趋向高阶，极值与间断附近则主动降为一阶。从低阶通量构造高阶版本时，必须保留低阶格式的一致性。如果左右状态相等，重构、预测和最终通量都应严格退化为同一物理通量，这是程序单元测试中最基本的一项。MUSCL 的边界外推值不是新的独立未知量，而是由单元平均和受限斜率重构得到。半步预测把单元内部通量差对两端值施加相同修正，使界面黎曼数据近似位于时间中点。

非线性系统的精确 GRP 极其复杂，实际方法往往放弃完整求解，而把边界外推值作为普通分段常数黎曼问题的数据：
\begin{equation}u(x,0)=\begin{cases}u_i^R,&x<0,\\u_{i+1}^L,&x>0.
\end{cases}\tag{13.26}\end{equation}
斜率取为左右差分的加权组合
\begin{equation}
\Delta_i=\tfrac12(1+\omega)\Delta u_{i-1/2}+
\tfrac12(1-\omega)\Delta u_{i+1/2},\quad -1\le\omega\le1,\tag{13.27}
\end{equation}
其中 $\Delta u_{i-1/2}=u_i^n-u_{i-1}^n$。$\omega=0$ 给出中心斜率。还可使用二次重构；适当选择参数（例如 $\kappa=1/3$）可得到三阶空间离散，PPM 则使用分段抛物线重构。

% [原PDF p447]
中心 SLIC 用低阶中心通量替代界面黎曼求解器，其高阶信息仍来自重构和半步预测。它在复杂方程组中实现简单，但耗散与波分辨率更依赖所选低阶通量及限制器。广义黎曼问题以分段多项式而非常状态为初值，因此波后的状态会随时间变化，波轨迹也可能弯曲。GRP 方法只提取界面通量所需的局部时间展开，避免求出整个时空解。PLM、GRP 与 MUSCL--Hancock 在线性平流上得到同一通量，并不意味着它们在非线性系统中完全等价。差别来自非线性波速、广义黎曼问题的时间导数以及所采用的近似闭合。

\subsection{MUSCL--Hancock 方法}
MUSCL--Hancock（MHM）以三步构造全离散二阶格式。

第一步，在每个单元按式(13.23)--(13.24)重构 $u_i^L,u_i^R$。第二步，把两个边界值在单元内部演化半个时间步：
\begin{equation}
\bar u_i^{L,R}=u_i^{L,R}+\frac{\Delta t}{2\Delta x}
[f(u_i^L)-f(u_i^R)].\tag{13.30}
\end{equation}
第三步，在界面以 $(\bar u_i^R,\bar u_{i+1}^L)$ 为左右数据求普通黎曼问题，并取自相似解在零速度处的状态。

% [原PDF p448]
MUSCL 的边界外推值不是新的独立未知量，而是由单元平均和受限斜率重构得到。半步预测把单元内部通量差对两端值施加相同修正，使界面黎曼数据近似位于时间中点。半离散线法把空间重构和时间积分解耦，适合高阶 Runge--Kutta、ENO 和 WENO；全离散 Hancock 预测则把二阶时间精度紧凑地嵌入一步更新。两者的 CFL 与边界处理需要分别分析。斜率参数控制左右差分的偏置。未限制时可选择参数获得二阶或三阶形式；限制后局部阶数会随数据变化，光滑单调区趋向高阶，极值与间断附近则主动降为一阶。

最终界面通量就是
\begin{equation}f_{i+1/2}=f[u_{i+1/2}(0)].\tag{13.32}\end{equation}
对 $f(u)=au$，半步演化后的外推值为
\begin{equation}
\bar u_i^L=u_i^n-\tfrac12(1+c)\Delta_i,qquad
\bar u_i^R=u_i^n+\tfrac12(1-c)\Delta_i.\tag{13.33}
\end{equation}
界面黎曼解按 $a$ 的符号选择 $\bar u_i^R$ 或 $\bar u_{i+1}^L$，从而自然保持迎风性。

% [原PDF p449]
PLM、GRP 与 MUSCL--Hancock 在线性平流上得到同一通量，并不意味着它们在非线性系统中完全等价。差别来自非线性波速、广义黎曼问题的时间导数以及所采用的近似闭合。隐式格式可以放宽显式稳定性时间步限制，却不会自动提高间断分辨率或保证熵正确。每步的非线性代数求解还可能占据主要成本，故它特别适合以时间推进求稳态而非所有瞬态问题。中心 SLIC 用低阶中心通量替代界面黎曼求解器，其高阶信息仍来自重构和半步预测。它在复杂方程组中实现简单，但耗散与波分辨率更依赖所选低阶通量及限制器。

用符号函数可统一写成
\begin{equation}
f^{\rm MHM}_{i+1/2}=\tfrac12[1+\operatorname{sign}(c)]f(\bar u_i^R)
+\tfrac12[1-\operatorname{sign}(c)]f(\bar u_{i+1}^L).\tag{13.35}
\end{equation}
将斜率式(13.27)展开后，通量系数依赖 $c$ 与 $\omega$。$\omega=0$ 恢复 Fromm 格式；$a>0$ 时 $\omega=1$ 为 Warming--Beam，$\omega=-1$ 为 Lax--Wendroff。任意 $\omega$ 都至少二阶；若
\begin{equation}\omega=\tfrac13[2c-\operatorname{sign}(c)],\tag{13.39}\end{equation}
则在线性平流下达到时空三阶。TVD 版本在第13.8节给出。

% [原PDF p450]
斜率参数控制左右差分的偏置。未限制时可选择参数获得二阶或三阶形式；限制后局部阶数会随数据变化，光滑单调区趋向高阶，极值与间断附近则主动降为一阶。从低阶通量构造高阶版本时，必须保留低阶格式的一致性。如果左右状态相等，重构、预测和最终通量都应严格退化为同一物理通量，这是程序单元测试中最基本的一项。MUSCL 的边界外推值不是新的独立未知量，而是由单元平均和受限斜率重构得到。半步预测把单元内部通量差对两端值施加相同修正，使界面黎曼数据近似位于时间中点。

\subsection{分段线性方法}
Colella 的分段线性方法（PLM）同样从 MUSCL 重构出发。它以梯形公式近似界面通量的时间平均：
\begin{equation}
f^{\rm PLM}_{i+1/2}=\tfrac12\{f[u^{\rm grp}_{i+1/2}(0,0)]+f[u^{\rm grp}_{i+1/2}(0,\Delta t)]\}.\tag{13.40}
\end{equation}
第一项由 $(u_i^R,u_{i+1}^L)$ 的普通黎曼问题直接得到。第二项通过从 $(x,t)=(0,\Delta t)$ 沿特征线反向追踪到初始重构函数取得；特征脚位于第 $i$ 或 $i+1$ 单元，取决于速度符号。这样无需显式求解完整 GRP，也能近似其界面时间演化。

% [原PDF p451]
中心 SLIC 用低阶中心通量替代界面黎曼求解器，其高阶信息仍来自重构和半步预测。它在复杂方程组中实现简单，但耗散与波分辨率更依赖所选低阶通量及限制器。广义黎曼问题以分段多项式而非常状态为初值，因此波后的状态会随时间变化，波轨迹也可能弯曲。GRP 方法只提取界面通量所需的局部时间展开，避免求出整个时空解。PLM、GRP 与 MUSCL--Hancock 在线性平流上得到同一通量，并不意味着它们在非线性系统中完全等价。差别来自非线性波速、广义黎曼问题的时间导数以及所采用的近似闭合。

对线性平流，$a>0$ 时特征脚为单元 $i$ 内的 $(1-c)\Delta x$，代入重构得 $u_i^n+(1/2-c)\Delta_i$；$a<0$ 时作镜像处理。PLM 通量最终为
\begin{equation}
f^{\rm PLM}_{i+1/2}=\begin{cases}
a[u_i^n+\tfrac12(1-c)\Delta_i],&a>0,\\
a[u_{i+1}^n-\tfrac12(1+c)\Delta_{i+1}],&a<0.
\end{cases}\tag{13.43}
\end{equation}
在线性方程上，这与 MUSCL--Hancock 通量完全相同，并具有相同的二阶或特定参数下的三阶精度。

% [原PDF p452]
MUSCL 的边界外推值不是新的独立未知量，而是由单元平均和受限斜率重构得到。半步预测把单元内部通量差对两端值施加相同修正，使界面黎曼数据近似位于时间中点。半离散线法把空间重构和时间积分解耦，适合高阶 Runge--Kutta、ENO 和 WENO；全离散 Hancock 预测则把二阶时间精度紧凑地嵌入一步更新。两者的 CFL 与边界处理需要分别分析。斜率参数控制左右差分的偏置。未限制时可选择参数获得二阶或三阶形式；限制后局部阶数会随数据变化，光滑单调区趋向高阶，极值与间断附近则主动降为一阶。

\subsection{广义黎曼问题方法}
GRP 方法直接以广义黎曼解在半时间层的值定义通量：
\begin{equation}f^{\rm GRP}_{i+1/2}=f[u^{\rm grp}_{i+1/2}(0,\Delta t/2)].\tag{13.44}\end{equation}
其步骤是先作单元重构，再在界面原点作 Taylor 展开
\begin{equation}
u^{\rm grp}(0,\Delta t/2)=u^{\rm grp}(0,0)+\frac{\Delta t}{2}u_t^{\rm grp}(0,0)
+O(\Delta t^2).\tag{13.45}
\end{equation}
第一项仅由界面两侧外推值决定，等于普通黎曼问题在 $x/t=0$ 处的状态；第二项需要获得界面处的时间导数。

% [原PDF p453]
PLM、GRP 与 MUSCL--Hancock 在线性平流上得到同一通量，并不意味着它们在非线性系统中完全等价。差别来自非线性波速、广义黎曼问题的时间导数以及所采用的近似闭合。隐式格式可以放宽显式稳定性时间步限制，却不会自动提高间断分辨率或保证熵正确。每步的非线性代数求解还可能占据主要成本，故它特别适合以时间推进求稳态而非所有瞬态问题。中心 SLIC 用低阶中心通量替代界面黎曼求解器，其高阶信息仍来自重构和半步预测。它在复杂方程组中实现简单，但耗散与波分辨率更依赖所选低阶通量及限制器。

对线性平流，$u_t=-au_x$。任意空间导数 $v=\partial_x^p u$ 仍满足 $v_t+av_x=0$，故可给斜率建立另一个分段常数黎曼问题：左、右数据分别为 $\Delta_i/\Delta x$ 与 $\Delta_{i+1}/\Delta x$。由其迎风解求得 $u_x$，再代回 Taylor 展开，恰好得到
\begin{equation}
f^{\rm GRP}_{i+1/2}=\begin{cases}
a[u_i^n+\tfrac12(1-c)\Delta_i],&a>0,\\
a[u_{i+1}^n-\tfrac12(1+c)\Delta_{i+1}],&a<0.
\end{cases}\tag{13.49}
\end{equation}
因此在线性问题上 GRP、PLM 和 MUSCL--Hancock 三种推导产生同一通量。

% [原PDF p454]
斜率参数控制左右差分的偏置。未限制时可选择参数获得二阶或三阶形式；限制后局部阶数会随数据变化，光滑单调区趋向高阶，极值与间断附近则主动降为一阶。从低阶通量构造高阶版本时，必须保留低阶格式的一致性。如果左右状态相等，重构、预测和最终通量都应严格退化为同一物理通量，这是程序单元测试中最基本的一项。MUSCL 的边界外推值不是新的独立未知量，而是由单元平均和受限斜率重构得到。半步预测把单元内部通量差对两端值施加相同修正，使界面黎曼数据近似位于时间中点。

把空间导数也作为黎曼变量的想法可推广到更高阶 GRP，还可能用于对流--扩散方程，因为黏性通量本身含空间导数；已有研究显示可由此扩大稳定域。GRP 也已用于欧拉方程和激波绕射问题。

\subsection{斜率限制中心格式}
MUSCL 重构不必只能提高 Godunov 迎风通量。设任意低阶两状态通量为 $f^{\rm LO}(u_L,u_R)$，例如 Lax--Friedrichs、Godunov 中心或 FORCE。斜率限制中心格式（SLIC）先作重构和半步演化，再直接计算
\begin{equation}f^{\rm SLIC}_{i+1/2}=f^{\rm LO}(\bar u_i^R,\bar u_{i+1}^L).\tag{13.52}\end{equation}
这样无需黎曼求解器，也能从低阶中心方法构造高阶格式。

% [原PDF p455]
中心 SLIC 用低阶中心通量替代界面黎曼求解器，其高阶信息仍来自重构和半步预测。它在复杂方程组中实现简单，但耗散与波分辨率更依赖所选低阶通量及限制器。广义黎曼问题以分段多项式而非常状态为初值，因此波后的状态会随时间变化，波轨迹也可能弯曲。GRP 方法只提取界面通量所需的局部时间展开，避免求出整个时空解。PLM、GRP 与 MUSCL--Hancock 在线性平流上得到同一通量，并不意味着它们在非线性系统中完全等价。差别来自非线性波速、广义黎曼问题的时间导数以及所采用的近似闭合。

在线性平流下，若低阶通量写为
\begin{equation}f^{\rm LO}_{i+1/2}=\beta_0(au_L)+\beta_1(au_R),\tag{13.54}\end{equation}
则 SLIC 高阶通量为
\begin{equation}
f^{\rm SLIC}_{i+1/2}=\beta_0a[u_i^n+\tfrac12(1-c)\Delta_i]
+\beta_1a[u_{i+1}^n-\tfrac12(1+c)\Delta_{i+1}].\tag{13.55}
\end{equation}
代入守恒更新会得到从 $i-2$ 到 $i+2$ 的五点中心格式。其系数由 $\beta_0,\beta_1,c,\omega$ 显式决定，既包含基本低阶通量的偏置，也包含 MUSCL 重构的斜率信息。

% [原PDF p456]
MUSCL 的边界外推值不是新的独立未知量，而是由单元平均和受限斜率重构得到。半步预测把单元内部通量差对两端值施加相同修正，使界面黎曼数据近似位于时间中点。半离散线法把空间重构和时间积分解耦，适合高阶 Runge--Kutta、ENO 和 WENO；全离散 Hancock 预测则把二阶时间精度紧凑地嵌入一步更新。两者的 CFL 与边界处理需要分别分析。斜率参数控制左右差分的偏置。未限制时可选择参数获得二阶或三阶形式；限制后局部阶数会随数据变化，光滑单调区趋向高阶，极值与间断附近则主动降为一阶。

只要低阶通量一致，SLIC 对任意 $\omega$ 都达到时空二阶。若该低阶通量本身为一阶，则取
\begin{equation}
\omega=\frac{3c(\beta_0-\beta_1)-2c^2-1}{3(\beta_0-\beta_1-c)}\tag{13.58}
\end{equation}
可达三阶；分母为零的唯一情形是低阶通量实际已变成二阶 Lax--Wendroff。对 Godunov 迎风、Godunov 中心、FORCE 与 Lax--Friedrichs，三阶参数分别化为 $(2c-\operatorname{sign}c)/3$、$(4c^2-1)/(3c)$、$1/(3c)$ 和 $2c/3$。这些中心格式的 TVD 版本将在第13.8节构造。

% [原PDF p457]
PLM、GRP 与 MUSCL--Hancock 在线性平流上得到同一通量，并不意味着它们在非线性系统中完全等价。差别来自非线性波速、广义黎曼问题的时间导数以及所采用的近似闭合。隐式格式可以放宽显式稳定性时间步限制，却不会自动提高间断分辨率或保证熵正确。每步的非线性代数求解还可能占据主要成本，故它特别适合以时间推进求稳态而非所有瞬态问题。中心 SLIC 用低阶中心通量替代界面黎曼求解器，其高阶信息仍来自重构和半步预测。它在复杂方程组中实现简单，但耗散与波分辨率更依赖所选低阶通量及限制器。

\subsection{半离散格式}
除全离散显式方法外，还可采用线法先离散空间：
\begin{equation}
\frac{du_i}{dt}=\frac1{\Delta x}(f_{i-1/2}-f_{i+1/2}),\qquad
f_{i+1/2}=f_{i+1/2}(\{u_j(t)\},t).\tag{13.59--13.60}
\end{equation}
界面通量可用 Godunov 一阶通量，也可先 MUSCL 重构后求黎曼问题得到高阶空间离散；更高阶重构自然给出更高阶空间精度。所得常微分方程可用 Euler 或 Runge--Kutta 方法推进。空间与时间分开带来极大灵活性，尤其适合 UNO、ENO 等很高阶格式。

% [原PDF p458]
斜率参数控制左右差分的偏置。未限制时可选择参数获得二阶或三阶形式；限制后局部阶数会随数据变化，光滑单调区趋向高阶，极值与间断附近则主动降为一阶。从低阶通量构造高阶版本时，必须保留低阶格式的一致性。如果左右状态相等，重构、预测和最终通量都应严格退化为同一物理通量，这是程序单元测试中最基本的一项。MUSCL 的边界外推值不是新的独立未知量，而是由单元平均和受限斜率重构得到。半步预测把单元内部通量差对两端值施加相同修正，使界面黎曼数据近似位于时间中点。

\subsection{隐式格式}
隐式守恒格式仍写成
\begin{equation}u_i^{n+1}=u_i^n+\frac{\Delta t}{\Delta x}(f_{i-1/2}-f_{i+1/2}),\tag{13.61}\end{equation}
但 $f_{i+1/2}$ 同时依赖已知的 $u^n$ 和未知的 $u^{n+1}$。每步必须求解代数方程组，增加计算与存储成本；优点是理论上时间步不再受显式稳定性条件严格限制，尤其适合通过伪时间推进求稳态流动。航空稳态计算是其重要应用。

以上高阶格式若不加修改，都会在大梯度附近产生振荡。后续目标是在光滑区保存高阶精度，同时消除这些不需要的极值。

\section{单调格式与精度}
\subsection{单调格式}
若格式
\begin{equation}u_i^{n+1}=H(u_{i-k_L+1}^n,\ldots,u_{i+k_R}^n)\tag{13.63}\end{equation}
满足对每个自变量 $\partial H/\partial u_j^n\ge0$，则称为单调格式。

% [原PDF p459]
单调性是对更新算子每个参数的序关系要求，TVD 则约束相邻差绝对值之和。非线性标量问题中单调格式必为TVD，但TVD允许更宽的非线性格式集合，因而有机会在光滑区达到高阶。数据比值在分母接近零时需要数值上谨慎处理。通常按上风差与局部差的符号和相对大小直接实现限制器，并为零局部差规定一致极限，避免任意小量改变波形。Godunov 定理禁止的是固定系数线性格式同时具有单调性和二阶精度。限制器使系数成为数据比值的函数，格式即使求解线性方程也已非线性，因此不与定理矛盾。

等价地，若两组数据逐点满足 $v_i^n\ge u_i^n$，推进后仍有 $v_i^{n+1}\ge u_i^{n+1}$。这离散地复制了标量守恒律精确解的比较原理。单调格式满足离散最大值原理：
\begin{equation}
\max_i u_i^{n+1}\le\max_i u_i^n,qquad
\min_i u_i^{n+1}\ge\min_i u_i^n.\tag{13.66}
\end{equation}
因此不会产生新极值或伪振荡，并有 $\min_j u_j^n\le u_i^{n+1}\le\max_j u_j^n$。代价是原有峰值被削低、谷值被抬高，即极值削顶，反映出明显数值耗散。

% [原PDF p460]
离散最大值原理保证新解不会超出旧数据的全局范围，却不直接限定间断宽度。一阶单调格式常以较大数值黏性换取这一性质，高分辨率方法则试图只在需要的局部加入耗散。光滑线性区域内相邻斜率近似相等，因而比值趋近一；限制器通过点 $(1,1)$ 才能恢复二阶高阶通量。局部极值附近两侧差分异号，比值为负，TVD限制器通常取零。总变差把所有网格跳跃的绝对值相加，新生过冲和欠冲都会增加这一量。TVD条件比目测曲线更精确，能够把无振荡要求转化为可验证的系数不等式。

对三点守恒格式，单调性要求两状态数值通量对左参数非减、对右参数非增：
\begin{equation}
\frac{\partial f_{i+1/2}}{\partial u_i^n}\ge0,qquad
\frac{\partial f_{i+1/2}}{\partial u_{i+1}^n}\le0.\tag{13.70}
\end{equation}
对线性格式式(13.5)，单调性的充要条件就是全部 $b_k\ge0$。Godunov 一阶迎风格式满足该条件，Lax--Wendroff 及表中其他二阶常系数格式则不满足。Godunov 一阶中心格式只在其稳定区的后半段 $1/2\le|c|\le1/(2\sqrt2)$ 单调。

% [原PDF p461]
Godunov 定理禁止的是固定系数线性格式同时具有单调性和二阶精度。限制器使系数成为数据比值的函数，格式即使求解线性方程也已非线性，因此不与定理矛盾。高分辨率并不意味着处处使用最大允许反扩散。过于压缩的函数会把光滑峰顶方化，长期积分还会放大相位误差；限制器的选择应由间断锐度与光滑精度的实际需求平衡。单调性是对更新算子每个参数的序关系要求，TVD 则约束相邻差绝对值之和。非线性标量问题中单调格式必为TVD，但TVD允许更宽的非线性格式集合，因而有机会在光滑区达到高阶。

\subsection{一个启发性例子}
考虑参数化 WAF 型通量
\begin{equation}
f_{i+1/2}^{\alpha}=\frac1{\Delta x}\int_{-\Delta x/2}^{\Delta x/2}
f[u_{i+1/2}(x,\alpha\Delta t)]dx,
\tag{13.72}
\end{equation}
其中 $\alpha\ge0$。在线性平流下可精确算得
\begin{equation}
f_{i+1/2}^{\alpha}=\tfrac12(1+2\alpha c)au_i^n+
\tfrac12(1-2\alpha c)au_{i+1}^n.\tag{13.73}
\end{equation}
$\alpha=1/2$ 恢复 WAF/Lax--Wendroff；增大 $\alpha$ 会增加迎风贡献并减小下风贡献。

% [原PDF p462]
总变差把所有网格跳跃的绝对值相加，新生过冲和欠冲都会增加这一量。TVD条件比目测曲线更精确，能够把无振荡要求转化为可验证的系数不等式。Harten 条件是充分条件而非所有TVD格式的必要条件，但它形式局部、便于构造。将更新写成左右增量的非负组合后，只要总系数不超过一，相邻差求和时附加项就会望远镜式相消。离散最大值原理保证新解不会超出旧数据的全局范围，却不直接限定间断宽度。一阶单调格式常以较大数值黏性换取这一性质，高分辨率方法则试图只在需要的局部加入耗散。

不同参数给出熟知格式：Lax--Wendroff 为 $\alpha=1/2$，Godunov 迎风为 $1/(2|c|)$，Godunov 中心为 $1$，FORCE 为 $(1+c^2)/(4c^2)$，Lax--Friedrichs 为 $1/(2c^2)$。更新系数是
\begin{equation}
b_{-1}=\tfrac12(1+2\alpha c)c,\quad b_0=1-2\alpha c^2,\quad
b_1=-\tfrac12(1-2\alpha c)c.\tag{13.76}
\end{equation}
稳定区间为
\begin{equation}\tfrac12\le\alpha\le\frac1{2c^2},\tag{13.77}\end{equation}
单调区间则为 $1/(2|c|)\le\alpha\le1/(2c^2)$。只有 $\alpha=1/2$ 是二阶，其余稳定参数均为一阶。

% [原PDF p463]
单调性是对更新算子每个参数的序关系要求，TVD 则约束相邻差绝对值之和。非线性标量问题中单调格式必为TVD，但TVD允许更宽的非线性格式集合，因而有机会在光滑区达到高阶。数据比值在分母接近零时需要数值上谨慎处理。通常按上风差与局部差的符号和相对大小直接实现限制器，并为零局部差规定一致极限，避免任意小量改变波形。Godunov 定理禁止的是固定系数线性格式同时具有单调性和二阶精度。限制器使系数成为数据比值的函数，格式即使求解线性方程也已非线性，因此不与定理矛盾。

因此 $1/2<\alpha<1/(2|c|)$ 存在无穷多个稳定、一阶、却非单调的格式。方波算例取 $M=1000,c=1/4,t=1$：$\alpha=1/2$ 的 Lax--Wendroff 间断过渡窄但振荡；$\alpha=3/4$ 的一阶非单调格式和 $\alpha=1$ 的 Godunov 中心格式仍振荡；$\alpha=2$ 的迎风 Godunov、$\alpha=4.25$ 的 FORCE 与 $\alpha=8$ 的 Lax--Friedrichs 无振荡，但越来越扩散。

这个例子显示，高阶分辨率和单调性对固定系数线性格式互相冲突。自然的出路是让 $\alpha$ 随局部数据变化：光滑区靠近二阶值，陡峭区进入单调区。这样即使方程线性，格式本身也成为非线性。

% [原PDF p464]
离散最大值原理保证新解不会超出旧数据的全局范围，却不直接限定间断宽度。一阶单调格式常以较大数值黏性换取这一性质，高分辨率方法则试图只在需要的局部加入耗散。光滑线性区域内相邻斜率近似相等，因而比值趋近一；限制器通过点 $(1,1)$ 才能恢复二阶高阶通量。局部极值附近两侧差分异号，比值为负，TVD限制器通常取零。总变差把所有网格跳跃的绝对值相加，新生过冲和欠冲都会增加这一量。TVD条件比目测曲线更精确，能够把无振荡要求转化为可验证的系数不等式。

图13.10的六组曲线把这种权衡直观展示出来。随着 $\alpha$ 增大，间断附近的过冲与欠冲逐渐消失，但过渡层同步加宽。Lax--Friedrichs 的总变差最小，却对方波最为涂抹；Lax--Wendroff 分辨率最高，却增加总变差。设计限制器的目标不是选择一个固定折中参数，而是逐界面、逐时间步适配局部光滑性。

% [原PDF p465]
Godunov 定理禁止的是固定系数线性格式同时具有单调性和二阶精度。限制器使系数成为数据比值的函数，格式即使求解线性方程也已非线性，因此不与定理矛盾。高分辨率并不意味着处处使用最大允许反扩散。过于压缩的函数会把光滑峰顶方化，长期积分还会放大相位误差；限制器的选择应由间断锐度与光滑精度的实际需求平衡。单调性是对更新算子每个参数的序关系要求，TVD 则约束相邻差绝对值之和。非线性标量问题中单调格式必为TVD，但TVD允许更宽的非线性格式集合，因而有机会在光滑区达到高阶。

\subsection{单调格式与 Godunov 定理}
计算经验表明，二阶线性格式对光滑解明显优于一阶格式，但在间断附近产生伪振荡；单调格式没有振荡，却精度有限。Godunov 定理把这一经验提升为严格结论：

\textbf{定理13.51（Godunov 定理）}\quad 对线性平流方程，不存在二阶或更高阶的单调线性格式。

证明使用精度矩条件。记 $s_q=\sum k^q b_k$；二阶精度要求 $s_0=1,s_1=-c,s_2=c^2$。将 $s_2$ 改写为
\begin{equation}s_2=\sum_k(k+c)^2b_k-2cs_1-c^2s_0.\tag{13.82}\end{equation}

% [原PDF p466]
总变差把所有网格跳跃的绝对值相加，新生过冲和欠冲都会增加这一量。TVD条件比目测曲线更精确，能够把无振荡要求转化为可验证的系数不等式。Harten 条件是充分条件而非所有TVD格式的必要条件，但它形式局部、便于构造。将更新写成左右增量的非负组合后，只要总系数不超过一，相邻差求和时附加项就会望远镜式相消。离散最大值原理保证新解不会超出旧数据的全局范围，却不直接限定间断宽度。一阶单调格式常以较大数值黏性换取这一性质，高分辨率方法则试图只在需要的局部加入耗散。

若格式单调，则 $b_k\ge0$，于是
\begin{equation}\sum_k(k+c)^2b_k+c^2\ge c^2.\tag{13.83}\end{equation}
要满足二阶等式，必须所有有非零权重的位置都满足 $k=-c$。除整数 Courant 数下精确平移这一退化情形外，这是不可能的。因此一般单调线性格式至多一阶。对非线性系统，要求所有特征波同时具有整数 Courant 数更不现实。

定理的关键假设是“固定系数”。所以绕开限制的方向不是继续寻找新的线性模板，而是让系数依解的局部形态变化。

\subsection{伪振荡与高分辨率}
高阶线性格式在陡峭波前后产生新的局部极值。Lax--Wendroff 的振荡多位于波后，Warming--Beam 多位于波前，这与各自截断误差首项的色散性质有关。

% [原PDF p467]
单调性是对更新算子每个参数的序关系要求，TVD 则约束相邻差绝对值之和。非线性标量问题中单调格式必为TVD，但TVD允许更宽的非线性格式集合，因而有机会在光滑区达到高阶。数据比值在分母接近零时需要数值上谨慎处理。通常按上风差与局部差的符号和相对大小直接实现限制器，并为零局部差规定一致极限，避免任意小量改变波形。Godunov 定理禁止的是固定系数线性格式同时具有单调性和二阶精度。限制器使系数成为数据比值的函数，格式即使求解线性方程也已非线性，因此不与定理矛盾。

所谓高分辨率格式应同时具有：光滑区二阶或更高阶；数值解不含伪振荡；间断过渡区比一阶单调方法窄。TVD 格式通过局部非线性限制近似实现这三项要求。

\subsection{数据相容性}
Roe 提出的早期思路是让算法适应局部数据。对线性平流，若新值 $u_i^{n+1}$ 位于其迎风数据对 $(u_{i-s}^n,u_i^n)$ 之间，其中 $s=\operatorname{sign}(c)$，则称格式与该数据相容。几何上，经过新时间层点的特征线脚恰位于这两个迎风网格值之间。

% [原PDF p468]
离散最大值原理保证新解不会超出旧数据的全局范围，却不直接限定间断宽度。一阶单调格式常以较大数值黏性换取这一性质，高分辨率方法则试图只在需要的局部加入耗散。光滑线性区域内相邻斜率近似相等，因而比值趋近一；限制器通过点 $(1,1)$ 才能恢复二阶高阶通量。局部极值附近两侧差分异号，比值为负，TVD限制器通常取零。总变差把所有网格跳跃的绝对值相加，新生过冲和欠冲都会增加这一量。TVD条件比目测曲线更精确，能够把无振荡要求转化为可验证的系数不等式。

相容条件为
\begin{equation}
\min(u_{i-s}^n,u_i^n)\le u_i^{n+1}\le\max(u_{i-s}^n,u_i^n),\tag{13.84}
\end{equation}
等价地
\begin{equation}0\le\frac{u_i^{n+1}-u_i^n}{u_{i-s}^n-u_i^n}\le1.\tag{13.85}\end{equation}
Godunov 一阶迎风更新可写成 $u_i^{n+1}=u_i^n+s c(u_{i-s}^n-u_i^n)$，故在 $|c|\le1$ 时与任意数据相容。一般三点格式可写成
\begin{equation}u_i^{n+1}=u_i^n-b_{-1}\Delta u_{i-1/2}+b_1\Delta u_{i+1/2}.\tag{13.88}\end{equation}

% [原PDF p469]
Godunov 定理禁止的是固定系数线性格式同时具有单调性和二阶精度。限制器使系数成为数据比值的函数，格式即使求解线性方程也已非线性，因此不与定理矛盾。高分辨率并不意味着处处使用最大允许反扩散。过于压缩的函数会把光滑峰顶方化，长期积分还会放大相位误差；限制器的选择应由间断锐度与光滑精度的实际需求平衡。单调性是对更新算子每个参数的序关系要求，TVD 则约束相邻差绝对值之和。非线性标量问题中单调格式必为TVD，但TVD允许更宽的非线性格式集合，因而有机会在光滑区达到高阶。

令
\begin{equation}r_i=\frac{\Delta u_{i-1/2}}{\Delta u_{i+1/2}},\tag{13.90}\end{equation}
则数据相容条件变为：$c>0$ 时 $0\le b_{-1}-b_1/r_i\le1$；$c<0$ 时 $0\le b_1-b_{-1}r_i\le1$。Lax--Wendroff 和 Warming--Beam 只对 $r_i$ 的某些区间相容。Roe 的思想是组合若干格式，使各自相容区的并集覆盖全部数据；所得算法系数依赖 $r_i$，因而为非线性，能够越过 Godunov 定理的固定系数限制。

\section{总变差递减方法}
非线性守恒律收敛证明通常依赖紧性，而总变差有界的函数族具有所需紧性。总变差稳定格式的网格近似处于紧函数集中，因而可以建立收敛性。TVD 是其中更强、也更易设计的一类：要求每一步的总变差不增加。

% [原PDF p470]
总变差把所有网格跳跃的绝对值相加，新生过冲和欠冲都会增加这一量。TVD条件比目测曲线更精确，能够把无振荡要求转化为可验证的系数不等式。Harten 条件是充分条件而非所有TVD格式的必要条件，但它形式局部、便于构造。将更新写成左右增量的非负组合后，只要总系数不超过一，相邻差求和时附加项就会望远镜式相消。离散最大值原理保证新解不会超出旧数据的全局范围，却不直接限定间断宽度。一阶单调格式常以较大数值黏性换取这一性质，高分辨率方法则试图只在需要的局部加入耗散。

\subsection{总变差}
函数 $u(x)$ 的总变差定义为
\begin{equation}
TV(u)=\lim_{\delta\to0}\sup\frac1\delta\int_{-\infty}^{\infty}|u(x+\delta)-u(x)|dx.\tag{13.91}
\end{equation}
光滑函数可写成 $TV(u)=\int|u'(x)|dx$；对间断函数，该式在分布意义下仍成立。网格函数的离散总变差为
\begin{equation}TV(u^n)=\sum_{i=-\infty}^{\infty}|u_{i+1}^n-u_i^n|.\tag{13.93}\end{equation}
为使其有限，远场应趋于零或常数。振荡越强，总变差越大；数值耗散越强，总变差通常下降越快。线性平流精确解的总变差保持不变，而 Lax--Friedrichs 解会随时间逐渐减小。

% [原PDF p471]
单调性是对更新算子每个参数的序关系要求，TVD 则约束相邻差绝对值之和。非线性标量问题中单调格式必为TVD，但TVD允许更宽的非线性格式集合，因而有机会在光滑区达到高阶。数据比值在分母接近零时需要数值上谨慎处理。通常按上风差与局部差的符号和相对大小直接实现限制器，并为零局部差规定一致极限，避免任意小量改变波形。Godunov 定理禁止的是固定系数线性格式同时具有单调性和二阶精度。限制器使系数成为数据比值的函数，格式即使求解线性方程也已非线性，因此不与定理矛盾。

标量守恒律的精确熵解不会产生新的局部极值；局部极小值随时间不降，局部极大值随时间不升。因此
\begin{equation}TV[u(t_2)]\le TV[u(t_1)],\qquad t_2\ge t_1.\tag{13.94}\end{equation}
若离散格式满足
\begin{equation}TV(u^{n+1})\le TV(u^n)\quad\text{对所有 }n,\tag{13.96}\end{equation}
则称为总变差递减或总变差不增格式。另一个相关概念是保持单调性：若输入序列单调递增或递减，输出仍保持同一单调方向。一般有集合包含关系
\begin{equation}\mathcal S_{\rm mon}\subseteq\mathcal S_{\rm TVD}\subseteq\mathcal S_{\rm mpr}.\tag{13.98}\end{equation}

% [原PDF p472]
离散最大值原理保证新解不会超出旧数据的全局范围，却不直接限定间断宽度。一阶单调格式常以较大数值黏性换取这一性质，高分辨率方法则试图只在需要的局部加入耗散。光滑线性区域内相邻斜率近似相等，因而比值趋近一；限制器通过点 $(1,1)$ 才能恢复二阶高阶通量。局部极值附近两侧差分异号，比值为负，TVD限制器通常取零。总变差把所有网格跳跃的绝对值相加，新生过冲和欠冲都会增加这一量。TVD条件比目测曲线更精确，能够把无振荡要求转化为可验证的系数不等式。

对线性平流的线性格式，上述三类实际上等价。若格式保持单调性，由相邻更新差
\begin{equation}u_{i+1}^{n+1}-u_i^{n+1}=\sum_kb_k(u_{i+1+k}^n-u_{i+k}^n)\tag{13.99}\end{equation}
可用只有一个跳跃的阶梯数据反证任何 $b_K<0$，故必有 $b_k\ge0$。反之，若所有系数非负，输入相邻差同号时输出差也同号。因此线性问题中“单调”“TVD”“保持单调性”是同一个条件。TVD 的价值在于它还有可用于非线性格式的明确代数判据和收敛理论。

% [原PDF p473]
Godunov 定理禁止的是固定系数线性格式同时具有单调性和二阶精度。限制器使系数成为数据比值的函数，格式即使求解线性方程也已非线性，因此不与定理矛盾。高分辨率并不意味着处处使用最大允许反扩散。过于压缩的函数会把光滑峰顶方化，长期积分还会放大相位误差；限制器的选择应由间断锐度与光滑精度的实际需求平衡。单调性是对更新算子每个参数的序关系要求，TVD 则约束相邻差绝对值之和。非线性标量问题中单调格式必为TVD，但TVD允许更宽的非线性格式集合，因而有机会在光滑区达到高阶。

Harten 考虑增量形式
\begin{equation}
u_i^{n+1}=u_i^n-C_{i-1/2}\Delta u_{i-1/2}+D_{i+1/2}\Delta u_{i+1/2}.\tag{13.101}
\end{equation}
一个重要充分条件是
\begin{equation}
C_{i+1/2}\ge0,qquad D_{i+1/2}\ge0,qquad
0\le C_{i+1/2}+D_{i+1/2}\le1.\tag{13.102}
\end{equation}
证明把相邻两个更新式相减，对所得三项取绝对值，再对 $i$ 求和。第一项系数 $1-C-D$ 非负，另两项的错位求和恰好相消，最终得到 $TV(u^{n+1})\le TV(u^n)$。

% [原PDF p474]
总变差把所有网格跳跃的绝对值相加，新生过冲和欠冲都会增加这一量。TVD条件比目测曲线更精确，能够把无振荡要求转化为可验证的系数不等式。Harten 条件是充分条件而非所有TVD格式的必要条件，但它形式局部、便于构造。将更新写成左右增量的非负组合后，只要总系数不超过一，相邻差求和时附加项就会望远镜式相消。离散最大值原理保证新解不会超出旧数据的全局范围，却不直接限定间断宽度。一阶单调格式常以较大数值黏性换取这一性质，高分辨率方法则试图只在需要的局部加入耗散。

$C,D$ 可以依赖局部数据，因此 Harten 定理直接适用于非线性格式，正好绕开 Godunov 定理只针对线性固定系数格式的限制。它是构造高分辨率格式最实用的工具之一。若允许总变差受控地略有增长，则得到总变差有界（TVB）格式。

\section{通量限制器方法}
下面先构造 WAF 的 TVD 版本，再给出一般高、低阶通量混合框架，并分别讨论迎风型和中心型限制器。

\subsection{WAF 的 TVD 版本}
把 WAF 通量改写为
\begin{equation}
f_{i+1/2}=\tfrac12(1+\phi)au_i^n+\tfrac12(1-\phi)au_{i+1}^n.\tag{13.105}
\end{equation}
$\phi=|c|$ 给出 Lax--Wendroff，$\phi=s=\operatorname{sign}(c)$ 给出 Godunov 迎风，$\phi=-s$ 给出无条件不稳定的下风格式。令 $-1\le\phi\le1$，并使其随局部数据变化。

% [原PDF p475]
通量限制器在低阶单调通量与高阶通量之间调节校正量。零值完全关闭高阶校正，一值恢复目标高阶格式；超出一的值属于反扩散，可使间断更锐利，但只能在TVD区域允许的范围内使用。Godunov点表示中心低阶通量经特定比例高阶校正后恰好达到一阶迎风通量。它为中心限制器与迎风限制器建立仿射映射，也清楚标出正耗散和较低耗散区域的分界。同名的迎风、WAF和中心限制器只有映射关系，不能把函数值直接互换。实现时应明确当前限制的是通量校正、WAF耗散参数还是重构斜率，并使用对应的比值定义。

对 $a>0$，将界面参数 $\phi_{i\pm1/2}$ 代入更新，并按数据相容性整理得
\begin{equation}
\frac{u_i^{n+1}-u_i^n}{u_{i-1}^n-u_i^n}
=\frac c2\left[\frac{1-\phi_{i+1/2}}{r_{i+1/2}}+phi_{i-1/2}+1\right],\tag{13.108}
\end{equation}
其中
\begin{equation}r_{i+1/2}=\frac{u_i^n-u_{i-1}^n}{u_{i+1}^n-u_i^n}.\tag{13.109}\end{equation}
要求左端位于 $[0,1]$，得到
\begin{equation}-1\le\frac{1-\phi_{i+1/2}}{r_{i+1/2}}+\phi_{i-1/2}
\le\frac{2-c}{c}.\tag{13.110}\end{equation}

% [原PDF p476]
Sweby 图把所有可接受限制器表示在比值—限制系数平面内。上边界决定最大压缩性，下边界保证不产生新极值，二阶点决定光滑区精度；这三个几何条件比逐个推导格式更易复用。CFL相关TVD区能利用当前时间步得到更宽的允许范围；CFL无关区域则对所有合法时间步统一安全。前者分辨率可能更高，后者实现简单且不易因局部波速更新滞后而越界。SUPERBEE 靠近TVD区上边界，适合保持接触面和方波；MINBEE靠近下界，更耗散但稳健；VANLEER与VANALBADA连续光滑，常在非线性迭代和复杂流场中提供较柔和响应。

对 $a<0$，结论相同，只需把 $c$ 换成 $|c|$，并把上风差取在界面右侧：
\begin{equation}
r_{i+1/2}=\begin{cases}
(u_i-u_{i-1})/(u_{i+1}-u_i),&a>0,\\
(u_{i+2}-u_{i+1})/(u_{i+1}-u_i),&a<0.
\end{cases}\tag{13.113}
\end{equation}
所以 $r=\Delta u_{\rm upw}/\Delta u_{\rm loc}$ 是上风变化与当前界面局部变化之比。限制器取 $\phi_{i+1/2}=\phi(r_{i+1/2},|c|)$，让数据本身决定附加耗散。

为构造 TVD 区域，引入下界 $L\in[-1,|c|]$，分别约束 $(1-\phi)/r$ 和相邻界面的 $\phi$。$L=-1$ 允许完全下风化，$L=0$ 禁止下风化。

% [原PDF p477]
同名的迎风、WAF和中心限制器只有映射关系，不能把函数值直接互换。实现时应明确当前限制的是通量校正、WAF耗散参数还是重构斜率，并使用对应的比值定义。数值试验表明，限制器性能具有问题依赖性。纯间断测试偏爱压缩函数，光滑波长期传播更需要低色散和温和限制；实际系统通常同时含两类结构，因此应综合多种基准题评估。通量限制器在低阶单调通量与高阶通量之间调节校正量。零值完全关闭高阶校正，一值恢复目标高阶格式；超出一的值属于反扩散，可使间断更锐利，但只能在TVD区域允许的范围内使用。

对 $r>0$，两条边界为
\begin{equation}
\phi\le\phi_L=1+(1+L)r,qquad
\phi\ge\phi_R=1-\frac{2(1-|c|)}{|c|}r;\tag{13.117,13.119}
\end{equation}
对 $r<0$ 不等号方向反转。局部变化与上风变化的位置随速度符号镜像：正速度取界面左侧上风差，负速度取右侧上风差。只要每个界面都按这一约定计算 $r$，同一限制器函数可同时用于两种传播方向。

% [原PDF p478]
SUPERBEE 靠近TVD区上边界，适合保持接触面和方波；MINBEE靠近下界，更耗散但稳健；VANLEER与VANALBADA连续光滑，常在非线性迭代和复杂流场中提供较柔和响应。中心低阶通量含有下风项，因此迎风格式的经典数据相容条件会过度收紧。扩展中心TVD判据分别处理单调区与局部极值，既保留原中心低阶格式，又允许向Lax--Wendroff加入受控高阶校正。Sweby 图把所有可接受限制器表示在比值—限制系数平面内。上边界决定最大压缩性，下边界保证不产生新极值，二阶点决定光滑区精度；这三个几何条件比逐个推导格式更易复用。

在 $r$--$\phi$ 平面中，上述直线与 $L\le\phi\le1$ 围成 TVD 区。$\phi\ge|c|$ 对应正数值黏性，会展宽间断、削平极值；$\phi\le|c|$ 对应负数值黏性，会使间断变陡；$\phi=|c|$ 就是没有额外黏性的 Lax--Wendroff。任何位于 TVD 区内的函数都无伪振荡，但为在光滑区保持二阶还必须满足
\begin{equation}\phi(1,|c|)=|c|.\tag{13.121}\end{equation}
$r\approx1$ 表示相邻差分近似相等，即局部解光滑。

% [原PDF p479]
通量限制器在低阶单调通量与高阶通量之间调节校正量。零值完全关闭高阶校正，一值恢复目标高阶格式；超出一的值属于反扩散，可使间断更锐利，但只能在TVD区域允许的范围内使用。Godunov点表示中心低阶通量经特定比例高阶校正后恰好达到一阶迎风通量。它为中心限制器与迎风限制器建立仿射映射，也清楚标出正耗散和较低耗散区域的分界。同名的迎风、WAF和中心限制器只有映射关系，不能把函数值直接互换。实现时应明确当前限制的是通量校正、WAF耗散参数还是重构斜率，并使用对应的比值定义。

常用 WAF 限制器可写为与标准通量限制器相对应的形式。ULTRAA 取 TVD 区最压缩的边界，SUPERA 为分段线性函数，VANLEER 与 VANALBADA 给出光滑有理函数，MINA 则更耗散。代表性表达式为
\begin{align}
\phi_{\rm VL}(r,|c|)&=\begin{cases}1,&r\le0,\\1-\dfrac{2(1-|c|)r}{1+r},&r\ge0,\end{cases}\tag{13.124}\\
\phi_{\rm VA}(r,|c|)&=\begin{cases}1,&r\le0,\\1-\dfrac{(1-|c|)r(1+r)}{1+r^2},&r\ge0,\end{cases}\tag{13.125}\\
\phi_{\rm MINA}(r,|c|)&=\begin{cases}1,&r\le0,\\1-(1-|c|)r,&0\le r\le1,\\|c|,&r\ge1.
\end{cases}\tag{13.126}
\end{align}
ULTRAA 不通过二阶点，因而不适合一般光滑解。

% [原PDF p480]
Sweby 图把所有可接受限制器表示在比值—限制系数平面内。上边界决定最大压缩性，下边界保证不产生新极值，二阶点决定光滑区精度；这三个几何条件比逐个推导格式更易复用。CFL相关TVD区能利用当前时间步得到更宽的允许范围；CFL无关区域则对所有合法时间步统一安全。前者分辨率可能更高，后者实现简单且不易因局部波速更新滞后而越界。SUPERBEE 靠近TVD区上边界，适合保持接触面和方波；MINBEE靠近下界，更耗散但稳健；VANLEER与VANALBADA连续光滑，常在非线性迭代和复杂流场中提供较柔和响应。

数值试验采用 $a=1,C_{\rm CFL}=0.8$。测试1是高斯初值 $u=\alpha e^{-\beta x^2}$，取 $\alpha=1,\beta=8$、区间 $[-1,1]$、$\Delta x=0.02$；测试2是 $[0.3,0.7]$ 内为1、其余为0的方波，$\Delta x=0.01$。两题均用100个单元，比较 $t=1$（125步）和 $t=10$（1250步）。参照方法为 Lax--Wendroff、Godunov 一阶迎风以及四种 WAF--TVD 限制器。

% [原PDF p481]
同名的迎风、WAF和中心限制器只有映射关系，不能把函数值直接互换。实现时应明确当前限制的是通量校正、WAF耗散参数还是重构斜率，并使用对应的比值定义。数值试验表明，限制器性能具有问题依赖性。纯间断测试偏爱压缩函数，光滑波长期传播更需要低色散和温和限制；实际系统通常同时含两类结构，因此应综合多种基准题评估。通量限制器在低阶单调通量与高阶通量之间调节校正量。零值完全关闭高阶校正，一值恢复目标高阶格式；超出一的值属于反扩散，可使间断更锐利，但只能在TVD区域允许的范围内使用。

对短时间光滑高斯，Lax--Wendroff 很准确，TVD 修正没有明显改进；Godunov 迎风则因数值黏性削平峰值。ULTRAA 最差，因为它不满足二阶点并会用负黏性错误地使轮廓变方。SUPERA 也有一定方波化倾向，VANLEER 与 MINA 较为平滑。

演化到 $t=10$ 后，Lax--Wendroff 的色散误差使数值波相对精确解滞后；Godunov 的扩散随步数积累。SUPERA、VANLEER、MINA 仍可接受，但 SUPERA 的削方更明显，VANLEER 和 MINA 在极值处也开始显出黏性。综合光滑解结果，SUPERA 与 VANLEER 表现较好。

% [原PDF p482]
SUPERBEE 靠近TVD区上边界，适合保持接触面和方波；MINBEE靠近下界，更耗散但稳健；VANLEER与VANALBADA连续光滑，常在非线性迭代和复杂流场中提供较柔和响应。中心低阶通量含有下风项，因此迎风格式的经典数据相容条件会过度收紧。扩展中心TVD判据分别处理单调区与局部极值，既保留原中心低阶格式，又允许向Lax--Wendroff加入受控高阶校正。Sweby 图把所有可接受限制器表示在比值—限制系数平面内。上边界决定最大压缩性，下边界保证不产生新极值，二阶点决定光滑区精度；这三个几何条件比逐个推导格式更易复用。

对含两个间断的方波，短时结果恰好相反：Lax--Wendroff 虽锐利却振荡，Godunov 虽无振荡却严重涂抹。TVD 结果结合了两者优点。ULTRAA 对纯间断最锐利，偶数时间步只有一个中间过渡点，奇数步甚至无中间点；MINA 最耗散，间断最宽。长时间后所有 TVD 结果都优于两个线性格式，排序大致为 ULTRAA、SUPERA、VANLEER、MINA。但实际解同时含光滑区和间断，故总体更推荐 SUPERA 与 VANLEER。

\subsection{一般通量限制器框架}
给定二阶或更高阶通量 $f^{\rm HI}$ 与单调一阶通量 $f^{\rm LO}$，定义
\begin{equation}
f^{\rm TVD}_{i+1/2}=f^{\rm LO}_{i+1/2}
+\psi_{i+1/2}(f^{\rm HI}_{i+1/2}-f^{\rm LO}_{i+1/2}).\tag{13.130}
\end{equation}
限制器 $\psi=0$ 退回鲁棒低阶，$\psi=1$ 恢复高阶；其余值在两者之间或允许受控反扩散。

% [原PDF p483]
通量限制器在低阶单调通量与高阶通量之间调节校正量。零值完全关闭高阶校正，一值恢复目标高阶格式；超出一的值属于反扩散，可使间断更锐利，但只能在TVD区域允许的范围内使用。Godunov点表示中心低阶通量经特定比例高阶校正后恰好达到一阶迎风通量。它为中心限制器与迎风限制器建立仿射映射，也清楚标出正耗散和较低耗散区域的分界。同名的迎风、WAF和中心限制器只有映射关系，不能把函数值直接互换。实现时应明确当前限制的是通量校正、WAF耗散参数还是重构斜率，并使用对应的比值定义。

在线性平流中设
\begin{equation}f^{\rm LO}=\alpha_0au_i+\alpha_1au_{i+1},\qquad
f^{\rm HI}=\beta_0au_i+\beta_1au_{i+1}.\tag{13.131}\end{equation}
Godunov 迎风、FORCE 和 Lax--Friedrichs 分别对应不同的 $\alpha_0,\alpha_1$；二点支撑的二阶高阶通量只能是 Lax--Wendroff，故 $\beta_0=(1+c)/2,\beta_1=(1-c)/2$。混合后两项系数分别为 $\alpha_j+(\beta_j-\alpha_j)\psi$。

% [原PDF p484]
Sweby 图把所有可接受限制器表示在比值—限制系数平面内。上边界决定最大压缩性，下边界保证不产生新极值，二阶点决定光滑区精度；这三个几何条件比逐个推导格式更易复用。CFL相关TVD区能利用当前时间步得到更宽的允许范围；CFL无关区域则对所有合法时间步统一安全。前者分辨率可能更高，后者实现简单且不易因局部波速更新滞后而越界。SUPERBEE 靠近TVD区上边界，适合保持接触面和方波；MINBEE靠近下界，更耗散但稳健；VANLEER与VANALBADA连续光滑，常在非线性迭代和复杂流场中提供较柔和响应。

图13.15给出高斯波在 $t=1$ 的比较。Lax--Wendroff、SUPERA、VANLEER、MINA 都较好地保持幅值和宽度，Godunov 已明显削峰，ULTRAA 则把圆滑峰顶压成近似平台。这说明“更压缩”并不等于对所有解更准确，限制器必须兼顾光滑极值处的退化行为。

% [原PDF p485]
同名的迎风、WAF和中心限制器只有映射关系，不能把函数值直接互换。实现时应明确当前限制的是通量校正、WAF耗散参数还是重构斜率，并使用对应的比值定义。数值试验表明，限制器性能具有问题依赖性。纯间断测试偏爱压缩函数，光滑波长期传播更需要低色散和温和限制；实际系统通常同时含两类结构，因此应综合多种基准题评估。通量限制器在低阶单调通量与高阶通量之间调节校正量。零值完全关闭高阶校正，一值恢复目标高阶格式；超出一的值属于反扩散，可使间断更锐利，但只能在TVD区域允许的范围内使用。

图13.16把同一高斯推进到 $t=10$。线性格式的相位误差和耗散误差均累积可见；TVD 限制器也不可能完全消除长期误差。VANLEER 在光滑性和耗散之间给出较均衡结果，SUPERA 略显方化，MINA 略显展宽，而 ULTRAA 已严重失真。

% [原PDF p486]
SUPERBEE 靠近TVD区上边界，适合保持接触面和方波；MINBEE靠近下界，更耗散但稳健；VANLEER与VANALBADA连续光滑，常在非线性迭代和复杂流场中提供较柔和响应。中心低阶通量含有下风项，因此迎风格式的经典数据相容条件会过度收紧。扩展中心TVD判据分别处理单调区与局部极值，既保留原中心低阶格式，又允许向Lax--Wendroff加入受控高阶校正。Sweby 图把所有可接受限制器表示在比值—限制系数平面内。上边界决定最大压缩性，下边界保证不产生新极值，二阶点决定光滑区精度；这三个几何条件比逐个推导格式更易复用。

图13.17、13.18对应方波的短时和长时结果。TVD 限制器没有 Lax--Wendroff 的过冲与欠冲，同时比 Godunov 的过渡带窄得多。长时演化把微小的数值黏性差异放大：ULTRAA 最锐利但不适合光滑区，MINA 最稳健但最扩散，SUPERA 与 VANLEER 是通常应用中的折中。

% [原PDF p487]
通量限制器在低阶单调通量与高阶通量之间调节校正量。零值完全关闭高阶校正，一值恢复目标高阶格式；超出一的值属于反扩散，可使间断更锐利，但只能在TVD区域允许的范围内使用。Godunov点表示中心低阶通量经特定比例高阶校正后恰好达到一阶迎风通量。它为中心限制器与迎风限制器建立仿射映射，也清楚标出正耗散和较低耗散区域的分界。同名的迎风、WAF和中心限制器只有映射关系，不能把函数值直接互换。实现时应明确当前限制的是通量校正、WAF耗散参数还是重构斜率，并使用对应的比值定义。

把式(13.130)代入守恒更新，可写成 Harten 增量形式
\begin{equation}
u_i^{n+1}=u_i^n-C\Delta u_{i-1/2}+D\Delta u_{i+1/2},\tag{13.137}
\end{equation}
其中
\begin{equation}
C=c[\alpha_0+(\beta_0-\alpha_0)\psi_{i-1/2}],\quad
D=-c[\alpha_1+(\beta_1-\alpha_1)\psi_{i+1/2}].
\end{equation}
由此可以直接应用 Harten 条件。下面先令低阶通量为 Godunov 迎风、高阶通量为 Lax--Wendroff，并把限制器统一记作 $\psi(r)$。

% [原PDF p488]
Sweby 图把所有可接受限制器表示在比值—限制系数平面内。上边界决定最大压缩性，下边界保证不产生新极值，二阶点决定光滑区精度；这三个几何条件比逐个推导格式更易复用。CFL相关TVD区能利用当前时间步得到更宽的允许范围；CFL无关区域则对所有合法时间步统一安全。前者分辨率可能更高，后者实现简单且不易因局部波速更新滞后而越界。SUPERBEE 靠近TVD区上边界，适合保持接触面和方波；MINBEE靠近下界，更耗散但稳健；VANLEER与VANALBADA连续光滑，常在非线性迭代和复杂流场中提供较柔和响应。

对 $a>0$，$\alpha_0=1,\alpha_1=0$。把更新合并为单一迎风差
\begin{equation}
u_i^{n+1}=u_i^n-\widehat C\Delta u_{i-1/2},qquad
\widehat C=C-D/r,\quad r=\Delta u_{i-1/2}/\Delta u_{i+1/2}.\tag{13.139}
\end{equation}
Harten 条件要求 $0\le\widehat C\le1$。再设所有限制器值落在全局区间 $\psi_B\le\psi\le\psi_T$，可把条件分解成对当前界面函数的上下斜线约束，避免相邻两个未知限制器耦合。

% [原PDF p489]
同名的迎风、WAF和中心限制器只有映射关系，不能把函数值直接互换。实现时应明确当前限制的是通量校正、WAF耗散参数还是重构斜率，并使用对应的比值定义。数值试验表明，限制器性能具有问题依赖性。纯间断测试偏爱压缩函数，光滑波长期传播更需要低色散和温和限制；实际系统通常同时含两类结构，因此应综合多种基准题评估。通量限制器在低阶单调通量与高阶通量之间调节校正量。零值完全关闭高阶校正，一值恢复目标高阶格式；超出一的值属于反扩散，可使间断更锐利，但只能在TVD区域允许的范围内使用。

对 $a<0$，迎风方向在右，写成 $u_i^{n+1}=u_i^n+\widehat D\Delta u_{i+1/2}$，并将 $r$ 仍定义为上风差与局部差之比。用 $|c|$ 统一两种符号后，TVD 条件为
\begin{align}
\psi_B&\le\psi(r)\le\psi_T,\tag{13.152}\\
r>0:&\quad \psi_L(r)\le\psi(r)\le\psi_R(r),\\
r<0:&\quad \psi_R(r)\le\psi(r)\le\psi_L(r),\tag{13.153--13.154}
\end{align}
其中
\begin{equation}
\psi_L=\left(\psi_T-\frac{2}{1-|c|}\right)r,qquad
\psi_R=\left(\psi_B+\frac2{|c|}\right)r.\tag{13.155--13.156}
\end{equation}

% [原PDF p490]
SUPERBEE 靠近TVD区上边界，适合保持接触面和方波；MINBEE靠近下界，更耗散但稳健；VANLEER与VANALBADA连续光滑，常在非线性迭代和复杂流场中提供较柔和响应。中心低阶通量含有下风项，因此迎风格式的经典数据相容条件会过度收紧。扩展中心TVD判据分别处理单调区与局部极值，既保留原中心低阶格式，又允许向Lax--Wendroff加入受控高阶校正。Sweby 图把所有可接受限制器表示在比值—限制系数平面内。上边界决定最大压缩性，下边界保证不产生新极值，二阶点决定光滑区精度；这三个几何条件比逐个推导格式更易复用。

取 $\psi_B=0,\psi_T=2/(1-|c|)$ 时，$r<0$ 的 TVD 区退化为 $\psi=0$，即局部极值附近回到一阶迎风。Sweby 的常用 CFL 无关区域进一步取 $\psi_B=0,\psi_T=2$，并把右边界斜率 $2/|c|$ 收紧为2。于是
\begin{equation}0\le\psi(r)\le\min(2r,2),\qquad r\ge0;\qquad \psi(r)=0, r<0.\tag{13.160a}\end{equation}
限制器应在该阴影区域内，并通过光滑二阶点 $(r,\psi)=(1,1)$。

% [原PDF p491]
通量限制器在低阶单调通量与高阶通量之间调节校正量。零值完全关闭高阶校正，一值恢复目标高阶格式；超出一的值属于反扩散，可使间断更锐利，但只能在TVD区域允许的范围内使用。Godunov点表示中心低阶通量经特定比例高阶校正后恰好达到一阶迎风通量。它为中心限制器与迎风限制器建立仿射映射，也清楚标出正耗散和较低耗散区域的分界。同名的迎风、WAF和中心限制器只有映射关系，不能把函数值直接互换。实现时应明确当前限制的是通量校正、WAF耗散参数还是重构斜率，并使用对应的比值定义。

Sweby 区中的五个常用迎风通量限制器为
\begin{align}
\psi_{\rm SB}(r)&=\max\{0,\min(2r,1),\min(r,2)\},\tag{13.162}\\
\psi_{\rm VL}(r)&=\begin{cases}0,&r\le0,\\2r/(1+r),&r>0,\end{cases}\tag{13.163}\\
\psi_{\rm VA}(r)&=\begin{cases}0,&r\le0,\\r(1+r)/(1+r^2),&r>0,\end{cases}\tag{13.164}\\
\psi_{\rm MB}(r)&=\max[0,\min(1,r)].\tag{13.165}
\end{align}
ULTRABEE 沿 CFL 相关 TVD 区最外边界，压缩性最强；SUPERBEE 是常用分段线性折中；VANLEER、VANALBADA 光滑；MINBEE（MINMOD）最保守。WAF 限制器与普通通量限制器一一对应：
\begin{equation}\phi(r)=1-(1-|c|)\psi(r).\tag{13.166}\end{equation}

% [原PDF p492]
Sweby 图把所有可接受限制器表示在比值—限制系数平面内。上边界决定最大压缩性，下边界保证不产生新极值，二阶点决定光滑区精度；这三个几何条件比逐个推导格式更易复用。CFL相关TVD区能利用当前时间步得到更宽的允许范围；CFL无关区域则对所有合法时间步统一安全。前者分辨率可能更高，后者实现简单且不易因局部波速更新滞后而越界。SUPERBEE 靠近TVD区上边界，适合保持接触面和方波；MINBEE靠近下界，更耗散但稳健；VANLEER与VANALBADA连续光滑，常在非线性迭代和复杂流场中提供较柔和响应。

\subsection{TVD 中心通量限制器}
现在低阶通量取一般中心通量，高阶通量仍为 Lax--Wendroff。由于中心低阶格式含下风贡献，普通迎风数据相容条件可能过严，甚至会把原本单调的低阶中心格式排除。因此需要更一般的 TVD 判据。

定义上风与下风变化之比
\begin{equation}
r_i=\frac{\Delta u_{i-s/2}}{\Delta u_{i+s/2}},\qquad s=\operatorname{sign}(a),\tag{13.168}
\end{equation}
以及归一化更新量
\begin{equation}
R_i^{(s)}=\frac{u_i^{n+1}-u_i^n}{u_{i-s}^n-u_i^n}
=\begin{cases}C-D/r_i,&a>0,\\D-C/r_i,&a<0.
\end{cases}\tag{13.169}
\end{equation}

% [原PDF p493]
同名的迎风、WAF和中心限制器只有映射关系，不能把函数值直接互换。实现时应明确当前限制的是通量校正、WAF耗散参数还是重构斜率，并使用对应的比值定义。数值试验表明，限制器性能具有问题依赖性。纯间断测试偏爱压缩函数，光滑波长期传播更需要低色散和温和限制；实际系统通常同时含两类结构，因此应综合多种基准题评估。通量限制器在低阶单调通量与高阶通量之间调节校正量。零值完全关闭高阶校正，一值恢复目标高阶格式；超出一的值属于反扩散，可使间断更锐利，但只能在TVD区域允许的范围内使用。

一个适用于中心格式的充分条件是：当 $r_i>0$ 时
\begin{equation}-\frac{\epsilon_R}{r_i}\le R_i^{(s)}\le\epsilon_L,\tag{13.170}\end{equation}
当 $r_i<0$ 时
\begin{equation}0\le R_i^{(s)}\le1-\epsilon_0-\frac{\epsilon_0}{r_i},\tag{13.171}\end{equation}
其中 $0\le\epsilon_L,\epsilon_0,\epsilon_R\le1$。第一式处理单调数据，约束新值处于相邻数据的加权包络；第二式处理局部极值，保证极值不被扩大。证明分别对单调增、单调减、局部极小和局部极大四种构型应用离散最大值原则。

% [原PDF p494]
SUPERBEE 靠近TVD区上边界，适合保持接触面和方波；MINBEE靠近下界，更耗散但稳健；VANLEER与VANALBADA连续光滑，常在非线性迭代和复杂流场中提供较柔和响应。中心低阶通量含有下风项，因此迎风格式的经典数据相容条件会过度收紧。扩展中心TVD判据分别处理单调区与局部极值，既保留原中心低阶格式，又允许向Lax--Wendroff加入受控高阶校正。Sweby 图把所有可接受限制器表示在比值—限制系数平面内。上边界决定最大压缩性，下边界保证不产生新极值，二阶点决定光滑区精度；这三个几何条件比逐个推导格式更易复用。

离开极值时，式(13.170)比通常的 $0\le R_i^{(s)}\le1$ 更宽松；在极值处，式(13.171)更严格。由于极值附近 TVD 格式本来就会退化到一阶，组合后的整体条件反而更适合中心格式。

中心限制器的 TVD 区由六条直线围成：全局上下界 $\phi_B,\phi_T$，以及在 $r>0$、$r<0$ 两侧的四条斜线 $\phi_L^+,\phi_R^+,\phi_L^-,\phi_R^-$。这些直线的斜率依赖低阶通量系数、Lax--Wendroff 系数和 Courant 数。对 FORCE、Lax--Friedrichs 与 Godunov 迎风，它们可由统一公式计算。

% [原PDF p495]
通量限制器在低阶单调通量与高阶通量之间调节校正量。零值完全关闭高阶校正，一值恢复目标高阶格式；超出一的值属于反扩散，可使间断更锐利，但只能在TVD区域允许的范围内使用。Godunov点表示中心低阶通量经特定比例高阶校正后恰好达到一阶迎风通量。它为中心限制器与迎风限制器建立仿射映射，也清楚标出正耗散和较低耗散区域的分界。同名的迎风、WAF和中心限制器只有映射关系，不能把函数值直接互换。实现时应明确当前限制的是通量校正、WAF耗散参数还是重构斜率，并使用对应的比值定义。

通常取 $\phi_B=0$，并选择 $\phi_T$ 使左边界斜率非正，从而可令 $\phi_L^+(r)=0$、$\phi_L^-(r)=\phi_g$。这里 $\phi_g$ 是“Godunov 点”：限制器取此值时，中心低阶通量与 Lax--Wendroff 的混合恰好等于 Godunov 一阶迎风通量。

推导中合适的参数为 $\epsilon_L=1$、$\epsilon_0=\epsilon_R=-c\alpha_1$（以 $a>0$ 表示）。若机械使用迎风数据相容条件，相当于令后两者为零，会删去 $\phi<\phi_g$ 的区域，甚至把 $\phi=0$ 对应的单调低阶中心格式排除；这说明中心格式必须使用扩展判据。

% [原PDF p496]
Sweby 图把所有可接受限制器表示在比值—限制系数平面内。上边界决定最大压缩性，下边界保证不产生新极值，二阶点决定光滑区精度；这三个几何条件比逐个推导格式更易复用。CFL相关TVD区能利用当前时间步得到更宽的允许范围；CFL无关区域则对所有合法时间步统一安全。前者分辨率可能更高，后者实现简单且不易因局部波速更新滞后而越界。SUPERBEE 靠近TVD区上边界，适合保持接触面和方波；MINBEE靠近下界，更耗散但稳健；VANLEER与VANALBADA连续光滑，常在非线性迭代和复杂流场中提供较柔和响应。

中心 TVD 通量为
\begin{equation}
f_{i+1/2}^{(c)}=f_{i+1/2}^{\rm LO}+\phi_{i+1/2}(r)
[f_{i+1/2}^{\rm LW}-f_{i+1/2}^{\rm LO}].\tag{13.186}
\end{equation}
当 $\phi=\phi_g$ 时，它等于 Godunov 迎风通量。中心限制器 $\phi$ 与普通迎风限制器 $\psi$ 的关系为
\begin{equation}\phi(r)=\phi_g+(1-\phi_g)\psi(r).\tag{13.189}\end{equation}
因此可借用成熟的迎风限制器形状构造中心限制器，但二者并不等价。对非线性标量律，可用局部特征速度定义 $c_{i+1/2}=\Delta t|a_{i+1/2}|/\Delta x$，无需真正求解黎曼问题也能构造类似迎风的中心通量。

% [原PDF p497]
同名的迎风、WAF和中心限制器只有映射关系，不能把函数值直接互换。实现时应明确当前限制的是通量校正、WAF耗散参数还是重构斜率，并使用对应的比值定义。数值试验表明，限制器性能具有问题依赖性。纯间断测试偏爱压缩函数，光滑波长期传播更需要低色散和温和限制；实际系统通常同时含两类结构，因此应综合多种基准题评估。通量限制器在低阶单调通量与高阶通量之间调节校正量。零值完全关闭高阶校正，一值恢复目标高阶格式；超出一的值属于反扩散，可使间断更锐利，但只能在TVD区域允许的范围内使用。

因 TVD 区随 $|c|$ 改变，实际实现可由 $c_{\min},c_{\max}$ 取最严格包络。若按模型方程取 $c_{\min}=0,c_{\max}=1$，得到 CFL 无关区域，其中最直接的限制器只有 MINBEE 型
\begin{equation}\phi_{\rm MB}(r)=\begin{cases}0,&r\le0,\\r,&0<r<1,\\1,&r\ge1.
\end{cases}\tag{13.193}\end{equation}
它稳健但偏扩散。为保留 $\phi>1$ 的较低耗散区域，可令
\begin{equation}\phi(r)=\widehat\phi_g+(1-\widehat\phi_g)\psi(r),\quad
\widehat\phi_g=0\ (r\le1),\quad \widehat\phi_g=\phi_g\ (r>1).\tag{13.194--13.195}
\end{equation}

% [原PDF p498]
SUPERBEE 靠近TVD区上边界，适合保持接触面和方波；MINBEE靠近下界，更耗散但稳健；VANLEER与VANALBADA连续光滑，常在非线性迭代和复杂流场中提供较柔和响应。中心低阶通量含有下风项，因此迎风格式的经典数据相容条件会过度收紧。扩展中心TVD判据分别处理单调区与局部极值，既保留原中心低阶格式，又允许向Lax--Wendroff加入受控高阶校正。Sweby 图把所有可接受限制器表示在比值—限制系数平面内。上边界决定最大压缩性，下边界保证不产生新极值，二阶点决定光滑区精度；这三个几何条件比逐个推导格式更易复用。

选择 ULTRABEE、SUPERBEE、VANLEER 或 VANALBADA 的 $\psi$，便得到对应的中心限制器。例如
\begin{equation}
\phi_{\rm VL}(r)=\begin{cases}
0,&r\le0,\\2r/(1+r),&0<r\le1,\\
\phi_g+2(1-\phi_g)r/(1+r),&r>1.
\end{cases}\tag{13.198}
\end{equation}
VANALBADA 只需把 $2r/(1+r)$ 换为 $r(1+r)/(1+r^2)$。这些中心限制器只是与同名迎风限制器相似，不能在程序中混用其公式。

\section{斜率限制器方法}
MUSCL 的未限制斜率会在大梯度处振荡。斜率限制器把 $\Delta_i$ 替换为非线性的 $\bar\Delta_i$，使演化后的左右外推值满足 TVD 包络。

% [原PDF p499]
斜率限制直接约束重构几何。两侧差分同号时保留不超过许可上界的斜率，异号时把斜率置零，这使单元边界外推值不会越出相邻平均值包络。应限制半步演化后的外推值，而非只限制初始重构值，因为时间预测会按通量差再次移动两端状态。忽略这一步可能使初始重构看似单调，预测后却重新越界。迎风斜率限制可使用特征传播方向定义比值，中心格式没有天然的单一上风方向，通常采用对称而保守的系数。非线性系统还需决定在守恒变量、原始变量或特征变量中实施限制。

直接限制初始外推值并不充分；应限制半步演化后的值：
\begin{equation}
\min(u_{i-1}^n,u_i^n)\le\bar u_i^L\le\max(u_{i-1}^n,u_i^n),
\quad
\min(u_i^n,u_{i+1}^n)\le\bar u_i^R\le\max(u_i^n,u_{i+1}^n).\tag{13.201}
\end{equation}
若随后使用任意单调三点低阶格式，新值必落在 $u_{i-2}^n,ldots,u_{i+2}^n$ 的局部最小--最大包络内，因而所得 MUSCL 格式为 TVD。这一结论同时适用于迎风和中心低阶通量。

% [原PDF p500]
应限制半步演化后的外推值，而非只限制初始重构值，因为时间预测会按通量差再次移动两端状态。忽略这一步可能使初始重构看似单调，预测后却重新越界。迎风斜率限制可使用特征传播方向定义比值，中心格式没有天然的单一上风方向，通常采用对称而保守的系数。非线性系统还需决定在守恒变量、原始变量或特征变量中实施限制。光滑极值处严格TVD方法会把斜率降为零，因此形式阶数可能局部下降。实际高阶方法有时改用TVB、ENO或WENO以减轻这一退化，但必须接受更复杂的参数或光滑性判据。

满足上述约束的一种限制斜率为
\begin{equation}
\bar\Delta_i=\frac12[\operatorname{sign}(\Delta u_{i-1/2})+
\operatorname{sign}(\Delta u_{i+1/2})]
\min\left[\frac{2|\Delta u_{i-1/2}|}{1+c},
\frac{2|\Delta u_{i+1/2}|}{1-c}\right].\tag{13.207}
\end{equation}
相邻差同号时，选择两侧允许幅值中较小者；异号时表明存在局部极值，符号和为零，故斜率自动置零。这正是避免新极值的核心机制。

% [原PDF p501]
迎风斜率限制可使用特征传播方向定义比值，中心格式没有天然的单一上风方向，通常采用对称而保守的系数。非线性系统还需决定在守恒变量、原始变量或特征变量中实施限制。光滑极值处严格TVD方法会把斜率降为零，因此形式阶数可能局部下降。实际高阶方法有时改用TVB、ENO或WENO以减轻这一退化，但必须接受更复杂的参数或光滑性判据。斜率限制直接约束重构几何。两侧差分同号时保留不超过许可上界的斜率，异号时把斜率置零，这使单元边界外推值不会越出相邻平均值包络。

把原始斜率写为
\begin{equation}\Delta_i=\tfrac12(1+\omega)\Delta u_{i-1/2}+
\tfrac12(1-\omega)\Delta u_{i+1/2},\qquad
\bar\Delta_i=\xi_i\Delta_i,\tag{13.213--13.214}\end{equation}
可得到斜率限制器 TVD 区。令 $r=\Delta u_{i-1/2}/\Delta u_{i+1/2}$，则 $r\le0$ 时必须 $\xi=0$；$r>0$ 时
\begin{equation}0\le\xi(r)\le\min[\xi_L(r),\xi_R(r)],\tag{13.215}\end{equation}
其中
\begin{equation}
\xi_L=\frac{2\beta_{i-1/2}r}{1-\omega+(1+\omega)r},\quad
\xi_R=\frac{2\beta_{i+1/2}}{1-\omega+(1+\omega)r}.\tag{13.216}
\end{equation}

% [原PDF p502]
光滑极值处严格TVD方法会把斜率降为零，因此形式阶数可能局部下降。实际高阶方法有时改用TVB、ENO或WENO以减轻这一退化，但必须接受更复杂的参数或光滑性判据。斜率限制直接约束重构几何。两侧差分同号时保留不超过许可上界的斜率，异号时把斜率置零，这使单元边界外推值不会越出相邻平均值包络。应限制半步演化后的外推值，而非只限制初始重构值，因为时间预测会按通量差再次移动两端状态。忽略这一步可能使初始重构看似单调，预测后却重新越界。

沿 TVD 区边界取值相当于 ULTRABEE，压缩性强；SUPERBEE 型在 $0<r<1/2$ 取 $2r$，$1/2<r<1$ 取1，$r>1$ 再由 $r,2,\xi_R$ 的最小值限制。VANLEER、VANALBADA 与 MINBEE 型分别为
\begin{align}
\xi_{\rm VL}&=\min[2r/(1+r),\xi_R],\\
\xi_{\rm VA}&=\min[r(1+r)/(1+r^2),\xi_R],\\
\xi_{\rm MB}&=\begin{cases}r,&0<r<1,\\\min(1,\xi_R),&r\ge1,
\end{cases}
\end{align}
并均在 $r\le0$ 时取零。迎风方法可利用 $c$ 中的传播方向；中心方法缺少这类方向信息，推荐保守取 $\beta_{i-1/2}=\beta_{i+1/2}=1$，得到对全部 $|c|\le1$ 有效但较小的 TVD 区。

% [原PDF p503]
斜率限制直接约束重构几何。两侧差分同号时保留不超过许可上界的斜率，异号时把斜率置零，这使单元边界外推值不会越出相邻平均值包络。应限制半步演化后的外推值，而非只限制初始重构值，因为时间预测会按通量差再次移动两端状态。忽略这一步可能使初始重构看似单调，预测后却重新越界。迎风斜率限制可使用特征传播方向定义比值，中心格式没有天然的单一上风方向，通常采用对称而保守的系数。非线性系统还需决定在守恒变量、原始变量或特征变量中实施限制。

斜率限制与通量限制虽可构造形状相似的函数，但二者作用位置不同，不能简单等同。迎风 MUSCL 常用受限斜率可写为
\begin{equation}
\bar\Delta_i=\psi(r_i)\Delta u_{i+1/2}\quad(a>0),
\qquad
\bar\Delta_i=\psi(r_i)\Delta u_{i-1/2}\quad(a<0),\tag{13.229}
\end{equation}
其中 $\psi$ 可取 MINBEE、VANLEER、VANALBADA 或 SUPERBEE；例如
\begin{equation}\psi_{\rm SB}(r)=\max[0,\min(2r,1),\min(r,2)].\tag{13.231}\end{equation}
在线性平流下，这类受限 MUSCL 的未限制基准二阶格式是 Lax--Wendroff，而不是 Fromm。实现时应明确斜率定义和 $r$ 的方向约定，否则同名限制器会产生不同结果。

% [原PDF p504]
应限制半步演化后的外推值，而非只限制初始重构值，因为时间预测会按通量差再次移动两端状态。忽略这一步可能使初始重构看似单调，预测后却重新越界。迎风斜率限制可使用特征传播方向定义比值，中心格式没有天然的单一上风方向，通常采用对称而保守的系数。非线性系统还需决定在守恒变量、原始变量或特征变量中实施限制。光滑极值处严格TVD方法会把斜率降为零，因此形式阶数可能局部下降。实际高阶方法有时改用TVB、ENO或WENO以减轻这一退化，但必须接受更复杂的参数或光滑性判据。

\section{TVD 方法的扩展}
\subsection{含源项的方程}
对
\begin{equation}u_t+au_x=s(u),\tag{13.232}\end{equation}
可采用算子分裂：齐次对流部分用 TVD 格式，常微分源项部分用合适的 ODE 求解器。某些特殊问题可通过变量变换写成新的齐次方程，再直接应用 TVD。必须认识到，一般源项本身可能使精确解的总变差增加，因此强求完整非齐次离散解总变差不增并不物理；这种情况下应控制各分裂子问题的稳定性和守恒性质，而不是机械套用 TVD 定义。

\subsection{含扩散项的方程}
对对流--扩散方程 $u_t+au_x=\alpha_{\rm phy}u_{xx}$，可把对流与扩散分裂计算；这一策略在 Navier--Stokes 求解中很常见。

% [原PDF p505]
迎风斜率限制可使用特征传播方向定义比值，中心格式没有天然的单一上风方向，通常采用对称而保守的系数。非线性系统还需决定在守恒变量、原始变量或特征变量中实施限制。光滑极值处严格TVD方法会把斜率降为零，因此形式阶数可能局部下降。实际高阶方法有时改用TVB、ENO或WENO以减轻这一退化，但必须接受更复杂的参数或光滑性判据。斜率限制直接约束重构几何。两侧差分同号时保留不超过许可上界的斜率，异号时把斜率置零，这使单元边界外推值不会越出相邻平均值包络。

TVD 对流离散等效地加入数值黏性 $\alpha_{\rm num}u_{xx}$。即使原方程已有物理黏性，也未必足以消除离散振荡。定义扩散数 $d=\alpha_{\rm phy}\Delta t/\Delta x^2$ 和单元 Reynolds 数 $R_{\rm cell}=|c|/d$，当
\begin{equation}R_{\rm cell}\le\frac{2}{1-|c|}\tag{13.235}\end{equation}
时，某些 TVD 离散中物理黏性本身足够；否则还需附加数值黏性。黏性限制器的理想作用是充分利用已有物理扩散，只补入保证无振荡所需的最小人工黏性，并在条件(13.235)满足时把它完全关闭。

\section{线性平流数值结果}
再次采用 $a=1,C_{\rm CFL}=0.8$，将高斯与方波推进到 $t=10$，共1250步。高斯取 $\alpha=1,\beta=8,\Delta x=0.02$；方波取 $\Delta x=0.01$；两者均用100个单元，并以精确平移解为参照。

% [原PDF p506]
光滑极值处严格TVD方法会把斜率降为零，因此形式阶数可能局部下降。实际高阶方法有时改用TVB、ENO或WENO以减轻这一退化，但必须接受更复杂的参数或光滑性判据。斜率限制直接约束重构几何。两侧差分同号时保留不超过许可上界的斜率，异号时把斜率置零，这使单元边界外推值不会越出相邻平均值包络。应限制半步演化后的外推值，而非只限制初始重构值，因为时间预测会按通量差再次移动两端状态。忽略这一步可能使初始重构看似单调，预测后却重新越界。

比较两类 TVD 方法：通量限制器与斜率限制器。每类又含一个基于 Godunov 的迎风版本和一个基于 FORCE 的中心版本。中心通量限制器称 FLIC，中心斜率限制器称 SLIC，重构参数取 $\omega=0$。分别测试 SUPERBEE、VANLEER、VANALBADA、MINBEE；注意中心限制器只与同名迎风限制器相似，并非同一个函数。

对这个标量问题，中心 TVD 方法与迎风 TVD 方法总体相当，未显示明显劣势；然而后续非线性系统算例将表明，迎风方法通常更准确。中心格式的优势是无需黎曼问题细节，程序更简单。

% [原PDF p507]
斜率限制直接约束重构几何。两侧差分同号时保留不超过许可上界的斜率，异号时把斜率置零，这使单元边界外推值不会越出相邻平均值包络。应限制半步演化后的外推值，而非只限制初始重构值，因为时间预测会按通量差再次移动两端状态。忽略这一步可能使初始重构看似单调，预测后却重新越界。迎风斜率限制可使用特征传播方向定义比值，中心格式没有天然的单一上风方向，通常采用对称而保守的系数。非线性系统还需决定在守恒变量、原始变量或特征变量中实施限制。

图13.25比较光滑高斯在 $t=10$ 时的通量限制迎风与 FLIC--FORCE。SUPERBEE 保持轮廓较锐但有轻微平台化；VANLEER 和 VANALBADA 给出较光滑峰顶；MINBEE 的峰值最低、宽度最大。左右两列的整体差异小于不同限制器之间的差异，说明标量常速问题中限制器选择是主导因素。

% [原PDF p508]
应限制半步演化后的外推值，而非只限制初始重构值，因为时间预测会按通量差再次移动两端状态。忽略这一步可能使初始重构看似单调，预测后却重新越界。迎风斜率限制可使用特征传播方向定义比值，中心格式没有天然的单一上风方向，通常采用对称而保守的系数。非线性系统还需决定在守恒变量、原始变量或特征变量中实施限制。光滑极值处严格TVD方法会把斜率降为零，因此形式阶数可能局部下降。实际高阶方法有时改用TVB、ENO或WENO以减轻这一退化，但必须接受更复杂的参数或光滑性判据。

图13.26比较方波的通量限制结果。四种方法都避免过冲与欠冲；SUPERBEE 的两侧过渡最窄，MINBEE 最宽，VANLEER 与 VANALBADA 居中。中心 FORCE 版本并未因缺少显式迎风黎曼解而失去 TVD 性质，在这一简单问题上可获得与迎风版本相近的间断分辨率。

% [原PDF p509]
迎风斜率限制可使用特征传播方向定义比值，中心格式没有天然的单一上风方向，通常采用对称而保守的系数。非线性系统还需决定在守恒变量、原始变量或特征变量中实施限制。光滑极值处严格TVD方法会把斜率降为零，因此形式阶数可能局部下降。实际高阶方法有时改用TVB、ENO或WENO以减轻这一退化，但必须接受更复杂的参数或光滑性判据。斜率限制直接约束重构几何。两侧差分同号时保留不超过许可上界的斜率，异号时把斜率置零，这使单元边界外推值不会越出相邻平均值包络。

图13.27给出斜率限制 MUSCL--Hancock 对光滑测试的结果。迎风与中心列均显示 SUPERBEE 较压缩、MINBEE 较耗散的规律；VANLEER、VANALBADA 在光滑峰值附近更自然。与通量限制法相比，差别主要来自限制器如何作用于重构斜率，而不是基本更新公式。

% [原PDF p510]
光滑极值处严格TVD方法会把斜率降为零，因此形式阶数可能局部下降。实际高阶方法有时改用TVB、ENO或WENO以减轻这一退化，但必须接受更复杂的参数或光滑性判据。斜率限制直接约束重构几何。两侧差分同号时保留不超过许可上界的斜率，异号时把斜率置零，这使单元边界外推值不会越出相邻平均值包络。应限制半步演化后的外推值，而非只限制初始重构值，因为时间预测会按通量差再次移动两端状态。忽略这一步可能使初始重构看似单调，预测后却重新越界。

图13.28给出方波的斜率限制结果。所有方案均保持无振荡，间断过渡宽度与同名通量限制器大体一致。数值比较再次强调：不存在对所有解都最优的限制器。强压缩限制器适合清晰间断，却可能损坏光滑极值；保守限制器鲁棒，但长期演化会累积扩散。工程计算通常在 VANLEER、VANALBADA 或 SUPERBEE 中结合问题特征选择。

% [原PDF p511]
\chapter{非线性系统的高阶与 TVD 格式}
从标量方程推广到系统时，单一波速被完整特征谱取代。时间步必须覆盖所有单元和边界幽灵状态中的最大信号速度；边界填充若产生更快状态，也必须在CFL估计中体现。守恒变量重构直接保持守恒结构，原始变量重构常给出更平滑的密度和压力剖面，特征变量重构则更能分离不同波族。后续格式将在精度、成本和鲁棒性之间比较这些选择。二阶边界需要两层幽灵单元，是因为重构与限制器会访问更宽模板。透射边界复制内部状态，反射壁反转相对壁面的法向速度，同时保持密度和压力镜像。

本章研究依赖时间 $t$ 与一维空间 $x$ 的非线性守恒律系统。迎风 TVD 格式是第6章 Godunov 一阶方法的高阶扩展，可与第4章精确黎曼求解器、第9--12章近似求解器以及第8章通量矢量分裂结合。中心 TVD 格式则扩展第7章 FORCE 方法。它们汇总了前面各章的工具，特别是第13章在标量问题中建立的 TVD 理论。

\section{引言}
以下以理想气体一维非定常欧拉方程说明算法，但推广到其他双曲系统并不困难。含反应等源项的系统见第15章，多维系统见第16章。迎风 TVD 通常精度较高，代价是黎曼求解器带来的复杂度；对理想气体，采用高效求解器后额外成本通常不大。中心 TVD 更易实现、略省计算，尤其适合黎曼结构未知、过于复杂或求解代价过高的系统。

% [原PDF p512]
守恒变量重构直接保持守恒结构，原始变量重构常给出更平滑的密度和压力剖面，特征变量重构则更能分离不同波族。后续格式将在精度、成本和鲁棒性之间比较这些选择。二阶边界需要两层幽灵单元，是因为重构与限制器会访问更宽模板。透射边界复制内部状态，反射壁反转相对壁面的法向速度，同时保持密度和压力镜像。WAF系统通量对整个黎曼扇作加权平均，权重由各波在半时间层的位置决定。若使用近似黎曼求解器，必须保证所提供的波速和中间状态足以重建这一分段结构。

初边值问题写成
\begin{equation}
\bm U_t+\bm F(\bm U)_x=0,\quad
\bm U(x,0)=\bm U^{(0)}(x),\quad
\bm U(0,t)=\bm U_l(t),\quad\bm U(L,t)=\bm U_r(t).\tag{14.1}
\end{equation}
对欧拉方程
\begin{equation}
\bm U=\begin{bmatrix}\rho\\\rho u\\E\end{bmatrix},\qquad
\bm F=\begin{bmatrix}\rho u\\\rho u^2+p\\u(E+p)\end{bmatrix}.\tag{14.2}
\end{equation}
统一的全离散显式守恒更新为
\begin{equation}\bm U_i^{n+1}=\bm U_i^n+\frac{\Delta t}{\Delta x}
(\bm F_{i-1/2}-\bm F_{i+1/2}).\tag{14.3}\end{equation}
本章将给出 WAF、MUSCL--Hancock 两种迎风 TVD，FLIC、SLIC 两种中心 TVD，以及基于原始变量和自适应原始--守恒变量的变体。

% [原PDF p513]
二阶边界需要两层幽灵单元，是因为重构与限制器会访问更宽模板。透射边界复制内部状态，反射壁反转相对壁面的法向速度，同时保持密度和压力镜像。WAF系统通量对整个黎曼扇作加权平均，权重由各波在半时间层的位置决定。若使用近似黎曼求解器，必须保证所提供的波速和中间状态足以重建这一分段结构。从标量方程推广到系统时，单一波速被完整特征谱取代。时间步必须覆盖所有单元和边界幽灵状态中的最大信号速度；边界填充若产生更快状态，也必须在CFL估计中体现。

\section{CFL 条件与边界条件}
所有方法的线性化稳定性要求 $|c|\le1$。非线性系统采用
\begin{equation}\Delta t=C_{\rm CFL}\frac{\Delta x}{S_{\max}^{(n)}},\qquad0<C_{\rm CFL}\le1,\tag{14.4}\end{equation}
其中 $S_{\max}^{(n)}$ 是当前时间层的最大波速。对欧拉方程可取
\begin{equation}S_{\max}^{(n)}=\max_i(|u_i^n|+a_i^n),\tag{14.5}\end{equation}
并须把边界条件产生的幽灵单元状态纳入最大值。若波速估计漏掉真正的快速波，即使 $C_{\rm CFL}$ 很小也可能失稳。

二阶格式在每个边界需要两个幽灵单元。透射边界按内部状态零梯度外推：左侧 $\bm W_0=\bm W_1,\bm W_{-1}=\bm W_2$；右侧 $\bm W_{M+1}=\bm W_M,\bm W_{M+2}=\bm W_{M-1}$。$\bm W$ 可取守恒变量或原始变量，但应与格式重构变量保持一致。

% [原PDF p514]
WAF系统通量对整个黎曼扇作加权平均，权重由各波在半时间层的位置决定。若使用近似黎曼求解器，必须保证所提供的波速和中间状态足以重建这一分段结构。从标量方程推广到系统时，单一波速被完整特征谱取代。时间步必须覆盖所有单元和边界幽灵状态中的最大信号速度；边界填充若产生更快状态，也必须在CFL估计中体现。守恒变量重构直接保持守恒结构，原始变量重构常给出更平滑的密度和压力剖面，特征变量重构则更能分离不同波族。后续格式将在精度、成本和鲁棒性之间比较这些选择。

对位于 $x=L$、速度为 $u_{\rm wall}$ 的反射壁，右侧两层幽灵单元满足
\begin{align}
\rho_{M+1}&=\rho_M,&u_{M+1}&=-u_M+2u_{\rm wall},&p_{M+1}&=p_M,\\
\rho_{M+2}&=\rho_{M-1},&u_{M+2}&=-u_{M-1}+2u_{\rm wall},&p_{M+2}&=p_{M-1}.
\tag{14.7}
\end{align}
左反射壁完全类似。确定时间步与边界填充后，格式只由界面通量决定。

\section{加权平均通量格式}
WAF 是 Godunov 一阶迎风法的确定性二阶扩展。原始版本把局部黎曼解的物理通量在半时间层上作空间平均：
\begin{equation}
\bm F_{i+1/2}=\frac1{\Delta x}\int_{-\Delta x/2}^{\Delta x/2}
\bm F[\bm U_{i+1/2}(x,\Delta t/2)]\,dx.\tag{14.8}
\end{equation}
局部黎曼解以 $(\bm U_i^n,\bm U_{i+1}^n)$ 为数据，可由精确或近似求解器给出。欧拉解通常有三条波，分隔四个常状态 $\bm U_i^n,\bm U_{*L},\bm U_{*R},\bm U_{i+1}^n$；若声波为稀疏波，扇内还含非均匀自相似状态。后续将把这些状态和波速转化为可计算的加权通量。
% ===== 译文分片 6（已合并）=====

% [原PDF p515]
因此，公式中的每个速度、跳跃和限制系数都应与同一局部黎曼问题对应；这种逐波一致性是获得正确激波速度和接触分辨率的前提。

WAF 构造的几何意义是：在一个时间步内，从界面黎曼扇中截取长度为 $\Delta x$ 的区间，再对其中的物理通量作平均。每个权重都由波速与时间步共同确定，因此当某条波移出该区间时，相应常状态的贡献自然消失。由此得到的并非任意加权，而是局部时空控制体上的守恒平均。

实现时应把各波按从左到右的速度排序，并统一保存波前后的状态、通量跳跃和 Courant 数。若波速次序、状态编号或界面朝向不一致，公式虽然仍能计算，却会把错误的跳跃分配给某一波族。对强激波，还须确认黎曼解给出的密度和压力保持物理可容许。

在欧拉黎曼解中，三条波把半时间层积分区间分成四个常状态区。若第 $k$ 区状态为 $\bm U^{(k)}$，则 WAF 通量为
\begin{equation}
\bm F_{i+1/2}=\sum_{k=1}^{N+1}\beta_k\bm F^{(k)}_{i+1/2},\qquad
\beta_k=\tfrac12(c_k-c_{k-1}),\tag{14.10--14.11}
\end{equation}
其中 $c_k=S_k\Delta t/\Delta x$，并约定 $c_0=-1,c_{N+1}=1$。权重就是各状态区在长度 $\Delta x$ 内的归一化占比，故其和为一。

% [原PDF p516]
从有限体积角度看，所有高阶操作最终都只能改变界面通量的构造，而不能改变相邻单元共享该通量的守恒更新关系。

通量平均型和状态平均型 WAF 在光滑区具有相同的二阶目标，但对非线性系统并不严格等价。前者直接组合各常状态的物理通量，后者先组合状态再调用非线性通量函数；由于 $\bm F(\sum\beta_k\bm W^{(k)})\ne\sum\beta_k\bm F(\bm W^{(k)})$，两者的耗散和鲁棒性可能不同。

稀疏波包含连续扇区，不能简单当成一条跨越两个常状态的间断。非跨声速扇可用邻近星区状态近似其贡献；若扇跨过 $x/t=0$，则必须采用声速点状态，使界面通量包含正确的熵满足解。这一处理也是避免产生非物理解压激波的关键。

消去权重后可写成波跳跃形式
\begin{equation}
\bm F_{i+1/2}=\tfrac12(\bm F_i+\bm F_{i+1})-
\tfrac12\sum_{k=1}^{N}c_k\Delta\bm F^{(k)}_{i+1/2},\quad
\Delta\bm F^{(k)}=\bm F^{(k+1)}-\bm F^{(k)}.\tag{14.12--14.13}
\end{equation}
另一版本先平均状态
\begin{equation}\bar{\bm W}_{i+1/2}=\tfrac12(\bm W_i+\bm W_{i+1})-
\tfrac12\sum_k c_k(\bm W^{(k+1)}-\bm W^{(k)}),\tag{14.16}\end{equation}
再取 $\bm F_{i+1/2}=\bm F(\bar{\bm W}_{i+1/2})$。推荐 $\bm W=(\rho,u,p)^T$，这样每个界面只需一次物理通量求值。

% [原PDF p517]
数值结果应同时检查波的位置、剖面宽度、极值和物理状态；单看一条曲线的平滑程度，不能判断格式是否正确。

TVD 修正按波族工作，而不是给整个欧拉通量只乘一个公共系数。密度或内能的逐波跳跃用来衡量上风侧与当前界面的平滑程度；平滑区限制器允许较高分辨率，局部极值或强间断附近则增大耗散。多维计算中，接触波携带的切向速度变化也应单独限制。

限制比的分子必须取在该波传播方向的上风界面：正波速读取左侧相邻界面，负波速读取右侧相邻界面。分母很小时要用明确的极限定义，而不应让数值噪声决定限制器。这样才能保证左右传播波的处理对称，并维持守恒界面通量的唯一性。

\subsection{状态型黎曼求解器中的稀疏波}
精确求解器和近似状态求解器给出用于计算物理通量的状态；若局部解含稀疏扇，WAF 空间平均还须包括扇内连续状态。欧拉方程的第1、3波可能为激波或稀疏波，第2波始终是速度 $u_*$ 的接触间断。以右稀疏波为例，尾部速度 $S_{TR}=u_*+a_{*R}$，头部速度 $S_{HR}=u_{i+1}+a_{i+1}$，理论上可对扇内解析解精确积分。

% [原PDF p518]
变量选择会影响限制器看到的波动幅度，却不应改变状态方程的定义；每次变量转换后都需保持密度、动量、能量和压力相容。

一个完整的 WAF 时间步由边界填充、全局稳定时间步、界面黎曼求解、逐波限制和守恒更新组成。黎曼求解器不仅返回中心通量，还要返回按波序排列的中间状态与速度；这些量随后被相邻界面的平滑度计算共同使用，所以数据布局必须保持一致。

在初始强间断附近，简单的最大波速估计可能低估尚未形成的星区声速。短暂减小 CFL 数能够避免第一步过大，但这只是时间离散的谨慎选择，不能替代对状态合法性和波速估计的检查。更新后应重新验证 $\rho>0$、$p>0$ 以及总能量与内能的一致性。

实践中可把稀疏扇与靠近时间轴的常状态合并。非跨声速右稀疏波用星形右状态，跨声速时用黎曼解在 $x/t=0$ 的扇内状态；左稀疏波同理。近似式为
\begin{equation}
\frac1{\Delta x}\int_{A_2}^{A_3}\bm F\,dx\approx\tfrac12(c_3-c_2)\bm F^{(3)}.\tag{14.19}
\end{equation}
跨声速时出现的 $\bm F[\bm U_{i+1/2}(0)]$ 正是 Godunov 通量，因此它始终是 WAF 通量的一部分。

% [原PDF p519]
因此，公式中的每个速度、跳跃和限制系数都应与同一局部黎曼问题对应；这种逐波一致性是获得正确激波速度和接触分辨率的前提。

WAF 构造的几何意义是：在一个时间步内，从界面黎曼扇中截取长度为 $\Delta x$ 的区间，再对其中的物理通量作平均。每个权重都由波速与时间步共同确定，因此当某条波移出该区间时，相应常状态的贡献自然消失。由此得到的并非任意加权，而是局部时空控制体上的守恒平均。

实现时应把各波按从左到右的速度排序，并统一保存波前后的状态、通量跳跃和 Courant 数。若波速次序、状态编号或界面朝向不一致，公式虽然仍能计算，却会把错误的跳跃分配给某一波族。对强激波，还须确认黎曼解给出的密度和压力保持物理可容许。

HLL、HLLC 以及通量矢量分裂已把稀疏结构表示成少数波与状态，不需要上述额外处理；第9章近似状态求解器则仍需区分跨声速扇。

\subsection{WAF 的 TVD 版本}
非线性系统的 TVD 推广带有经验性：标量方程只有一个波，而欧拉系统有三个波族。做法是选择一个跨每个波都会变化的标量 $q$，对每个波分别构造限制器。TVD--WAF 通量为
\begin{equation}
\bm F_{i+1/2}=\tfrac12(\bm F_i+\bm F_{i+1})-
\tfrac12\sum_k\operatorname{sign}(c_k)\phi^{(k)}(r^{(k)})\Delta\bm F^{(k)}.\tag{14.23}
\end{equation}

% [原PDF p520]
从有限体积角度看，所有高阶操作最终都只能改变界面通量的构造，而不能改变相邻单元共享该通量的守恒更新关系。

通量平均型和状态平均型 WAF 在光滑区具有相同的二阶目标，但对非线性系统并不严格等价。前者直接组合各常状态的物理通量，后者先组合状态再调用非线性通量函数；由于 $\bm F(\sum\beta_k\bm W^{(k)})\ne\sum\beta_k\bm F(\bm W^{(k)})$，两者的耗散和鲁棒性可能不同。

稀疏波包含连续扇区，不能简单当成一条跨越两个常状态的间断。非跨声速扇可用邻近星区状态近似其贡献；若扇跨过 $x/t=0$，则必须采用声速点状态，使界面通量包含正确的熵满足解。这一处理也是避免产生非物理解压激波的关键。

波族比值按传播方向取上风跳跃与局部跳跃之比：
\begin{equation}
r^{(k)}=\begin{cases}
\Delta q^{(k)}_{i-1/2}/\Delta q^{(k)}_{i+1/2},&c_k>0,\\
\Delta q^{(k)}_{i+3/2}/\Delta q^{(k)}_{i+1/2},&c_k<0.
\end{cases}\tag{14.25}
\end{equation}
欧拉方程通常取 $q=\rho$ 或单位质量内能 $e$；多维问题还必须对切向速度跳跃限制。状态平均版本用同样的 $\operatorname{sign}(c_k)\phi_k$ 替代裸 Courant 数，再对所得平均原始量计算通量。

% [原PDF p521]
数值结果应同时检查波的位置、剖面宽度、极值和物理状态；单看一条曲线的平滑程度，不能判断格式是否正确。

TVD 修正按波族工作，而不是给整个欧拉通量只乘一个公共系数。密度或内能的逐波跳跃用来衡量上风侧与当前界面的平滑程度；平滑区限制器允许较高分辨率，局部极值或强间断附近则增大耗散。多维计算中，接触波携带的切向速度变化也应单独限制。

限制比的分子必须取在该波传播方向的上风界面：正波速读取左侧相邻界面，负波速读取右侧相邻界面。分母很小时要用明确的极限定义，而不应让数值噪声决定限制器。这样才能保证左右传播波的处理对称，并维持守恒界面通量的唯一性。

可采用与 SUPERBEE、VANLEER、VANALBADA、MINBEE 对应的 WAF 限制器，二者关系为
\begin{equation}\phi(r)=1-(1-|c|)\psi(r).\tag{14.30a}\end{equation}
推荐的黎曼求解器包括精确解、自适应近似状态解和 HLLC；WAF 也可推广到非均匀网格与 Courant 数2的大时间步版本。

\subsection{WAF 算法小结}
每个时间步先填充边界，再由最大波速计算 $\Delta t$。对激波管型强间断，简化波速估计初期可能偏小，宜在头几步把 $C_{\rm CFL}$ 暂降至约0.2。随后求各界面黎曼问题，保存状态、物理通量、波速以及所选变量的逐波跳跃。

% [原PDF p522]
变量选择会影响限制器看到的波动幅度，却不应改变状态方程的定义；每次变量转换后都需保持密度、动量、能量和压力相容。

一个完整的 WAF 时间步由边界填充、全局稳定时间步、界面黎曼求解、逐波限制和守恒更新组成。黎曼求解器不仅返回中心通量，还要返回按波序排列的中间状态与速度；这些量随后被相邻界面的平滑度计算共同使用，所以数据布局必须保持一致。

在初始强间断附近，简单的最大波速估计可能低估尚未形成的星区声速。短暂减小 CFL 数能够避免第一步过大，但这只是时间离散的谨慎选择，不能替代对状态合法性和波速估计的检查。更新后应重新验证 $\rho>0$、$p>0$ 以及总能量与内能的一致性。

对每个界面、每个波族计算 $r_k,c_k,\phi_k$，由式(14.23)或状态平均形式得到唯一界面通量；再用守恒式(14.3)更新所有单元，进入下一时间层。界面黎曼信息必须在相邻界面之间一致组织，才能正确取得上风跳跃。

\section{MUSCL--Hancock 格式}
MUSCL--Hancock 通过分段线性重构、半步预测和普通黎曼问题，把 Godunov 一阶迎风法提升为全离散二阶方法。对一般系统
\begin{equation}\bm U_t+\bm F(\bm U)_x=0\tag{14.31}\end{equation}
算法仍沿用第13章标量形式，但斜率、外推值和通量均为矢量。

% [原PDF p523]
因此，公式中的每个速度、跳跃和限制系数都应与同一局部黎曼问题对应；这种逐波一致性是获得正确激波速度和接触分辨率的前提。

MUSCL--Hancock 方法首先在每个单元中重构一个分段线性剖面，使左右端值的平均仍等于原单元平均，从而不破坏守恒性。随后用单元内部左右端通量差进行半步预测，把空间重构和时间演化耦合起来。最终界面黎曼问题看到的是半时间层的左右状态，因此整体达到二阶精度。

斜率限制是该方法处理间断的核心。若左右差分异号，局部可能包含极值或间断，受限斜率应收缩至零；若同号，则在中心斜率与单边斜率之间选取满足 TVD 条件的值。限制过弱会在激波附近产生振荡，限制过强则会把接触间断和光滑波峰显著抹宽。

第一步在单元内重构
\begin{equation}
\bm U_i(x)=\bm U_i^n+\frac{x-x_i}{\Delta x}\bm\Delta_i,\qquad
\bm U_i^L=\bm U_i^n-\tfrac12\bm\Delta_i,\quad
\bm U_i^R=\bm U_i^n+\tfrac12\bm\Delta_i.\tag{14.32--14.33}
\end{equation}
第二步把两端值在单元内演化半步：
\begin{equation}
\bar{\bm U}_i^{L,R}=\bm U_i^{L,R}+\frac{\Delta t}{2\Delta x}
[\bm F(\bm U_i^L)-\bm F(\bm U_i^R)].\tag{14.34}
\end{equation}

% [原PDF p524]
从有限体积角度看，所有高阶操作最终都只能改变界面通量的构造，而不能改变相邻单元共享该通量的守恒更新关系。

守恒变量重构直接对应有限体积平均，形式最自然；原始变量重构通常更容易保持压力剖面的规则性；特征变量重构则先把扰动分解到三个波族，能够避免不同物理波相互污染。三种选择在光滑区均可二阶，但面对强激波、近真空或物质界面时表现可能显著不同。

半步预测后必须检查左右外推状态是否物理可容许。出现负密度或负压力说明当前斜率、网格分辨率或时间步不适合，不能仅在通量计算阶段静默截断。可靠实现会报告单元位置、预测状态及所用限制器，并通过收紧重构或明确减小时间步来恢复有效状态。

第三步以 $(\bar{\bm U}_i^R,\bar{\bm U}_{i+1}^L)$ 为数据求常规黎曼问题，并取
\begin{equation}\bm F_{i+1/2}=\bm F[\bm U_{i+1/2}(0)].\tag{14.36}\end{equation}
未限制斜率可取
\begin{equation}\bm\Delta_i=\tfrac12(1+\omega)(\bm U_i-\bm U_{i-1})+
\tfrac12(1-\omega)(\bm U_{i+1}-\bm U_i),\quad -1\le\omega\le1.\tag{14.37}\end{equation}
不同黎曼波型的处理与一阶 Godunov 方法相同，差别仅在左右数据已经过高阶预测。

% [原PDF p525]
数值结果应同时检查波的位置、剖面宽度、极值和物理状态；单看一条曲线的平滑程度，不能判断格式是否正确。

特征限制在局部平均状态上构造左右特征矩阵，将相邻差分投影为声学波、接触波等独立幅值。逐波限制后再映射回守恒变量，既可对激波施加强耗散，又能对线性退化波保留较高分辨率。代价是每个单元都要稳定地计算特征基及其逆变换。

算法的最终守恒性来自每个界面只计算一次通量，并由左右单元以相反符号共享。重构与预测都可以采用不同变量，但送入黎曼求解器前必须转换成同一种完整物理状态。边界幽灵单元也要参与同样的重构流程，否则靠近边界的空间阶数会无意降低。

也可在原始变量 $\bm W=(\rho,u,p)^T$ 中重构和预测。若 $\bm W_t+\bm A(\bm W)\bm W_x=0$，则
\begin{align}
\bar{\bm W}_i^L&=\bm W_i^n-\tfrac12[\bm I+(\Delta t/\Delta x)\bm A_i]\bm\Delta_i,\\
\bar{\bm W}_i^R&=\bm W_i^n+\tfrac12[\bm I-(\Delta t/\Delta x)\bm A_i]\bm\Delta_i.\tag{14.41}
\end{align}
欧拉原始变量矩阵为
\begin{equation}\bm A=\begin{bmatrix}u&\rho&0\\0&u&1/\rho\\0&\rho a^2&u\end{bmatrix}.\tag{14.43}\end{equation}

% [原PDF p526]
变量选择会影响限制器看到的波动幅度，却不应改变状态方程的定义；每次变量转换后都需保持密度、动量、能量和压力相容。

MUSCL--Hancock 方法首先在每个单元中重构一个分段线性剖面，使左右端值的平均仍等于原单元平均，从而不破坏守恒性。随后用单元内部左右端通量差进行半步预测，把空间重构和时间演化耦合起来。最终界面黎曼问题看到的是半时间层的左右状态，因此整体达到二阶精度。

斜率限制是该方法处理间断的核心。若左右差分异号，局部可能包含极值或间断，受限斜率应收缩至零；若同号，则在中心斜率与单边斜率之间选取满足 TVD 条件的值。限制过弱会在激波附近产生振荡，限制过强则会把接触间断和光滑波峰显著抹宽。

\subsection{MUSCL--Hancock 的 TVD 限制}
对系统可直接逐分量限制守恒量或原始量，也可先投影到局部特征变量。常见做法令每个分量的受限斜率由左右差分的同号最小值确定；左右差异号时斜率为零。直接强制方程给出的限制斜率简单稳健，但系统级 TVD 仍是经验推广。

在强激波问题中，原始变量限制通常避免由守恒量线性组合引起的压力振荡；特征限制能更清楚地区分声波、接触波和剪切波，但需在每个单元计算左右特征矢量。

% [原PDF p527]
因此，公式中的每个速度、跳跃和限制系数都应与同一局部黎曼问题对应；这种逐波一致性是获得正确激波速度和接触分辨率的前提。

守恒变量重构直接对应有限体积平均，形式最自然；原始变量重构通常更容易保持压力剖面的规则性；特征变量重构则先把扰动分解到三个波族，能够避免不同物理波相互污染。三种选择在光滑区均可二阶，但面对强激波、近真空或物质界面时表现可能显著不同。

半步预测后必须检查左右外推状态是否物理可容许。出现负密度或负压力说明当前斜率、网格分辨率或时间步不适合，不能仅在通量计算阶段静默截断。可靠实现会报告单元位置、预测状态及所用限制器，并通过收紧重构或明确减小时间步来恢复有效状态。

特征限制先在平均状态处构造左特征矩阵 $\bm L$，把差分投影为 $\Delta\bm w=\bm L\Delta\bm U$，逐特征分量应用 MINBEE、VANLEER、VANALBADA 或 SUPERBEE，再用右特征矩阵还原。欧拉三波分别对应 $u-a,u,u+a$，因此限制器可按不同波族的局部平滑性独立调节。

% [原PDF p528]
从有限体积角度看，所有高阶操作最终都只能改变界面通量的构造，而不能改变相邻单元共享该通量的守恒更新关系。

特征限制在局部平均状态上构造左右特征矩阵，将相邻差分投影为声学波、接触波等独立幅值。逐波限制后再映射回守恒变量，既可对激波施加强耗散，又能对线性退化波保留较高分辨率。代价是每个单元都要稳定地计算特征基及其逆变换。

算法的最终守恒性来自每个界面只计算一次通量，并由左右单元以相反符号共享。重构与预测都可以采用不同变量，但送入黎曼求解器前必须转换成同一种完整物理状态。边界幽灵单元也要参与同样的重构流程，否则靠近边界的空间阶数会无意降低。

若希望消除 Courant 数对限制边界的依赖，可取对称保守系数，从而得到对全部 $|c|\le1$ 有效的较小 TVD 区。算法流程仍是：边界填充、时间步计算、受限重构、半步预测、界面黎曼求解、守恒更新。任何一步产生非物理密度或压力都应被视为重构/时间步失效，而非静默截断。

\section{中心 TVD 格式}
中心方法不求黎曼问题，常以 FORCE 为低阶通量，并通过通量限制或斜率限制获得二阶。它们更易推广到复杂系统，但通常比迎风 TVD 稍耗散。

% [原PDF p529]
数值结果应同时检查波的位置、剖面宽度、极值和物理状态；单看一条曲线的平滑程度，不能判断格式是否正确。

中心高分辨率格式不显式追踪波族，而以局部双状态预测构造界面通量。FORCE 通量把 Lax--Friedrichs 的强耗散与 Richtmyer 通量的较低耗散平均起来，因而无需特征分解也能给出稳定的一阶基准。高阶修正只应在局部解足够光滑时释放。

FLIC 在 FORCE 与 Richtmyer 两种通量之间进行受限反扩散，SLIC 则先限制单元斜率、预测左右界面状态，再用 FORCE 计算通量。两者的限制位置不同，但都要保持单一共享界面通量。中心格式的限制器区域与迎风格式不同，不能把迎风公式不加区分地直接移植。

每步先按式(14.4)--(14.5)计算时间步并设置边界。中心格式同样需要对强初始间断谨慎估计最大波速。随后分别构造 FLIC 通量或 SLIC 重构，最后通过同一个守恒更新式推进。

% [原PDF p530]
变量选择会影响限制器看到的波动幅度，却不应改变状态方程的定义；每次变量转换后都需保持密度、动量、能量和压力相容。

原始变量方程是守恒律的非保守改写。在光滑解中，两种形式通过链式法则等价；跨越激波时，乘积和导数的弱解解释不再唯一，非守恒离散通常不能满足正确的 Rankine--Hugoniot 跳跃条件。因此纯原始变量更新适合光滑声波、接触界面和被动输运，而不应独立承担激波传播。

原始变量方法的优势在于它直接演化密度、速度、压力和材料标记，可减轻多介质接触面上的压力振荡。为了同时获得正确激波速度，可用探测器在受压缩且压力跳跃显著的区域切换到守恒格式，在其余区域继续使用原始变量格式；切换必须保持界面通量与更新变量的一致。

\subsection{FORCE 通量回顾}
对系统 $\bm U_t+\bm F_x=0$，Lax--Friedrichs 与 Richtmyer 两步通量分别为
\begin{align}
\bm U_{i+1/2}^{RI}&=\tfrac12(\bm U_i+\bm U_{i+1})-
\frac{\Delta t}{2\Delta x}(\bm F_{i+1}-\bm F_i),\\
\bm F_{i+1/2}^{RI}&=\bm F(\bm U_{i+1/2}^{RI}),\\
\bm F_{i+1/2}^{FORCE}&=\tfrac12(\bm F_{i+1/2}^{LF}+\bm F_{i+1/2}^{RI}).
\end{align}
FORCE 是单调一阶中心通量，不需要波分解。

% [原PDF p531]
因此，公式中的每个速度、跳跃和限制系数都应与同一局部黎曼问题对应；这种逐波一致性是获得正确激波速度和接触分辨率的前提。

WAF 型原始变量方法仍以局部黎曼扇的半时间平均作为界面状态，只是更新项由系数矩阵 $\bm A$ 乘状态差给出。若只在界面中心采样而忽略扇区占比，线性平流上会退回一阶；对左右黎曼问题的影响作对称平均，才能恢复预期的时空二阶精度。

MUSCL--Hancock 型原始变量方法同样包含受限重构和半步预测。矩阵预测式清楚展示了各特征速度对左右外推值的作用。实际程序中，应先确定变量顺序和状态方程，再统一实现 $\bm W\leftrightarrow\bm U$ 转换，避免压力、总能量和被动标量在不同步骤使用不一致的定义。

FLIC（Flux-Limiter Centred）以 FORCE 为低阶、Richtmyer 为高阶：
\begin{equation}
\bm F_{i+1/2}^{FLIC}=\bm F_{i+1/2}^{FORCE}+
\phi_{i+1/2}(\bm F_{i+1/2}^{RI}-\bm F_{i+1/2}^{FORCE}).\tag{14.64}
\end{equation}
限制器由某个代表性变量的上风/局部变化比得到。因系统有多个波，严格 TVD 证明不再成立，但这一经验构造在欧拉算例中表现可靠。

% [原PDF p532]
从有限体积角度看，所有高阶操作最终都只能改变界面通量的构造，而不能改变相邻单元共享该通量的守恒更新关系。

自适应守恒—原始变量算法的关键是激波判据。一个有效判据通常同时考察速度散度的压缩性、压力相对跳跃和局部声速尺度；单纯以密度梯度判定，会把平滑材料界面误认作激波。标记后的区域还应适当扩展若干单元，覆盖高阶模板和激波移动的距离。

格式切换不能让相邻单元各自计算不同的界面作用。应先为每个界面选定唯一更新机制，再把同一数值通量或等效界面贡献施加给两侧单元。这样才能保留全局守恒；若在原始变量区采用非守恒更新，还应独立监测质量、动量和总能量误差。

欧拉方程建议用总能量 $q=E$ 构造左右比值，并以局部最大特征速度决定上风方向。中心限制器需使用第13章专门推导的中心版本，不能直接照搬迎风限制器。SLIC 则先用中心型受限斜率重构、半步预测，再把两侧预测值送入 FORCE 通量。

% [原PDF p533]
数值结果应同时检查波的位置、剖面宽度、极值和物理状态；单看一条曲线的平滑程度，不能判断格式是否正确。

中心高分辨率格式不显式追踪波族，而以局部双状态预测构造界面通量。FORCE 通量把 Lax--Friedrichs 的强耗散与 Richtmyer 通量的较低耗散平均起来，因而无需特征分解也能给出稳定的一阶基准。高阶修正只应在局部解足够光滑时释放。

FLIC 在 FORCE 与 Richtmyer 两种通量之间进行受限反扩散，SLIC 则先限制单元斜率、预测左右界面状态，再用 FORCE 计算通量。两者的限制位置不同，但都要保持单一共享界面通量。中心格式的限制器区域与迎风格式不同，不能把迎风公式不加区分地直接移植。

对 SLIC，界面左右输入为 $\bm W_L=\bar{\bm W}_i^R$、$\bm W_R=\bar{\bm W}_{i+1}^L$，计算 $\bm F_{i+1/2}^{FORCE}(\bm W_L,\bm W_R)$。中心斜率应取对称限制系数，逐原始量或守恒量处理；使用迎风特征斜率公式一般不能保证中心格式的 TVD 性质。

\section{原始变量格式}
\subsection{方程形式与适用范围}
原始变量格式基于非守恒形式
\begin{equation}\bm W_t+\bm A(\bm W)\bm W_x=0.\tag{14.77}\end{equation}
它在含激波时会给出错误激波强度与传播速度，因而不能单独用于有激波的流动；但无激波、物质界面、剪切波、声学弱波及缺少已知守恒形式的系统中可能更合适。

% [原PDF p534]
变量选择会影响限制器看到的波动幅度，却不应改变状态方程的定义；每次变量转换后都需保持密度、动量、能量和压力相容。

原始变量方程是守恒律的非保守改写。在光滑解中，两种形式通过链式法则等价；跨越激波时，乘积和导数的弱解解释不再唯一，非守恒离散通常不能满足正确的 Rankine--Hugoniot 跳跃条件。因此纯原始变量更新适合光滑声波、接触界面和被动输运，而不应独立承担激波传播。

原始变量方法的优势在于它直接演化密度、速度、压力和材料标记，可减轻多介质接触面上的压力振荡。为了同时获得正确激波速度，可用探测器在受压缩且压力跳跃显著的区域切换到守恒格式，在其余区域继续使用原始变量格式；切换必须保持界面通量与更新变量的一致。

某些守恒格式在多组分物质界面会产生压力振荡，而原始变量格式能避免这类误差。可通过自适应算法仅在激波附近切换到守恒更新，在其他区域保留原始变量优势。

对增广一维欧拉方程取 $\bm W=(\rho,u,p,q)^T$，其中 $q$ 是被动标量，则
\begin{equation}
\bm A=\begin{bmatrix}
u&\rho&0&0\\0&u&1/\rho&0\\0&\rho a^2&u&0\\0&0&0&u
\end{bmatrix}.\tag{14.78}
\end{equation}
$q$ 可表示反应进度、组分浓度、物质界面标记或多维气动中的切向速度。

% [原PDF p535]
因此，公式中的每个速度、跳跃和限制系数都应与同一局部黎曼问题对应；这种逐波一致性是获得正确激波速度和接触分辨率的前提。

WAF 型原始变量方法仍以局部黎曼扇的半时间平均作为界面状态，只是更新项由系数矩阵 $\bm A$ 乘状态差给出。若只在界面中心采样而忽略扇区占比，线性平流上会退回一阶；对左右黎曼问题的影响作对称平均，才能恢复预期的时空二阶精度。

MUSCL--Hancock 型原始变量方法同样包含受限重构和半步预测。矩阵预测式清楚展示了各特征速度对左右外推值的作用。实际程序中，应先确定变量顺序和状态方程，再统一实现 $\bm W\leftrightarrow\bm U$ 转换，避免压力、总能量和被动标量在不同步骤使用不一致的定义。

在交错网格上采用两步差分格式
\begin{equation}
\bm W_i^{n+1}=\bm W_i^n+\frac{\Delta t}{\Delta x}\bm A_i
(\bm W_{i-1/2}^{n+1/2}-\bm W_{i+1/2}^{n+1/2}).\tag{14.79}
\end{equation}
关键是确定半时间层界面状态与点 $i$ 的系数矩阵。它们分别对应守恒方法中的界面通量和局部线性化状态。

% [原PDF p536]
从有限体积角度看，所有高阶操作最终都只能改变界面通量的构造，而不能改变相邻单元共享该通量的守恒更新关系。

自适应守恒—原始变量算法的关键是激波判据。一个有效判据通常同时考察速度散度的压缩性、压力相对跳跃和局部声速尺度；单纯以密度梯度判定，会把平滑材料界面误认作激波。标记后的区域还应适当扩展若干单元，覆盖高阶模板和激波移动的距离。

格式切换不能让相邻单元各自计算不同的界面作用。应先为每个界面选定唯一更新机制，再把同一数值通量或等效界面贡献施加给两侧单元。这样才能保留全局守恒；若在原始变量区采用非守恒更新，还应独立监测质量、动量和总能量误差。

\subsection{WAF 型原始变量格式}
半时间界面状态定义为黎曼解的空间平均
\begin{equation}
\bm W_{i+1/2}^{n+1/2}=\sum_{k=1}^{N+1}\beta_k\bm W_{i+1/2}^{(k)}
=\tfrac12(\bm W_i+\bm W_{i+1})-\tfrac12\sum_kc_k\Delta\bm W^{(k)}.\tag{14.81--14.83}
\end{equation}
系数矩阵取 $\bm A_i=\bm A(\bar{\bm W}_i)$。简单取 $\bar{\bm W}_i=\bm W_i^n$ 鲁棒性不足，更好的做法是平均左右两个黎曼问题在半时间层对单元中心的影响。

% [原PDF p537]
数值结果应同时检查波的位置、剖面宽度、极值和物理状态；单看一条曲线的平滑程度，不能判断格式是否正确。

中心高分辨率格式不显式追踪波族，而以局部双状态预测构造界面通量。FORCE 通量把 Lax--Friedrichs 的强耗散与 Richtmyer 通量的较低耗散平均起来，因而无需特征分解也能给出稳定的一阶基准。高阶修正只应在局部解足够光滑时释放。

FLIC 在 FORCE 与 Richtmyer 两种通量之间进行受限反扩散，SLIC 则先限制单元斜率、预测左右界面状态，再用 FORCE 计算通量。两者的限制位置不同，但都要保持单一共享界面通量。中心格式的限制器区域与迎风格式不同，不能把迎风公式不加区分地直接移植。

精细平均可对左右黎曼扇各取半区积分；实用近似为
\begin{equation}\bar{\bm W}_i=\tfrac12(\bm W_{i-1/2}^{n+1/2}+\bm W_{i+1/2}^{n+1/2}).\tag{14.88}\end{equation}
在线性平流上，完整状态平均恢复二阶 Lax--Wendroff；若对界面状态只用中点采样，则退化成一阶迎风。TVD 版本用 $\operatorname{sign}(c_k)\phi_k$ 替换 $c_k$，与守恒 WAF 完全平行。

% [原PDF p538]
变量选择会影响限制器看到的波动幅度，却不应改变状态方程的定义；每次变量转换后都需保持密度、动量、能量和压力相容。

原始变量方程是守恒律的非保守改写。在光滑解中，两种形式通过链式法则等价；跨越激波时，乘积和导数的弱解解释不再唯一，非守恒离散通常不能满足正确的 Rankine--Hugoniot 跳跃条件。因此纯原始变量更新适合光滑声波、接触界面和被动输运，而不应独立承担激波传播。

原始变量方法的优势在于它直接演化密度、速度、压力和材料标记，可减轻多介质接触面上的压力振荡。为了同时获得正确激波速度，可用探测器在受压缩且压力跳跃显著的区域切换到守恒格式，在其余区域继续使用原始变量格式；切换必须保持界面通量与更新变量的一致。

算法先求各界面原始变量黎曼问题，计算受限半时间状态；再由左右状态平均构造 $\bm A_i$；最后应用式(14.79)。若改为计算 $\bm F_{i+1/2}=\bm F(\bm W_{i+1/2}^{n+1/2})$ 并用守恒更新，就立即得到守恒 WAF 版本。

\subsection{MUSCL--Hancock 原始变量格式}
先重构 $\bm W_i^{L,R}=\bm W_i^n\mp\bm\Delta_i/2$，再用
\begin{equation}\bar{\bm W}_i^{L,R}=\bm W_i^{L,R}+\frac{\Delta t}{2\Delta x}
\bm A(\bm W_i^n)(\bm W_i^L-\bm W_i^R)\tag{14.92}\end{equation}
半步演化，最后对 $(\bar{\bm W}_i^R,\bar{\bm W}_{i+1}^L)$ 求黎曼问题并取 $x/t=0$ 状态。

% [原PDF p539]
因此，公式中的每个速度、跳跃和限制系数都应与同一局部黎曼问题对应；这种逐波一致性是获得正确激波速度和接触分辨率的前提。

WAF 型原始变量方法仍以局部黎曼扇的半时间平均作为界面状态，只是更新项由系数矩阵 $\bm A$ 乘状态差给出。若只在界面中心采样而忽略扇区占比，线性平流上会退回一阶；对左右黎曼问题的影响作对称平均，才能恢复预期的时空二阶精度。

MUSCL--Hancock 型原始变量方法同样包含受限重构和半步预测。矩阵预测式清楚展示了各特征速度对左右外推值的作用。实际程序中，应先确定变量顺序和状态方程，再统一实现 $\bm W\leftrightarrow\bm U$ 转换，避免压力、总能量和被动标量在不同步骤使用不一致的定义。

矩阵状态可取
\begin{equation}\bar{\bm W}_i=\tfrac12(\bar{\bm W}_i^L+\bar{\bm W}_i^R)
=\bm W_i^n-\frac{\Delta t}{2\Delta x}\bm A_i\bm\Delta_i,\tag{14.94--14.95}\end{equation}
也可平均相邻两个界面半步状态。在线性平流上，此法为二阶 Fromm；斜率全零时退化为一阶原始变量迎风格式。用第14.4节受限斜率即可构造 TVD 版本；若改以半步状态计算物理通量并守恒更新，则恢复守恒 MUSCL--Hancock。

% [原PDF p540]
从有限体积角度看，所有高阶操作最终都只能改变界面通量的构造，而不能改变相邻单元共享该通量的守恒更新关系。

自适应守恒—原始变量算法的关键是激波判据。一个有效判据通常同时考察速度散度的压缩性、压力相对跳跃和局部声速尺度；单纯以密度梯度判定，会把平滑材料界面误认作激波。标记后的区域还应适当扩展若干单元，覆盖高阶模板和激波移动的距离。

格式切换不能让相邻单元各自计算不同的界面作用。应先为每个界面选定唯一更新机制，再把同一数值通量或等效界面贡献施加给两侧单元。这样才能保留全局守恒；若在原始变量区采用非守恒更新，还应独立监测质量、动量和总能量误差。

\subsection{自适应原始量--守恒量格式}
基本原则是：只有当相邻黎曼问题中存在朝向当前网格点传播的激波时，才用守恒方法；其余位置用原始变量方法。令 $S_{str}=1+P_{TOL}$，若
\begin{equation}
\frac{p^*_{i-1/2}}{p_i^n}>S_{str},\ S_{i-1/2}>0,
\quad\text{或}\quad
\frac{p^*_{i+1/2}}{p_i^n}>S_{str},\ S_{i+1/2}<0,\tag{14.98}
\end{equation}
则点 $i$ 用守恒更新，否则用原始量更新。$P_{TOL}\in(0,0.1)$ 通常不敏感；过大可能漏掉中等激波，过小则几乎处处使用守恒格式。

% [原PDF p541]
数值结果应同时检查波的位置、剖面宽度、极值和物理状态；单看一条曲线的平滑程度，不能判断格式是否正确。

激波管算例同时包含左行稀疏波、接触间断和右行激波，是比较高分辨率格式的基本测试。数值结果应与精确黎曼解在同一时刻对齐，并分别检查密度、速度、压力和内能。只观察密度曲线不足以发现激波速度误差、接触面压力振荡或稀疏扇中的熵错误。

不同限制器体现分辨率与非振荡性的权衡。较压缩的限制器能用更少网格点捕捉接触间断，但更容易在光滑极值附近形成平台或折角；较耗散的限制器使曲线平顺，却会扩大激波和接触层。评价时既要看过冲，也要看间断厚度与波位置。

\section{数值结果}
采用修正 Sod 问题：$\bm W_L=(1,0,1)^T$、$\bm W_R=(0.125,0,0.1)^T$，初始间断 $x_0=0.3$，$\gamma=1.4$，区间 $[0,1]$，100单元，$C_{\rm CFL}=0.9$，透射边界。精确解含右激波、右行接触面和左侧跨声速稀疏波。比较 HLLC--WAF、HLLC--MUSCL--Hancock、FLIC 与 SLIC，并测试多种限制器。

SUPERBEE 的间断最锐利但可能出现小振荡，MINBEE 无明显振荡却扩散接触面，VANLEER 与 VANALBADA 通常给出居中结果。

% [原PDF p542]
变量选择会影响限制器看到的波动幅度，却不应改变状态方程的定义；每次变量转换后都需保持密度、动量、能量和压力相容。

WAF 与 MUSCL--Hancock 都可达到二阶，但其误差来源不同：前者直接平均黎曼扇并逐波限制，后者在单元内重构后进行半步预测。对于相同黎曼求解器和 CFL 数，二者在光滑区可能非常接近，而接触波、跨声速稀疏波及强激波附近的耗散分布会不同。

收敛判断应排除间断点附近不可避免的一阶局部误差，并在光滑区域计算范数。对于含激波的整体误差，$L_1$ 收敛率通常低于纯光滑问题；这并不意味着格式失去守恒。还应记录时间步、边界位置和精确解采样方式，避免把对齐误差误判为数值色散。

MUSCL--Hancock 使用 SUPERBEE 时接触和激波分辨率高，但内能曲线可见振荡；MINBEE 基本消除振荡，同时扩大过渡区。用直接受限斜率与用斜率限制函数所得结果还会受变量选择和黎曼求解器影响，因此应在密度、速度、压力、内能四个分量上共同检查。

% [原PDF p543]
因此，公式中的每个速度、跳跃和限制系数都应与同一局部黎曼问题对应；这种逐波一致性是获得正确激波速度和接触分辨率的前提。

图中的局部振荡若只出现在初始间断后的头几个时间步，首先应核对 CFL 估计与边界填充；若长期附着在激波后方，则更可能与限制器、黎曼波速或变量重构有关。接触间断上的压力应近似常数，任何明显压力波纹都提示能量与状态方程处理不一致。

强稀疏问题还要检查密度和压力是否接近真空。高阶外推可能在单元平均仍为正时产生非法界面值，因此仅检查更新后的单元平均并不充分。对比算法时必须采用相同网格、终止时刻、CFL 数和边界条件，才能把差别归因于格式本身。

图14.9、14.10比较 WAF 的 SUPERBEE 与 MINBEE：前者在接触面附近只跨少量网格点，后者则平滑得多。两者都能正确给出稀疏波头尾、接触传播速度和激波位置，差别主要是局部耗散量。

% [原PDF p544]
从有限体积角度看，所有高阶操作最终都只能改变界面通量的构造，而不能改变相邻单元共享该通量的守恒更新关系。

激波管算例同时包含左行稀疏波、接触间断和右行激波，是比较高分辨率格式的基本测试。数值结果应与精确黎曼解在同一时刻对齐，并分别检查密度、速度、压力和内能。只观察密度曲线不足以发现激波速度误差、接触面压力振荡或稀疏扇中的熵错误。

不同限制器体现分辨率与非振荡性的权衡。较压缩的限制器能用更少网格点捕捉接触间断，但更容易在光滑极值附近形成平台或折角；较耗散的限制器使曲线平顺，却会扩大激波和接触层。评价时既要看过冲，也要看间断厚度与波位置。

图14.11、14.12给出 MUSCL--Hancock 的两种直接受限斜率结果。SUPERBEE 类限制在间断处保留更大斜率，因而锐利但可能越过系统级严格TVD边界；MINBEE 类把更多单元降为一阶，曲线单调但接触扩散明显。

% [原PDF p545]
数值结果应同时检查波的位置、剖面宽度、极值和物理状态；单看一条曲线的平滑程度，不能判断格式是否正确。

WAF 与 MUSCL--Hancock 都可达到二阶，但其误差来源不同：前者直接平均黎曼扇并逐波限制，后者在单元内重构后进行半步预测。对于相同黎曼求解器和 CFL 数，二者在光滑区可能非常接近，而接触波、跨声速稀疏波及强激波附近的耗散分布会不同。

收敛判断应排除间断点附近不可避免的一阶局部误差，并在光滑区域计算范数。对于含激波的整体误差，$L_1$ 收敛率通常低于纯光滑问题；这并不意味着格式失去守恒。还应记录时间步、边界位置和精确解采样方式，避免把对齐误差误判为数值色散。

图14.13、14.14使用限制器函数形式构造斜率。VANLEER 与 VANALBADA 在光滑稀疏扇中较连续，在激波处自动减小斜率；与直接 minmod 型公式相比，它们通常能改善光滑区相位和幅值误差。

% [原PDF p546]
变量选择会影响限制器看到的波动幅度，却不应改变状态方程的定义；每次变量转换后都需保持密度、动量、能量和压力相容。

图中的局部振荡若只出现在初始间断后的头几个时间步，首先应核对 CFL 估计与边界填充；若长期附着在激波后方，则更可能与限制器、黎曼波速或变量重构有关。接触间断上的压力应近似常数，任何明显压力波纹都提示能量与状态方程处理不一致。

强稀疏问题还要检查密度和压力是否接近真空。高阶外推可能在单元平均仍为正时产生非法界面值，因此仅检查更新后的单元平均并不充分。对比算法时必须采用相同网格、终止时刻、CFL 数和边界条件，才能把差别归因于格式本身。

图14.15、14.16展示 FLIC 中心通量限制结果。无需黎曼求解器仍能捕捉正确波序，但接触面通常比 HLLC 迎风方法更宽。中心 SUPERBEE 可提高锐度，代价是某些变量上的轻微波纹。

% [原PDF p547]
因此，公式中的每个速度、跳跃和限制系数都应与同一局部黎曼问题对应；这种逐波一致性是获得正确激波速度和接触分辨率的前提。

激波管算例同时包含左行稀疏波、接触间断和右行激波，是比较高分辨率格式的基本测试。数值结果应与精确黎曼解在同一时刻对齐，并分别检查密度、速度、压力和内能。只观察密度曲线不足以发现激波速度误差、接触面压力振荡或稀疏扇中的熵错误。

不同限制器体现分辨率与非振荡性的权衡。较压缩的限制器能用更少网格点捕捉接触间断，但更容易在光滑极值附近形成平台或折角；较耗散的限制器使曲线平顺，却会扩大激波和接触层。评价时既要看过冲，也要看间断厚度与波位置。

图14.17、14.18展示 SLIC 中心斜率限制结果。总体趋势与 FLIC 相同：MINBEE 最耗散，VANLEER 或 SUPERBEE 更锐利。对这一标量式激波管，迎风方法精度略高；中心格式的主要优势仍是实现简单与无需详细特征分解。

% [原PDF p548]
\chapter{带源项偏微分方程的分裂格式}
在程序验证中，可关闭输运或关闭源项分别测试两个子算子，再用有解析解的耦合模型测量分裂误差和整体时间阶。

含源项的守恒方程可写成齐次输运算子与局部源算子的组合。Lie 分裂先完整执行一个算子再执行另一个，时间精度通常为一阶；Strang 分裂采用半步源项、整步输运、半步源项的对称次序，在算子足够光滑时达到二阶。两算子不对易造成的交换子误差决定了分裂精度。

分裂后每个子问题都应使用与其自身稳定性相符的求解器。输运步受 CFL 条件限制，源项步则可能由化学、松弛或几何时间尺度控制。把两者机械地使用同一显式步长可能导致刚性失稳，即使齐次有限体积更新本身完全稳定。

\section{引言}
考虑非齐次双曲守恒律
\begin{equation}\bm U_t+\bm F(\bm U)_x=\bm S(\bm U).\tag{15.1}\end{equation}
齐次部分 $\bm U_t+\bm F_x=0$ 描述输运与波传播；忽略空间变化则得到源项常微分方程 $d\bm U/dt=\bm S(\bm U)$。几何降维、变底浅水、反应流、外力和能量交换都会产生源项。

% [原PDF p549]
若源项含守恒交换，应把各组分的增减以成对方式离散，并在每个子步后检查总量，避免 ODE 容差积累成系统漂移。

精确模型问题把平流与指数增长或衰减结合起来，可明确分离空间离散误差、时间积分误差和算子分裂误差。由于常系数平流与线性反应在某些情形下对易，分裂可能表现得异常准确；这类特例不能代替对一般非线性源项的检验。

实现 Strang 分裂时，连续两个主时间步交界处的半源步可在某些自治问题中合并，但输出时刻、可变步长和边界更新会改变这一优化的成立条件。最清楚的实现是显式保留半步—整步—半步结构，并记录每个子步采用的实际时间间隔。

反应气体往往同时含极快化学时间尺度和较慢流体时间尺度，形成刚性 ODE。隐式 ODE 求解虽可缓解刚性稳定限制，但若时空网格不能解析反应区，耦合系统仍可能出现以错误速度传播的伪波。

处理非齐次系统有两条路线：在单一格式中耦合通量与源项，例如源项迎风或预测--校正；或把一时间步拆成齐次 PDE 与局部 ODE，分别调用最合适的成熟求解器。分裂易实现且可达到高阶，但必须关注算子不交换误差、刚性和守恒平衡。

% [原PDF p550]
刚性判据来自源项 Jacobian 的特征尺度；它描述稳定性困难，不能仅由源项数值大小或单个变量的瞬时变化判断。

源项常微分方程若具有相差悬殊的特征时间，就属于刚性问题。显式欧拉的稳定区很小，必须取远小于流动 CFL 步长的子步；隐式方法或适合刚性系统的多步、Runge--Kutta 方法虽然每步更昂贵，却允许沿慢时间尺度推进。稳定性与精度要求应分别评估。

源项求解器的误差容限需要与有限体积空间误差协调。把 ODE 解得远比网格误差精确只会浪费计算量；容限过松则会污染守恒量和波速。对放热反应等强耦合问题，还应在源步后重算温度、声速和状态方程，并检查更新状态处于物理可容许域。

\section{模型方程的分裂}
模型问题
\begin{equation}u_t+au_x=\lambda u,qquad u(x,0)=u_0(x)\tag{15.4--15.5}\end{equation}
的精确解为
\begin{equation}u(x,t)=u_0(x-at)e^{\lambda t}.\tag{15.6}\end{equation}
特征线满足 $dx/dt=a$，沿特征则有 $du/dt=\lambda u$。这已经提示可先把初值平移，再在每一点乘以源项增长因子。

% [原PDF p551]
分裂次序、边界刷新和输出时刻都是离散算法的一部分，复现实验时应与通量格式和 ODE 求解器一起明确记录。

算子分裂便于复用成熟的齐次黎曼求解器和独立 ODE 模块，但也可能破坏稳态平衡。例如重力与压力梯度若被分开处理，会在原本静止的平衡解上产生伪速度。需要保持良好平衡时，应把相互抵消的通量和源项按同一离散关系设计，而不能仅依靠提高 Strang 分裂阶数。

边界条件应在每个输运子步前按当前状态更新；若源项改变了边界所依赖的热力学量，源步后也要重新构造幽灵状态。守恒诊断则要区分边界通量造成的变化与体源项积分造成的变化，这样质量、动量或能量偏差才有明确来源。

精确地先解
\begin{equation}r_t+ar_x=0,\quad r(x,0)=u_0(x),\end{equation}
再解
\begin{equation}ds/dt=\lambda s,\quad s(0)=r(x,t),\end{equation}
便得到原精确解。因此以 $C(t)$ 表示对流算子、$S(t)$ 表示源算子，模型方程满足
\begin{equation}u(x,t)=S(t)C(t)u_0.\tag{15.10}\end{equation}
该例中两个作用可交换，故分裂无误差。也可把源步骤写成空间 ODE $d f(u)/dx=\lambda u$，但一般问题较少采用。

% [原PDF p552]
在程序验证中，可关闭输运或关闭源项分别测试两个子算子，再用有解析解的耦合模型测量分裂误差和整体时间阶。

含源项的守恒方程可写成齐次输运算子与局部源算子的组合。Lie 分裂先完整执行一个算子再执行另一个，时间精度通常为一阶；Strang 分裂采用半步源项、整步输运、半步源项的对称次序，在算子足够光滑时达到二阶。两算子不对易造成的交换子误差决定了分裂精度。

分裂后每个子问题都应使用与其自身稳定性相符的求解器。输运步受 CFL 条件限制，源项步则可能由化学、松弛或几何时间尺度控制。把两者机械地使用同一显式步长可能导致刚性失稳，即使齐次有限体积更新本身完全稳定。

\section{基于分裂的数值方法}
对非线性标量方程 $u_t+f(u)_x=s(u)$，用近似算子替代精确算子。一阶 Lie 分裂为
\begin{equation}u^{n+1}=S(\Delta t)C(\Delta t)u^n,\tag{15.17}\end{equation}
即先把齐次守恒律推进完整一步，再以所得状态为初值，在每个单元解源项 ODE 完整一步。若两个子算子至少一阶，这一组合对时间为一阶。

% [原PDF p553]
若源项含守恒交换，应把各组分的增减以成对方式离散，并在每个子步后检查总量，避免 ODE 容差积累成系统漂移。

精确模型问题把平流与指数增长或衰减结合起来，可明确分离空间离散误差、时间积分误差和算子分裂误差。由于常系数平流与线性反应在某些情形下对易，分裂可能表现得异常准确；这类特例不能代替对一般非线性源项的检验。

实现 Strang 分裂时，连续两个主时间步交界处的半源步可在某些自治问题中合并，但输出时刻、可变步长和边界更新会改变这一优化的成立条件。最清楚的实现是显式保留半步—整步—半步结构，并记录每个子步采用的实际时间间隔。

对称 Strang 分裂为
\begin{equation}u^{n+1}=S(\Delta t/2)C(\Delta t)S(\Delta t/2)u^n,\tag{15.18}\end{equation}
当子算子至少二阶时，整体达到二阶。对系统完全相同：
\begin{equation}\bm U^{n+1}=S(\Delta t/2)C(\Delta t)S(\Delta t/2)\bm U^n.\tag{15.23}\end{equation}
源项步骤在每个网格单元独立，天然适合并行；输运步骤则可使用第6--14章任一守恒格式。

% [原PDF p554]
刚性判据来自源项 Jacobian 的特征尺度；它描述稳定性困难，不能仅由源项数值大小或单个变量的瞬时变化判断。

源项常微分方程若具有相差悬殊的特征时间，就属于刚性问题。显式欧拉的稳定区很小，必须取远小于流动 CFL 步长的子步；隐式方法或适合刚性系统的多步、Runge--Kutta 方法虽然每步更昂贵，却允许沿慢时间尺度推进。稳定性与精度要求应分别评估。

源项求解器的误差容限需要与有限体积空间误差协调。把 ODE 解得远比网格误差精确只会浪费计算量；容限过松则会污染守恒量和波速。对放热反应等强耦合问题，还应在源步后重算温度、声速和状态方程，并检查更新状态处于物理可容许域。

另一分裂把第二步写成空间 ODE $d\bm F/dx=\bm S$，对无显式时间依赖或含空间导数的源项可能有利，但经验较少。分裂的吸引力在于可分别选择最适合的 PDE 与 ODE 算法，而不必重新设计一个庞大耦合格式。

\section{ODE 求解器说明}
一阶系统写作
\begin{equation}\bm U'=\bm S(t,\bm U),\qquad
\bm A=\frac{\partial\bm S}{\partial\bm U}.\tag{15.26--15.28}\end{equation}
Jacobian 特征值决定局部稳定性与时间尺度。若所有特征值实部为负，平衡解渐近稳定。

% [原PDF p555]
分裂次序、边界刷新和输出时刻都是离散算法的一部分，复现实验时应与通量格式和 ODE 求解器一起明确记录。

算子分裂便于复用成熟的齐次黎曼求解器和独立 ODE 模块，但也可能破坏稳态平衡。例如重力与压力梯度若被分开处理，会在原本静止的平衡解上产生伪速度。需要保持良好平衡时，应把相互抵消的通量和源项按同一离散关系设计，而不能仅依靠提高 Strang 分裂阶数。

边界条件应在每个输运子步前按当前状态更新；若源项改变了边界所依赖的热力学量，源步后也要重新构造幽灵状态。守恒诊断则要区分边界通量造成的变化与体源项积分造成的变化，这样质量、动量或能量偏差才有明确来源。

刚性系统通常满足全部特征值实部为负，同时
\begin{equation}R_{stiff}=\frac{\max_j|\operatorname{Re}\lambda_j|}
{\min_j|\operatorname{Re}\lambda_j|}\gg1.
\end{equation}
即使比值约20也会给显式算法造成困难，实际化学反应可达 $10^6$。刚性还与所考察时间区间有关：初始短暂层变化剧烈，之后解可能缓慢贴近平衡。选择 ODE 求解器前应先分析或估计 Jacobian 谱。

% [原PDF p556]
在程序验证中，可关闭输运或关闭源项分别测试两个子算子，再用有解析解的耦合模型测量分裂误差和整体时间阶。

含源项的守恒方程可写成齐次输运算子与局部源算子的组合。Lie 分裂先完整执行一个算子再执行另一个，时间精度通常为一阶；Strang 分裂采用半步源项、整步输运、半步源项的对称次序，在算子足够光滑时达到二阶。两算子不对易造成的交换子误差决定了分裂精度。

分裂后每个子问题都应使用与其自身稳定性相符的求解器。输运步受 CFL 条件限制，源项步则可能由化学、松弛或几何时间尺度控制。把两者机械地使用同一显式步长可能导致刚性失稳，即使齐次有限体积更新本身完全稳定。

对初值 $\bm U(t_0)=\bm U_0$，显式 Euler、后向 Euler 与梯形法分别为
\begin{align}
\bm U^{n+1}&=\bm U^n+\Delta t\bm S(t_n,\bm U^n),\tag{15.33}\\
\bm U^{n+1}&=\bm U^n+\Delta t\bm S(t_{n+1},\bm U^{n+1}),\tag{15.34}\\
\bm U^{n+1}&=\bm U^n+\tfrac12\Delta t[\bm S_n+\bm S_{n+1}].\tag{15.35}
\end{align}
显式 Euler 为一阶；后两式隐式，梯形法二阶。常用显式 RK2 与经典 RK4 分别用两级和四级源项评估提高精度。

% [原PDF p557]
若源项含守恒交换，应把各组分的增减以成对方式离散，并在每个子步后检查总量，避免 ODE 容差积累成系统漂移。

精确模型问题把平流与指数增长或衰减结合起来，可明确分离空间离散误差、时间积分误差和算子分裂误差。由于常系数平流与线性反应在某些情形下对易，分裂可能表现得异常准确；这类特例不能代替对一般非线性源项的检验。

实现 Strang 分裂时，连续两个主时间步交界处的半源步可在某些自治问题中合并，但输出时刻、可变步长和边界更新会改变这一优化的成立条件。最清楚的实现是显式保留半步—整步—半步结构，并记录每个子步采用的实际时间间隔。

模型 ODE $u'=\lambda u$ 的显式 Euler 放大因子为 $1+\lambda\Delta t$，稳定要求
\begin{equation}\Delta t\le2/|\lambda|\quad(\lambda<0).\tag{15.39}\end{equation}
刚性越强，允许步长越小。隐式梯形法的放大因子
\begin{equation}\frac{1+\lambda\Delta t/2}{1-\lambda\Delta t/2}\tag{15.40}\end{equation}
对任意 $\Delta t$ 且 $\lambda\le0$ 均稳定。隐式法每步需解非线性代数系统，成本更高，但刚性问题通常别无实用选择。

在分裂实现中，每个网格点都要解一组只持续一个流体时间步的 ODE。先由对流 CFL 得 $\Delta t_c$；若 ODE 稳定步长 $\Delta t_s\ge\Delta t_c$，源项一步完成；否则把源项子循环成 $k$ 个步长 $\Delta t_c/k$ 的稳定小步。

% [原PDF p558]
刚性判据来自源项 Jacobian 的特征尺度；它描述稳定性困难，不能仅由源项数值大小或单个变量的瞬时变化判断。

源项常微分方程若具有相差悬殊的特征时间，就属于刚性问题。显式欧拉的稳定区很小，必须取远小于流动 CFL 步长的子步；隐式方法或适合刚性系统的多步、Runge--Kutta 方法虽然每步更昂贵，却允许沿慢时间尺度推进。稳定性与精度要求应分别评估。

源项求解器的误差容限需要与有限体积空间误差协调。把 ODE 解得远比网格误差精确只会浪费计算量；容限过松则会污染守恒量和波速。对放热反应等强耦合问题，还应在源步后重算温度、声速和状态方程，并检查更新状态处于物理可容许域。

Strang 分裂的两个半源项步骤也应按同样原则子循环。若使用隐式源项求解器，虽通常不受稳定性步长限制，仍须控制精度误差并保证温度、组分等物理约束。

\section{结语}
目前没有对所有源项系统都明显优于分裂的通用替代。源项迎风、良平衡方法等在特定问题中很有效，但必须针对源项结构设计。应用前应优先用精确解验证；最简单的是式(15.6)，更复杂的有 Fickett 爆轰类比和带精确解的反应欧拉问题。

% [原PDF p559]
分裂次序、边界刷新和输出时刻都是离散算法的一部分，复现实验时应与通量格式和 ODE 求解器一起明确记录。

算子分裂便于复用成熟的齐次黎曼求解器和独立 ODE 模块，但也可能破坏稳态平衡。例如重力与压力梯度若被分开处理，会在原本静止的平衡解上产生伪速度。需要保持良好平衡时，应把相互抵消的通量和源项按同一离散关系设计，而不能仅依靠提高 Strang 分裂阶数。

边界条件应在每个输运子步前按当前状态更新；若源项改变了边界所依赖的热力学量，源步后也要重新构造幽灵状态。守恒诊断则要区分边界通量造成的变化与体源项积分造成的变化，这样质量、动量或能量偏差才有明确来源。

爆轰等强耦合反应流可能需要极细网格才能避免错误传播速度，现实多维计算成本很高。分裂也可用来处理扩散项：把双曲对流与抛物扩散分别推进。它还自然引出下一章的方向分裂多维算法。

% [原PDF p560]
\chapter{多维偏微分方程数值方法}
多维离散的正确性可先用均匀流、线性平流和旋转对称问题验证，再进入强激波算例；这些测试分别暴露几何、稳定性和方向偏置错误。

方向分裂把多维守恒律拆成沿各坐标方向的一维更新。常系数线性平流中，各方向算子对易，顺序对结果无影响；非线性系统中，局部波速随状态变化，先做哪个方向会改变后续子步看到的数据，因此一阶顺序分裂产生与步长同阶的分裂误差。

二阶 Strang 方向分裂采用半个 $x$ 步、完整 $y$ 步、再半个 $x$ 步；下一主步可以交换方向，以减小长期方向偏置。三维时需构造对称的子步序列。每个子步都必须使用相应方向的网格间距和法向通量，而不能重复套用未经置换的一维变量顺序。

\section{引言}
笛卡尔坐标下的一般三维守恒律为
\begin{equation}\bm U_t+\bm F(\bm U)_x+\bm G(\bm U)_y+\bm H(\bm U)_z=0.\tag{16.1}\end{equation}
本章讨论两条路线：方向分裂把各坐标方向依次交给一维格式；非分裂有限体积法在同一步中汇总控制体所有面的通量。前者几乎可复用第6--14章任何一维方法，后者更直接体现多维守恒，但高阶构造和稳定性更复杂。

% [原PDF p561]
不论采用分裂还是非分裂形式，单个内部面都只能产生一个法向通量；两侧重复求解而结果不一致会直接破坏离散守恒。

欧拉方程沿某一方向求解时，法向速度参与声学特征值，两个切向速度则作为随接触波输运的量。实现中可通过变量置换复用同一一维黎曼求解器：进入求解器前把目标方向速度放到法向分量，返回通量后再映射回全局动量顺序。旋转或置换必须同时作用于左右状态和输出通量。

多维显式时间步要依据所有方向的传播速度确定。对分裂格式，常用限制可按各方向最小 CFL 步长取全局最小值；对非分裂格式，稳定条件往往包含各方向谱半径之和。无论采用哪种表达，都要以实际网格尺度和局部声速计算，不能只检查某一个速度分量。

\section{方向分裂}
对三维常速平流
\begin{equation}u_t+a_1u_x+a_2u_y+a_3u_z=0,\tag{16.2}\end{equation}
依次作 $x,y,z$ 扫掠，每个扫掠都推进完整 $\Delta t$。以 $X,Y,Z$ 表示三个精确一维传播算子，则由于常系数平移彼此交换，方向分裂可精确恢复三维解。

% [原PDF p562]
高阶预测所用的状态必须全部来自同一旧时间层，通量全部形成后再统一更新；这可避免隐含的扫描顺序依赖。

分裂算法的优点是代码结构清晰，一维高分辨率模块可直接复用；缺点是斜向波和旋涡在一个子步中被人为看成沿坐标轴传播，可能形成网格方向偏置。交替扫描次序能够降低但不能完全消除这种误差。检验时应比较旋转后的同一问题，而不只运行与网格轴对齐的激波管。

每个方向扫掠后，守恒变量已经更新，下一方向必须读取这个新状态并重新计算压力、声速及边界幽灵单元。若缓存了一维黎曼数据，缓存不能跨方向或跨子步复用。对 WAF 方法，切向速度属于线性退化波携带的量，也应接受与接触波一致的限制。

具体地
\begin{equation}u^{n+1}=Z(\Delta t)Y(\Delta t)X(\Delta t)u^n.\tag{16.6}\end{equation}
对非线性二维系统，最简单的 $YX$ 或 $XY$ 分裂只有一阶时间精度。对称 Strang 组合
\begin{equation}
\bm U^{n+1}=X(\Delta t/2)Y(\Delta t)X(\Delta t/2)\bm U^n\tag{16.13}
\end{equation}
在子算子至少二阶且解光滑时达到二阶；交换 $X,Y$ 得等价形式。

% [原PDF p563]
切向速度虽不进入法向声学特征值，却随接触波输运并参与动能，因此其旋转、限制和回写都不能遗漏。

三维数组的扫掠通常把固定的两条索引视为一条一维线。为保持性能，可以把一条线的状态复制到连续工作区，完成边界、重构、黎曼和更新后再写回；但拷贝范围必须包含高阶模板所需的幽灵层。并行分区的 halo 交换也应在每个方向子步前满足当前模板宽度。

方向分裂仍保持每个一维子步的守恒性，因为相邻单元共享同一法向界面通量。然而总精度还取决于分裂次序与时间步。若使用可变步长，两个半步必须精确等于本主步的一半；跨输出时刻截断步长后，应重新构造完整的对称序列。

每个 $x$ 扫掠对所有固定 $j$ 的行求一维问题，每个 $y$ 扫掠对所有固定 $i$ 的列求一维问题。两个方向使用同一全局时间步。平均 $XY$ 与 $YX$ 也可达二阶，但工作量翻倍；对称半步组合只增加约50%。

连续多个时间步时，相邻的两个同方向半步可合并成一个整步，使 Strang 分裂几乎达到一阶分裂的效率。非线性问题的 $\Delta t$ 随波速变化，合并时应确保外侧半步使用一致的时间长度。

% [原PDF p564]
多维离散的正确性可先用均匀流、线性平流和旋转对称问题验证，再进入强激波算例；这些测试分别暴露几何、稳定性和方向偏置错误。

方向分裂把多维守恒律拆成沿各坐标方向的一维更新。常系数线性平流中，各方向算子对易，顺序对结果无影响；非线性系统中，局部波速随状态变化，先做哪个方向会改变后续子步看到的数据，因此一阶顺序分裂产生与步长同阶的分裂误差。

二阶 Strang 方向分裂采用半个 $x$ 步、完整 $y$ 步、再半个 $x$ 步；下一主步可以交换方向，以减小长期方向偏置。三维时需构造对称的子步序列。每个子步都必须使用相应方向的网格间距和法向通量，而不能重复套用未经置换的一维变量顺序。

另一种廉价二阶策略在两个时间步内使用对称扫掠序列 $XYYX$ 或 $YXXY$，每隔一步达到二阶。三维最简单分裂为
\begin{equation}\bm U^{n+1}=Z(\Delta t)Y(\Delta t)X(\Delta t)\bm U^n,\tag{16.24}\end{equation}
而对称二阶形式为
\begin{equation}\bm U^{n+1}=X(\Delta t/2)Y(\Delta t/2)Z(\Delta t)
Y(\Delta t/2)X(\Delta t/2)\bm U^n.\tag{16.26}\end{equation}

% [原PDF p565]
不论采用分裂还是非分裂形式，单个内部面都只能产生一个法向通量；两侧重复求解而结果不一致会直接破坏离散守恒。

欧拉方程沿某一方向求解时，法向速度参与声学特征值，两个切向速度则作为随接触波输运的量。实现中可通过变量置换复用同一一维黎曼求解器：进入求解器前把目标方向速度放到法向分量，返回通量后再映射回全局动量顺序。旋转或置换必须同时作用于左右状态和输出通量。

多维显式时间步要依据所有方向的传播速度确定。对分裂格式，常用限制可按各方向最小 CFL 步长取全局最小值；对非分裂格式，稳定条件往往包含各方向谱半径之和。无论采用哪种表达，都要以实际网格尺度和局部声速计算，不能只检查某一个速度分量。

分裂阶数推导依赖光滑解的 Taylor 展开。间断解中，方向算子不交换与一维格式的间断处理会产生额外误差。能把激波保持为真间断的随机选择法反而可能在方向分裂中完全失败；通常带适度涂抹的一维守恒格式更稳健。接触面在某些特殊构型下例外。

% [原PDF p566]
高阶预测所用的状态必须全部来自同一旧时间层，通量全部形成后再统一更新；这可避免隐含的扫描顺序依赖。

分裂算法的优点是代码结构清晰，一维高分辨率模块可直接复用；缺点是斜向波和旋涡在一个子步中被人为看成沿坐标轴传播，可能形成网格方向偏置。交替扫描次序能够降低但不能完全消除这种误差。检验时应比较旋转后的同一问题，而不只运行与网格轴对齐的激波管。

每个方向扫掠后，守恒变量已经更新，下一方向必须读取这个新状态并重新计算压力、声速及边界幽灵单元。若缓存了一维黎曼数据，缓存不能跨方向或跨子步复用。对 WAF 方法，切向速度属于线性退化波携带的量，也应接受与接触波一致的限制。

\section{三维方向分裂的实际实现}
三维欧拉变量为 $\bm U=(\rho,\rho u,\rho v,\rho w,E)^T$。三个方向通量分别以 $u,v,w$ 为法向速度，法向动量通量含压力，两个切向动量按法向速度输运，总能量通量为法向速度乘 $E+p$。

为复用一个一维子程序，每次扫掠重排变量：$x$ 扫掠顺序 $(\rho,\rho u,E,\rho v,\rho w)$，$y$ 扫掠把 $v$ 放到法向位置，$z$ 扫掠把 $w$ 放到法向位置。

% [原PDF p567]
切向速度虽不进入法向声学特征值，却随接触波输运并参与动能，因此其旋转、限制和回写都不能遗漏。

三维数组的扫掠通常把固定的两条索引视为一条一维线。为保持性能，可以把一条线的状态复制到连续工作区，完成边界、重构、黎曼和更新后再写回；但拷贝范围必须包含高阶模板所需的幽灵层。并行分区的 halo 交换也应在每个方向子步前满足当前模板宽度。

方向分裂仍保持每个一维子步的守恒性，因为相邻单元共享同一法向界面通量。然而总精度还取决于分裂次序与时间步。若使用可变步长，两个半步必须精确等于本主步的一半；跨输出时刻截断步长后，应重新构造完整的对称序列。

在 $x$ 扫掠中，切向速度满足
\begin{equation}v_t+uv_x=0,qquad w_t+uw_x=0,\tag{16.32}\end{equation}
即以法向速度被动输运。三个守恒更新分别为沿相应方向的界面通量差；$x$ 后得到 $\bm U^{n+1/3}$，$y$ 后得到 $\bm U^{n+2/3}$，$z$ 后得到 $\bm U^{n+1}$。总能量仍必须包含三个速度分量的动能。

% [原PDF p568]
多维离散的正确性可先用均匀流、线性平流和旋转对称问题验证，再进入强激波算例；这些测试分别暴露几何、稳定性和方向偏置错误。

方向分裂把多维守恒律拆成沿各坐标方向的一维更新。常系数线性平流中，各方向算子对易，顺序对结果无影响；非线性系统中，局部波速随状态变化，先做哪个方向会改变后续子步看到的数据，因此一阶顺序分裂产生与步长同阶的分裂误差。

二阶 Strang 方向分裂采用半个 $x$ 步、完整 $y$ 步、再半个 $x$ 步；下一主步可以交换方向，以减小长期方向偏置。三维时需构造对称的子步序列。每个子步都必须使用相应方向的网格间距和法向通量，而不能重复套用未经置换的一维变量顺序。

单一子程序调用可概括为：$x$ 向传入 $(\rho,u,v,w,p,\Delta x)$，$y$ 向传入 $(\rho,v,u,w,p,\Delta y)$，$z$ 向传入 $(\rho,w,u,v,p,\Delta z)$。高阶对称分裂只需改变调用顺序和半步/整步时间长度。

全局稳定时间步可取
\begin{equation}
\Delta t=C_{\rm CFL}\min_{i,j,k}\left(
\frac{\Delta x}{S_x},\frac{\Delta y}{S_y},\frac{\Delta z}{S_z}\right).\tag{16.37}
\end{equation}
最简单 $S_d=|u_d|+a$ 对强初始间断可能低估真实波速，需更可靠的界面波速或初期较小 CFL。

% [原PDF p569]
不论采用分裂还是非分裂形式，单个内部面都只能产生一个法向通量；两侧重复求解而结果不一致会直接破坏离散守恒。

欧拉方程沿某一方向求解时，法向速度参与声学特征值，两个切向速度则作为随接触波输运的量。实现中可通过变量置换复用同一一维黎曼求解器：进入求解器前把目标方向速度放到法向分量，返回通量后再映射回全局动量顺序。旋转或置换必须同时作用于左右状态和输出通量。

多维显式时间步要依据所有方向的传播速度确定。对分裂格式，常用限制可按各方向最小 CFL 步长取全局最小值；对非分裂格式，稳定条件往往包含各方向谱半径之和。无论采用哪种表达，都要以实际网格尺度和局部声速计算，不能只检查某一个速度分量。

\subsection{界面通量与 TVD 条件}
一维欧拉加两个切向动量后有五个波：两条声波、接触波和两条剪切波。WAF 限制时，密度可限制声波和接触波，切向速度跳跃则分别限制两条剪切波。若不同线性波得到不同耗散，原本重合于速度 $u_*$ 的波会在“耗散波结构”中分开，并产生新的中间状态。

% [原PDF p570]
高阶预测所用的状态必须全部来自同一旧时间层，通量全部形成后再统一更新；这可避免隐含的扫描顺序依赖。

分裂算法的优点是代码结构清晰，一维高分辨率模块可直接复用；缺点是斜向波和旋涡在一个子步中被人为看成沿坐标轴传播，可能形成网格方向偏置。交替扫描次序能够降低但不能完全消除这种误差。检验时应比较旋转后的同一问题，而不只运行与网格轴对齐的激波管。

每个方向扫掠后，守恒变量已经更新，下一方向必须读取这个新状态并重新计算压力、声速及边界幽灵单元。若缓存了一维黎曼数据，缓存不能跨方向或跨子步复用。对 WAF 方法，切向速度属于线性退化波携带的量，也应接受与接触波一致的限制。

三条线性波的限制系数可有六种排序，对应不同的中间密度与切向速度组合。简单地把三者都设为最大耗散虽然安全，却会过分涂抹某些波，不推荐。更合理的是先按一维三方程系统限制声波与接触波，再把每个切向动量当作被动标量单独处理。

% [原PDF p571]
切向速度虽不进入法向声学特征值，却随接触波输运并参与动能，因此其旋转、限制和回写都不能遗漏。

三维数组的扫掠通常把固定的两条索引视为一条一维线。为保持性能，可以把一条线的状态复制到连续工作区，完成边界、重构、黎曼和更新后再写回；但拷贝范围必须包含高阶模板所需的幽灵层。并行分区的 halo 交换也应在每个方向子步前满足当前模板宽度。

方向分裂仍保持每个一维子步的守恒性，因为相邻单元共享同一法向界面通量。然而总精度还取决于分裂次序与时间步。若使用可变步长，两个半步必须精确等于本主步的一半；跨输出时刻截断步长后，应重新构造完整的对称序列。

对任一被动标量 $\eta=v,w$，写成
\begin{equation}(\rho\eta)_t+(f_1\eta)_x=0,\tag{16.43}\end{equation}
其中 $f_1$ 是已经受限的质量通量。其动量通量可取
\begin{equation}
(\rho\eta)^{TVD}=\left[\tfrac12(1+\phi_\eta)\eta_i+
\tfrac12(1-\phi_\eta)\eta_{i+1}\right]f_1.\tag{16.44}
\end{equation}
这种做法能用同一逻辑处理任意多个组分或被动标量。需要注意，总能量含切向动能，仅以密度限制前三式并不完全理想。

% [原PDF p572]
多维离散的正确性可先用均匀流、线性平流和旋转对称问题验证，再进入强激波算例；这些测试分别暴露几何、稳定性和方向偏置错误。

非分裂有限体积法在一个更新式中同时汇总单元所有面的通量，不引入方向子步。二维笛卡尔网格上的更新由左右两个 $x$ 向通量和上下两个 $y$ 向通量组成；每个内部面只计算一次，并由相邻单元以相反符号使用。这样离散守恒与控制体积分形式直接一致。

与方向分裂相比，非分裂方法能在同一时间层处理斜向传播和多维波相互作用，但界面预测必须包含横向影响。若仅沿面法向做彼此独立的一维预测，形式上虽是一次更新，实际仍会遗漏交叉导数，并可能降低时间精度或稳定范围。

\section{非分裂有限体积法}
非分裂法在一次更新中加入所有面的通量，与方向分裂逐向推进不同。二者都源于控制体积分守恒；若分裂子格式本身是一维有限体积法，可把分裂理解成某种预测--校正多维有限体积过程。

对二维系统 $\bm U_t+\bm F_x+\bm G_y=0$，矩形单元平均位于中心，四个面分别有 $x$、$y$ 数值通量。

% [原PDF p573]
不论采用分裂还是非分裂形式，单个内部面都只能产生一个法向通量；两侧重复求解而结果不一致会直接破坏离散守恒。

二维 MUSCL--Hancock 在单元内同时重构 $x$、$y$ 两个方向的斜率。四个面中心状态由同一个单元平均外推得到，半步预测则使用两个方向的通量散度，使每个面状态都感受到横向传播。随后在每个面法向上求一维黎曼问题，得到完整的守恒通量。

斜率可以逐守恒量、逐原始量或在各方向特征空间中限制。为了避免角点附近不同方向的外推共同产生负压力，需检查全部预测面状态，而不只是单元中心。复杂强激波下，多维限制器还应考虑梯度合成后的幅值，防止各方向单独合法而合成状态越界。

显式非分裂更新为
\begin{equation}
\bm U_{ij}^{n+1}=\bm U_{ij}^n+\frac{\Delta t}{\Delta x}
(\bm F_{i-1/2,j}-\bm F_{i+1/2,j})+\frac{\Delta t}{\Delta y}
(\bm G_{i,j-1/2}-\bm G_{i,j+1/2}).\tag{16.46}
\end{equation}
一阶 Godunov 有限体积法在每个面独立求法向一维黎曼问题，并取 $x/t=0$ 或 $y/t=0$ 的物理通量。一个面只使用该面两侧的单元平均。

% [原PDF p574]
高阶预测所用的状态必须全部来自同一旧时间层，通量全部形成后再统一更新；这可避免隐含的扫描顺序依赖。

非分裂 WAF 方法从界面附近的多维时空控制体平均出发。法向黎曼扇决定主要波结构，横向通量差则修正半时间层状态；所得公式可理解为一维 WAF 与横向预测的结合。若省略横向项，斜向接触和旋涡会表现出更强的网格依赖。

多维 Courant 条件不应简单照搬一维 $|c|\le1$。在笛卡尔网格上，常见稳定约束涉及 $\Delta t(|u|+a)/\Delta x$ 与 $\Delta t(|v|+a)/\Delta y$ 的组合。实际允许系数还取决于通量、限制器与预测方式，应以理论范围为上限并通过标准算例验证。

二维常速平流且 $a_1,a_2>0$ 时
\begin{equation}
u_{ij}^{n+1}=u_{ij}^n+c_1(u_{i-1,j}^n-u_{ij}^n)+
c_2(u_{i,j-1}^n-u_{ij}^n),\tag{16.51}
\end{equation}
$c_1=a_1\Delta t/\Delta x,c_2=a_2\Delta t/\Delta y$。模板只含两个坐标方向的上风邻点，没有直接使用二维速度方向上的对角单元 $(i-1,j-1)$，因此只是逐面迎风而非完整多维迎风。

% [原PDF p575]
切向速度虽不进入法向声学特征值，却随接触波输运并参与动能，因此其旋转、限制和回写都不能遗漏。

有限体积重构的对象是单元平均，不是单元中心点值。在线性函数上二者相同，但高阶推导中必须区分；面值应由保持单元平均的多项式积分关系得到。二维情况下，斜率和交叉项的矩条件决定空间精度，随意把有限差分点值公式套到单元平均上可能引入系统偏差。

一个健壮的非分裂实现先重构并预测所有单元的面状态，再统一求各方向的界面通量，最后一次性更新单元平均。若边算一个面边更新相邻单元，后续界面会读取混合时间层状态，破坏显式格式的定义。因此旧状态、预测状态、通量和新状态需要清楚分离。

三维更新再加 $z$ 面通量差。常速正向平流形式为
\begin{equation}
u_{ijk}^{n+1}=u_{ijk}^n+\sum_{d=1}^3c_d(u_{\mathrm{upw},d}^n-u_{ijk}^n).\tag{16.56}
\end{equation}
一般常系数二维格式的精度判据为
\begin{equation}\sum_{\alpha,\beta}\alpha^q\beta^r b_{\alpha\beta}=(-c_1)^q(-c_2)^r,
\quad q+r\le m,\tag{16.59}\end{equation}
三维则增加 $\gamma^s$ 与 $(-c_3)^s$。

% [原PDF p576]
多维离散的正确性可先用均匀流、线性平流和旋转对称问题验证，再进入强激波算例；这些测试分别暴露几何、稳定性和方向偏置错误。

边界附近的非分裂预测同时需要法向和切向信息。固壁上法向速度按反射条件处理，切向速度和热力学量依物理边界确定；角点则必须使两个边界规则相容。仅逐方向填充幽灵单元可能在角点产生冲突，应明确规定角点状态或使用几何一致的镜像构造。

守恒检查应按面通量逐项核对：封闭或周期区域的所有内部面贡献必须成对抵消，剩余变化只能来自物理边界。对多维实现，这一检查比单独比较解析剖面更容易发现面编号、法向符号或动量分量置换错误。

上述矩条件证明逐面 Godunov 格式为时空一阶。线性稳定性可用 von Neumann 放大因子数值扫描：对每组方向 Courant 数，在大量相位角组合上计算 $|A|$，若全部不超过一，则判为稳定。这样的离散扫描给出证据，但不是严格解析证明。

% [原PDF p577]
不论采用分裂还是非分裂形式，单个内部面都只能产生一个法向通量；两侧重复求解而结果不一致会直接破坏离散守恒。

数值实验应包含轴对齐和平移旋转后的问题。若旋转后激波速度、接触厚度或对称性明显恶化，通常说明横向预测、限制器或 CFL 条件存在方向依赖。光滑涡输运可用于测量二阶收敛，二维黎曼问题则检验多波相互作用和格式鲁棒性。

分辨率比较必须保持相同物理终止时间和近似一致的 CFL 数。高阶格式在粗网格上通常能减少耗散，但计算代价包括更多重构、状态转换与黎曼求解；公平评价应同时报告误差、网格数、时间步数和运行成本。

二维 Godunov 非分裂格式的稳定条件为
\begin{equation}c_1+c_2\le1,\tag{16.62}\end{equation}
三维为
\begin{equation}c_1+c_2+c_3\le1.\tag{16.63}\end{equation}
若各方向 Courant 数相等，上限分别降为 $1/2$ 和 $1/3$。直接把稳定的一维通量放入多维更新，通常缩小稳定域，某些组合甚至会无条件失稳。

% [原PDF p578]
高阶预测所用的状态必须全部来自同一旧时间层，通量全部形成后再统一更新；这可避免隐含的扫描顺序依赖。

非分裂有限体积法在一个更新式中同时汇总单元所有面的通量，不引入方向子步。二维笛卡尔网格上的更新由左右两个 $x$ 向通量和上下两个 $y$ 向通量组成；每个内部面只计算一次，并由相邻单元以相反符号使用。这样离散守恒与控制体积分形式直接一致。

与方向分裂相比，非分裂方法能在同一时间层处理斜向传播和多维波相互作用，但界面预测必须包含横向影响。若仅沿面法向做彼此独立的一维预测，形式上虽是一次更新，实际仍会遗漏交叉导数，并可能降低时间精度或稳定范围。

\section{MUSCL--Hancock 非分裂有限体积格式}
二维算法仍分三步。首先沿 $x,y$ 独立重构四个面外推值
\begin{equation}
\bm U_{ij}^{\mp x}=\bm U_{ij}^n\mp\tfrac12\bm\Delta_i,qquad
\bm U_{ij}^{\mp y}=\bm U_{ij}^n\mp\tfrac12\bm\Delta_j.\tag{16.64}
\end{equation}
其次所有面值都用两个方向的单元内通量差共同预测半步：
\begin{equation}
\widehat{\bm U}_{ij}^{l}=\bm U_{ij}^{l}+\frac{\Delta t}{2\Delta x}(\bm F^{-x}-\bm F^{+x})+
\frac{\Delta t}{2\Delta y}(\bm G^{-y}-\bm G^{+y}).\tag{16.65}
\end{equation}

% [原PDF p579]
切向速度虽不进入法向声学特征值，却随接触波输运并参与动能，因此其旋转、限制和回写都不能遗漏。

二维 MUSCL--Hancock 在单元内同时重构 $x$、$y$ 两个方向的斜率。四个面中心状态由同一个单元平均外推得到，半步预测则使用两个方向的通量散度，使每个面状态都感受到横向传播。随后在每个面法向上求一维黎曼问题，得到完整的守恒通量。

斜率可以逐守恒量、逐原始量或在各方向特征空间中限制。为了避免角点附近不同方向的外推共同产生负压力，需检查全部预测面状态，而不只是单元中心。复杂强激波下，多维限制器还应考虑梯度合成后的幅值，防止各方向单独合法而合成状态越界。

然后在每个 $x$ 面以相邻预测值求 $x$ 向黎曼问题，在每个 $y$ 面求 $y$ 向黎曼问题，并把两组通量同时代入式(16.46)。对线性正速平流，$x$ 通量除本向斜率外还含横向预测项 $-c_2\Delta_j/2$，$y$ 通量同样含 $-c_1\Delta_i/2$；这些交叉项使格式真正非分裂。

采用中心差分斜率可得十点模板、时空二阶。稳定域仍约为 $c_1+c_2\le1$：高阶重构提高精度，但没有扩大一阶 Godunov 的稳定区。非线性系统的面通量可用精确或第9--12章任一近似黎曼求解器。

% [原PDF p580]
多维离散的正确性可先用均匀流、线性平流和旋转对称问题验证，再进入强激波算例；这些测试分别暴露几何、稳定性和方向偏置错误。

非分裂 WAF 方法从界面附近的多维时空控制体平均出发。法向黎曼扇决定主要波结构，横向通量差则修正半时间层状态；所得公式可理解为一维 WAF 与横向预测的结合。若省略横向项，斜向接触和旋涡会表现出更强的网格依赖。

多维 Courant 条件不应简单照搬一维 $|c|\le1$。在笛卡尔网格上，常见稳定约束涉及 $\Delta t(|u|+a)/\Delta x$ 与 $\Delta t(|v|+a)/\Delta y$ 的组合。实际允许系数还取决于通量、限制器与预测方式，应以理论范围为上限并通过标准算例验证。

\section{WAF 型有限体积格式}
多维 WAF 先在标量平流上建立。沿某一面评估半时间层空间平均时，横向运动会使积分区域同时受邻近行/列数据影响，因此模板比一维 WAF 更宽。所得二阶格式仍需满足各方向联合 CFL 条件。

% [原PDF p581]
不论采用分裂还是非分裂形式，单个内部面都只能产生一个法向通量；两侧重复求解而结果不一致会直接破坏离散守恒。

有限体积重构的对象是单元平均，不是单元中心点值。在线性函数上二者相同，但高阶推导中必须区分；面值应由保持单元平均的多项式积分关系得到。二维情况下，斜率和交叉项的矩条件决定空间精度，随意把有限差分点值公式套到单元平均上可能引入系统偏差。

一个健壮的非分裂实现先重构并预测所有单元的面状态，再统一求各方向的界面通量，最后一次性更新单元平均。若边算一个面边更新相邻单元，后续界面会读取混合时间层状态，破坏显式格式的定义。因此旧状态、预测状态、通量和新状态需要清楚分离。

二维常速平流的 WAF 面通量可视为先沿横向作 Godunov 半步预测，再沿法向对局部黎曼结构取加权平均。这种“横向一阶、法向二阶”的组合在完整更新中恢复时空二阶，并自然产生角向耦合项。

% [原PDF p582]
高阶预测所用的状态必须全部来自同一旧时间层，通量全部形成后再统一更新；这可避免隐含的扫描顺序依赖。

边界附近的非分裂预测同时需要法向和切向信息。固壁上法向速度按反射条件处理，切向速度和热力学量依物理边界确定；角点则必须使两个边界规则相容。仅逐方向填充幽灵单元可能在角点产生冲突，应明确规定角点状态或使用几何一致的镜像构造。

守恒检查应按面通量逐项核对：封闭或周期区域的所有内部面贡献必须成对抵消，剩余变化只能来自物理边界。对多维实现，这一检查比单独比较解析剖面更容易发现面编号、法向符号或动量分量置换错误。

对 $x$ 面，半时间层积分范围跨越相邻两个单元的中线；正 $a_1$ 时法向状态由左侧贡献，横向 $a_2$ 决定来自上方或下方的部分。几何权重由 $c_1,c_2$ 给出，所有权重之和保持一，保证常数状态精确保存。

% [原PDF p583]
切向速度虽不进入法向声学特征值，却随接触波输运并参与动能，因此其旋转、限制和回写都不能遗漏。

数值实验应包含轴对齐和平移旋转后的问题。若旋转后激波速度、接触厚度或对称性明显恶化，通常说明横向预测、限制器或 CFL 条件存在方向依赖。光滑涡输运可用于测量二阶收敛，二维黎曼问题则检验多波相互作用和格式鲁棒性。

分辨率比较必须保持相同物理终止时间和近似一致的 CFL 数。高阶格式在粗网格上通常能减少耗散，但计算代价包括更多重构、状态转换与黎曼求解；公平评价应同时报告误差、网格数、时间步数和运行成本。

一阶多维 WAF 退化为逐面 Godunov 格式，其模板只含本单元与两个方向上风单元。二阶 WAF 增加次上风及横向相邻点，模板覆盖更完整的上风区域，因此能比简单逐面方法更准确地输运斜向波形。

% [原PDF p584]
多维离散的正确性可先用均匀流、线性平流和旋转对称问题验证，再进入强激波算例；这些测试分别暴露几何、稳定性和方向偏置错误。

非分裂有限体积法在一个更新式中同时汇总单元所有面的通量，不引入方向子步。二维笛卡尔网格上的更新由左右两个 $x$ 向通量和上下两个 $y$ 向通量组成；每个内部面只计算一次，并由相邻单元以相反符号使用。这样离散守恒与控制体积分形式直接一致。

与方向分裂相比，非分裂方法能在同一时间层处理斜向传播和多维波相互作用，但界面预测必须包含横向影响。若仅沿面法向做彼此独立的一维预测，形式上虽是一次更新，实际仍会遗漏交叉导数，并可能降低时间精度或稳定范围。

二维二阶 WAF 模板包含法向两个上风层和横向耦合点。用多维精度矩条件可验证其二阶性；稳定性仍需对 $c_1,c_2$ 联合分析，不能分别满足一维 CFL 后就假定多维稳定。

% [原PDF p585]
不论采用分裂还是非分裂形式，单个内部面都只能产生一个法向通量；两侧重复求解而结果不一致会直接破坏离散守恒。

二维 MUSCL--Hancock 在单元内同时重构 $x$、$y$ 两个方向的斜率。四个面中心状态由同一个单元平均外推得到，半步预测则使用两个方向的通量散度，使每个面状态都感受到横向传播。随后在每个面法向上求一维黎曼问题，得到完整的守恒通量。

斜率可以逐守恒量、逐原始量或在各方向特征空间中限制。为了避免角点附近不同方向的外推共同产生负压力，需检查全部预测面状态，而不只是单元中心。复杂强激波下，多维限制器还应考虑梯度合成后的幅值，防止各方向单独合法而合成状态越界。

三维 WAF 更新使用 $x,y,z$ 三组面通量。一般 WAF 算子可记作 $L^{waf}_{x,\Delta t/2}(\bm U_L,\bm U_R)$，表示在指定方向以半时间信息处理一对状态；通过组合横向 Godunov 预测，可构造真正多维界面输入。

% [原PDF p586]
高阶预测所用的状态必须全部来自同一旧时间层，通量全部形成后再统一更新；这可避免隐含的扫描顺序依赖。

非分裂 WAF 方法从界面附近的多维时空控制体平均出发。法向黎曼扇决定主要波结构，横向通量差则修正半时间层状态；所得公式可理解为一维 WAF 与横向预测的结合。若省略横向项，斜向接触和旋涡会表现出更强的网格依赖。

多维 Courant 条件不应简单照搬一维 $|c|\le1$。在笛卡尔网格上，常见稳定约束涉及 $\Delta t(|u|+a)/\Delta x$ 与 $\Delta t(|v|+a)/\Delta y$ 的组合。实际允许系数还取决于通量、限制器与预测方式，应以理论范围为上限并通过标准算例验证。

最简单的三维构造先用另外两个方向的一阶 Godunov 算子把相邻单元推进半步，再把得到的两个预测状态送入当前方向 WAF。精度定理表明，若缺少某些交叉项只得到一阶；完整对称组合才能达到二阶。

% [原PDF p587]
切向速度虽不进入法向声学特征值，却随接触波输运并参与动能，因此其旋转、限制和回写都不能遗漏。

有限体积重构的对象是单元平均，不是单元中心点值。在线性函数上二者相同，但高阶推导中必须区分；面值应由保持单元平均的多项式积分关系得到。二维情况下，斜率和交叉项的矩条件决定空间精度，随意把有限差分点值公式套到单元平均上可能引入系统偏差。

一个健壮的非分裂实现先重构并预测所有单元的面状态，再统一求各方向的界面通量，最后一次性更新单元平均。若边算一个面边更新相邻单元，后续界面会读取混合时间层状态，破坏显式格式的定义。因此旧状态、预测状态、通量和新状态需要清楚分离。

对二维线性平流，二阶 WAF 通量展开后满足 $q+r\le2$ 的全部矩条件。练习要求逐项核对常数、一阶以及 $x^2,xy,y^2$ 矩；其中混合矩 $xy$ 正是检验横向耦合是否正确的关键。

% [原PDF p588]
多维离散的正确性可先用均匀流、线性平流和旋转对称问题验证，再进入强激波算例；这些测试分别暴露几何、稳定性和方向偏置错误。

边界附近的非分裂预测同时需要法向和切向信息。固壁上法向速度按反射条件处理，切向速度和热力学量依物理边界确定；角点则必须使两个边界规则相容。仅逐方向填充幽灵单元可能在角点产生冲突，应明确规定角点状态或使用几何一致的镜像构造。

守恒检查应按面通量逐项核对：封闭或周期区域的所有内部面贡献必须成对抵消，剩余变化只能来自物理边界。对多维实现，这一检查比单独比较解析剖面更容易发现面编号、法向符号或动量分量置换错误。

对非线性二维系统，$x$ 面通量的构造可写成：先沿 $y$ 方向对面两侧单元分别作半步 Godunov 更新，得到 $\bm U_{L}^{god},\bm U_R^{god}$；再用 $x$ 向 WAF 算子计算
\begin{equation}\bm F_{i+1/2,j}=L^{waf}_{x,\Delta t/2}(\bm U_L^{god},\bm U_R^{god}).\tag{16.91a}\end{equation}
$y$ 面作对称处理。这样面通量包含横向传播信息，同时保持局部一维黎曼求解器可复用。

% [原PDF p589]
不论采用分裂还是非分裂形式，单个内部面都只能产生一个法向通量；两侧重复求解而结果不一致会直接破坏离散守恒。

数值实验应包含轴对齐和平移旋转后的问题。若旋转后激波速度、接触厚度或对称性明显恶化，通常说明横向预测、限制器或 CFL 条件存在方向依赖。光滑涡输运可用于测量二阶收敛，二维黎曼问题则检验多波相互作用和格式鲁棒性。

分辨率比较必须保持相同物理终止时间和近似一致的 CFL 数。高阶格式在粗网格上通常能减少耗散，但计算代价包括更多重构、状态转换与黎曼求解；公平评价应同时报告误差、网格数、时间步数和运行成本。

横向 Godunov 半步本身也通过相邻横向面黎曼通量计算。几何上，最终 $x$ 面 WAF 平均使用的左右状态已包含 $y$ 向影响，因而不再是原始单元平均。若忽略这一步，格式退化为逐方向独立的简化扩展。

% [原PDF p590]
高阶预测所用的状态必须全部来自同一旧时间层，通量全部形成后再统一更新；这可避免隐含的扫描顺序依赖。

非分裂有限体积法在一个更新式中同时汇总单元所有面的通量，不引入方向子步。二维笛卡尔网格上的更新由左右两个 $x$ 向通量和上下两个 $y$ 向通量组成；每个内部面只计算一次，并由相邻单元以相反符号使用。这样离散守恒与控制体积分形式直接一致。

与方向分裂相比，非分裂方法能在同一时间层处理斜向传播和多维波相互作用，但界面预测必须包含横向影响。若仅沿面法向做彼此独立的一维预测，形式上虽是一次更新，实际仍会遗漏交叉导数，并可能降低时间精度或稳定范围。

TVD 版本在 WAF 算子内部对每个波族应用限制器；横向 Godunov 预测使用一阶模式，即令 WAF 参数取对应迎风值。由于限制器比值需要相邻界面波跳跃，非分裂多维实现必须组织更宽的临时数据与边界幽灵层。

% [原PDF p591]
切向速度虽不进入法向声学特征值，却随接触波输运并参与动能，因此其旋转、限制和回写都不能遗漏。

二维 MUSCL--Hancock 在单元内同时重构 $x$、$y$ 两个方向的斜率。四个面中心状态由同一个单元平均外推得到，半步预测则使用两个方向的通量散度，使每个面状态都感受到横向传播。随后在每个面法向上求一维黎曼问题，得到完整的守恒通量。

斜率可以逐守恒量、逐原始量或在各方向特征空间中限制。为了避免角点附近不同方向的外推共同产生负压力，需检查全部预测面状态，而不只是单元中心。复杂强激波下，多维限制器还应考虑梯度合成后的幅值，防止各方向单独合法而合成状态越界。

三维推广采用嵌套横向预测：计算 $x$ 面前，先由 $y,z$ 两方向的一阶 Godunov 通量修正左右状态半步；再执行 $x$ 向 WAF。三个方向循环对称，最终同时代入非分裂守恒更新，而不是逐向覆盖状态。

% [原PDF p592]
几何量和流动物理量应分层存储并分别验证；自由流不能保持时，应优先检查面积向量闭合、体积符号和共享面方向。

在非笛卡尔网格上，有限体积守恒式仍由控制体体积和各面通量积分组成。几何信息进入面面积、外法向和单元体积，而不是改变守恒定律本身。每个内部面必须由两个相邻单元共享完全相同的面积向量，仅符号相反，才能使内部通量严格抵消。

广义坐标变换把物理空间方程映射到计算空间，但度量系数必须满足离散几何守恒关系。若常数状态在静止曲线网格上不能保持不变，误差来源通常不是黎曼求解器，而是面向量、Jacobian 或度量导数的不一致离散。自由流保持应作为首要单元测试。

\section{非笛卡尔几何}
规则矩形域可直接使用笛卡尔网格；曲线边界域若用阶梯状笛卡尔近似，边界误差往往过大。常见选择是广义坐标变换或直接在任意控制体上实施有限体积积分。前者适合结构曲线网格，后者更自然支持一般多边形/多面体。

% [原PDF p593]
局部法向坐标只用于黎曼求解，最终更新仍在全局守恒变量中完成；旋转前后动量通量的单位和面积缩放必须一致。

任意方向界面的黎曼问题可在局部正交坐标中求解。先把速度和动量旋转到法向—切向基，调用一维欧拉求解器取得法向质量、动量和能量通量，再把动量通量旋回全局坐标。面法向应取单位向量，面积大小另行乘入，避免重复缩放。

局部基的切向方向可以有多种选择，但同一面的两侧必须使用一致基。若交换面左右单元，法向与通量符号要同时改变；这种反对称性可用随机状态自动测试。边界面则由内部状态和物理边界构造外侧状态，再沿外法向求同样的局部黎曼问题。

\subsection{广义坐标变换}
建立从物理坐标 $(t,x,y,z)$ 到计算坐标 $(\tau,\xi,\eta,\zeta)$ 的可逆映射。链式法则把各物理导数表示成网格度量与计算导数的组合；守恒形式要求把 Jacobian 与度量因子一并吸收到变换后的守恒量和通量中。

% [原PDF p594]
非正交网格上的精度不仅取决于通量函数，还取决于梯度重构和面中心插值；网格质量因此是数值误差的一部分。

曲线坐标中的网格尺度不能只用坐标差代替。稳定时间步应根据单元体积与所有面上的最大法向信号速度组合估计；狭小、倾斜或高度扭曲的控制体可能成为全局限制。只取最短中心距会忽略面面积和法向投影，导致过于乐观或过分保守的步长。

重构梯度可由最小二乘或 Green--Gauss 方法获得，并在面中心外推。限制器应基于相邻单元可容许范围，同时考虑非正交几何；否则在歪斜网格上，即使一维限制器满足条件，面值仍可能越过局部极值或进入非物理状态。

变换后系统仍写成
\begin{equation}\widehat{\bm U}_{\tau}+\widehat{\bm F}_{\xi}+
\widehat{\bm G}_{\eta}+\widehat{\bm H}_{\zeta}=0.\tag{16.110}\end{equation}
帽变量包含物理守恒量、映射 Jacobian 和逆变换度量。可在规则计算网格上应用前述方法，但必须满足离散几何守恒律；否则即使均匀流也会因度量误差产生伪扰动。

% [原PDF p595]
任意网格实现最有力的守恒测试是逐面反对称性：交换左右状态并反转法向后，所得全局面通量应严格变号。

几何数据应在计算开始前进行一致性审计，包括正体积、闭合面积向量和共享面匹配。对任意封闭单元，所有外向面积向量之和应接近零；这一恒等式保证常通量的离散散度为零。若网格随时间运动，还必须满足相应的离散几何守恒律。

程序中最好以“面”为中心存储通量，而不是让每个单元独立重复计算。一次局部黎曼求解产生带固定法向的面通量，然后加到 owner、减到 neighbour。该数据结构天然保证守恒，也便于并行通信中明确共享面所有权和避免重复更新。

\subsection{非笛卡尔域的有限体积法}
有限体积法直接对任意控制体 $V$ 积分：单元平均的变化等于所有边界面外向通量之和。二维多边形每条边有长度、外法向和局部左右状态；三维多面体则使用面面积与法向。几何只通过这些面量进入守恒更新。

% [原PDF p596]
几何量和流动物理量应分层存储并分别验证；自由流不能保持时，应优先检查面积向量闭合、体积符号和共享面方向。

在非笛卡尔网格上，有限体积守恒式仍由控制体体积和各面通量积分组成。几何信息进入面面积、外法向和单元体积，而不是改变守恒定律本身。每个内部面必须由两个相邻单元共享完全相同的面积向量，仅符号相反，才能使内部通量严格抵消。

广义坐标变换把物理空间方程映射到计算空间，但度量系数必须满足离散几何守恒关系。若常数状态在静止曲线网格上不能保持不变，误差来源通常不是黎曼求解器，而是面向量、Jacobian 或度量导数的不一致离散。自由流保持应作为首要单元测试。

对二维单元，更新可写成
\begin{equation}
\bm U_V^{n+1}=\bm U_V^n-\frac{\Delta t}{|V|}
\sum_{s=1}^{N_f}|A_s|\bm F_{n,s},\tag{16.121a}
\end{equation}
其中 $\bm F_{n,s}$ 是第 $s$ 边沿外法向的数值通量。只要相邻单元对共享面使用同一个通量且法向相反，全局守恒便自动成立。

% [原PDF p597]
局部法向坐标只用于黎曼求解，最终更新仍在全局守恒变量中完成；旋转前后动量通量的单位和面积缩放必须一致。

任意方向界面的黎曼问题可在局部正交坐标中求解。先把速度和动量旋转到法向—切向基，调用一维欧拉求解器取得法向质量、动量和能量通量，再把动量通量旋回全局坐标。面法向应取单位向量，面积大小另行乘入，避免重复缩放。

局部基的切向方向可以有多种选择，但同一面的两侧必须使用一致基。若交换面左右单元，法向与通量符号要同时改变；这种反对称性可用随机状态自动测试。边界面则由内部状态和物理边界构造外侧状态，再沿外法向求同样的局部黎曼问题。

为复用一维黎曼求解器，在每个面建立局部正交坐标。旋转矩阵 $\bm T_s$ 把全局速度和动量变换为法向、切向分量；面左状态总取控制体内部，右状态取相邻单元或边界幽灵状态。

% [原PDF p598]
非正交网格上的精度不仅取决于通量函数，还取决于梯度重构和面中心插值；网格质量因此是数值误差的一部分。

曲线坐标中的网格尺度不能只用坐标差代替。稳定时间步应根据单元体积与所有面上的最大法向信号速度组合估计；狭小、倾斜或高度扭曲的控制体可能成为全局限制。只取最短中心距会忽略面面积和法向投影，导致过于乐观或过分保守的步长。

重构梯度可由最小二乘或 Green--Gauss 方法获得，并在面中心外推。限制器应基于相邻单元可容许范围，同时考虑非正交几何；否则在歪斜网格上，即使一维限制器满足条件，面值仍可能越过局部极值或进入非物理状态。

旋转后的守恒量 $\widehat{\bm U}=\bm T_s\bm U$ 满足局部法向一维欧拉结构。先求局部黎曼通量 $\widehat{\bm F}_s$，再用逆旋转
\begin{equation}\bm F_{n,s}=\bm T_s^{-1}\widehat{\bm F}_s\tag{16.128a}\end{equation}
还原到全局分量。压力和密度不因旋转改变，动量矢量按基变换处理。

% [原PDF p599]
任意网格实现最有力的守恒测试是逐面反对称性：交换左右状态并反转法向后，所得全局面通量应严格变号。

几何数据应在计算开始前进行一致性审计，包括正体积、闭合面积向量和共享面匹配。对任意封闭单元，所有外向面积向量之和应接近零；这一恒等式保证常通量的离散散度为零。若网格随时间运动，还必须满足相应的离散几何守恒律。

程序中最好以“面”为中心存储通量，而不是让每个单元独立重复计算。一次局部黎曼求解产生带固定法向的面通量，然后加到 owner、减到 neighbour。该数据结构天然保证守恒，也便于并行通信中明确共享面所有权和避免重复更新。

在矩形网格上，这一任意面公式严格退化为式(16.129)的 $x,y$ 通量差更新。书中示例给出带斜边多边形的面积、边长和外法向，说明每个边界贡献都按“面积乘法向通量”进入单元残差。

% [原PDF p600]
几何量和流动物理量应分层存储并分别验证；自由流不能保持时，应优先检查面积向量闭合、体积符号和共享面方向。

在非笛卡尔网格上，有限体积守恒式仍由控制体体积和各面通量积分组成。几何信息进入面面积、外法向和单元体积，而不是改变守恒定律本身。每个内部面必须由两个相邻单元共享完全相同的面积向量，仅符号相反，才能使内部通量严格抵消。

广义坐标变换把物理空间方程映射到计算空间，但度量系数必须满足离散几何守恒关系。若常数状态在静止曲线网格上不能保持不变，误差来源通常不是黎曼求解器，而是面向量、Jacobian 或度量导数的不一致离散。自由流保持应作为首要单元测试。

物理边界处的右状态由边界条件构造。固壁通常在局部坐标中反射法向速度并保留切向速度、密度和压力；远场或透射边界则按入射特征数目给定或外推变量。保持法向方向约定一致是避免通量符号错误的关键。

% [原PDF p601]
局部法向坐标只用于黎曼求解，最终更新仍在全局守恒变量中完成；旋转前后动量通量的单位和面积缩放必须一致。

任意方向界面的黎曼问题可在局部正交坐标中求解。先把速度和动量旋转到法向—切向基，调用一维欧拉求解器取得法向质量、动量和能量通量，再把动量通量旋回全局坐标。面法向应取单位向量，面积大小另行乘入，避免重复缩放。

局部基的切向方向可以有多种选择，但同一面的两侧必须使用一致基。若交换面左右单元，法向与通量符号要同时改变；这种反对称性可用随机状态自动测试。边界面则由内部状态和物理边界构造外侧状态，再沿外法向求同样的局部黎曼问题。

任意网格界面算法可归纳为三步：用 $\bm T_s$ 旋转左右状态；调用任一一维精确或近似黎曼求解器得到局部法向通量；逆旋转并乘面面积累加到左右单元，符号相反。该框架把复杂几何与一维物理解耦，便于复用 HLLC、Roe 或 Osher 等求解器。

% [原PDF p602]
\chapter{多维测试问题}
爆炸波的径向参考解包含几何源项，比较时须用相同的初始半径、状态方程和终止时刻，并把笛卡尔采样点转换为半径。

二维爆炸波问题把高压高密度圆区置于低压环境，初始间断向外产生激波，向内产生稀疏波，中间由接触面分隔。虽然解析解可由径向一维欧拉方程获得，笛卡尔网格计算并不显式使用径向对称，因此它是检验多维格式方向误差和对称保持能力的严格算例。

结果应沿通过中心的多条方向取样，并与径向参考解比较。若坐标轴与对角线上的激波位置不同，说明数值耗散或稳定条件具有网格方向偏置；若总质量与能量漂移，则应先检查边界通量和共享面符号，而不是把误差归因于圆形界面的阶梯化表示。

本章给出检验多维数值解的基准题。评估不能只看图像是否平滑，还应核对对称性、激波和接触面位置、守恒误差、网格方向效应以及随加密的收敛趋势。二维、三维问题常无完整解析解，因此可使用对称降维解、高分辨率参考解和不变量进行交叉验证。

% [原PDF p603]
中心附近、接触面和外行激波代表三种不同误差机制，应分别评估对称性、数值耗散和波速，而不能用单一指标概括。

爆炸初态的压力比很大，第一步预测最容易出现非物理状态。时间步应由全域最大法向信号速度确定，并在高阶重构后检查各面密度和压力。反射边界要翻转法向速度而保持切向速度，透射边界则延拓出流状态；边界距离必须足够大，避免返回波污染比较时刻。

二维曲线图之外，还应比较径向密度、速度、压力和内能。接触面通常比激波更易被数值耗散抹宽，而稀疏扇中的光滑区域可用于观察高阶精度。中心处速度应保持接近零，任意持续的非零偏置都可能来自网格不对称、边界处理或方向扫掠顺序。

\section{爆炸与内爆}
圆柱或球对称爆炸由中心高压区向外驱动激波，随后是接触面和向内传播的稀疏波；内爆则使强波向中心汇聚，对网格各向同性和极端压缩鲁棒性更敏感。用笛卡尔网格计算径向对称题，可以直接观察坐标方向偏置。

% [原PDF p604]
网格加密研究要保持物理区域与边界位置不变，并至少比较两个取样方向；这样才能判断圆对称误差是否真正收敛。

不同格式的比较要采用相同网格、CFL 数、终止时刻和色标。等值图若各自自动缩放，会掩盖过冲与耗散差异。除了视觉图像，还应报告径向剖面、守恒误差以及最大最小状态；这些量能区分真正的物理解结构和由限制器产生的数值伪影。

当使用分裂格式时，可在连续时间步交替扫掠次序并比较对称性；非分裂格式则重点检查横向预测。随着网格加密，圆形激波的轴向—对角线差异应减小，激波位置应收敛到径向参考解。若差异不随加密下降，需要重新检查算法而非仅增加分辨率。

\subsection{二维爆炸测试}
在 $2\times2$ 正方形域内，以中心为圆心给定内外两个常状态。典型数据在半径 $r_0$ 内取高压高密度，外部取低压低密度，初始速度为零，理想气体 $\gamma=1.4$。边界按问题设置透射或对称条件。

计算结果应呈同心圆形的外行激波、接触间断和稀疏结构。若数值耗散或网格各向异性过强，圆形波会沿坐标轴与对角线出现不同半径。

% [原PDF p605]
高压比测试首先考验物理可容许性，其次才是分辨率；若出现负状态，必须定位到预测或通量步骤并修正离散选择。

二维爆炸波问题把高压高密度圆区置于低压环境，初始间断向外产生激波，向内产生稀疏波，中间由接触面分隔。虽然解析解可由径向一维欧拉方程获得，笛卡尔网格计算并不显式使用径向对称，因此它是检验多维格式方向误差和对称保持能力的严格算例。

结果应沿通过中心的多条方向取样，并与径向参考解比较。若坐标轴与对角线上的激波位置不同，说明数值耗散或稳定条件具有网格方向偏置；若总质量与能量漂移，则应先检查边界通量和共享面符号，而不是把误差归因于圆形界面的阶梯化表示。

图17.2给出 $t=0.25$ 的密度分布：外侧为圆形激波，其后是圆形接触面，内部为稀疏波。可把穿过中心的径向剖面与等价一维圆柱对称问题比较，检验激波位置和峰值。二维等值线还显示方法对圆对称性的保持程度。

% [原PDF p606]
爆炸波的径向参考解包含几何源项，比较时须用相同的初始半径、状态方程和终止时刻，并把笛卡尔采样点转换为半径。

爆炸初态的压力比很大，第一步预测最容易出现非物理状态。时间步应由全域最大法向信号速度确定，并在高阶重构后检查各面密度和压力。反射边界要翻转法向速度而保持切向速度，透射边界则延拓出流状态；边界距离必须足够大，避免返回波污染比较时刻。

二维曲线图之外，还应比较径向密度、速度、压力和内能。接触面通常比激波更易被数值耗散抹宽，而稀疏扇中的光滑区域可用于观察高阶精度。中心处速度应保持接近零，任意持续的非零偏置都可能来自网格不对称、边界处理或方向扫掠顺序。

书中还比较随机选择法与等价激波管数据。随机选择法在一维能保持真间断，但多维方向分裂可能导致明显噪声和方向误差；适度耗散的守恒格式通常给出更稳定的圆形波。分辨率、扫掠顺序和随机采样都会影响结果。

% [原PDF p607]
中心附近、接触面和外行激波代表三种不同误差机制，应分别评估对称性、数值耗散和波速，而不能用单一指标概括。

不同格式的比较要采用相同网格、CFL 数、终止时刻和色标。等值图若各自自动缩放，会掩盖过冲与耗散差异。除了视觉图像，还应报告径向剖面、守恒误差以及最大最小状态；这些量能区分真正的物理解结构和由限制器产生的数值伪影。

当使用分裂格式时，可在连续时间步交替扫掠次序并比较对称性；非分裂格式则重点检查横向预测。随着网格加密，圆形激波的轴向—对角线差异应减小，激波位置应收敛到径向参考解。若差异不随加密下降，需要重新检查算法而非仅增加分辨率。

\subsection{三维爆炸测试}
三维测试在 $2\times2\times2$ 立方体中设置球形高压核心与低压外区，初始速度为零。精确解保持球对称，包含外行球形激波、球形接触面与内部稀疏波。三维计算应检查不同坐标轴、面对角线和体对角线上的径向剖面是否一致，并同时监测总质量与总能量守恒。
% ===== 译文分片 7（已合并）=====

% [原PDF p608]
球形爆炸也是检查多维算法旋转对称性的严格试验。笛卡尔网格并不具有连续球对称，若方向分裂误差、网格各向异性或边界处理过强，沿坐标轴和对角线提取的径向曲线会出现系统差异。把三维结果沿 $x$ 轴同独立的一维径向方程比较，不仅核对激波幅值和位置，也检验几何源项、球面积增长以及三方向通量组合是否正确。间断分辨率保持一维水平，说明主要耗散仍由基本通量决定，而非维数本身。
图 17.6 给出了球形爆炸在 $t=0.25$ 时、$z=0$ 平面上的压力分布。解由向外传播的球形激波和其后的球形稀疏波组成。图 17.7 又把三维 WAF 解沿与 $x$ 轴重合的径向直线取样，并同一维径向 RCM 解比较；符号与实线吻合良好。与二维情形一样，三维计算对间断的分辨率仍接近相应的一维问题，说明维数增加本身并未迫使格式额外涂抹波面。

\section{楔面上的激波反射}
激波研究涉及实验、理论和计算三个方面，国际激波研讨会及其论文集提供了大量资料。这里考察一列平面激波撞击倾斜楔面的经典问题。初始激波垂直于 $x$ 方向，位于 $x=x_S$，激波马赫数为 $M_S$；它沿 $x$ 方向传播，楔面与传播方向夹角为 $\phi$。这个构型看似简单，却能产生不同的多波相互作用结构，也是检验多维激波捕捉格式的重要标准算例。

% [原PDF p609]
激波反射形态通常分为四类：规则反射、单马赫反射、复杂马赫反射和双马赫反射。具体出现哪一种，主要由入射激波马赫数 $M_S$ 与楔角 $\phi$ 决定。规则反射中入射激波与反射激波在壁面相交；马赫反射则出现近似竖直的马赫杆、三重点和由三重点发出的滑移面。Ben--Dor 的专著系统讨论了这些波系，Heilig、Reichenbach、Dewey 与 Takayama 等人的工作则给出了实验和诊断资料。

图 17.8 的初始几何中，楔尖位于激波前方。边界上必须正确实现无穿透反射条件，入口和远场边界应保持激波前、后的给定状态。由于滑移面横跨密度与切向速度间断，而压力和法向速度近似连续，这个算例不仅检验激波位置，也能区分不同黎曼求解器对中间波的耗散程度。

% [原PDF p610]
这个算例的初值预处理十分关键。解析跳跃满足连续 Rankine--Hugoniot 条件，但离散 WAF--HLLC 激波通常占据若干网格单元，其剖面还受限制器影响。若让数值激波到达楔尖时才开始形成，启动声波会与物理反射混在一起。先预推进并重置，相当于把算法自身的稳态激波剖面作为初值，使后续差异主要来自真正的多维相互作用。AMR 计算还必须在粗细层之间保持同一激波位置，并通过通量修正保证守恒。
这里选取 $M_S=1.7$、$\phi=25^\circ$。计算域在 $x$--$y$ 平面上为 $25.0\times16.5$，楔尖位于 $x=4.69$，初始激波位于 $x=4.0$。激波前为环境状态
\[
\rho_a=1.225\ \mathrm{kg/m^3},\qquad p_a=1.01325\times10^5\ \mathrm{Pa},
\qquad u_a=v_a=0.
\]
给定 $M_S$ 后，激波后状态由 Rankine--Hugoniot 跳跃条件确定。

直接用解析跳跃初始化时，离散格式会在最初若干步建立自身的数值激波剖面，从而产生启动误差并污染细小结构。常用做法是先让激波传播一段距离，再用已经形成的数值剖面重置初值。这里采用二维预推进方案，让激波走过其初始位置到楔尖距离的 $90\%$ 后重置激波后状态。

计算采用按维分裂的 WAF 方法、HLLC 黎曼求解器和 van Leer 限制器，$C_{\mathrm{CFL}}=0.9$，最大波速按第 16 章的方法估计。规则曲线网格为 $800\times528$。另一个结果结合 CHIMERA 重叠网格、自适应网格加密和笛卡尔网格；基础网格为 $50\times33$，再加两级加密，每级加密因子为 4。

% [原PDF p611]
\begin{sourceblock}
\sourcecrop{611}{50}{130}{385}{615}{0.62}
\captionof{figure}{$M_S=1.7$、楔角 $25^\circ$ 的激波反射。上：Takayama 教授提供的实验纹影；下：细规则网格上 WAF--HLLC 方法得到的密度等值线。}
\end{sourceblock}
实验图与密度等值线都显示单马赫反射：入射激波、反射激波和马赫杆在三重点相交，滑移面由三重点向楔面延伸。数值解准确再现了主激波的曲率、马赫杆的高度以及反射波走向。比较时不应只看最外层激波，还应检查滑移面是否过度扩散、三重点附近是否产生非物理振荡，以及楔面边界是否保持几何一致性。

% [原PDF p612]
\begin{sourceblock}
\sourcecrop{612}{48}{130}{386}{615}{0.62}
\captionof{figure}{同一激波反射问题的 AMR--CHIMERA 结果。上：输出时刻的加密与重叠网格结构；下：WAF--HLLC 方法得到的密度等值线。}
\end{sourceblock}
自适应网格把高分辨率集中在入射激波、反射激波、马赫杆和滑移面附近，远离波系处保持粗网格。该结果说明守恒通量与网格间插值必须协同：若重叠区传递不守恒，即使局部网格很细，也可能在接缝处生成弱伪波。两级四倍加密使最细局部分辨率可与高密度规则网格比较，而单元总数显著减少。

% [原PDF p613]
马赫反射中的滑移面是接触与剪切的组合：法向质量通量和压力满足相容关系，密度及切向速度可以跳跃。HLL 只保留最左、最右两条包络波，把整个中间波扇压成一个状态，因此会直接消除滑移面的特征自由度；HLLC 增加接触波后可显著改善。Roe 和精确求解器也属于完整波模型，但在极端状态下还要关注熵修正、正性与真实气体模型。由此可见，算法选择与物理模型选择必须分开评估。
该算例的波型属于单马赫反射。三重点连接入射激波的一段、反射激波和沿楔面升起的马赫杆；从三重点发出的滑移面以尖锐角度接近楔面。实验与计算在定性上符合得很好。不过在很高的激波马赫数下，真实气体效应与黏性耗散变得重要，理想气体 Euler 方程将不足以描述实验；此时应改用 Navier--Stokes 方程并配合适当的真实气体状态方程。

现代激波捕捉方法通常都能得到总体波系，但滑移面是更敏感的判据。HLL 求解器的近似波模型不含接触波和剪切波，会把滑移面涂抹到不可接受的程度。完整黎曼求解器显式保留所有特征族，一般比不完整求解器、中心格式和通量矢量分裂方法更准确。因而选择求解器时，不能只比较强激波的位置，还要考察慢速接触与剪切结构。

% [原PDF p614]
\chapter{多维空间中的 FORCE 通量}
本章把第 7 章的一维中心 FORCE 数值通量推广到二维和三维的结构与非结构网格。所得一步守恒格式在相当宽的 Courant 数范围内具有单调性和线性稳定性。除基本一阶格式外，还将展示借助 ADER 思路在空间与时间上提高精度的结果。

\section{引言}
Lax--Wendroff 定理以及 Hou--LeFloch 的进一步结果表明，非线性双曲守恒律的数值方法必须保持守恒。因此，构造单调的单元界面数值通量是有限体积法和间断 Galerkin 有限元法的核心。数值通量可以用于全离散或半离散框架，也可以部署在结构或非结构网格上。

Lax--Friedrichs 通量与 Godunov 通量代表两条基本路线：前者是中心或对称方法，后者是迎风、基于黎曼问题或基于特征的方法。迎风格式沿单元面的法向求局部一维黎曼问题，显式利用方程中的波传播信息；中心格式虽不按传统方式求黎曼问题，却仍使用相同的微分方程与左右初值。

% [原PDF p615]
一般而言，Godunov 型迎风法有更高的潜在精度，中心法则更简单、更通用。迎风优势主要体现在相对网格移动缓慢的中间波，尤其是接触波和剪切波。然而这种优势高度依赖所用黎曼求解器。若近似波模型包含精确问题中的全部特征场，称为完整求解器；否则称为不完整求解器。

Rusanov 是单波不完整求解器，HLL 是双波不完整求解器；当系统含两个以上特征场时，HLL 无法表示全部波。精确求解器显然完整；HLLC 对一至三维 Euler 方程完整，Roe 与 Osher--Solomon 求解器对其目标系统也完整。传统通量矢量分裂在中间波分辨率方面类似不完整求解器，往往严重涂抹未被波模型表示的波。若没有合适的完整求解器，与其采用复杂但同样耗散的不完整迎风法，不如采用能在一般网格上保持守恒并可高阶扩展的中心方法。

FORCE 最初为一维双曲问题提出，与 Nessyahu--Tadmor 中心格式及 Arminjon 等人的工作相关。两者都在交错控制体上应用积分平均，但 Nessyahu--Tadmor 是交错网格上的两步过程，一维 CFL 上限为 $1/2$；FORCE 是非交错网格上的一步守恒格式，具有明确数值通量，一维 CFL 上限为 1。

% [原PDF p616]
中心格式的“通用”并不是说它完全不受波速影响。稳定时间步仍需最大传播速度，正性保护也可能依赖局部物理状态；它省去的是在每个界面构造完整特征分解和波型。对于磁流体、弹塑性、多相混合等特征结构复杂的系统，这种简化可显著降低实现难度。另一方面，若关心慢接触或材料界面，完整迎风求解器仍可能以同样网格得到更锐利结果，高阶重构也不能完全消除基础通量的波模型差异。
Nessyahu--Tadmor 格式已推广到规则笛卡尔网格以及非结构三角形、四面体网格，并已有若干收敛性证明。对于 FORCE，Chen 与 Toro 证明了它在一维等熵气体动力学方程和含底坡源项的非线性浅水方程上的收敛性。相关中心格式文献十分丰富。

FORCE 与 Liska--Wendroff 复合格式之间也有值得注意的类比。复合格式在时间上交替使用 $L_a$ 步 Lax--Friedrichs 和 $L_b$ 步 Lax--Wendroff，以经验方式兼顾准确性与单调性；FORCE 则在通量层面组合二者，并严格取算术平均权重 $1/2$。这个权重并非随意选择：在不显式使用波传播信息的条件下，它给出了单调性区域最大的组合。

本章推广适用于二维、三维结构与非结构网格，并可在有限体积和间断 Galerkin 框架下用 ADER 方法扩展到任意时空阶。后续先回顾一维 FORCE，再构造笛卡尔网格上的二维和三维通量，分析单调性与线性稳定性，最后推广到一般网格并给出算例。

% [原PDF p617]
有限体积更新只要求同一界面通量以相反符号进入左右单元，因此局部守恒可直接推出全局守恒。通量还必须满足一致性：左右状态相同时退化为物理通量。Godunov 通量通过完整波扇自动满足这些要求，但代价取决于黎曼解的复杂度。FORCE 的目标是在保留两点一致性与守恒结构的同时，用显式代数组合替代特征分解和非线性迭代，使一个通量模板能服务多种双曲系统。
\section{一维 FORCE 回顾}
考虑 $m\times m$ 一维双曲守恒律
\begin{equation}
\partial_t\bm Q+\partial_x\bm F(\bm Q)=0, \tag{18.1}
\end{equation}
有限体积更新为
\begin{equation}
\bm Q_i^{n+1}=\bm Q_i^n-\frac{\Delta t}{\Delta x}
\left(\bm F_{i+1/2}-\bm F_{i-1/2}\right). \tag{18.2}
\end{equation}

\subsection{FORCE 及相关通量}
Godunov 方法以左右常状态 $\bm Q_i^n,\bm Q_{i+1}^n$ 构造局部黎曼问题。其相似解记为 $\bm Q_{i+1/2}(x/t)$，由 $m$ 个波族分隔 $m+1$ 个常状态。界面沿 $x/t=0$，所以 Godunov 通量为
\begin{equation}
\bm F_{i+1/2}^{\mathrm{God}}=\bm F\!\left(\bm Q_{i+1/2}(0)\right). \tag{18.4}
\end{equation}
复杂系统的精确黎曼解往往需要迭代求解非线性方程，实际计算通常使用近似求解器。

% [原PDF p618]
LF 项的跳跃系数与网格速度 $\Delta x/\Delta t$ 成正比，以较强人工黏性换取稳健性；LW 项先预测半时间界面状态，再计算物理通量，提供低耗散的二阶修正。算术平均不需要波速符号，却能在标量问题上得到单调范围。实现时必须以同一左右状态分别计算两项再平均，不能把非线性物理通量错误地移出平均，因为一般有 $\bm F((\bm Q_L+\bm Q_R)/2)\ne[\bm F(\bm Q_L)+\bm F(\bm Q_R)]/2$。
中心两点通量写成 $\bm F_{i+1/2}=\bm H(\bm Q_i^n,\bm Q_{i+1}^n)$，并满足一致性 $\bm H(\bm V,\bm V)=\bm F(\bm V)$。经典 Lax--Friedrichs 通量为
\begin{equation}
\bm F_{i+1/2}^{\mathrm{LF}}=\frac12[\bm F(\bm Q_i^n)+\bm F(\bm Q_{i+1}^n)]
-\frac12\frac{\Delta x}{\Delta t}(\bm Q_{i+1}^n-\bm Q_i^n). \tag{18.6}
\end{equation}
它对模型方程在 Courant 数不超过 1 时线性稳定，标量情形还具有单调性。

两步 Lax--Wendroff 通量为
\begin{align}
\bm F_{i+1/2}^{\mathrm{LW}}&=\bm F(\bm Q_{i+1/2}^{\mathrm{LW}}),\\
\bm Q_{i+1/2}^{\mathrm{LW}}&=\frac12(\bm Q_i^n+\bm Q_{i+1}^n)
-\frac12\frac{\Delta t}{\Delta x}[\bm F(\bm Q_{i+1}^n)-\bm F(\bm Q_i^n)]. \tag{18.7}
\end{align}
它线性稳定到 Courant 数 1，时空二阶，但并不单调。FORCE 通量恰为两者的算术平均
\begin{equation}
\bm F_{i+1/2}^{\mathrm{FORCE}}=\frac12
(\bm F_{i+1/2}^{\mathrm{LW}}+\bm F_{i+1/2}^{\mathrm{LF}}). \tag{18.8}
\end{equation}
这种结构把低耗散但可能振荡的 LW 通量与稳健而耗散的 LF 通量结合起来。

% [原PDF p619]
数值黏性式可由修正方程得到。固定 $c$ 时，增大 $\omega$ 会线性削弱二阶扩散，却逐步接近可能产生过冲的 LW 极限。若 $c\to0$，LF 黏性按 $1/c$ 增长，说明在固定空间网格上过度减小时间步不一定改善中心格式的总误差，反而可能增加单位物理时间内累计的扩散。比较算法必须使用相同网格和相近 CFL 数，不能只看单步结果或名义时间阶。
\subsection{单调性与数值黏性}
对常系数平流方程
\[
q_t+\lambda q_x=0,\qquad c=\lambda\Delta t/\Delta x,
\]
考虑加权组合
\begin{equation}
f_{i+1/2}^{(\omega)}=\omega f_{i+1/2}^{\mathrm{LW}}
+(1-\omega)f_{i+1/2}^{\mathrm{LF}},\qquad0\le\omega\le1. \tag{18.14}
\end{equation}
它产生的三点格式数值黏性系数为
\begin{equation}
\alpha=\frac12\frac{(1-c^2)(1-\omega)}{c}\lambda\Delta x. \tag{18.16}
\end{equation}
因此黏性随 $\omega$ 线性下降：$\omega=0$ 是黏性最大的 LF，$\omega=1$ 是黏性为零的 LW，$\omega=1/2$ 则恢复 FORCE。这个参数平面清楚揭示了“中心耗散”与“迎风信息”之间的定量关系。

% [原PDF p620]
图 18.1 的横轴是局部 Courant 数，纵轴是 LW 权重。单调边界随 $|c|$ 增大而下降。FORCE 的水平线 $\omega=1/2$ 完全位于单调区闭包内，因而不必逐界面估计特征速度即可保持标量单调性。Godunov 曲线则给出每个 $c$ 下仍保持单调的最大低耗散权重。这个几何解释说明固定权重 $1/2$ 的选择并非经验折中，而是无显式迎风信息时的最优边界值。
在 $(c,\omega)$ 单位正方形中，常权重 $0,1/2,1$ 分别对应 LF、FORCE 和 LW。若允许权重依赖 Courant 数，取
\begin{equation}
\omega_g(c)=\frac{1}{1+|c|}, \tag{18.17}
\end{equation}
加权通量就与 Godunov 迎风通量完全一致。曲线 $\omega_g(c)$ 把参数平面分成两区：曲线上方为非单调格式，下方为单调格式。所有不依赖 $c$ 且 $0\le\omega\le1/2$ 的常权重格式均单调；其中 FORCE 位于边界，具有这类纯中心组合中最小的数值黏性。

\begin{sourceblock}
\sourcecrop{620}{132}{435}{305}{615}{0.27}
\captionof{figure}{加权 FORCE 型通量在 $c$--$\omega$ 平面中的单调区。Godunov 曲线以上为非单调区，以下为单调区。}
\end{sourceblock}

% [原PDF p621]
所谓额外 $19\%$ 的黏性降低，是在整个 Courant 数区间对参数平面面积积分所得，并非每个 $c$ 上都固定减少同一比例。对非线性系统同时存在接触波、声波和材料波，没有唯一局部 Courant 数；把一个最大或代表速度代入 $\omega_g$ 只能近似 Godunov。MUSTA 通过若干局部预测--校正阶段逐渐逼近迎风通量，但仍保持无需显式特征分解的设计目标，适合传统黎曼求解器难以构造的系统。
当 $1/2<\omega\le1$ 时，常权重格式只能在越来越小的 CFL 数下保持单调。定义
\[
T_g=\{(c,\omega):0\le c\le1,;1/2<\omega\le(1+|c|)^{-1}\},
\]
并把采用 $\omega_g$ 的方法称为广义 FORCE（GFORCE 或 Godunov--FORCE）。对线性平流，它就是 Godunov 方法；对非线性标量方程和系统，则需选择代表性的局部 Courant 数，因而使用少量迎风信息。

若把单位正方形面积解释为 LF 与 LW 之间的总数值黏性差，FORCE 相对 LF 减少 $50\%$ 黏性。Godunov 的黏性系数
\begin{equation}
\alpha_{\mathrm{God}}=\frac12(1-|c|)\lambda\Delta x \tag{18.19}
\end{equation}
又比 FORCE 平均减少约 $19\%$，这是单调方法类中可达到的最大额外收益，来自迎风信息。GFORCE 试图在非线性问题中逼近这部分收益；对系统只能部分实现，还可用 MUSTA 多级策略进一步降低黏性。

% [原PDF p622]
含源项的总方程在本节先取齐次形式，是为了独立构造对流通量。若直接保留源项，半步交错平均与主网格更新都要加入一致的体积积分；只在最终式末端加一项会破坏高阶平衡。笛卡尔单元四个面分别沿本地法向计算，$x$ 面使用 $\bm F$，$y$ 面使用 $\bm G$。推导借助交错体，但最终仍是标准非交错有限体积更新，便于同边界条件和并行面循环结合。
FORCE 的关键优点是完全中心：不需要特征值、特征矢量或传统黎曼求解器。本节把它推广到含源项的一般多维系统
\begin{equation}
\partial_t\bm Q+\partial_x\bm F(\bm Q)+\partial_y\bm G(\bm Q)
+\partial_z\bm H(\bm Q)=\bm S(\bm Q), \tag{18.20}
\end{equation}
并首先令 $\bm S=0$，在二维、三维笛卡尔网格上构造守恒格式。

\section{笛卡尔网格上的 FORCE 格式}
\subsection{二维情形}
二维更新写为
\begin{equation}
\bm Q_{i,j}^{n+1}=\bm Q_{i,j}^{n}
-\frac{\Delta t}{\Delta x}(\bm F_{i+1/2,j}-\bm F_{i-1/2,j})
-\frac{\Delta t}{\Delta y}(\bm G_{i,j+1/2}-\bm G_{i,j-1/2}). \tag{18.21}
\end{equation}
任务是为单元四条边构造通量。FORCE 同时使用主网格和交错网格：主网格保存数值单元平均，交错控制体跨越每条边，用积分守恒律定义半时间状态。这里采用边基交错网格，而不是顶点基交错网格。

% [原PDF p623]
二维主单元 $C_{i,j}$ 的四条边按逆时针定向，使 $C_{i,j}$ 始终位于有向边左侧。把本单元重心、相邻单元重心与公共边两个顶点连接，就得到跨越该边的四边形交错体。记它在 $C_{i,j}$ 内、外的部分为 $S^-_{i+1/2,j}$ 与 $S^+_{i+1/2,j}$，公共边为
\[
S^0_{i+1/2,j}=S^-_{i+1/2,j}\cap S^+_{i+1/2,j},
\]
其余三个方向作类似定义。上标“负”表示有向边左侧，“正”表示右侧，这个约定使一般网格上的面法向与左右状态定义保持一致。

构造分三步：第一步在每个跨边交错体内推进半个时间步，得到交错平均；第二步以这些半时刻平均为数据，在主单元内再推进半步；第三步用代数整理把两步结果写成非交错的一步守恒形式，并识别出每个界面的 FORCE 数值通量。

% [原PDF p624]
图中主网格保存单元平均，虚线交错体只承担局部半步演化。每个主单元值同时进入四个交错预测，相邻单元仅通过共同面耦合。笛卡尔网格上左右半体面积相等，所以初始交错平均是算术平均；一般网格必须改为体积加权。逆时针面定向保证法向符号统一，也是把界面通量装配为守恒残差的基本约定。
对含 $n_k$ 条边的任意多边形控制体 $S_k$，积分守恒律在时间 $\delta t$ 上给出
\begin{equation}
\bm Q_k^1=\bm Q_k^0-\frac{\delta t}{|S_k|}
\sum_{l=1}^{n_k}\int_{v_l}^{v_{l+1}}
[\cos\theta_l\bm F(\bm Q)+\sin\theta_l\bm G(\bm Q)]\,ds, \tag{18.23}
\end{equation}
其中 $\bm n_l=(\cos\theta_l,\sin\theta_l)$ 是第 $l$ 边外法向，$v_{n_k+1}=v_1$，$|S_k|$ 为面积。把这个积分式应用到四个边基交错体，并令 $\delta t=\Delta t/2$，即可得到半时刻的边状态。几何平均项来自交错体左右两半面积相等，通量差项来自外边界积分。

\begin{sourceblock}
\sourcecrop{624}{50}{285}{385}{615}{0.36}
\captionof{figure}{二维笛卡尔主网格与边基交错网格。每个交错体跨越主单元的一条公共边。}
\end{sourceblock}

% [原PDF p625]
半步预测的方向结构来自积分几何，而不是简单忽略横向传播。$x$ 面交错体中切向贡献在规则笛卡尔几何上成对抵消，四个面状态再由主单元更新共同耦合。若网格非正交或面切向长度不匹配，这种简化不再成立，必须回到完整边界积分，逐面使用真实法向与几何测度。最终非交错公式无需全局存储交错网格，这显著降低了内存和通信复杂度。
四个半时刻平均为
\begin{align}
\bm Q_{i+1/2,j}^{n+1/2}&=\tfrac12(\bm Q_{i,j}^n+\bm Q_{i+1,j}^n)
-\frac{\Delta t}{\Delta x}[\bm F(\bm Q_{i+1,j}^n)-\bm F(\bm Q_{i,j}^n)],\\
\bm Q_{i,j+1/2}^{n+1/2}&=\tfrac12(\bm Q_{i,j}^n+\bm Q_{i,j+1}^n)
-\frac{\Delta t}{\Delta y}[\bm G(\bm Q_{i,j+1}^n)-\bm G(\bm Q_{i,j}^n)], \tag{18.24}
\end{align}
另两条边由指标反向得到。第二阶段在 $C_{i,j}$ 中再推进 $\Delta t/2$：四个交错平均先以 $1/4$ 相加，再减去半步的 $x$、$y$ 面通量差。这一两阶段表达式虽然直观，却还不是最终实现形式。

经整理后恢复式（18.21），例如右侧面的通量由半时刻物理通量、左右旧时刻物理通量平均和跳跃耗散项组成。重要的是，最终只需旧层相邻单元值，不必存储全局交错网格；交错体只是推导局部界面通量的几何工具。

% [原PDF p626]
维数因子 $\alpha=2$ 来自主单元由四个半交错体重新平均，而不是人为稳定参数。LW2 预测使用有效尺度 $2\Delta t$，LF2 耗散使用相应的 $\Delta x/(2\Delta t)$，组合仍与原守恒律量纲一致。若仅在二维外层分别调用保持 $\alpha=1$ 的一维 FORCE 通量，得到的是按维分裂算法，不等同于这里的无分裂多维 FORCE；两者的稳定区域、横向耦合和数值耗散均不同。
二维 FORCE 通量仍可写成两个分量通量的平均：
\begin{equation}
\bm F_{i+1/2,j}=\frac12(\bm F_{i+1/2,j}^{\mathrm{LW2}}+
\bm F_{i+1/2,j}^{\mathrm{LF2}}). \tag{18.27}
\end{equation}
其中
\begin{align}
\bm F^{\mathrm{LW2}}_{i+1/2,j}&=\bm F(\bm Q^{\mathrm{LW2}}_{i+1/2,j}),\\
\bm Q^{\mathrm{LW2}}_{i+1/2,j}&=\tfrac12(\bm Q_{i,j}^n+\bm Q_{i+1,j}^n)
-\tfrac12\frac{2\Delta t}{\Delta x}[\bm F(\bm Q_{i+1,j}^n)-\bm F(\bm Q_{i,j}^n)], \tag{18.29}\\
\bm F^{\mathrm{LF2}}_{i+1/2,j}&=\tfrac12[\bm F(\bm Q_{i,j}^n)+\bm F(\bm Q_{i+1,j}^n)]
-\tfrac12\frac{\Delta x}{2\Delta t}(\bm Q_{i+1,j}^n-\bm Q_{i,j}^n). \tag{18.30}
\end{align}
与一维式相比，唯一差别是在时间步中出现空间维数因子 $\alpha=2$。其他各面通量按相应方向完全类似。非笛卡尔但结构化的二维网格也可沿同一积分几何直接推广。

\subsection{三维情形}
三维守恒更新再加入 $z$ 方向通量差，维数参数变为 $\alpha=3$。

% [原PDF p627]
三维主单元有六个面，每面都对应一个跨面交错体。$y$、$z$ 通量由 $x$ 面公式循环置换坐标和物理通量获得。各向异性网格中 $\Delta x,\Delta y,\Delta z$ 必须分别进入对应方向，而时间步全局共享。边界面应先由物理条件构造外侧虚状态，再调用同一 FORCE 面函数，以维持内部面与边界面的离散结构一致；不能在更新后直接覆盖守恒变量来冒充通量边界。
一般维数参数 $\alpha$ 下，某个 $x$ 面的通量为
\begin{align}
\bm F_{i+1/2}^{\mathrm{FORCE}\alpha}&=\tfrac12
(\bm F_{i+1/2}^{\mathrm{LW}\alpha}+\bm F_{i+1/2}^{\mathrm{LF}\alpha}), \tag{18.32}\\
\bm Q_{i+1/2}^{\mathrm{LW}\alpha}&=\tfrac12(\bm Q_i^n+\bm Q_{i+1}^n)
-\tfrac12\frac{\alpha\Delta t}{\Delta x}[\bm F(\bm Q_{i+1}^n)-\bm F(\bm Q_i^n)], \tag{18.34}\\
\bm F_{i+1/2}^{\mathrm{LF}\alpha}&=\tfrac12[\bm F(\bm Q_i^n)+\bm F(\bm Q_{i+1}^n)]
-\tfrac12\frac{\Delta x}{\alpha\Delta t}(\bm Q_{i+1}^n-\bm Q_i^n). \tag{18.35}
\end{align}
三维时分别对 $x,y,z$ 六个面计算该通量。

\section{FORCE 格式的性质}
为理解多维通量，先把 $\alpha$ 当作参数放回一维守恒更新。这产生一族有趣的一维 FORCE$\alpha$ 格式；$\alpha=1,2,3$ 分别对应从一、二、三维推导得到的通量。随后再直接分析真正二维和三维格式的单调性与线性稳定性。

% [原PDF p628]
\subsection{一维解释}
当 $\alpha=1$ 时，分量通量就是经典 LF 与 LW，平均后为标准 FORCE。$\alpha=2$ 时，LW 型通量对应 Godunov 中心格式，其线性稳定范围为 $|c|\le1/(2\sqrt2)$，却仅在 $1/2\le|c|\le1/(2\sqrt2)$ 的矛盾窄域内单调，因而单独使用价值有限；相应 LF 型通量在完整稳定区 $|c|\le1/2$ 内单调，且耗散小于经典 LF。

模型平流方程上的 FORCE$\alpha$ 三点格式满足统一限制
\begin{equation}
|c|\le \frac{\sqrt{2\alpha-1}}{\alpha}. \tag{18.41}
\end{equation}
在此范围内它既单调又线性稳定。上限随维数参数 $\alpha$ 增大而下降，但同时单个方向通量的数值耗散也发生变化。这个一维解释有助于把真正的多维稳定条件同方向 Courant 数联系起来。

\subsection{数值试验}
下面用理想气体一维 Euler 方程、$\gamma=1.4$ 比较不同 $\alpha$ 的效果。

% [原PDF p629]
静止接触计算到 $t=5$，很长的输出时间会放大每一步的累计耗散，因此比短时激波管更能区分不同通量。Woodward--Colella 问题则产生强激波、接触波与稀疏波的反复作用，压力跨越五个数量级，主要检验强非线性稳健性。参考 Godunov 解虽为一阶，但完整精确黎曼求解器保留每个波族，在单调一阶方法中具有最低黏性，因而是合理基准。
第一个试验是第 10 章测试 6 的静止接触间断。定义域为 $[0,1]$，在 $x=1/2$ 处密度从左侧 $1.4$ 跳到右侧 $1$，速度处处为零，压力处处为 1；边界透射，网格 100 个单元，计算到 $t=5$。除 van Leer 格式用 $C_{\mathrm{CFL}}=0.6$ 外，各 FORCE$\alpha$ 取式（18.41）上限的 $90\%$。

随着 $\alpha=1,2,3,15$ 增大，接触面数值黏性明显降低。更值得注意的是，不完整的 Rusanov 与 HLL 求解器在这个慢速中间波试验中，比各 FORCE$\alpha$ 中心格式更耗散；van Leer 通量矢量分裂虽优于 $\alpha=1$，仍不及 $\alpha=2,3$。这再次说明“迎风”本身不保证精细分辨率，近似波模型是否包含接触场才是关键。

第二个 Woodward--Colella 试验在 $[0,1]$ 上设置三段压力 $1000,0.01,100$，分界点为 $0.1$ 和 $0.9$，初始密度为 1、速度为零，两端反射；使用 3000 个单元，输出 $t=0.038$。它无闭式精确解，以精确黎曼求解器的一阶 Godunov 解为参考。

% [原PDF p630]
强波相互作用还揭示单调性和物理可容许性并非同一概念。标量线性模型上的非负系数能阻止新极值，但 Euler 系统中还要求密度和内能保持正值；高压比碰撞可能把中间预测状态推到可容许集之外。实际 FORCE$\alpha$ 实现应在计算物理通量前验证状态，并通过满足理论条件的时间步或专门的保正高阶限制器处理，而不应把负密度简单截断。参考曲线的差异要同时从峰值、波位置和守恒量误差判断。
Woodward--Colella 问题含两个强激波的多次碰撞，是检验稳健性与耗散的严苛算例。FORCE$\alpha$ 均比参考 Godunov 解更扩散，但当 $\alpha$ 增大时，密度峰值和窄结构逐渐接近参考结果。这种趋势与静止接触试验一致：多维推导引入的参数改变了 LF 与 LW 两部分的相对耗散尺度。

不过不能仅凭一维参数试验任意增大 $\alpha$。真正多维计算中，$\alpha$ 由空间维数决定，稳定范围受各方向 Courant 数共同约束；若把它当作自由调参，必须重新验证单调性、正性和非线性稳定性。图 18.4、18.5 的主要目的，是说明多维 FORCE 通量在一维投影下的性质，而不是建议普遍使用极大的人工维数。

\subsection{多维分析}
考虑常系数三维平流方程
\begin{equation}
q_t+\lambda_1q_x+\lambda_2q_y+\lambda_3q_z=0, \tag{18.42}
\end{equation}
并定义 $c_x=\lambda_1\Delta t/\Delta x$、$c_y=\lambda_2\Delta t/\Delta y$、$c_z=\lambda_3\Delta t/\Delta z$。

% [原PDF p631]
\begin{sourceblock}
\sourcecrop{631}{75}{145}{330}{600}{0.58}
\captionof{figure}{静止接触试验。上：不同 $\alpha$ 的 FORCE$\alpha$；下：同 Rusanov、HLL 与 van Leer 通量分裂的比较。}
\end{sourceblock}
上图的精确解是保持在 $x=0.5$ 的垂直密度跳跃。曲线展宽直接度量接触耗散：FORCE1 最宽，FORCE2、FORCE3 依次变陡，FORCE15 最接近精确间断。下图显示 Rusanov 与 HLL 因波模型缺失接触场而扩散更强。图中的 distance、density 分别表示位置和密度；符号与线型只用于区分算法。

% [原PDF p632]
\begin{sourceblock}
\sourcecrop{632}{65}{345}{375}{590}{0.32}
\captionof{figure}{Woodward--Colella 强波相互作用：不同 FORCE$\alpha$ 解与精确黎曼求解器的一阶 Godunov 参考解。}
\end{sourceblock}
把三个方向 FORCE 通量代入守恒式，可写成邻域值的线性组合
\begin{equation}
q_{i,j,k}^{n+1}=\sum_{l=-1}^{1}\sum_{m=-1}^{1}\sum_{r=-1}^{1}
\beta_{lmr}q_{i+l,j+m,k+r}. \tag{18.45}
\end{equation}
实际只有中心点和六个轴向相邻点的系数非零。要求全部系数非负，就得到单调性条件
\begin{equation}
c_x^2+c_y^2\le\frac12\quad(\text{二维}),\qquad
c_x^2+c_y^2+c_z^2\le\frac13\quad(\text{三维}). \tag{18.47}
\end{equation}

% [原PDF p633]
\begin{sourceblock}
\sourcecrop{633}{105}{385}{325}{590}{0.27}\\[0.45em]
\sourcecrop{633}{105}{125}{325}{325}{0.27}
\captionof{figure}{二维与三维 FORCE 格式的线性稳定及单调区域。三维图显示 $c_x$--$c_y$ 平面上的截面。}
\end{sourceblock}
二维允许区域是半径 $1/\sqrt2$ 的圆盘；三维允许区域是半径 $1/\sqrt3$ 的球，其平面截面如图所示。条件不仅限制每个方向的 Courant 数，还限制它们的平方和。因而在对角传播时，不能让各方向都取一维上限。经验 Fourier 稳定性分析得到与单调性完全相同的区域，这意味着这些一阶格式在该区域内同时具备线性稳定和凸组合性质。

% [原PDF p634]
第 16 章介绍的经验线性稳定性方法进一步确认式（18.47）。图 18.6 对应二维 $c_x$--$c_y$ 平面，图 18.7 是三维稳定球的一个截面。实际程序可先按各单元的局部特征速度计算方向 Courant 数，再取满足平方和限制的全局时间步；对非均匀网格则应使用每个单元的几何尺度。

\section{一般网格上的 FORCE 格式}
先考虑二维非结构三角形网格。主单元 $C_i$ 有三个共享边邻居 $C_{i,1},C_{i,2},C_{i,3}$。与笛卡尔情形相同，需要保存平均值的主网格，以及用于定义面通量的边基交错网格。对每条公共边，连接边的两个端点与左右单元重心，形成跨边四边形；其左右两部分分别记作 $S^-_{i,j}$ 和 $S^+_{i,j}$。

% [原PDF p635]
非结构推广的关键在几何一致性。一个公共面在左右单元中必须具有相同面积和相反单位法向，交错半体体积也要按同一顶点顺序计算。高度扭曲或退化单元会使半体总体积很小，放大半步通量差，因此时间步应结合局部内切尺度限制。边基构造保持每个通量只依赖面两侧单元，模板紧致，且天然适合网格分区后的面通信。
有向面按逆时针遍历，$C_i$ 位于左侧。公共面及交错体定义为
\begin{equation}
S_{i,j}=S^-_{i,j}\cup S^+_{i,j},\qquad
S^0_{i,j}=C_i\cap C_{i,j},\qquad j=1,2,3. \tag{18.48--18.49}
\end{equation}
对 $\alpha$ 维计算域 $\Omega\subset\mathbb R^\alpha$，设共形剖分为 $\mathcal C_\Omega=\bigcup_iC_i$。这一构造不依赖三角形：二维多边形、三维四面体或一般多面体均可对每个平面面建立相应的双半体。重心连接只用于定义局部几何体积与法向，不改变主网格的守恒拓扑。

\begin{sourceblock}
\sourcecrop{635}{85}{390}{350}{615}{0.35}
\captionof{figure}{三角形主网格及边基交错体。虚线四边形跨越每条公共边，由左右单元重心与边端点组成。}
\end{sourceblock}

% [原PDF p636]
一般网格中，式（18.54）的分母是 $|S^-_{i,j}|+|S^+_{i,j}|$，左右旧状态按各自半体体积贡献加权。法向通量差写成 $[\overline{\bm F}(\bm Q_j)-\overline{\bm F}(\bm Q_i)]\cdot\bm n_j$，所以旋转网格只需改变法向而无需重写方程。第二阶段各半体面积加权和应精确恢复 $|C_i|$，这是一致保持常状态的离散几何恒等式；若几何数据不闭合，均匀自由流也会产生伪残差。
若 $C_i$ 有 $n_i$ 个平面面，则
\[
\partial C_i=\bigcup_{j=1}^{n_i}S^0_{i,j},\qquad
|C_i|=\sum_{j=1}^{n_i}|S^-_{i,j}|.
\]
积分守恒律给出主单元更新
\begin{equation}
\bm Q_i^{n+1}=\bm Q_i^n-\frac{\delta t}{|C_i|}
\sum_{j=1}^{n_i}\int_{S^0_{i,j}}\overline{\bm F}(\bm Q)\cdot\bm n_j\,dA,
\quad \overline{\bm F}=[\bm F,\bm G,\bm H]. \tag{18.53}
\end{equation}
第一阶段以左右单元平均为初值，在每个交错体中推进半步。若左右半体体积不等，加权平均必须使用 $|S^-|$ 与 $|S^+|$；通量差则乘公共面面积和外法向。第二阶段用这些半时刻状态在主单元中再推进半步，得到显式两步交错格式。

% [原PDF p637]
一般网格 LF 型耗散项的几何系数可理解为跨面交错体的有效“网格速度”。左右半体不对称时，物理通量平均也不再是简单的二分之一，而由相反半体体积配权。这确保在交换左右单元并反转法向后，数值法向通量严格变号。实现层面最好把每个内部面作为唯一对象，保存面积、单位法向、左右单元和半体几何；一次计算通量后同时累加两个残差，可从数据结构上杜绝不守恒的舍入差异。
经 Gauss 定理和面法向归一化整理后，得到一步守恒式
\begin{equation}
\bm Q_i^{n+1}=\bm Q_i^n-\frac{\Delta t}{|C_i|}
\sum_{j=1}^{n_i}|S^0_{i,j}|\,
\overline{\bm F}^{\mathrm{FORCE}\alpha}_{i,j}\cdot\bm n_j, \tag{18.56}
\end{equation}
其中
\begin{equation}
\overline{\bm F}^{\mathrm{FORCE}\alpha}_{i,j}=\frac12
(\overline{\bm F}^{\mathrm{LW}\alpha}_{i,j}+
\overline{\bm F}^{\mathrm{LF}\alpha}_{i,j}). \tag{18.57}
\end{equation}
LW 型通量是在第一阶段得到的界面中间状态上计算物理通量张量；LF 型通量则由左右物理通量的几何加权平均减去与 $\bm Q_j^n-\bm Q_i^n$ 成正比的耗散项。所有几何因子都由公共面面积和两侧半体体积确定。

这一形式十分适合非结构数据结构：对每个内部面只计算一次左右状态、法向和 FORCE 通量，再以相反符号累加到两个单元，天然保证成对守恒。边界面则由物理边界条件构造外侧虚状态。无需特征分解使它适合复杂系统，但时间步仍应按局部网格尺度与最大传播速度限制。

% [原PDF p638]
从一阶 FORCE 到三阶 ADER--FORCE，需要的不只是高阶插值。每个时间步先以 WENO 从非结构邻域恢复空间多项式，再在单元内作局部时空预测，最后在每个面上以足够阶的时空求积计算 FORCE 通量。若只提高空间重构而仍用一阶时间通量，整体会受时间误差限制；反之，高阶预测配低阶几何求积也会降阶。非结构网格还要求曲边或面法向积分与解多项式使用相容阶数。
\section{数值结果}
用第 17 章的 $M_S=1.7$、$\phi=25^\circ$ 马赫反射检验非结构网格 FORCE。求解 $\gamma=1.4$ 的二维理想气体 Euler 方程，采用三角形网格；既计算基本一阶 FORCE，也在有限体积框架中结合 WENO 非线性重构与 ADER 时间推进形成高阶格式。

图 18.10 使用网格尺度 $1/20$ 和 $1/50$，均为三阶 ADER--FORCE；图 18.11 使用更细的 $1/100$ 网格，对比一阶 FORCE 与三阶扩展。高阶解在相同网格上显著改善滑移面、三重点和反射激波的清晰度，同时维持非结构网格上的守恒。

\section{结语}
因此，验收高阶实现时既要检查形式阶数，也要检查逐界面守恒、导数尺度、求积精度和源项耦合。只有这些环节同时一致，局部高阶预测才能真正转化为全局高阶数值解。
多维 FORCE 是一般网格上简单的一阶守恒通量，在较宽 Courant 数范围内单调且线性稳定。它可作为有限体积或间断 Galerkin 高阶方法的底层通量，既支持半离散也支持全离散实现。最大的工程优势是不需要系统特征结构或传统黎曼求解器，因此尤其适合特征分解困难的复杂双曲系统。

% [原PDF p639]
\begin{sourceblock}
\sourcecrop{639}{55}{175}{385}{600}{0.58}
\captionof{figure}{$M_S=1.7$、楔角 $25^\circ$ 的激波反射：三阶 ADER--FORCE 在尺度 $1/20$（上）和 $1/50$（下）的非结构网格与密度等值线。}
\end{sourceblock}
粗网格已能恢复主波系，细网格则使马赫杆、反射激波和滑移面更加清晰。三角形方向变化没有在光滑区留下明显网格印迹，说明面法向通量与几何权重配合正确。观察非结构结果时还应核查公共面的通量是否只计算一次并反号累加，否则极小的不守恒误差会沿长时间传播积累。

% [原PDF p640]
\begin{sourceblock}
\sourcecrop{640}{55}{175}{385}{600}{0.58}
\captionof{figure}{尺度 $1/100$ 网格上的密度解。上：一阶 FORCE；下：三阶 ADER--FORCE。}
\end{sourceblock}
一阶格式保持稳定，却明显扩散滑移面及三重点附近细节；三阶格式在相同几何与基本通量上通过高阶重构和时空预测恢复锐利结构。比较显示，高阶扩展不是简单减小耗散系数，而是在光滑区提高时空一致精度，并由非线性重构在间断附近抑制振荡。至此，第 18 章的推导、性质分析和非结构算例形成完整闭环。

% [原PDF p641]
\chapter{广义黎曼问题}
本章研究含源项非线性双曲偏微分方程组的广义黎曼问题。它从两方面推广经典黎曼问题：界面两侧初值不再是常量，而是各自光滑的任意向量场；控制方程还允许代数源项。文献中也称其为导数黎曼问题或高阶黎曼问题。本书记为 $\mathrm{GRP}_K$，非负整数 $K$ 表示近似阶数。

$\mathrm{GRP}_K$ 的界面解可用于构造时空均为 $K+1$ 阶的全离散一步有限体积或间断 Galerkin 方法。考虑平衡律
\begin{equation}
\partial_t\bm Q+\partial_x\bm F(\bm Q)=\bm S(\bm Q),\qquad
\bm Q(x,0)=\bm Q_0(x). \tag{19.1}
\end{equation}
$\bm Q$ 为守恒变量，$\bm F$ 为通量，$\bm S$ 为不含未知量微分项的源项。

\section{引言}
在空间单元 $I_i=[x_{i-1/2},x_{i+1/2}]$ 和时间区间 $[t^n,t^{n+1}]$ 上积分方程，是高阶有限体积构造的起点。

% [原PDF p642]
式（19.2）是高阶方法的守恒骨架，而不是某一种具体格式。只要左右单元对公共面的 $\bm F_{i+1/2}$ 使用完全相同的数值，内部通量在全域求和时成对抵消；源项则保留为体积贡献。空间重构必须保持式（19.3）的平均，否则更新的未知量语义会发生改变。界面通量和源项分别是时间一维积分与时空二维积分，它们的求积阶数都至少应达到目标格式阶数。

高阶单元分布并不作为新的独立自由度保存，而是每步从邻域平均重构。因此重构、局部预测、求积和守恒更新形成一个闭环：重构恢复空间信息，广义黎曼问题把它转成界面时间信息，单元预测给出源项所需的内部时空信息，最后只更新平均值。这种组织既保持有限体积的鲁棒守恒性质，又允许任意时空阶。
空间--时间控制体积分给出精确关系
\begin{equation}
\bm Q_i^{n+1}=\bm Q_i^n-\frac{\Delta t}{\Delta x}
(\bm F_{i+1/2}-\bm F_{i-1/2})+\Delta t\,\bm S_i, \tag{19.2}
\end{equation}
其中
\begin{align}
\bm Q_i^n&=\frac1{\Delta x}\int_{x_{i-1/2}}^{x_{i+1/2}}\bm Q(x,t^n)\,dx, \tag{19.3}\\
\bm F_{i+1/2}&=\frac1{\Delta t}\int_{t^n}^{t^{n+1}}
\bm F(\bm Q(x_{i+1/2},t))\,dt,\\
\bm S_i&=\frac1{\Delta t\Delta x}\int_{t^n}^{t^{n+1}}
\int_{x_{i-1/2}}^{x_{i+1/2}}\bm S(\bm Q(x,t))\,dx\,dt. \tag{19.4}
\end{align}
单元平均虽是常量，积分中的解在单元内可以有高阶分布；这种分布由一组单元平均经守恒重构得到。

数值方法要给出时间平均界面通量和空间--时间平均源项的近似。为此需要两类预测：界面上的时间函数 $\bm Q_{LR}(\tau)$，以及单元内部的时空函数 $\bm Q_i(x,t)$。再用适当阶数的求积计算式（19.4），就能更新单元平均。

% [原PDF p643]
经典 Godunov 通量只取 $t\to0^+$ 的界面状态，因而每步只利用分片常数信息；即使黎曼问题精确求解，整体仍是一阶。要提高时间阶，不能只把左右状态换成高阶外推值后仍取瞬时通量，还必须知道整个 $[0,\Delta t]$ 上界面状态怎样变化。$\mathrm{GRP}_K$ 正是提供这条时间迹线。它允许波扇内部状态随时间改变，曲线特征由初值斜率和源项共同决定。

求完整半平面解通常远比数值通量所需信息昂贵。局部界面展开把问题压缩为有限个时间导数，因此计算量随目标阶数增长，却不随空间波扇采样点增长。二阶方法只需首个时间导数，更高阶继续加入局部系数。这个“只求所需函数值”的原则也是 ADER 能够成为一步方法的基础。
经典黎曼问题使用由一个间断分隔的左右常状态。第 4 章给出了 Euler 方程精确解，第 7 至 12 章讨论近似求解器，Godunov 一阶方法正建立在这个局部问题上。计算一阶界面通量只需知道相似解在 $x/t=0$ 的状态，而不必存储整个半平面解。

广义问题则允许界面两侧为任意光滑函数并包含源项。一般无法获得整个 $t>0$ 半平面的闭式解，但高阶一步 Godunov 方法只要求界面处随时间变化的局部解。分片线性初值对应 $\mathrm{GRP}_1$，可产生二阶格式；Kolgan、van Leer、Ben--Artzi 与 Falcovitz 等提出了不同的二阶构造。更高阶的分片光滑初值问题也由 Harten 及其他研究者从不同角度处理。

本书采用 $K$ 表示界面时间展开的最高阶，因此 $\mathrm{GRP}_K$ 产生 $K+1$ 阶的通量近似。必须区分“求整个波扇”和“只求界面短时间迹线”：后者信息量少得多，正是任意阶方法可实现的关键。

% [原PDF p644]
高阶方法的文献路线虽然表面不同，核心都在恢复界面短时间演化。Harten 方法先独立演化左右重构，再在时间节点求经典黎曼问题；Toro--Titarev 则先以经典黎曼解确定波拓扑，然后用导数黎曼问题直接构造界面 Taylor 系数。二者在线性常系数问题上等价，非线性时因冻结与演化次序不同而只保证同阶，不保证数值通量逐点相同。

含源项时还需处理平衡和刚性。显式 CK 递推把源项及其导数写入高阶系数，适合解析推导；局部隐式预测器则把反应过程直接包含在单元时空解中，更适合强刚性。无论采用何法，首项仍由齐次经典黎曼问题控制，因为有限源项在无穷小时间内不能改变初始跳跃的瞬时迎风选取。
Toro--Titarev 方法把含源项的完整 $\mathrm{GRP}_K$ 用于 ADER 有限体积格式。界面解先写成 $K$ 阶时间幂级数。首项由一个通常非线性的经典齐次黎曼问题给出；高阶时间导数通过 Cauchy--Kowalewski 过程改写成空间导数的函数。随后为各阶空间导数建立演化方程，并分别求经典线性黎曼问题，从而定义界面处的单侧导数状态。

于是一个 $K$ 阶广义问题归结为：求一个非线性经典黎曼问题得到首项，再求 $K$ 个空间导数的线性黎曼问题得到高阶系数。若初值在整个邻域光滑，界面没有跳跃，这套过程退化为普通 Cauchy--Kowalewski 构造；这在单元内部计算高阶源项时尤其有用。

界面时间函数一旦得到，就能对物理通量作高阶时间积分，形成时空任意阶的一步 Godunov 格式。全离散意味着每个时间步只进行一次重构和一次局部预测，不再依赖多级 Runge--Kutta 外推。

% [原PDF p645]
阶数定义强调“非平凡导数”的最高次数。例如左右均为常数时 $K=0$；若一侧线性、另一侧常数，则仍为 $K=1$；若使用 WENO 得到四次多项式，则局部问题属于 $\mathrm{GRP}_4$。左右函数不要求在原点连续，函数值和每阶导数都可有跳跃。每个跳跃在不同特征场中的传播方向由局部 Jacobian 决定，不能用一个标量平均代替。

短时间内波拓扑与 $\mathrm{GRP}_0$ 相同是一项局部假设：在下一次波相互作用、特征交叉或真空形成之前成立。有限体积 CFL 条件把时间步限制在相邻波尚未跨越多个单元的尺度，正好与这一假设匹配。每步重新重构并重启广义黎曼问题，也避免把局部级数越过其解析有效区间。
\section{问题陈述}
广义黎曼问题定义为
\begin{equation}
\left\{
\begin{aligned}
&\partial_t\bm Q+\partial_x\bm F(\bm Q)=\bm S(\bm Q),\quad x\in\mathbb R, t>0,\\
&\bm Q(x,0)=
\begin{cases}\bm Q_L(x),&x<0,\\\bm Q_R(x),&x>0.
\end{cases}
\end{aligned}\right. \tag{19.5}
\end{equation}
两侧向量函数在各自半轴上光滑，并具有至多 $K$ 阶非平凡空间导数。若 $K_L,K_R$ 分别为左右多项式的最高阶，则 $K=\max(K_L,K_R)$。$\mathrm{GRP}_0$ 的一阶及以上导数均为零，即经典分片常数黎曼问题；$\mathrm{GRP}_1$ 对应分片线性初值。

经典问题的解只依赖 $x/t$，特征线为直线，波之间是常状态扇形区域。$K>0$ 时波间状态随空间和时间变化，特征线弯曲，通常不再具有全局相似性。不过在足够短时间内，假设其波拓扑仍与首项的经典黎曼解相同，可把广义解视为经典波扇的高阶扰动。

% [原PDF p646]
经典波扇的每个常状态区域由直特征或激波轨迹分隔，整个结构随 $x/t$ 自相似。广义数据中，界面外的初始梯度使波两侧状态随时间演化，特征速度也随状态变化，因此轨迹弯曲。尽管如此，在第一瞬间每条波仍从经典黎曼问题的相同特征族发出。半解析方法正利用这一点：用 $\mathrm{GRP}_0$ 固定零阶拓扑，再把斜率、曲率和源项视为沿该拓扑传播的高阶扰动。
\begin{sourceblock}
\sourcecrop{646}{120}{390}{315}{610}{0.30}
\captionof{figure}{典型 $3\times3$ 双曲系统的经典黎曼问题 $\mathrm{GRP}_0$：上为分片常数初值，下为直特征组成的相似波系。}
\end{sourceblock}
本章目标不是求整个半平面，而是求原点界面上的短时间函数 $\bm Q_{LR}(\tau)$。对高阶 Godunov 有限体积法，它足以通过时间积分生成 $K+1$ 阶数值通量。另一方面，源项需要单元时空控制体内部的高阶近似 $\bm Q_i(x,t)$，再对 $\bm S(\bm Q_i)$ 作空间--时间积分。

经典分片常数情形会恢复 Godunov 一阶迎风法。广义情形同时承担两项任务：界面局部问题控制通量，单元内部光滑 Cauchy 问题控制源项。二者都用同一偏微分方程把时间导数转换为空间导数，因此可在同一时空阶上匹配，避免通量高阶而源项低阶造成整体降阶。

% [原PDF p647]
Cauchy--Kowalewski 构造同时回答存在性与系数可计算性。理论要求数据解析，以保证形式级数确实收敛到唯一局部解；数值算法只保留有限项，并把它作为一个时间步内的局部预测。对间断界面不能直接套用单一解析展开，因为左右极限不同；广义黎曼方法通过首项非线性波扇和导数黎曼问题把两个单侧解析展开连接起来。单元内部没有跳跃，则可直接使用普通 CK 过程计算源项预测。
\begin{sourceblock}
\sourcecrop{647}{120}{385}{315}{610}{0.30}
\captionof{figure}{广义黎曼问题 $\mathrm{GRP}_K$：上为界面两侧分片光滑初值，下为弯曲特征组成的短时间波系。}
\end{sourceblock}
\section{Cauchy--Kowalewski 定理}
该定理说明：若偏微分方程、初值及所涉及函数在展开点附近均解析，则初值问题存在唯一解析解。证明同时给出构造方法，即在原点作幂级数展开，反复使用控制方程和初值计算所有系数。这一递推称为 Cauchy--Kowalewski 过程；在数值分析语境中也常称 Lax--Wendroff 过程。

“无穷次可微”并不自动等于“解析”。解析要求 Taylor 级数在某邻域收敛且确实等于原函数。数值应用通常使用有限阶多项式重构，因此局部解析条件自然满足；真正需要谨慎的是间断界面，因为那里必须以左右单侧数据和黎曼问题代替单一 Taylor 展开。

% [原PDF p648]
多指标 $\bm\alpha=(\alpha_1,\ldots,\alpha_n)$ 同时记录每个自变量的微分次数，总阶为 $|\bm\alpha|=\sum_j\alpha_j$。截断到 $|\bm\alpha|\le K$ 才得到完整 $K$ 阶多变量 Taylor 多项式，不能只保留每个方向的纯导数而遗漏混合导数。时空预测中 $xt$、$x^2t$ 等混合项正是多维展开的一部分，它们通过 PDE 被换成初始空间导数。

解析性比 $C^\infty$ 更强，存在无穷可微但 Taylor 级数不等于原函数的反例。数值方法并不需要无穷级数真实收敛到全局解，只需有限阶局部截断在网格细化时给出正确余项。因此多项式重构为算法提供了可计算的有限导数，理论解析条件则解释这些递推在局部为何具有唯一意义。
\subsection{级数展开与解析函数}
若 $f\in C^\infty(I)$，在 $x_0\in I$ 的 Taylor 级数为
\begin{equation}
\sum_{k=0}^{\infty}\frac{f^{(k)}(x_0)}{k!}(x-x_0)^k. \tag{19.6}
\end{equation}
只有当该级数在 $x_0$ 的邻域内收敛到 $f(x)$ 时，$f$ 才在 $x_0$ 解析。指数函数 $e^x=\sum_{k\ge0}x^k/k!$ 和多项式是熟悉的解析函数。

对 $n$ 变量函数，多指标展开写为
\begin{equation}
f(\bm x)=\sum_{\bm\alpha\ge0}
\frac{\partial_{x_1}^{\alpha_1}\cdots\partial_{x_n}^{\alpha_n}f(\bm x^{(0)})}
{\alpha_1!\cdots\alpha_n!}
\prod_{j=1}^n(x_j-x_j^{(0)})^{\alpha_j}. \tag{19.7}
\end{equation}
二维二阶展开包含常数、两个一阶项、两个纯二阶项和一个混合项；混合项系数不带 $1/2$，因为其多指标阶乘为 $1!1!$。

% [原PDF p649]
定理中的 $G(x,t,q,q_x)$ 可显含空间和时间，因此也涵盖变系数和某些非自治源项。构造系数时，先由初值确定所有纯空间导数，再由 PDE 得到 $q_t$；继续对 PDE 求 $x$ 或 $t$ 导数，逐层消去新的时间导数。每一层只依赖已经确定的低阶项，形成三角递推。对方程组，标量乘法变成矩阵与高阶张量收缩，但逻辑不变。

数值实现必须保持混合偏导交换的一致性，例如 $q_{xt}=q_{tx}$。手写高阶公式容易漏掉链式法则产生的组合系数，通常应使用符号代数、自动微分或局部弱形式预测。无论工具如何，最终系数都应通过线性平流、常状态和已知解析解做单元测试。
\subsection{定理的一个说明}
考虑一阶方程
\begin{equation}
q_t=G(x,t,q,q_x),\qquad q(x,0)=q_0(x). \tag{19.10--19.11}
\end{equation}
若 $q_0$ 在原点附近解析，$G$ 在 $(0,0,q_0(0),q_0'(0))$ 附近解析，则原点附近唯一解析解为
\begin{equation}
q(x,t)=\sum_{\alpha_1,\alpha_2\ge0}
\frac{\partial_x^{\alpha_1}\partial_t^{\alpha_2}q(0,0)}
{\alpha_1!\alpha_2!}x^{\alpha_1}t^{\alpha_2}. \tag{19.12}
\end{equation}
所有系数通过反复微分 PDE 并在 $t=0$ 代入初值获得。对高阶数值方法而言，这把未知的时间导数系统地换成重构多项式可直接提供的空间导数。

\subsection{Cauchy--Kowalewski 方法}
下面先在标量守恒律上演示二阶构造，再把同一思想用于单元内部的时空预测和源项积分。

% [原PDF p650]
考虑
\[
q_t+\partial_x f(q)=0,qquad q(x,0)=q_0(x).
\]
二阶局部展开为
\begin{equation}
q=q(0,0)+q_t(0,0)t+q_x(0,0)x
+\tfrac12q_{xx}(0,0)x^2+q_{tx}(0,0)xt+\tfrac12q_{tt}(0,0)t^2. \tag{19.15}
\end{equation}
记 $f^{(k)}=d^kf/dq^k$。由方程和链式法则，
\begin{align}
q_t&=-f'(q)q_x, \tag{19.19}\\
q_{tx}&=-[f'(q)q_{xx}+f''(q)q_x^2], \tag{19.20}
\end{align}
而 $q_{tt}$ 再对第一式求时间导数并代回 $q_t,q_{tx}$ 得到。最终每个时间或混合导数都只依赖 $q_0(0)$ 及其空间导数。这种符号替换在高阶时迅速复杂，实际实现通常用自动代数或局部预测器生成。

这里 $f'(q)$ 是标量特征速度；对方程组则变成通量 Jacobian。应注意，书中印刷式（19.21）的组合项必须按链式法则实现，不能把 $f$ 与 $f'$ 的角色混淆。数值代码最好通过递推或自动微分验证每一阶导数。

% [原PDF p651]
图 19.3 所列的左右导数并不是把同一函数在原点作双值延拓，而是两段独立光滑数据的单侧极限。对守恒变量向量，每阶都是一个完整向量；若系统有 $m$ 个分量，$K$ 阶问题包含 $2m(K+1)$ 个单侧标量数据。首项通过非线性波扇耦合所有分量，高阶导数再在首项 Jacobian 的特征基中传播。

界面时间多项式只需要在固定空间位置 $x=0$ 评估。移动激波可能很快离开该位置，随后界面状态落入波扇的另一常区；但 CFL 足够小时，局部展开在一次求积区间内仍可准确描述这段变化。若某个波在步内正好穿过界面，求积节点和波拓扑处理必须与基础黎曼求解器一致。
\section{一种求解方法}
回到 $\mathrm{GRP}_K$。左右初值 $\bm Q_L(x),\bm Q_R(x)$ 在界面外光滑，例如各分量为最高 $K$ 次多项式。界面 $x=0$ 处函数及其空间导数一般不连续，对每个分量 $q$ 有
\[
\Delta^{(k)}q=q_R^{(k)}(0)-q_L^{(k)}(0),\qquad k=0,\ldots,K.
\]
每一对单侧导数将作为一个新的经典黎曼问题的左右常状态。

\begin{sourceblock}
\sourcecrop{651}{90}{150}{340}{320}{0.30}
\captionof{figure}{$\mathrm{GRP}_K$ 的分片光滑初值及界面两侧从函数值到 $K$ 阶空间导数的数据列。}
\end{sourceblock}
求解目标是仅在固定界面上构造时间幂级数，而不是追踪完整弯曲波系。短时间展开的适用半径由附近波相互作用、非线性强度和解析性共同限制，所以每个有限体积时间步都重新构造局部展开。

% [原PDF p652]
首项忽略源项并不意味着算法是算子分裂。它只是利用 $\tau=0^+$ 的渐近事实：有限 $\bm S$ 对状态的改变量为 $O(\tau)$，所以常数项由齐次跳跃问题决定；从线性项开始，$\bm S$ 便通过 CK 关系进入。这样可避免为每种源项重新构造非齐次黎曼求解器，同时仍在统一时间多项式中耦合通量与反应。

$\bm D^{(0)}(0)$ 可以来自精确或近似经典求解器。若近似求解器缺失某些中间波，首项的相应分辨率缺陷也会传递到高阶方法；提高 $K$ 不能弥补错误的波模型。因此 ADER 的高阶时间精度与基础黎曼求解器的物理完备性是两个独立维度。
界面解写为
\begin{equation}
\bm Q_{LR}(\tau)=\bm Q(0,0^+)+
\sum_{k=1}^{K}\left[\partial_t^{(k)}\bm Q(0,0^+)\right]\frac{\tau^k}{k!}. \tag{19.22}
\end{equation}
\subsection{首项}
首项只反映左右界面外推值的瞬时相互作用。忽略源项，以 $\bm Q_L(0),\bm Q_R(0)$ 为常初值求齐次经典问题
\begin{equation}
\bm Q_t+\bm F(\bm Q)_x=0,qquad
\bm Q(x,0)=\begin{cases}\bm Q_L(0),&x<0,\\\bm Q_R(0),&x>0.
\end{cases} \tag{19.23}
\end{equation}
若相似解为 $\bm D^{(0)}(x/t)$，则
\begin{equation}
\bm Q(0,0^+)=\bm D^{(0)}(0). \tag{19.24}
\end{equation}
这个界面值称为 Godunov 状态。源项在零时间极限不改变首项，但会进入所有高阶时间系数。

% [原PDF p653]
CK 映射 $\bm P^{(k)}$ 一般是非线性的：除最高阶空间导数外，还包含低阶导数的乘积、通量 Jacobian 的导数以及源项导数。只知道 $\partial_x^k\bm Q$ 并不足够，必须同时保存从零阶到 $k$ 阶的全部 Godunov 状态。递推实现可按总阶逐层构造，避免重复计算公共子表达式。

对线性平流，所有阶只乘 $(-\lambda)^k$，这既是理论简化，也是代码校验基准。若程序在该问题上没有精确恢复同一上游多项式的平移，通常意味着导数阶乘、局部坐标尺度或左右特征投影出现错误。特别是重构若使用归一化坐标，转换回物理导数时要乘正确的 $\Delta x^{-k}$。
\subsection{高阶项}
计算 $\partial_t^{(k)}\bm Q(0,0^+)$ 分三步。第一步用 Cauchy--Kowalewski 过程写成
\begin{equation}
\partial_t^{(k)}\bm Q=\bm P^{(k)}
(\partial_x^{(0)}\bm Q,\partial_x^{(1)}\bm Q,\ldots,
\partial_x^{(k)}\bm Q), \tag{19.26}
\end{equation}
其中 $\bm P^{(k)}$ 已包含源项。线性平流 $q_t+\lambda q_x=0$ 的结果特别简单：$\partial_t^kq=(-\lambda)^k\partial_x^kq$。

问题转化为给界面空间导数赋值。初始时刻只有单侧导数
\[
\bm Q_L^{(k)}(0)=\lim_{x\to0^-}\frac{d^k\bm Q_L}{dx^k},\qquad
\bm Q_R^{(k)}(0)=\lim_{x\to0^+}\frac{d^k\bm Q_R}{dx^k}.
\]
不能简单平均这两个值，因为不同特征场携带的信息方向不同；必须为导数建立演化问题并在 $x/t=0$ 取迎风状态。

% [原PDF p654]
导数方程右端 $\bm H^{(k)}$ 包含两类项：原始物理源项求导产生的贡献，以及非线性通量求导后低阶导数之间的乘积。它们在有限时间演化中不可忽略，但为决定 $t=0^+$ 的迎风空间导数状态，方法只保留冻结后的齐次主部。随后 $\bm P^{(k)}$ 又把这些影响加入时间系数，因而该简化是界面状态选择而非物理源项删除。

常系数线性系统是最透明的情况：$\bm A$ 不随状态变化，$\bm H^{(k)}=0$，每阶导数确实满足原系统。Burgers 例子中的 $-(q_x)^2$ 说明非线性会产生导数反应项；它可能使梯度有限时间爆破，对应原方程中激波形成，也限制局部解析展开的时间半径。
对原系统求 $k$ 次空间导数并整理，得到
\begin{equation}
\partial_t(\partial_x^{(k)}\bm Q)+\bm A(\bm Q)
\partial_x(\partial_x^{(k)}\bm Q)=\bm H^{(k)},\qquad
\bm A=\frac{\partial\bm F}{\partial\bm Q}. \tag{19.29}
\end{equation}
$\bm H^{(k)}$ 依赖 $\bm Q$ 及至多 $k$ 阶空间导数，并包含源项。常系数齐次线性系统中 $\bm H^{(k)}=0$，所有阶导数满足原方程。线性平流同样如此；但非线性 Burgers 方程的一阶导数满足
\[
(q_x)_t+q(q_x)_x=-(q_x)^2,
\]
可见导数方程一般带有由非线性产生的源项。

为了形成实用界面算法，在第一瞬间先忽略 $\bm H^{(k)}$，再把 Jacobian 冻结在首项 Godunov 状态上。这个简化不等于从最终时间导数中删除物理源项；源项仍通过 $\bm P^{(k)}$ 进入时间展开。

% [原PDF p655]
冻结在 Godunov 状态而不是简单左右平均，是方法保留非线性波相互作用信息的关键。首项已经由完整或近似非线性黎曼问题选择了界面所处的波区，在该状态计算 $\bm A$，高阶小扰动便沿正确的局部特征方向传播。若用算术平均，强激波两侧可能得到错误特征符号，导致导数从错误一侧进入界面。

线性导数黎曼问题的解析解可用特征投影表示：把左右导数分解到 $\bm A_{LR}^{(0)}$ 的特征向量上，正特征值取左系数，负特征值取右系数，零速场采用统一约定，再合成 $\bm D^{(k)}(0)$。所有 $k$ 共用同一矩阵分解，因此增加阶数主要增加向量投影，而非新的非线性求解。
冻结矩阵定义为
\begin{equation}
\bm A_{LR}^{(0)}=\bm A(\bm Q(0,0^+)). \tag{19.31}
\end{equation}
对每个 $k=1,ldots,K$ 求线性齐次黎曼问题
\begin{equation}
\left\{
\begin{aligned}
&\partial_t(\partial_x^{(k)}\bm Q)+\bm A_{LR}^{(0)}
\partial_x(\partial_x^{(k)}\bm Q)=0,\\
&\partial_x^{(k)}\bm Q(x,0)=
\begin{cases}\bm Q_L^{(k)}(0),&x<0,\\\bm Q_R^{(k)}(0),&x>0.
\end{cases}
\end{aligned}\right. \tag{19.32}
\end{equation}
所有导数问题共享同一特征分解，因此矩阵只需分解一次。若其相似解为 $\bm D^{(k)}(x/t)$，则
\begin{equation}
\partial_x^{(k)}\bm Q(0,0^+)=\bm D^{(k)}(0). \tag{19.33}
\end{equation}
这也称第 $k$ 阶导数的 Godunov 状态。把它们代入 $\bm P^{(k)}$，即可得到全部时间系数。

% [原PDF p656]
源项体积积分必须体现空间变化。若只在单元中心建立时间展开，通常只能准确积分低阶或特殊对称数据；高阶方案需要多个空间求积点，或者一个能在整个单元评价的时空多项式。求积权重之和应精确恢复常源项，预测器应保持常状态。对具有稳态平衡的方程，还需让通量梯度与源项在离散层面抵消，即所谓良平衡性质；一般 CK 高阶一致性本身并不自动保证精确良平衡。
\subsection{源项}
计算式（19.4）的体积源项时，单元内部没有界面跳跃，重构多项式提供所有初始空间导数。可以选择若干空间求积点 $x_j$，在每点作时间 Taylor 展开，再用 Cauchy--Kowalewski 过程把时间导数换成空间导数；也可直接构造完整的空间--时间幂级数，并消去全部时间和混合导数。随后用相应阶数的求积计算 $\bm S(\bm Q)$ 的时空平均。

\subsection{数值通量与数值源项：步骤汇总}
第一步，以 $\bm Q_L(0),\bm Q_R(0)$ 求经典齐次黎曼问题，得到首项 $\bm D^{(0)}(0)$。第二步对每阶时间导数作 CK 替换，建立空间导数方程，忽略其瞬时右端并在 $\bm D^{(0)}(0)$ 处线性化；随后求 $K$ 个导数黎曼问题。第三步组装界面时间多项式。第四步对物理通量作时间求积。第五步在各单元内构造光滑时空预测并对源项作体积求积。

% [原PDF p657]
更明确地，线性化导数问题给出 $\bm D^{(k)}(0)$ 后，界面解为
\begin{equation}
\bm Q_{LR}(\tau)=\bm C_0+\bm C_1\tau+\cdots+\bm C_K\tau^K, \tag{19.40}
\end{equation}
其中
\begin{equation}
\bm C_0=\bm D^{(0)}(0),\qquad
\bm C_k=\frac1{k!}\bm P^{(k)}
(\partial_x^{(0)}\bm Q,\ldots,\partial_x^{(k)}\bm Q)\big|_{(0,0^+)}. \tag{19.41}
\end{equation}
因此 $\mathrm{GRP}_K$ 总共需要 $K+1$ 个经典齐次黎曼问题：首项问题一般非线性，高阶导数问题则共享一个常系数矩阵的线性系统。该分解把“非线性波拓扑”和“高阶局部变化”分开处理，既保留首项的正确波向，也避免每阶都解非线性问题。

实现时可缓存 $\bm A_{LR}^{(0)}$ 的左右特征向量。每个导数跳跃先投影到特征空间，再按特征值正负选择左右分量，最后变换回守恒变量。若有零特征值，需要采用与基础黎曼求解器一致的界面约定，避免不同阶对静止波作相反选择。

% [原PDF p658]
算法成本可粗略分成一次首项非线性黎曼求解、一次 Jacobian 特征分解、$K$ 组导数投影以及通量求积。与每阶重复求非线性问题相比，共享特征结构显著节省成本。对于没有解析特征分解的系统，可用数值线性求解或局部 DG 预测替代，但必须确保结果仍达到所需阶数。若基础求解器本身只返回通量而不返回界面状态，还需要扩展接口，因为高阶系数依赖 $\bm Q(0,0^+)$ 及其 Jacobian。
数值通量由
\[
\bm F_{i+1/2}=\frac1{\Delta t}\int_0^{\Delta t}
\bm F(\bm Q_{LR}(\tau))\,d\tau
\]
按所需阶数求积；数值源项则由单元内时空预测对 $\bm S$ 的体积积分得到。

\subsection{若干说明}
首项可使用第 4 章的精确 Euler 黎曼求解器，也可使用第 7 至 12 章的近似求解器或更一般的 MUSTA 类方法。多数常见系统的 $K$ 个线性导数问题可解析求解，而且只需一次特征结构。常系数线性系统上该过程有严格依据；非线性系统目前主要依靠数值证据支持，但包括源项的试验表现良好。

该方法也能作为分片光滑平流--反应方程的短时间半解析参考解。在每个初值间断 $x=x_j$ 建立局部广义黎曼问题，计算小时间 $\tau$ 的界面值，可用于验证高阶代码。必须牢记幂级数只在局部有效，时间过长、波相互作用或解析半径耗尽后，应重新初始化而不能继续外推。

\section{例子}
以下将方法应用于标量平流--反应方程 $q_t+f(q)_x=s(q)$，从无源线性问题逐步过渡到含非线性源项的 Burgers 方程。

% [原PDF p659]
线性平流例子说明整个广义构造的本质。经典首项决定信息来自哪一侧，高阶导数问题对每一阶作相同的迎风选择，CK 关系再把空间 Taylor 多项式沿特征方向转成时间 Taylor 多项式。由于速度恒定，不同阶之间没有非线性耦合，也没有源项组合，所得式子应与精确平移完全一致。

当 $\lambda=0$ 时界面状态不会由平流决定，需按所采用的零波速约定处理；对守恒通量 $f(q)=\lambda q$，通量本身为零，所以任何一致状态都不影响更新。方程组中零特征场却可能对应静止接触，应采用能保持该波的基础求解器约定。
\subsection{线性平流方程}
考虑常速、无源广义黎曼问题
\begin{equation}
q_t+\lambda q_x=0,qquad q(x,0)=
\begin{cases}q_L(x),&x<0,\\q_R(x),&x>0,
\end{cases} \tag{19.43}
\end{equation}
左右函数光滑，例如最高 $K$ 次多项式。界面解写成
\[
q_{LR}(\tau)=q(0,0^+)+\sum_{k=1}^{K}\partial_t^kq(0,0^+)\frac{\tau^k}{k!}.
\]
首项的经典黎曼问题仍是常速平流，因此
\begin{equation}
q(0,0^+)=
\begin{cases}q_L(0),&\lambda>0,\\q_R(0),&\lambda<0.
\end{cases} \tag{19.47}
\end{equation}
这就是最直接的迎风选择：正速度从左侧取信息，负速度从右侧取信息。

% [原PDF p660]
式（19.48）可由归纳法验证：每对时间求导一次，就用 $q_t=-\lambda q_x$ 多引入一个 $-\lambda$ 和一个空间导数。导数黎曼问题之所以与原方程相同，是因为常数 $\lambda$ 与微分算子交换。若速度依赖 $x$ 或 $q$，这一交换性质消失，导数方程会出现额外项，也就需要第 19.4 节的冻结线性化。

数值通量需要积分的是 $f(q_{LR})=\lambda q_{LR}$，因此多项式可逐项精确积分。若使用高斯求积，节点数应足以积分目标次数；用不足阶数的求积会使本来精确的局部平移降阶。这是验证 ADER 通量模块最简单也最严格的算例之一。
CK 递推给出
\begin{equation}
\partial_t^kq=(-1)^k\lambda^k\partial_x^kq. \tag{19.48}
\end{equation}
而每个空间导数都满足同一个平流方程
\[
\partial_t(\partial_x^kq)+\lambda\partial_x(\partial_x^kq)=0.
\]
所以第 $k$ 个导数黎曼问题的界面状态为
\begin{equation}
d^{(k)}(0)=
\begin{cases}q_L^{(k)}(0),&\lambda>0,\\q_R^{(k)}(0),&\lambda<0.
\end{cases} \tag{19.52}
\end{equation}
所有阶数都从同一迎风侧取值。时间多项式系数为 $c_k=(-1)^k\lambda^kd^{(k)}(0)/k!$，这说明广义求解器在常系数问题上只是沿特征回溯后对上游多项式作 Taylor 展开。

% [原PDF p661]
式（19.55）实际上是 $q_L(-\lambda\tau)$ 或 $q_R(-\lambda\tau)$ 在原点的截断 Taylor 展开。正速度时特征从界面向后回溯到负半轴，故取左函数；负速度时回溯到正半轴，故取右函数。这里的左右选择由特征回溯而不是位置符号本身决定。

加入线性反应后，精确解多出 $e^{\beta t}$。若 $\beta<0$，幅值沿特征衰减；若 $\beta>0$，则增长并可能使显式时间步受到反应尺度限制。即使平流 CFL 很小，$|\beta|\Delta t$ 很大时显式 Taylor 截断也会失稳，因此刚性问题应使用局部隐式预测器。
完整界面解为
\begin{equation}
q_{LR}(\tau)=
\begin{cases}
q_L(0)+\displaystyle\sum_{k=1}^{K}(-1)^k\lambda^kq_L^{(k)}(0)\tau^k/k!,&\lambda>0,\\
q_R(0)+\displaystyle\sum_{k=1}^{K}(-1)^k\lambda^kq_R^{(k)}(0)\tau^k/k!,&\lambda<0.
\end{cases} \tag{19.55}
\end{equation}
这与精确解 $q(x,t)=q_0(x-\lambda t)$ 在 $x=0$ 的展开完全一致，从而验证了方法。若左右数据本身是 $K$ 次多项式，该式在波尚未遇到其他间断前就是精确界面值，通量时间积分也可解析完成。

\subsection{含源项的线性平流}
进一步考虑
\begin{equation}
q_t+\lambda q_x=\beta q,\qquad q(x,0)=q_0(x), \tag{19.57}
\end{equation}
其精确解为 $q(x,t)=q_0(x-\lambda t)e^{\beta t}$。通常取 $\beta<0$ 表示衰减，但推导不依赖其符号。首项仍由齐次平流问题决定，源项从一阶及更高时间导数进入。

% [原PDF p662]
式（19.61）就是算子 $(-\lambda\partial_x+\beta)^k$ 的二项展开，因为常系数平流算子与乘法反应算子可交换。每个 $n$ 表示在 $k$ 次时间递推中选择了 $n$ 次反应和 $k-n$ 次平流。变系数或非线性源项下算子不再交换，必须保留链式法则产生的更多交叉项。

界面导数仍由齐次平流方向选择，而反应通过组合系数作用于所选导数。这种做法在常系数例子中可与精确 $q_0(-\lambda t)e^{\beta t}$ 展开逐项核对。数值源项则不能只用界面值，必须在整个单元时空预测上积分，否则空间变化会被遗漏。
最初两阶时间导数为
\begin{align}
q_t&=-\lambda q_x+\beta q,\\
q_{tt}&=\lambda^2q_{xx}-2\lambda\beta q_x+\beta^2q. \tag{19.60}
\end{align}
一般式是二项展开
\begin{equation}
\partial_t^kq=\sum_{n=0}^{k}\binom{k}{n}\beta^n(-\lambda)^{k-n}
\partial_x^{k-n}q. \tag{19.61}
\end{equation}
空间导数满足原来的平流--反应方程，但构造界面瞬时导数状态时忽略右端，只求与无源情形相同的齐次导数黎曼问题。随后把这些迎风空间导数代回式（19.61），源项影响便完整进入时间系数。

这种分工十分重要：并不直接求带源黎曼问题，而是用齐次波传播决定界面导数的上游来源，用 CK 关系加入反应。对常系数问题，展开与精确因子 $e^{\beta t}$ 的 Taylor 级数一致。数值通量对 $q_{LR}(\tau)$ 作时间积分，数值源项则按前述单元时空积分计算。

% [原PDF p663]
非线性情形中，首项相似解可能包含激波或稀疏波，$d^{(0)}(0)$ 要由界面所在波区确定。高阶导数随后只描述这一首项附近的短时间变化。$f'(q)$ 是局部特征速度，$f''(q)$ 衡量波速随状态的变化；式（19.69）中的 $f''q_x^2$ 正是特征汇聚或发散对梯度的影响。$s'(q)q_x$ 则描述反应如何放大或衰减空间梯度。

若首项恰位于稀疏扇内部，冻结速度等于该界面状态的特征速度；若位于激波一侧，则由黎曼解选择相应状态。基础求解器必须在 sonic 点和零速度波附近给出稳定一致的状态，否则高阶导数的左右选择会敏感跳变。
\subsection{含源项的非线性方程}
考虑
\begin{equation}
q_t+\partial_xf(q)=s(q),\qquad
q(x,0)=\begin{cases}q_L(x),&x<0,\\q_R(x),&x>0.
\end{cases} \tag{19.63}
\end{equation}
首项由齐次非线性黎曼问题得到，记相似解为 $d^{(0)}(x/t)$，则 $q(0,0^+)=d^{(0)}(0)$。高阶项首先用 CK 过程计算，例如
\begin{equation}
q_t=-f'(q)q_x+s(q), \tag{19.67}
\end{equation}
以及
\begin{equation}
q_{tx}=-f'(q)q_{xx}-f''(q)q_x^2+s'(q)q_x. \tag{19.69}
\end{equation}
非线性使系数依赖首项状态，源项及其导数则与空间导数发生耦合。

% [原PDF p664]
二阶式包含四类物理机制：$[f']^2q_{xx}$ 是两次平流，$2f'f''q_x^2$ 来自非线性波速，含 $s f''$ 与 $f's'$ 的项是通量--源项耦合，$ss'$ 是两次局部反应。遗漏任何一类都会使含源问题的时间阶下降。高阶时组合数迅速增加，适合用符号生成后再进行公共子表达式消除。

导数黎曼问题使用冻结速度 $\lambda_{LR}^{(0)}$，所以标量情形只需按其符号取左或右的 $q^{(k)}$。速度非常接近零时，数值舍入可能改变选择；程序应采用与基础首项求解器一致的容差和零速规则，保证阶数之间不产生不连续切换。
第二阶时间导数为
\begin{equation}
q_{tt}=[f'(q)]^2q_{xx}+2f'(q)f''(q)q_x^2
-[s(q)f''(q)+2f'(q)s'(q)]q_x+s(q)s'(q). \tag{19.70}
\end{equation}
更高阶表达式迅速增长，宜用符号代数生成。空间导数的演化方程具有
\[
\partial_t(\partial_x^kq)+\lambda(q)\partial_x(\partial_x^kq)=h^{(k)},
\qquad\lambda(q)=f'(q).
\]
忽略 $h^{(k)}$ 并在首项处冻结速度 $\lambda_{LR}^{(0)}=f'(q(0,0^+))$，各阶导数状态按该速度符号从左或右选取。再代入 $p^{(k)}$，得到界面时间多项式、数值通量和数值源项。

\subsection{含非线性源项的 Burgers 方程}
以下具体例子检验高阶短时间展开的有效范围。

% [原PDF p665]
左初值在界面附近约为较高的高斯峰，右初值幅值仅为其一小部分，因此原点存在强下降跳跃。Burgers 通量的特征速度就是 $q$，正状态把信息向右输运；正源项 $e^{-q}$ 又不断增加状态，但对较小 $q$ 增长更强，使激波速度随时间改变。这个构型因此产生向右加速的弯曲激波轨迹，而不是齐次 Burgers 方程中的常速 Rankine--Hugoniot 激波。

细网格 WENO 参考解用于展示全局结构，局部 $\mathrm{GRP}_K$ 只预测 $x=0$ 的时间迹线。二者比较前必须把单元平均高阶插值到精确界面位置；直接拿相邻单元中心值会把空间偏移误差误认为广义黎曼误差。
问题为
\begin{equation}
q_t+\partial_x\left(\frac12q^2\right)=e^{-q},\qquad
q(x,0)=\begin{cases}
2e^{-2(x-1/5)^2},&x<0,\\
\frac18e^{-2(x+1/5)^2},&x>0.
\end{cases} \tag{19.76}
\end{equation}
目标是求界面 $x=0$ 上的 $q_{LR}(\tau)$。参考全局解用五阶 WENO 在 $[-5,5]$ 的 1000 个单元上计算。初始跳跃发出向右传播并不断加速的激波，源项 $e^{-q}$ 改变其两侧状态和传播速度。

\begin{sourceblock}
\sourcecrop{665}{110}{315}{330}{490}{0.28}
\captionof{figure}{含非线性源项 Burgers 广义黎曼问题的 $x$--$t$ 等值线。主要结构是由原点产生并向右加速的激波。}
\end{sourceblock}
这个算例同时包含非线性通量、非线性反应和非多项式初值，能检验 CK 系数、导数黎曼问题及源项耦合是否一致。

% [原PDF p666]
\begin{sourceblock}
\sourcecrop{666}{115}{425}{325}{595}{0.27}
\captionof{figure}{界面解比较：六阶 $\mathrm{GRP}_K$ 短时间展开与由细网格 WENO 解高阶插值得到的参考值。}
\end{sourceblock}
参考有限体积解并不恰好位于 $x=0$，因此用 WENO 高阶插值获得界面值。对较短时间，两条曲线符合极好；时间增大后，有限阶幂级数的误差明显增加。非常接近 $t=0$ 时，参考数值解反而会因激波涂抹而不够准确，所以比较必须同时考虑参考解的空间离散误差。

表 19.1 进一步列出 $K=1$ 至 6 在四个时刻的误差。小时间下，增加展开阶数带来快速误差下降；到 $t=0.1$ 时，虽然高阶仍有改善，但已不再呈现同样规则的阶数收益。这正是局部时间 Taylor 预测的本质：它适合覆盖一个 CFL 时间步，不应作为长时间整体解析解使用。

% [原PDF p667]
表中的误差不是严格的网格收敛率，而是在固定若干时间上比较不同截断阶数。对足够小的 $t$，Taylor 余项近似按 $t^{K+1}$ 下降，所以升阶效果显著；时间接近解析半径后，高阶系数增长和参考插值误差都会破坏简单比例。三阶在 $t=0.1$ 反而比二阶误差大，也提醒我们有限级数并非对任意阶数和任意时间单调改进。

在实际 ADER 中，局部时间只覆盖一个 CFL 步，随后以新单元平均重新重构。因此该长时间退化不会逐步沿用同一多项式，而会通过每步重启得到控制。误差表的价值是确认短时间展开正确，并界定单次预测可安全使用的时间尺度。
\begin{table}[htbp]
\centering
\caption{$q_{LR}(\tau)$ 对含源 Burgers 问题的误差。}
\begin{tabular}{ccccc}
\toprule
阶数 & $t=0.010$ & $t=0.025$ & $t=0.050$ & $t=0.100$\\
\midrule
1 & $2.625656\!\times10^{-2}$ & $6.724719\!\times10^{-2}$ & $1.374329\!\times10^{-1}$ & $2.735556\!\times10^{-1}$\\
2 & $5.662654\!\times10^{-4}$ & $3.021457\!\times10^{-3}$ & $8.981471\!\times10^{-3}$ & $1.665263\!\times10^{-2}$\\
3 & $5.807203\!\times10^{-5}$ & $8.80652\!\times10^{-4}$ & $6.626964\!\times10^{-3}$ & $4.578111\!\times10^{-2}$\\
4 & $9.508504\!\times10^{-7}$ & $4.158045\!\times10^{-5}$ & $7.508961\!\times10^{-4}$ & $1.324177\!\times10^{-2}$\\
5 & $8.990528\!\times10^{-8}$ & $7.949777\!\times10^{-6}$ & $2.128054\!\times10^{-4}$ & $4.63232\!\times10^{-3}$\\
6 & $5.743972\!\times10^{-9}$ & $1.39097\!\times10^{-6}$ & $8.609855\!\times10^{-5}$ & $4.932606\!\times10^{-3}$\\
\bottomrule
\end{tabular}
\end{table}
误差表说明阶数与有效时间必须共同考察。六阶在 $t=0.01$ 的误差约为一阶的五百万分之一，但在 $t=0.1$ 只改善约两个数量级。实际 ADER 步长由 CFL 条件限制，局部预测会在下一步用新重构重新开始，因此不会长期外推同一个级数。

\section{其他求解器}
Toro--Titarev 求解器以界面时间 Taylor 展开为核心，是 Ben--Artzi--Falcovitz 二阶方法的任意阶推广。

% [原PDF p668]
不同 GRP 求解器最终都要在若干时间求积节点提供界面状态。Taylor 型方法先一次性计算系数，可廉价评价任意节点；演化--黎曼型方法则先获得左右时空预测，再在每个节点调用经典黎曼求解器。前者适合低维、可符号微分的系统，后者模块化更强，并自然处理非保守复杂源项。比较算法时应把重构和求积保持一致，才能把误差差别归因于局部预测器本身。
Harten--Osher--Engquist--Chakravarthy 方法也可重新解释为 $\mathrm{GRP}_K$ 求解器。它先分别在界面左右演化重构数据的边界外推值，再在所需求积时刻以这些演化状态求经典黎曼问题。与 Toro--Titarev 方法直接展开界面解不同，它把局部光滑演化和界面波相互作用的次序重新组织；两者在精度阶上等价，但非线性问题中的具体结果不完全相同。

Dumbser--Enaux--Toro 求解器是后一思路的数值化推广。它使用局部时空间断 Galerkin 预测器演化左右数据，避免显式 Cauchy--Kowalewski 替换；随后在时间求积点求经典黎曼问题。若局部预测隐式完成，就能正确处理刚性源项，同时保持全局格式显式。预测得到的单元时空多项式也可直接用于源项体积积分，因此通量与源项自然具有同阶精度。

这种局部隐式、全局显式结构把刚性限制封装在单元预测阶段，不需求解全局耦合系统。对于复杂方程组，它还避免手工推导高阶时间导数，是现代 ADER 实现更常用的路线。

% [原PDF p669]
从软件实现看，广义黎曼模块应明确区分五类输入输出。输入包括界面两侧重构多项式的函数值与各阶物理空间导数、基础方程的通量和源项模型；首项求解器返回界面 Godunov 状态；线性化模块在该状态构造 Jacobian 与特征结构；导数求解器返回每阶空间导数的界面状态；最后 CK 或局部预测模块组装时间多项式。这样的分层使基础黎曼求解器、重构阶数和源项模型可以独立替换，也便于用常状态、线性平流和已知标量解逐层验证。

误差控制还应覆盖物理与数值两个层面。数学上的 $K+1$ 阶只描述光滑区域和渐近网格细化；遇到激波时，全局范数收敛阶会降低，非线性限制器可能主动退阶。界面预测若产生不可容许状态，不能靠无声截断掩盖，而应检查时间步、重构和基础波模型。含源项系统还可能存在稳态平衡、刚性反应或几何奇点，需要相应的良平衡重构、局部隐式预测或坐标专门处理；一般 GRP 公式提供统一框架，但并不自动解决所有这些附加约束。

本章方法最重要的概念贡献，是把高阶时间精度转化为局部初值问题，而不是依赖全局多级时间迭代。空间重构给出界面两侧的高阶数据，物理方程决定这些数据在第一瞬间如何相互作用，时间求积再把局部解变成守恒通量。这个链条保持了 Godunov 方法“由局部波传播定义界面通量”的核心，同时把分片常数初值推广为分片光滑初值。下一章的 ADER 正是把该局部链条在每个网格面和每个时间步系统化执行。

\section{结语}
本章详细研究了含源项非线性双曲系统的 $\mathrm{GRP}_K$ 半解析求解法。初值由原点跳跃连接的两个任意光滑向量场构成。求解过程把完整非齐次非线性问题化为一个齐次非线性经典黎曼问题，以及一列空间导数的齐次线性黎曼问题。标量例子说明了首项、CK 替换、导数迎风状态和源项耦合的具体做法。

广义黎曼解的主要用途是构造高阶数值通量与数值源项，从而在有限体积或间断 Galerkin 框架下形成时空任意阶的全离散一步格式。不同求解器的差异主要在局部预测：显式 CK 展开便于分析，局部时空 DG 预测更适合复杂系统和刚性源项。无论采用哪条路线，都必须在一个短时间步内使用界面时间函数，并保证通量、源项与重构的阶数匹配。

下一章把这些部件组织成 ADER 有限体积方法，说明从单元平均出发，如何依次完成重构、局部预测、界面通量求积、源项体积求积和一步守恒更新。

% [原PDF p670]
\chapter{ADER 方法}
ADER 是一族时空任意阶的全离散一步格式，用于求解含源项的双曲方程。它推广了 Godunov 方法：数值通量是物理通量在第 19 章广义黎曼界面解上的时间平均，数值源项则是相应时空控制体内源项的高阶平均。ADER 可用于有限体积和间断 Galerkin 框架；这里先讨论一维模型的有限体积形式。

ADER 名称源于“任意高阶导数黎曼问题”思想，但现代实现未必显式形成所有导数。其不变的核心是：在一个时间步内构造高阶局部时空解，并用它直接完成通量与源项积分。与多步或多级方法相比，一步结构更容易保持几何守恒和局部时间信息，也更适合异构并行；代价是每单元预测器更复杂。选择阶数时应综合目标误差、解的光滑比例、通信成本和局部算力。

\section{引言}
二阶 TVD 方法已经成熟并广泛用于工业软件，但长时间波传播等问题常要求更高精度。若目标误差低于某一阈值，已有证据表明，在粗网格上采用高阶方法通常比在极细网格上采用低阶方法更节省内存和计算时间。真正的问题不是“阶数越高越好”，而是在给定误差容限下选择总成本最低的时空离散。

适合波动问题的高阶格式还必须满足严格条件：由 Lax--Wendroff 定理要求守恒，由 Godunov 定理要求高于一阶的无振荡格式必须是非线性的；同时还要处理源项和复杂多维几何，在光滑解上达到理论收敛阶，在间断附近避免或显著减弱伪振荡。

% [原PDF p671]
全离散一步与半离散多级方法的成本结构不同。半离散 WENO 每个 Runge--Kutta 阶段通常要重新计算重构或至少重新计算空间残差；ADER 每步只重构一次，在局部预测中生成整个时间区间的高阶信息，然后一次完成通量和源项求积。高阶时这种局部计算更密集，却减少了全局同步和邻域通信，对并行计算尤其有吸引力。

高阶并不自动等于高效。重构模板、求积点数、局部预测迭代和存储都会随阶数增加；只有当目标误差足够小、解具有足够光滑区域时，粗网格高阶方案的优势才会显现。间断附近非线性重构会局部降低有效阶数，这是保证无振荡的必要代价。
Harten 等提出了包含若干关键思想的全离散一步高阶方法；Shu 等发展出的半离散 WENO 加多级时间推进如今更常见。ADER 则在每个时间步只做一次高阶空间重构，并在每个界面求广义黎曼问题，以局部时空预测直接完成高阶时间积分。

早期 ADER 针对笛卡尔网格上线性问题，在一、二、三维均有构造，固定模板实现甚至达到时空十阶。常系数线性问题上，它与 Harten--Osher--Engquist--Chakravarthy 方法相同；非线性问题上两者虽同阶，却因局部演化组织不同而不完全相同。ADER 也可看作 Ben--Artzi--Falcovitz 二阶 GRP 方法的任意阶推广。

后续研究把 ADER 扩展到非线性标量方程、二维和三维系统、非笛卡尔及非结构网格，并发展出有限体积和 DG 两种框架。局部隐式时空 DG 预测器进一步使高阶精度与刚性源项处理相容。

% [原PDF p672]
现代 ADER 框架可把不同广义黎曼求解器视为可替换的局部预测模块。显式 CK 版本直接产生时间导数，局部时空 DG 版本则求一个单元内弱形式预测多项式；两者都必须向界面求积返回时间相关左右状态，并向源项求积返回单元内部时空状态。清晰的接口分离使同一守恒更新可支持不同方程组、刚性程度和网格类型。

本节的一维标量模型保留了全部逻辑：精确积分恒等式、守恒重构、界面 GRP、时间平均通量、体积平均源项。推广到系统和多维时，标量通量变成法向通量向量，空间求积变成面与体积求积，但步骤顺序不变。
本章先给出有限体积 ADER 框架，再构造含源标量方程的二阶与任意阶方案，随后展示数值结果。基本流程在不同阶数间保持不变，变化的是重构多项式次数、局部预测阶数和时空求积规则。

\section{有限体积方法}
\subsection{框架}
考虑
\begin{equation}
q_t+\partial_xf(q)=s(q). \tag{20.1}
\end{equation}
在 $[x_{i-1/2},x_{i+1/2}]\times[t^n,t^{n+1}]$ 上积分得到
\begin{equation}
q_i^{n+1}=q_i^n-\frac{\Delta t}{\Delta x}
(f_{i+1/2}-f_{i-1/2})+\Delta t\,s_i. \tag{20.2}
\end{equation}
其中 $q_i^n$ 是空间平均，$f_{i+1/2}$ 是界面物理通量的时间平均，$s_i$ 是源项的空间--时间平均。这一积分关系本身精确，误差全部来自三个平均量的数值近似。

% [原PDF p673]
式（20.2）中的三个平均各有不同几何含义。$q_i^n$ 是固定时间切片上的空间平均；$f_{i+1/2}$ 是固定空间界面上的时间平均；$s_i$ 是整个空间--时间控制体的平均。把它们混用点值会破坏积分推导。例如用单元中心源项代替时空平均只能得到低阶，除非源项与解在单元内恰为常量。

网格可非均匀，此时每个单元使用自己的 $\Delta x_i$，界面仍必须共享同一时间平均通量。初始化连续数据时应采用与目标阶数相称的空间求积生成真实单元平均；否则第一步之前就已经引入点值到平均值的低阶投影误差。
具体定义为
\begin{align}
q_i^n&=\frac1{\Delta x}\int_{x_{i-1/2}}^{x_{i+1/2}}q(x,t^n)\,dx, \tag{20.3}\\
f_{i+1/2}&=\frac1{\Delta t}\int_{t^n}^{t^{n+1}}f(q(x_{i+1/2},t))\,dt, \tag{20.4}\\
s_i&=\frac1{\Delta t\Delta x}\int_{t^n}^{t^{n+1}}
\int_{x_{i-1/2}}^{x_{i+1/2}}s(q(x,t))\,dx\,dt. \tag{20.5}
\end{align}
在数值算法中，$q_i^n$ 是已知单元平均。首先从邻域平均重构单元内空间函数；然后预测它在整个时间步中的局部演化。界面预测为数值通量提供积分函数，单元内部预测为数值源项提供积分函数。完成两个求积后，式（20.2）一次性推进到新时间层。

初始化时应把连续初值真正积分成单元平均，而不是简单取单元中心点值；后续更新自动保持平均量语义。若初始投影低阶，即使演化格式很高阶，也可能在收敛试验中污染初始误差。

\subsection{数值通量}

% [原PDF p674]
界面两侧多项式 $p_i,p_{i+1}$ 各自守恒重构，却一般在 $x=0$ 有函数值和各阶导数跳跃，因此恰好构成 $\mathrm{GRP}_K$ 数据。求得 $q_{i+1/2}(\tau)$ 后，非线性通量必须先逐求积点计算 $f(q)$ 再积分，不能简单把平均状态代入 $f$。高斯求积的节点数要按 $f(q(\tau))$ 的有效次数或目标误差选择。

当 $K=0$ 时，时间函数只有 Godunov 首项，式（20.8）退化为经典一阶黎曼通量；当 $K=1$ 时出现线性时间变化并得到二阶 GRP 方法。ADER 因而不是另一套孤立格式，而是沿同一界面积分定义连续推广 Godunov 阶数。
在界面局部坐标中令 $x_{i+1/2}=0,t^n=0$。ADER 通量为
\begin{equation}
f_{i+1/2}=\frac1{\Delta t}\int_0^{\Delta t}f(q_{i+1/2}(\tau))\,d\tau, \tag{20.8}
\end{equation}
其中 $q_{i+1/2}(\tau)$ 是广义黎曼问题
\begin{equation}
q_t+f(q)_x=s(q),\qquad q(x,0)=
\begin{cases}p_i(x),&x<0,\\p_{i+1}(x),&x>0
\end{cases} \tag{20.9}
\end{equation}
在固定界面上的解。$p_i,p_{i+1}$ 是守恒、非振荡重构得到的 $K$ 次多项式。界面时间函数写成
\begin{equation}
q_{i+1/2}(\tau)=q(0,0^+)+\sum_{k=1}^{K}
\partial_t^kq(0,0^+)\frac{\tau^k}{k!}. \tag{20.10}
\end{equation}
无源时，由 $K$ 次重构和相应界面预测得到的有限体积格式在空间、时间均为 $K+1$ 阶。

\subsection{数值源项}
源项必须与界面通量同阶，否则整体格式会被最低阶环节限制。

% [原PDF p675]
源项预测可在空间高斯点分别构造时间多项式，也可先构造一个完整局部时空多项式。前者概念直接，后者便于在所有面和体积求积中复用。刚性源项下显式 Taylor 系数可能很大，局部隐式预测能在不形成全局耦合矩阵的情况下稳定处理反应，再由全局显式守恒更新连接各单元。

重构的守恒条件比插值通过某个点更重要：有限体积未知量是平均值，$p_i$ 的积分必须严格恢复 $q_i^n$。固定模板线性重构在光滑区易于达到高阶，但跨越间断时会产生 Gibbs 振荡。ENO 通过选择较平滑模板，WENO 通过非线性组合多个候选模板，在保持守恒的同时抑制振荡。
ADER 源项为
\begin{equation}
s_i=\frac1{\Delta t\Delta x}\int_0^{\Delta t}
\int_{x_{i-1/2}}^{x_{i+1/2}}s(q_i(x,t))\,dx\,dt. \tag{20.11}
\end{equation}
在若干空间求积点 $\hat x$ 上，可由 $p_i(x)$ 提供空间导数，再通过 CK 过程构造时间 Taylor 级数；也可直接用局部时空预测器。对 $s(q_i)$ 采用足够阶数的张量积求积，即得到与通量匹配的数值源项。

\subsection{重构}
每个时间层只有单元平均 $\{q_i^n\}$，它们表示分片常数数据，单元内变化在平均过程中丢失。重构要由邻域平均恢复局部光滑函数 $p_i(x)$，并严格满足守恒约束
\[
\frac1{\Delta x}\int_{I_i}p_i(x)\,dx=q_i^n.
\]
一般相邻多项式在单元面处不连续，恰好为广义黎曼问题提供左右初值。

固定模板重构称为线性重构。由 Godunov 定理，高于一阶的线性格式在间断附近必然振荡；因此要求无振荡时，必须让模板或权重依赖局部数据，形成非线性重构。

% [原PDF p676]
图 20.1 中水平短线表示单元平均，并不是多项式在中心的必经点。守恒重构只约束面积，中心值可高于或低于平均。相邻多项式在公共面产生两个外推值和两组导数，正是 GRP 的左右数据。在线性固定模板上，这些量是邻域平均的固定线性组合；在 ENO/WENO 中，模板或权重随光滑性变化，使整个格式成为非线性，从而绕开 Godunov 线性高阶单调性障碍。
\begin{sourceblock}
\sourcecrop{676}{65}{480}{330}{610}{0.25}
\captionof{figure}{三个相邻有限体积中的单元平均与重构多项式。多项式保持各自平均，但在界面一般具有跳跃。}
\end{sourceblock}
最简单的非线性重构基于第 13 章的 TVD 准则；更高阶 ADER 常用 ENO 或 WENO。ENO 不在全部样本上构造一个高次全局多项式，而是在每个单元局部选择较平滑的模板。其三个特征是：局部性，只在 $I_i$ 内定义；守恒性，积分平均严格等于 $q_i^n$；自适应性，系数随数据光滑性改变。由此在含间断或大梯度的数据上保持“本质无振荡”。

最低阶重构只是常多项式
\begin{equation}
p_i(x)=q_i^n,\qquad x\in I_i, \tag{20.14}
\end{equation}
它把 ADER 框架退化为一阶 Godunov 型有限体积法。提高阶数时，重构、广义黎曼预测、界面时间求积和源项时空求积必须同步升级；其中任何一个环节低一阶，都会限制最终收敛阶。
% ===== 译文分片 8（已合并）=====

% [原PDF p677]
在每个区间 $I_i$ 上取常数多项式，就得到全局分片常数重构；它显然满足前述守恒性、精度和局部性要求。若希望得到一次重构，在单元
$I_i=[x_{i-1/2},x_{i+1/2}]$ 内写成
\begin{equation}
 p_i^{(1)}(x)=q_i^n+(x-x_i)\Delta_i^{(1)}, \tag{20.15}
\end{equation}
其中斜率由相邻两侧梯度决定：
\begin{align}
 \Delta_i^{(1)}&=\Delta_i^{(1)}(\Delta_{i-1/2},\Delta_{i+1/2}), \tag{20.16}\\
 \Delta_{i-1/2}&=\frac{q_i^n-q_{i-1}^n}{\Delta x},\qquad
 \Delta_{i+1/2}=\frac{q_{i+1}^n-q_i^n}{\Delta x}. \tag{20.17}
\end{align}
最简单的中心选择是
\begin{equation}
 \Delta_i^{(1)}=\frac12(\Delta_{i-1/2}+\Delta_{i+1/2}). \tag{20.18}
\end{equation}
ENO 方法则自适应选择较平滑的一侧：
\begin{equation}
 \Delta_i^{(1)}=\begin{cases}
 \Delta_{i-1/2},&|\Delta_{i-1/2}|\le |\Delta_{i+1/2}|,\\
 \Delta_{i+1/2},&|\Delta_{i-1/2}|>|\Delta_{i+1/2}|.
 \end{cases} \tag{20.19}
\end{equation}
由此定义全局分片线性函数
\begin{equation}
 h(x)=\begin{cases}p_i(x),&x\in I_i,\\0,&\text{其他位置。}\end{cases} \tag{20.20}
\end{equation}
这一过程等价于在两个候选模板
\begin{equation}
 S_L=\{q_{i-1}^n,q_i^n\},\qquad S_R=\{q_i^n,q_{i+1}^n\} \tag{20.21}
\end{equation}
上分别构造 $q_L(x)$、$q_R(x)$，再选取一阶差商绝对值较小者。更高阶 ENO 多项式也按同一原则逐级扩展模板。WENO（加权本质无振荡）并不舍弃候选模板，而是以光滑度相关的非线性权重组合所有 ENO 候选多项式。其系统论述见文献 [331,270,493]；适用于二维、三维非结构网格的改进 WENO 重构见 [175,176]。

% [原PDF p678]
\section{模型方程的二阶格式}
ADER 是一个统一的任意阶框架。为把全部步骤展示清楚，本节仍按一般高阶构造的顺序，针对带源项的线性平流模型推导二阶格式：
\begin{equation}
 \partial_t q+\lambda\partial_x q=\beta q, \tag{20.22}
\end{equation}
其中通量 $f(q)=\lambda q$，源项 $s(q)=\beta q$，常数 $\lambda$ 已知且 $\beta\le0$。有限体积更新式 (20.2) 所需的两个量，是界面数值通量 $f_{i+1/2}$ 与单元数值源项 $s_i$。

\subsection{数值通量与数值源项}
先以一次多项式重构相邻单元数据：
\begin{equation}
 p_i(x)=q_i^n+\Delta_i^{(1)}(x-x_i),\qquad
 p_{i+1}(x)=q_{i+1}^n+\Delta_{i+1}^{(1)}(x-x_{i+1}), \tag{20.23}
\end{equation}
斜率例如可按式 (20.19) 选取。为计算式 (20.8) 中的 $f_{i+1/2}$，把界面移到 $x=0$，求解一阶广义黎曼问题
\begin{equation}
 \left.\begin{aligned}
 &\partial_tq+\lambda\partial_xq=\beta q,\\
 &q(x,0)=\begin{cases}p_i(x),&x<0,\\p_{i+1}(x),&x>0.
 \end{cases}
 \end{aligned}\right\} \tag{20.24}
\end{equation}
界面解在一个时间步内写成
\begin{equation}
 q_{i+1/2}(\tau)=q(0,0^+)+\tau\partial_tq(0,0^+). \tag{20.25}
\end{equation}
第一步确定首项。以左右边界外推值为初态，求齐次经典黎曼问题：
\begin{equation}
 q(x,0)=\begin{cases}
 p_i(0)=q_i^n+\frac12\Delta x\Delta_i^{(1)},&x<0,\\
 p_{i+1}(0)=q_{i+1}^n-\frac12\Delta x\Delta_{i+1}^{(1)},&x>0.
 \end{cases} \tag{20.26}
\end{equation}
这一步只处理由传播方向决定的界面瞬时状态；源项将在高阶时间导数及体积分中同时进入。

% [原PDF p679]
式 (20.26) 的自相似解为
\begin{equation}
 d_{i+1/2}^{(0)}(x/t)=\begin{cases}
 q_i^n+\frac12\Delta x\Delta_i^{(1)},&x/t<\lambda,\\
 q_{i+1}^n-\frac12\Delta x\Delta_{i+1}^{(1)},&x/t>\lambda,
 \end{cases} \tag{20.27}
\end{equation}
故首项就是 $x/t=0$ 处的 Godunov 状态：
\begin{equation}
 q(0,0^+)=d_{i+1/2}^{(0)}(0)=\begin{cases}
 q_i^n+\frac12\Delta x\Delta_i^{(1)},&\lambda>0,\\
 q_{i+1}^n-\frac12\Delta x\Delta_{i+1}^{(1)},&\lambda<0.
 \end{cases} \tag{20.28}
\end{equation}
第二步计算式 (20.25) 的高阶项。Cauchy--Kowalewski 过程直接利用偏微分方程，把时间导数改写为空间导数和源项：
\begin{equation}
 \partial_tq(x,t)=-\lambda\partial_xq(x,t)+\beta q(x,t), \tag{20.29}
\end{equation}
于是
\begin{equation}
 q_{i+1/2}(\tau)=q(0,0^+)+[-\lambda\partial_xq(0,0^+)+\beta q(0,0^+)]\tau. \tag{20.30}
\end{equation}
还需确定界面空间导数。对原方程关于 $x$ 求导可见，$\partial_xq$ 满足
\begin{equation}
 \partial_t(\partial_xq)+\lambda\partial_x(\partial_xq)=\beta\partial_xq. \tag{20.31}
\end{equation}
为求导数的界面迎风值，先忽略梯度方程的源项，建立齐次黎曼问题
\begin{equation}
 \left.\begin{aligned}
 &\partial_t(\partial_xq)+\lambda\partial_x(\partial_xq)=0,\\
 &\partial_xq(x,0)=\begin{cases}\Delta_i^{(1)},&x<0,\\
 \Delta_{i+1}^{(1)},&x>0.
 \end{cases}
 \end{aligned}\right\} \tag{20.32}
\end{equation}
其解仍只由 $\lambda$ 的符号选择左右斜率。该步骤说明 ADER 并非简单在通量上加一个时间校正，而是在界面同时演化状态及其导数。

% [原PDF p680]
导数黎曼问题的自相似解及其界面值分别是
\begin{align}
 d_{i+1/2}^{(1)}(x/t)&=\begin{cases}\Delta_i^{(1)},&x/t<\lambda,\\
 \Delta_{i+1}^{(1)},&x/t>\lambda,
 \end{cases} \tag{20.33}\\
 \partial_xq(0,0^+)&=d_{i+1/2}^{(1)}(0)=\begin{cases}\Delta_i^{(1)},&\lambda>0,\\
 \Delta_{i+1}^{(1)},&\lambda<0.
 \end{cases} \tag{20.34}
\end{align}
第三步把首项和导数项合并，得到广义黎曼问题的完整二阶界面解：
\begin{equation}
q_{i+1/2}(\tau)=\begin{cases}
q_i^n+\frac12\Delta x\Delta_i^{(1)}+
\tau\!\left[-\lambda\Delta_i^{(1)}+\beta\left(q_i^n+\frac12\Delta x\Delta_i^{(1)}\right)\right],&\lambda>0,\\[2mm]
q_{i+1}^n-\frac12\Delta x\Delta_{i+1}^{(1)}+
\tau\!\left[-\lambda\Delta_{i+1}^{(1)}+\beta\left(q_{i+1}^n-\frac12\Delta x\Delta_{i+1}^{(1)}\right)\right],&\lambda<0.
\end{cases}\tag{20.35}
\end{equation}
第四步在 $[0,\Delta t]$ 上积分界面通量。令 Courant 数 $c=\lambda\Delta t/\Delta x$，无量纲反应数 $r=\beta\Delta t$，可写成
\begin{equation}
f_{i+1/2}=\begin{cases}
\lambda\!\left[q_i^n+\frac12(1-c)\Delta x\Delta_i^{(1)}+
\frac12r\left(q_i^n+\frac12\Delta x\Delta_i^{(1)}\right)\right],&\lambda>0,\\
\lambda\!\left[q_{i+1}^n-\frac12(1+c)\Delta x\Delta_{i+1}^{(1)}+
\frac12r\left(q_{i+1}^n-\frac12\Delta x\Delta_{i+1}^{(1)}\right)\right],&\lambda<0.
\end{cases}\tag{20.36}
\end{equation}
值得强调的是，数值通量已经“知道”源项的存在；但这并不替代单元内部的数值源项贡献。第五步以时空中点求积计算体积分。对单元中心的重构解作 Taylor 展开并应用 Cauchy--Kowalewski 关系：
\begin{equation}
q_i(x_i,\tau)=q_i(x_i,0^+)+\tau[-\lambda\partial_xq_i(x_i,0^+)+\beta q_i(x_i,0^+)].\tag{20.37}
\end{equation}
这一推导还揭示了守恒通量与体源项之间的分工。界面预测器沿特征方向选择上游重构，并在时间积分中加入反应造成的状态变化；单元源项则近似控制体内的平均生成或衰减。两部分必须使用相容的时间展开，否则即使各自看似二阶，组合更新也可能只剩一阶，或者破坏本应保持的平衡状态。因而式 (20.36) 中的源项修正不能随意删除，也不能与式 (20.38) 重复计入。

% [原PDF p681]
在本例中 $q_i(x_i,0^+)=q_i^n=p_i(x_i)$，且 $\partial_xq_i(x_i,0^+)=\Delta_i^{(1)}$，因此式 (20.11) 的二阶数值源项为
\begin{equation}
s_i=\beta\left[q_i^n+\frac12\Delta t(-\lambda\Delta_i^{(1)}+\beta q_i^n)\right]. \tag{20.38}
\end{equation}
\subsection{完整格式}
式 (20.36) 与 (20.38) 已完全确定模型方程的二阶有限体积格式。对 $\lambda>0$，更新式为
\begin{align}
q_i^{n+1}={}&q_i^n-c\left[(q_i^n-q_{i-1}^n)+\frac12(1-c)
(\Delta_i^{(1)}-\Delta_{i-1}^{(1)})\right]\nonumber\\
&-cr\left[\frac12(q_i^n-q_{i-1}^n)+\frac12\Delta x
(\Delta_i^{(1)}-\Delta_{i-1}^{(1)})\right]
+r\left[(1+\tfrac12r)q_i^n-\tfrac12c\Delta x\Delta_i^{(1)}\right]. \tag{20.39}
\end{align}
若所有斜率置零，同时删去源项对数值通量的贡献，就退化为平流部分的一阶 Godunov 迎风格式与源项迎风离散的组合：
\begin{equation}
q_i^{n+1}=q_i^n-c(q_i^n-q_{i-1}^n)-r\left[\frac12cq_{i-1}^n+(1-\tfrac12c)q_i^n\right]. \tag{20.40}
\end{equation}
这与文献 [412,51] 分析的一类带源平流迎风格式相对应。当 $\beta=0$ 且 $\lambda>0$ 时，不同斜率选择会精确恢复三个经典二阶格式：
\begin{equation}
\Delta_i^{(1)}=\begin{cases}
\Delta_{i-1/2},&\text{Warming--Beam，}\\
\Delta_{i+1/2},&\text{Lax--Wendroff，}\\
\frac12(\Delta_{i-1/2}+\Delta_{i+1/2}),&\text{Fromm。}
\end{cases}\tag{20.41}
\end{equation}
所以，采用 ENO 斜率 (20.19) 的二阶 ADER，本质上是在 Warming--Beam 与 Lax--Wendroff 两个格式之间按局部光滑性切换。

% [原PDF p682]
\section{任意阶精度格式}
任意阶 ADER 的构造程序与前述二阶例子完全相同，只是重构多项式、界面时间展开和体积分求积的阶数随之提高。
\subsection{数值通量}
在界面求解 $K$ 次重构数据的广义黎曼问题
\begin{equation}
\left.\begin{aligned}
&\partial_tq+\lambda\partial_xq=\beta q,\\
&q(x,0)=\begin{cases}p_i(x),&x<0,\\p_{i+1}(x),&x>0,
\end{cases}
\end{aligned}\right\}\tag{20.42}
\end{equation}
其中 $p_i,p_{i+1}$ 是任意 $K$ 次多项式。按照第 19 章，界面解展开为
\begin{equation}
q_{i+1/2}(\tau)=q(0,0^+)+\sum_{k=1}^{K}[\partial_t^{(k)}q(0,0^+)]\frac{\tau^k}{k!}.\tag{20.43}
\end{equation}
首项仍由左右外推界面值组成的经典黎曼问题给出。对每个高阶项，Cauchy--Kowalewski 过程将时间导数化为空间导数：
\begin{equation}
\partial_t^{(k)}q(x,t)=\sum_{n=0}^{k}\frac{k!}{(k-n)!n!}
\beta^n(-\lambda)^{k-n}\partial_x^{(k-n)}q(x,t).\tag{20.44}
\end{equation}
无源时简化为
\begin{equation}
\partial_t^{(k)}q=(-1)^k\lambda^k\partial_x^{(k)}q.\tag{20.45}
\end{equation}
像式 (20.32)--(20.34) 那样逐阶求空间导数的经典黎曼问题，便得到任意阶 $\partial_x^{(l)}q(0,0^+)$，继而确定全部时间导数和式 (20.43)。数值通量是该界面预测解的时间平均：
\begin{equation}
f_{i+1/2}=\frac1{\Delta t}\int_0^{\Delta t}\lambda\left[q(0,0^+)+
\sum_{k=1}^{K}\partial_t^{(k)}q(0,0^+)\frac{\tau^k}{k!}\right]\,d\tau.\tag{20.46}
\end{equation}

% [原PDF p683]
本线性模型可精确完成时间积分，得到
\begin{equation}
f_{i+1/2}=\lambda\left[q(0,0^+)+\sum_{k=1}^{K}
\partial_t^{(k)}q(0,0^+)\frac{(\Delta t)^k}{(k+1)!}\right].\tag{20.47}
\end{equation}
齐次情形 $s(q)=0$ 的显式形式为
\begin{equation}
f_{i+1/2}=\begin{cases}
\displaystyle\lambda\left[p_i(0)+\sum_{k=1}^{K}(-1)^k\lambda^k
\partial_x^{(k)}p_i(0)\frac{(\Delta t)^k}{(k+1)!}\right],&\lambda>0,\\[2mm]
\displaystyle\lambda\left[p_{i+1}(0)+\sum_{k=1}^{K}(-1)^k\lambda^k
\partial_x^{(k)}p_{i+1}(0)\frac{(\Delta t)^k}{(k+1)!}\right],&\lambda<0.
\end{cases}\tag{20.48}
\end{equation}
\subsection{数值源项}
为计算有限体积格式 (20.2) 的 $s_i$，在单元 $i$ 内选取 $L$ 个求积点 $x_i^{(l)}$。每点的时域预测解写成
\begin{equation}
q(x_i^{(l)},\tau)=q(x_i^{(l)},0^+)+\sum_{k=1}^{K}
[\partial_t^{(k)}q(x_i^{(l)},0^+)]\frac{\tau^k}{k!}.\tag{20.49}
\end{equation}
首项由重构多项式直接给出：$q(x_i^{(l)},0^+)=p_i(x_i^{(l)})$。其余时间导数由
\begin{equation}
\partial_t^{(k)}q(x_i^{(l)},0)=\sum_{n=0}^{k}\frac{k!}{(k-n)!n!}
\beta^n(-\lambda)^{k-n}\partial_x^{(k-n)}q(x_i^{(l)},0^+)\tag{20.50}
\end{equation}
计算，而右侧空间导数均可从 $p_i(x)$ 求得。最后以与目标阶数相匹配的时空求积计算式 (20.11) 的完整体积分，即得 $s_i$。
在实际实现中，$L$ 个空间求积点与时间 Taylor 项数必须共同保证 $K+1$ 阶精度。重构只在 $t^n$ 执行一次，但同一重构多项式会被界面预测和单元体积分反复求值。这样既保持单步结构，也避免多级时间推进中每一级重新限制和重构。对线性常系数模型，所有导数都能显式写出；对非线性系统，则需在局部预测器中递归计算 Jacobian、源项导数及其耦合项。

% [原PDF p684]
\subsection{构造小结}
以上步骤给出了带源标量双曲方程的统一求解器。首先用 TVD、ENO 或 WENO 非线性重构产生任意 $K$ 次多项式 $p_i(x)$；其次在每个界面求解式 (20.42) 的广义黎曼问题，并用式 (20.43) 的时域预测解计算数值通量；最后在预定内部求积点上应用 Cauchy--Kowalewski 过程，用重构多项式提供全部空间导数，再完成时空体积分得到数值源项。

最终方法是空间、时间均为 $K+1$ 阶的全离散单步格式，而且每个时间步只需执行一次非线性重构。迄今讨论的显式 ADER 适合非刚性源项。文献 [173] 提出了一种处理刚性源项的新变体，在保持时空高阶精度的同时协调刚性；另见 [171]。

\section{数值算例}
以下算例用一维线性平流和二维非结构网格上的线性化 Euler 方程，说明高阶 ADER 的表现。
\subsection{光滑剖面的长时间平流}
取式 (20.22) 中 $\beta=0$、$\lambda=1$，计算域为 $[0,10\pi]$，周期边界，光滑初值
\begin{equation}
q(x,0)=q_0(x)=\sin(\pi x),\tag{20.51}
\end{equation}
精确解为 $q(x,t)=q_0(x-\lambda t)$。把一个光滑但快速变化的波形推进极长时间，对数值耗散和相位误差都是严苛考验。传统 TVD 方法尤其难以在很长演化时间后保持波峰与波谷。图 20.2 比较 WAF--MINBEE 在 50 与 100 个单元、$C_{\mathrm{cfl}}=0.95$、$t=1000\pi$ 时的结果。

% [原PDF p685]
\begin{figure}[htbp]
\centering
\fbox{\parbox{0.78\textwidth}{\centering 图示数据摘要：$t=1000\pi$ 时，分别采用 50 和 100 个单元的 WAF--MINBEE 解；横轴 $x\in[0,10\pi]$，纵轴为解 $q$，实线表示精确正弦波，符号表示数值解。}}
\caption{线性平流：MINBEE 限制器的 TVD 结果，$C_{\mathrm{cfl}}=0.95$。原图由 V. A. Titarev 博士提供。}
\end{figure}
结果清楚显示二阶 TVD 方法在长时间演化中的局限。MINBEE 或 MINMOD 这类耗散型限制器会强烈削平极值；时间越长，峰值衰减越明显。增加网格数只能延缓而不能消除这种累积耗散。

若改用压缩性更强的限制器，总体误差通常会下降，但代价是把本来光滑的轮廓逐渐“方波化”。图 20.3 给出 WAF--SUPERBEE 在同一输出时间、50 与 100 个单元上的结果，宜与图 20.2 对照。这个例子说明：TVD 性质负责禁止新的总变差增长，却不保证光滑波形的高阶相位精度；限制器越积极地保护极值，越可能牺牲光滑形状。

% [原PDF p686]
一致高阶的 ADER 方法更适合长时间问题，特别是解虽然光滑却包含快速空间变化时。真实物理解完全可能高度振荡而仍然光滑；这些具有物理意义的振荡不应被数值方法人为压制。第五阶 ADER 在 $t=1000\pi$ 的结果见图 20.4，同样采用 50 与 100 个单元。与 MINBEE 和 SUPERBEE 的 TVD 结果相比，ADER 对振幅和相位的保持都明显更好。

\begin{figure}[htbp]
\centering
\fbox{\parbox{0.78\textwidth}{\centering 图示数据摘要：WAF--SUPERBEE 在 50、100 个单元上的长时间平流解。实线为精确解，离散符号为数值解；压缩限制器减小峰值损失，却使圆滑的正弦轮廓趋于方形。}}
\caption{线性平流：SUPERBEE 限制器的 TVD 结果，$t=1000\pi$，$C_{\mathrm{cfl}}=0.95$。}
\end{figure}
这里真正重要的并不是某一个网格上的瞬时误差，而是每个周期中很小的耗散和频散误差经过上千个周期后的积累。高阶时空一致性使 ADER 能显著推迟这种退化。

% [原PDF p687]
\begin{figure}[htbp]
\centering
\fbox{\parbox{0.78\textwidth}{\centering 图示数据摘要：第五阶 ADER 在 50、100 个单元上的结果。$t=1000\pi$ 后，离散点仍紧贴精确正弦曲线，峰谷和相位均明显优于二阶 TVD 方法。}}
\caption{光滑剖面平流：第五阶 ADER，$C_{\mathrm{cfl}}=0.95$。原图由 V. A. Titarev 博士提供。}
\end{figure}

\subsection{收敛阶}
验证一个方法是否真正达到宣称的阶数，是数值研究中不可省略的工作。关于高阶精度的表述常常只来自形式推导，却没有可靠的网格收敛证据。即使对非齐次方程，构造一个确实二阶且非线性的格式也绝非琐事；换到非结构网格后更具挑战。

下一组结果取自文献 [178]，将 ADER 实现在间断 Galerkin 有限元框架中，求解二维非结构三角形网格上的线性化 Euler 方程。它同时检验空间多项式、时间预测、数值通量以及不规则几何映射的耦合误差，因此比规则一维算例更能揭示实现中的阶数损失。

% [原PDF p688]
表 20.1 汇总空间和时间二阶至十阶格式的经验收敛率。研究刻意采用图 20.6 所示的畸变网格。表中第一列是网格尺度 $h$，第二列是单元数 $N_d$，随后三列为 $L_\infty$、$L_1$、$L_2$ 误差，最后三列给出相应观测阶。实测收敛率与理论阶数一致。

图 20.5 上图把误差对 $1/h$ 作图，下图把误差对 CPU 时间作图，覆盖二阶到十阶 ADER--DG。进入小误差区以后，高阶方法通常比低阶方法更经济；虽然单元和单步成本较高，但达到同一误差所需的网格和总计算量显著减少。

\section{本章小结}
本章介绍了求解带源双曲方程的 ADER 方法族。它是经典一阶 Godunov 迎风方法的自然延伸，核心组成是：用高阶空间多项式代替分片常数数据；在单元界面求解广义而非经典黎曼问题；对界面时域解积分得到数值通量。Godunov 方法的经典自相似界面状态恒定，因此最后一步是平凡的，而 ADER 必须显式表示界面随时间的高阶演化。源项以两条途径进入格式：一条是高阶数值源项体积分，另一条是广义黎曼问题使数值通量本身含有源项信息。

ADER 是全离散、单步的方法，每步只进行一次重构，既可置于有限体积框架，也可用于间断 Galerkin 有限元。方法也不限于一阶双曲系统；已有工作把黎曼问题概念推广到反应--扩散型抛物方程，并报告时空最高十阶的格式 [535]，相关工作见 [224]。任意阶框架使使用者能按需求选择复杂度，一阶 Godunov 方法则位于有限体积 ADER 层次的底端。

% [原PDF p689]
本章仅给出了有限体积 ADER 的基础入门，处理一维、含非刚性源项的双曲系统。长时间平流结果展示了二阶 TVD 的局限；ADER--DG 的畸变非结构网格试验则表明，直至十阶的格式仍能达到理论预期的收敛率。这一事实尤其重要，因为高阶方法在复杂网格上的实际阶数容易受到几何、求积和边界实现的共同限制。

双曲守恒律 ADER 的近期进展见 [89,175--177,419,173,171]。当前研究正在把它推广到带源的非线性双曲系统，尤其是主微分部分无法写成守恒律形式的所谓非守恒系统。相关的最新工作与通信见 [85,174]。在这些系统中，广义黎曼问题不仅要预测高阶导数，还必须处理路径依赖的非守恒乘积，因而其弱解定义与离散一致性都比守恒系统更微妙。

从算法观点看，ADER 的优势是把空间高阶和时间高阶放入同一个局部时空预测器，而不是用多级 Runge--Kutta 反复重构。其代价是局部预测问题更复杂。实际实现应同时核验：重构阶、求积阶、Cauchy--Kowalewski 变换阶数、界面求解器精度以及源项平衡；任一环节降阶都会限制全局精度。
进一步说，形式上的高阶展开并不足以证明整个程序高阶。边界闭合、非结构网格几何表示、限制器启动位置以及源项求积都可能成为误差瓶颈。书中采用故意畸变的网格和多个误差范数，就是为了排除规则网格上的偶然抵消。对新实现，至少应报告连续网格序列、观测收敛阶、误差--成本关系和长期相位误差；仅展示一幅视觉上光滑的结果，不能替代这些定量验证。

% [原PDF p690]
\begin{table}[htbp]
\centering\small
\caption{ADER--DG 求解线性化 Euler 方程的代表性收敛结果。$O2$--$O10$ 分别表示二阶至十阶；完整原表按各网格列出 $h,N_d,L_\infty,L_1,L_2$ 及三个观测阶。}
\begin{tabular}{c r c c c}
\toprule
阶数 & 最细网格单元数 & $L_\infty$ & $L_1$ & $L_2$\\
\midrule
$O2$  & 190610 & $8.2223\times10^{-3}$ & $6.3435\times10^{-1}$ & $3.7319\times10^{-2}$\\
$O4$  & 158842 & $1.2070\times10^{-4}$ & $2.8285\times10^{-3}$ & $1.8062\times10^{-4}$\\
$O6$  & 83328  & $1.5737\times10^{-5}$ & $1.0614\times10^{-3}$ & $5.6039\times10^{-5}$\\
$O8$  & 142848 & $1.3771\times10^{-7}$ & $1.4681\times10^{-5}$ & $7.7184\times10^{-7}$\\
$O9$  & 178560 & $1.7872\times10^{-8}$ & $1.3131\times10^{-6}$ & $7.0996\times10^{-8}$\\
$O10$ & 218240 & $1.1050\times10^{-9}$ & $9.6184\times10^{-8}$ & $5.2430\times10^{-9}$\\
\bottomrule
\end{tabular}
\end{table}
完整数据表显示，随着网格加密，各误差范数的观测阶逐渐接近设计阶。十阶格式最后两层在 $L_1$ 范数上的观测阶约为 $10.1$，在 $L_2$ 范数上约为 $10.0$；二阶格式则在细网格上约为 $2.4$--$2.7$。这些结果说明，即使网格高度不规则，高阶局部时空离散仍能维持预期的渐近收敛行为。

% [原PDF p691]
\begin{figure}[htbp]
\centering
\fbox{\parbox{0.80\textwidth}{\centering 对数坐标结果摘要：上图横轴为 $1/h$、纵轴为 $L_\infty$ 误差；下图横轴为 CPU 时间、纵轴仍为 $L_\infty$ 误差。曲线对应 $P0$ 至 $P9$ 元，即二阶至十阶 ADER--DG。}}
\caption{ADER--DG 的误差--网格尺度与误差--CPU 时间关系。原图由 M. Dumbser 博士提供。}
\end{figure}
随着多项式阶数提高，误差--网格曲线的斜率显著增大。误差--CPU 图则揭示更有实际意义的效率交叉：在较粗容差下，低阶格式可能因单步便宜而占优；当目标误差足够小时，高阶格式用远少于低阶格式的计算时间达到同等精度。这一比较必须在相同误差范数、相同物理问题和同等实现优化下进行，不能只比较每单元成本。

图中 $P0$ 到 $P9$ 表示逐渐升高的间断 Galerkin 多项式次数。ADER 的局部时空预测与 DG 空间表示配合后，每条曲线都呈现其设计阶数对应的渐近斜率，因而也为前述任意阶构造提供了独立数值证据。

% [原PDF p692]
\begin{figure}[htbp]
\centering
\fbox{\parbox{0.80\textwidth}{\centering 网格拓扑摘要：（a）极粗、强畸变的非结构三角网格；（b）同一几何上的加密网格。计算域约为 $[-50,50]\times[-50,50]$，节点分布刻意偏离规则排列。}}
\caption{线性化 Euler 方程 ADER--DG 收敛研究所用网格拓扑。原图由 M. Dumbser 博士提供。}
\end{figure}
采用畸变网格而非规则网格，可以避免把特殊对称性或误差抵消误认为方法本身的高阶性质。粗、细两套拓扑保持同一类不规则性，从而让测得的阶数反映离散在一般非结构网格上的真实表现。表 20.1 与图 20.5、20.6 共同构成验证链：网格序列已知，误差随 $h$ 的斜率匹配设计阶，误差随成本的曲线又说明高阶精度在小误差区具有计算收益。

% [原PDF p693]
\chapter{结语}
本章首先概括全书的主要脉络，指出未展开的相关主题并补充较新的参考资料；随后讨论这些方法的潜在工程应用与当时活跃的研究方向，最后简要介绍 NUMERICA 程序库。

\section{数值方法要点回顾}
本书研究了一族用于双曲守恒律系统的有限体积方法，重点是数值通量。所讨论的大量通量都围绕经典黎曼问题展开，并形成两条基本路线：迎风方法精确或近似求解黎曼问题，显式提取微分方程中的波传播信息；中心方法避免直接求黎曼问题，构造通量时不显式使用特征波信息。

第 4、9、10、11、12 章的大部分内容用于不同黎曼求解器。第 8 章的通量矢量分裂（FVS）则提供另一种迎风通量路线，并比较三种分裂。第 7 章的随机选择法较为特殊：它使用第 4 章的经典黎曼精确解，对间断具有无限分辨率，但难以推广到两个以上独立变量的一般系统。中心路线从 Lax--Wendroff、Lax--Friedrichs 和 Godunov 中心格式出发，第 7 章又引入一维 FORCE 通量。

这些方法的共同问题不是“是否使用有限体积”，而是如何在界面上编码传播方向、耗散强度与守恒收支。黎曼求解器把这些信息集中在局部波结构中，中心通量则通过对称平均和稳定化耗散间接体现它们。

% [原PDF p694]
FORCE 与 Tadmor 及其合作者倡导的中心方法有关，第 17 章已将它推广到结构和非结构网格的多维问题。

显式利用黎曼解的数值通量可分为完整与不完整两类。完整求解器的近似波结构包含精确黎曼解的全部特征场；遗漏任一特征场时即为不完整求解器。对孤立激波，完整、不完整、FVS 和中心通量的表现往往相近；但对接触间断、剪切波等中间波，其分辨率差异很大。完整求解器通常最准确。除 Liou--Steffen 分裂外，其余 FVS 或中心通量常含过量人工黏性，这也影响 Navier--Stokes 计算中的剪切层分辨率。

一阶层面强调单调通量。在掌握若干单调数值通量后，本书主要讨论了二阶推广。Godunov 定理（第 13 章）表明，有实用价值且高于一阶的格式必须是非线性的，因此第 13 章引入 TVD 准则，第 14 章再以经验方式把它推广到非线性系统。

第 19 章研究广义黎曼问题，第 20 章以此建立一维模型方程的 ADER。ADER 是显式、全离散、单步、非线性且时空任意阶的方法，可以在有限体积或间断 Galerkin 框架中实现。高阶非线性版本采用 ENO/WENO 重构，二阶版本也可施加 TVD 条件。它与经典 ENO/WENO 有许多共同特征，相关资料已在第 20 章列出。

尚未纳入本书的重要主题包括 Colella--Woodward 的分片抛物线方法 PPM [137,584]、隐式方法 [239,595,593] 以及真正的多维迎风技术 [411,152,563,153]。这些方法均可在本书建立的守恒、重构和界面通量框架上继续理解。

% [原PDF p695]
多维迎风的进一步资料见 [197,22]，较新进展见 [113,159,2] 及其引用。书中迎风与中心通量不仅可用于有限体积，也能嵌入间断 Galerkin 有限元；后者的背景资料可参阅 Cockburn 及合作者的 [125,123,121,126,128]。

与本书主题相关的教材已经很多，包括 Godlewski--Raviart [215]、Laney [298]、Kulikovskii--Pogorelov--Semenov [292]、Oran--Boris [366]、Toro [520]、LeVeque [311]、Guinot [229]、Ben--Artzi--Falcovitz [39]、Drikakis--Rider [165]。综述文集 [545] 收录 97 篇关于 Godunov 方法及应用的同行评审文章；间断 Galerkin 和有限元方向还可参阅 [122,248,581,189]。

\section{潜在应用}
本书方法可用于范围很广的实际问题，Roe [415] 展示了若干颇具雄心的迎风方法应用。

浅水方程是环境流动建模中最常见的双曲系统之一，本书的通量、重构和源项处理均可直接用于它。代表性工作包括 [339,207,504,576,4,51,197,50,194,73]；Toro--Garc\'ia--Navarro 的综述 [534] 给出大量更新资料，另见教材 [520]。

Navier--Stokes 方程是另一类重要应用。定常不可压缩方程可用 Chorin [109] 的人工可压缩性改写成第 1 章所述系统，其无黏部分为双曲型，因而可使用本书方法；同一思路也可扩展到非定常不可压缩流。早期工作包括 Marx [345,346]、Drikakis [162,163]、Weinan--Shu [579]、Turkel--Arnone [550]、H\"anel--Sharma [235]、Tamamidis--Zhang--Assanis [481]、Toro [514] 等。

% [原PDF p696]
更新的不可压缩高阶方法见 Bassi 等 [30] 及其参考文献，也可参阅 Drikakis--Rider [165]。可压缩 Navier--Stokes 是更自然的下一步推广，H\"anel 的 VKI 讲义 [234] 很有价值；经典工作包括 [565,82,529,166,5]。

把 TVD 用于黏性问题时必须谨慎。TVD 本质上通过数值黏性控制大梯度附近的伪振荡，而标准 TVD 条件并未计入物理黏性，因而可能造成过量总耗散。多数问题仍能获得满意结果，但在关注精细流动物理时，这一局限会显现。标准 TVD 是为无黏问题构造的；黏性问题需要重新设计与物理扩散相协调的 TVD 条件，可能的推广见 [503]，更先进方法见 [31,127,129]。

Euler 方程还能描述一般可压缩介质中的波传播。极端热力学条件下必须采用真实气体状态方程，Godunov 型格式因而需要求解一般状态方程的黎曼问题。介质也可能是高压下可压缩的液体；相关求解器见 Ivings--Causon--Toro [264]。Menikoff--Plohr 的综述 [349] 尤其值得推荐，另见 [135,208,209,429,398,356]。同样的技术还可用于固体波传播 [284,320--323,199,491]。

气体和固体中的爆轰波是另一大应用领域。经典理论资料见 [365, 191, 583, 117, 116, 180, 201, 204]，数值计算工作见 [479, 367, 368, 276, 136, 78, 77, 448, 449, 118, 363]，以及 Oran--Boris 的专著 [366]。

% [原PDF p697]
天体物理大量采用本书讨论的数值方法，常见方程包括相对论 Euler、理想和相对论磁流体力学（MHD）。Mart\'i 与 M\"uller [342] 首次报告相对论气体动力学黎曼问题的精确解，扩展了 Thompson [488] 的激波管解，并说明它对 Godunov 型相对论流体计算的重要性；另见 [431,386,338,341--343,187]。天体流模拟的系统资料可参阅 LeVeque、Mihalas、Dorfi、M\"uller [312]。

早期 MHD 工作包括 [80, 566, 146, 147, 482, 29, 188, 26, 27]，较新工作包括 [154, 230, 255, 318, 351, 285, 28, 425, 68, 171]。Giacomazzo--Rezzolla [202] 得到了相对论 MHD 黎曼问题的精确解，为评估数值格式提供了宝贵基准。数值相对论也已成为迅速发展的应用方向，综述见 Lehner [306]。

多相流通常由不同材料组成，例如气相或液相载体中的固体颗粒。“相”也可广义理解为具有不同材料性质的场或流体，如空气与蒸汽。多场模型可包含许多相，每一相通常拥有自己的速度矢量；这不同于多组分流，后者只需载体流体的一套速度。多相流背景教材包括 Drew--Gulliver [160]、Drew--Marsden--Sirovich [161]、Jackson [265]、Gidaspow [203]、Ishii--Hibiki [261]、St\"adtke [460]、Brennen [79]、Prosperetti--Tryggvason [388]。

% [原PDF p698]
多相模型一般具有三个显著特点：方程含源项；可能含非守恒乘积；忽略耗散时，系统甚至可能呈椭圆--双曲混合型。Stewart--Wendroff [464] 的经典论文讨论了其中若干问题。非守恒乘积的研究相对较新；与此同时，也已有守恒形式的多相模型，例如 Romenski 等 [417]。

将 Godunov 方法和黎曼问题用于可压缩多相流的早期尝试见 [501,525]，其后工作包括 [421,427,428,549,547,548,142,426,461,90,102]。大量研究围绕 Baer--Nunziato 方程 [21] 展开。Embid--Baer [183] 给出全面的数学分析；Andrianov--Warnecke [14] 求解间接黎曼问题，即先假定解再反求初态，这类解适合检验数值方法。

Schwendeman--Wahle--Kapila [438] 似乎最早求得 Baer--Nunziato 方程的直接黎曼问题，因此既能作基准，也可构造 Godunov 通量。Deledicque--Miltiadis [156] 也求解了类似模型。基于 [438] 的经典黎曼精确解，Toro--Castro [531] 又求解该系统的广义黎曼问题。Dumbser、Hidalgo、Castro、Par\'es、Toro [174] 最近报告了适用于三维非结构网格的新路径守恒格式。

自由表面两相流模型也在发展，起点之一是 Pitman--Le [385]；相应数值方法见 Pelanti--Bouchut--Mangeney [376]、Rhebergen--Bokhove--van der Vegt [401] 以及 Dumbser--Castro--Par\'es--Toro [172]。

% [原PDF p699]
\section{当前研究课题}
非守恒乘积近年来受到数值分析界的广泛关注。它意味着方程主部不能写成守恒或散度形式，从数学与物理层面都带来如何定义弱解的问题。Dal Maso--LeFloch--Murat [148] 给出了数学处理，早期工作另见 Volpert [571]。数值上正确离散非守恒项同样困难。Toumi [546] 和 Toumi--Kumbaro [549] 基于 [148] 的开创工作推动了适合非守恒系统的新方法；近期路径守恒研究见 Par\'es 及合作者 [375,93,94,85,172,174]，另见 [401]。

源项几乎存在于所有实际系统中，多相流的各相尤其通过源项耦合。虽然近年已有显著进展，正确处理源项仍是严肃挑战。Roe [412] 提出利用波传播信息对源项迎风离散，是重要进展；更基本的要求是格式应保持平衡。Glimm--Marshall--Plohr [214] 可能最早明确处理这一问题，早期研究还包括 [51,569,225,226,310,220]，较新工作见 [196,92,94,85,172,174]。

许多应用中的源项具有刚性。Dumbser--Enaux--Toro [173] 在高阶 ADER 框架中提出局部隐式的广义黎曼求解器，同时保持整体显式的单步更新，从而协调高阶精度与刚性。

新黎曼求解器仍是活跃课题。理想目标是构造一个包含全部特征场、能准确分辨中间波、又兼具鲁棒性和通用性的完整求解器。Toro [496] 的 MUSTA 预测--校正多级方法具有这种潜力：先用若干级把初始数据局部演化，再由校正级生成所需界面通量。

% [原PDF p700]
MUSTA 的进一步发展见 [523,524,521,495,357]。名为 EVILIN 的一个有前景变体 [522] 以简单中心型通量作预测器，在校正步使用精确黎曼解的简单线性化，从而获得完整黎曼求解器。EVILIN 具有一定通用性，但其鲁棒性仍需进一步研究。

\section{NUMERICA 程序库}
NUMERICA [518] 是面向教学、研究和应用的源代码库，其内容以本书和 [520] 为基础。程序首先可作为学习工具：自学者能通过阅读实现理解数值方法细节；课堂上也可用它演示算法的主要步骤及一种可行编码方式。其次，这些程序可作为开发原型，既容易改造成实际应用工具，也能成为更大型工程程序的起点。

有关 NUMERICA，可直接联系作者 E. F. Toro 教授：意大利特伦托大学土木与环境工程系应用数学实验室；电子邮件 \texttt{toro@ing.unitn.it}；网站 \texttt{http://www.ing.unitn.it/toro}。该库对学术用户免费，可从网站下载；非学术用户可联系 \texttt{inquiries@numeritek.com}。上述地址按原书保留，使用时应注意网站和电子邮件可能已经变更。
作为全书的收束，NUMERICA 体现了作者反复强调的思想：数值方法应当能从数学结构追溯到清晰、可检验的实现。教学代码虽然不等同于生产软件，却应明确展示状态重构、黎曼求解、通量装配、源项离散、稳定时间步和边界处理之间的接口。研究者可在这些小程序上替换单个模块并作受控比较，再逐步推广到多维、复杂状态方程、多相和刚性源项问题。
这种由简入繁的路线也便于建立回归测试：先以精确黎曼解和标准算例验证基本通量，再逐层加入高阶重构、源项、复杂网格和多物理耦合。每次扩展都应保留可重复的守恒误差、收敛阶和鲁棒性记录，使代码演化与数学假设保持一致。
\chapter*{参考文献}
\addcontentsline{toc}{chapter}{参考文献}
\markboth{参考文献}{参考文献}
\begingroup\small\sloppy
\setlength{\parindent}{0pt}
\setlength{\parskip}{0.25em}
% [原PDF p701]
\par\hangindent=2em\hangafter=1 1. R. Abgrall. How to Prevent Pressure Oscillations in Multicomponent Flow Calculations: A Quasiconservative Approach. J. Comput. Phys., 125:150–160, 1996.
\par\hangindent=2em\hangafter=1 2. R. Abgrall and M. Mezine. Construction of Second–Order Accurate Monotone and Stable Residual Distribution Schemes for Steady Problems. J. Comput. Phys., 195:474–507, 2004.
\par\hangindent=2em\hangafter=1 3. C. M. Albone. Report on the AGCFM Working Party on Software Qual- ity Assurance. Technical Report AERO 2105, Royal Aircraft Establishment, Farnborough, UK, 1987.
\par\hangindent=2em\hangafter=1 4. A. Alcrudo and P. Garc´ıa-Navarro. A High Resolution Godunov–Type Scheme in Finite Volumes for the 2D Shallow Water Equations. Int. J. Num. Meth. Fluids, 16:489–505, 1993.
\par\hangindent=2em\hangafter=1 5. S. R. Allmaras. Contamination of Laminar Boundary Layers by Artificial Dissi- pation in Navier–Stokes Solutions. In Numerical Methods for Fluid Dynamics, 4, pages 443–449. Oxford University Press, 1993.
\par\hangindent=2em\hangafter=1 6. V. R. Ambati and O. Bokhove. Space–Time Discontinuous Galerkin Discretiza- tion of Rotating Shallow Water Equations. J. Comput. Phys. , 225(2):1233– 1261, 2007.
\par\hangindent=2em\hangafter=1 7. D. A. Anderson, J. C. Tannehill, and R. H. Pletcher. Computational Fluid Mechanics and Heat Transfer. Hemisphere Publishing Corporation, 1984.
\par\hangindent=2em\hangafter=1 8. D. S. Anderson and J. J. Gottlieb. On Random Numbers for the Random Choice Method. Technical Report 258, UTIAS, University of Toronto, Canada, 1987.
\par\hangindent=2em\hangafter=1 9. J. D. Anderson. Hypersonic and High Temperature Gas Dynamics. Mc Graw– Hill, 1989.
\par\hangindent=2em\hangafter=1 10. J. D. Anderson. Modern Compressible Flow. Mc Graw–Hill, 1990.
\par\hangindent=2em\hangafter=1 11. J. L. Anderson, S. Preiser, and E. L. Rubin. Conservation Form of the Equa- tions of Hydrodynamics in Curvilinear Coordinate Systems.J. Comput. Phys., 2:279–287, 1968.
\par\hangindent=2em\hangafter=1 12. W. Anderson, J. L. Thomas, and B. van Leer. Comparison of Finite Volume Flux Vector Splittings for the Euler Equations. AIAA Journal, 24:1453–1460, 1986.
% [原PDF p702]
\par\hangindent=2em\hangafter=1 13. W. K. Anderson, J. L. Thomas, and L. Rumsey. Extension and Application of Flux–Vector Splitting to Calculations on Dynamic Meshes. AIAA Journal, 27:673–674, 1989.
\par\hangindent=2em\hangafter=1 14. N. Andrianov and G. Warnecke. The Riemann Problem for the Baer–Nunziato Two–Phase Flow Model. J. Comput. Phys., 212:434–464, 2004.
\par\hangindent=2em\hangafter=1 15. T. M. Apostol. Mathematical Analysis. Addison–Wesley Publishing Company, 1974.
\par\hangindent=2em\hangafter=1 16. P. Arminjon and A. St-Cyr. Nessyahu–Tadmor Type Central Finite Volume Methods Without Predictor for 3D Cartesian and Unstructured Tetrahedral Grids. Applied Numerical Mathematics, 46:135–155, 2003.
\par\hangindent=2em\hangafter=1 17. P. Arminjon, A. St-Cyr, and A. Madrane. New Two– and Three Dimensional Non-Oscillatory Central Finite Volume Methods on Staggered Cartesian Grids. Applied Numerical Mathematics, 40:367–390, 2002.
\par\hangindent=2em\hangafter=1 18. P. Arminjon and R. Touma. Central Finite Volume Methods with Constrained Transport Divergence Treatment for Ideal MHD. J. Comput. Phys., 204:737– 759, 2005.
\par\hangindent=2em\hangafter=1 19. M. Arora and P. L. Roe. On Postshock Oscillations Due to Shock Capturing Schemes in Unsteady Flows. J. Comput. Phys., 130:25–40, 1997.
\par\hangindent=2em\hangafter=1 20. E. H. Atta and J. Vadyak. A Grid Overlapping Scheme for Flowfield Computa- tions About Multicomponent Configurations. AIAA Journal, 21(9):1271–1277, 1983.
\par\hangindent=2em\hangafter=1 21. M. R. Baer and J. W. Nunziato. A Two–Phase Mixture Theory for the Deflagration–to–Detonation Transition (DDT) in Reactive Granular Materi- als. J. Multiphase Flow, 12:861–889, 1986.
\par\hangindent=2em\hangafter=1 22. M. J. Baines and M. E. Hubbard. Multidimensional Upwinding with Grid Adaptation. In Numerical Methods for Wave Propagation. Toro, E. F. and Clarke, J. F. (Editors), pages 33–54. Kluwer Academic Publishers, 1998.
\par\hangindent=2em\hangafter=1 23. M. J. Baines and P. K. Sweby. A Flux Limiter Based on Fromm’s Scheme. Technical Report 9/84, Department of Mathematics, University of Reading, UK, 1984.
\par\hangindent=2em\hangafter=1 24. G. J. Ball and R. A. East. Shock and Blast Attenuation by Aqueous Foam Barriers: Influences of Barrier Geometry. J. Shock Waves, 9(1):37–47, 1999.
\par\hangindent=2em\hangafter=1 25. D. Balsara and C. W. Shu. Monotonicity Preserving Weighted Essentially Nonoscillatory Schemes with Increasingly High Order of Accuracy.J. Comput. Phys., 160:405–452, 2000.
\par\hangindent=2em\hangafter=1 26. D Balsara and D. Spicer. A Staggered Mesh Algorithm Using High Order Go- dunov Fluxes to Ensure Solenoidal Magnetic Fields in Magnetohydrodynamic Simulations. J. Comput. Phys., 149:270–292, 1999.
\par\hangindent=2em\hangafter=1 27. D. Balsara and D. Spicer. Maintaining Pressure Positivity in Magneto– hydrodynamic Simulations. J. Comput. Phys., 149:133–148, 1999.
\par\hangindent=2em\hangafter=1 28. D. S. Balsara, T. Rumpf, M. Dumbser, and C. D. Munz. Efficient, High Accu- racy ADER–WENO Schemes for Hydrodynamics and Divergence–Free Mag- netohydrodynamics. J. Comput. Phys., 2008 (Submitted).
\par\hangindent=2em\hangafter=1 29. A. A. Barmin, A. G. Kulikovskii, and N. V. Pogorelov. Shock–Capturing Ap- proach and Nonevolutionary Solutions in Magnetohydrodynamics.J. Comput. Phys., 126:77–90, 1996.
\par\hangindent=2em\hangafter=1 30. F. Bassi, A. Crivellini, D. A. Di Pietro, and S. S. Rebay. An Artificial Com- pressibility Flux for the Discontinuous Galerkin Solution of the Incompressible Navier–Stokes Equations. J. Comput. Phys., 218:794–815, 2006.
% [原PDF p703]
\par\hangindent=2em\hangafter=1 31. F. Bassi and S. S. Rebay. High–Order Accurate Discontinuous Finite Ele- ment Method for the Numerical Solution of the Compressible Navier–Stokes Equations. J. Comput. Phys., 131:267–279, 1997.
\par\hangindent=2em\hangafter=1 32. P. Batten, N. Clarke, C. Lambert, and D. Causon. On the Choice of Wave Speeds for the HLLC Riemann Solver.SIAM J. Sci. and Stat. Comp., 18:1553– 1570, 1997.
\par\hangindent=2em\hangafter=1 33. P. Batten, M. A. Leschziner, and U. C. Goldberg. Average–State Jacobians and Implicit Methods for Compressible Viscous and Turbulent Flows Flows. J. Comput. Phys., 137:38–78, 1997.
\par\hangindent=2em\hangafter=1 34. A. Bayliss and E. Turkel. Radiation–Boundary Conditions for Wave–Like Equations. Comm. Pure and Appl. Maths., 33:708–725, 1980.
\par\hangindent=2em\hangafter=1 35. E. Becker. Gas Dynamics. Academic Press, 1968.
\par\hangindent=2em\hangafter=1 36. M. Ben-Artzi. Application of the Generalised Riemann Problem Method to 1–D Compressible Flows with Interfaces. J. Comput. Phys., 65:170–178, 1986.
\par\hangindent=2em\hangafter=1 37. M. Ben-Artzi and J. Falcovitz. A Second Order Godunov–Type Scheme for Compressible Fluid Dynamics. J. Comput. Phys., 55:1–32, 1984.
\par\hangindent=2em\hangafter=1 38. M. Ben-Artzi and J. Falcovitz. A High Resolution Upwind Scheme for Quasi 1–D Flows. In Numerical Methods for the Euler Equations of Fluid Dynamics, pages 66–83. INRIA, SIAM, 1985.
\par\hangindent=2em\hangafter=1 39. M. Ben-Artzi and J. Falcovitz. Generalized Riemann Problems in Computa- tional Fluid Dynamics. Cambridge University Press, 2003.
\par\hangindent=2em\hangafter=1 40. M. Ben-Artzi and J. Li. Hyperbolic Balance Laws: Riemann Invariants and Hyperbolic Balance Laws. Numerische Mathematik, 106:69–425, 2007.
\par\hangindent=2em\hangafter=1 41. M. Ben-Artzi, J. Li, and G. Warnecke. A Direct Eulerian GRP Scheme for Compressible Fluid Flows. J. Comput. Phys., 218(19–43), 2006.
\par\hangindent=2em\hangafter=1 42. G. Ben-Dor. Shock Wave Reflection Phenomena. Springer-Verlag, 1992.
\par\hangindent=2em\hangafter=1 43. J. A. Benek, J. L. Steger, and F. C. Dougherty. A Flexible Grid Embedding Technique with Application to the Euler Equations. AIAA Paper 83–1944, 1983.
\par\hangindent=2em\hangafter=1 44. F. Benkhaldoun and A. Chalabi. Characteristic Based Scheme for Conserva- tion Laws with Source Terms. Technical Report 95.15, Centre National de la Recherche Scientifique, Universit´e Paul Sabatier, UFR-MIG, 118 route de Narbone, 31062, Toulouse Cedex, France, 1995.
\par\hangindent=2em\hangafter=1 45. M. J. Berger. Adaptive Mesh Refinement for Hyperbolic Partial Differential Equations. PhD thesis, Dept. of Computer Science, Stanford University, Cali- fornia, 1982.
\par\hangindent=2em\hangafter=1 46. M. J. Berger and P. Colella. Local Adaptive Mesh Refinement for Shock Hy- drodynamics. J. Comput. Phys., 82:64–84, 1989.
\par\hangindent=2em\hangafter=1 47. M. J. Berger and R. J. LeVeque. Stable Boundary Conditions for Cartesian Grid Calculations. Technical Report ICASE 90–37, NASA Langley Research Center, USA, 1990.
\par\hangindent=2em\hangafter=1 48. M. J. Berger and J. Oliger. Adaptive Mesh Refinement for Hyperbolic Partial Differential Equations. J. Comput. Phys., 53:484–512, 1984.
\par\hangindent=2em\hangafter=1 49. A. C. Berkenbosch, E. F. Kaasschieter, and J. H. M. ten Thije Boonkkamp. The Numerical Computation of One–dimensional Detonation Waves. Techni- cal Report RANA 94–21, Eindhoven University of Technology, Department of Mathematics and Computer Science, The Netherlands, 1994.
\par\hangindent=2em\hangafter=1 50. L. Berm´u d e z ,A .D e r v i e u x ,J .A .D e s i d e r i ,a n dM .E .V ´azquez. Upwind Schemes for the Two–Dimensional Shallow Water Equations with Variable Depth Using
% [原PDF p704]
\par\hangindent=2em\hangafter=1 Unstructured Meshes. Technical Report 2738, INRIA, Rocquencourt, 78153 Le Chesnay, France, 1995.
\par\hangindent=2em\hangafter=1 51. L. Berm´udez and M. E. V´azquez. Upwind Methods for Hyperbolic Conserva- tion Laws with Source Terms. Computers and Fluids, 23:1049–1071, 1994.
\par\hangindent=2em\hangafter=1 52. H. Bershader and R. Hanson (Editors). Shock Waves and Shock Tubes (Proc. 15th Int. Symp. on Shock Waves and Shock Tubes, Berkeley, California, USA 1985). Stanford University Press, 1986.
\par\hangindent=2em\hangafter=1 53. C. Berthon, P. Charrier, and B. Dubroca. An HLLC Scheme to Solve The M1 Model of Radiative Transfer in Two Space Dimensions. Journal of Scientific Computing, 31(3):347–389, 2007.
\par\hangindent=2em\hangafter=1 54. C. Berthon, F. Coquel, J. M. H´ erard, and M. Uhlmann. An Approximate Solution of the Riemann Problem for a Realisable Second–Moment Turbulent Closure. J. Shock Waves, 11(2):245–269, 2002.
\par\hangindent=2em\hangafter=1 55. M. Berzins, P. M. Dew, and R. M. Fuzerland. Developing Software for Time– Dependent Problems Using the Method of Lines and Differential Algebraic Integrators. Applied Numerical Mathematics, 5:375–397, 1989.
\par\hangindent=2em\hangafter=1 56. M. Berzins and J. M. Ware. Solving Convection and Convection Reaction Problems using M.O.L. Applied Numerical Mathematics, 20:83–89, 1996.
\par\hangindent=2em\hangafter=1 57. F. Bianco, G. Puppo, and G. Russo. High Order Central Schemes for Hyper- bolic Systems of Conservation Laws. SIAM J. Scientific Computing , 21:294– 322, 1999.
\par\hangindent=2em\hangafter=1 58. S. J. Billett. Numerical Aspects of Convection Diffusion Problems. MSc. Thesis, College of Aeronautics, Cranfield Institute of Technology, UK, 1991.
\par\hangindent=2em\hangafter=1 59. S. J. Billett. A Class of Upwind Methods for Conservation Laws. Ph.D thesis, College of Aeronautics, Cranfield University, UK, 1994.
\par\hangindent=2em\hangafter=1 60. S. J. Billett and E. F. Toro. Restoring Monotonicity of Slowly–Moving Shocks Computed with Godunov–Type Schemes.CoA Report 9207, Cranfield Institute of Technology, UK, 1992.
\par\hangindent=2em\hangafter=1 61. S. J. Billett and E. F. Toro. Implementing a Three–Dimensional Finite Volume WAF–Type Scheme for the Euler Equations. InProc. Third ECCOMAS Com- putational Fluid Dynamics Conference, 9–13 September 1996, Paris, France , pages 732–738. Wiley, 1996.
\par\hangindent=2em\hangafter=1 62. S. J. Billett and E. F. Toro. On the Accuracy and Stability of Explicit Schemes for Multidimensional Linear Homogeneous Advection Equations. J. Comp. Phys., 131:247–250, 1997.
\par\hangindent=2em\hangafter=1 63. S. J. Billett and E. F. Toro. WAF–Type Schemes for Multidimensional Hyper- bolic Conservation Laws. J. Comp. Phys., 130:1–24, 1997.
\par\hangindent=2em\hangafter=1 64. S. J. Billett and E. F. Toro. Unsplit WAF–Type Schemes for Three Dimensional Hyperbolic Conservation Laws. In Numerical Methods for Wave Propagation. Toro, E. F. and Clarke, J. F. (Editors), pages 75–124. Kluwer Academic Pub- lishers, 1998.
\par\hangindent=2em\hangafter=1 65. S. J. Billett and E. F. Toro. TVB Limiter Functions for Multidimensional Schemes. In Proc. 6th Intern. Symposium on Computational Fluid Dynamics, Lake Tahoe, Nevada, USA, pages 117–122, September 4-8, 1995.
\par\hangindent=2em\hangafter=1 66. E. P. Boden. An Adaptive Gridding Technique for Conservation Laws on Com- plex Domains . PhD thesis, Department of Aerospace Science, Cranfield Uni- versity, UK, 1997.
\par\hangindent=2em\hangafter=1 67. E. P. Boden and E. F. Toro. A Combined Chimera–AMR Technique for Hyper- bolic PDEs. In Proc. Fifth Annual Conference of the CFD Society of Canada, pages 5.13–5.18. University of Victoria, Canada, 1997.
% [原PDF p705]
\par\hangindent=2em\hangafter=1 68. G. Bodo, S. Masaglia, A. Mignone, and P. Rossi (Editors). Jets from Young Stars III. Numerical MHD and Instabilities. Lecture Notes in Physics 754 . Springer, 2008.
\par\hangindent=2em\hangafter=1 69. D. L. Book, J. P. Boris, and K. Hain. Flux Corrected Transport II: General- izations of the Method. J. Comput. Phys., 18:248–283, 1975.
\par\hangindent=2em\hangafter=1 70. J. P. Boris and D. L. Book. Flux–Corrected Transport. I. SHASTA, A Fluid Transport Algorithm that Works. J. Comput. Phys., 11:38–69, 1973.
\par\hangindent=2em\hangafter=1 71. J. P. Boris and D. L. Book. Flux–Corrected Transport III: Minimal–Error FCT Algorithms. J. Comput. Phys., 20:397–431, 1976.
\par\hangindent=2em\hangafter=1 72. J. P. Boris and E. S. Oran. Numerical Simulation of Reactive Flow. Elsevier, Amsterdam–New York, 1987.
\par\hangindent=2em\hangafter=1 73. A. G. L. Borthwick, M. Fujihara, and B. D. Rogers. Godunov Solutions of Shallow Water Equations on Curvilinear and Quadtree Grids. In Godunov Methods: Theory and Applications. Edited Review, E. F. Toro (Editor), pages 141–148. Kluwer Academic/Plenum Publishers, 2001.
\par\hangindent=2em\hangafter=1 74. F. Bouchut. Nonlinear Stability of Finite Volume Methods for Hyperbolic Con- servation laws and Well–Balanced Schemes. Birkh¨auser, 2004.
\par\hangindent=2em\hangafter=1 75. F. Bouchut, C. Klingenberg, and K. Waagan. A Multiwave Approximate Rie- mann Solver for Ideal MHD Based on Relaxation I - Theoretical Framework. Numer. Math., 108(1):7–42, 2007.
\par\hangindent=2em\hangafter=1 76. A. Bourgeade, P. Le Floch, and Raviart P. A. An Asymptotic Expansion for the Solution of the Generalized Riemann Problem. Part 2: Application to the Euler Equations of Gas Dynamics. Ann. Inst. Henri Poincar´ e. Analyse non Line´are, 6(6):437–480, 1989.
\par\hangindent=2em\hangafter=1 77. A. Bourlioux and A. J. Majda. Theoretical and Numerical Structure for Unsta- ble Two–Dimensional Detonations. Combustion and Flame, 90:211–229, 1992.
\par\hangindent=2em\hangafter=1 78. A. Bourlioux, A. J. Majda, and V. Roytburd. Theoretical and Numerical Structure for Unstable One–Dimensional Detonations. SIAM J. Appl. Math., 51:303–343, 1991.
\par\hangindent=2em\hangafter=1 79. C. E. Brennen. Fundamentals of Multiphase Flow. Cambridge University Press, 2005.
\par\hangindent=2em\hangafter=1 80. M. Brio and C. C. Wu. An Upwind Differencing Scheme for the Equations of Ideal Magnetohydrodynamics. J. Comput. Phys., 75:400–422, 1988.
\par\hangindent=2em\hangafter=1 81. M. Brown. Differential Equations and their Applications. Springer–Verlag, 1975.
\par\hangindent=2em\hangafter=1 82. R. E. Brown. Numerical Solution of the 2D Unsteady Navier–Stokes Equations Using Viscous–Convective Operator Splitting. MSc. Thesis, Department of Aerospace Science, Cranfield University, UK, 1990.
\par\hangindent=2em\hangafter=1 83. R. Brun and L. Z. Dumitrescu (Editors). Shock Waves @ Marseille (Proc. 19th Int. Symp. on Shock Waves, Marseille, France, 1993), Four Volumes . Springer–Verlag, 1995.
\par\hangindent=2em\hangafter=1 84. M. C´aceres. Development of a Third–Order Accurate Scheme of the MUSCL Type for the Time–Dependent One–Dimensional Euler Equations. MSc. The- sis, Department of Aerospace Science, Cranfield University, UK, 1993.
\par\hangindent=2em\hangafter=1 85. A. Canestrelli, A. Siviglia, M. Dumbser, and E. F. Toro. Well–Balanced High– Order Centred Schemes for Non–Conservative Hyperbolic Systems. Applica- tion to Shallow Water with Fixed and Mobile Bed. (Submitted), 2008.
\par\hangindent=2em\hangafter=1 86. G. Capdeville. Towards a Compact High–Order Method for Non–linear Hyperbolic Systems, II. The Hermite–HLLC Scheme. J. Comput. Phys. , 227(22):9428–9462, 2008.
% [原PDF p706]
\par\hangindent=2em\hangafter=1 87. Carasso, Raviart, and Serre. Proc. First International Conference on Hyper- bolic Problems. Springer, 1986.
\par\hangindent=2em\hangafter=1 88. J. Casper and H. L. Atkins. A Finite Volume High–Order ENO Scheme for Two-Dimensional Hyperbolic Systems. J. Comput. Phys., 106:62–76, 1993.
\par\hangindent=2em\hangafter=1 89. C. E. Castro. High–Order ADER FV/DG Numerical Methods for Hyperbolic Equations. PhD thesis, Department of Civil and Environmental Engineering, University of Trento, Italy, 2007.
\par\hangindent=2em\hangafter=1 90. C. E. Castro and E. F. Toro. A Riemann Solver and Upwind Methods for a Two–Phase Flow Model in Non–Conservative Form. Int. J. for Numer. Meth. in Fluids, 60:275–307, 2006.
\par\hangindent=2em\hangafter=1 91. C. E. Castro and E. F. Toro. Solvers for the High–Order Riemann Problem for Hyperbolic Balance Laws. J. Comput. Phys., 227:2481–2513, 2008.
\par\hangindent=2em\hangafter=1 92. M. J. Castro, J. M. Gallardo, J. A. L´opez, and C. Par´es. Well–Balanced High Order Extensions of Godunov’s Method for Semilinear Balance Laws. SIAM Journal of Numerical Analysis, 46:1012–1039, 2008.
\par\hangindent=2em\hangafter=1 93. M. J. Castro, P. G. LeFloch, M. L. Mu˜ noz-Ruiz, and C. Par´es. Why Many Theories of Shock Waves are Necessary: Convergence Error in Formally Path– Consistent Schemes. J. Comput. Phys., 227:8107–8129, 2008.
\par\hangindent=2em\hangafter=1 94. M. J. Castro, A. Pardo, C. Par´es, and E. F. Toro. On Some Fast Well–Balanced First Order Solvers for Nonconservative Systems. Mathematics of Computa- tion, 2008 (submitted).
\par\hangindent=2em\hangafter=1 95. V. Casulli and E. F Toro. Preliminary Results on a Semi–Lagrangian Version of the WAF Method. Technical Report 9018, College of Aeronautics, Cranfield Institute of Technology, UK, 1990.
\par\hangindent=2em\hangafter=1 96. A. Chakraborty and E. F. Toro. Development of an Approximate Riemann Solver for the Steady Supersonic Euler Equations. The Aeronautical Journal, 98:325–339, 1994.
\par\hangindent=2em\hangafter=1 97. S. R. Chakravarthy, D. A. Anderson, and M. D. Salas. The Split–Coefficient Matrix Method for Hyperbolic Systems of Gas Dynamics Equations. AIAA Paper 80–0268, AIAA 19th Aerospace Science Meeting, 1980.
\par\hangindent=2em\hangafter=1 98. A. Chalabi. Stable Upwind Schemes for Conservation Laws with Source Term. IMA J. Numer. Anal., 12:217–241, 1992.
\par\hangindent=2em\hangafter=1 99. A. Chalabi. On Convergence of Numerical Schemes for Hyperbolic Conserva- tion Laws with Stiff Source Terms. Technical Report 95.22, Centre National de la Recherche Scientifique, Universit´e Paul Sabatier, UFR-MIG, 118 route de Narbone, 31062, Toulouse Cedex, France, 1995.
\par\hangindent=2em\hangafter=1 100. A. Chalabi. An Error Bound for the Polygonal Approximation of Hyperbolic C o n s e r v a t i o nL a w sw i t hS o u r c eT e r m s .Computers and Mathematics with Ap- plications, 32:59–63, 1996.
\par\hangindent=2em\hangafter=1 101. A. Chalabi. On Convergence of Numerical Schemes for Hyperbolic Conserva- tion Laws with Stiff Source Terms. Math. of Comput., 66:527–545, 1997.
\par\hangindent=2em\hangafter=1 102. C. H. Chang and M. S. Liou. A Robust and Accurate Approach to Computing Compressible Multiphase Flow: Stratified Flow Model and AUSM+ up Scheme. J. Comput. Phys., 225:840–873, 2007.
\par\hangindent=2em\hangafter=1 103. G. Q. Chen and E. F. Toro. Centred Schemes for Non–Linear Hyperbolic Equations. Journal of Hyperbolic Differential Equations, 1(1):531–566, 2004.
\par\hangindent=2em\hangafter=1 104. D. N. Cheney. Upwinding Convective and Viscous Terms via a Modified GRP Approach. MSc. Thesis, Department of Aerospace Science, Cranfield Univer- sity, UK, 1994.
% [原PDF p707]
\par\hangindent=2em\hangafter=1 105. G. Chesshire and W. D. Henshaw. Composite Overlapping Meshes for the Solution of Partial Differential Equations. J. Comput. Phys., 90:1–64, 1990.
\par\hangindent=2em\hangafter=1 106. G. Chesshire and W. D. Henshaw. A Scheme for Conservative Interpolation on Overlapping Grids. SIAM J. Sci. Stat., 15(4):819–845, 1994.
\par\hangindent=2em\hangafter=1 107. B. L. Chessor, Jr. Evaluation of Computed Steady and Unsteady Quasi–One– Dimensional Viscous/Inviscid Interacting Internal Flow Fields Through Com- parison with Two–Dimensional Navier–Stokes Solutions. MSc. Thesis, 1992, Mississippi State University, USA, 1992.
\par\hangindent=2em\hangafter=1 108. T. Chiang and Hoffmann K. Determination of computational time step for chemically reacting flows. AIAA 20th Fluid Dynamics, Plasma Dynamics and Lasers Conference, Buffalo, New York, June 12–14, USA, 1989, Paper AIAA 89–1855, 1989.
\par\hangindent=2em\hangafter=1 109. A. J. Chorin. A Numerical Method for Solving Viscous Flow Problems. J. Comput. Phys., 2:12–26, 1967.
\par\hangindent=2em\hangafter=1 110. A. J. Chorin. Random Choice Solutions of Hyperbolic Systems. J. Comput. Phys., 22:517–533, 1976.
\par\hangindent=2em\hangafter=1 111. A. J. Chorin. Random Choice Methods with Applications to Reacting Gas Flow. J. Comput. Phys., 25:253–272, 1977.
\par\hangindent=2em\hangafter=1 112. A. J. Chorin and J. E. Marsden. A Mathematical Introduction to Fluid Me- chanics. Springer–Verlag, 1993.
\par\hangindent=2em\hangafter=1 113. C. S. Chu and C. W. Shu. High Order Residual Distribution Conservative Fi- nite Difference WENO Schemes for Convection–Diffusion Steady State Prob- lems on Non–Smooth Meshes. J. Comput. Phys., 224:992–1020, 2007.
\par\hangindent=2em\hangafter=1 114. D. K. Clarke, M. D. Salas, and H. A. Hassan. Euler Calculations for Multiele- ment Airfoils Using Cartesian Grids. AIAA J., 24:353–, 1987.
\par\hangindent=2em\hangafter=1 115. J. F. Clarke. Compressible Flow Produced by Distributed Sources of Mass: An Exact Solution. Technical Report CoA Report 8710, College of Aeronautics, Cranfield Institute of Technology, Cranfield, UK, 1987.
\par\hangindent=2em\hangafter=1 116. J. F. Clarke. Chemical Reactions in High Speed Flows. Phil. Trans. Roy. Soc. London, A 335:161–199, 1989.
\par\hangindent=2em\hangafter=1 117. J. F. Clarke. Fast Flames, Waves and Detonations. Prog. Energy Comb. Sci., 15:241–271, 1989.
\par\hangindent=2em\hangafter=1 118. J. F. Clarke, S. Karni, J. J. Quirk, L. G. Simmons, P. L. Roe, and E. F. Toro. Numerical Computation of Two–Dimensional, Unsteady Detonation Waves in High Energy Solids. J. Comput. Phys., 106:215–233, 1993.
\par\hangindent=2em\hangafter=1 119. J. F. Clarke and M. McChesney. Dynamics of Relaxing Gases. Butterworths, 1976.
\par\hangindent=2em\hangafter=1 120. J. F. Clarke and E. F. Toro. Gas Flows Generated by Solid–Propellant Burning. In R. Glowinski, B. Larrouturou, and R. Teman, editors,Proc. Intern. Confer. on Numerical Simulation of Combustion Phenomena, number 241 in Lecture Notes in Physics, pages 192–205, INRIA, Sophia Antipolis, France, May 21–24 1985.
\par\hangindent=2em\hangafter=1 121. B. Cockburn, S. Hou, and C. W. Shu. The Runge–Kutta Local Projection Discontinuous Galerkin Finite Element Method for Conservation Laws IV: the Multidimensional Case. Mathematics of Computation, 54:545–581, 1990.
\par\hangindent=2em\hangafter=1 122. B. Cockburn, G. E. Karniadakis, and C. W. Shu. Discontinuous Galerkin Methods. Lecture Notes in Computational Science and Engineering. Springer, 2000.
% [原PDF p708]
\par\hangindent=2em\hangafter=1 123. B. Cockburn, S. Y. Li, and C. W. Shu. TVB Runge–Kutta Local Projection Discontinuous Galerkin Finite Element Method for Conservation Laws III: One Dimensional Systems. J. Comput. Phys., 84:90–113, 1989.
\par\hangindent=2em\hangafter=1 124. B. Cockburn and C. Shu. Nonlinearly Stable Compact Schemes for Shock Calculations. SIAM J. Numer. Anal., 31:607–627, 1994.
\par\hangindent=2em\hangafter=1 125. B. Cockburn and C. W. Shu. TVB Runge–Kutta Local Projection Discon- tinuous Galerkin Finite Element Method for Conservation Laws II: General Framework. Mathematics of Computation, 52:411–435, 1989.
\par\hangindent=2em\hangafter=1 126. B. Cockburn and C. W. Shu. The Runge–Kutta Local Projection P1– Discontinuous Galerkin Finite Element Method for Scalar Conservation Laws. Mathematical Modelling and Numerical Analysis, 25:337–361, 1991.
\par\hangindent=2em\hangafter=1 127. B. Cockburn and C. W. Shu. The Local Discontinuous Galerkin Method for Time–Dependent Convection Diffusion Systems. SIAM Journal on Numerical Analysis, 35:2440–2463, 1998.
\par\hangindent=2em\hangafter=1 128. B. Cockburn and C. W. Shu. The Runge–Kutta Discontinuous Galerkin Method for Conservation Laws V: Multidimensional Systems. J. Comput. Phys., 141:199–224, 1998.
\par\hangindent=2em\hangafter=1 129. B. Cockburn and C. W. Shu. Runge–Kutta Discontinuous Galerkin Methods for Convection–Dominated Problems. Journal of Scientific Computing, 16:173– 261, 2001.
\par\hangindent=2em\hangafter=1 130. E. Coddington and N. Levinson. Theory of Ordinary Differential Equations. McGraw–Hill, New York, 1955.
\par\hangindent=2em\hangafter=1 131. P. Colella. An Analysis of the Effect of Operator Splitting and of the Sampling Procedure on the Accuracy of Glimm’s Method. PhD thesis, Department of Mathematics, University of California, USA, 1978.
\par\hangindent=2em\hangafter=1 132. P. Colella. Glimm’s Method for Gas Dynamics. SIAM J. Sci. Stat. Comput., 3(1):76–110, 1982.
\par\hangindent=2em\hangafter=1 133. P. Colella. A Direct Eulerian MUSCL Scheme for Gas Dynamics. SIAM J. Sci. Stat. Comput., 6:104–117, 1985.
\par\hangindent=2em\hangafter=1 134. P. Colella. Multidimensional Upwind Methods for Hyperbolic Conservation Laws. J. Comput. Phys., 87:171–200, 1990.
\par\hangindent=2em\hangafter=1 135. P. Colella and H. H. Glaz. Efficient Solution Algorithms for the Riemann Problem for Real Gases. J. Comput. Phys., 59:264–289, 1985.
\par\hangindent=2em\hangafter=1 136. P. Colella, A. Majda, and V. Roytburd. Theoretical and Numerical Structure for Reacting Shock Waves. SIAM J. Sci. Stat. Comput., 7:1059–1080, 1986.
\par\hangindent=2em\hangafter=1 137. P. Colella and P. R. Woodward. The Piecewise Parabolic Method (PPM) for Gas Dynamical Simulation. J. Comput. Phys., 54:174–201, 1984.
\par\hangindent=2em\hangafter=1 138. P. Concus and W. Proskurowski. Numerical Solution of a Nonlinear Hyperbolic Equation by the Random Choice Method.J. Comput. Phys., 30:153–166, 1979.
\par\hangindent=2em\hangafter=1 139. S. D. Conte and C. de Boor. Elementary Numerical Analyis. McGraw–Hill Kogakusha Ltd., 1980.
\par\hangindent=2em\hangafter=1 140. J. M. Corber´an and M. Ll. Gasc´ on. TVD Schemes for the Calculation of Flow in Pipes of Variable Cross–Section. Mathl. Comput. Modelling, 21:85–92, 1995.
\par\hangindent=2em\hangafter=1 141. J. Corner. Theory of the Interior Ballistics of Guns. Wiley, 1950.
\par\hangindent=2em\hangafter=1 142. J. Cortes, A. Debussche, and I. Toumi. A Density Perturbation Method to Study the Eigenstructure of Two–Phase Flow Equation Systems. J. Comput. Phys., 147:1–22, 1998.
\par\hangindent=2em\hangafter=1 143. R. Courant and K. O. Friedrichs.Supersonic Flow and Shock Waves. Springer- Verlag, 1985.
% [原PDF p709]
\par\hangindent=2em\hangafter=1 144. R. Courant, E. Isaacson, and M. Rees. On the Solution of Nonlinear Hyperbolic Differential Equations by Finite Differences.Comm. Pure. Appl. Math., 5:243– 255, 1952.
\par\hangindent=2em\hangafter=1 145. M. G. Crandall and A. Majda. Monotone Difference Approximations for Scalar Conservation Laws. Mathematics of Computations, 34:1–34, 1980.
\par\hangindent=2em\hangafter=1 146. W. Dai and P. R. Woodward. An Approximate Riemann Solver for Ideal Magnetohydrodynamics. J. Comput. Phys., 111:354–372, 1994.
\par\hangindent=2em\hangafter=1 147. W. Dai and P. R. Woodward. Extension of the Piecewise Parabolic Method to Multidimensional Ideal Magnetohydrodynamics. J. Comput. Phys., 115:485–, 1994.
\par\hangindent=2em\hangafter=1 148. G. Dal Maso, P. G. LeFloch, and F. Murat. Definition and Weak Stability of Nonconservative Products. J. Math. Pure Appl., 74:483–548, 1995.
\par\hangindent=2em\hangafter=1 149. S. F. Davis. TVD Finite Difference Schemes and Artificial Viscosity. Technical Report ICASE 84–20, NASA Langley Research Center, Hampton, USA, 1984.
\par\hangindent=2em\hangafter=1 150. S. F. Davis. Simplified Second–Order Godunov–Type Methods. SIAM J. Sci. Stat. Comput., 9:445–473, 1988.
\par\hangindent=2em\hangafter=1 151. A. S. Dawes. Natural Coordinates and High Speeed Flows. A Numerical Method for Reactive Gases. PhD thesis, College of Aeronautics, Cranfield Institute of Technology, UK, 1992.
\par\hangindent=2em\hangafter=1 152. H. Deconinck, K. G. Powell, P. L. Roe, and R. Struijs. Multi–dimensional Schemes for Scalar Advection. AIAA, 91(1532), 1991.
\par\hangindent=2em\hangafter=1 153. H. Deconinck, P. L. Roe, and R. Struijs. Multi–dimensional Generalisation of Roe’s Flux Difference Splitter for the Euler Equations.Computers and Fluids, 22:215–222, 1993.
\par\hangindent=2em\hangafter=1 154. A. Dedner, F. Kemm, D. Kr¨ oner, C. D. Munz, T. Schnitzer, and M. Wesenberg. Hyperbolic Divergence Cleaning for the MHD Equations. J. Comput. Phys., 175:645–673, 2002.
\par\hangindent=2em\hangafter=1 155. L. Del Zanna, N. Bucciantini, and P. Londrillo. An Efficient Shock–Capturing Central Scheme for Multi–dimensional Relativistic Flows II. Magnetohydrody- namics. Astronomy and Astrophysics, 400:397–413, 2003.
\par\hangindent=2em\hangafter=1 156. V. Deledicque and M. V. Papalexandris. An Exact Riemann Solver for Com- pressible Two–Phase Flow Models Containing Nononservative Products. J. Comput. Phys., 222:217–245, 2007.
\par\hangindent=2em\hangafter=1 157. T. DeNeff and G. Moretti. Shock Fitting for Everyone.Computers and Fluids, 8:327–334, 1980.
\par\hangindent=2em\hangafter=1 158. J. M. Dewey. Mach Reflection Research–Paradox and Progress. InShock Waves (Proc. 18th Int. Symp. on Shock Waves, Sendai, Japan, 1991. Takayama, Ed- itor), Vol. I, pages 113–120. Springer–Verlag, 1992.
\par\hangindent=2em\hangafter=1 159. J. Dobes and H. Deconinck. Second Order Blended Multidimensional Upwind Residual Distribution Scheme for Steady and Unsteady Computations.Journal of Computational and Applied Mathematics, 215:378–389, 2008.
\par\hangindent=2em\hangafter=1 160. D. A. Drew and R. Gulliver. Particulate Flows: Processing and Rheology . Springer, 1997.
\par\hangindent=2em\hangafter=1 161. D. A. Drew, J. E. Marsden, and L. Sirovich.Theory of Multicomponent Fluids. Springer, 1998.
\par\hangindent=2em\hangafter=1 162. D. Drikakis. A Characteristic–Based Method for Incompressible Flow. Int. J. Numer. Meth. Fluids, 19:667–685, 1994.
\par\hangindent=2em\hangafter=1 163. D. Drikakis. A Parallel Multiblock Characteristic–Based Method for Three– Dimensional Incompressible Flows. Advances in Engineering Software, 26:111– 119, 1996.
% [原PDF p710]
\par\hangindent=2em\hangafter=1 164. D. Drikakis, O. P. Iliev, and D. P. Vassileva. A Nonlinear Multigrid Method for the Three–Dimensional Incompressible Navier–Stokes Equations. J. Comput. Phys., 146:301–321, 1998.
\par\hangindent=2em\hangafter=1 165. D. Drikakis and W. Rider. High–Resolution Methods for Incompressible and Low–Speed Flows. Springer, 2005.
\par\hangindent=2em\hangafter=1 166. D. Drikakis and S. Tsangaris. On the Solution of the Compressible Navier– Stokes Equations Using Improved Flux Vector Splitting Methods.Appl. Math. Modelling, 17:282–297, 1993.
\par\hangindent=2em\hangafter=1 167. F. Dubois and G. Mehlman. A Non–Parameterized Entropy Fix for Roe’s Method. AIAA Journal, 31 (1):199–200, 1993.
\par\hangindent=2em\hangafter=1 168. J. K. Dukowicz. A General, Non–Iterative Riemann Solver for Godunov’s Method. J. Comput. Phys., 61:119–137, 1985.
\par\hangindent=2em\hangafter=1 169. M. Dumbser. Arbitrary High Order Schemes for the Solution of Hyperbolic Conservation Laws in Complex Domains. PhD thesis, PhD thesis, Institut f¨ur Aero- un Gasdynamik, Universit¨ at Stuttgart, Germany, 2005.
\par\hangindent=2em\hangafter=1 170. M. Dumbser. Building Blocks for Arbitrary High Order Discontinuous Galerkin Schemes. Journal of Scientific Computing , 27:215–230, 2006.
\par\hangindent=2em\hangafter=1 171. M. Dumbser, D. Balsara, E. F. Toro, and C. D. Munz. A Unified Framework for the Construction of One–Step Finite–Volume and Discontinuous Galerkin Schemes. J. Comput. Phys., 227:8209–8253, 2008.
\par\hangindent=2em\hangafter=1 172. M. Dumbser, M. J. Castro, C. Par´ es, and E. F. Toro. ADER Schemes on Unstructured Meshes for Nonconservative Hyperbolic Systems: Applications to Geophysical Flows. Computers and Fluids, 2008 (submitted).
\par\hangindent=2em\hangafter=1 173. M. Dumbser, C. Enaux, and E. F. Toro. Finite Volume Schemes of Very High Order of Accuracy for Stiff Hyperbolic Balance Laws.J. Comput. Phys., 227(8):3971–4001, 2008.
\par\hangindent=2em\hangafter=1 174. M. Dumbser, A. Hidalgo, M. J. Castro, C. Par´ es, and E. F. Toro. FORCE Schemes on Unstructured Meshes II: Nonconservative Hyperbolic Systems. Isaac Newton Institute for Mathematical Sciences, University of Cambridge, UK. Preprint NI09053–NPA, December 2008. Submitted to J. Comput. Phys., 2008.
\par\hangindent=2em\hangafter=1 175. M. Dumbser and M. K¨ aser. Arbitrary High Order Non-Oscillatory Finite Vol- ume Schemes on Unstructured Meshes for Linear Hyperbolic Systems.J. Com- put. Phys., 221(2):693–723, 2007.
\par\hangindent=2em\hangafter=1 176. M. Dumbser, M. K¨aser, V. A. Titarev, and E. F. Toro. Quadrature-Free Non- Oscillatory Finite Volume Schemes on Unstructured Meshes for Nonlinear Hy- perbolic Systems. J. Comput. Phys., 226(8):204–243, 2007.
\par\hangindent=2em\hangafter=1 177. M. Dumbser, M. K¨aser, and E. F. Toro. An Arbitrary High Order Discontinu- ous Galerkin Method for Elastic Waves on Unstructured Meshes V: Local Time Stepping and p–Adaptivity. Geophysical Journal International, 171:695–717, 2007.
\par\hangindent=2em\hangafter=1 178. M. Dumbser and C. D. Munz. ADER Discontinuous Galerkin Schemes for Aeroacoustics. Comptes Rendus M´ecanique, 333:683–687, 2005.
\par\hangindent=2em\hangafter=1 179. P. Dutt. A Riemann Solver Based on a Global Existence Proof for the Riemann Problem. Technical Report ICASE 86–3, NASA Langley Research Center, USA, 1986.
\par\hangindent=2em\hangafter=1 180. D. H. Edwards, G. O. Thomas, and T. L. Williams. Initiation of Detonation by Steady Planar Incident Shock Waves. Combustion and Flame, 43:187–198, 1981.
% [原PDF p711]
\par\hangindent=2em\hangafter=1 181. B. Einfeldt. On Godunov–Type Methods for Gas Dynamics. SIAM J. Numer. Anal., 25(2):294–318, 1988.
\par\hangindent=2em\hangafter=1 182. B. Einfeldt, C. D. Munz, P. L. Roe, and B. Sj¨ ogreen. On Godunov–Type Methods near Low Densities. J. Comput. Phys., 92:273–295, 1991.
\par\hangindent=2em\hangafter=1 183. P. Embid and M. Baer. Mathematical Analysis of a Two–Phase Continuum Mixture Theory. Continuum Mech. Thermodyn., 4:279–312, 1992.
\par\hangindent=2em\hangafter=1 184. B. Engquist and B. Gustafsson. Proc. Third International Conference on Hy- perbolic Problems, Vol. I and II. Chartwell–Bratt, 1991.
\par\hangindent=2em\hangafter=1 185. B. Engquist and S. Osher. One Sided Difference Approximations for Nonlinear Conservation Laws. Math. Comp., 36(154):321–351, 1981.
\par\hangindent=2em\hangafter=1 186. L. C. Evans. Partial Differential Equations. American Mathematical Society, 2002.
\par\hangindent=2em\hangafter=1 187. S. A. E. G. Falle and S. S. Komissarov. An Upwind Numerical Scheme for Relativistic Hydrodynamics with a General Equation of State.Monthly Notices of the Royal Astronomical Society, 278:586–602, 1996.
\par\hangindent=2em\hangafter=1 188. S. A. E. G. Falle and S. S. Komissarov. A Multidimensional Upwind Scheme for Magnetohydrodynamics. Monthly Notices of the Royal Astronomical Society, 297:265–277, 1997.
\par\hangindent=2em\hangafter=1 189. M. Feistauer, J. Felcman, and I. Straskraba. Mathematical and Computational Methods for Compressible Flow. Oxford University Press, 2003.
\par\hangindent=2em\hangafter=1 190. M. Fey, R. Jeltsch, and S M¨ uller. The Influence of a Source Term, an Ex- ample: Chemically Reacting Hypersonic Flow. Technical Report 92–06, SAM, Eidgen¨ossische Technische Hochschule (ETH), Z¨ urich, Switzerland, 1992.
\par\hangindent=2em\hangafter=1 191. W. Fickett and W. Davis. Detonation. University of California Press, 1979.
\par\hangindent=2em\hangafter=1 192. C. A. J. Fletcher. Computational Techniques for Fluid Dynamics, Vols. I and II. Springer–Verlag, 1988.
\par\hangindent=2em\hangafter=1 193. L. Fraccarollo and E. F. Toro. A Shock-Capturing Method for Two Dimensional Dam–Break Problems. Proceedings of the Fifth International Symposium in Computational Fluid Dynamics, Sendai, Japan, 1993.
\par\hangindent=2em\hangafter=1 194. L. Fraccarollo and E. F. Toro. Experimental and Numerical Assessment of the Shallow Water Model for Two–Dimensional Dam–Break Type problems. Journal of Hydraulic Research, 33:843–864, 1995.
\par\hangindent=2em\hangafter=1 195. J. E. Fromm. A Method for Reducing Dispersion in Convective Difference Schemes. J. Comput. Phys., 3:176–189, 1968.
\par\hangindent=2em\hangafter=1 196. M. Gallardo, C. Par´es, and M. J. Castro. On a Well–Balanced High–Order Finite Volume Scheme for Shallow Water Equations with Topography and Dry Areas. J. Comput. Phys., 227:574–601, 2007.
\par\hangindent=2em\hangafter=1 197. P. Garc´ıa-Navarro, M. E. Hubbard, and A. Priestley. Genuinely Multidimen- sional Upwinding for the 2D Shallow Water Equations. J. Comput. Phys. , 121:79–93, 1995.
\par\hangindent=2em\hangafter=1 198. P. H. Gaskell and A. K. C. Lau. Curvature Compensated Convective Transport: SMART, a New Boundedness Preserving Transport Algorithm.Int. J. Numer. Meth. Fluids, 8:617–641, 1988.
\par\hangindent=2em\hangafter=1 199. S. L. Gavrilyuk, N. Favrie, and R. Saurel. Modelling Wave Dynamics of Com- pressible Elastic Materials. J. Comput. Phys., 227(5):2941–2969, 2007.
\par\hangindent=2em\hangafter=1 200. C. W. Gear. Numerical Initial–Value Problems in Ordinary Differential Equa- tions. Prentice–Hall, Englewood Cliffs, NJ, 1971.
\par\hangindent=2em\hangafter=1 201. B. E. Gelfand, S. M. Frolov, and M. A. Nettleton. Gaseous Detonations–A Selective Review. Prog. Energy Comb. Sci., 17:327–371, 1991.
% [原PDF p712]
\par\hangindent=2em\hangafter=1 202. B. Giacomazzo and L. Rezzolla. The Exact Solution of the Riemann Problem in Relativistic Magnetohydrodynamics. Journal of Fluid Mechanics, 562:223– 259, 2006.
\par\hangindent=2em\hangafter=1 203. D. Gidaspow. Multiphase Flow and Fluidization. Academic Press, 1994.
\par\hangindent=2em\hangafter=1 204. R. B. Gilbert and R. A. Strehlow. Theory of Detonation Initiation Behind Reflected Shock Waves. AIAA Journal, 4:1777–1783, 1966.
\par\hangindent=2em\hangafter=1 205. M. B. Giles. Non-Reflecting Boundary Conditions for Euler Equation Calcu- lations. AIAA J., 28:2050–2058, 1990.
\par\hangindent=2em\hangafter=1 206. L. Giraud and G. Manzini. Parallel Implementations of 2D Explicit Euler Solvers. J. Comput. Phys., 123:111–118, 1996.
\par\hangindent=2em\hangafter=1 207. P. Glaister. Difference Schemes for the Shallow Water Equations. Technical Report 9/87, Department of Mathematics, University of Reading, England, 1987.
\par\hangindent=2em\hangafter=1 208. P. Glaister. An Approximate Linearised Riemann Solver for the Euler Equa- tions of Gas Dynamics. J. Comput. Phys., 74:382–408, 1988.
\par\hangindent=2em\hangafter=1 209. P. Glaister. An Approximate Linearised Riemann Solver for the Three– Dimensional Euler Equations for Real Gases Using Operator Splitting. J. Comput. Phys., 77:361–383, 1988.
\par\hangindent=2em\hangafter=1 210. H. M. Glaz, P. Colella, I. I. Glass, and R. L. Deschambault. A Detailed Numer- ical, Graphical and Experimental Study of Oblique Shock Wave Reflections. Technical Report 285, Institute for Aerospace Science, University of Toronto (UTIAS), 1985.
\par\hangindent=2em\hangafter=1 211. H. M. Glaz, P. Colella, I. I. Glass, and R. L. Deschambault. A Numerical Study of Oblique Shock–Wave Reflections with Experimental Comparisons. Proc. Roy. Soc. London, A 398:117–140, 1985.
\par\hangindent=2em\hangafter=1 212. J. Glimm. Solution in the Large for Nonlinear Hyperbolic Systems of Equa- tions. Comm. Pure. Appl. Math., 18:697–715, 1965.
\par\hangindent=2em\hangafter=1 213. J. Glimm, M. J. Graham, J. W Grove, and B. J. Plohr.Proc. Fifth International Conference on Hyperbolic Problems. World Scientific, 1996.
\par\hangindent=2em\hangafter=1 214. J. Glimm, G. Marshall, and B. Plohr. A Generalized Riemann Problem for Quasi–One–Dimensional Gas Flows. Advances in Applied Mathematics, 5:1–30, 1984.
\par\hangindent=2em\hangafter=1 215. E. Godlewski and P. A. Raviart. Numerical Approximation of Hyperbolic Systems of Conservation Laws . Applied Mathematical Sciences, Vol. 118, Springer–Verlag, New York, 1996.
\par\hangindent=2em\hangafter=1 216. S. K. Godunov. A Finite Difference Method for the Computation of Discon- tinuous Solutions of the Equations of Fluid Dynamics. Mat. Sb., 47:357–393, 1959.
\par\hangindent=2em\hangafter=1 217. S. K. Godunov, A. V. Zabrodin, and G. P. Prokopov. –. J. Comp. Math. Phys. USSR 1 (1962), pages 1187–, 1962.
\par\hangindent=2em\hangafter=1 218. S. K. Godunov (Ed.). Numerical Solution of Multi–Dimensional Problems in Gas Dynamics. Nauka Press, Moscow, 1976.
\par\hangindent=2em\hangafter=1 219. J. B. Goodman and R. J. LeVeque. On the Accuracy of Stable Schemes for 2D Scalar Conservation Laws. Math. Comp., 45(21):15–21, 1985.
\par\hangindent=2em\hangafter=1 220. L. Gosse. A Well–Balanced Flux–Vector Splitting Scheme Designed for Hy- perbolic Systems of Conservation Laws with Source Terms. Computers and Mathematics with Applications, 39:135–159, 2000.
\par\hangindent=2em\hangafter=1 221. J. J. Gottlieb. Staggered and Non–Staggered Grids with Variable Node Spac- ing and Local Timestepping for Random Choice Method. J. Comput. Phys., 78:160–177, 1988.
% [原PDF p713]
\par\hangindent=2em\hangafter=1 222. J. J. Gottlieb and C. P. T. Groth. Assessment of Riemann Solvers for Unsteady One–Dimensional Inviscid Flows of Perfect Gases. J. Comput. Phys., 78:437– 458, 1988.
\par\hangindent=2em\hangafter=1 223. S. Gottlieb and C. W. Shu. Total Variation Diminishing Runge–Kutta Schemes. Technical Report ICASE 96–50, NASA Langley Research Center, Hampton, USA, 1996.
\par\hangindent=2em\hangafter=1 224. G. Grassner, F. L¨orcher, and C. D. Munz. A Contribution to the Construction of Diffusion Fluxes for Finite Volume and Discontinuous Galerkin Schemes.J. Comput. Phys., 224:1049–1063, 2007.
\par\hangindent=2em\hangafter=1 225. J. M. Greenberg and A. Y. LeRoux. A Well–Balanced Scheme for the Numer- ical Processing of Source Terms in Hyperbolic Equations. SIAM J. Numerical Analysis, 33:1–16, 1996.
\par\hangindent=2em\hangafter=1 226. J. M. Greenberg, A. Y. LeRoux, R. Baraille, and A Noussair. Analysis and Approximation of Conservation Laws with Source Terms.SIAM J. Numerical Analysis, 34:1980–2007, 1997.
\par\hangindent=2em\hangafter=1 227. D. F. Griffiths, A. M. Stuart, and H. C. Yee. Numerical Wave Propagation in an Advection Equation with a Nonlinear Source Term. SIAM J. Numer. Anal., 29:1244–1260, 1992.
\par\hangindent=2em\hangafter=1 228. H. Gr¨onig (Editor). Shock Tubes and Waves (Proc. 16th Int. Symp. on Shock Tubes and Waves, Aachen Germany 1987) . VCH Verlagsgessellschaft mbH, 1988.
\par\hangindent=2em\hangafter=1 229. V. Guinot. Godunov–Type Schemes. An Introduction for Engineers. Elsevier, 2003.
\par\hangindent=2em\hangafter=1 230. K. F. Gurski. An HLLC–Type Approximate Riemann Solver for Ideal Magne- tohydrodynamics. SIAM J. Sci. Comput., 25(6):2165–2187, 2004.
\par\hangindent=2em\hangafter=1 231. B. Haasdonk, D. Kr¨ oner, and C. Rhode. Convergence of a Staggered Lax– Friedrichs Scheme on Unstructured Grids in 2D. Numerische Mathematik, 88:459–484, 2001.
\par\hangindent=2em\hangafter=1 232. R. Haberman. Mathematical Models. SIAM, 1998.
\par\hangindent=2em\hangafter=1 233. J. M. Hammersley and D. C. Handscombe. Monte Carlo Methods. Chapman and Hall, 1964.
\par\hangindent=2em\hangafter=1 234. D. H¨anel. On the Numerical Solution of the Navier–Stokes Equations for Com- pressible Flow. VKI Lecture Series 1989–04 on Computational Fluid Dynamics, 1989.
\par\hangindent=2em\hangafter=1 235. D. H¨anel and S. D. Sharma. Simulation of Hydrodynamical Free–Surface Flow. In Proc. Third ECCOMAS Computational Fluid Dynamics Conference, 9–13 September 1996, Paris, France, pages 1019–1024. Wiley, 1996.
\par\hangindent=2em\hangafter=1 236. E. Harabetian. A Convergent Series Expansion for Hyperbolic Systems of Conservation Laws. Trans Amer. Math. Soc., 294:383–424,, 1986.
\par\hangindent=2em\hangafter=1 237. R. Harris, Z. J. Wang, and Y. Liu. Efficient Quadrature–Free High–Order Spectral Volume Method on Unstructured Grids: Theory and 2D. J. Comput. Phys., 227(3):1620–1642, 2008.
\par\hangindent=2em\hangafter=1 238. A. Harten. High Resolution Schemes for Hyperbolic Conservation Laws. J. Comput. Phys., 49:357–393, 1983.
\par\hangindent=2em\hangafter=1 239. A. Harten. On a Class of High Resolution Total Variation Stable Finite Dif- ference Schemes. SIAM J. Numer. Anal., 21(1):1–23, 1984.
\par\hangindent=2em\hangafter=1 240. A. Harten. Preliminary Results on the Extension of ENO Schemes to Two– Dimensional Problems. In International Conference on Non–Linear Hyperbolic Problems, pages 23–51. Springer–Verlag, Berlin, 1987.
% [原PDF p714]
\par\hangindent=2em\hangafter=1 241. A. Harten. ENO Schemes with Subcell Resolution. J. Comput. Phys., 83:148– 184, 1989.
\par\hangindent=2em\hangafter=1 242. A. Harten, B. Engquist, S. Osher, and S. R. Chakravarthy. Uniformly High Order Accuracy Essentially Non–oscillatory Schemes III. J. Comput. Phys., 71:231–303, 1987.
\par\hangindent=2em\hangafter=1 243. A. Harten and J. M. Hyman. Self Adjusting Grid Methods for One– Dimensional Hyperbolic Conservation Laws. J. Comput. Phys. , 50:235–269, 1983.
\par\hangindent=2em\hangafter=1 244. A. Harten, P. D. Lax, and B. van Leer. On Upstream Differencing and Godunov–Type Schemes for Hyperbolic Conservation Laws. SIAM Review, 25(1):35–61, 1983.
\par\hangindent=2em\hangafter=1 245. A. Harten and S. Osher. Uniformly High–Order Accurate Nonoscillatory Schemes I. SIAM J. Numer. Anal., 24(2):279–309, 1987.
\par\hangindent=2em\hangafter=1 246. W. Heilig and H. Reichenbach. High Speed Interferometric Study of Unstation- ary Shok Wave Processes. Technical report, Fraunhofer Institut f¨ur Kurzzeit- dynamik, Ernst–Mach–Institut, Freiburg im Breisgau, Germany, 1982.
\par\hangindent=2em\hangafter=1 247. P. W. Hemker and S. P. Spekreijse. Multiple Grid and Osher’s Scheme for the Efficient Solution of the Steady Euler Equations. Applied Numerical Mathe- matics, 2:475–493, 1986.
\par\hangindent=2em\hangafter=1 248. J. S. Hesthaven and T. Warburton. Nodal Discontinuous Galerkin Methods: Algorithms, Analysis and Applications. Springer, 2008.
\par\hangindent=2em\hangafter=1 249. F. B. Hildebrand. Introduction to Numerical Analysis. Tata McGraw–Hill Publishing Co. Limited, New Delhi, 1974.
\par\hangindent=2em\hangafter=1 250. R. Hillier. Numerical Modelling of Shock Wave Diffraction. In Shock Waves @ Marseille, Vol. 4 (Brun and Dumitrescu, Editors) , pages 17–26. Springer– Verlag, 1995.
\par\hangindent=2em\hangafter=1 251. C. Hirsch. Numerical Computation of Internal and External Flows, Vol. I: Fundamentals of Numerical Discretization. Wiley, 1988.
\par\hangindent=2em\hangafter=1 252. C. Hirsch. Numerical Computation of Internal and External Flows, Vol. II: Computational Methods for Inviscid and Viscous Flows. Wiley, 1990.
\par\hangindent=2em\hangafter=1 253. K. A. Hoffmann. Computational Fluid Dynamics for Engineers. Engineering Education Systems, Austin, Texas, USA, 1989.
\par\hangindent=2em\hangafter=1 254. M. Holt. Numerical Methods in Fluid Dynamics. Springer–Verlag, 1984.
\par\hangindent=2em\hangafter=1 255. V. Honkkila and P. Janhunen. HLLC Solver for Ideal Relativistic MHD. J. Comput. Phys., 203(1):344–357, 2007.
\par\hangindent=2em\hangafter=1 256. H. Honma and I. I. Glass. Random Choice Solutions for Weak Spherical Shock– Wave Transitions of N–Waves in Air with Vibrational Excitation. Technical Report 253, UTIAS, University of Toronto, Canada, 1983.
\par\hangindent=2em\hangafter=1 257. H. Honma, M. Wada, and K. Inomata. Random Choice Solutions for Two– Dimensional and Axi–Symetric Supersonic Flows. Technical Report S. P. Num- ber 2, The Institute of Space and Astronautical Science, 1984.
\par\hangindent=2em\hangafter=1 258. L. H¨ ormander. Lectures on Nonlinear Hyperbolic Differential Equations, Math´ematiques et Applications 26. Springer–Verlag, 1997.
\par\hangindent=2em\hangafter=1 259. T. Y. Hou and LeFloch P. Why Non–Conservative Schemes Converge to the Wrong Solutions: Error Analysis. M a t h .o fC o m p u t ., 62:497–530, 1994.
\par\hangindent=2em\hangafter=1 260. E. L. Ince and I. N. Sneddon. The Solution of Ordinary Differential Equations. Longman Scientific and Technical, 1987.
\par\hangindent=2em\hangafter=1 261. M. Ishii and T. Hibiki. Thermo–fluid Dynamics of Two–Phase Flow . Springer, 2006.
% [原PDF p715]
\par\hangindent=2em\hangafter=1 262. M. J. Ivings, D. M. Causon, and E. F. Toro. On Hybrid High–Resolution Upwind Methods for Multicomponent Flows. ZAMM, 77, Issue 9:645–668, 1997.
\par\hangindent=2em\hangafter=1 263. M. J. Ivings, D. M. Causon, and E. F. Toro. On Riemann Solvers for Com- pressible Liquids. Technical Report 97–4, Department of Mathematics and Physics, Manchester Metropolitan University, UK, 1997.
\par\hangindent=2em\hangafter=1 264. M. J. Ivings, D. M. Causon, and E. F. Toro. On Riemann Solvers for Com- pressible Liquids. Int. J. Numer. Meth. Fluids, 28:395–418, 1998.
\par\hangindent=2em\hangafter=1 265. R. Jackson. The Dynamics of Fluidized Particles. Cambridge University Press, 2000.
\par\hangindent=2em\hangafter=1 266. G. Jagadeesh, E. Arunan, and K. P. J Reddy (Editors).Shock Waves: Proceed- ings of the 25th International Symposium on Shock Waves–ISSW25, Bangalore, India. Universities Press, 2006.
\par\hangindent=2em\hangafter=1 267. A. Jameson and P. D. Lax. Conditions for the Construction of Multi–Point Total Variation Diminishing Difference Schemes. Technical Report MAE 1650, Department of Mechanical and Aerospace Engineering, University of Prince- ton, USA, 1984.
\par\hangindent=2em\hangafter=1 268. A. Jameson, W. Schmidt, and E. Turkel. Numerical Solution of the Euler Equations by Finite Volume Methods using Runge–Kutta Stepping Schemes. Technical Report 81–1259, AIAA, 1981.
\par\hangindent=2em\hangafter=1 269. A. Jeffrey. Quasilinear Hyperbolic Systems and Waves. Pitman, 1976.
\par\hangindent=2em\hangafter=1 270. G. S. Jiang and C. W. Shu. Efficient Implementation of Weighted ENO Schemes. J. Comp. Phys., 126:202–228, 1996.
\par\hangindent=2em\hangafter=1 271. G. S. Jiang and E. Tadmor. Non–oscillatory Central Schemes for Multi– dimensional Hyperbolic Conservation Laws. SIAM Journal of Scientific Com- puting, 19 (6):1892–1917, 1998.
\par\hangindent=2em\hangafter=1 272. F. John. Partial Differential Equations. Springer–Verlag, 1982.
\par\hangindent=2em\hangafter=1 273. L. W. Johnson and R. D. Riess. Numerical Analysis. Addison–Wesley Pub- lishing Company, 1982.
\par\hangindent=2em\hangafter=1 274. P. Jorgenson and E. Turkel. Central–Difference TVD and TVB Schemes for Time–Dependent and Steady State Problems. Technical Report ICOMP–91– 27; AIAA–92–0053, NASA Lewis Research Center, Cleveland, Ohio, USA, 1992.
\par\hangindent=2em\hangafter=1 275. D. Kahaner, C. Moler, and S. Nash.Numerical Methods and Software.P r e n t i c e Hall, Englewood Cliffs, New Jersey, 1989.
\par\hangindent=2em\hangafter=1 276. K. Kailasanath, E. S. Oran, J. P. Boris, and T. R. Young. Determination of the Detonation Cell Size and the Role of Transverse Waves in Two–Dimensional Detonations. Combustion and Flame, 61:199–209, 1985.
\par\hangindent=2em\hangafter=1 277. S. Karni. Accelerated Convergence to Steady State by Gradual Far–Field Damping. AIAA J., 30:1220–1228, 1992.
\par\hangindent=2em\hangafter=1 278. S. Karni. Multicomponent Flow Calculations Using a Consistent Primitive Algorithm. J. Comput. Phys., 112(1):31–43, 1994.
\par\hangindent=2em\hangafter=1 279. S. Karni. Hybrid Multifluid Algorithms. Technical Report 95–001, Courant Mathematics and Computing Laboratory, 1995.
\par\hangindent=2em\hangafter=1 280. S. Karni and S. ˆCani´c. Computations of Slowly Moving Shocks. J. Comput. Phys., 136:132–139, 1997.
\par\hangindent=2em\hangafter=1 281. M. K¨aser. Adaptive Methods for the Numerical Simulation of Transport Pro- cesses. PhD thesis, Institute of Numerical Mathematics and Scientific Com- puting, University of Munich, Germany, 2003.
% [原PDF p716]
\par\hangindent=2em\hangafter=1 282. M. K¨aser. ADER Schemes for the Solution of Conservation Laws on Adap- tive Triangulations. Mathematical Methods and Modelling in Hydrocarbon Exploration and Production. Springer–Verlag (to appear), 2004.
\par\hangindent=2em\hangafter=1 283. M. K¨aser and A. Iske. Adaptive ADER Schemes for the Solution of Scalar Non–Linear Hyperbolic Problems. J. Comput. Phys., -:–, 2005.
\par\hangindent=2em\hangafter=1 284. K. S. Kim and J. Ballmann. Numerische Simulation Mechanischer Wellen in Geschichteten Elastischen K¨orpern. Z. Angew. Math. Mech., 70:204–206, 1990.
\par\hangindent=2em\hangafter=1 285. C. Klingenberg, W. Schmidt, and K. Waagan. Numerical Comparison of Rie- mann Solvers for Astrophysical Hydrodynamics. J. Comput. Phys., 227(1):12– 35, 2007.
\par\hangindent=2em\hangafter=1 286. L. Knut-Andreas and S. Noelle. An Improved Quadrature Rule for the Flux Computation in Staggered Central Difference Schemes in Multidimensions. J. Scientific Computing , 18(1):69–81, 2003.
\par\hangindent=2em\hangafter=1 287. V. P. Kolgan. Application of the Principle of Minimum Derivatives to the Construction of Difference Schemes for Computing Discontinuous Solutions of Gas Dynamics (in Russian). Uch. Zap. TsaGI, Russia, 3(6):68–77, 1972.
\par\hangindent=2em\hangafter=1 288. V. P. Kolgan. Finite Difference Schemes for the Computation of Three– dimensional Solutions of Gas Dynamics and Calculation of the Flow over a Body under an Angle of Attack (in Russian). Technical Report Vol. 6, No 2, pp 1–6, Uchenye Zapiski) TsaGI (Scientific Notes of the Central Institute of Aerodynamics, Russia), 1975.
\par\hangindent=2em\hangafter=1 289. V. P. Kolgan. Finite Difference Schemes for the Computation of Two– dimensional Discontinuous Solutions of Non–stationary Gas Dynamics (in Rus- sian). Technical Report Vol. 6, No 1, pp 9–14, Uchenye Zapiski) TsaGI (Sci- entific Notes of the Central Institute of Aerodynamics, Russia), 1975.
\par\hangindent=2em\hangafter=1 290. B. Koren and S. P. Spekreijse. Multigrid and Defect Correction for the Ef- ficient Solution of the Steady Euler Equations. In Notes on Numerical Fluid Mechanics, Vol. 17, pages 87–100. Vieweg, 1987.
\par\hangindent=2em\hangafter=1 291. D. Kr¨ oner.Numerical Schemes for Conservation Laws. Wiley Teubner, 1997.
\par\hangindent=2em\hangafter=1 292. A. G. Kulikovskii, N. V. Pogorelov, and A. Y. Semenov.Mathematical Aspects of Numerical Solution of Hyperbolic Systems. Chapman and Hall, 2001.
\par\hangindent=2em\hangafter=1 293. A. Kurganov and G. Petrova. A Third–Order Semi-Discrete Genuinely Mul- tidimensional Central Scheme for Hyperbolic Conservation Laws and Related Problems. Numerische Mathematik, 88(3):683–729, 2001.
\par\hangindent=2em\hangafter=1 294. A. Lafon and H. C. Yee. Dynamical Approach Study of Spurious Steady– State Numerical Solutions of Nonlinear Differential Equations, Part III. The Effects of Nonlinear Source Terms in Reaction–Convection Equations. Intern. J. Computational Fluid Dynamics, 6:1–36, 1996.
\par\hangindent=2em\hangafter=1 295. A. Lafon and H. C. Yee. Dynamical Approach Study of Spurious Steady–State Numerical Solutions of Nonlinear Differential Equations, Part IV. Stability vs. Methods of Discretizing Nonlinear Source Terms in Reaction–Convection Equations. Intern. J. Computational Fluid Dynamics, 6:89–123, 1996.
\par\hangindent=2em\hangafter=1 296. J. D. Lambert. Computational Methods in Ordinary Differential Equations . John Wiley and Sons, 1973.
\par\hangindent=2em\hangafter=1 297. L. D. Landau and E. M. Lifshitz. Fluid Mechanics. Pergamon Press, 1982.
\par\hangindent=2em\hangafter=1 298. C. B. Laney. Computational Gasdynamics. Cambridge University Press, 1998.
\par\hangindent=2em\hangafter=1 299. P. D. Lax. Weak Solutions of Nonlinear Hyperbolic Equations and Their Nu- merical Computation. Comm. Pure. Appl. Math., VII:159–193, 1954.
\par\hangindent=2em\hangafter=1 300. P. D. Lax. The Formation and Decay of Shock Waves.American Mathematical Monthly, 79:227–241, 1972.
% [原PDF p717]
\par\hangindent=2em\hangafter=1 301. P. D. Lax. Hyperbolic Systems of Conservation Laws and the Mathematical Theory of Shock Waves . Society for Industrial and Applied Mathematics, Philadelphia, 1990.
\par\hangindent=2em\hangafter=1 302. P. D. Lax and B. Wendroff. Systems of Conservation Laws.Comm. Pure Appl. Math., 13:217–237, 1960.
\par\hangindent=2em\hangafter=1 303. P. D. Lax and B. Wendroff. Difference Schemes for Hyperbolic Equations with High Order Accuracy. Comm. Pure Appl. Math., XVII:381–393, 1964.
\par\hangindent=2em\hangafter=1 304. P. Le Floch and P. A. Raviart. An Asymptotic Expansion for the Solution of the Generalized Riemann Problem. Part 1: General Theory. Ann. Inst. Henri Poincar´e. Analyse non Line´are, 5(2):179–207, 1988.
\par\hangindent=2em\hangafter=1 305. P. Le Floch and L. Tatsien. A Global Asymptotic Expansion for the Solution of the Generalized Riemann Problem.Ann. Inst. Henri Poincar´ e. Analyse non Line´are, 3:321–340, 1991.
\par\hangindent=2em\hangafter=1 306. L. Lehner. Numerical Relativity: a Review. Class. Quantum Grav., 18:R25– R86, 2001.
\par\hangindent=2em\hangafter=1 307. B. P. Leonard. Simple High–Accuracy Resolution Program for Convective Modelling of Discontinuities. Int. J. Numer. Methods in Fluids, 8:1291–1318, 1988.
\par\hangindent=2em\hangafter=1 308. R. J. LeVeque. Numerical Methods for Conservation Laws. Birkh¨ auser Verlag, 1992.
\par\hangindent=2em\hangafter=1 309. R. J. LeVeque. Simplified Multidimensional Flux Limiter Methods. In M. J. Baines and K. W. Morton, editors, Numerical Methods in Fluid Dynamics 4: Proceedings of the 1992 International Conference on Numerical Methods in Fluids, pages 173–189, Reading, 1993.
\par\hangindent=2em\hangafter=1 310. R. J. LeVeque. Balancing Source Terms and Flux Gradients in High–Resolution Godunov Methods. J. Comput. Phys., 146:346–365, 1998.
\par\hangindent=2em\hangafter=1 311. R. J. LeVeque. Finite Volume Methods for Hyperbolic Problems. Cambridge University Press, 2002.
\par\hangindent=2em\hangafter=1 312. R. J. LeVeque, D. Mihalas, E. A. Dorfi, and E. M¨ uller.Computational Methods for Astrophysical Flow. Springer–Verlag, 1998.
\par\hangindent=2em\hangafter=1 313. R. J. LeVeque and K. M. Shyue. Two–Dimensional Front Tracking Based on High Resolution Wave Propagation Methods. J. Comput. Phys., 123:354–368, 1996.
\par\hangindent=2em\hangafter=1 314. R. J. LeVeque and H. C. Yee. A Study of Numerical Methods for Hyperbolic Conservation Laws with Stiff Source Terms. J. Comput. Phys. , 86:187–210, 1990.
\par\hangindent=2em\hangafter=1 315. D. W. Levy, G. Puppo, and G. Russo. Central WENO Schemes for Hyperbolic Systems of Conservation Laws. Mathematical Models and Numerical Analysis, 33:547–571, 1999.
\par\hangindent=2em\hangafter=1 316. D. W. Levy, G. Puppo, and G. Russo. A Fourth Order Central WENO Scheme for Multidimensional Systems of Hyperbolic Conservation Laws.SIAM J. Sci- entific Computing , 24:480–506, 2002.
\par\hangindent=2em\hangafter=1 317. K. M. Li and M. Holt. Numerical Solutions to Water Waves Generated by Shallow Underwater Explosions. Phys. Fluids, 24:816–824, 1981.
\par\hangindent=2em\hangafter=1 318. S. Li. An HLLC Riemann Solver for Magneto–Hydrodynamics. J. Comput. Phys., 203(1):344–357, 2005.
\par\hangindent=2em\hangafter=1 319. NAG Library. NAG Library Routines, Mark 18, Routines D03PWF and D03PXF, 1996.
% [原PDF p718]
\par\hangindent=2em\hangafter=1 320. X. Lin and J. Ballmann. A Riemann Solver and a Second–Order Godunov Method for Elastic–Plastic Wave Propagation in Solids.Int. J. Impact Engng., 13(3):463–478, 1993.
\par\hangindent=2em\hangafter=1 321. X. Lin and Ballmann J. Numerical Method for Elastic–Plastic Waves in Cracked Solids, Part 1: Anti–Plane Shear Wave.Archive of Applied Mechanics, 63:261–282, 1993.
\par\hangindent=2em\hangafter=1 322. X. Lin and Ballmann J. Numerical Method for Elastic–Plastic Waves in Cracked Solids, Part 2: Plain Strain Problem. Archive of Applied Mechan- ics, 63:283–295, 1993.
\par\hangindent=2em\hangafter=1 323. X. Lin and Ballmann J. Reconsideration of Chen’s Problem by Finite Difference Method. Engineering Fracture and Mechanics, 44(5):735–739, 1993.
\par\hangindent=2em\hangafter=1 324. T. J. Linde. A Three–Dimensional Adaptive Multifluid MHD Model for the Heliosphere. PhD thesis, Aerospace Engineering, University of Michigan, USA, 1998.
\par\hangindent=2em\hangafter=1 325. T. J. Linde, T. I. Gombosi, P. L. Roe, K. G. Powell, and D. L. DeZeeuw. The Heliosphere in Magnetized Local Interstellar Medium: Results of a 3D MHD Simulation. J. Geophys. Res., 103(A2):1989–1904, 1998.
\par\hangindent=2em\hangafter=1 326. W. B. Lindquist. Current Progress in Hyperbolic Systems: Riemann Problems and Computations. American Mathematical Society, 1989.
\par\hangindent=2em\hangafter=1 327. M. S. Liou. Recent Progress and Applications of AUSM+. In Sixteenth Inter- national Conference on Numerical Methods in Fluid Dynamics, Lecture Notes in Physics 515, pages 302–307. Springer–Verlag, 1998.
\par\hangindent=2em\hangafter=1 328. M. S. Liou and C. J. Steffen. A New Flux Splitting Scheme.J. Comput. Phys., 107:23–39, 1993.
\par\hangindent=2em\hangafter=1 329. R. Liska and B. Wendroff. Composite Schemes for Conservation Laws. SIAM J. Numerical Analysis, 35(6):2250–2271, 1998.
\par\hangindent=2em\hangafter=1 330. X. D. Liu and P. D. Lax. Positive Schemes for Solving Multi–Dimensional Hy- perbolic Conservation Laws. Technical Report 95–003, Courant Mathematics and Computing Laboratory, New York University, USA, 1995.
\par\hangindent=2em\hangafter=1 331. X. D. Liu and T. Osher, S. Chan. Weighted Essentially Non–oscillatory Schemes. J. Comput. Phys., 115:200–212, 1994.
\par\hangindent=2em\hangafter=1 332. C. Y. Loh, M. S. Liou, and W. H. Hui. An Investigation of Random Choice Method for Three Dimensional Steady Supersonic Flow. In First Asian Com- putational Fluid Dynamics Conference, pages 963–971. The Hong Kong Uni- versity of Science and Technology, 1995.
\par\hangindent=2em\hangafter=1 333. J. Lorenz and H. J. Schroll. Stiff Well–Posedness for Hyperbolic Systems with Large Relaxation Terms. Technical Report 125, Institut f¨ ur Geometrie un Praktische Mathematik, RWTH Aachen, Germany, 1996.
\par\hangindent=2em\hangafter=1 334. H. Luo, J. D. Baum, and R. L¨ ohner. A Hermite WENO–Based Limiter for Discontinuous Galerkin Method on Unstructured Grids. J. Comput. Phys. , 225(1):686–713, 2007.
\par\hangindent=2em\hangafter=1 335. A. Majda and S. Osher. Numerical Viscosity and Entropy Condition. Comm. Pure and Applied Mathematics, 32:797–838, 1979.
\par\hangindent=2em\hangafter=1 336. J. C. Mandal and S. M. Desphande. Kinetic Flux Vector Splitting for the Euler Equations. Computers and Fluids, 23:447–, 1994.
\par\hangindent=2em\hangafter=1 337. M. J. Maron and R. J. Lopez. Numerical Analysis. Wadsworth, 1991.
\par\hangindent=2em\hangafter=1 338. A. Marquina, J. M. Mart´ı, J. M. Iba˜nez, Miralles J. A., and R. Donat. Ultrarel- ativistic Hydrodynamics: High–Resolution Shock–Capturing Methods. Astron. Astrophys., 258:566–, 1992.
% [原PDF p719]
\par\hangindent=2em\hangafter=1 339. E. Marshall and R. M´endez. Computational Aspects of the Random Choice Method for Shallow Water Equations. J. Comput. Phys., 39:1–21, 1981.
\par\hangindent=2em\hangafter=1 340. G. Marshall and B. Plohr. A Random Choice Method for Two–Dimensional Steady Supersonic Shock Wave Diffraction Problems. J. Comput. Phys. , 56:410–427, 1984.
\par\hangindent=2em\hangafter=1 341. J. M. Mart´ı, J. M. Iba˜nez, and J. A. Miralles. Numerical Relativistic Hydro- dynamics: Local Characteristic Approach. Phys. Rev., 43:3794–, 1991.
\par\hangindent=2em\hangafter=1 342. J. M. Mart´ıa n dE .M ¨uller. The Analytical Solution of the Riemann Problem for Relativistic Hydrodynamics. Journal of Fluid Mechanics, 258:317–333, 1994.
\par\hangindent=2em\hangafter=1 343. J. M. Mart´ ı, E. M¨u l l e r ,J .A .F o n t ,J .M .I b a ˜nez, and A. Marquina. Morphology and Dynamics of Relativistic Jets. Technical Report MPA 895, Max–Planck– Institut f¨ ur Astrophysik, Garching, Germany, 1995.
\par\hangindent=2em\hangafter=1 344. Y. P. Marx. A Practical Implementation of Turbulence Models for the Com- putation of Three–Dimensional Separated Flows. Int. J. Numer. Meth. Fluids, 13:775–796, 1991.
\par\hangindent=2em\hangafter=1 345. Y. P. Marx. Evaluation of the Artificial Compressibility Method for the So- lution of the Incompressible Navier–Stokes Equations. In Notes on Numerical Fluid Mechanics, Vol. 35, pages 201–210. Vieweg, 1991.
\par\hangindent=2em\hangafter=1 346. Y. P. Marx. Time Integration Schemes for the Unsteady Incompressible Navier–Stokes Equations. J. Comput. Phys., 112:182–209, 1994.
\par\hangindent=2em\hangafter=1 347. J. H. Mathews. Numerical Methods. Prentice–Hall International, Inc., 1987.
\par\hangindent=2em\hangafter=1 348. C. Y. McNeil. Efficient Upwind Algorithms for Solution of the Euler and Navier–Stokes Equations. PhD thesis, Department of Aerospace Science, Cran- field University, UK, 1995.
\par\hangindent=2em\hangafter=1 349. R. Menikoff and B. J. Plohr. The Riemann Problem for Fluid Flow of Real Materials. Reviews of Modern Physics, 61:75–130, 1989.
\par\hangindent=2em\hangafter=1 350. I. S. Men’shov. Increasing the Order of Approximation of Godunov’s Scheme Using the Generalized Riemann Problem. USSR Comput. Math. Phys. , 30(5):54–65, 1990.
\par\hangindent=2em\hangafter=1 351. A. Mignone and G. Bodo. An HLLC Riemann Solver for Relativistic Flows. Monthly Notices of the Royal Astronomical Society, 368(3):1040–1054, 2006.
\par\hangindent=2em\hangafter=1 352. A. R. Mitchell and D. F. Griffiths. The Finite Difference Method in Partial Differential Equations. John Wiley and Sons, 1980.
\par\hangindent=2em\hangafter=1 353. G. Moretti. The λ–scheme. Computers and Fluids, 7:191–205, 1979.
\par\hangindent=2em\hangafter=1 354. G. Moretti. Computation of Flows with Strong Shocks.Ann. Rev. Fluid Mech., 19:313–337, 1987.
\par\hangindent=2em\hangafter=1 355. G. Moretti and G. Bleich. A Time–Dependent Computational Method for Blunt–Body Flows. AIAA J., 4:2136–2141, 1966.
\par\hangindent=2em\hangafter=1 356. S. M¨uller and A. Voss. The Riemann Problem for the Euler Equations with Nonconvex and Nonsmooth Equation of State: Construction of Wave Curves. SIAM Journal on Scientific Computing, 28(2):651–681, 2006.
\par\hangindent=2em\hangafter=1 357. S. T. Munkejord, S. Evje, and T. Flatten. A Multi–Stage Centred Scheme Approach Applied to a Drift Flux Two–Phase Flow Model. Int. J. Numer. Meth. Fluids, 52:679–705, 2006.
\par\hangindent=2em\hangafter=1 358. C. D. Munz. A Tracking Method for Gas Flow into Vacuum Based on the Vac- uum Riemann Problem. Mathematical Methods in Applied Sciences, 17:597– 612, 1994.
\par\hangindent=2em\hangafter=1 359. C. D. Munz, R. Schneider, and O. Gerlinger. The Numerical Approximation of a Free Gas–Vacuum Interface. In Numerical Methods for the Navier–Stokes Equations, pages 181–190. Vieweg–Verlag, 1994.
% [原PDF p720]
\par\hangindent=2em\hangafter=1 360. S. Navarro-Martinez and O. R. Tutty. Numerical Simulation of G¨ortler Vortices in Hypersonic Compression Ramps. Computers and Fluids , 34(2):225–247, 2005.
\par\hangindent=2em\hangafter=1 361. A. Nemec and M. J. Aftosmis. Adjoint Sensitivity Computations for an Embedded–Boundary Cartesian Mesh Method. J. Comput. Phys. , 227(4):2724–2742, 2008.
\par\hangindent=2em\hangafter=1 362. H. Nessyahu and E. Tadmor. Non–Oscillatory Central Differencing for Hyper- bolic Conservation Laws. J. Comput. Phys., 87:408–463, 1990.
\par\hangindent=2em\hangafter=1 363. N. Nikiforakis and J. F. Clarke. Quasi–Steady Structures in the Two Dimen- sional Initiation of Detonations.Proc. Roy. Soc. Lond. A, 452:2023–2042, 1996.
\par\hangindent=2em\hangafter=1 364. H. Olivier and H. Gr¨ onig. The Random Choice Method Applied to Two- Dimensional Shock Focusing and Diffraction. J. Comput. Phys. , 63:85–106, 1986.
\par\hangindent=2em\hangafter=1 365. A. K. Oppenheim and R. I. Soloukhin. Experiments in Gas Dynamics of Explosions. Annual Review of Fluid Mechanics, 5:31–58, 1973.
\par\hangindent=2em\hangafter=1 366. E. Oran and J. P. Boris. Numerical Simulation of Reactive Flow. Cambridge University Press, 2001.
\par\hangindent=2em\hangafter=1 367. E. S. Oran and J. P. Boris. Detailed Modelling of Combustion Systems. Prog. Energy. Comb. Sci., 7:1–72, 1981.
\par\hangindent=2em\hangafter=1 368. E. S. Oran and J. P. Boris. Numerical Simulation of Reactive Flow. Elsevier, New York, 1987.
\par\hangindent=2em\hangafter=1 369. S. Osher. Riemann Solvers, the Entropy Condition and Difference Approxima- tions. SIAM J. Numer. Anal., 21(2):217–235, 1984.
\par\hangindent=2em\hangafter=1 370. S. Osher and S. Chakravarthy. Upwind Schemes and Boundary Conditions with Applications to Euler Equations in General Geometries. J. Comput. Phys. , 50:447–481, 1983.
\par\hangindent=2em\hangafter=1 371. S. Osher and S. Chakravarthy. High Resolution Schemes and the Entropy Condition. SIAM J. Numer. Anal., 21(5):955–984, 1984.
\par\hangindent=2em\hangafter=1 372. S. Osher and F. Solomon. Upwind Difference Schemes for Hyperbolic Conser- vation Laws. Math. Comp., 38,158:339–374, 1982.
\par\hangindent=2em\hangafter=1 373. Liu T. P. and J. A. Smoller. On the Vacuum State for the Isentropic Gas Dynamics Equations. Advances in Applied Mathematics, 1:345–359, 1980.
\par\hangindent=2em\hangafter=1 374. M. Pandolfi. Computation of Steady Supersonic Flows by a Flux–Difference Splitting Method. Computers and Fluids, 10:37–46, 1985.
\par\hangindent=2em\hangafter=1 375. C. Par´es. Numerical Methods for Nonconservative Hyperbolic Systems: a The- oretical Framework. SIAM Journal on Numerical Analysis, 44:300–321, 2006.
\par\hangindent=2em\hangafter=1 376. M. Pelanti, F. Bouchut, and A. Mangeney. A Roe–Type Scheme for Two–Phase Shallow Granular Flows over Variable Topography. Mathematical Modelling and Numerical Analysis, 42:851–885, 2008.
\par\hangindent=2em\hangafter=1 377. R. B. Pember. Numerical Methods for Hyperbolic Conservation Laws with Stiff Relaxation I: Spurious Solutions. SIAM J. Appl. Math. , 53:1293–1330, 1993.
\par\hangindent=2em\hangafter=1 378. R. B. Pember. Numerical Methods for Hyperbolic Conservation Laws with Stiff Relaxation II: Higher–Order Godunov Methods. SIAM J. Sci. Comput., 14:824–859, 1993.
\par\hangindent=2em\hangafter=1 379. L. L. Pennisi. Elements of Complex Analysis . Holt, Rinehart and Winston, 1963.
\par\hangindent=2em\hangafter=1 380. B. Perthame. Boltzmann Type Schemes for Gas Dynamics and the Entropy Property. SIAM J. Numer. Anal., 27:1405–1421, 1990.
% [原PDF p721]
\par\hangindent=2em\hangafter=1 381. B. Perthame. Second–Order Boltzman Schemes for Compressible Euler Equa- tions in one and two Space Dimensions.SIAM J. Numer. Anal., 29:1992, 1992.
\par\hangindent=2em\hangafter=1 382. L. Pesch and J. J. W. van der Vegt. A Discontinuous Galerkin Finite Element Discretization of the Euler Equations for Compressible and Incompressible Fluids. J. Comput. Phys., 227(11):5426–5446, 2008.
\par\hangindent=2em\hangafter=1 383. I. G. Petrovskii. Partial Differential Equations. London Cliffe Books Ltd., 1967.
\par\hangindent=2em\hangafter=1 384. J. Pike. Riemann Solvers for Perfect and Near–Perfect Gases. AIAA Journal, 31(10):1801–1808, 1993.
\par\hangindent=2em\hangafter=1 385. B. Pitman and L. Le. A Two–Fluid Model for Avalanche and Debris Flows. Phil. Trans. R. Soc. A, 363:1573–1601, 2005.
\par\hangindent=2em\hangafter=1 386. B. J. Plohr and D. H. Sharp. Riemann Problems and their Application to Ultra–Relativistic Heavy Ion Collisions. In The VIII International Confer- ence on Mathematical Physics (Mebkhout and S´en´eor, editors), pages 708–713. World Scientific Pub. Co., Singapore, 1987.
\par\hangindent=2em\hangafter=1 387. K. H. Prendergast and K. Xu. Numerical Hydrodynamics from Gas–Kinetic Theory. J. Comput. Phys., 109:53–66, 1993.
\par\hangindent=2em\hangafter=1 388. A. Prosperetti and G. Tryggvason. Fundamentals of Multiphase Flow. Cam- bridge University Press, 2007.
\par\hangindent=2em\hangafter=1 389. D. I. Pullin. Direct Simulation Methods for Compressible Inviscid Ideal Gas Flow. J. Comput. Phys., 34:231–244, 1980.
\par\hangindent=2em\hangafter=1 390. N. Qin. A Comparative Study of Two Upwind Schemes as Applied to Navier– Stokes Solutions for Resolving Boundary Layers in Hypersonic Viscous Flow. Technical Report GU Aero Report 9120, Department of Aerospace Engineer- ing, University of Glasgow, UK, 1992.
\par\hangindent=2em\hangafter=1 391. N. Qin and G. W. Foster. Study of Flow Interactions due to a Supersonic Lateral Jet Using High Resolution Navier–Stokes Solutions. InAIAA–95–2151 Paper, 27th AIAA Fluid Dynamics Conference, June 1995, San Diego, USA, 1995.
\par\hangindent=2em\hangafter=1 392. N. Qin and A. Redlich. Massively Separated Flows due to Transverse Sonic Jet in Hypersonic Laminar Stream. In AIAA–96–1934 Paper, 27th AIAA Fluid Dynamics Conference, June 1996, New Orleans, USA, 1996.
\par\hangindent=2em\hangafter=1 393. N. Qin and B. E. Richards. Sparse Quasi–Newton Method for High Resolution Schemes. In Notes on Numerical Fluid Mechanics, Vol. 20, pages 310–317. Vieweg, 1988.
\par\hangindent=2em\hangafter=1 394. N. Qin and B. E. Richards. Sparse Quasi–Newton Method for Navier–Stokes Solution. In Notes on Numerical Fluid Mechanics, Vol. 29, pages 474–483. Vieweg, 1990.
\par\hangindent=2em\hangafter=1 395. N. Qin, K. W. Scriba, and B. E. Richards. Shock–Shock, Shock–Vortex Interac- tion and Aerodynamic Heating In Hypersonic Corner Flow. The Aeronautical Journal, pages 152–160, 1991.
\par\hangindent=2em\hangafter=1 396. N. Qin, X. Xu, and B. E. Richards. Newton–Like Methods for Fast High Resolution Simulation of Hypersonic Viscous Flows. Computing Systems in Engineering, 3:429–435, 1992.
\par\hangindent=2em\hangafter=1 397. J. Qiu, B. C. Khoo, and C. W. Shu. A Numerical Study for the Performance of the Runge–Kutta Discontinuous Galerkin Method Based on Different Nu- merical Fluxes. J. Comput. Phys., 212(2):540–565, 2006.
\par\hangindent=2em\hangafter=1 398. L. Quartapelle, L. Castelletti, A. Guardone, and G. Quaranta. Solution of the Riemann Problem of Classical Gasdynamics. J. Comput. Phys., 190:118–140, 2003.
% [原PDF p722]
\par\hangindent=2em\hangafter=1 399. J. J. Quirk. An Adaptive Grid Algorithm for Computational Shock Hydrody- namics. PhD Thesis, College of Aeronautics, Cranfield Institute of Technology, UK, 1991.
\par\hangindent=2em\hangafter=1 400. J. J. Quirk. An Alternative to Unstructured Grids for Computing Gas Dynamic Flows Around Arbitrarily Complex Two Dimensional Bodies. Computers and Fluids, 23(1):125–142, 1994.
\par\hangindent=2em\hangafter=1 401. S. Rhebergen, O. Bokhove, and J. J. W. van der Vegt. Discontinuous Galerkin Finite Element Methods for Hyperbolic Nonconservative Partial Differential Equations. J. Comput. Phys., 227:1887–1922, 2008.
\par\hangindent=2em\hangafter=1 402. R. D. Richtmyer and K. W. Morton. Difference Methods for Initial Value Problems. Interscience-Wiley, New York, 1967.
\par\hangindent=2em\hangafter=1 403. N. H. Risebro and A. Tveito. Front Tracking Applied to Nonstrictly Hyperbolic Systems of Conservation Laws. SIAM J. Sci. Stat. Comput. , 12:1401–1419, 1991.
\par\hangindent=2em\hangafter=1 404. A. Rizzi. Numerical Implementation of Solid-Body Boundary Conditions for the Euler Equations. ZAMM, 58:301–304, 1978.
\par\hangindent=2em\hangafter=1 405. P. J. Roache. Computational Fluid Dynamics. Hermosa Publishers, 1982.
\par\hangindent=2em\hangafter=1 406. T. W. Roberts. The Behavior of Flux Difference Splitting Schemes near Slowly Moving Shock Waves. J. Comput. Phys., 90:141–160, 1990.
\par\hangindent=2em\hangafter=1 407. P. L. Roe. Approximate Riemann Solvers, Parameter Vectors, and Difference Schemes. J. Comput. Phys., 43:357–372, 1981.
\par\hangindent=2em\hangafter=1 408. P. L. Roe. Numerical Algorithms for the Linear Wave Equation. Technical Report 81047, Royal Aircraft Establishment, Bedford, UK, 1981.
\par\hangindent=2em\hangafter=1 409. P. L. Roe. Some Contributions to the Modelling of Discontinuous Flows. In Proceedings of the SIAM/AMS Seminar, San Diego, 1983.
\par\hangindent=2em\hangafter=1 410. P. L. Roe. Generalized Formulation of TVD Lax–Wendroff Schemes. Technical Report ICASE 84–53, NASA Langley Research Center, Hampton, USA, 1984.
\par\hangindent=2em\hangafter=1 411. P. L. Roe. Discrete Models for the Numerical Analysis of Time–Dependent Multidimensional Gas Dynamics. J. Comput. Phys., 63:458–476, 1986.
\par\hangindent=2em\hangafter=1 412. P. L. Roe. Upwind Differencing Schemes for Hyperbolic Conservation Laws with Source Terms. In Proc. First International Conference on Hyperbolic Problems, Carasso, Raviart and Serre (Editors), pages 41–51. Springer, 1986.
\par\hangindent=2em\hangafter=1 413. P. L. Roe. Remote Boundary Conditions for Unsteady Multidimensional Aero- dynamic Computations. Computers and Fluids, 17:221–231, 1989.
\par\hangindent=2em\hangafter=1 414. P. L. Roe. Sonic Flux Formulae. SIAM J. Sci. Stat. Comput., 13(2):611–630, 1992.
\par\hangindent=2em\hangafter=1 415. P. L. Roe. The Harten Memorial Lecture–New Applications of Upwinding. In Numerical Methods for Wave Propagation. Toro, E. F. and Clarke, J. F. (Editors), pages 1–31. Kluwer Academic Publishers, 1998.
\par\hangindent=2em\hangafter=1 416. P. L. Roe and J. Pike. Efficient Construction and Utilisation of Approximate Riemann Solutions. In Computing Methods in Applied Science and Engineer- ing. North–Holland, 1984.
\par\hangindent=2em\hangafter=1 417. E. Romenski, E. D. Resnyanski, and E. F. Toro. Conservative Hyperbolic Formulation for Compressible Two–Phase Flow with Different Phase Pressures and Temperatures. Quarterly of Applied Mathematics, 65:259–279, 2007.
\par\hangindent=2em\hangafter=1 418. V. V. Rusanov. Calculation of Interaction of Non–Steady Shock Waves with Obstacles. J. Comput. Math. Phys. USSR, 1:267–279, 1961.
\par\hangindent=2em\hangafter=1 419. G. Russo, E. F. Toro, and V. A. Titarev. ADER–Runge–Kutta Schemes for Conservation Laws in One Space Dimension. In Hyperbolic Problems: Theory,
% [原PDF p723]
\par\hangindent=2em\hangafter=1 Numerics and Applications. Benzone–Gavage S. and Serre D. (Editors), Vol. 1, pages 929–936. Springer, 2008.
\par\hangindent=2em\hangafter=1 420. A. V. Safranov. Difference Method for Gasdynamical Equations Based on the Jump Conditions (in Russian). Mathematical Modelling, 20:76–84, 2008.
\par\hangindent=2em\hangafter=1 421. L. Sainsaulieu. Finite Volume Approximation of Two Phase–Fluid Flows Based on an Approximate Roe–Type Riemann Solver. J. Comput. Phys., 121:1–28, 1995.
\par\hangindent=2em\hangafter=1 422. T. Saito and I. I. Glass. Application of Random Choice to Problems in Shock and Detonation Wave Dynamics. Technical Report UTIAS 240, Institute for Aerospace Studies, University of Toronto, 1979.
\par\hangindent=2em\hangafter=1 423. D. A. S´anchez. Ordinary Differential Equations and Stability Theory.D o v e r Publications, Inc., 1968.
\par\hangindent=2em\hangafter=1 424. R. H. Sanders and K. H. Prendergast. On the Origin the Three Kiloparsec Arm. Astrophys. J., 188:489–500, 1974.
\par\hangindent=2em\hangafter=1 425. A. Sangan. An HLLC Scheme for Ten–Moments Approximation Coupled with Magnetic Field. Internat. J. Comput. Science and Math., 2(1/2):73–109, 2008.
\par\hangindent=2em\hangafter=1 426. R. Saurel and R. Abgrall. A Multiphase Godunov Method for Compressible Multifluid and Multiphase Flows. J. Comput. Phys., 150:425–467, 1999.
\par\hangindent=2em\hangafter=1 427. R. Saurel, E. Daniel, and J. C. Loraud. Two–Phase Flows: Second Order Schemes and Boundary Conditions. AIAA Journal, 32(6):1214–1221, 1994.
\par\hangindent=2em\hangafter=1 428. R. Saurel, A. Forestier, D. Veyret, and J.-C. Loraud. A Finite Volume Scheme for Two–Phase Compressible Flows. Int. J. Numer. Meth. Fluids, 18:803–819, 1994.
\par\hangindent=2em\hangafter=1 429. R. Saurel, M. Larini, and J. C. Loraud. Exact and Approximate Riemann Solvers for Real Gases. J. Comput. Phys., 112:126–137, 1994.
\par\hangindent=2em\hangafter=1 430. M. Schleicher. Ein Einfaches und Effizientes Verfahren zur Loesung des Riemann–Problems. Z. Flugwiss. Weltraumforsch., 17:265–269, 1993.
\par\hangindent=2em\hangafter=1 431. V. Schneider, U. Katscher, D. H. Rischke, B. Waldhauser, J. A. Maruhn, and C. D. Munz. New Algorithms for Ultra–Relativistic Numerical Hydrodynamics. J. Comput. Phys., 105:92–, 1993.
\par\hangindent=2em\hangafter=1 432. H. J. Schroll and A. Tveito. A System of Conservation Laws with a Relaxation Term. Technical Report 103, Institut f¨ ur Geometrie un Praktische Mathematik, RWTH Aachen, Germany, 1994.
\par\hangindent=2em\hangafter=1 433. H. J. Schroll, A. Tveito, and R. Winther. An Error Bound for a Finite Differ- ence Scheme Applied to a Stiff System of Conservation Laws. Technical Report 1994–3, Department of Informatics, University of Oslo, Norway, 1994.
\par\hangindent=2em\hangafter=1 434. H. J. Schroll and R. Winther. Finite Difference Schemes for Scalar Conservation Laws with Source Terms. Technical Report 107, Institut f¨ ur Geometrie un Praktische Mathematik, RWTH Aachen, Germany, 1994.
\par\hangindent=2em\hangafter=1 435. T. Schwartzkopff. Finite–Volumen Verfahren Hoher Ordnung und Heterogene Gebietszerlegung f¨ur die Numerische Aeroakustik. PhD thesis, Institut f¨ur Aero- un Gasdynamik, Universit¨ at Stuttgart, Germany, 2005.
\par\hangindent=2em\hangafter=1 436. T. Schwartzkopff, M. Dumbser, and C. D. Munz. Fast High Order ADER Schemes for Linear Hyperbolic Equations. J. Comput. Phys. , 197:532–539, 2004.
\par\hangindent=2em\hangafter=1 437. T. Schwartzkopff, C. D. Munz, and E. F. Toro. ADER: High–Order Approach for Linear Hyperbolic Systems in 2D. J. Scientific Computing , 17:231–240, 2002.
% [原PDF p724]
\par\hangindent=2em\hangafter=1 438. D. W. Schwendeman, C. W. Wahle, and A. K. Kapila. The Riemann Problem and a High–Resolution Godunov Method for a Model of Compressible Two– Phase Flow. J. Comput. Phys., 212:490–526, 2006.
\par\hangindent=2em\hangafter=1 439. F. W. Sears and G. L. Salinger. Thermodynamics, Kinetic Theory and Statis- tical Thermodynamics. Addison–Wesley, 1986.
\par\hangindent=2em\hangafter=1 440. C. C. L. Sells. . Technical Report RAE TR 80065, Royal Aircraft Establish- ment, Bedford, UK, 1980.
\par\hangindent=2em\hangafter=1 441. L. F. Shampine. Numerical Solution of Ordinary Differential Equations. Chap- man and Hall, London and New York, 1994.
\par\hangindent=2em\hangafter=1 442. J. S. Shang. A Fractional–Step Method for Solving 3D, Time–Domain Maxwell Equations. J. Comput. Phys., 118:109–119, 1995.
\par\hangindent=2em\hangafter=1 443. Z. C. Shi and J. J. Gottlieb. Random Choice Method for Two Dimensional Planar and Axisymmetric Steady Supersonic Flows. Technical Report 297, UTIAS, University of Toronto, Canada, 1985.
\par\hangindent=2em\hangafter=1 444. C. Shu and S. Osher. Efficient Implementation of Essentially Non–oscillatory Shock–Capturing Schemes II. J. Comput. Phys., 83:32–78, 1989.
\par\hangindent=2em\hangafter=1 445. C. W. Shu. TVB Boundary Treatment for Numerical Solutions of Conservation Laws. Math. Comp., 49(179):123–134, 1987.
\par\hangindent=2em\hangafter=1 446. C. W. Shu. TVB Uniformly High–Order Schemes for Conservation Laws.Math. Comput., 49(179):105–121, 1987.
\par\hangindent=2em\hangafter=1 447. C. W. Shu and S. Osher. Efficient Implementation of Essentially Non– oscillatory Shock–Capturing Schemes. J. Comput. Phys., 77:439–471, 1988.
\par\hangindent=2em\hangafter=1 448. G. Singh. A Numerical Study of the Evolution of Piston Driven Detonation . PhD thesis, Department of Arospace Science, Cranfield University, UK, 1990.
\par\hangindent=2em\hangafter=1 449. G. Singh and J. F. Clarke. Transient Phenomena in the Initiation of a Me- chanically Driven Plane Detonation. Proc. Roy. Soc. London A, 438:23–46, 1992.
\par\hangindent=2em\hangafter=1 450. G. D. Smith. Numerical Solution of Partial Differential Equations. Clarendon Press, Oxford, 1978.
\par\hangindent=2em\hangafter=1 451. J. Smoller. Shock Waves and Reaction–Diffusion Equations. Springer–Verlag, 1994.
\par\hangindent=2em\hangafter=1 452. G. A. Sod. A Numerical Study of a Converging Cylindrical Shock. J. Fluid Mechanics, 83:785–794, 1977.
\par\hangindent=2em\hangafter=1 453. G. A. Sod. A Survey of Several Finite Difference Methods for Systems of Nonlinear Hyperbolic Conservation Laws. J. Comput. Phys., 27:1–31, 1978.
\par\hangindent=2em\hangafter=1 454. G. A. Sod. Numerical Methods in Fluid Dynamics . Cambridge University Press, 1985.
\par\hangindent=2em\hangafter=1 455. G. A. Sod. A Random Choice Method for the Stefan Problem. Computers Math. Applic., 19:1–9, 1990.
\par\hangindent=2em\hangafter=1 456. W. Speares. A Finite Volume Approach to the Weighted Average Flux Method. MSc thesis, College of Aeronautics, Cranfield Institute of Technology, UK, 1991.
\par\hangindent=2em\hangafter=1 457. W. Speares and E. F. Toro. A High-Resolution Algorithm for Time–Dependent Shock Dominated Problems with Adaptive Mesh Refinement. Z. Flugwiss. Weltraumforsch. (Journal of Flight Sciences and Space Research), 19:267–281, 1995.
\par\hangindent=2em\hangafter=1 458. S. P. Spekreijse. Multigrid Solution of Monotone Second–Order Discretizations of Hyperbolic Conservation Laws. Math. Comp., 49(179):135–155, 1987.
% [原PDF p725]
\par\hangindent=2em\hangafter=1 459. S. P. Spekreijse. Multigrid Solution of the Steady Euler Equations.P h Dt h e s i s , Centre for Mathematics and Computer Science, Amsterdam, The Netherlands, 1988.
\par\hangindent=2em\hangafter=1 460. H. St¨adtke. Gasdynamic Aspects of Two–Phase Flow . Wiley–VCH, Berlin, 2006.
\par\hangindent=2em\hangafter=1 461. H. St¨a d t k e ,G .F r a n c h e l l o ,B .W o r t h ,U .G r a f ,P .R o m s t e d t ,A .K u m b a r o , J. Garc´ıa-Cascales, H. Pailliere, H. Deconinck, M. Ricchiuto, B. Smith, F. De Cachard, E. F. Toro, E. Romenski, and S. Mimouni. Advanced Three– Dimensional Two–Phase Flow Simulation Tools for Application to Reactor Safety (ASTAR). Nuclear Engineering and Design, 235:379–400, 2005.
\par\hangindent=2em\hangafter=1 462. J. L. Steger, F. C. Dougherty, and J. A. Benek. A Chimera Grid Scheme. In Advances in Grid Generation, ASME FED, Vol. 5, 1983.
\par\hangindent=2em\hangafter=1 463. J. L. Steger and R. F. Warming. Flux Vector Splitting of the Inviscid Gasdy- namic Equations with Applications to Finite Difference Methods. J. Comput. Phys., 40:263–293, 1981.
\par\hangindent=2em\hangafter=1 464. H. B. Stewart and B. Wendroff. Two–Phase Flow: Models and Methods. J. Comput. Phys., 56:363–409, 1984.
\par\hangindent=2em\hangafter=1 465. J. J. Stoker. Water Waves. The Mathematical Theory with Applications. John Wiley and Sons, 1992.
\par\hangindent=2em\hangafter=1 466. G. Strang. On the Construction and Comparison of Difference Schemes.SIAM J. Numer. Anal., 5(3):506–517, 1968.
\par\hangindent=2em\hangafter=1 467. R. C. Swanson and E. Turkel. On Central–Difference and Upwind Methods. Technical Report ICASE 90–44, NASA Langley Research Center, Hampton, Virginia, USA, 1990.
\par\hangindent=2em\hangafter=1 468. B. K. Swartz and B. Wendroff. AZTEC: A Front Tracking Code Based on Godunov’s Method. Applied Numerical Mathematics, 2:385–397, 1986.
\par\hangindent=2em\hangafter=1 469. P. K. Sweby. Shock Capturing Schemes. PhD thesis, Department of Mathe- matics, University of Reading, UK, 1982.
\par\hangindent=2em\hangafter=1 470. P. K. Sweby. High Resolution Schemes Using Flux Limiters for Hyperbolic Conservation Laws. SIAM J. Numer. Anal., 21:995–1011, 1984.
\par\hangindent=2em\hangafter=1 471. P. K. Sweby. TVD Schemes for Inhomogeneous Conservation Laws. In Notes on Numerical Fluid Mechanics, Vol. 24, Non–Linear Hyperbolic Equations– Theory, Computation Methods and Applications, pages 599–607. Vieweg, 1989.
\par\hangindent=2em\hangafter=1 472. P. K. Sweby and M. J. Baines. On Convergence of Roe’s Scheme for the General Non–linear Wave Equation. J. Comput. Phys., 56(1):135–148, 1984.
\par\hangindent=2em\hangafter=1 473. E. Tadmor. Convenient Total Variation Diminishing Conditions for Non– Linear Difference Schemes. SIAM J. Numer. Anal., 25:1002–1014, 1988.
\par\hangindent=2em\hangafter=1 474. M. Takahiro and K. Kanya. A Multi–State HLL Approximate Riemann Solver for Ideal Magnetohydrodynamics. J. Comput. Phys., 208(1):315–344, 2005.
\par\hangindent=2em\hangafter=1 475. Y. Takakura and E. F. Toro. Arbitrarily Accurate Non–Oscillatory Schemes for a Non–Linear Conservation Law. J. Computational Fluid Dynamics, 11:7–18, 2002.
\par\hangindent=2em\hangafter=1 476. Y. Takano. An Application of the Random Choice Method to a Reactive Gas with Many Chemical Species. J. Comput. Phys., 67:173–187, 1986.
\par\hangindent=2em\hangafter=1 477. K. Takayama. Optical Flow Visualization of Shock Wave Phenomena (Paul Vieille Memorial Lecture). In Shock Waves @ Marseille, Vol. 4 (Brun and Dumitrescu, Editors), pages 7–16. Springer–Verlag, 1995.
\par\hangindent=2em\hangafter=1 478. K. Takayama (Editor). Shock Waves (Proc. 18th Int. Symp. on Shock Waves, Sendai, Japan, 1991), Two Volumes. Springer–Verlag, 1992.
% [原PDF p726]
\par\hangindent=2em\hangafter=1 479. S. Taki and T. Fujiwara. Numerical Analysis of Two–Dimensional Nonsteady Detonations. AIAA Journal, 16:73–77, 1978.
\par\hangindent=2em\hangafter=1 480. C. K. W. Tam and J. C. Webb. Dispersion–Relation–Preserving Finite Differ- ence Schemes for Computational Acoustics. J. Comput. Phys., 107:262–281, 1993.
\par\hangindent=2em\hangafter=1 481. P. Tamamidis, G. Zhang, and D. N. Assanis. Comparison of Pressure–Based and Artificial Compressibility Methods for Solving 3D Steady Incompressible Viscous Flows. J. Comput. Phys., 124(1):1–13, 1996.
\par\hangindent=2em\hangafter=1 482. T. Tanaka. Finite Volume TVD Scheme on an Unstructured Grid System for Three–Dimensional MHD Simulation of Inhomogeneous Systems Including Strong Background Potential Fields. J. Comput. Phys., 111:381–, 1994.
\par\hangindent=2em\hangafter=1 483. L. Tatsien and Y. Wenci. Boundary–Value Problems for Quasi–Linear Hyper- bolic Systems. Duke University Mathematics Series, 1985.
\par\hangindent=2em\hangafter=1 484. J. W. Thomas. Numerical Partial Differential Equations, Texts in Applied Mathematics 22. Springer–Verlag, 1998.
\par\hangindent=2em\hangafter=1 485. J. F. Thompson, F. C. Thames, and C. W. Mastin. Automatic Numerical Grid Generation of Body–Fitted Curvilinear Coordinate System for Field Contain- ing any Number of Arbitrary Two–Dimensional Bodies. J. Comput. Phys. , 15:299–319, 1974.
\par\hangindent=2em\hangafter=1 486. J. F. Thompson, J. F. Warsi, and C. W. Mastin.Grid Generation: Foundations and Applications. North–Holland, 1982.
\par\hangindent=2em\hangafter=1 487. J. F. (Editor) Thompson. Numerical Grid Generation. North–Holland, 1982.
\par\hangindent=2em\hangafter=1 488. K. W. Thompson. The Special Relativistic Shock Tube. J. Fluid Mechanics, 171:365–, 1986.
\par\hangindent=2em\hangafter=1 489. K. W. Thompson. Time Dependent Boundary Conditions for Hyperbolic Sys- tems. J. Comput. Phys., 68:1–24, 1987.
\par\hangindent=2em\hangafter=1 490. V. A. Titarev. Derivative Riemann Problem and ADER Schemes.P h Dt h e s i s , Department of Mathematics, University of Trento, Italy, 2005.
\par\hangindent=2em\hangafter=1 491. V. A. Titarev, E. Romenski, and E. F. Toro. MUSTA–Type Upwind Fluxes for Non–Linear Elasticity. Int. J. Numer. Meth. Engineering, 73:897–926, 2008.
\par\hangindent=2em\hangafter=1 492. V. A. Titarev and E. F. Toro. ADER: Arbitrary High Order Godunov Ap- proach. J. Scientific Computing , 17:609–618, 2002.
\par\hangindent=2em\hangafter=1 493. V. A. Titarev and E. F. Toro. Finite Volume WENO Schemes for Three- Dimensional Conservation Laws. J. Comput. Phys., 201(1):238–260, 2004.
\par\hangindent=2em\hangafter=1 494. V. A. Titarev and E. F. Toro. ADER Schemes for Three–Dimensional Hyper- bolic Systems. J. Comput. Phys., 204:715–736, 2005.
\par\hangindent=2em\hangafter=1 495. V. A. Titarev and E. F. Toro. MUSTA Schemes for Multi–Dimensional Hy- perbolic Systems: Analysis and Improvements. Int. J. Numer. Meth. Fluids , 49:117–147, 2005.
\par\hangindent=2em\hangafter=1 496. E. F. Toro. Multi–Stage Predictor–Corrector Fluxes for Hyperbolic Equations. Technical Report NI03037–NPA, Isaac Newton Institute for Mathematical Sci- ences, University of Cambridge, UK, 17th June, 2003.
\par\hangindent=2em\hangafter=1 497. E. F. Toro. A New Numerical Technique for Quasi-Linear Hyperbolic Systems of Conservation Laws. Technical Report CoA-8608, College of Aeronautics, Cranfield Institute of Technology, UK, 1986.
\par\hangindent=2em\hangafter=1 498. E. F. Toro. A Fast Riemann Solver with Constant Covolume Applied to the Random Choice Method. Int. J. Numer. Meth. Fluids, 9:1145–1164, 1989.
\par\hangindent=2em\hangafter=1 499. E. F. Toro. A Weighted Average Flux Method for Hyperbolic Conservation Laws. Proc. Roy. Soc. London, A423:401–418, 1989.
% [原PDF p727]
\par\hangindent=2em\hangafter=1 500. E. F. Toro. Random Choice Based Hybrid Schemes for one and Two– Dimensional Gas Dynamics. In Non–linear Hyperbolic Equations – Theory, Computation Methods and Applications. Notes on Numerical Fluid Mechanics, pages 630–639. Vieweg, Braunschweig, 1989.
\par\hangindent=2em\hangafter=1 501. E. F. Toro. Riemann–Problem Based Techniques for Computing Reactive Two– Phase Flows. In Dervieux and Larrouturrou, editors, Proc. Third. Intern. Confer. on Numerical Combustion, number 351 in Lecture Notes in Physics, pages 472–481, Antibes, France, May 1989.
\par\hangindent=2em\hangafter=1 502. E. F. Toro. A Linearised Riemann Solver for the Time–Dependent Euler Equa- tions of Gas Dynamics. Proc. Roy. Soc. London, A434:683–693, 1991.
\par\hangindent=2em\hangafter=1 503. E. F Toro. Some Aspects of Shock Capturing Methods for Gas Dynamics. Technical Report COA–9112, College of Aeronautics, Cranfield Institute of Technology, UK, 1991.
\par\hangindent=2em\hangafter=1 504. E. F. Toro. Riemann Problems and the WAF Method for Solving Two– Dimensional Shallow Water Equations. Phil. Trans. Roy. Soc. London Series A: Physical Sciences and Engineering, A338:43–68, 1992.
\par\hangindent=2em\hangafter=1 505. E. F. Toro. The Weighted Average Flux Method Applied to the Euler Equa- tions. Phil. Trans. Roy. Soc. London Series A: Physical Sciences and Engi- neering, A341:499–530, 1992.
\par\hangindent=2em\hangafter=1 506. E. F. Toro. The Weighted Average Flux Method Applied to the Time– Dependent Euler Equations. Phil. Trans. Roy. Soc. London, A341:499–530, 1992.
\par\hangindent=2em\hangafter=1 507. E. F. Toro. Viscous Flux Limiters. In J. B. Vos, A. Rizzi, and I. L. Ryhming, editors, Notes on Numerical Fluid Dynamics, Vol. 35, pages 592–600. Vieweg, 1992.
\par\hangindent=2em\hangafter=1 508. E. F. Toro. Defects of Conservative Approaches and Adaptive Primitive– Conservative Schemes for Computing Solutions to Hyperbolic Conservation Laws. Technical Report MMU 9401, Department of Mathematics and Physics, Manchester Metropolitan University, UK, 1994.
\par\hangindent=2em\hangafter=1 509. E. F. Toro. Direct Riemann Solvers for the Time–Dependent Euler Equations. Shock Waves, 5:75–80, 1995.
\par\hangindent=2em\hangafter=1 510. E. F. Toro. MUSCL–Type Primitive Variable Schemes. Technical Report MMU–9501, Department of Mathematics and Physics, Manchester Metropoli- tan University, UK, 1995.
\par\hangindent=2em\hangafter=1 511. E. F. Toro. On Adaptive Primitive–Conservative Schemes for Conservation Laws. In M. M. Hafez, editor, Sixth International Symposium on Compu- tational Fluid Dynamics: A Collection of Technical Papers , volume 3, pages 1288–1293, Lake Tahoe, Nevada, USA, September 4–8, 1995.
\par\hangindent=2em\hangafter=1 512. E. F Toro. Some IVPs for Which Conservative Methods Fail Miserably. In M. M. Hafez, editor, Sixth International Symposium on Computational Fluid Dynamics: A Collection of Technical Papers, volume 3, pages 1294–1299, Lake Tahoe, Nevada, USA, September 4-8, 1995.
\par\hangindent=2em\hangafter=1 513. E. F. Toro. On Glimm–Related Schemes for Conservation Laws. Techni- cal Report MMU–9602, Department of Mathematics and Physics, Manchester Metropolitan University, UK, 1996.
\par\hangindent=2em\hangafter=1 514. E. F. Toro. Exact and Approximate Riemann Solvers for the Artificial Com- pressibility Equations. Technical Report 97–02, Department of Mathematics and Physics, Manchester Metropolitan University, UK, 1997.
% [原PDF p728]
\par\hangindent=2em\hangafter=1 515. E. F. Toro. On Two Glimm–Related Schemes for Hyperbolic Conservation Laws. In Proc. Fifth Annual Conference of the CFD Society of Canada, pages 3.49–3.54. University of Victoria, Canada, 1997.
\par\hangindent=2em\hangafter=1 516. E. F. Toro. Riemann Solvers and Numerical Methods for Fluid Dynamics . Springer–Verlag, 1997.
\par\hangindent=2em\hangafter=1 517. E. F. Toro. Primitive, Conservative and Adaptive Schemes for Hyperbolic Conservation Laws. In Numerical Methods for Wave Propagation. Toro, E. F. and Clarke (Editors), J. F. , pages 323–385. Kluwer Academic Publishers, 1998.
\par\hangindent=2em\hangafter=1 518. E. F. Toro. NUMERICA: A Library of Source Codes for Teaching, Research and Applications. Numeritek Ltd., www.numeritek.com, 1999.
\par\hangindent=2em\hangafter=1 519. E. F. Toro. Riemann Solvers and Numerical Methods for Fluid Dynamics, Second Edition. Springer–Verlag, 1999.
\par\hangindent=2em\hangafter=1 520. E. F. Toro. Shock–Capturing Methods for Free–Surface Shallow Flows. Wiley and Sons Ltd., 2001.
\par\hangindent=2em\hangafter=1 521. E. F. Toro. A Multi–Stage Numerical Flux.Appl. Numer. Math., 56:1464–1479, 2006.
\par\hangindent=2em\hangafter=1 522. E. F. Toro. Riemann Solvers with Evolved Initial Conditions. Int. J. Numer. Meth. in Fluids, 52:433–453, 2006.
\par\hangindent=2em\hangafter=1 523. E. F Toro and Titarev V. A. Derivative Riemann Solvers for Systems of Con- servation Laws and ADER Methods. J. Comput. Phys., 212(1):150–165, 2006.
\par\hangindent=2em\hangafter=1 524. E. F Toro and Titarev V. A. MUSTA Schemes for Systems of Conservation Laws. J. Comput. Phys., 216:403–429, 2006.
\par\hangindent=2em\hangafter=1 525. E. F. Toro, S. J. Billet, and E. P. Boden. Advanced Modelling Techniques for Propulsion Systems. In Proceedings of Conference on Energetic Materials and Propulsion Technology, Salsbury, Australia, April 18–19, 1996,p a g e s– .– , 1996.
\par\hangindent=2em\hangafter=1 526. E. F. Toro and S. J. Billett. Centred TVD Schemes for Hyperbolic Conservation Laws. Technical Report MMU–9603, Department of Mathematics and Physics, Manchester Metropolitan University, UK, 1996.
\par\hangindent=2em\hangafter=1 527. E. F. Toro and S. J. Billett. A Unified Riemann–Problem Based Extension of the Warming–Beam and Lax–Wendroff Schemes. IMA J. Numer. Anal., 17:61–102, 1997.
\par\hangindent=2em\hangafter=1 528. E. F. Toro and S. J. Billett. Centred TVD Schemes for Hyperbolic Conservation Laws. IMA J. Numerical Analysis, 20:47–79, 2000.
\par\hangindent=2em\hangafter=1 529. E. F. Toro and R. E. Brown. The WAF Method and Splitting Procedures for Viscous, Shocked Flows. In Proceedings of the 18th International Symposium on Shock Waves, Tohoku University, Sendai, Japan (K. Takayama, Editor) , pages 1119–1126. Springer–Verlag, 1992.
\par\hangindent=2em\hangafter=1 530. E. F. Toro and C. E. Castro. The Derivative Riemann Problem for the Baer– Nunziato Equations. In Hyperbolic Problems: Theory, Numerics and Applica- tions. Vol. 1. Benzoni–Gavage and Serre (editors), pages 1045–1052. Springer, 2008.
\par\hangindent=2em\hangafter=1 531. E. F. Toro and C. E. Castro. The Derivative Riemann Problem for the Baer– Nunziato Equations. In Hyperbolic Problems: Theory, Numerics and Applica- tions HYPE2006, Vol. 1. S. Benzoni–Gavage and D. Serre (Editors) , pages 1045–1052. Springer, 2008.
\par\hangindent=2em\hangafter=1 532. E. F. Toro and A. Chakraborty. Development of an Approximate Riemann Solver for the Steady Supersonic Euler Equations. The Aeronautical Journal, 98:325–339, 1994.
% [原PDF p729]
\par\hangindent=2em\hangafter=1 533. E. F. Toro and C. C. Chou. A Linearised Riemann Solver for the Steady Supersonic Euler Equations. Int. J. Numer. Meth. Fluids, 16:173–186, 1993.
\par\hangindent=2em\hangafter=1 534. E. F. Toro and P. Garc´ıa-Navarro. Godunov–Type Methods for Free–Surface Shallow Flows: A Review. J. Hydraulic Research, 45(6):736–751, 2007.
\par\hangindent=2em\hangafter=1 535. E. F. Toro and A. Hidalgo. ADER Finite Volume Schemes for Diffusion– Reaction Equations. Applied Numerical Mathematics, 59:73–100, 2009.
\par\hangindent=2em\hangafter=1 536. E. F. Toro, A. Hidalgo, and M. Dumbser. FORCE Schemes on Unstruc- tured Meshes I: Conservative Hyperbolic Systems. Isaac Newton Institute for Mathematical Sciences, University of Cambridge, UK. Preprint NI09005-NPA, September 2008. Submitted to J. Comput. Phys., 2008.
\par\hangindent=2em\hangafter=1 537. E. F. Toro, R. C. Millington, and L. A. M. Nejad. Towards Very High–Order Godunov Schemes. In Godunov Methods: Theory and Applications. Edited R e v i e w ,E .F .T o r o( E d i t o r ), pages 905–937. Kluwer Academic/Plenum Pub- lishers, 2001.
\par\hangindent=2em\hangafter=1 538. E. F. Toro, M. Olim, and K. Takayama. Unusual Increase in Tsunami Wave Amplitude at the Okushiri Island: Mach Reflection of Shallow Water Waves. In Proceedings of the 22nd International Symposium on Shock Waves, Vol. II. Ball, Hillier and Roberts(Editors), pages 1207–1212. University of Southamp- ton, UK, 1999.
\par\hangindent=2em\hangafter=1 539. E. F. Toro and P. L. Roe. A Hybridised High–Order Random Choice Method for Quasi-Linear Hyperbolic Systems. In Gr¨ onig, editor, Proc. 16th Intern. Symp. on Shock Tubes and Waves , pages 701–708, Aachen, Germany, July 1987.
\par\hangindent=2em\hangafter=1 540. E. F. Toro and P. L. Roe. A Hybrid Scheme for the Euler Equations Using the Random Choice and Roe’s Methods. InNumerical Methods for Fluid Dynamics III. The Institute of Mathematics and its Applications Conference Series, New Series No. 17, Morton and Baines (Editors), pages 391–402. Oxford University Press, New York, 1988.
\par\hangindent=2em\hangafter=1 541. E. F. Toro, M. Spruce, and W. Speares. Restoration of the Contact Sur- face in the HLL–Riemann Solver. Technical Report CoA–9204, Department of Aerospace Science, Collegue of Aeronautics, Cranfield Institute of Technology, UK, 1992.
\par\hangindent=2em\hangafter=1 542. E. F. Toro, M. Spruce, and W. Speares. Restoration of the Contact Surface in the HLL–Riemann Solver. Shock Waves, 4:25–34, 1994.
\par\hangindent=2em\hangafter=1 543. E. F. Toro and V. A. Titarev. Solution of the Generalised Riemann Problem for Advection–Reaction Equations. Proc. Roy. Soc. London A, 458:271–281, 2002.
\par\hangindent=2em\hangafter=1 544. E. F. Toro and V. A. Titarev. ADER Schemes for Scalar Hyperbolic Conser- vation Laws with Source Terms in Three Space Dimensions.J. Comput. Phys., 202(1):196–215, 2005.
\par\hangindent=2em\hangafter=1 545. E. F. Toro (Editor). Godunov Methods: Theory and Applications. Edited Re- view. Kluwer Academic/Plenum Publishers, 2001.
\par\hangindent=2em\hangafter=1 546. I. Toumi. A Weak Formulation of Roe’s Approximate Riemann Solver. J. Comput. Phys., 102:360–373, 1992.
\par\hangindent=2em\hangafter=1 547. I. Toumi. An Upwind Numerical Method for Two–Fluid Two–Phase Flow Models. Nuclear Science and Engineering, 123:147–168, 1996.
\par\hangindent=2em\hangafter=1 548. I. Toumi and D. Caruge. An Implicit Second–Order Numerical Method for Three–Dimensional Two–Phase Flow Calculations. Nuclear Science and Engi- neering, 130:1–13, 1998.
% [原PDF p730]
\par\hangindent=2em\hangafter=1 549. I. Toumi and A. Kumbaro. An Approximate Linearized Riemann Solver for a Two–Fluid Model. J. Comput. Phys., 124:286–300, 1996.
\par\hangindent=2em\hangafter=1 550. E. Turkel and A. Arnone. Pseudo–Compressibility Methods for the Incom- pressible Flow Equations. Technical Report ICASE 93–66, NASA Langley Research Center, USA, 1993.
\par\hangindent=2em\hangafter=1 551. A. Tveito and R. Winther. Introduction to Partial Differential Equations. Springer–Verlag, 1998.
\par\hangindent=2em\hangafter=1 552. G. D. van Albada, B. van Leer, and W. W. Roberts. A Comparative Study of Computational Methods in Cosmic Gas Dynamics. Astron. Astrophysics, 108:76–84, 1982.
\par\hangindent=2em\hangafter=1 553. J. J. W. van der Vegt and H. van der Ven. Space–Time Discontinuous Galerkin Finite Element Method with Dynamic Grid Motion for Inviscid Compressible Flows: General Formulation. J. Comput. Phys., 182(2):546–585, 2002.
\par\hangindent=2em\hangafter=1 554. B. van Leer. Towards the Ultimate Conservative Difference Scheme I. The Quest for Monotonicity. Lecture Notes in Physics, 18:163–168, 1973.
\par\hangindent=2em\hangafter=1 555. B. van Leer. Towards the Ultimate Conservative Difference Scheme II. Mono- tonicity and Conservation Combined in a Second Order Scheme. J. Comput. Phys., 14:361–370, 1974.
\par\hangindent=2em\hangafter=1 556. B. van Leer. MUSCL, A New Approach to Numerical Gas Dynamics. InCom- puting in Plasma Physics and Astrophysics, Max–Planck–Institut f¨ur Plasma Physik, Garchung, Germany, April 1976.
\par\hangindent=2em\hangafter=1 557. B. van Leer. Towards the Ultimate Conservative Difference Scheme III. Upstream–Centered Finite Difference Schemes for Ideal Compressible Flow. J. Comput. Phys., 23:263–275, 1977.
\par\hangindent=2em\hangafter=1 558. B. van Leer. Towards the Ultimate Conservative Difference Scheme IV. A New Approach to Numerical Convection. J. Comput. Phys., 23:276–299, 1977.
\par\hangindent=2em\hangafter=1 559. B. van Leer. Towards the Ultimate Conservative Difference Scheme V. A Second Order Sequel to Godunov’s Method. J. Comput. Phys., 32:101–136, 1979.
\par\hangindent=2em\hangafter=1 560. B. van Leer. Flux–Vector Splitting for the Euler Equations. Technical Report ICASE 82–30, NASA Langley Research Center, USA, 1982.
\par\hangindent=2em\hangafter=1 561. B. van Leer. Flux Vector Splitting for the Euler Equations. In Proceedings of the 8th International Conference on Numerical Methods in Fluid Dynamics, pages 507–512. Springer–Verlag, 1982.
\par\hangindent=2em\hangafter=1 562. B. van Leer. On the Relation Between the Upwind–Differencing Schemes of Godunov, Enguist-Osher and Roe. SIAM J. Sci. Stat. Comput. , 5(1):1–20, 1985.
\par\hangindent=2em\hangafter=1 563. B. van Leer. Progress in Multi–Dimensional Upwind Differencing. Technical Report CR-189708/ICASE 92–43, NASA, Sept. 1992.
\par\hangindent=2em\hangafter=1 564. B. van Leer. Upwind and High–Resolution Methods for Compressible Flow: from Donor Cell to Residual Distrubution Schemes. Review Article.Commun. Comput. Phys., 1(2):192–206, 2006.
\par\hangindent=2em\hangafter=1 565. B. van Leer, J. L. Thomas, and P. L. Roe. A comparison of numerical flux formulas for the euler and navier–stokes equations. AIAA 8th Computational Fluid Dynamics Conference, June 9–11, 1987, Honolulu, Hawaii. Paper AIAA– 87–1104–CP, 1987.
\par\hangindent=2em\hangafter=1 566. M. H. P. M van Puten. A Numerical Implementation of MHD in Divergence Form. J. Comput. Phys., 105:339–353, 1993.
% [原PDF p731]
\par\hangindent=2em\hangafter=1 567. M. O. Varner, W. R. Martindale, W. J. Phares, and J. C. Adams. Large Per- turbation Flow Field Analysis and Simulation for Supersonic Inlets. Technical Report NASA CR–174676, NASA, USA, 1984.
\par\hangindent=2em\hangafter=1 568. M. E. V´azquez-Cend´on. Estudio de Esquemas Descentrados para su Aplicaci´on a las Leyes de Conservaci´ on Hiperb´olicas con T´erminos de Fuente.P h Dt h e s i s , Departamento de Matem´aticas Aplicadas, Universidad de Santiago de Com- postela, Espa˜na, 1994.
\par\hangindent=2em\hangafter=1 569. M. E. V´ azquez-Cend´on. Improved Treatment of Source Terms in Upwind Schemes for the Shallow Water Equations in Channels with Irregular Geome- try. J. Comput. Phys., 148:497–526, 1999.
\par\hangindent=2em\hangafter=1 570. M. Vinokur. An Analysis of Finite–Difference and Finite–Volume Formulations of Conservation Laws. J. Comput. Phys., 81:1–52, 1989.
\par\hangindent=2em\hangafter=1 571. A. I. Volpert. The Space BV and Quaslinear Equations. Math. USSR Sbornik, 73:225–267, 1967.
\par\hangindent=2em\hangafter=1 572. L. Wang and D. J. Mavriplis. Implicit Solution of the Unsteady Euler Equations for High–Order Accurate Discontinuous Galerkin Discretizations. J. Comput. Phys., 225(2):1994–2015, 2007.
\par\hangindent=2em\hangafter=1 573. S. Y. Wang, P. B. Buttler, and H. Krier. Non–Ideal Equations of State for Combusting and Detonating Explosives. Prog. Energy Combust. Sci., 11:311– 331, 1985.
\par\hangindent=2em\hangafter=1 574. R. F. Warming and R. W. Beam. Upwind Second Order Difference Schemes with Applications in Aerodynamic Flows. AIAA Journal, 24:1241–1249, 1976.
\par\hangindent=2em\hangafter=1 575. Z. U. A. Warsi. Fluid Dynamics: Theoretical and Computational Approaches. CRC Press Inc., 1992.
\par\hangindent=2em\hangafter=1 576. G. Watson, D. H. Peregrine, and E. F. Toro. Numerical Solution of the Shallow Water Equations on a Beach Using the Weighted Average Flux Method. In Computational Fluid Dynamics 1992, Vol. 1, pages 495–502. Elsevier, 1992.
\par\hangindent=2em\hangafter=1 577. N. P. Weatherill. Numerical Grid Generation. Lecture series 1990–06, von Karman Institute for Fluid Dynamics, 1990.
\par\hangindent=2em\hangafter=1 578. N. P. Weatherill. An Upwind Kinetic Flux Vector Splitting Method on General Mesh Topologies. Intern. J. Numer. Meth. Engineering, 37:623–643, 1994.
\par\hangindent=2em\hangafter=1 579. E. Weinan and C. W. Shu. A Numerical Resolution Study of High–Order Es- sentially Non–Oscillatory Schemes Applied to Incompressible Flow. Technical Report ICASE 92–39, NASA Langley Research Center, USA, 1992.
\par\hangindent=2em\hangafter=1 580. B. V. Wells, M. J. Baines, and P. Glaister. Generation of Arbitrary Lagrangian– Eulerian (ALE) Velocities, Based on Monitor Functions, for the Solution of Compressible Fluid Equations. Int. J. Numer. Methods in Fluids , 47:1375– 1381, 2005.
\par\hangindent=2em\hangafter=1 581. P. Wesseling. Principles of Computational Fluid Dynamics. Springer, 2001.
\par\hangindent=2em\hangafter=1 582. G. B. Whitam. Linear and Non–linear Waves. John Wiley and Sons, 1974.
\par\hangindent=2em\hangafter=1 583. F. A. Williams. Combustion Theory. Benjamin Cummings Publishing Com- pany, Inc Merlo Park, California, 1985.
\par\hangindent=2em\hangafter=1 584. P. Woodward and P. Colella. The Numerical Simulation of Two–Dimensional Fluid Flow with Strong Shocks. J. Comput. Phys., 54:115–173, 1984.
\par\hangindent=2em\hangafter=1 585. K. Xu. Gas Kinetic Schemes for Unsteady Compressible Flow Simulations. VKI Fluid Dynamics Lecture Series, 1998–3, 1998.
\par\hangindent=2em\hangafter=1 586. K. Xu, L. Martinelli, and A. Jameson. Gas–Kinetic Finite Volume Methods, Flux–Vector Splitting and Artificial Diffusion. J. Comput. Phys., 120:48–65, 1995.
% [原PDF p732]
\par\hangindent=2em\hangafter=1 587. K. Xu and K. H. Prendergast. Numerical Navier–Stokes Solutions from Gas– Kinetic Theory. J. Comput. Phys., 114:9–17, 1994.
\par\hangindent=2em\hangafter=1 588. N. N. Yanenko. The Method of Fractional Steps. Springer–Verlag, New York, 1971.
\par\hangindent=2em\hangafter=1 589. G. Yang, D. M. Causon, D. M. Ingram, R. Saunders, and P. Batten. A Carte- sian Cut Cell Method for Compressible Flows Part A: Static Body Problems. Aeronautical Journal, 101:47–56, 1997.
\par\hangindent=2em\hangafter=1 590. G. Yang, D. M. Causon, D. M. Ingram, R. Saunders, and P. Batten. A Carte- sian Cut Cell Method for Compressible Flows Part B: Moving Body Problems. Aeronautical Journal, 101:57–65, 1997.
\par\hangindent=2em\hangafter=1 591. J. Y. Yang, M. H. Chen, I. N. Tsai, and J. W. Chang. A Kinetic Beam Scheme for Relativistic Gas Dynamics. J. Comput. Phys., 136:19–40, 1997.
\par\hangindent=2em\hangafter=1 592. H. C. Yee. Construction of Explicit and Implicit Symmetric TVD Schemes and their Applications. J. Comput. Phys., 68:151–179, 1987.
\par\hangindent=2em\hangafter=1 593. H. C. Yee. A Class of High–Resolution Explicit and Implicit Shock–Capturing Methods. Von Karman Institute of Fluid Dynamics, Lecture Series 1989–04, 1989.
\par\hangindent=2em\hangafter=1 594. H. C. Yee and P. K. Sweby. Non–Linear Dynamics and Numerical Uncertainties in CFD. Technical Report TM 110398, NASA Ames Research Center, Moffett Field, USA, 1996.
\par\hangindent=2em\hangafter=1 595. H. C. Yee, R. F. Warming, and A. Harten. Implicit Total Variation Diminishing (TVD) Schemes for Steady State Problems. J. Comput. Phys. , 57:327–360, 1985.
\par\hangindent=2em\hangafter=1 596. E. C. Zachmanoglou and D. W. Thoe. Introduction to Partial Differential Equations. Dover Publications, Inc. New York, 1986.
\par\hangindent=2em\hangafter=1 597. S. T. Zalesak. Fully Multidimensional Flux–Corrected Transport Algorithms for Fluids. J. Comput. Phys., 31:335–362, 1979.
\par\hangindent=2em\hangafter=1 598. D. D. Zeeuw and K. G. Powell. An Adaptively Refined Cartesian Mesh Solver for the Euler Equations. J. Comput. Phys., 104:56–68, 1993.
\par\hangindent=2em\hangafter=1 599. Y. B. Zeldovich and Y. P. Raizer. Physics of Shock Waves and High Temper- ature Hydrodynamic Phenomena. Academic Press, 1966.
\par\hangindent=2em\hangafter=1 600. G-C. Zha and E. Bilgen. Numerical Solution of Euler Equations by a New Flux Vector Splitting Scheme. Int. J. Numer. Meth. Fluids, 17:115–144, 1993.
\par\hangindent=2em\hangafter=1 601. T. Zhang and Y. Zheng. Conjecture on the Structure of Solutions of the Riemann Problem for Two Dimensional Gas Dynamics.SIAM J. Math. Anal., 21(3):593–630, 1990.
\par\hangindent=2em\hangafter=1 602. H. W. Zheng, C. Shu, and Chew Y. T. An Object–Oriented and Quadrilateral– Mesh Based Solution Adaptive Algorithm for Compressible Multi–Fluid Flows. J. Comput. Phys., 227(14):6895–6921, 2008.
\endgroup
\chapter*{中英双语索引}
\addcontentsline{toc}{chapter}{中英双语索引}
\markboth{中英双语索引}{中英双语索引}
\begingroup\small\sloppy
\setlength{\parindent}{0pt}
\setlength{\parskip}{0.10em}
% [原PDF p733]
\par\hangindent=1.5em 123 问题（123 problem）\dotfill 129, 281
\par\hangindent=1.5em 自适应格式（adaptive scheme）\dotfill 522
\par\hangindent=1.5em ADER 方法（ADER method）\dotfill 655–659, 661, 663, 666, 667, 669, 671, 673, 674
\par\hangindent=1.5em 绝热指数（adiabatic exponent）\dotfill 12
\par\hangindent=1.5em 平流（advection）\dotfill 18, 66
\par\hangindent=1.5em 放大因子（amplification factor）\dotfill 540, 559
\par\hangindent=1.5em 振幅（amplitude）\dotfill 167
\par\hangindent=1.5em 自适应网格加密（AMR）（AMR）\dotfill 593
\par\hangindent=1.5em 近似黎曼求解器（approximate Riemann solver）\dotfill 112
\par\hangindent=1.5em 算术平均（arithmetic mean）\dotfill 251
\par\hangindent=1.5em 人工可压缩性方程（artificial compressibility equations）\dotfill 39, 40, 681
\par\hangindent=1.5em 人工可压缩性因子（artificial compressibility factor）\dotfill 40
\par\hangindent=1.5em 人工黏性（artificial viscosity）\dotfill 171, 487
\par\hangindent=1.5em 渐近稳定系统（asymptotically stable system）\dotfill 538
\par\hangindent=1.5em 阿伏伽德罗常数（Avogadro Number）\dotfill 11
\par\hangindent=1.5em 后向差分（backward difference）\dotfill 165
\par\hangindent=1.5em 后向 Euler 方法（backward Euler method）\dotfill 539
\par\hangindent=1.5em Boltzmann 方法（Boltzmann approach）\dotfill 265
\par\hangindent=1.5em Boltzmann 常数（Boltzmann Constant）\dotfill 11, 12
\par\hangindent=1.5em 边界条件（boundary conditions）\dotfill 42, 182, 222, 223
\par\hangindent=1.5em 边界黎曼问题（boundary Riemann problem）\dotfill 223
\par\hangindent=1.5em Buckley--Leverett 方程（Buckley–Leverett equation）\dotfill 64
\par\hangindent=1.5em 体积黏性（bulk viscosity）\dotfill 16
\par\hangindent=1.5em Burgers 方程（Burgers’s equation）\dotfill 30, 71
\par\hangindent=1.5em 量热状态方程（caloric equation of state）\dotfill 7
\par\hangindent=1.5em 量热完全气体（calorically ideal gas）\dotfill 7
\par\hangindent=1.5em 标准状态方程（canonical equations of state）\dotfill 8
\par\hangindent=1.5em 标准形（canonical form）\dotfill 52
\par\hangindent=1.5em Cartesian 网格（Cartesian mesh）\dotfill 599, 605, 609, 610, 618
\par\hangindent=1.5em Castro--Toro 的 $\mathrm{GRP}_K$ 求解器（Castro–Toro GRP K solver）\dotfill 652
\par\hangindent=1.5em Cauchy 问题（Cauchy problem）\dotfill 626, 627, 629, 633, 635
\par\hangindent=1.5em Cauchy--Kowalewski 过程（Cauchy–Kowalewski）\dotfill 628, 630–635, 637, 640, 643, 646, 647, 651, 652, 659, 660, 665, 667–669
\par\hangindent=1.5em Cauchy--Riemann 方程（Cauchy–Riemann equations）\dotfill 46
\par\hangindent=1.5em 单元平均值（cell average）\dotfill 176
\par\hangindent=1.5em 单元中心方法（cell centred methods）\dotfill 176, 555
\par\hangindent=1.5em 单元 Reynolds 数（cell Reynolds number）\dotfill 487
\par\hangindent=1.5em 单元（cells）\dotfill 176
\par\hangindent=1.5em 中心稀疏波（centred rarefaction）\dotfill 76
\par\hangindent=1.5em 中心格式（centred scheme）\dotfill 416, 597, 598, 601, 603, 604, 611, 612, 679–681
\par\hangindent=1.5em CFL 系数（CFL coefficient）\dotfill 183, 221, 495
\par\hangindent=1.5em CFL 条件（CFL condition）\dotfill 217, 218
\par\hangindent=1.5em CFL 数（CFL number）\dotfill 167, 168
\par\hangindent=1.5em 链式法则（chain rule）\dotfill 43
\par\hangindent=1.5em 特征曲线（characteristic curves）\dotfill 47
\par\hangindent=1.5em 特征方程（characteristic equation）\dotfill 92, 93, 299
\par\hangindent=1.5em 特征场（characteristic field）\dotfill 77, 95, 111
\par\hangindent=1.5em 特征形式（characteristic form）\dotfill 52
\par\hangindent=1.5em 特征限制（characteristic limiting）\dotfill 293, 508
\par\hangindent=1.5em 特征多项式（characteristic polynomial）\dotfill 44, 90, 538
\par\hangindent=1.5em 特征速度（characteristic speed）\dotfill 48, 52, 66, 77
\par\hangindent=1.5em 特征变量（characteristic variables）\dotfill 51, 54, 188, 189
\par\hangindent=1.5em 特征线（characteristics）\dotfill 533
\par\hangindent=1.5em 化学组分（chemical species）\dotfill 326
\par\hangindent=1.5em 卡方统计（chi–square statistics）\dotfill 251
\par\hangindent=1.5em 嵌套网格 CHIMERA（CHIMERA）\dotfill 593
\par\hangindent=1.5em CIR 格式（CIR scheme）\dotfill 168, 171, 172, 177
\par\hangindent=1.5em 经典黎曼问题（classical Riemann problem）\dotfill 625, 629, 636–639, 642–644, 646, 647, 651–653
% [原PDF p734]
\par\hangindent=1.5em 极值削平（clipping of extrema）\dotfill 171, 460, 670
\par\hangindent=1.5em 紧支撑（compact support）\dotfill 63
\par\hangindent=1.5em 完整黎曼求解器（complete Riemann solver）\dotfill 316, 596, 598, 680
\par\hangindent=1.5em 复合格式（composite scheme）\dotfill 599
\par\hangindent=1.5em compressive 区域（compressive region）\dotfill 70
\par\hangindent=1.5em 计算空间（computational space）\dotfill 576
\par\hangindent=1.5em 凹通量（concave flux）\dotfill 66
\par\hangindent=1.5em 守恒律（conservation laws）\dotfill 3, 43, 61, 66
\par\hangindent=1.5em 能量守恒（conservation of energy）\dotfill 4, 24, 25
\par\hangindent=1.5em 质量守恒（conservation of mass）\dotfill 2, 4, 20, 21
\par\hangindent=1.5em 动量守恒（conservation of momentum）\dotfill 4, 21, 22
\par\hangindent=1.5em 守恒格式（conservative scheme）\dotfill 2, 175
\par\hangindent=1.5em 守恒变量（conserved variables）\dotfill 2, 43, 61
\par\hangindent=1.5em 一致性条件（Consistency Condition）\dotfill 175, 350, 417
\par\hangindent=1.5em 一致性条件（consistency condition）\dotfill 319–324, 601
\par\hangindent=1.5em 接触间断（contact discontinuity）\dotfill 85, 96, 107, 118, 140
\par\hangindent=1.5em 接触波（contact wave）\dotfill 85
\par\hangindent=1.5em 对流通量分量（convective flux component）\dotfill 279
\par\hangindent=1.5em 收敛阶（convergence rates）\dotfill 672, 673
\par\hangindent=1.5em 凸通量（convex flux）\dotfill 66
\par\hangindent=1.5em Courant 数（Courant number）\dotfill 167
\par\hangindent=1.5em Courant 数系数（Courant number coefficient）\dotfill 183
\par\hangindent=1.5em 协体积状态方程（covolume equation of state）\dotfill 13, 223
\par\hangindent=1.5em 协体积气体（covolume gases）\dotfill 117, 143
\par\hangindent=1.5em 数据相容性（data compatibility）\dotfill 450
\par\hangindent=1.5em 解耦系统（decoupled system）\dotfill 52
\par\hangindent=1.5em 密度（density）\dotfill 2
\par\hangindent=1.5em 因变量（dependent variables）\dotfill 42
\par\hangindent=1.5em 爆轰（detonation）\dotfill 682, 683
\par\hangindent=1.5em 爆轰类比（detonation analogue）\dotfill 78
\par\hangindent=1.5em 爆轰波（detonation waves）\dotfill 78
\par\hangindent=1.5em 对角矩阵（diagonal matrix）\dotfill 106
\par\hangindent=1.5em 可对角化矩阵（diagonalisable matrix）\dotfill 106
\par\hangindent=1.5em 可对角化系统（diagonalisable system）\dotfill 51
\par\hangindent=1.5em 扩散（diffusion）\dotfill 18
\par\hangindent=1.5em 扩散数（diffusion number）\dotfill 487
\par\hangindent=1.5em 维数分裂（dimensional splitting）\dotfill 543
\par\hangindent=1.5em 间断 Galerkin 有限元（discontinuous Galerkin finite elements）\dotfill 673
\par\hangindent=1.5em 间断 Galerkin 方法（discontinuous Galerkin method）\dotfill 655, 657, 673, 674, 680, 681
\par\hangindent=1.5em 色散方程（dispersive equation）\dotfill 174
\par\hangindent=1.5em 离解（dissociation）\dotfill 16
\par\hangindent=1.5em 分布理论（Distribution Theory）\dotfill 452
\par\hangindent=1.5em 散度算子（divergence operator）\dotfill 2
\par\hangindent=1.5em 依赖域（domain of dependence）\dotfill 60
\par\hangindent=1.5em 决定域（domain of determinacy）\dotfill 60
\par\hangindent=1.5em 点积（dot product）\dotfill 2
\par\hangindent=1.5em Dumbser--Enaux--Toro 的 $\mathrm{GRP}_K$ 求解器（Dumbser–Enaux–Toro GRP K solver）\dotfill 652
\par\hangindent=1.5em 特征值（eigenvalues）\dotfill 44, 65, 92, 104, 110, 298, 538
\par\hangindent=1.5em 基本波（elementary waves）\dotfill 94
\par\hangindent=1.5em 椭圆系统（elliptic system）\dotfill 45, 46
\par\hangindent=1.5em 本质无振荡方法（ENO）（ENO）\dotfill 439, 661, 662, 666, 669, 680
\par\hangindent=1.5em 焓（enthalpy）\dotfill 8, 89
\par\hangindent=1.5em 熵（entropy）\dotfill 5, 7, 93
\par\hangindent=1.5em 熵条件（entropy condition）\dotfill 72, 96
\par\hangindent=1.5em 熵形式（entropy formulation）\dotfill 94
\par\hangindent=1.5em 熵异常（entropy glitch）\dotfill 227
\par\hangindent=1.5em 违反熵条件的激波（entropy–violating shock）\dotfill 73
\par\hangindent=1.5em 状态方程（equation of state）\dotfill 67, 140
\par\hangindent=1.5em 逃逸速度（escape velocity）\dotfill 140
\par\hangindent=1.5em Euler 方法（Euler method）\dotfill 539
\par\hangindent=1.5em expansive 区域（expansive region）\dotfill 70
\par\hangindent=1.5em 显式格式（explicit scheme）\dotfill 167, 539
\par\hangindent=1.5em 爆炸（explosion）\dotfill 586, 587, 590
\par\hangindent=1.5em 远场边界条件（far–field boundary conditions）\dotfill 224
\par\hangindent=1.5em 虚单元（fictitious cell）\dotfill 182, 222, 495
\par\hangindent=1.5em 有限体积（finite volume）\dotfill 19, 176
\par\hangindent=1.5em 有限体积方法（finite volume method）\dotfill 543, 655, 657–660, 665, 666, 668, 669, 673, 674, 679–681
\par\hangindent=1.5em 流动参数（flow parameter）\dotfill 501, 552
\par\hangindent=1.5em 通量差分分裂（Flux Difference Splitting）\dotfill 265
\par\hangindent=1.5em 通量函数（flux function）\dotfill 66, 175
\par\hangindent=1.5em 通量矢量（flux vector）\dotfill 43, 61
\par\hangindent=1.5em 通量矢量分裂（Flux Vector Splitting）\dotfill 90, 265
\par\hangindent=1.5em FORCE 通量（FORCE flux）\dotfill 598–601, 603, 605, 606, 608–610, 612, 613, 617, 620
\par\hangindent=1.5em 强迫项（forcing term）\dotfill 486
\par\hangindent=1.5em 前向差分（forward difference）\dotfill 164, 167
\par\hangindent=1.5em 分步法（fractional step method）\dotfill 535, 543, 544
\par\hangindent=1.5em 自由能（free energy）\dotfill 9
\par\hangindent=1.5em 自由表面（free surface）\dotfill 33, 35
\par\hangindent=1.5em 自由来流 Mach 数（free-stream Mach number）\dotfill 46
\par\hangindent=1.5em Fromm 格式（Fromm scheme）\dotfill 174
\par\hangindent=1.5em 界面追踪格式（front tracking schemes）\dotfill 294
\par\hangindent=1.5em 气体常数（Gas Constant）\dotfill 11
\par\hangindent=1.5em 一般状态方程（general equation of state）\dotfill 682
\par\hangindent=1.5em 广义黎曼不变量（Generalised Riemann Invariant）\dotfill 81, 97, 98, 111, 122–124, 135, 141, 377,
% [原PDF p735]
\par\hangindent=1.5em 广义黎曼不变量（续）（Generalised Riemann Invariant, continued）\dotfill 385–387, 389, 390, 393, 395, 398, 399, 404
\par\hangindent=1.5em 广义黎曼问题（generalized Riemann problem）\dotfill 625, 627–630, 635, 640–643, 648, 649, 651–653, 655, 656, 659, 660, 663, 665, 667, 669, 673
\par\hangindent=1.5em 广义解（generalized solution）\dotfill 63
\par\hangindent=1.5em 真正非线性场（genuinely non–linear field）\dotfill 77, 84, 85, 95, 107, 111, 382, 395
\par\hangindent=1.5em 位势（geopotential）\dotfill 37
\par\hangindent=1.5em GFORCE 通量（GFORCE flux）\dotfill 602–604
\par\hangindent=1.5em Godunov 中心通量（Godunov centred flux）\dotfill 609, 611
\par\hangindent=1.5em Godunov 中心格式（Godunov centred scheme）\dotfill 443–445
\par\hangindent=1.5em Godunov 通量（Godunov flux）\dotfill 180, 181, 191, 219
\par\hangindent=1.5em Godunov 格式（Godunov scheme）\dotfill 216, 218, 597, 600, 603, 604, 612
\par\hangindent=1.5em 梯度算子（gradient operator）\dotfill 2
\par\hangindent=1.5em 网格生成（grid generation）\dotfill 574
\par\hangindent=1.5em 网格速度（grid speed）\dotfill 167, 171
\par\hangindent=1.5em Harten--Osher--Engquist--Chakravarthy 的 $\mathrm{GRP}_K$ 求解器（Harten–Osher–Engquist–Chakravarthy GRP K solver）\dotfill 651
\par\hangindent=1.5em 稀疏波头（head of rarefaction）\dotfill 75
\par\hangindent=1.5em 热容（heat capacity）\dotfill 9
\par\hangindent=1.5em 热传导（heat conduction）\dotfill 18, 66
\par\hangindent=1.5em 反应热（heats of reaction）\dotfill 78
\par\hangindent=1.5em 高分辨率（high resolution）\dotfill 445
\par\hangindent=1.5em 高阶黎曼问题（high–order Riemann problem）\dotfill 625
\par\hangindent=1.5em 高分辨率方法（high–resolution methods）\dotfill 449
\par\hangindent=1.5em HLL 黎曼求解器（HLL Riemann solver）\dotfill 315, 320, 321, 329, 333
\par\hangindent=1.5em HLLC 通量（HLLC flux）\dotfill 322–327, 331–333, 335, 336
\par\hangindent=1.5em HLLC 黎曼求解器（HLLC Riemann solver）\dotfill 315, 593
\par\hangindent=1.5em HLLE 黎曼求解器（HLLE Riemann solver）\dotfill 315, 328, 329
\par\hangindent=1.5em HLLEM 黎曼求解器（HLLEM Riemann solver）\dotfill 329
\par\hangindent=1.5em 均熵流（homentropic flow）\dotfill 94
\par\hangindent=1.5em 齐次性质（homogeneity property）\dotfill 89, 269, 270
\par\hangindent=1.5em 齐次方程（homogeneous equation）\dotfill 533
\par\hangindent=1.5em 齐次系统（homogeneous system）\dotfill 18, 26, 42, 531
\par\hangindent=1.5em 双曲守恒律（hyperbolic conservation laws）\dotfill 2
\par\hangindent=1.5em 双曲型偏微分方程（hyperbolic PDE）\dotfill 55
\par\hangindent=1.5em 双曲系统（hyperbolic system）\dotfill 45, 61, 91
\par\hangindent=1.5em 时间方向双曲性（hyperbolicity in time）\dotfill 105, 106
\par\hangindent=1.5em 理想气体（ideal gases）\dotfill 117
\par\hangindent=1.5em implicit 方法（implicit method）\dotfill 539
\par\hangindent=1.5em 内爆（implosion）\dotfill 586, 587, 589
\par\hangindent=1.5em 初生空化算例（incipient cavitation case）\dotfill 403
\par\hangindent=1.5em 不完整黎曼求解器（incomplete Riemann solver）\dotfill 316, 596, 598, 680
\par\hangindent=1.5em 格式增量形式（incremental form of a scheme）\dotfill 450
\par\hangindent=1.5em 自变量（independent variables）\dotfill 42
\par\hangindent=1.5em 无限分辨率（infinite resolution）\dotfill 335
\par\hangindent=1.5em 非齐次方程（inhomogeneous equation）\dotfill 535, 536
\par\hangindent=1.5em 非齐次问题（inhomogeneous problem）\dotfill 486
\par\hangindent=1.5em 非齐次系统（inhomogeneous system）\dotfill 26, 78
\par\hangindent=1.5em 初始条件（initial condition）\dotfill 42
\par\hangindent=1.5em 初值问题（initial–value problem）\dotfill 47
\par\hangindent=1.5em 不稳定性（instability）\dotfill 74, 221
\par\hangindent=1.5em 积分曲线（integral curves）\dotfill 377, 381
\par\hangindent=1.5em 积分形式（integral form）\dotfill 62
\par\hangindent=1.5em 积分路径（integration path）\dotfill 381, 384, 385, 389
\par\hangindent=1.5em 单元界面（intercell boundary）\dotfill 177
\par\hangindent=1.5em 单元间通量（intercell flux）\dotfill 180, 184
\par\hangindent=1.5em 内能（internal energy）\dotfill 3, 8, 88
\par\hangindent=1.5em 交点（intersection point）\dotfill 377, 384, 385, 389, 397, 403
\par\hangindent=1.5em Jacobian 矩阵的逆（inverse of Jacobian matrix）\dotfill 271
\par\hangindent=1.5em inviscid Burgers’s 方程（inviscid Burgers’s equation）\dotfill 42, 64, 181
\par\hangindent=1.5em 电离（ionisation）\dotfill 16
\par\hangindent=1.5em 等熵方程（isentropic equations）\dotfill 30
\par\hangindent=1.5em 等熵流（isentropic flow）\dotfill 94
\par\hangindent=1.5em 等熵气体动力学（isentropic gas dynamics）\dotfill 65
\par\hangindent=1.5em 等熵律（isentropic law）\dotfill 94
\par\hangindent=1.5em 等温压缩系数（isothermal compressibility）\dotfill 9, 11
\par\hangindent=1.5em 等温方程（isothermal equations）\dotfill 30, 64
\par\hangindent=1.5em 各向同性介质（isotropic medium）\dotfill 15
\par\hangindent=1.5em Jacobian 矩阵（Jacobian matrix）\dotfill 43, 61, 64, 88, 104, 106, 108, 270, 271, 275, 354, 392, 538
\par\hangindent=1.5em 变换 Jacobian（Jacobian of the transformation）\dotfill 577
\par\hangindent=1.5em 跳跃间断（jump discontinuity）\dotfill 55
\par\hangindent=1.5em 运动黏度（kinematic viscosity）\dotfill 38
\par\hangindent=1.5em 动能（kinetic energy）\dotfill 3
\par\hangindent=1.5em 大时间步格式（large–time step scheme）\dotfill 424
\par\hangindent=1.5em Lax 等价定理（Lax Equivalence Theorem）\dotfill 418
\par\hangindent=1.5em Lax--Friedrichs 通量（Lax–Friedrichs flux）\dotfill 185, 329, 599, 601–605, 609, 610
\par\hangindent=1.5em Lax--Friedrichs 格式（Lax–Friedrichs scheme）\dotfill 172, 184, 226, 244, 249
\par\hangindent=1.5em Lax--Wendroff 通量（Lax–Wendroff flux）\dotfill 186, 599, 601–605, 609–611, 620
\par\hangindent=1.5em Lax--Wendroff 格式（Lax–Wendroff scheme）\dotfill 173
\par\hangindent=1.5em 左特征向量（left eigenvector）\dotfill 44, 92
% [原PDF p736]
\par\hangindent=1.5em 线性平流方程（linear advection equation）\dotfill 31, 33, 42, 47, 244
\par\hangindent=1.5em 线性格式（linear schemes）\dotfill 415
\par\hangindent=1.5em 线性系统（linear system）\dotfill 31, 42
\par\hangindent=1.5em 线性退化场（linearly degenerate field）\dotfill 77, 83, 95, 107, 111, 382, 395
\par\hangindent=1.5em 长时间演化（long time evolution）\dotfill 669, 671
\par\hangindent=1.5em 低密度（low density）\dotfill 227
\par\hangindent=1.5em Mach 数（Mach number）\dotfill 32, 101, 266, 277, 591–593
\par\hangindent=1.5em Mach 反射（Mach reflection）\dotfill 592, 593
\par\hangindent=1.5em Mach 杆（Mach stem）\dotfill 593
\par\hangindent=1.5em 磁流体力学（magnetohydrodynamics）\dotfill 683
\par\hangindent=1.5em 质量通量（mass flux）\dotfill 121, 124, 145, 146
\par\hangindent=1.5em 物质导数（material derivative）\dotfill 19
\par\hangindent=1.5em 平均自由程（mean–free path）\dotfill 70
\par\hangindent=1.5em 网格函数（mesh function）\dotfill 452
\par\hangindent=1.5em 网格生成（mesh generation）\dotfill 574
\par\hangindent=1.5em 线方法（method of lines）\dotfill 439
\par\hangindent=1.5em 度量系数（metrics）\dotfill 576
\par\hangindent=1.5em 磁流体力学方程（MHD）（MHD equations）\dotfill 683
\par\hangindent=1.5em 修正方程（modified equation）\dotfill 171, 172, 249
\par\hangindent=1.5em 摩尔（mole）\dotfill 11
\par\hangindent=1.5em 动量（momentum）\dotfill 2
\par\hangindent=1.5em 单原子气体（monatomic gases）\dotfill 13
\par\hangindent=1.5em 单调通量（monotone flux）\dotfill 680
\par\hangindent=1.5em 单调格式（monotone scheme）\dotfill 172, 173, 226, 250, 521, 568
\par\hangindent=1.5em 移动边界（moving boundary）\dotfill 223
\par\hangindent=1.5em 多组分流（multicomponent flow）\dotfill 326
\par\hangindent=1.5em 多维流动（multidimensional flow）\dotfill 326
\par\hangindent=1.5em 多维迎风（multidimensional upwinding）\dotfill 680
\par\hangindent=1.5em 多相流（multiphase flow）\dotfill 683–685
\par\hangindent=1.5em MUSTA 方法（MUSTA approach）\dotfill 604
\par\hangindent=1.5em Navier--Stokes 方程（Navier–Stokes equations）\dotfill 17, 680–682
\par\hangindent=1.5em Nessyahu--Tadmor 格式（Nessyahu–Tadmor scheme）\dotfill 598, 599
\par\hangindent=1.5em Newton--Raphson 迭代（Newton–Raphson）\dotfill 127
\par\hangindent=1.5em Newton 流体（Newtonian fluid）\dotfill 15
\par\hangindent=1.5em 非守恒方法（non–conservative method）\dotfill 174
\par\hangindent=1.5em 非守恒乘积（non–conservative products）\dotfill 685
\par\hangindent=1.5em 非守恒系统（non–conservative system）\dotfill 674
\par\hangindent=1.5em 非凸通量（non–convex flux）\dotfill 67
\par\hangindent=1.5em 非线性格式（non–linear scheme）\dotfill 447
\par\hangindent=1.5em 非线性系统（non–linear system）\dotfill 42
\par\hangindent=1.5em 无反射边界条件（non–reflecting boundary conditions）\dotfill 224
\par\hangindent=1.5em 法向速度（normal velocity）\dotfill 107, 550
\par\hangindent=1.5em NUMERICA 程序库（NUMERICA）\dotfill 334, 679, 686
\par\hangindent=1.5em 数值耗散（numerical dissipation）\dotfill 333
\par\hangindent=1.5em 数值依赖域（numerical domain of dependence）\dotfill 169, 171
\par\hangindent=1.5em 数值通量（numerical flux）\dotfill 175, 177, 180, 626, 630, 636, 640, 642, 645, 646, 648, 651–653, 655, 658, 663, 665–669, 679–681
\par\hangindent=1.5em 数值源项（numerical source）\dotfill 626, 628, 630, 633, 640, 642, 653, 655, 658–660, 663, 665, 666, 668, 669
\par\hangindent=1.5em 数值黏性（numerical viscosity）\dotfill 171, 172, 460, 602–604
\par\hangindent=1.5em 单边导数（one–sided derivative）\dotfill 164
\par\hangindent=1.5em 单边差分（one–sided differencing）\dotfill 168, 188, 190
\par\hangindent=1.5em 开口端边界条件（open–end boundary conditions）\dotfill 224
\par\hangindent=1.5em 常微分方程（ordinary differential equation）\dotfill 47
\par\hangindent=1.5em 无振荡格式（oscillation–free scheme）\dotfill 563
\par\hangindent=1.5em 参数矢量（parameter vector）\dotfill 352, 355
\par\hangindent=1.5em 质点速度（particle velocity）\dotfill 101, 102
\par\hangindent=1.5em 被动标量（passive scalar）\dotfill 296, 317, 327, 333
\par\hangindent=1.5em 路径守恒格式（path–conservative scheme）\dotfill 684
\par\hangindent=1.5em 相角（phase angle）\dotfill 167, 559
\par\hangindent=1.5em 相平面（phase plane）\dotfill 77
\par\hangindent=1.5em 相空间（phase space）\dotfill 77, 381
\par\hangindent=1.5em 物理通量（physical flux）\dotfill 175
\par\hangindent=1.5em 物理空间（physical space）\dotfill 576
\par\hangindent=1.5em 物理变量（physical variables）\dotfill 2
\par\hangindent=1.5em Poisson 方程（Poisson equation）\dotfill 39
\par\hangindent=1.5em 多原子气体（polyatomic gases）\dotfill 13
\par\hangindent=1.5em 多方气体（polytropic gas）\dotfill 12
\par\hangindent=1.5em 正性保持守恒求解器（positively conservative solver）\dotfill 329
\par\hangindent=1.5em 分片抛物线方法（PPM）（PPM method）\dotfill 680
\par\hangindent=1.5em Prandtl 数（Prandtl number）\dotfill 17
\par\hangindent=1.5em 压力（pressure）\dotfill 2
\par\hangindent=1.5em 压力正性条件（pressure positivity condition）\dotfill 127, 143, 403
\par\hangindent=1.5em 原始变量（primitive variables）\dotfill 2, 91, 109
\par\hangindent=1.5em 原始变量形式（primitive–variable formulation）\dotfill 114
\par\hangindent=1.5em 拟线性系统（quasi–linear system）\dotfill 42
\par\hangindent=1.5em 辐射边界条件（radiation boundary conditions）\dotfill 224
\par\hangindent=1.5em 随机采样（random sampling）\dotfill 240, 241
\par\hangindent=1.5em 影响域（range of influence）\dotfill 61
\par\hangindent=1.5em Rankine--Hugoniot 条件（Rankine–Hugoniot condition）\dotfill 71, 79, 80, 99, 112, 123, 139, 321, 323, 593
\par\hangindent=1.5em 稀疏激波（rarefaction shock）\dotfill 73, 75, 366
% [原PDF p737]
\par\hangindent=1.5em 稀疏波（rarefaction wave）\dotfill 74, 76, 85, 96, 98, 107, 111, 118
\par\hangindent=1.5em 比热比（ratio of specific heats）\dotfill 12, 13, 88
\par\hangindent=1.5em 反应进度变量（reaction progress variable）\dotfill 78
\par\hangindent=1.5em 反应--扩散方程（reaction–diffusion equation）\dotfill 673
\par\hangindent=1.5em 真实气体（real gases）\dotfill 14
\par\hangindent=1.5em 重构（reconstruction）\dotfill 626, 652, 656, 659–663, 666–669
\par\hangindent=1.5em 规则反射（regular reflection）\dotfill 592
\par\hangindent=1.5em 相对速度（relative velocity）\dotfill 99, 102, 120
\par\hangindent=1.5em 相对论气体动力学（relativistic gas dynamics）\dotfill 683
\par\hangindent=1.5em Richtmyer 格式（Richtmyer scheme）\dotfill 186, 226, 249
\par\hangindent=1.5em 黎曼方法（Riemann approach）\dotfill 265
\par\hangindent=1.5em 黎曼问题（Riemann problem）\dotfill 49, 55, 79, 83, 94, 116, 140, 177, 315, 317–320, 322, 327–329
\par\hangindent=1.5em 右特征向量（right eigenvector）\dotfill 65
\par\hangindent=1.5em 右特征向量组（right eigenvectors）\dotfill 44, 92, 94, 104, 110, 298
\par\hangindent=1.5em 鲁棒性（robustness）\dotfill 226, 227, 282, 307, 334, 372, 405
\par\hangindent=1.5em Roe 矩阵（Roe matrix）\dotfill 351, 353
\par\hangindent=1.5em 旋转矩阵（rotation matrix）\dotfill 105, 109
\par\hangindent=1.5em 旋转不变性（rotational invariance）\dotfill 105, 108, 109
\par\hangindent=1.5em Runge--Kutta 方法（Runge–Kutta method）\dotfill 539
\par\hangindent=1.5em Rusanov 通量（Rusanov flux）\dotfill 329
\par\hangindent=1.5em 缩放因子（scaling factor）\dotfill 45, 92
\par\hangindent=1.5em 浅水方程（shallow water equations）\dotfill 35, 681
\par\hangindent=1.5em 剪切黏性（shear viscosity）\dotfill 16, 38
\par\hangindent=1.5em 剪切波（shear wave）\dotfill 107, 111
\par\hangindent=1.5em 激波拟合（shock fitting）\dotfill 174
\par\hangindent=1.5em 激波 Mach 数（shock Mach number）\dotfill 101, 103
\par\hangindent=1.5em 激波厚度（shock thickness）\dotfill 71
\par\hangindent=1.5em 激波（shock wave）\dotfill 66, 70, 76, 84, 96, 98, 101, 107, 111, 112, 118, 174
\par\hangindent=1.5em 激波管问题（shock–tube problem）\dotfill 95, 117
\par\hangindent=1.5em 信号速度（signal velocities）\dotfill 318
\par\hangindent=1.5em 自相似解（similarity solution）\dotfill 83, 95, 111, 215, 629, 636, 639, 641, 643, 644, 646
\par\hangindent=1.5em 慢速移动激波（slowly moving shock）\dotfill 228, 405
\par\hangindent=1.5em 小扰动方程（small perturbation equations）\dotfill 32, 46
\par\hangindent=1.5em 涂抹（smearing）\dotfill 171
\par\hangindent=1.5em Sod 测试问题（Sod’s test problem）\dotfill 129
\par\hangindent=1.5em 声爆（sonic boom）\dotfill 70
\par\hangindent=1.5em 声速流（sonic flow）\dotfill 32, 296, 333, 366, 377
\par\hangindent=1.5em 声速点（sonic point）\dotfill 182, 281, 377, 382, 384, 385, 389, 395, 397
\par\hangindent=1.5em 声速稀疏波（sonic rarefaction）\dotfill 118, 182, 296, 335, 366, 368, 369, 500
\par\hangindent=1.5em 声速（sound speed）\dotfill 10, 32, 44, 64, 65, 82, 83, 88, 93, 101, 266, 270
\par\hangindent=1.5em 源项（source term）\dotfill 23, 25, 47, 78, 486, 531, 684, 685
\par\hangindent=1.5em 组分方程（species equation）\dotfill 336
\par\hangindent=1.5em 比内能（specific internal energy）\dotfill 5, 13
\par\hangindent=1.5em 分裂黎曼问题（split Riemann problem）\dotfill 111, 112
\par\hangindent=1.5em 分裂系数矩阵格式（Split–Coefficient Matrix Scheme）\dotfill 267
\par\hangindent=1.5em 伪振荡（spurious oscillations）\dotfill 173, 226, 440, 445, 447, 501, 573
\par\hangindent=1.5em 稳定性（stability）\dotfill 167, 171
\par\hangindent=1.5em 稳定性条件（stability condition）\dotfill 168
\par\hangindent=1.5em 常微分方程稳定性（stability of ODEs）\dotfill 538
\par\hangindent=1.5em 标准差（standard deviation）\dotfill 251
\par\hangindent=1.5em 星形区（Star Region）\dotfill 56, 58, 59, 80, 118, 129, 139, 144, 390
\par\hangindent=1.5em 启动误差（start up errors）\dotfill 593
\par\hangindent=1.5em 静止接触间断（stationary contact）\dotfill 282
\par\hangindent=1.5em 稳态（steady state）\dotfill 40
\par\hangindent=1.5em 定常超声速流（steady supersonic flow）\dotfill 151
\par\hangindent=1.5em 格式模板（stencil of a scheme）\dotfill 169
\par\hangindent=1.5em 刚性常微分方程（stiff ODEs）\dotfill 538
\par\hangindent=1.5em 刚性源项（stiff source term）\dotfill 685
\par\hangindent=1.5em 刚性比（stiffness ratio）\dotfill 538
\par\hangindent=1.5em 流函数（stream function）\dotfill 38
\par\hangindent=1.5em 严格双曲性（strict hyperbolicity）\dotfill 45
\par\hangindent=1.5em 结构网格（structured mesh）\dotfill 597, 599, 605, 609
\par\hangindent=1.5em 亚声速流（subsonic flow）\dotfill 32, 46
\par\hangindent=1.5em 随体导数（substantial derivative）\dotfill 19
\par\hangindent=1.5em 超声速流（supersonic flow）\dotfill 32, 46
\par\hangindent=1.5em 格式支撑域（support of a scheme）\dotfill 416
\par\hangindent=1.5em Sutherland 公式（Sutherland formula）\dotfill 16
\par\hangindent=1.5em 稀疏波尾（tail of rarefaction）\dotfill 75
\par\hangindent=1.5em 切向速度（tangential velocity）\dotfill 107, 111, 550
\par\hangindent=1.5em 张量积（tensor product）\dotfill 4
\par\hangindent=1.5em 张量（tensors）\dotfill 4
\par\hangindent=1.5em 测试函数（test functions）\dotfill 63
\par\hangindent=1.5em 热导率（thermal conductivity）\dotfill 17
\par\hangindent=1.5em 热状态方程（thermal equation of state）\dotfill 6
\par\hangindent=1.5em 时间层（time level）\dotfill 166, 177
\par\hangindent=1.5em 时间步（time step）\dotfill 182, 183, 221
\par\hangindent=1.5em Toro--Titarev 的 $\mathrm{GRP}_K$ 求解器（Toro–Titarev GRP K solver）\dotfill 651
\par\hangindent=1.5em 全导数（total derivative）\dotfill 67
\par\hangindent=1.5em 总能量（total energy）\dotfill 2, 3
\par\hangindent=1.5em 交通流方程（traffic flow equation）\dotfill 64, 66
\par\hangindent=1.5em 跨声速稀疏波（transonic rarefaction）\dotfill 182
% [原PDF p738]
\par\hangindent=1.5em 跨声速激波（transonic shock）\dotfill 384
\par\hangindent=1.5em 透明边界条件（transparent boundary conditions）\dotfill 224
\par\hangindent=1.5em 梯形法（Trapezoidal Method）\dotfill 539, 540
\par\hangindent=1.5em 试探解（trial solution）\dotfill 167
\par\hangindent=1.5em 真实依赖域（true domain of dependence）\dotfill 171
\par\hangindent=1.5em 截断误差（truncation error）\dotfill 417
\par\hangindent=1.5em TVB 格式（TVB schemes）\dotfill 456
\par\hangindent=1.5em 总变差递减（TVD）（TVD）\dotfill 413–415, 425, 431, 438, 440, 451–458, 460, 462–464, 469–484, 486–488, 655, 661, 669–671, 674, 680, 682
\par\hangindent=1.5em 通用气体常数（Universal Gas Constant）\dotfill 11
\par\hangindent=1.5em UNO 格式（UNO schemes）\dotfill 439
\par\hangindent=1.5em 非结构网格（unstructured mesh）\dotfill 597, 599, 605, 617, 620, 674
\par\hangindent=1.5em 迎风差分（upwind differencing）\dotfill 168
\par\hangindent=1.5em 迎风通量（upwind flux）\dotfill 679–681, 685
\par\hangindent=1.5em 迎风格式（upwind scheme）\dotfill 165, 175, 597, 598, 600, 603, 604, 612
\par\hangindent=1.5em 源项迎风离散（upwinding the source terms）\dotfill 532
\par\hangindent=1.5em 真空（vacuum）\dotfill 116, 118, 127, 139, 142, 304
\par\hangindent=1.5em 真空条件（vacuum condition）\dotfill 143
\par\hangindent=1.5em 真空问题（vacuum problem）\dotfill 345
\par\hangindent=1.5em van der Corput 数列（van der Corput numbers）\dotfill 250
\par\hangindent=1.5em 矢势（vector potential）\dotfill 38
\par\hangindent=1.5em 速度（velocity）\dotfill 2
\par\hangindent=1.5em 黏性（viscosity）\dotfill 33
\par\hangindent=1.5em 黏性系数（viscosity coefficient）\dotfill 66
\par\hangindent=1.5em 黏性耗散（viscous dissipation）\dotfill 66
\par\hangindent=1.5em 黏性通量限制器（viscous flux limiters）\dotfill 487
\par\hangindent=1.5em 体膨胀系数（volume expansivity）\dotfill 9, 11
\par\hangindent=1.5em 涡量输运方程（vorticity transport equation）\dotfill 39
\par\hangindent=1.5em 涡量矢量（vorticity vector）\dotfill 38
\par\hangindent=1.5em WAF 方法（WAF method）\dotfill 587, 593, 609, 669
\par\hangindent=1.5em Warming--Beam 方法（Warming–Beam method）\dotfill 173
\par\hangindent=1.5em 波（wave）\dotfill 49
\par\hangindent=1.5em 波破碎（wave breaking）\dotfill 70
\par\hangindent=1.5em 波动方程（wave equation）\dotfill 31
\par\hangindent=1.5em 波跳跃（wave jump）\dotfill 58, 191
\par\hangindent=1.5em 波长（wave length）\dotfill 167
\par\hangindent=1.5em 波数（wave number）\dotfill 167
\par\hangindent=1.5em 波速估计（wave speed estimate）\dotfill 326–331, 335
\par\hangindent=1.5em 波陡化（wave steepening）\dotfill 66
\par\hangindent=1.5em 波强（wave strength）\dotfill 58
\par\hangindent=1.5em 逐波限制（wave–by–wave limiting）\dotfill 508
\par\hangindent=1.5em 固体中的波（waves in solids）\dotfill 682
\par\hangindent=1.5em 弱激波（weak shock）\dotfill 70
\par\hangindent=1.5em 弱解（weak solution）\dotfill 63
\par\hangindent=1.5em 良平衡格式（well–balanced scheme）\dotfill 685
\par\hangindent=1.5em 加权本质无振荡方法（WENO）（WENO）\dotfill 649, 661, 662, 669, 680
\par\hangindent=1.5em $x$ 向扫描（x sweep）\dotfill 546, 548
\par\hangindent=1.5em $y$ 向扫描（y sweep）\dotfill 546, 548
\par\hangindent=1.5em $z$ 向扫描（z sweep）\dotfill 548
\endgroup
\end{document}
